\input{preamble}

% OK, start here.
%
\begin{document}

\title{Cohomology of Sheaves}


\maketitle

\phantomsection
\label{section-phantom}

\tableofcontents

\section{Introduction}
\label{section-introduction}

\noindent
In this document we work out some topics on cohomology of sheaves
on topological spaces. We mostly work in the generality of modules
over a sheaf of rings and we work with morphisms of ringed spaces.
To see what happens for sheaves on sites take a look at the chapter
Cohomology on Sites, Section \ref{sites-cohomology-section-introduction}.
Basic references are \cite{Godement} and \cite{Iversen}.





\section{Cohomology of sheaves}
\label{section-cohomology-sheaves}

\noindent
Let $X$ be a topological space. Let $\mathcal{F}$ be an abelian sheaf.
We know that the category of abelian sheaves on $X$ has enough injectives, see
Injectives, Lemma \ref{injectives-lemma-abelian-sheaves-space}.
Hence we can choose an injective resolution
$\mathcal{F}[0] \to \mathcal{I}^\bullet$. As is customary we define
\begin{equation}
\label{equation-cohomology}
H^i(X, \mathcal{F}) = H^i(\Gamma(X, \mathcal{I}^\bullet))
\end{equation}
to be the {\it $i$th cohomology group of the abelian sheaf $\mathcal{F}$}.
The family of functors $H^i(X, -)$ forms a universal $\delta$-functor
from $\textit{Ab}(X) \to \textit{Ab}$.

\medskip\noindent
Let $f : X \to Y$ be a continuous map of topological spaces. With
$\mathcal{F}[0] \to \mathcal{I}^\bullet$ as above
we define
\begin{equation}
\label{equation-higher-direct-image}
R^if_*\mathcal{F} = H^i(f_*\mathcal{I}^\bullet)
\end{equation}
to be the {\it $i$th higher direct image of $\mathcal{F}$}.
The family of functors $R^if_*$ forms a universal $\delta$-functor
from $\textit{Ab}(X) \to \textit{Ab}(Y)$.

\medskip\noindent
Let $(X, \mathcal{O}_X)$ be a ringed space. Let $\mathcal{F}$ be an
$\mathcal{O}_X$-module. We know that the category of $\mathcal{O}_X$-modules
on $X$ has enough injectives, see
Injectives, Lemma \ref{injectives-lemma-sheaves-modules-space}.
Hence we can choose an injective resolution
$\mathcal{F}[0] \to \mathcal{I}^\bullet$. As is customary we define
\begin{equation}
\label{equation-cohomology-modules}
H^i(X, \mathcal{F}) = H^i(\Gamma(X, \mathcal{I}^\bullet))
\end{equation}
to be the {\it $i$th cohomology group of $\mathcal{F}$}.
The family of functors $H^i(X, -)$ forms a universal $\delta$-functor
from $\textit{Mod}(\mathcal{O}_X) \to \text{Mod}_{\mathcal{O}_X(X)}$.
For the classical algebraic-variety formulation of this module structure,
and hence its restriction to the ground field, see
\cite[Chapter II, \S 2, no. 38, p. 231]{FAC}.

\medskip\noindent
Let $f : (X, \mathcal{O}_X) \to (Y, \mathcal{O}_Y)$ be a morphism of ringed
spaces. With $\mathcal{F}[0] \to \mathcal{I}^\bullet$ as above
we define
\begin{equation}
\label{equation-higher-direct-image-modules}
R^if_*\mathcal{F} = H^i(f_*\mathcal{I}^\bullet)
\end{equation}
to be the {\it $i$th higher direct image of $\mathcal{F}$}.
The family of functors $R^if_*$ forms a universal $\delta$-functor
from $\textit{Mod}(\mathcal{O}_X) \to \textit{Mod}(\mathcal{O}_Y)$.





\section{Derived functors}
\label{section-derived-functors}

\noindent
We briefly explain how to get right derived functors using resolution
functors. For the unbounded derived functors, please see
Section \ref{section-unbounded}.

\medskip\noindent
Let $(X, \mathcal{O}_X)$ be a ringed space. The category
$\textit{Mod}(\mathcal{O}_X)$ is abelian, see
Modules, Lemma \ref{modules-lemma-abelian}.
In this chapter we will write
$$
K(\mathcal{O}_X) = K(\textit{Mod}(\mathcal{O}_X))
\quad
\text{and}
\quad
D(\mathcal{O}_X) = D(\textit{Mod}(\mathcal{O}_X)).
$$
and similarly for the bounded versions for the triangulated categories
introduced in
Derived Categories, Definition \ref{derived-definition-complexes-notation} and
Definition \ref{derived-definition-unbounded-derived-category}.
By
Derived Categories, Remark \ref{derived-remark-big-abelian-category}
there exists a resolution functor
$$
j = j_X :
K^{+}(\textit{Mod}(\mathcal{O}_X))
\longrightarrow
K^{+}(\mathcal{I})
$$
where $\mathcal{I}$ is the strictly full additive subcategory of
$\textit{Mod}(\mathcal{O}_X)$ consisting of injective sheaves.
For any left exact functor
$F : \textit{Mod}(\mathcal{O}_X) \to \mathcal{B}$
into any abelian category $\mathcal{B}$ we will denote $RF$ the
right derived functor described in
Derived Categories, Section \ref{derived-section-right-derived-functor}
and constructed using the resolution functor $j_X$ just described:
\begin{equation}
\label{equation-RF}
RF = F \circ j_X' : D^{+}(X) \longrightarrow D^{+}(\mathcal{B})
\end{equation}
see
Derived Categories, Lemma \ref{derived-lemma-right-derived-functor}
for notation. Note that we may think of $RF$ as defined on
$\textit{Mod}(\mathcal{O}_X)$,
$\text{Comp}^{+}(\textit{Mod}(\mathcal{O}_X))$,
$K^{+}(X)$, or $D^{+}(X)$
depending on the situation. According to
Derived Categories, Definition \ref{derived-definition-higher-derived-functors}
we obtain the $i$th right derived functor
\begin{equation}
\label{equation-RFi}
R^iF = H^i \circ RF : \textit{Mod}(\mathcal{O}_X) \longrightarrow \mathcal{B}
\end{equation}
so that $R^0F = F$ and $\{R^iF, \delta\}_{i \geq 0}$ is universal
$\delta$-functor, see
Derived Categories, Lemma \ref{derived-lemma-higher-derived-functors}.

\medskip\noindent
Here are two special cases of this construction.
Given a ring $R$ we write $K(R) = K(\text{Mod}_R)$ and
$D(R) = D(\text{Mod}_R)$ and similarly for bounded versions.
For any open $U \subset X$ we have a left exact functor
$
\Gamma(U, -) :
\textit{Mod}(\mathcal{O}_X)
\longrightarrow
\text{Mod}_{\mathcal{O}_X(U)}
$
which gives rise to
\begin{equation}
\label{equation-total-derived-cohomology}
R\Gamma(U, -) :
D^{+}(X)
\longrightarrow
D^{+}(\mathcal{O}_X(U))
\end{equation}
by the discussion above. We set $H^i(U, -) = R^i\Gamma(U, -)$.
If $U = X$ we recover (\ref{equation-cohomology-modules}).
If $f : X \to Y$ is a morphism of ringed spaces, then we have
the left exact functor
$
f_* :
\textit{Mod}(\mathcal{O}_X)
\longrightarrow
\textit{Mod}(\mathcal{O}_Y)
$
which gives rise to the {\it derived pushforward}
\begin{equation}
\label{equation-total-derived-direct-image}
Rf_* :
D^{+}(X)
\longrightarrow
D^{+}(Y)
\end{equation}
The $i$th cohomology sheaf of $Rf_*\mathcal{F}^\bullet$ is denoted
$R^if_*\mathcal{F}^\bullet$ and called the $i$th {\it higher direct image}
in accordance with (\ref{equation-higher-direct-image-modules}).
The two displayed functors above are exact functors
of derived categories.

\medskip\noindent
{\bf Abuse of notation:} When the functor $Rf_*$, or any other
derived functor, is applied to a sheaf $\mathcal{F}$ on $X$ or a complex
of sheaves it is understood that $\mathcal{F}$ has been replaced by a
suitable resolution of $\mathcal{F}$. To facilitate this kind of
operation we will say, given an object
$\mathcal{F}^\bullet \in D(\mathcal{O}_X)$,
that a bounded below complex $\mathcal{I}^\bullet$ of injectives of
$\textit{Mod}(\mathcal{O}_X)$
{\it represents $\mathcal{F}^\bullet$ in the derived category}
if there exists a quasi-isomorphism
$\mathcal{F}^\bullet \to \mathcal{I}^\bullet$. In the same vein the phrase
``let $\alpha : \mathcal{F}^\bullet \to \mathcal{G}^\bullet$ be
a morphism of $D(\mathcal{O}_X)$''
does not mean that $\alpha$ is represented by a
morphism of complexes. If we have an actual morphism of complexes we will
say so.









\section{First cohomology and torsors}
\label{section-h1-torsors}

\begin{definition}
\label{definition-torsor}
Let $X$ be a topological space.
Let $\mathcal{G}$ be a sheaf of (possibly non-commutative) groups on $X$.
A {\it pseudo torsor}, or more precisely a {\it pseudo $\mathcal{G}$-torsor},
is a sheaf of sets $\mathcal{F}$ on $X$ endowed with an action
$\mathcal{G} \times \mathcal{F} \to \mathcal{F}$ such that
\begin{enumerate}
\item whenever $\mathcal{F}(U)$ is nonempty the action
$\mathcal{G}(U) \times \mathcal{F}(U) \to \mathcal{F}(U)$
is simply transitive
\end{enumerate}
A {\it morphism of pseudo $\mathcal{G}$-torsors} $\mathcal{F} \to \mathcal{F}'$
is a morphism of sheaves of sets compatible with the
$\mathcal{G}$-actions. A {\it torsor}, or more precisely a
{\it $\mathcal{G}$-torsor}, is a pseudo $\mathcal{G}$-torsor such
that in addition
\begin{enumerate}
\item[(2)] for every $x \in X$ the stalk $\mathcal{F}_x$ is nonempty.
\end{enumerate}
A {\it morphism of $\mathcal{G}$-torsors} is a morphism of
pseudo $\mathcal{G}$-torsors. The {\it trivial $\mathcal{G}$-torsor}
is the sheaf $\mathcal{G}$ endowed with the obvious left
$\mathcal{G}$-action.
\end{definition}

\noindent
It is clear that a morphism of torsors is automatically an isomorphism.

\begin{lemma}
\label{lemma-trivial-torsor}
Let $X$ be a topological space.
Let $\mathcal{G}$ be a sheaf of (possibly non-commutative) groups on $X$.
A $\mathcal{G}$-torsor $\mathcal{F}$ is trivial if and only if
$\mathcal{F}(X) \not = \emptyset$.
\end{lemma}

\begin{proof}
Omitted.
\end{proof}

\begin{lemma}[Crossed homomorphisms]
\label{lemma-crossed-homomorphism-affine-action}
\begin{reference}
\cite[Chapter III, \S 3.4, p. 164]{Tohoku}
\end{reference}
Let $X$ be a topological space. Let $\mathcal P$ and $\mathcal A$ be
sheaves of groups on $X$, and let $\mathcal P$ act on $\mathcal A$ by
group automorphisms. Write $\sigma_p(a)$ for the action on local sections.
Let
$$
u : \mathcal P \longrightarrow \mathcal A
$$
be a morphism of sheaves of sets such that
$$
u(1)=1
\quad\text{and}\quad
u(pp')=u(p)\sigma_p(u(p'))
$$
for all local sections $p,p'$ of $\mathcal P$. Then the following hold.
\begin{enumerate}
\item The inverse image $\mathcal G=u^{-1}(1)$ is a subsheaf of groups of
$\mathcal P$.
\item For local sections $p,p'$ over the same open,
$$
u(p)=u(p')
\quad\Longleftrightarrow\quad
p^{-1}p'\in\mathcal G.
$$
Thus every nonempty fibre of $u$ is an orbit for right multiplication by
$\mathcal G$, namely a subset $p\mathcal G$.
\item The formula
$$
\rho_p(a)=u(p)\sigma_p(a)
$$
defines an action of $\mathcal P$ on the underlying sheaf of sets of
$\mathcal A$. It satisfies
$$
\rho_p(ab)=\rho_p(a)\sigma_p(b),
\qquad
u(p)=\rho_p(1),
$$
and $\mathcal G$ is the stabilizer subsheaf of $1$ for this action.
\end{enumerate}
\end{lemma}

\begin{proof}
The crossed homomorphism identity gives
$$
u(p^{-1})=\sigma_{p^{-1}}(u(p)^{-1})
$$
and hence
$$
u(p^{-1}p')
=\sigma_{p^{-1}}(u(p)^{-1}u(p')).
$$
The first formula shows that $\mathcal G$ is closed under inverses; the
crossed homomorphism identity shows that it is closed under products. The
second formula proves the equivalence in (2).

For local sections $p,p'$ and $a$ we have
\begin{align*}
\rho_p(\rho_{p'}(a))
&=u(p)\sigma_p(u(p'))\sigma_{pp'}(a) \\
&=u(pp')\sigma_{pp'}(a) \\
&=\rho_{pp'}(a).
\end{align*}
Also $\rho_1$ is the identity, so the maps $\rho_p$ define the asserted
action. Finally,
$$
\rho_p(ab)
=u(p)\sigma_p(a)\sigma_p(b)
=\rho_p(a)\sigma_p(b).
$$
The remaining statements follow by setting $a=1$. All constructions commute
with restriction, so they define the indicated subsheaves and sheaf actions.
\end{proof}

\begin{lemma}[The torsor obstruction of a crossed homomorphism]
\label{lemma-crossed-homomorphism-torsor-obstruction}
\begin{reference}
\cite[Chapter III, \S 3.4, p. 164]{Tohoku}
\end{reference}
In the situation of
Lemma \ref{lemma-crossed-homomorphism-affine-action},
let $\rho$ be the action constructed there and let
$\mathcal{G}=u^{-1}(1)$. Let $\mathcal{E}$ be a
$\mathcal{P}$-torsor. If $\tau$ is an action of $\mathcal{P}$ on a
sheaf of sets $\mathcal{B}$, denote by
$\mathcal{B}(\tau)^{\mathcal{E}}$ the sheaf associated to the
presheaf
$$
U \longmapsto
\mathcal{P}(U)\backslash
\bigl(\mathcal{E}(U)\times\mathcal{B}(U)\bigr),
$$
where
$$
p\cdot(e,b)=(p\cdot e,\tau_p(b)).
$$
Then the following hold.
\begin{enumerate}
\item The sheaf $\mathcal{A}(\sigma)^{\mathcal{E}}$ is naturally a
sheaf of groups, and multiplication in $\mathcal{A}$ induces a right
action
$$
\mathcal{A}(\rho)^{\mathcal{E}}\times
\mathcal{A}(\sigma)^{\mathcal{E}}
\longrightarrow
\mathcal{A}(\rho)^{\mathcal{E}}.
$$
This action is locally simply transitive and
$\mathcal{A}(\rho)^{\mathcal{E}}$ is locally nonempty.
\item Set
$$
\mathcal{H}=
\bigl(\mathcal{A}(\sigma)^{\mathcal{E}}\bigr)^{\mathrm{opp}}.
$$
The formula
$$
h^{\mathrm{opp}}\cdot s=s\cdot h
$$
makes $\mathcal{A}(\rho)^{\mathcal{E}}$ an
$\mathcal{H}$-torsor. Denote its isomorphism class in the pointed set
of isomorphism classes of $\mathcal{H}$-torsors by
$d_u(\mathcal{E})$. This class is neutral if and only if
$$
\mathcal{A}(\rho)^{\mathcal{E}}(X)\neq\emptyset.
$$
If $\mathcal{A}(\sigma)^{\mathcal{E}}$ is abelian, then
Lemma \ref{lemma-torsors-h1} identifies this class with an element
$$
d_u(\mathcal{E})\in
H^1\bigl(X,\mathcal{A}(\sigma)^{\mathcal{E}}\bigr).
$$
\item Let $\mathcal{P}/\mathcal{G}$ be the sheafification of the
presheaf of orbits for right multiplication by $\mathcal{G}$, whose
local representatives are $p\mathcal{G}$, with its canonical left
$\mathcal{P}$-action. The rule
$$
p\mathcal{G}\longmapsto u(p)
$$
defines a $\mathcal{P}$-equivariant monomorphism
$$
\mathcal{P}/\mathcal{G}\longrightarrow\mathcal{A},
$$
where $\mathcal{P}$ acts on $\mathcal{A}$ through $\rho$. It is an
isomorphism if and only if $u$ is an epimorphism. Consequently there
is a natural monomorphism
$$
(\mathcal{P}/\mathcal{G})^{\mathcal{E}}
\longrightarrow
\mathcal{A}(\rho)^{\mathcal{E}},
$$
which is an isomorphism if and only if $u$ is an epimorphism.
\item The torsor $\mathcal{E}$ admits a reduction of structure group
from $\mathcal{P}$ to $\mathcal{G}$ if and only if
$$
(\mathcal{P}/\mathcal{G})^{\mathcal{E}}(X)\neq\emptyset.
$$
Thus a reduction implies that $d_u(\mathcal{E})$ is neutral. If $u$
is an epimorphism, the converse holds as well.
\end{enumerate}
\end{lemma}

\begin{proof}
The action $\sigma$ is by group automorphisms. Hence the multiplication,
identity, and inverse maps of $\mathcal{A}$ are
$\mathcal{P}$-equivariant for $\sigma$, and they induce the asserted
group structure on $\mathcal{A}(\sigma)^{\mathcal{E}}$. The identity
$$
\rho_p(ab)=\rho_p(a)\sigma_p(b)
$$
from Lemma \ref{lemma-crossed-homomorphism-affine-action} shows that
multiplication induces the displayed right action.

Locally on $X$ we may choose a section $e$ of $\mathcal{E}$. On such
an open the rule $[e,a]\mapsto a$ identifies both associated sheaves
with the corresponding copy of $\mathcal{A}$, and the right action
with regular right multiplication. It is therefore locally simply
transitive and its source is locally nonempty. If $h_1,h_2$ are local
sections of $\mathcal{A}(\sigma)^{\mathcal{E}}$, then
$$
h_1^{\mathrm{opp}}\cdot(h_2^{\mathrm{opp}}\cdot s)
=(s\cdot h_2)\cdot h_1
=s\cdot(h_2h_1)
=(h_1^{\mathrm{opp}}h_2^{\mathrm{opp}})\cdot s.
$$
Thus this is a left action of the opposite group and it makes
$\mathcal{A}(\rho)^{\mathcal{E}}$ an $\mathcal{H}$-torsor. The
assertion about neutrality follows from Lemma \ref{lemma-trivial-torsor};
the assertion in the abelian case follows from
Lemma \ref{lemma-torsors-h1}.

If $g$ is a local section of $\mathcal{G}$, then
$$
u(pg)=u(p)\sigma_p(u(g))=u(p).
$$
Conversely, the preceding lemma shows that $u(p)=u(p')$ if and only if
$p^{-1}p'$ lies in $\mathcal{G}$. Thus $p\mathcal{G}\mapsto u(p)$ is
well defined and is a monomorphism. It is $\mathcal{P}$-equivariant
because
$$
u(qp)=u(q)\sigma_q(u(p))=\rho_q(u(p)).
$$
On stalks it is surjective exactly when $u$ is surjective. This proves
the assertions about this map. After passing to an open on which
$\mathcal{E}$ has a section, its associated map identifies with the
original map, proving the corresponding assertions for the twists.

For completeness, suppose first that $\mathcal{E}$ has a reduction
$\mathcal{F}$ to $\mathcal{G}$. Locally choosing a section $f$ of
$\mathcal{F}$ gives the section $[f,\mathcal{G}]$ of
$(\mathcal{P}/\mathcal{G})^{\mathcal{E}}$; this is independent of the
choice of $f$ and hence glues. Conversely, given a global section $s$
of $(\mathcal{P}/\mathcal{G})^{\mathcal{E}}$, define
$$
\mathcal{F}(U)=
\{e\in\mathcal{E}(U)\mid [e,\mathcal{G}]=s|_U\}.
$$
Locally every value of $s$ has the form $[e,p\mathcal{G}]$, and then
$p^{-1}e$ lies in $\mathcal{F}$. It follows that $\mathcal{F}$ is
locally nonempty. The simply transitive action on $\mathcal{E}$ shows
that $\mathcal{F}$ is a $\mathcal{G}$-torsor, and the natural map from
the induced $\mathcal{P}$-torsor
$\mathcal{P}\times^{\mathcal{G}}\mathcal{F}$ to $\mathcal{E}$ is an
isomorphism. This proves the reduction criterion. The final assertions
follow by applying the comparison monomorphism, which is an isomorphism
when $u$ is an epimorphism.
\end{proof}

\begin{lemma}
\label{lemma-torsors-h1}
Let $X$ be a topological space.
Let $\mathcal{H}$ be an abelian sheaf on $X$.
There is a canonical bijection between the set of isomorphism
classes of $\mathcal{H}$-torsors and $H^1(X, \mathcal{H})$.
\end{lemma}

\begin{proof}
Let $\mathcal{F}$ be a $\mathcal{H}$-torsor.
Consider the free abelian sheaf $\mathbf{Z}[\mathcal{F}]$
on $\mathcal{F}$. It is the sheafification of the rule
which associates to $U \subset X$ open the collection of finite
formal sums $\sum n_i[s_i]$ with $n_i \in \mathbf{Z}$
and $s_i \in \mathcal{F}(U)$. There is a natural map
$$
\sigma : \mathbf{Z}[\mathcal{F}] \longrightarrow \underline{\mathbf{Z}}
$$
which to a local section $\sum n_i[s_i]$ associates $\sum n_i$.
The kernel of $\sigma$ is generated by the local section of the form
$[s] - [s']$. There is a canonical map
$a : \Ker(\sigma) \to \mathcal{H}$
which maps $[s] - [s'] \mapsto h$ where $h$ is the local section of
$\mathcal{H}$ such that $h \cdot s' = s$. Consider the pushout diagram
$$
\xymatrix{
0 \ar[r] &
\Ker(\sigma) \ar[r] \ar[d]^a &
\mathbf{Z}[\mathcal{F}] \ar[r] \ar[d] &
\underline{\mathbf{Z}} \ar[r] \ar[d] &
0 \\
0 \ar[r] &
\mathcal{H} \ar[r] &
\mathcal{E} \ar[r] &
\underline{\mathbf{Z}} \ar[r] &
0
}
$$
Here $\mathcal{E}$ is the extension obtained by pushout.
From the long exact cohomology sequence associated to the lower
short exact sequence we obtain an element
$\xi = \xi_\mathcal{F} \in H^1(X, \mathcal{H})$
by applying the boundary operator to $1 \in H^0(X, \underline{\mathbf{Z}})$.

\medskip\noindent
Conversely, given $\xi \in H^1(X, \mathcal{H})$ we can associate to
$\xi$ a torsor as follows. Choose an embedding $\mathcal{H} \to \mathcal{I}$
of $\mathcal{H}$ into an injective abelian sheaf $\mathcal{I}$. We set
$\mathcal{Q} = \mathcal{I}/\mathcal{H}$ so that we have a short exact
sequence
$$
\xymatrix{
0 \ar[r] &
\mathcal{H} \ar[r] &
\mathcal{I} \ar[r] &
\mathcal{Q} \ar[r] &
0
}
$$
The element $\xi$ is the image of a global section $q \in H^0(X, \mathcal{Q})$
because $H^1(X, \mathcal{I}) = 0$ (see
Derived Categories, Lemma \ref{derived-lemma-higher-derived-functors}).
Let $\mathcal{F} \subset \mathcal{I}$ be the subsheaf (of sets) of sections
that map to $q$ in the sheaf $\mathcal{Q}$. It is easy to verify that
$\mathcal{F}$ is a torsor.

\medskip\noindent
We omit the verification that the two constructions given
above are mutually inverse.
\end{proof}







\section{First cohomology and extensions}
\label{section-h1-extensions}

\begin{lemma}
\label{lemma-h1-extensions}
Let $(X, \mathcal{O}_X)$ be a ringed space. Let $\mathcal{F}$ be a sheaf of
$\mathcal{O}_X$-modules. There is a canonical bijection
$$
\Ext^1_{\textit{Mod}(\mathcal{O}_X)}(\mathcal{O}_X, \mathcal{F})
\longrightarrow
H^1(X, \mathcal{F})
$$
which associates to the extension
$$
0 \to \mathcal{F} \to \mathcal{E} \to \mathcal{O}_X \to 0
$$
the image of $1 \in \Gamma(X, \mathcal{O}_X)$ in $H^1(X, \mathcal{F})$.
\end{lemma}

\begin{proof}
Let us construct the inverse of the map given in the lemma. Let
$\xi \in H^1(X, \mathcal{F})$. Choose an injection
$\mathcal{F} \subset \mathcal{I}$ with $\mathcal{I}$ injective in
$\textit{Mod}(\mathcal{O}_X)$.
Set $\mathcal{Q} = \mathcal{I}/\mathcal{F}$.
By the long exact sequence of cohomology, we see that
$\xi$ is the image of a section
$\tilde \xi \in \Gamma(X, \mathcal{Q}) =
\Hom_{\mathcal{O}_X}(\mathcal{O}_X, \mathcal{Q})$.
Now, we just form the pullback
$$
\xymatrix{
0 \ar[r] &
\mathcal{F} \ar[r] \ar@{=}[d] &
\mathcal{E} \ar[r] \ar[d] &
\mathcal{O}_X \ar[r] \ar[d]^{\tilde \xi} &
0 \\
0 \ar[r] &
\mathcal{F} \ar[r] &
\mathcal{I} \ar[r] &
\mathcal{Q} \ar[r] &
0
}
$$
see Homology, Section \ref{homology-section-extensions}.
\end{proof}








\section{First cohomology and invertible sheaves}
\label{section-invertible-sheaves}

\noindent
The Picard group of a ringed space is defined in
Modules, Section \ref{modules-section-invertible}.

\begin{lemma}
\label{lemma-h1-invertible}
Let $(X, \mathcal{O}_X)$ be a ringed space. If all stalks
$\mathcal{O}_{X, x}$ are local rings, then there is a canonical isomorphism
$$
H^1(X, \mathcal{O}_X^*) = \Pic(X).
$$
of abelian groups.
\end{lemma}

\begin{proof}
Let $\mathcal{L}$ be an invertible $\mathcal{O}_X$-module.
Consider the presheaf $\mathcal{L}^*$ defined by the rule
$$
U \longmapsto \{s \in \mathcal{L}(U)
\text{ such that } \mathcal{O}_U \xrightarrow{s \cdot -} \mathcal{L}_U
\text{ is an isomorphism}\}
$$
This presheaf satisfies the sheaf condition. Moreover, if
$f \in \mathcal{O}_X^*(U)$ and $s \in \mathcal{L}^*(U)$, then clearly
$fs \in \mathcal{L}^*(U)$. By the same token, if $s, s' \in \mathcal{L}^*(U)$
then there exists a unique $f \in \mathcal{O}_X^*(U)$ such that
$fs = s'$. Moreover, the sheaf $\mathcal{L}^*$ has sections locally
by Modules, Lemma \ref{modules-lemma-invertible-is-locally-free-rank-1}.
In other words we
see that $\mathcal{L}^*$ is a $\mathcal{O}_X^*$-torsor. Thus we get
a map
$$
\begin{matrix}
\text{invertible sheaves on }(X, \mathcal{O}_X) \\
\text{ up to isomorphism}
\end{matrix}
\longrightarrow
\begin{matrix}
\mathcal{O}_X^*\text{-torsors} \\
\text{ up to isomorphism}
\end{matrix}
$$
We omit the verification that this is a homomorphism of abelian groups.
By
Lemma \ref{lemma-torsors-h1}
the right hand side is canonically
bijective to $H^1(X, \mathcal{O}_X^*)$.
Thus we have to show this map is injective and surjective.

\medskip\noindent
Injective. If the torsor $\mathcal{L}^*$ is trivial, this means by
Lemma \ref{lemma-trivial-torsor}
that $\mathcal{L}^*$ has a global section.
Hence this means exactly that $\mathcal{L} \cong \mathcal{O}_X$ is
the neutral element in $\Pic(X)$.

\medskip\noindent
Surjective. Let $\mathcal{F}$ be an $\mathcal{O}_X^*$-torsor.
Consider the presheaf of sets
$$
\mathcal{L}_1 : U \longmapsto
(\mathcal{F}(U) \times \mathcal{O}_X(U))/\mathcal{O}_X^*(U)
$$
where the action of $f \in \mathcal{O}_X^*(U)$ on
$(s, g)$ is $(fs, f^{-1}g)$. Then $\mathcal{L}_1$ is a presheaf
of $\mathcal{O}_X$-modules by setting
$(s, g) + (s', g') = (s, g + (s'/s)g')$ where $s'/s$ is the local
section $f$ of $\mathcal{O}_X^*$ such that $fs = s'$, and
$h(s, g) = (s, hg)$ for $h$ a local section of $\mathcal{O}_X$.
We omit the verification that the sheafification
$\mathcal{L} = \mathcal{L}_1^\#$ is an invertible $\mathcal{O}_X$-module
whose associated $\mathcal{O}_X^*$-torsor $\mathcal{L}^*$ is isomorphic
to $\mathcal{F}$.
\end{proof}













\section{Locality of cohomology}
\label{section-locality}

\noindent
The following lemma says there is no ambiguity in defining the cohomology
of a sheaf $\mathcal{F}$ over an open.

\begin{lemma}
\label{lemma-cohomology-of-open}
Let $X$ be a ringed space.
Let $U \subset X$ be an open subspace.
\begin{enumerate}
\item If $\mathcal{I}$ is an injective $\mathcal{O}_X$-module
then $\mathcal{I}|_U$ is an injective $\mathcal{O}_U$-module.
\item For any sheaf of $\mathcal{O}_X$-modules $\mathcal{F}$ we have
$H^p(U, \mathcal{F}) = H^p(U, \mathcal{F}|_U)$.
\end{enumerate}
\end{lemma}

\begin{proof}
Denote $j : U \to X$ the open immersion.
Recall that the functor $j^{-1}$ of restriction to $U$ is a right adjoint
to the functor $j_!$ of extension by $0$, see
Sheaves, Lemma \ref{sheaves-lemma-j-shriek-modules}.
Moreover, $j_!$ is exact. Hence (1) follows from
Homology, Lemma \ref{homology-lemma-adjoint-preserve-injectives}.

\medskip\noindent
By definition $H^p(U, \mathcal{F}) = H^p(\Gamma(U, \mathcal{I}^\bullet))$
where $\mathcal{F} \to \mathcal{I}^\bullet$ is an injective resolution
in $\textit{Mod}(\mathcal{O}_X)$.
By the above we see that $\mathcal{F}|_U \to \mathcal{I}^\bullet|_U$
is an injective resolution in $\textit{Mod}(\mathcal{O}_U)$.
Hence $H^p(U, \mathcal{F}|_U)$ is equal to
$H^p(\Gamma(U, \mathcal{I}^\bullet|_U))$.
Of course $\Gamma(U, \mathcal{F}) = \Gamma(U, \mathcal{F}|_U)$ for
any sheaf $\mathcal{F}$ on $X$.
Hence the equality
in (2).
\end{proof}

\noindent
Let $X$ be a ringed space.
Let $\mathcal{F}$ be a sheaf of $\mathcal{O}_X$-modules.
Let $U \subset V \subset X$ be open subsets.
Then there is a canonical {\it restriction mapping}
\begin{equation}
\label{equation-restriction-mapping}
H^n(V, \mathcal{F})
\longrightarrow
H^n(U, \mathcal{F}), \quad
\xi \longmapsto \xi|_U
\end{equation}
functorial in $\mathcal{F}$. Namely, choose any injective
resolution $\mathcal{F} \to \mathcal{I}^\bullet$. The restriction
mappings of the sheaves $\mathcal{I}^p$ give a morphism of complexes
$$
\Gamma(V, \mathcal{I}^\bullet)
\longrightarrow
\Gamma(U, \mathcal{I}^\bullet)
$$
The LHS is a complex representing $R\Gamma(V, \mathcal{F})$
and the RHS is a complex representing $R\Gamma(U, \mathcal{F})$.
We get the map on cohomology groups by applying the functor $H^n$.
As indicated we will use the notation $\xi \mapsto \xi|_U$ to denote this map.
Thus the rule $U \mapsto H^n(U, \mathcal{F})$ is a presheaf of
$\mathcal{O}_X$-modules. This presheaf is customarily denoted
$\underline{H}^n(\mathcal{F})$. We will give another interpretation
of this presheaf in Lemma \ref{lemma-include}.

\begin{lemma}
\label{lemma-kill-cohomology-class-on-covering}
Let $X$ be a ringed space.
Let $\mathcal{F}$ be a sheaf of $\mathcal{O}_X$-modules.
Let $U \subset X$ be an open subspace.
Let $n > 0$ and let $\xi \in H^n(U, \mathcal{F})$.
Then there exists an open covering
$U = \bigcup_{i\in I} U_i$ such that $\xi|_{U_i} = 0$ for
all $i \in I$.
\end{lemma}

\begin{proof}
Let $\mathcal{F} \to \mathcal{I}^\bullet$ be an injective resolution.
Then
$$
H^n(U, \mathcal{F}) =
\frac{\Ker(\mathcal{I}^n(U) \to \mathcal{I}^{n + 1}(U))}
{\Im(\mathcal{I}^{n - 1}(U) \to \mathcal{I}^n(U))}.
$$
Pick an element $\tilde \xi \in \mathcal{I}^n(U)$ representing the
cohomology class in the presentation above. Since $\mathcal{I}^\bullet$
is an injective resolution of $\mathcal{F}$ and $n > 0$ we see that
the complex $\mathcal{I}^\bullet$ is exact in degree $n$. Hence
$\Im(\mathcal{I}^{n - 1} \to \mathcal{I}^n) =
\Ker(\mathcal{I}^n \to \mathcal{I}^{n + 1})$ as sheaves.
Since $\tilde \xi$ is a section of the kernel sheaf over $U$
we conclude there exists an open covering $U = \bigcup_{i \in I} U_i$
such that $\tilde \xi|_{U_i}$ is the image under $d$ of a section
$\xi_i \in \mathcal{I}^{n - 1}(U_i)$. By our definition of the
restriction $\xi|_{U_i}$ as corresponding to the class of
$\tilde \xi|_{U_i}$ we conclude.
\end{proof}

\begin{lemma}
\label{lemma-describe-higher-direct-images}
Let $f : X \to Y$ be a morphism of ringed spaces.
Let $\mathcal{F}$ be a $\mathcal{O}_X$-module.
The sheaves $R^if_*\mathcal{F}$ are the sheaves
associated to the presheaves
$$
V \longmapsto H^i(f^{-1}(V), \mathcal{F})
$$
with restriction mappings as in Equation (\ref{equation-restriction-mapping}).
There is a similar statement for $R^if_*$ applied to a
bounded below complex $\mathcal{F}^\bullet$.
\end{lemma}

\begin{proof}
Let $\mathcal{F} \to \mathcal{I}^\bullet$ be an injective resolution.
Then $R^if_*\mathcal{F}$ is by definition the $i$th cohomology sheaf
of the complex
$$
f_*\mathcal{I}^0 \to f_*\mathcal{I}^1 \to f_*\mathcal{I}^2 \to \ldots
$$
By definition of the abelian category structure on $\mathcal{O}_Y$-modules
this cohomology sheaf is the sheaf associated to the presheaf
$$
V
\longmapsto
\frac{\Ker(f_*\mathcal{I}^i(V) \to f_*\mathcal{I}^{i + 1}(V))}
{\Im(f_*\mathcal{I}^{i - 1}(V) \to f_*\mathcal{I}^i(V))}
$$
and this is obviously equal to
$$
\frac{\Ker(\mathcal{I}^i(f^{-1}(V)) \to \mathcal{I}^{i + 1}(f^{-1}(V)))}
{\Im(\mathcal{I}^{i - 1}(f^{-1}(V)) \to \mathcal{I}^i(f^{-1}(V)))}
$$
which is equal to $H^i(f^{-1}(V), \mathcal{F})$
and we win.
\end{proof}

\begin{lemma}
\label{lemma-derived-functor-sheafification}
\begin{reference}
\cite[Chapter III, \S 3.7, Lemma 3.7.2, p. 173]{Tohoku}
\end{reference}
Let $X$ be a topological space, let $\mathcal{A}$ be an abelian category
with enough injectives, and let
$$
T : \mathcal{A} \longrightarrow \textit{Ab}(X)
$$
be an additive functor. For $A \in \mathcal{A}$ and $i \geq 0$, the
sheaf $R^iT(A)$ is the sheaf associated to the presheaf
$$
U \longmapsto R^i(\Gamma(U, -) \circ T)(A)
$$
on $X$.
\end{lemma}

\begin{proof}
Choose an injective resolution $A[0] \to I^\bullet$. This resolution
computes both right derived functors. Hence
$$
R^iT(A) = H^i(T(I^\bullet))
$$
and, for every open $U \subset X$,
$$
R^i(\Gamma(U, -) \circ T)(A)
= H^i(\Gamma(U, T(I^\bullet))).
$$
The cohomology sheaf of a complex of abelian sheaves is the sheaf
associated to the presheaf obtained by taking the cohomology of its
complex of sections over every open subset. Applying this observation to
$T(I^\bullet)$ proves the lemma.
\end{proof}

\begin{lemma}
\label{lemma-localize-higher-direct-images}
Let $f : X \to Y$ be a morphism of ringed spaces.
Let $\mathcal{F}$ be an $\mathcal{O}_X$-module.
Let $V \subset Y$ be an open subspace.
Denote $g : f^{-1}(V) \to V$ the restriction of $f$.
Then we have
$$
R^pg_*(\mathcal{F}|_{f^{-1}(V)}) = (R^pf_*\mathcal{F})|_V
$$
There is a similar statement for the
derived image $Rf_*\mathcal{F}^\bullet$ where $\mathcal{F}^\bullet$
is a bounded below complex of $\mathcal{O}_X$-modules.
\end{lemma}

\begin{proof}
First proof. Apply Lemmas \ref{lemma-describe-higher-direct-images}
and \ref{lemma-cohomology-of-open} to see the displayed equality.
Second proof. Choose an injective resolution
$\mathcal{F} \to \mathcal{I}^\bullet$
and use that $\mathcal{F}|_{f^{-1}(V)} \to \mathcal{I}^\bullet|_{f^{-1}(V)}$
is an injective resolution also.
\end{proof}

\begin{remark}
\label{remark-daniel}
Here is a different approach to the proofs of
Lemmas \ref{lemma-kill-cohomology-class-on-covering} and
\ref{lemma-describe-higher-direct-images} above.
Let $(X, \mathcal{O}_X)$ be a ringed space.
Let $i_X : \textit{Mod}(\mathcal{O}_X) \to \textit{PMod}(\mathcal{O}_X)$
be the inclusion functor and let $\#$ be the sheafification functor.
Recall that $i_X$ is left exact and $\#$ is exact.
\begin{enumerate}
\item First prove Lemma \ref{lemma-include} below which says that the
right derived functors of $i_X$ are given by
$R^pi_X\mathcal{F} = \underline{H}^p(\mathcal{F})$.
Here is another proof: The equality is clear for $p = 0$.
Both $(R^pi_X)_{p \geq 0}$ and $(\underline{H}^p)_{p \geq 0}$
are delta functors vanishing on injectives, hence both are universal,
hence they are isomorphic. See Homology,
Section \ref{homology-section-cohomological-delta-functor}.
\item A restatement of Lemma \ref{lemma-kill-cohomology-class-on-covering}
is that $(\underline{H}^p(\mathcal{F}))^\# = 0$, $p > 0$ for any sheaf of
$\mathcal{O}_X$-modules $\mathcal{F}$.
To see this is true, use that ${}^\#$ is exact so
$$
(\underline{H}^p(\mathcal{F}))^\# =
(R^pi_X\mathcal{F})^\# =
R^p(\# \circ i_X)(\mathcal{F}) = 0
$$
because $\# \circ i_X$ is the identity functor.
\item Let $f : X \to Y$ be a morphism of ringed spaces.
Let $\mathcal{F}$ be an $\mathcal{O}_X$-module. The presheaf
$V \mapsto H^p(f^{-1}V, \mathcal{F})$ is equal to
$R^p (i_Y \circ f_*)\mathcal{F}$. You can prove this by noticing that
both give universal delta functors as in the argument of (1) above.
Hence Lemma \ref{lemma-describe-higher-direct-images}
says that $R^p f_* \mathcal{F}= (R^p (i_Y \circ f_*)\mathcal{F})^\#$.
Again using that $\#$ is exact a that $\# \circ i_Y$ is the identity
functor we see that
$$
R^p f_* \mathcal{F} =
R^p(\# \circ i_Y \circ f_*)\mathcal{F} =
(R^p (i_Y \circ f_*)\mathcal{F})^\#
$$
as desired.
\end{enumerate}
\end{remark}






\section{Mayer-Vietoris}
\label{section-mayer-vietoris}

\noindent
Below will construct the {\v C}ech-to-cohomology spectral sequence, see
Lemma \ref{lemma-cech-spectral-sequence}.
A special case of that spectral sequence is the Mayer-Vietoris
long exact sequence. Since it is such a basic, useful and easy to understand
variant of the spectral sequence we treat it here separately.

\begin{lemma}
\label{lemma-injective-restriction-surjective}
\begin{slogan}
Injectives are flasque.
\end{slogan}
Let $X$ be a ringed space.
Let $U' \subset U \subset X$ be open subspaces.
For any injective $\mathcal{O}_X$-module $\mathcal{I}$ the
restriction mapping
$\mathcal{I}(U) \to \mathcal{I}(U')$ is surjective.
\end{lemma}

\begin{proof}
Let $j : U \to X$ and $j' : U' \to X$ be the open immersions.
Recall that $j_!\mathcal{O}_U$ is the extension by zero of
$\mathcal{O}_U = \mathcal{O}_X|_U$, see
Sheaves, Section \ref{sheaves-section-open-immersions}.
Since $j_!$ is a left adjoint to restriction we see that
for any sheaf $\mathcal{F}$ of $\mathcal{O}_X$-modules
$$
\Hom_{\mathcal{O}_X}(j_!\mathcal{O}_U, \mathcal{F})
=
\Hom_{\mathcal{O}_U}(\mathcal{O}_U, \mathcal{F}|_U)
=
\mathcal{F}(U)
$$
see Sheaves, Lemma \ref{sheaves-lemma-j-shriek-modules}.
Similarly, the sheaf $j'_!\mathcal{O}_{U'}$ represents the
functor $\mathcal{F} \mapsto \mathcal{F}(U')$.
Moreover there
is an obvious canonical map of $\mathcal{O}_X$-modules
$$
j'_!\mathcal{O}_{U'} \longrightarrow j_!\mathcal{O}_U
$$
which corresponds to the restriction mapping
$\mathcal{F}(U) \to \mathcal{F}(U')$ via Yoneda's lemma
(Categories, Lemma \ref{categories-lemma-yoneda}). By the description
of the stalks of the sheaves
$j'_!\mathcal{O}_{U'}$, $j_!\mathcal{O}_U$
we see that the displayed map above is injective (see lemma cited above).
Hence if $\mathcal{I}$ is an injective $\mathcal{O}_X$-module,
then the map
$$
\Hom_{\mathcal{O}_X}(j_!\mathcal{O}_U, \mathcal{I})
\longrightarrow
\Hom_{\mathcal{O}_X}(j'_!\mathcal{O}_{U'}, \mathcal{I})
$$
is surjective, see
Homology, Lemma \ref{homology-lemma-characterize-injectives}.
Putting everything together we obtain the lemma.
\end{proof}

\begin{lemma}[Mayer-Vietoris]
\label{lemma-mayer-vietoris}
Let $X$ be a ringed space. Suppose that $X = U \cup V$ is a
union of two open subsets. For every $\mathcal{O}_X$-module $\mathcal{F}$
there exists a long exact cohomology sequence
$$
0 \to
H^0(X, \mathcal{F}) \to
H^0(U, \mathcal{F}) \oplus H^0(V, \mathcal{F}) \to
H^0(U \cap V, \mathcal{F}) \to
H^1(X, \mathcal{F}) \to \ldots
$$
This long exact sequence is functorial in $\mathcal{F}$.
\end{lemma}

\begin{proof}
The sheaf condition says that the kernel of
$(1, -1) : \mathcal{F}(U) \oplus \mathcal{F}(V) \to \mathcal{F}(U \cap V)$
is equal to the image of $\mathcal{F}(X)$ by the first map
for any abelian sheaf $\mathcal{F}$.
Lemma \ref{lemma-injective-restriction-surjective} above implies that the map
$(1, -1) : \mathcal{I}(U) \oplus \mathcal{I}(V) \to \mathcal{I}(U \cap V)$
is surjective whenever $\mathcal{I}$ is an injective $\mathcal{O}_X$-module.
Hence if $\mathcal{F} \to \mathcal{I}^\bullet$ is an injective resolution
of $\mathcal{F}$, then we get a short exact sequence of complexes
$$
0 \to
\mathcal{I}^\bullet(X) \to
\mathcal{I}^\bullet(U) \oplus \mathcal{I}^\bullet(V) \to
\mathcal{I}^\bullet(U \cap V) \to
0.
$$
Taking cohomology gives the result (use
Homology, Lemma \ref{homology-lemma-long-exact-sequence-cochain}).
We omit the proof of the functoriality of the sequence.
\end{proof}

\begin{lemma}[Relative Mayer-Vietoris]
\label{lemma-relative-mayer-vietoris}
Let $f : X \to Y$ be a morphism of ringed spaces.
Suppose that $X = U \cup V$ is a union of two open subsets.
Denote $a = f|_U : U \to Y$, $b = f|_V : V \to Y$, and
$c = f|_{U \cap V} : U \cap V \to Y$.
For every $\mathcal{O}_X$-module $\mathcal{F}$
there exists a long exact sequence
$$
0 \to
f_*\mathcal{F} \to
a_*(\mathcal{F}|_U) \oplus b_*(\mathcal{F}|_V) \to
c_*(\mathcal{F}|_{U \cap V}) \to
R^1f_*\mathcal{F} \to \ldots
$$
This long exact sequence is functorial in $\mathcal{F}$.
\end{lemma}

\begin{proof}
Let $\mathcal{F} \to \mathcal{I}^\bullet$ be an injective resolution
of $\mathcal{F}$. We claim that we
get a short exact sequence of complexes
$$
0 \to
f_*\mathcal{I}^\bullet \to
a_*\mathcal{I}^\bullet|_U \oplus b_*\mathcal{I}^\bullet|_V \to
c_*\mathcal{I}^\bullet|_{U \cap V} \to
0.
$$
Namely, for any open $W \subset Y$, and for any $n \geq 0$ the
corresponding sequence of groups of sections over $W$
$$
0 \to
\mathcal{I}^n(f^{-1}(W)) \to
\mathcal{I}^n(U \cap f^{-1}(W))
\oplus \mathcal{I}^n(V \cap f^{-1}(W)) \to
\mathcal{I}^n(U \cap V \cap f^{-1}(W)) \to
0
$$
was shown to be short exact in the proof of Lemma \ref{lemma-mayer-vietoris}.
The lemma follows by taking cohomology sheaves and using the fact that
$\mathcal{I}^\bullet|_U$ is an injective resolution of $\mathcal{F}|_U$
and similarly for $\mathcal{I}^\bullet|_V$, $\mathcal{I}^\bullet|_{U \cap V}$
see Lemma \ref{lemma-cohomology-of-open}.
\end{proof}




















\section{The {\v C}ech complex and {\v C}ech cohomology}
\label{section-cech}

\noindent
Let $X$ be a topological space.
Let $\mathcal{U} : U = \bigcup_{i \in I} U_i$ be an open covering,
see Topology, Basic notion (\ref{topology-item-covering}).
As is customary we denote
$U_{i_0\ldots i_p} = U_{i_0} \cap \ldots \cap U_{i_p}$ for the
$(p + 1)$-fold intersection of members of $\mathcal{U}$.
Let $\mathcal{F}$ be an abelian presheaf on $X$.
Set
$$
\check{\mathcal{C}}^p(\mathcal{U}, \mathcal{F})
=
\prod\nolimits_{(i_0, \ldots, i_p) \in I^{p + 1}}
\mathcal{F}(U_{i_0\ldots i_p}).
$$
This is an abelian group. For
$s \in \check{\mathcal{C}}^p(\mathcal{U}, \mathcal{F})$ we denote
$s_{i_0\ldots i_p}$ its value in $\mathcal{F}(U_{i_0\ldots i_p})$.
Note that if $s \in \check{\mathcal{C}}^1(\mathcal{U}, \mathcal{F})$
and $i, j \in I$ then $s_{ij}$ and $s_{ji}$ are both elements
of $\mathcal{F}(U_i \cap U_j)$ but there is no imposed
relation between $s_{ij}$ and $s_{ji}$. In other words, we are {\it not}
working with alternating cochains (these will be defined
in Section \ref{section-alternating-cech}).

\begin{remark}
\label{remark-FAC-simplicial-operators}
\begin{reference}
\cite[Chapter I, \S 3, no. 19, pp. 212--213]{FAC}
\end{reference}
The source packages operations on covering cochains as follows. Let
$K_p(I)$ be the free abelian group on sequences of $p + 1$ elements of
$I$. A homomorphism $h : K_p(I) \to K_q(I)$ is called simplicial there
if the coefficient of a sequence $s'$ in $h(s)$ is zero unless every
vertex of $s'$ occurs in $s$. Restriction maps then define a transpose
$$
{}^t h :
\check{\mathcal{C}}^q(\mathcal{U}, \mathcal{F})
\longrightarrow
\check{\mathcal{C}}^p(\mathcal{U}, \mathcal{F}).
$$
This construction is additive, reverses composition, and preserves the
identity. We do not use this full operator calculus. The differential
below is written directly as the alternating sum of the face maps; see
also Simplicial Methods, Section
\ref{simplicial-section-dold-kan-cosimplicial}.
\end{remark}

We define
$$
d : \check{\mathcal{C}}^p(\mathcal{U}, \mathcal{F})
\longrightarrow
\check{\mathcal{C}}^{p + 1}(\mathcal{U}, \mathcal{F})
$$
by the formula
\begin{equation}
\label{equation-d-cech}
d(s)_{i_0\ldots i_{p + 1}}
=
\sum\nolimits_{j = 0}^{p + 1}
(-1)^j
s_{i_0\ldots \hat i_j \ldots i_{p + 1}}|_{U_{i_0\ldots i_{p + 1}}}
\end{equation}
It is straightforward to see that $d \circ d = 0$. In other words
$\check{\mathcal{C}}^\bullet(\mathcal{U}, \mathcal{F})$ is a complex.

\begin{definition}
\label{definition-cech-complex}
\begin{reference}
\cite[Chapter I, \S 3, nos. 18 and 20, pp. 212--213]{FAC}
\end{reference}
Let $X$ be a topological space.
Let $\mathcal{U} : U = \bigcup_{i \in I} U_i$ be an open covering.
Let $\mathcal{F}$ be an abelian presheaf on $X$.
The complex $\check{\mathcal{C}}^\bullet(\mathcal{U}, \mathcal{F})$
is the {\it {\v C}ech complex} associated to $\mathcal{F}$ and the
open covering $\mathcal{U}$. Its cohomology groups
$H^i(\check{\mathcal{C}}^\bullet(\mathcal{U}, \mathcal{F}))$ are
called the {\it {\v C}ech cohomology groups} associated to
$\mathcal{F}$ and the covering $\mathcal{U}$.
They are denoted $\check H^i(\mathcal{U}, \mathcal{F})$.
\end{definition}

\begin{lemma}
\label{lemma-cech-h0}
\begin{reference}
\cite[Chapter I, \S 3, no. 20, Proposition 1, p. 213]{FAC}
\end{reference}
Let $X$ be a topological space.
Let $\mathcal{F}$ be an abelian presheaf on $X$.
The following are equivalent
\begin{enumerate}
\item $\mathcal{F}$ is an abelian sheaf and
\item for every open covering $\mathcal{U} : U = \bigcup_{i \in I} U_i$
the natural map
$$
\mathcal{F}(U) \to \check{H}^0(\mathcal{U}, \mathcal{F})
$$
is bijective.
\end{enumerate}
\end{lemma}

\begin{proof}
This is true since the sheaf condition is exactly that
$\mathcal{F}(U) \to \check{H}^0(\mathcal{U}, \mathcal{F})$
is bijective for every open covering.
\end{proof}

\begin{lemma}
\label{lemma-cech-trivial}
\begin{reference}
\cite[Chapter I, \S 4, no. 29, Lemma 1, p. 221]{FAC}
\end{reference}
Let $X$ be a topological space. Let $\mathcal{F}$ be an abelian presheaf on $X$.
Let $\mathcal{U} : U = \bigcup_{i \in I} U_i$ be an open covering. If
$U_i = U$ for some $i \in I$, then the extended {\v C}ech complex
$$
\mathcal{F}(U) \to \check{\mathcal{C}}^\bullet(\mathcal{U}, \mathcal{F})
$$
obtained by putting $\mathcal{F}(U)$ in degree $-1$ with differential given by
the canonical map of $\mathcal{F}(U)$ into
$\check{\mathcal{C}}^0(\mathcal{U}, \mathcal{F})$
is homotopy equivalent to $0$.
\end{lemma}

\begin{proof}
Fix an element $i \in I$ with $U = U_i$. Observe that
$U_{i_0 \ldots i_p} = U_{i_0 \ldots \hat i_j \ldots i_p}$ if $i_j = i$.
Let us define a homotopy
$$
h :
\prod\nolimits_{i_0 \ldots i_{p + 1}} \mathcal{F}(U_{i_0 \ldots i_{p + 1}})
\longrightarrow
\prod\nolimits_{i_0 \ldots i_p} \mathcal{F}(U_{i_0 \ldots i_p})
$$
by the rule
$$
h(s)_{i_0 \ldots i_p} = s_{i i_0 \ldots i_p}
$$
In other words, $h : \prod_{i_0} \mathcal{F}(U_{i_0}) \to \mathcal{F}(U)$
is projection onto the factor $\mathcal{F}(U_i) = \mathcal{F}(U)$
and in general the map $h$ equals the projection onto the factors
$\mathcal{F}(U_{i i_1 \ldots i_{p + 1}}) =
\mathcal{F}(U_{i_1 \ldots i_{p + 1}})$.
We compute
\begin{align*}
(dh + hd)(s)_{i_0 \ldots i_p}
& =
\sum\nolimits_{j = 0}^p
(-1)^j
h(s)_{i_0 \ldots \hat i_j \ldots i_p}
+
d(s)_{i i_0 \ldots i_p}\\
& =
\sum\nolimits_{j = 0}^p
(-1)^j
s_{i i_0 \ldots \hat i_j \ldots i_p}
+
s_{i_0 \ldots i_p}
+
\sum\nolimits_{j = 0}^p
(-1)^{j + 1}
s_{i i_0 \ldots \hat i_j \ldots i_p} \\
& =
s_{i_0 \ldots i_p}
\end{align*}
This proves the identity map is homotopic to zero as desired.
\end{proof}






\section{{\v C}ech cohomology as a functor on presheaves}
\label{section-cech-functor}

\noindent
Warning: In this section we work almost exclusively with presheaves and
categories of presheaves and the results are completely wrong in the
setting of sheaves and categories of sheaves!

\medskip\noindent
Let $X$ be a ringed space.
Let $\mathcal{U} : U = \bigcup_{i \in I} U_i$ be an open covering.
Let $\mathcal{F}$ be a presheaf of $\mathcal{O}_X$-modules.
We have the {\v C}ech complex
$\check{\mathcal{C}}^\bullet(\mathcal{U}, \mathcal{F})$
of $\mathcal{F}$ just by thinking of $\mathcal{F}$
as a presheaf of abelian groups. However, each term
$\check{\mathcal{C}}^p(\mathcal{U}, \mathcal{F})$ has a natural
structure of a $\mathcal{O}_X(U)$-module and the differential is given by
$\mathcal{O}_X(U)$-module maps. Moreover, it is clear that the
construction
$$
\mathcal{F} \longmapsto \check{\mathcal{C}}^\bullet(\mathcal{U}, \mathcal{F})
$$
is functorial in $\mathcal{F}$. In fact, it is a functor
\begin{equation}
\label{equation-cech-functor}
\check{\mathcal{C}}^\bullet(\mathcal{U}, -) :
\textit{PMod}(\mathcal{O}_X)
\longrightarrow
\text{Comp}^{+}(\text{Mod}_{\mathcal{O}_X(U)})
\end{equation}
see
Derived Categories, Definition \ref{derived-definition-complexes-notation}
for notation. Recall that the category of bounded below complexes
in an abelian category is an abelian category, see
Homology, Lemma \ref{homology-lemma-cat-cochain-abelian}.

\begin{lemma}
\label{lemma-cech-exact-presheaves}
The functor given by Equation (\ref{equation-cech-functor})
is an exact functor (see Homology, Lemma \ref{homology-lemma-exact-functor}).
\end{lemma}

\begin{proof}
For any open $W \subset U$ the functor
$\mathcal{F} \mapsto \mathcal{F}(W)$ is an additive exact functor
from $\textit{PMod}(\mathcal{O}_X)$ to $\text{Mod}_{\mathcal{O}_X(U)}$.
The terms
$\check{\mathcal{C}}^p(\mathcal{U}, \mathcal{F})$
of the complex are products of these exact functors and hence exact.
Moreover a sequence of complexes is exact if and only if the sequence
of terms in a given degree is exact. Hence the lemma follows.
\end{proof}

\begin{lemma}
\label{lemma-cech-cohomology-delta-functor-presheaves}
Let $X$ be a ringed space.
Let $\mathcal{U} : U = \bigcup_{i \in I} U_i$ be an open covering.
The functors $\mathcal{F} \mapsto \check{H}^n(\mathcal{U}, \mathcal{F})$
form a $\delta$-functor from the abelian category of
presheaves of $\mathcal{O}_X$-modules to the category
of $\mathcal{O}_X(U)$-modules (see
Homology, Definition \ref{homology-definition-cohomological-delta-functor}).
\end{lemma}

\begin{proof}
By
Lemma \ref{lemma-cech-exact-presheaves}
a short exact sequence of presheaves of
$\mathcal{O}_X$-modules
$0 \to \mathcal{F}_1 \to \mathcal{F}_2 \to \mathcal{F}_3 \to 0$
is turned into a short exact sequence of complexes of
$\mathcal{O}_X(U)$-modules. Hence we can use
Homology, Lemma \ref{homology-lemma-long-exact-sequence-cochain}
to get the boundary maps
$\delta_{\mathcal{F}_1 \to \mathcal{F}_2 \to \mathcal{F}_3} :
\check{H}^n(\mathcal{U}, \mathcal{F}_3) \to
\check{H}^{n + 1}(\mathcal{U}, \mathcal{F}_1)$
and a corresponding long exact sequence. We omit the verification
that these maps are compatible with maps between short exact
sequences of presheaves.
\end{proof}


\noindent
In the formulation of the following lemma we use the functor $j_{p!}$ of
extension by $0$ for presheaves of modules
relative to an open immersion $j : U \to X$.
See Sheaves, Section \ref{sheaves-section-open-immersions}. For any open
$W \subset X$ and any presheaf $\mathcal{G}$ of $\mathcal{O}_X|_U$-modules
we have
$$
(j_{p!}\mathcal{G})(W) =
\left\{
\begin{matrix}
\mathcal{G}(W) & \text{if } W \subset U \\
0 & \text{else.}
\end{matrix}
\right.
$$
Moreover, the functor $j_{p!}$ is a left adjoint to the restriction functor
see Sheaves, Lemma \ref{sheaves-lemma-j-shriek-modules}.
In particular we have the following formula
$$
\Hom_{\mathcal{O}_X}(j_{p!}\mathcal{O}_U, \mathcal{F})
=
\Hom_{\mathcal{O}_U}(\mathcal{O}_U, \mathcal{F}|_U)
=
\mathcal{F}(U).
$$
Since the functor $\mathcal{F} \mapsto \mathcal{F}(U)$ is an exact functor
on the category of presheaves we conclude that the presheaf
$j_{p!}\mathcal{O}_U$ is a projective object in the category
$\textit{PMod}(\mathcal{O}_X)$, see
Homology, Lemma \ref{homology-lemma-characterize-projectives}.

\medskip\noindent
Note that if we are given open subsets $U \subset V \subset X$
with associated open immersions $j_U, j_V$, then we have a canonical
map $(j_U)_{p!}\mathcal{O}_U \to (j_V)_{p!}\mathcal{O}_V$. It is the
identity on sections over any open $W \subset U$ and $0$ else.
In terms of the identification
$\Hom_{\mathcal{O}_X}((j_U)_{p!}\mathcal{O}_U, (j_V)_{p!}\mathcal{O}_V) =
(j_V)_{p!}\mathcal{O}_V(U) = \mathcal{O}_V(U)$ it corresponds to
the element $1 \in \mathcal{O}_V(U)$.

\begin{lemma}
\label{lemma-cech-map-into}
Let $X$ be a ringed space.
Let $\mathcal{U} : U = \bigcup_{i \in I} U_i$ be an open covering.
Denote $j_{i_0\ldots i_p} : U_{i_0 \ldots i_p} \to X$ the open immersion.
Consider the chain complex $K(\mathcal{U})_\bullet$
of presheaves of $\mathcal{O}_X$-modules
$$
\ldots
\to
\bigoplus_{i_0i_1i_2} (j_{i_0i_1i_2})_{p!}\mathcal{O}_{U_{i_0i_1i_2}}
\to
\bigoplus_{i_0i_1} (j_{i_0i_1})_{p!}\mathcal{O}_{U_{i_0i_1}}
\to
\bigoplus_{i_0} (j_{i_0})_{p!}\mathcal{O}_{U_{i_0}}
\to 0 \to \ldots
$$
where the last nonzero term is placed in degree $0$
and where the map
$$
(j_{i_0\ldots i_{p + 1}})_{p!}\mathcal{O}_{U_{i_0\ldots i_{p + 1}}}
\longrightarrow
(j_{i_0\ldots \hat i_j \ldots i_{p + 1}})_{p!}
\mathcal{O}_{U_{i_0\ldots \hat i_j \ldots i_{p + 1}}}
$$
is given by $(-1)^j$ times the canonical map.
Then there is an isomorphism
$$
\Hom_{\mathcal{O}_X}(K(\mathcal{U})_\bullet, \mathcal{F})
=
\check{\mathcal{C}}^\bullet(\mathcal{U}, \mathcal{F})
$$
functorial in $\mathcal{F} \in \Ob(\textit{PMod}(\mathcal{O}_X))$.
\end{lemma}

\begin{proof}
We saw in the discussion just above the lemma that
$$
\Hom_{\mathcal{O}_X}(
(j_{i_0\ldots i_p})_{p!}\mathcal{O}_{U_{i_0\ldots i_p}},
\mathcal{F})
=
\mathcal{F}(U_{i_0\ldots i_p}).
$$
Hence we see that it is indeed the case that the direct sum
$$
\bigoplus\nolimits_{i_0 \ldots i_p}
(j_{i_0 \ldots i_p})_{p!}\mathcal{O}_{U_{i_0 \ldots i_p}}
$$
represents the functor
$$
\mathcal{F}
\longmapsto
\prod\nolimits_{i_0\ldots i_p} \mathcal{F}(U_{i_0\ldots i_p}).
$$
Hence by Categories, Yoneda Lemma \ref{categories-lemma-yoneda}
we see that there is a complex $K(\mathcal{U})_\bullet$ with terms
as given. It is a simple matter to see that the maps are as given
in the lemma.
\end{proof}

\begin{lemma}
\label{lemma-homology-complex}
Let $X$ be a ringed space.
Let $\mathcal{U} : U = \bigcup_{i \in I} U_i$ be an open covering.
Let $\mathcal{O}_\mathcal{U} \subset \mathcal{O}_X$
be the image presheaf of the map
$\bigoplus j_{p!}\mathcal{O}_{U_i} \to \mathcal{O}_X$.
The chain complex $K(\mathcal{U})_\bullet$ of presheaves
of Lemma \ref{lemma-cech-map-into} above has homology presheaves
$$
H_i(K(\mathcal{U})_\bullet) =
\left\{
\begin{matrix}
0 & \text{if} & i \not = 0 \\
\mathcal{O}_\mathcal{U} & \text{if} & i = 0
\end{matrix}
\right.
$$
\end{lemma}

\begin{proof}
Consider the extended complex $K^{ext}_\bullet$ one gets by putting
$\mathcal{O}_\mathcal{U}$ in degree $-1$ with the obvious map
$K(\mathcal{U})_0 =
\bigoplus_{i_0} (j_{i_0})_{p!}\mathcal{O}_{U_{i_0}} \to
\mathcal{O}_\mathcal{U}$.
It suffices to show that taking sections of this extended complex over
any open $W \subset X$ leads to an acyclic complex.
In fact, we claim that for every $W \subset X$ the complex
$K^{ext}_\bullet(W)$ is homotopy equivalent to the zero complex.
Write $I = I_1 \amalg I_2$ where $W \subset U_i$ if and only
if $i \in I_1$.

\medskip\noindent
If $I_1 = \emptyset$, then the complex $K^{ext}_\bullet(W) = 0$ so there is
nothing to prove.

\medskip\noindent
If $I_1 \not = \emptyset$, then
$\mathcal{O}_\mathcal{U}(W) = \mathcal{O}_X(W)$
and
$$
K^{ext}_p(W) =
\bigoplus\nolimits_{i_0 \ldots i_p \in I_1} \mathcal{O}_X(W).
$$
This is true because of the simple description of the presheaves
$(j_{i_0 \ldots i_p})_{p!}\mathcal{O}_{U_{i_0 \ldots i_p}}$.
Moreover, the differential of the complex $K^{ext}_\bullet(W)$
is given by
$$
d(s)_{i_0 \ldots i_p} =
\sum\nolimits_{j = 0, \ldots, p + 1} \sum\nolimits_{i \in I_1}
(-1)^j s_{i_0 \ldots i_{j - 1} i i_j \ldots i_p}.
$$
The sum is finite as the element $s$ has finite support.
Fix an element $i_{\text{fix}} \in I_1$. Define a map
$$
h : K^{ext}_p(W) \longrightarrow K^{ext}_{p + 1}(W)
$$
by the rule
$$
h(s)_{i_0 \ldots i_{p + 1}} =
\left\{
\begin{matrix}
0 & \text{if} & i_0 \not = i_{\text{fix}} \\
s_{i_1 \ldots i_{p + 1}} & \text{if} & i_0 = i_{\text{fix}}
\end{matrix}
\right.
$$
We will use the shorthand
$h(s)_{i_0 \ldots i_{p + 1}} = (i_0 = i_{\text{fix}}) s_{i_1 \ldots i_p}$
for this. Then we compute
\begin{eqnarray*}
& & (dh + hd)(s)_{i_0 \ldots i_p} \\
& = &
\sum_j \sum_{i \in I_1} (-1)^j h(s)_{i_0 \ldots i_{j - 1} i i_j \ldots i_p}
+
(i = i_0) d(s)_{i_1 \ldots i_p} \\
& = &
s_{i_0 \ldots i_p} +
\sum_{j \geq 1}\sum_{i \in I_1}
(-1)^j (i_0 = i_{\text{fix}}) s_{i_1 \ldots i_{j - 1} i i_j \ldots i_p}
+
(i_0 = i_{\text{fix}}) d(s)_{i_1 \ldots i_p}
\end{eqnarray*}
which is equal to $s_{i_0 \ldots i_p}$ as desired.
\end{proof}

\begin{lemma}
\label{lemma-cech-cohomology-derived-presheaves}
Let $X$ be a ringed space.
Let $\mathcal{U} : U = \bigcup_{i \in I} U_i$
be an open covering of $U \subset X$.
The {\v C}ech cohomology functors $\check{H}^p(\mathcal{U}, -)$
are canonically isomorphic as a $\delta$-functor to
the right derived functors of the functor
$$
\check{H}^0(\mathcal{U}, -) :
\textit{PMod}(\mathcal{O}_X)
\longrightarrow
\text{Mod}_{\mathcal{O}_X(U)}.
$$
Moreover, there is a functorial quasi-isomorphism
$$
\check{\mathcal{C}}^\bullet(\mathcal{U}, \mathcal{F})
\longrightarrow
R\check{H}^0(\mathcal{U}, \mathcal{F})
$$
where the right hand side indicates the right derived functor
$$
R\check{H}^0(\mathcal{U}, -) :
D^{+}(\textit{PMod}(\mathcal{O}_X))
\longrightarrow
D^{+}(\mathcal{O}_X(U))
$$
of the left exact functor $\check{H}^0(\mathcal{U}, -)$.
\end{lemma}

\begin{proof}
Note that the category of presheaves of $\mathcal{O}_X$-modules
has enough injectives, see
Injectives, Proposition \ref{injectives-proposition-presheaves-modules}.
Note that $\check{H}^0(\mathcal{U}, -)$ is a left exact functor
from the category of presheaves of $\mathcal{O}_X$-modules
to the category of $\mathcal{O}_X(U)$-modules.
Hence the derived functor and the right derived functor exist, see
Derived Categories, Section \ref{derived-section-right-derived-functor}.

\medskip\noindent
Let $\mathcal{I}$ be an injective presheaf of $\mathcal{O}_X$-modules.
In this case the functor $\Hom_{\mathcal{O}_X}(-, \mathcal{I})$
is exact on $\textit{PMod}(\mathcal{O}_X)$. By
Lemma \ref{lemma-cech-map-into} we have
$$
\Hom_{\mathcal{O}_X}(K(\mathcal{U})_\bullet, \mathcal{I})
=
\check{\mathcal{C}}^\bullet(\mathcal{U}, \mathcal{I}).
$$
By Lemma \ref{lemma-homology-complex} we have that $K(\mathcal{U})_\bullet$ is
quasi-isomorphic to $\mathcal{O}_\mathcal{U}[0]$. Hence by
the exactness of Hom into $\mathcal{I}$ mentioned above we see
that $\check{H}^i(\mathcal{U}, \mathcal{I}) = 0$ for all
$i > 0$. Thus the $\delta$-functor $(\check{H}^n, \delta)$
(see Lemma \ref{lemma-cech-cohomology-delta-functor-presheaves})
satisfies the assumptions of
Homology, Lemma \ref{homology-lemma-efface-implies-universal},
and hence is a universal $\delta$-functor.

\medskip\noindent
By
Derived Categories, Lemma \ref{derived-lemma-higher-derived-functors}
also the sequence $R^i\check{H}^0(\mathcal{U}, -)$
forms a universal $\delta$-functor. By the uniqueness of universal
$\delta$-functors, see
Homology, Lemma \ref{homology-lemma-uniqueness-universal-delta-functor}
we conclude that
$R^i\check{H}^0(\mathcal{U}, -) = \check{H}^i(\mathcal{U}, -)$.
This is enough for most applications
and the reader is suggested to skip the rest of the proof.

\medskip\noindent
Let $\mathcal{F}$ be any presheaf of $\mathcal{O}_X$-modules.
Choose an injective resolution $\mathcal{F} \to \mathcal{I}^\bullet$
in the category $\textit{PMod}(\mathcal{O}_X)$.
Consider the double complex
$\check{\mathcal{C}}^\bullet(\mathcal{U}, \mathcal{I}^\bullet)$ with terms
$\check{\mathcal{C}}^p(\mathcal{U}, \mathcal{I}^q)$.
Consider the associated total complex
$\text{Tot}(\check{\mathcal{C}}^\bullet(\mathcal{U}, \mathcal{I}^\bullet))$,
see Homology, Definition \ref{homology-definition-associated-simple-complex}.
There is a map of complexes
$$
\check{\mathcal{C}}^\bullet(\mathcal{U}, \mathcal{F})
\longrightarrow
\text{Tot}(\check{\mathcal{C}}^\bullet(\mathcal{U}, \mathcal{I}^\bullet))
$$
coming from the maps
$\check{\mathcal{C}}^p(\mathcal{U}, \mathcal{F})
\to \check{\mathcal{C}}^p(\mathcal{U}, \mathcal{I}^0)$
and there is a map of complexes
$$
\check{H}^0(\mathcal{U}, \mathcal{I}^\bullet)
\longrightarrow
\text{Tot}(\check{\mathcal{C}}^\bullet(\mathcal{U}, \mathcal{I}^\bullet))
$$
coming from the maps
$\check{H}^0(\mathcal{U}, \mathcal{I}^q) \to
\check{\mathcal{C}}^0(\mathcal{U}, \mathcal{I}^q)$.
Both of these maps are quasi-isomorphisms by an application of
Homology, Lemma \ref{homology-lemma-double-complex-gives-resolution}.
Namely, the columns of the double complex are exact in positive degrees
because the {\v C}ech complex as a functor is exact
(Lemma \ref{lemma-cech-exact-presheaves})
and the rows of the double complex are exact in positive degrees
since as we just saw the higher {\v C}ech cohomology groups of the injective
presheaves $\mathcal{I}^q$ are zero.
Since quasi-isomorphisms become invertible
in $D^{+}(\mathcal{O}_X(U))$ this gives the last displayed morphism
of the lemma. We omit the verification that this morphism is
functorial.
\end{proof}





\section{{\v C}ech cohomology and cohomology}
\label{section-cech-cohomology-cohomology}

\begin{lemma}
\label{lemma-injective-trivial-cech}
Let $X$ be a ringed space.
Let $\mathcal{U} : U = \bigcup_{i \in I} U_i$ be an open covering.
Let $\mathcal{I}$ be an injective $\mathcal{O}_X$-module.
Then
$$
\check{H}^p(\mathcal{U}, \mathcal{I}) =
\left\{
\begin{matrix}
\mathcal{I}(U) & \text{if} & p = 0 \\
0 & \text{if} & p > 0
\end{matrix}
\right.
$$
\end{lemma}

\begin{proof}
An injective $\mathcal{O}_X$-module is also injective as an object in
the category $\textit{PMod}(\mathcal{O}_X)$ (for example since
sheafification is an exact left adjoint to the inclusion functor,
using Homology, Lemma \ref{homology-lemma-adjoint-preserve-injectives}).
Hence we can apply Lemma \ref{lemma-cech-cohomology-derived-presheaves}
(or its proof) to see the result.
\end{proof}

\begin{lemma}
\label{lemma-cech-cohomology}
Let $X$ be a ringed space.
Let $\mathcal{U} : U = \bigcup_{i \in I} U_i$ be an open covering.
There is a transformation
$$
\check{\mathcal{C}}^\bullet(\mathcal{U}, -)
\longrightarrow
R\Gamma(U, -)
$$
of functors
$\textit{Mod}(\mathcal{O}_X) \to D^{+}(\mathcal{O}_X(U))$.
In particular this provides canonical maps
$\check{H}^p(\mathcal{U}, \mathcal{F}) \to H^p(U, \mathcal{F})$ for
$\mathcal{F}$ ranging over $\textit{Mod}(\mathcal{O}_X)$.
\end{lemma}

\begin{proof}
Let $\mathcal{F}$ be an $\mathcal{O}_X$-module. Choose an injective resolution
$\mathcal{F} \to \mathcal{I}^\bullet$. Consider the double complex
$\check{\mathcal{C}}^\bullet(\mathcal{U}, \mathcal{I}^\bullet)$ with terms
$\check{\mathcal{C}}^p(\mathcal{U}, \mathcal{I}^q)$.
There is a map of complexes
$$
\alpha :
\Gamma(U, \mathcal{I}^\bullet)
\longrightarrow
\text{Tot}(\check{\mathcal{C}}^\bullet(\mathcal{U}, \mathcal{I}^\bullet))
$$
coming from the maps
$\mathcal{I}^q(U) \to \check{H}^0(\mathcal{U}, \mathcal{I}^q)$
and a map of complexes
$$
\beta :
\check{\mathcal{C}}^\bullet(\mathcal{U}, \mathcal{F})
\longrightarrow
\text{Tot}(\check{\mathcal{C}}^\bullet(\mathcal{U}, \mathcal{I}^\bullet))
$$
coming from the map $\mathcal{F} \to \mathcal{I}^0$.
We can apply
Homology, Lemma \ref{homology-lemma-double-complex-gives-resolution}
to see that $\alpha$ is a quasi-isomorphism.
Namely, Lemma \ref{lemma-injective-trivial-cech} implies that
the $q$th row of the double complex
$\check{\mathcal{C}}^\bullet(\mathcal{U}, \mathcal{I}^\bullet)$ is a
resolution of $\Gamma(U, \mathcal{I}^q)$.
Hence $\alpha$ becomes invertible in $D^{+}(\mathcal{O}_X(U))$ and
the transformation of the lemma is the composition of $\beta$
followed by the inverse of $\alpha$. We omit the verification
that this is functorial.
\end{proof}

\begin{lemma}
\label{lemma-cech-h1}
Let $X$ be a topological space. Let $\mathcal{H}$ be an abelian sheaf
on $X$. Let $\mathcal{U} : X = \bigcup_{i \in I} U_i$ be an open covering.
The map
$$
\check{H}^1(\mathcal{U}, \mathcal{H}) \longrightarrow H^1(X, \mathcal{H})
$$
is injective and identifies $\check{H}^1(\mathcal{U}, \mathcal{H})$ via
the bijection of Lemma \ref{lemma-torsors-h1}
with the set of isomorphism classes of $\mathcal{H}$-torsors
which restrict to trivial torsors over each $U_i$.
\end{lemma}

\begin{proof}
To see this we construct an inverse map. Namely, let $\mathcal{F}$ be a
$\mathcal{H}$-torsor whose restriction to $U_i$ is trivial. By
Lemma \ref{lemma-trivial-torsor} this means there
exists a section $s_i \in \mathcal{F}(U_i)$. On $U_{i_0} \cap U_{i_1}$
there is a unique section $s_{i_0i_1}$ of $\mathcal{H}$ such that
$s_{i_0i_1} \cdot s_{i_0}|_{U_{i_0} \cap U_{i_1}} =
s_{i_1}|_{U_{i_0} \cap U_{i_1}}$. A computation shows
that $s_{i_0i_1}$ is a {\v C}ech cocycle and that its class is well
defined (i.e., does not depend on the choice of the sections $s_i$).
The inverse maps the isomorphism class of $\mathcal{F}$ to the cohomology
class of the cocycle $(s_{i_0i_1})$.
We omit the verification that this map is indeed an inverse.
\end{proof}

\begin{lemma}
\label{lemma-include}
Let $X$ be a ringed space.
Consider the functor
$i : \textit{Mod}(\mathcal{O}_X) \to \textit{PMod}(\mathcal{O}_X)$.
It is a left exact functor with right derived functors given by
$$
R^pi(\mathcal{F}) = \underline{H}^p(\mathcal{F}) :
U \longmapsto H^p(U, \mathcal{F})
$$
see discussion in Section \ref{section-locality}.
\end{lemma}

\begin{proof}
It is clear that $i$ is left exact.
Choose an injective resolution $\mathcal{F} \to \mathcal{I}^\bullet$.
By definition $R^pi$ is the $p$th cohomology {\it presheaf}
of the complex $\mathcal{I}^\bullet$. In other words, the
sections of $R^pi(\mathcal{F})$ over an open $U$ are given by
$$
\frac{\Ker(\mathcal{I}^p(U) \to \mathcal{I}^{p + 1}(U))}
{\Im(\mathcal{I}^{p - 1}(U) \to \mathcal{I}^p(U))}.
$$
which is the definition of $H^p(U, \mathcal{F})$.
\end{proof}

\begin{lemma}
\label{lemma-cech-spectral-sequence}
Let $X$ be a ringed space.
Let $\mathcal{U} : U = \bigcup_{i \in I} U_i$ be an open covering.
For any sheaf of $\mathcal{O}_X$-modules $\mathcal{F}$ there
is a spectral sequence $(E_r, d_r)_{r \geq 0}$ with
$$
E_2^{p, q} = \check{H}^p(\mathcal{U}, \underline{H}^q(\mathcal{F}))
$$
converging to $H^{p + q}(U, \mathcal{F})$.
This spectral sequence is functorial in $\mathcal{F}$.
\end{lemma}

\begin{proof}
This is a Grothendieck spectral sequence
(see
Derived Categories, Lemma \ref{derived-lemma-grothendieck-spectral-sequence})
for the functors
$$
i :  \textit{Mod}(\mathcal{O}_X) \to \textit{PMod}(\mathcal{O}_X)
\quad\text{and}\quad
\check{H}^0(\mathcal{U}, - ) : \textit{PMod}(\mathcal{O}_X)
\to \text{Mod}_{\mathcal{O}_X(U)}.
$$
Namely, we have $\check{H}^0(\mathcal{U}, i(\mathcal{F})) = \mathcal{F}(U)$
by Lemma \ref{lemma-cech-h0}. We have that $i(\mathcal{I})$ is
{\v C}ech acyclic by Lemma \ref{lemma-injective-trivial-cech}. And we
have that $\check{H}^p(\mathcal{U}, -) = R^p\check{H}^0(\mathcal{U}, -)$
as functors on $\textit{PMod}(\mathcal{O}_X)$
by Lemma \ref{lemma-cech-cohomology-derived-presheaves}.
Putting everything together gives the lemma.
\end{proof}

\begin{lemma}
\label{lemma-cech-spectral-sequence-application}
Let $X$ be a ringed space.
Let $\mathcal{U} : U = \bigcup_{i \in I} U_i$ be an open covering.
Let $\mathcal{F}$ be an $\mathcal{O}_X$-module.
Assume that $H^i(U_{i_0 \ldots i_p}, \mathcal{F}) = 0$
for all $i > 0$, all $p \geq 0$ and all $i_0, \ldots, i_p \in I$.
Then $\check{H}^p(\mathcal{U}, \mathcal{F}) = H^p(U, \mathcal{F})$
as $\mathcal{O}_X(U)$-modules.
\end{lemma}

\begin{proof}
We will use the spectral sequence of
Lemma \ref{lemma-cech-spectral-sequence}.
The assumptions mean that $E_2^{p, q} = 0$ for all $(p, q)$ with
$q \not = 0$. Hence the spectral sequence degenerates at $E_2$
and the result follows.
\end{proof}

\begin{lemma}
\label{lemma-ses-cech-h1}
\begin{reference}
\cite[Chapter II, \S 3, no. 46, Theorem 3, Corollary 2,
p. 239]{FAC}
\end{reference}
\begin{history}
The cited corollary considers a short exact sequence on an affine classical
variety and uses vanishing of $H^1$ for the coherent kernel to lift global
sections of the quotient. The lemma below isolates the same obstruction at
the Cech level on an arbitrary ringed space: only the kernel requires the
vanishing hypothesis, while the middle and quotient sheaves need not be
quasi-coherent or coherent.
\end{history}
Let $X$ be a ringed space.
Let
$$
0 \to \mathcal{F} \to \mathcal{G} \to \mathcal{H} \to 0
$$
be a short exact sequence of $\mathcal{O}_X$-modules.
Let $U \subset X$ be an open subset.
If there exists a cofinal system of open coverings $\mathcal{U}$
of $U$ such that $\check{H}^1(\mathcal{U}, \mathcal{F}) = 0$,
then the map $\mathcal{G}(U) \to \mathcal{H}(U)$ is
surjective.
\end{lemma}

\begin{proof}
Take an element $s \in \mathcal{H}(U)$. Choose an open covering
$\mathcal{U} : U = \bigcup_{i \in I} U_i$ such that
(a) $\check{H}^1(\mathcal{U}, \mathcal{F}) = 0$ and (b)
$s|_{U_i}$ is the image of a section $s_i \in \mathcal{G}(U_i)$.
Since we can certainly find $\mathcal{U}$ such that (b) holds
it follows from the assumptions of the lemma that we can find
$\mathcal{U}$ such that (a) and (b) both hold.
Consider the sections
$$
s_{i_0i_1} = s_{i_1}|_{U_{i_0i_1}} - s_{i_0}|_{U_{i_0i_1}}.
$$
Since $s_i$ lifts $s$ we see that $s_{i_0i_1} \in \mathcal{F}(U_{i_0i_1})$.
By the vanishing of $\check{H}^1(\mathcal{U}, \mathcal{F})$ we can
find sections $t_i \in \mathcal{F}(U_i)$ such that
$$
s_{i_0i_1} = t_{i_1}|_{U_{i_0i_1}} - t_{i_0}|_{U_{i_0i_1}}.
$$
Then clearly the sections $s_i - t_i$ satisfy the sheaf condition
and glue to a section of $\mathcal{G}$ over $U$ which maps to $s$.
Hence we win.
\end{proof}

\begin{lemma}
\label{lemma-cech-vanish}
\begin{slogan}
If higher {\v C}ech cohomology of an abelian sheaf vanishes for all open covers,
then higher cohomology vanishes.
\end{slogan}
Let $X$ be a ringed space.
Let $\mathcal{F}$ be an $\mathcal{O}_X$-module such that
$$
\check{H}^p(\mathcal{U}, \mathcal{F}) = 0
$$
for all $p > 0$ and any open covering $\mathcal{U} : U = \bigcup_{i \in I} U_i$
of an open of $X$. Then $H^p(U, \mathcal{F}) = 0$ for all $p > 0$
and any open $U \subset X$.
\end{lemma}

\begin{proof}
Let $\mathcal{F}$ be a sheaf satisfying the assumption of the lemma.
We will indicate this by saying ``$\mathcal{F}$ has vanishing higher
{\v C}ech cohomology for any open covering''.
Choose an embedding $\mathcal{F} \to \mathcal{I}$ into an
injective $\mathcal{O}_X$-module.
By Lemma \ref{lemma-injective-trivial-cech} $\mathcal{I}$ has vanishing higher
{\v C}ech cohomology for any open covering.
Let $\mathcal{Q} = \mathcal{I}/\mathcal{F}$
so that we have a short exact sequence
$$
0 \to \mathcal{F} \to \mathcal{I} \to \mathcal{Q} \to 0.
$$
By Lemma \ref{lemma-ses-cech-h1} and our assumptions
this sequence is actually exact as a sequence of presheaves!
In particular we have a long exact sequence of {\v C}ech cohomology
groups for any open covering $\mathcal{U}$, see
Lemma \ref{lemma-cech-cohomology-delta-functor-presheaves}
for example. This implies that $\mathcal{Q}$ is also an $\mathcal{O}_X$-module
with vanishing higher {\v C}ech cohomology for all open coverings.

\medskip\noindent
Next, we look at the long exact cohomology sequence
$$
\xymatrix{
0 \ar[r] &
H^0(U, \mathcal{F}) \ar[r] &
H^0(U, \mathcal{I}) \ar[r] &
H^0(U, \mathcal{Q}) \ar[lld] \\
&
H^1(U, \mathcal{F}) \ar[r] &
H^1(U, \mathcal{I}) \ar[r] &
H^1(U, \mathcal{Q}) \ar[lld] \\
&
\ldots & \ldots & \ldots \\
}
$$
for any open $U \subset X$. Since $\mathcal{I}$ is injective we
have $H^n(U, \mathcal{I}) = 0$ for $n > 0$ (see
Derived Categories, Lemma \ref{derived-lemma-higher-derived-functors}).
By the above we see that $H^0(U, \mathcal{I}) \to H^0(U, \mathcal{Q})$
is surjective and hence $H^1(U, \mathcal{F}) = 0$.
Since $\mathcal{F}$ was an arbitrary $\mathcal{O}_X$-module with
vanishing higher {\v C}ech cohomology we conclude that also
$H^1(U, \mathcal{Q}) = 0$ since $\mathcal{Q}$ is another of these
sheaves (see above). By the long exact sequence this in turn implies
that $H^2(U, \mathcal{F}) = 0$. And so on and so forth.
\end{proof}

\begin{lemma}
\label{lemma-cech-vanish-basis}
(Variant of Lemma \ref{lemma-cech-vanish}.)
Let $X$ be a ringed space.
Let $\mathcal{B}$ be a basis for the topology on $X$.
Let $\mathcal{F}$ be an $\mathcal{O}_X$-module.
Assume there exists a set of open coverings $\text{Cov}$
with the following properties:
\begin{enumerate}
\item For every $\mathcal{U} \in \text{Cov}$
with $\mathcal{U} : U = \bigcup_{i \in I} U_i$ we have
$U, U_i \in \mathcal{B}$ and every $U_{i_0 \ldots i_p} \in \mathcal{B}$.
\item For every $U \in \mathcal{B}$ the open coverings of $U$
occurring in $\text{Cov}$ is a cofinal system of open coverings
of $U$.
\item For every $\mathcal{U} \in \text{Cov}$ we have
$\check{H}^p(\mathcal{U}, \mathcal{F}) = 0$ for all $p > 0$.
\end{enumerate}
Then $H^p(U, \mathcal{F}) = 0$ for all $p > 0$ and any $U \in \mathcal{B}$.
\end{lemma}

\begin{proof}
Let $\mathcal{F}$ and $\text{Cov}$ be as in the lemma.
We will indicate this by saying ``$\mathcal{F}$ has vanishing higher
{\v C}ech cohomology for any $\mathcal{U} \in \text{Cov}$''.
Choose an embedding $\mathcal{F} \to \mathcal{I}$ into an
injective $\mathcal{O}_X$-module.
By Lemma \ref{lemma-injective-trivial-cech} $\mathcal{I}$
has vanishing higher {\v C}ech cohomology for any $\mathcal{U} \in \text{Cov}$.
Let $\mathcal{Q} = \mathcal{I}/\mathcal{F}$
so that we have a short exact sequence
$$
0 \to \mathcal{F} \to \mathcal{I} \to \mathcal{Q} \to 0.
$$
By Lemma \ref{lemma-ses-cech-h1} and our assumption (2)
this sequence gives rise to an exact sequence
$$
0 \to \mathcal{F}(U) \to \mathcal{I}(U) \to \mathcal{Q}(U) \to 0.
$$
for every $U \in \mathcal{B}$. Hence for any $\mathcal{U} \in \text{Cov}$
we get a short exact sequence of {\v C}ech complexes
$$
0 \to
\check{\mathcal{C}}^\bullet(\mathcal{U}, \mathcal{F}) \to
\check{\mathcal{C}}^\bullet(\mathcal{U}, \mathcal{I}) \to
\check{\mathcal{C}}^\bullet(\mathcal{U}, \mathcal{Q}) \to 0
$$
since each term in the {\v C}ech complex is made up out of a product of
values over elements of $\mathcal{B}$ by assumption (1).
In particular we have a long exact sequence of {\v C}ech cohomology
groups for any open covering $\mathcal{U} \in \text{Cov}$.
This implies that $\mathcal{Q}$ is also an $\mathcal{O}_X$-module
with vanishing higher {\v C}ech cohomology for all
$\mathcal{U} \in \text{Cov}$.

\medskip\noindent
Next, we look at the long exact cohomology sequence
$$
\xymatrix{
0 \ar[r] &
H^0(U, \mathcal{F}) \ar[r] &
H^0(U, \mathcal{I}) \ar[r] &
H^0(U, \mathcal{Q}) \ar[lld] \\
&
H^1(U, \mathcal{F}) \ar[r] &
H^1(U, \mathcal{I}) \ar[r] &
H^1(U, \mathcal{Q}) \ar[lld] \\
&
\ldots & \ldots & \ldots \\
}
$$
for any $U \in \mathcal{B}$. Since $\mathcal{I}$ is injective we
have $H^n(U, \mathcal{I}) = 0$ for $n > 0$ (see
Derived Categories, Lemma \ref{derived-lemma-higher-derived-functors}).
By the above we see that $H^0(U, \mathcal{I}) \to H^0(U, \mathcal{Q})$
is surjective and hence $H^1(U, \mathcal{F}) = 0$.
Since $\mathcal{F}$ was an arbitrary $\mathcal{O}_X$-module with
vanishing higher {\v C}ech cohomology for all $\mathcal{U} \in \text{Cov}$
we conclude that also $H^1(U, \mathcal{Q}) = 0$ since $\mathcal{Q}$ is
another of these sheaves (see above). By the long exact sequence this in
turn implies that $H^2(U, \mathcal{F}) = 0$. And so on and so forth.
\end{proof}

\begin{lemma}
\label{lemma-pushforward-injective}
Let $f : X \to Y$ be a morphism of ringed spaces.
Let $\mathcal{I}$ be an injective $\mathcal{O}_X$-module.
Then
\begin{enumerate}
\item $\check{H}^p(\mathcal{V}, f_*\mathcal{I}) = 0$
for all $p > 0$ and any open covering
$\mathcal{V} : V = \bigcup_{j \in J} V_j$ of $Y$.
\item $H^p(V, f_*\mathcal{I}) = 0$ for all $p > 0$ and
every open $V \subset Y$.
\end{enumerate}
In other words, $f_*\mathcal{I}$ is right acyclic for $\Gamma(V, -)$
(see
Derived Categories, Definition \ref{derived-definition-derived-functor})
for any $V \subset Y$ open.
\end{lemma}

\begin{proof}
Set $\mathcal{U} : f^{-1}(V) = \bigcup_{j \in J} f^{-1}(V_j)$.
It is an open covering of $X$ and
$$
\check{\mathcal{C}}^\bullet(\mathcal{V}, f_*\mathcal{I}) =
\check{\mathcal{C}}^\bullet(\mathcal{U}, \mathcal{I}).
$$
This is true because
$$
f_*\mathcal{I}(V_{j_0 \ldots j_p})
= \mathcal{I}(f^{-1}(V_{j_0 \ldots j_p})) =
\mathcal{I}(f^{-1}(V_{j_0}) \cap \ldots \cap f^{-1}(V_{j_p}))
= \mathcal{I}(U_{j_0 \ldots j_p}).
$$
Thus the first statement of the lemma follows from
Lemma \ref{lemma-injective-trivial-cech}. The second statement
follows from the first and Lemma \ref{lemma-cech-vanish}.
\end{proof}

\noindent
The following lemma implies in particular that
$f_* : \textit{Ab}(X) \to \textit{Ab}(Y)$ transforms injective
abelian sheaves into injective abelian sheaves.

\begin{lemma}
\label{lemma-pushforward-injective-flat}
Let $f : X \to Y$ be a morphism of ringed spaces.
Assume $f$ is flat.
Then $f_*\mathcal{I}$ is an injective $\mathcal{O}_Y$-module
for any injective $\mathcal{O}_X$-module $\mathcal{I}$.
\end{lemma}

\begin{proof}
In this case the functor $f^*$ transforms injections into injections
(Modules, Lemma \ref{modules-lemma-pullback-flat}).
Hence the result follows from
Homology, Lemma \ref{homology-lemma-adjoint-preserve-injectives}.
\end{proof}

\begin{lemma}
\label{lemma-cohomology-products}
Let $(X, \mathcal{O}_X)$ be a ringed space. Let $I$ be a set.
For $i \in I$ let  $\mathcal{F}_i$ be an $\mathcal{O}_X$-module.
Let $U \subset X$ be open. The canonical map
$$
H^p(U, \prod\nolimits_{i \in I} \mathcal{F}_i)
\longrightarrow
\prod\nolimits_{i \in I} H^p(U, \mathcal{F}_i)
$$
is an isomorphism for $p = 0$ and injective for $p = 1$.
\end{lemma}

\begin{proof}
The statement for $p = 0$ is true because the product of sheaves
is equal to the product of the underlying presheaves, see
Sheaves, Section \ref{sheaves-section-limits-sheaves}.
Proof for $p = 1$. Set $\mathcal{F} = \prod \mathcal{F}_i$.
Let $\xi \in H^1(U, \mathcal{F})$ map to zero in
$\prod H^1(U, \mathcal{F}_i)$. By locality of cohomology, see
Lemma \ref{lemma-kill-cohomology-class-on-covering},
there exists an open covering $\mathcal{U} : U = \bigcup U_j$ such that
$\xi|_{U_j} = 0$ for all $j$. By Lemma \ref{lemma-cech-h1} this means
$\xi$ comes from an element
$\check \xi \in \check H^1(\mathcal{U}, \mathcal{F})$.
Since the maps
$\check H^1(\mathcal{U}, \mathcal{F}_i) \to H^1(U, \mathcal{F}_i)$
are injective for all $i$ (by Lemma \ref{lemma-cech-h1}), and since
the image of $\xi$ is zero in $\prod H^1(U, \mathcal{F}_i)$ we see
that the image
$\check \xi_i = 0$ in $\check H^1(\mathcal{U}, \mathcal{F}_i)$.
However, since $\mathcal{F} = \prod \mathcal{F}_i$ we see that
$\check{\mathcal{C}}^\bullet(\mathcal{U}, \mathcal{F})$ is the
product of the complexes
$\check{\mathcal{C}}^\bullet(\mathcal{U}, \mathcal{F}_i)$,
hence by
Homology, Lemma \ref{homology-lemma-product-abelian-groups-exact}
we conclude that $\check \xi = 0$ as desired.
\end{proof}







\section{Flasque sheaves}
\label{section-flasque}

\noindent
Here is the definition.

\begin{definition}
\label{definition-flasque}
Let $X$ be a topological space. We say a presheaf of sets
$\mathcal{F}$ is {\it flasque} or {\it flabby} if for every
$U \subset V$ open in $X$ the restriction map
$\mathcal{F}(V) \to \mathcal{F}(U)$ is surjective.
\end{definition}

\begin{lemma}[Pointwise flasque envelope]
\label{lemma-pointwise-flasque-envelope-groups}
\begin{reference}
\cite[Chapter III, \S 3.4, proof of Proposition 3.4.2, p. 162]{Tohoku}
\end{reference}
\begin{history}
For an arbitrary sheaf of sets the pointwise product below is a sheaf and
the germ map is a monomorphism, but the product need not be flasque when
some of the factors are empty. Thus the flasqueness assertion in the cited
source is stated here in the group-valued cases in which it is used. The
factors $\mathcal F_x$ are the stalks of $\mathcal F$; they are not asserted
to be the stalks of $\Pi(\mathcal F)$.
\end{history}
Let $X$ be a topological space and let $\mathcal F$ be a sheaf of sets on
$X$. The rule
$$
\Pi(\mathcal F)(U)=\prod_{x\in U}\mathcal F_x
$$
defines a sheaf of sets, functorially in $\mathcal F$, and the germ maps
define a canonical monomorphism
$$
\mathcal F\longrightarrow\Pi(\mathcal F).
$$
If $\mathcal F$ is a sheaf of possibly noncommutative groups, then
$\Pi(\mathcal F)$ is canonically a sheaf of groups, the displayed
monomorphism is a homomorphism, and $\Pi(\mathcal F)$ is flasque. The same
statements hold in the category of abelian sheaves.
\end{lemma}

\begin{proof}
The sheaf assertion is Sheaves, Example
\ref{sheaves-example-sheaf-product-pointwise}, and injectivity of the germ
map is Sheaves, Lemma \ref{sheaves-lemma-sheaf-subset-stalks}.
Functoriality follows from functoriality of stalks. If $\mathcal F$ is a
sheaf of groups, then the group structure follows from Sheaves, Example
\ref{sheaves-example-sheaf-product-pointwise-algebraic-structure}. For
$U\subset V$, the restriction
$$
\prod_{x\in V}\mathcal F_x\longrightarrow
\prod_{x\in U}\mathcal F_x
$$
is surjective: extend a tuple by putting the identity in every factor
indexed by $V\setminus U$. Hence $\Pi(\mathcal F)$ is flasque. This
construction preserves commutativity.
\end{proof}

\begin{lemma}
\label{lemma-flasque-group-torsors-trivial}
\begin{reference}
\cite[Chapter III, \S 3.4, proof of Proposition 3.4.2, p. 162]{Tohoku}
\end{reference}
Let $X$ be a topological space and let $\mathcal G$ be a sheaf of possibly
noncommutative groups on $X$ whose underlying sheaf of sets is flasque.
Every $\mathcal G$-torsor is trivial. Consequently the pointed set of
isomorphism classes of $\mathcal G$-torsors consists only of its neutral
element. If $\mathcal G$ is abelian, then
$$
H^1(X,\mathcal G)=0.
$$
\end{lemma}

\begin{proof}
Let $\mathcal T$ be a $\mathcal G$-torsor. Consider the set of pairs
$(V,s)$, where $V\subset X$ is open and $s\in\mathcal T(V)$, ordered by
extension. This set is nonempty, since a sheaf of sets has a unique section
over the empty open. Sections in a chain glue over the union of their
domains, so Zorn's lemma gives a maximal pair $(V,s)$.

Choose $x\in X$. Since $\mathcal T_x$ is nonempty, there are an open
neighbourhood $W$ of $x$ and a section $t\in\mathcal T(W)$. Simple
transitivity on $V\cap W$ gives a unique
$$
g\in\mathcal G(V\cap W)
\quad\text{such that}\quad
s|_{V\cap W}=g\cdot t|_{V\cap W}.
$$
By flasqueness, $g$ lifts to a section $\widetilde g\in\mathcal G(W)$.
The sections $s$ and $\widetilde g\cdot t$ agree on $V\cap W$, hence glue
to a section over $V\cup W$. Maximality gives $x\in V$. Since $x$ was
arbitrary, $V=X$. Thus $\mathcal T(X)$ is nonempty, and $\mathcal T$ is
trivial by Lemma \ref{lemma-trivial-torsor}. In the abelian case the final
assertion follows from Lemma \ref{lemma-torsors-h1}.
\end{proof}

\begin{lemma}[The crossed quotient of a pointwise envelope]
\label{lemma-pointwise-envelope-crossed-homomorphism}
\begin{reference}
\cite[Chapter III, \S 3.4, proof of Proposition 3.4.2,
pp. 162--163]{Tohoku}
\end{reference}
Let $X$ be a topological space and let
$\mathcal F \to \mathcal G$ be a monomorphism of sheaves of groups.
Assume that $\mathcal F$ is abelian and normal in $\mathcal G$. Put
$$
\mathcal A=\Pi(\mathcal F),\qquad
\mathcal B=\Pi(\mathcal G),
$$
and identify $\mathcal F$, $\mathcal G$, and $\mathcal A$ with their
images in $\mathcal B$. Let $\mathcal P\subset\mathcal B$ be the
subsheaf whose sections are locally of the form $ag$, with
$a\in\mathcal A$ and $g\in\mathcal G$, and let
$$
\mathcal F'=\mathcal A/\mathcal F
$$
be the quotient sheaf of abelian groups. Then the following hold.
\begin{enumerate}
\item We have $\mathcal A\cap\mathcal G=\mathcal F$. Conjugation in
$\mathcal B$ makes $\mathcal P$ act by group automorphisms on
$\mathcal A$, preserves $\mathcal F$, and hence induces an action
$\sigma$ of $\mathcal P$ on $\mathcal F'$.
\item There is a unique morphism of sheaves of sets
$$
u:\mathcal P\longrightarrow\mathcal F'
$$
characterized locally by
$$
u(ag)=a\bmod\mathcal F.
$$
It is an epimorphism and $u^{-1}(1)=\mathcal G$.
\item For local sections $p,p'$ over the same open we have
$$
u(1)=1,\qquad
u(pp')=u(p)\sigma_p(u(p')).
$$
Moreover,
$$
u(p)=u(p')\quad\Longleftrightarrow\quad
p^{-1}p'\in\mathcal G.
$$
Thus every nonempty fibre is an orbit for right multiplication by
$\mathcal G$.
\end{enumerate}
Here the epimorphism in (2) is an epimorphism of sheaves of sets; it need
not be surjective on sections over every open. In general $u$ is not a
homomorphism of sheaves of groups.
\end{lemma}

\begin{proof}
The stalk maps $\mathcal F_x\to\mathcal G_x$ are injective, so
functoriality of the pointwise envelope gives a monomorphism
$\mathcal A\to\mathcal B$. Normality of
$\mathcal F_x\subset\mathcal G_x$ for every $x$ shows that
$\mathcal A$ is normal in $\mathcal B$. Hence the local products defining
$\mathcal P$ form a subsheaf of groups.

Suppose that $g\in\mathcal G(U)$ belongs to $\mathcal A(U)$. Then
$g_x\in\mathcal F_x$ for every $x\in U$. Near each $x$, lift this germ to
a section of $\mathcal F$ and shrink until its image equals the restriction
of $g$. These local lifts agree because $\mathcal F\to\mathcal G$ is a
monomorphism, and hence they glue. Thus $g\in\mathcal F(U)$, which proves
$$
\mathcal A\cap\mathcal G=\mathcal F.
$$
The sheaf $\mathcal A$ is abelian, so it centralizes $\mathcal F$, while
$\mathcal G$ normalizes $\mathcal F$ by hypothesis. It follows that
$\mathcal P$ preserves both $\mathcal A$ and $\mathcal F$ under
conjugation, giving the asserted action on $\mathcal F'$.

If $ag=a'g'$ locally, then
$$
a^{-1}a'=gg'^{-1}\in\mathcal A\cap\mathcal G=\mathcal F.
$$
Thus $a$ and $a'$ have the same class in $\mathcal F'$. The local rules
defining $u$ therefore agree and glue. Its restriction to $\mathcal A$ is
the quotient map $\mathcal A\to\mathcal F'$, so $u$ is an epimorphism. Its
neutral fibre contains $\mathcal G$. Conversely, if locally $p=ag$ and
$u(p)=1$, then locally $a\in\mathcal F$, and hence locally
$p\in\mathcal G$. Since membership in a subsheaf is local, this proves
$u^{-1}(1)=\mathcal G$.

Finally, locally write $p=ag$ and $p'=bh$. Then
$$
pp'=a(gbg^{-1})gh.
$$
Conjugation by $p=ag$ on $\mathcal F'$ agrees with conjugation by $g$,
because $\mathcal A$ is abelian. Therefore
$$
u(pp')=\overline a\,\overline{gbg^{-1}}
=u(p)\sigma_p(u(p')).
$$
The remaining assertions follow from
Lemma \ref{lemma-crossed-homomorphism-affine-action}, applied using
$u^{-1}(1)=\mathcal G$.
\end{proof}

\begin{lemma}[Pointwise envelopes lift torsors uniquely]
\label{lemma-pointwise-envelope-torsor-bijection}
\begin{reference}
\cite[Chapter III, \S 3.4, proof of Proposition 3.4.2,
pp. 162--163]{Tohoku}
\end{reference}
\begin{history}
Flasqueness is not a local property in general. Thus we do not use the
assertion in the cited proof that a sheaf locally isomorphic to a flasque
sheaf is flasque. Instead, the twisted kernel below is canonically another
pointwise envelope and is therefore flasque. We state the conclusion for
isomorphism classes of torsors; no group structure is asserted when the
sheaves of groups are noncommutative.
\end{history}
Let
$$
1\longrightarrow\mathcal F\longrightarrow\mathcal G
\longrightarrow\mathcal H\longrightarrow 1
$$
be an exact sequence of sheaves of groups on a topological space $X$, and
assume that $\mathcal F$ is abelian. Put
$$
\mathcal A=\Pi(\mathcal F),\qquad
\mathcal B=\Pi(\mathcal G),\qquad
\mathcal C=\Pi(\mathcal H),
$$
and let $\mathcal P=\mathcal A\mathcal G\subset\mathcal B$ be the
local-product subsheaf of
Lemma \ref{lemma-pointwise-envelope-crossed-homomorphism}. Then the
following hold.
\begin{enumerate}
\item The map $\mathcal B\to\mathcal C$ is surjective on sections over
every open, with kernel $\mathcal A$, and
$$
\mathcal P=\mathcal B\times_{\mathcal C}\mathcal H.
$$
In particular,
$$
1\longrightarrow\mathcal A\longrightarrow\mathcal P
\longrightarrow\mathcal H\longrightarrow 1
$$
is exact, and $\mathcal P(U)\to\mathcal H(U)$ is surjective for every
open $U\subset X$.
\item Conjugation gives an action of $\mathcal H$ on $\mathcal F$, and
the conjugation action of $\mathcal P$ on $\mathcal A$ factors through
$\mathcal H$. If $\mathcal E$ is an $\mathcal H$-torsor and superscripts
$\mathcal E$ denote the associated twists, then there is a canonical
isomorphism
$$
\mathcal A^{\mathcal E}\longrightarrow
\Pi(\mathcal F^{\mathcal E}).
$$
Consequently $\mathcal A^{\mathcal E}$ is flasque.
\item Extension of structure group induces a bijection of pointed sets
$$
\left\{
\begin{matrix}
\mathcal P\text{-torsors}\\
\text{up to isomorphism}
\end{matrix}
\right\}
\longrightarrow
\left\{
\begin{matrix}
\mathcal H\text{-torsors}\\
\text{up to isomorphism}.
\end{matrix}
\right\}
$$
Thus every $\mathcal H$-torsor is induced from a $\mathcal P$-torsor,
and the isomorphism class of such a lift is unique.
\end{enumerate}
\end{lemma}

\begin{proof}
Exactness on stalks gives surjections
$\mathcal G_x\to\mathcal H_x$ with kernels $\mathcal F_x$. Hence, for
every open $U$, the map
$$
\mathcal B(U)=\prod_{x\in U}\mathcal G_x
\longrightarrow
\prod_{x\in U}\mathcal H_x=\mathcal C(U)
$$
is surjective and has kernel $\mathcal A(U)$.

View $\mathcal H$ as a subsheaf of $\mathcal C$ by the canonical germ
map. If $b\in\mathcal B(U)$ has image in $\mathcal H(U)$, then locally
that image lifts to a section $g$ of $\mathcal G$. On such an open,
$bg^{-1}$ lies in $\mathcal A$, and hence $b$ lies in $\mathcal P$.
The converse is immediate from the definition of $\mathcal P$. This
proves the fibre-product identity. The remaining assertions in (1) follow
from the sectionwise surjectivity of $\mathcal B\to\mathcal C$.

Conjugation by a local lift in $\mathcal G$ defines an automorphism of
$\mathcal F$. Two lifts differ by a section of $\mathcal F$, whose inner
action on $\mathcal F$ is trivial because $\mathcal F$ is abelian. Thus
the action factors through $\mathcal H$. Similarly, $\mathcal A$ is
abelian, so conjugation by $\mathcal P$ on $\mathcal A$ factors through
$\mathcal P/\mathcal A=\mathcal H$, and it acts coordinatewise on the
stalk factors of $\mathcal A$.

Locally choose a section of the $\mathcal H$-torsor $\mathcal E$. This
identifies the two twists in (2) with $\mathcal A$ and
$\Pi(\mathcal F)$, respectively. On an overlap the two identifications
change by the same coordinatewise conjugation by a section of
$\mathcal H$. The local identifications therefore glue to the asserted
canonical isomorphism. It is flasque by
Lemma \ref{lemma-pointwise-flasque-envelope-groups}.

We prove surjectivity in (3). Let $\mathcal E$ be an
$\mathcal H$-torsor and extend its structure group along
$\mathcal H\to\mathcal C$. The resulting $\mathcal C$-torsor is trivial
by Lemma \ref{lemma-flasque-group-torsors-trivial}. Choose a global
section of it, an open covering $(U_i)$, and sections
$e_i\in\mathcal E(U_i)$. Write
$$
e_i=h_{ij}e_j
$$
on $U_i\cap U_j$. If $\overline e_i$ denotes the induced local frame of
the $\mathcal C$-torsor, then
$\overline e_i=h_{ij}\overline e_j$. Thus the chosen global section has
unique expressions
$$
c_i\overline e_i,\qquad c_i\in\mathcal C(U_i),
$$
and these satisfy
$$
h_{ij}=c_i^{-1}c_j.
$$
Choose lifts $b_i\in\mathcal B(U_i)$ of $c_i$ and put
$$
p_{ij}=b_i^{-1}b_j.
$$
Then $(p_{ij})$ is a cocycle, its image is $(h_{ij})$, and the
fibre-product identity shows that $p_{ij}$ is a section of $\mathcal P$.
Gluing trivial $\mathcal P$-torsors with this cocycle gives a
$\mathcal P$-torsor whose extension to $\mathcal H$ is isomorphic to
$\mathcal E$.

Finally, let $\mathcal Q$ and $\mathcal Q'$ be $\mathcal P$-torsors and
fix an isomorphism between their induced $\mathcal H$-torsors
$\mathcal E$. Let $\mathcal I$ be the sheaf of
$\mathcal P$-equivariant isomorphisms
$\mathcal Q\to\mathcal Q'$ inducing this fixed isomorphism. It is locally
nonempty: after trivializing the two torsors, the required local element
of $\mathcal H$ lifts to $\mathcal P$ by (1). A change between two such
local isomorphisms is an element of $\mathcal A$, and changing a local
frame conjugates that element by the corresponding section of
$\mathcal H$. Hence the usual local calculation identifies $\mathcal I$
as an $\mathcal A^{\mathcal E}$-torsor. This sheaf of groups is flasque
by (2), so Lemma \ref{lemma-flasque-group-torsors-trivial} gives a global
section of $\mathcal I$. The two $\mathcal P$-torsors are therefore
isomorphic. This proves injectivity.
\end{proof}

\noindent
We will use this terminology also for abelian sheaves and
sheaves of modules if $X$ is a ringed space.
Clearly it suffices to assume the restriction maps
$\mathcal{F}(X) \to \mathcal{F}(U)$ is surjective for every
open $U \subset X$.

\begin{lemma}
\label{lemma-injective-flasque}
Let $(X, \mathcal{O}_X)$ be a ringed space.
Then any injective $\mathcal{O}_X$-module is flasque.
\end{lemma}

\begin{proof}
This is a reformulation of Lemma \ref{lemma-injective-restriction-surjective}.
\end{proof}

\begin{lemma}
\label{lemma-flasque-acyclic}
Let $(X, \mathcal{O}_X)$ be a ringed space. Any flasque $\mathcal{O}_X$-module
is acyclic for $R\Gamma(X, -)$ as well as $R\Gamma(U, -)$ for any
open $U$ of $X$.
\end{lemma}

\begin{proof}
We will prove this using
Derived Categories, Lemma \ref{derived-lemma-subcategory-right-acyclics}.
Since every injective module is flasque we see that we can embed
every $\mathcal{O}_X$-module into a flasque module, see
Injectives, Lemma \ref{injectives-lemma-abelian-sheaves-space}.
Thus it suffices to show that given a short exact sequence
$$
0 \to \mathcal{F} \to \mathcal{G} \to \mathcal{H} \to 0
$$
with $\mathcal{F}$, $\mathcal{G}$ flasque, then $\mathcal{H}$
is flasque and the sequence remains short exact after taking sections
on any open of $X$. In fact, the second statement implies the first.
Thus, let $U \subset X$ be an open subspace. Let $s \in \mathcal{H}(U)$.
We will show that we can lift $s$ to a section of $\mathcal{G}$
over $U$. To do this consider the set $T$ of pairs $(V, t)$
where $V \subset U$ is open and $t \in \mathcal{G}(V)$ is a section
mapping to $s|_V$ in $\mathcal{H}$.
We put a partial ordering on $T$ by setting
$(V, t) \leq (V', t')$ if and only if $V \subset V'$ and $t'|_V = t$.
If $(V_\alpha, t_\alpha)$, $\alpha \in A$
is a totally ordered subset of $T$, then $V = \bigcup V_\alpha$
is open and there is a unique section $t \in \mathcal{G}(V)$
restricting to $t_\alpha$ over $V_\alpha$ by the sheaf condition on
$\mathcal{G}$. Thus by Zorn's lemma there exists a maximal element
$(V, t)$ in $T$. We will show that $V = U$ thereby finishing the proof.
Namely, pick any $x \in U$. We can find a small open neighbourhood
$W \subset U$ of $x$ and $t' \in \mathcal{G}(W)$ mapping to $s|_W$
in $\mathcal{H}$. Then $t'|_{W \cap V} - t|_{W \cap V}$ maps to
zero in $\mathcal{H}$, hence comes from some section
$r' \in \mathcal{F}(W \cap V)$. Using that $\mathcal{F}$ is flasque
we find a section $r \in \mathcal{F}(W)$ restricting to $r'$
over $W \cap V$. Modifying $t'$ by the image of $r$ we may
assume that $t$ and $t'$ restrict to the same section over
$W \cap V$. By the sheaf condition of $\mathcal{G}$
we can find a section $\tilde t$ of $\mathcal{G}$ over
$W \cup V$ restricting to $t$ and $t'$.
By maximality of $(V, t)$ we see that $V \cup W = V$.
Thus $x \in V$ and we are done.
\end{proof}

\begin{lemma}[The degree-two boundary of a pointwise envelope]
\label{lemma-pointwise-envelope-degree-two-boundary}
\begin{reference}
\cite[Chapter III, \S 3.4, proof of Proposition 3.4.2 and Remark 1,
pp. 164--165]{Tohoku}
\end{reference}
\begin{history}
The cited proof says that a sheaf locally isomorphic to a flasque sheaf is
flasque, which is false in general. We avoid that assertion by identifying
the twisted middle term below with a pointwise envelope. There is also a
sign convention to translate: with the left-torsor convention of
Lemma \ref{lemma-torsors-h1}, the class agreeing with the usual abelian
connecting map is $+\delta(d_u)$, whereas the cited source records transition
functions in the inverse convention and writes the same comparison with a
minus sign. The final paragraph of the proof gives the explicit calculation.
No compatibility with arbitrary inverse image functors is asserted here.
\end{history}
Let
$$
1\longrightarrow\mathcal F\longrightarrow\mathcal G
\longrightarrow\mathcal H\longrightarrow 1
$$
be an exact sequence of sheaves of groups on a topological space $X$, and
assume that $\mathcal F$ is abelian. Put
$$
\mathcal A=\Pi(\mathcal F),\qquad
\mathcal P=\mathcal A\mathcal G,\qquad
\mathcal F'=\mathcal A/\mathcal F,
$$
and let $u:\mathcal P\to\mathcal F'$ be the crossed epimorphism of
Lemma \ref{lemma-pointwise-envelope-crossed-homomorphism}. Let $\mathcal E$
be an $\mathcal H$-torsor. Choose a $\mathcal P$-torsor $\mathcal Q$ and an
isomorphism $\alpha$ from the $\mathcal H$-torsor induced by $\mathcal Q$ to
$\mathcal E$, as provided by
Lemma \ref{lemma-pointwise-envelope-torsor-bijection}. Then the following
hold.
\begin{enumerate}
\item Twisting the conjugation-equivariant sequence
$$
0\longrightarrow\mathcal F\longrightarrow\mathcal A
\longrightarrow\mathcal F'\longrightarrow 0
$$
by $\mathcal Q$ gives an exact sequence of abelian sheaves
$$
0\longrightarrow\mathcal F^{\mathcal Q}
\longrightarrow\mathcal A^{\mathcal Q}
\longrightarrow(\mathcal F')^{\mathcal Q}\longrightarrow 0.
$$
The isomorphism $\alpha$ induces canonical identifications
$$
\mathcal F^{\mathcal Q}\cong\mathcal F^{\mathcal E},
\qquad
\mathcal A^{\mathcal Q}\cong\mathcal A^{\mathcal E}
\cong\Pi(\mathcal F^{\mathcal E}).
$$
In particular, $\mathcal A^{\mathcal Q}$ is flasque.
\item The connecting map, followed by the identification in (1), is an
isomorphism
$$
\delta_{\mathcal Q,\alpha}:
H^1\bigl(X,(\mathcal F')^{\mathcal Q}\bigr)
\longrightarrow H^2(X,\mathcal F^{\mathcal E}).
$$
The element
$$
\partial(\mathcal E)=
\delta_{\mathcal Q,\alpha}\bigl(d_u(\mathcal Q)\bigr)
\in H^2(X,\mathcal F^{\mathcal E})
$$
is independent of the choice of $(\mathcal Q,\alpha)$.
\item We have
$$
\partial(\mathcal E)=0
\quad\Longleftrightarrow\quad
[\mathcal E]\ \text{is induced by a }\mathcal G\text{-torsor}.
$$
The construction is functorial under isomorphisms of
$\mathcal H$-torsors, with the corresponding transport of the twisted
coefficient sheaves.
\item If $\mathcal G$ is abelian, then $\mathcal F$, $\mathcal G$, and
$\mathcal H$ are abelian, the twists are untwisted, and
$\partial(\mathcal E)$ is the image of $[\mathcal E]$ under the usual
connecting homomorphism
$$
H^1(X,\mathcal H)\longrightarrow H^2(X,\mathcal F).
$$
\end{enumerate}
\end{lemma}

\begin{proof}
The action of $\mathcal P$ on $\mathcal A$ is conjugation inside
$\Pi(\mathcal G)$. The kernel $\mathcal A$ of
$\mathcal P\to\mathcal H$ acts trivially on $\mathcal A$, and hence on
$\mathcal F$ and $\mathcal F'$, because $\mathcal A$ is abelian. Thus all
three actions factor through $\mathcal H$, and the maps in the first
displayed short exact sequence are $\mathcal P$-equivariant. After choosing
a local section of $\mathcal Q$, its twist is identified with the original
sequence. Exactness is local, so the twisted sequence is exact.

Since the actions factor through $\mathcal H$, the isomorphism $\alpha$
identifies the $\mathcal Q$-twists of $\mathcal F$ and $\mathcal A$ with
their $\mathcal E$-twists. The last isomorphism in (1), and its flasqueness,
are Lemma \ref{lemma-pointwise-envelope-torsor-bijection}. Regard these
abelian sheaves as modules on the ringed space
$(X,\underline{\mathbf Z})$. Lemma \ref{lemma-flasque-acyclic} gives
$$
H^i(X,\mathcal A^{\mathcal Q})=0\quad\text{for all }i>0.
$$
The long exact cohomology sequence therefore makes
$\delta_{\mathcal Q,\alpha}$ an isomorphism.

We next prove independence. Suppose that $(\mathcal Q',\alpha')$ is a
second choice. The last paragraph of the proof of
Lemma \ref{lemma-pointwise-envelope-torsor-bijection}, applied to the fixed
isomorphism $(\alpha')^{-1}\alpha$ of induced $\mathcal H$-torsors, gives a
$\mathcal P$-equivariant isomorphism
$\varphi:\mathcal Q\to\mathcal Q'$ inducing that isomorphism. It identifies
the two twisted short exact sequences and the two affine torsors defining
$d_u$. Hence it carries $d_u(\mathcal Q)$ to $d_u(\mathcal Q')$, and
naturality of connecting maps identifies the resulting degree-two classes.
Any two possible $\varphi$ differ locally by a section of the twisted
kernel $\mathcal A^{\mathcal E}$. Conjugation by $\mathcal A$ acts trivially
on $\mathcal F$, $\mathcal A$, and $\mathcal F'$, so the induced maps on
the three coefficient twists are the same. This also proves functoriality
under isomorphisms of $\mathcal H$-torsors.

The sheaf $\mathcal F'$ is abelian. Thus
Lemma \ref{lemma-crossed-homomorphism-torsor-obstruction} gives
$$
d_u(\mathcal Q)\in H^1\bigl(X,(\mathcal F')^{\mathcal Q}\bigr),
$$
and, because $u$ is an epimorphism, this class vanishes if and only if
$\mathcal Q$ admits a reduction of structure group to $\mathcal G$.
Since $\delta_{\mathcal Q,\alpha}$ is an isomorphism, the same is true of
$\partial(\mathcal E)$. Such a reduction induces $\mathcal E$ from a
$\mathcal G$-torsor. Conversely, a $\mathcal G$-torsor inducing
$\mathcal E$ induces a $\mathcal P$-torsor with a $\mathcal G$-reduction;
the injectivity assertion of
Lemma \ref{lemma-pointwise-envelope-torsor-bijection} identifies that
$\mathcal P$-torsor with $\mathcal Q$. This proves (3).

It remains to check the sign in (4). Write the abelian groups additively and
put
$$
\mathcal B=\Pi(\mathcal G),\qquad
\mathcal C=\Pi(\mathcal H),\qquad
\mathcal H'=\mathcal C/\mathcal H.
$$
There is a $3$ by $3$ diagram with exact rows and columns
$$
\xymatrix{
0 \ar[r] & \mathcal F \ar[r] \ar[d] &
\mathcal G \ar[r] \ar[d] & \mathcal H \ar[r] \ar[d] & 0 \\
0 \ar[r] & \mathcal A \ar[r] \ar[d] &
\mathcal B \ar[r] \ar[d] & \mathcal C \ar[r] \ar[d] & 0 \\
0 \ar[r] & \mathcal F' \ar[r] &
\mathcal B/\mathcal G \ar[r] & \mathcal H' \ar[r] & 0.
}
$$
Extend $\mathcal E$ to a $\mathcal C$-torsor. It is trivial by
Lemma \ref{lemma-flasque-group-torsors-trivial}. Choose local sections
$e_i$ of $\mathcal E$ and write
$$
e_i=h_{ij}+e_j.
$$
If a chosen global section of the induced $\mathcal C$-torsor is written
$c_i+e_i$, then $h_{ij}=c_j-c_i$. Choose lifts
$b_i\in\mathcal B(U_i)$ of $c_i$ and put
$$
p_{ij}=b_j-b_i.
$$
By the independence already proved, we may take $\mathcal Q$ to be the
lifting $\mathcal P$-torsor constructed from these transition functions. If
$q_i$ are its corresponding local sections and $r_i=[q_i,0]$ are the local
sections of the affine torsor defining $d_u(\mathcal Q)$, then the orbit
relation gives
$$
r_j=[-p_{ij}+q_i,0]
=[q_i,\rho_{p_{ij}}(0)]
=[q_i,u(p_{ij})].
$$
Consequently Lemma \ref{lemma-cech-h1} represents
$d_u(\mathcal Q)$ by the cocycle $(u(p_{ij}))$. On the other hand, that
lemma represents $[\mathcal E]$ by $(-h_{ij})$, because its convention is
that the transition $a_{ij}$ satisfies $a_{ij}+e_i=e_j$.

The classes of the $c_i$ glue to a section
$t\in H^0(X,\mathcal H')$. The right-hand connecting map in the diagram
sends local lifts $(c_i)$ to the cocycle $(c_i-c_j)=(-h_{ij})$, hence to
$[\mathcal E]$. Similarly, the bottom connecting map sends $t$ to the class
of $(-u(p_{ij}))=-d_u(\mathcal Q)$. The square of connecting maps is
anticommutative by Derived Categories, Proposition
\ref{derived-proposition-9}. It follows that
$$
-\delta_{\mathcal F,\mathcal A,\mathcal F'}
\bigl(d_u(\mathcal Q)\bigr)
=-\delta_{\mathcal F,\mathcal G,\mathcal H}([\mathcal E]),
$$
and hence
$$
\delta_{\mathcal F,\mathcal A,\mathcal F'}
\bigl(d_u(\mathcal Q)\bigr)
=\delta_{\mathcal F,\mathcal G,\mathcal H}([\mathcal E]).
$$
This proves (4). The cited source records $\mathcal E$ by $(h_{ij})$ rather
than by the inverse cocycle used by Lemma \ref{lemma-cech-h1}; in that
convention the last comparison is written $\partial(E)=-d(Q)$.
\end{proof}

\begin{lemma}[The holomorphic logarithmic-derivative obstruction]
\label{lemma-holomorphic-logarithmic-derivative-obstruction}
\begin{reference}
\cite[Chapter III, \S 3.4, Remark 3, pp. 165--166]{Tohoku}
\end{reference}
\begin{history}
The affine formula below uses the left-torsor convention of
Definition \ref{definition-torsor} and the right logarithmic derivative
$(dg)g^{-1}$. Thus neither an inverse adjoint action nor an extra minus sign
is present. The cited source credits Weil with first introducing
$d(\mathcal E)$ and Atiyah with its geometric characterization by holomorphic
connections. Here that characterization is proved directly from the local
connection-form transformation law; no comparison with a derived Atiyah
class is asserted. Surjectivity of the logarithmic derivative is
asserted only as an epimorphism of sheaves on a complex curve, not on global
sections or in higher dimension.
\end{history}
Let $X$ be a complex manifold and let $G$ be a complex Lie group with Lie
algebra $V$. Let $\mathcal P$ be the sheaf of holomorphic maps to $G$, and
let
$$
\mathcal A=\Omega_X^1\otimes_{\mathbf C}V
$$
be the additive sheaf of holomorphic $V$-valued $1$-forms. For a local
section $g$ of $\mathcal P$, set
$$
\sigma_g(\alpha)=\operatorname{Ad}_g(\alpha),
\qquad
u(g)=(dg)g^{-1}.
$$
Then the following hold.
\begin{enumerate}
\item We have
$$
u(1)=0,
\qquad
u(gg')=u(g)+\sigma_g\bigl(u(g')\bigr).
$$
Thus $u$ is a crossed homomorphism in the additive form of
Lemma \ref{lemma-crossed-homomorphism-affine-action}, and its affine action
is
$$
\rho_g(\alpha)=(dg)g^{-1}+\operatorname{Ad}_g(\alpha).
$$
The kernel of $u$ is the constant sheaf $\underline G_X$, whose sections
are the locally constant $G$-valued maps.
\item Let $\mathcal E$ be a $\mathcal P$-torsor. Define
$\operatorname{ad}(\mathcal E)$ using the same left-torsor twist convention
and the adjoint representation on $V$, and put
$$
\Omega_X^1\bigl(\operatorname{ad}(\mathcal E)\bigr)
=\Omega_X^1\otimes_{\mathcal O_X}\operatorname{ad}(\mathcal E).
$$
There is a canonical isomorphism
$$
\mathcal A(\sigma)^{\mathcal E}
\cong
\Omega_X^1\bigl(\operatorname{ad}(\mathcal E)\bigr).
$$
The affine twist $\mathcal A(\rho)^{\mathcal E}$ is the torsor of
holomorphic connection forms on $\mathcal E$ in this convention. Hence
Lemma \ref{lemma-crossed-homomorphism-torsor-obstruction} gives a class
$$
d(\mathcal E)=d_u(\mathcal E)
\in
H^1\bigl(X,
\Omega_X^1(\operatorname{ad}(\mathcal E))\bigr),
$$
and this class is zero if and only if $\mathcal E$ admits a holomorphic
connection.
\item A reduction of $\mathcal E$ from $\mathcal P$ to
$\underline G_X$ is equivalent to an integrable holomorphic connection on
$\mathcal E$. Either condition implies $d(\mathcal E)=0$.
\item If $X$ has complex dimension $1$, then $u$ is an epimorphism of
sheaves. Consequently
\begin{align*}
d(\mathcal E)=0
&\quad\Longleftrightarrow\quad
\mathcal E\text{ reduces to }\underline G_X, \\
&\quad\Longleftrightarrow\quad
\mathcal E\text{ admits an integrable holomorphic connection}.
\end{align*}
In this case every holomorphic connection is integrable.
\end{enumerate}
\end{lemma}

\begin{proof}
The adjoint action on a $V$-valued form is pointwise. The product rule and
right translation to the identity give
\begin{align*}
d(gg')(gg')^{-1}
&=(dg)g^{-1}
+g(dg')(g')^{-1}g^{-1} \\
&=u(g)+\operatorname{Ad}_g\bigl(u(g')\bigr).
\end{align*}
This proves the crossed homomorphism identity and the assertion about the
affine action. Since right translation induces an isomorphism on tangent
spaces, $u(g)=0$ if and only if $dg=0$. A holomorphic map with zero
differential is locally constant, which proves the assertion about the
kernel.

Choose an open covering $X=\bigcup U_i$ and sections
$e_i\in\mathcal E(U_i)$. Write
$$
e_i=g_{ij}e_j
$$
on overlaps. The linear twist is obtained by gluing $V$-valued holomorphic
$1$-forms by
$$
\alpha_i=\operatorname{Ad}_{g_{ij}}(\alpha_j).
$$
This is precisely the descent description of
$\Omega_X^1(\operatorname{ad}(\mathcal E))$, and proves the canonical
isomorphism in (2).

A section of the affine twist is similarly a collection of local forms
$\alpha_i$ satisfying
$$
\alpha_i=(dg_{ij})g_{ij}^{-1}
+\operatorname{Ad}_{g_{ij}}(\alpha_j).
$$
This is the local connection-form transformation law for the adopted
left-torsor and right-logarithmic-derivative convention. Thus global sections
of $\mathcal A(\rho)^{\mathcal E}$ are exactly holomorphic connections on
$\mathcal E$. The linear twist is abelian, so
Lemma \ref{lemma-crossed-homomorphism-torsor-obstruction} and
Lemma \ref{lemma-torsors-h1} identify the isomorphism class of this affine
torsor with the displayed element $d_u(\mathcal E)$, and identify its
neutrality with the existence of a connection.

There is no hidden sign in this identification. Indeed, if
$r_i=[e_i,0]$ denotes the local origin of the affine twist, then the orbit
relation gives
$$
r_j=[e_i,u(g_{ij})].
$$
Thus Lemma \ref{lemma-cech-h1} represents the torsor by the positive twisted
cocycle $([e_i,u(g_{ij})])$. This is the crossed-homomorphism class denoted
$d(E)$ in the cited source.

By the reduction implication in
Lemma \ref{lemma-crossed-homomorphism-torsor-obstruction}, a reduction to
$\underline G_X$ implies $d(\mathcal E)=0$. In frames for such a reduction
we have $u(g_{ij})=0$, so the zero local forms glue. The resulting connection
has zero curvature and is integrable.

Conversely, suppose that $\mathcal E$ has an integrable holomorphic
connection. Its horizontal distribution is holomorphic and involutive, so
the holomorphic Frobenius theorem gives local horizontal frames. In these
frames the connection forms vanish. The transformation law then gives
$(dg_{ij})g_{ij}^{-1}=0$, so the transition maps are locally constant and
define a reduction to $\underline G_X$.

It remains to consider a complex curve. In a local coordinate $z$, write a
holomorphic $V$-valued $1$-form as $a(z)\,dz$. The equation
$$
(dg)g^{-1}=a(z)\,dz
$$
is the holomorphic ordinary differential equation
$$
g'(z)=(dR_{g(z)})_1a(z)
$$
on $G$, and it has a local holomorphic solution through every prescribed
initial value. Hence $u$ is an epimorphism of sheaves. The converse reduction
implication now follows from
Lemma \ref{lemma-crossed-homomorphism-torsor-obstruction}. Finally, the
curvature of a holomorphic connection is an
$\operatorname{ad}(\mathcal E)$-valued holomorphic $2$-form, and such forms
vanish on a complex curve. This proves (4).
\end{proof}

\noindent
The following lemma does not hold for flasque presheaves.

\begin{lemma}
\label{lemma-flasque-acyclic-cech}
Let $(X, \mathcal{O}_X)$ be a ringed space.
Let $\mathcal{F}$ be a sheaf of $\mathcal{O}_X$-modules.
Let $\mathcal{U} : U = \bigcup U_i$ be an open covering.
If $\mathcal{F}$ is flasque, then
$\check{H}^p(\mathcal{U}, \mathcal{F}) = 0$ for $p > 0$.
\end{lemma}

\begin{proof}
The presheaves $\underline{H}^q(\mathcal{F})$ used in the statement
of Lemma \ref{lemma-cech-spectral-sequence} are zero by
Lemma \ref{lemma-flasque-acyclic}.
Hence $\check{H}^p(U, \mathcal{F}) = H^p(U, \mathcal{F}) = 0$
by Lemma \ref{lemma-flasque-acyclic} again.
\end{proof}

\begin{lemma}
\label{lemma-flasque-acyclic-pushforward}
Let $f : (X, \mathcal{O}_X) \to (Y, \mathcal{O}_Y)$ be a morphism
of ringed spaces. Let $\mathcal{F}$ be a sheaf of $\mathcal{O}_X$-modules.
If $\mathcal{F}$ is flasque, then $R^pf_*\mathcal{F} = 0$ for $p > 0$.
\end{lemma}

\begin{proof}
Immediate from
Lemma \ref{lemma-describe-higher-direct-images} and
Lemma \ref{lemma-flasque-acyclic}.
\end{proof}

\noindent
The following lemma can be proved by an elementary induction
argument for finite coverings, compare with the discussion
of {\v C}ech cohomology in \cite{FOAG}.

\begin{lemma}
\label{lemma-vanishing-ravi}
Let $X$ be a topological space. Let $\mathcal{F}$ be an abelian sheaf
on $X$. Let $\mathcal{U} : U = \bigcup_{i \in I} U_i$ be an
open covering. Assume the restriction mappings
$\mathcal{F}(U) \to \mathcal{F}(U')$ are surjective
for $U'$ an arbitrary union of opens of the form $U_{i_0 \ldots i_p}$.
Then $\check{H}^p(\mathcal{U}, \mathcal{F})$
vanishes for $p > 0$.
\end{lemma}

\begin{proof}
Let $Y$ be the set of nonempty subsets of $I$. We will use the letters
$A, B, C, \ldots$ to denote elements of $Y$, i.e., nonempty subsets of $I$.
For a finite nonempty subset $J \subset I$ let
$$
V_J = \{A \in Y \mid J \subset A\}
$$
This means that $V_{\{i\}} = \{A \in Y \mid i \in A\}$ and
$V_J = \bigcap_{j \in J} V_{\{j\}}$.
Then $V_J \subset V_K$ if and only if $J \supset K$.
There is a unique topology on $Y$ such that the collection of
subsets $V_J$ is a basis for the topology on $Y$. Any open is of the form
$$
V = \bigcup\nolimits_{t \in T} V_{J_t}
$$
for some family of finite subsets $J_t$. If $J_t \subset J_{t'}$
then we may remove $J_{t'}$ from the family without changing $V$.
Thus we may assume there are no inclusions among the $J_t$.
In this case the minimal elements of $V$ are the sets $A = J_t$.
Hence we can read off the family $(J_t)_{t \in T}$ from the open $V$.

\medskip\noindent
We can completely understand open coverings in $Y$. First, because
the elements $A \in Y$ are nonempty subsets of $I$ we have
$$
Y = \bigcup\nolimits_{i \in I} V_{\{i\}}
$$
To understand other coverings, let $V$ be as above and let $V_s \subset Y$
be an open corresponding to the family $(J_{s, t})_{t \in T_s}$. Then
$$
V = \bigcup\nolimits_{s \in S} V_s
$$
if and only if for each $t \in T$ there exists an $s \in S$ and
$t_s \in T_s$ such that $J_t = J_{s, t_s}$. Namely, as the family
$(J_t)_{t \in T}$ is minimal, the minimal element $A = J_t$
has to be in $V_s$ for some $s$, hence $A \in V_{J_{t_s}}$ for some
$t_s \in T_s$. But since $A$ is also minimal in $V_s$ we conclude
that $J_{t_s} = J_t$.

\medskip\noindent
Next we map the set of opens of $Y$ to opens of $X$. Namely, we send
$Y$ to $U$, we use the rule
$$
V_J \mapsto U_J = \bigcap\nolimits_{i \in J} U_i
$$
on the opens $V_J$, and we extend it to arbitrary opens $V$ by the rule
$$
V = \bigcup\nolimits_{t \in T} V_{J_t}
\mapsto
\bigcup\nolimits_{t \in T} U_{J_t}
$$
The classification of open coverings of $Y$ given above shows that
this rule transforms open coverings into open coverings. Thus we obtain
an abelian sheaf $\mathcal{G}$ on $Y$ by setting
$\mathcal{G}(Y) = \mathcal{F}(U)$ and for
$V = \bigcup\nolimits_{t \in T} V_{J_t}$ setting
$$
\mathcal{G}(V) = \mathcal{F}\left(\bigcup\nolimits_{t \in T} U_{J_t}\right)
$$
and using the restriction maps of $\mathcal{F}$.

\medskip\noindent
With these preliminaries out of the way we can prove our lemma as follows.
We have an open covering
$\mathcal{V} : Y = \bigcup_{i \in I} V_{\{i\}}$ of $Y$.
By construction we have an equality
$$
\check{C}^\bullet(\mathcal{V}, \mathcal{G}) =
\check{C}^\bullet(\mathcal{U}, \mathcal{F})
$$
of {\v C}ech complexes. Since the sheaf $\mathcal{G}$ is flasque on $Y$
(by our assumption on $\mathcal{F}$ in the statement of the lemma)
the vanishing follows from
Lemma \ref{lemma-flasque-acyclic-cech}.
\end{proof}




\section{The Leray spectral sequence}
\label{section-Leray}

\begin{lemma}
\label{lemma-before-Leray}
Let $f : X \to Y$ be a morphism of ringed spaces.
There is a commutative diagram
$$
\xymatrix{
D^{+}(X) \ar[rr]_-{R\Gamma(X, -)} \ar[d]_{Rf_*} & &
D^{+}(\mathcal{O}_X(X)) \ar[d]^{\text{restriction}} \\
D^{+}(Y) \ar[rr]^-{R\Gamma(Y, -)} & &
D^{+}(\mathcal{O}_Y(Y))
}
$$
More generally for any $V \subset Y$ open and $U = f^{-1}(V)$ there
is a commutative diagram
$$
\xymatrix{
D^{+}(X) \ar[rr]_-{R\Gamma(U, -)} \ar[d]_{Rf_*} & &
D^{+}(\mathcal{O}_X(U)) \ar[d]^{\text{restriction}} \\
D^{+}(Y) \ar[rr]^-{R\Gamma(V, -)} & &
D^{+}(\mathcal{O}_Y(V))
}
$$
See also Remark \ref{remark-elucidate-lemma} for more explanation.
\end{lemma}

\begin{proof}
Let
$\Gamma_{res} : \textit{Mod}(\mathcal{O}_X) \to \text{Mod}_{\mathcal{O}_Y(Y)}$
be the functor which associates to an $\mathcal{O}_X$-module $\mathcal{F}$
the global sections of $\mathcal{F}$ viewed as an $\mathcal{O}_Y(Y)$-module
via the map $f^\sharp : \mathcal{O}_Y(Y) \to \mathcal{O}_X(X)$. Let
$restriction : \text{Mod}_{\mathcal{O}_X(X)} \to \text{Mod}_{\mathcal{O}_Y(Y)}$
be the restriction functor induced by
$f^\sharp : \mathcal{O}_Y(Y) \to \mathcal{O}_X(X)$. Note that $restriction$
is exact so that
its right derived functor is computed by simply applying the restriction
functor, see
Derived Categories, Lemma \ref{derived-lemma-right-derived-exact-functor}.
It is clear that
$$
\Gamma_{res}
=
restriction \circ \Gamma(X, -)
=
\Gamma(Y, -) \circ f_*
$$
We claim that
Derived Categories, Lemma \ref{derived-lemma-compose-derived-functors}
applies to both compositions. For the first this is clear by our remarks
above. For the second, it follows from
Lemma \ref{lemma-pushforward-injective} which implies that
injective $\mathcal{O}_X$-modules are mapped to $\Gamma(Y, -)$-acyclic
sheaves on $Y$.
\end{proof}

\begin{remark}
\label{remark-elucidate-lemma}
Here is a down-to-earth explanation of the meaning of
Lemma \ref{lemma-before-Leray}. It says that given
$f : X \to Y$ and $\mathcal{F} \in \textit{Mod}(\mathcal{O}_X)$
and given an injective resolution $\mathcal{F} \to \mathcal{I}^\bullet$
we have
$$
\begin{matrix}
R\Gamma(X, \mathcal{F}) & \text{is represented by} &
\Gamma(X, \mathcal{I}^\bullet) \\
Rf_*\mathcal{F} & \text{is represented by} & f_*\mathcal{I}^\bullet \\
R\Gamma(Y, Rf_*\mathcal{F}) & \text{is represented by} &
\Gamma(Y, f_*\mathcal{I}^\bullet)
\end{matrix}
$$
the last fact coming from Leray's acyclicity lemma
(Derived Categories, Lemma \ref{derived-lemma-leray-acyclicity})
and Lemma \ref{lemma-pushforward-injective}.
Finally, it combines this with the trivial observation that
$$
\Gamma(X, \mathcal{I}^\bullet)
=
\Gamma(Y, f_*\mathcal{I}^\bullet).
$$
to arrive at the commutativity of the diagram of the lemma.
\end{remark}

\begin{lemma}
\label{lemma-modules-abelian}
Let $X$ be a ringed space.
Let $\mathcal{F}$ be an $\mathcal{O}_X$-module.
\begin{enumerate}
\item The cohomology groups $H^i(U, \mathcal{F})$ for $U \subset X$ open
of $\mathcal{F}$ computed as an $\mathcal{O}_X$-module, or computed as an
abelian sheaf are identical.
\item Let $f : X \to Y$ be a morphism of ringed spaces.
The higher direct images $R^if_*\mathcal{F}$ of $\mathcal{F}$
computed as an $\mathcal{O}_X$-module, or computed as an abelian sheaf
are identical.
\end{enumerate}
There are similar statements in the case of bounded below
complexes of $\mathcal{O}_X$-modules.
\end{lemma}

\begin{proof}
Consider the morphism of ringed spaces
$(X, \mathcal{O}_X) \to (X, \underline{\mathbf{Z}}_X)$ given
by the identity on the underlying topological space and by
the unique map of sheaves of rings
$\underline{\mathbf{Z}}_X \to \mathcal{O}_X$.
Let $\mathcal{F}$ be an $\mathcal{O}_X$-module.
Denote $\mathcal{F}_{ab}$ the same sheaf seen as an
$\underline{\mathbf{Z}}_X$-module, i.e., seen as a sheaf of
abelian groups. Let
$\mathcal{F} \to \mathcal{I}^\bullet$ be an injective resolution.
By Remark \ref{remark-elucidate-lemma} we see that
$\Gamma(X, \mathcal{I}^\bullet)$ computes both
$R\Gamma(X, \mathcal{F})$ and $R\Gamma(X, \mathcal{F}_{ab})$.
This proves (1).

\medskip\noindent
To prove (2) we use (1) and Lemma \ref{lemma-describe-higher-direct-images}.
The result follows immediately.
\end{proof}

\begin{lemma}[Leray spectral sequence]
\label{lemma-Leray}
Let $f : X \to Y$ be a morphism of ringed spaces.
Let $\mathcal{F}^\bullet$ be
a bounded below complex of $\mathcal{O}_X$-modules.
There is a spectral sequence
$$
E_2^{p, q} = H^p(Y, R^qf_*(\mathcal{F}^\bullet))
$$
converging to $H^{p + q}(X, \mathcal{F}^\bullet)$.
\end{lemma}

\begin{proof}
This is just the Grothendieck spectral sequence
Derived Categories, Lemma \ref{derived-lemma-grothendieck-spectral-sequence}
coming from the composition of functors
$\Gamma_{res} = \Gamma(Y, -) \circ f_*$ where $\Gamma_{res}$ is as
in the proof of Lemma \ref{lemma-before-Leray}.
To see that the assumptions of
Derived Categories, Lemma \ref{derived-lemma-grothendieck-spectral-sequence}
are satisfied, see the proof of Lemma \ref{lemma-before-Leray} or
Remark \ref{remark-elucidate-lemma}.
\end{proof}

\begin{remark}
\label{remark-Leray-ss-more-structure}
The Leray spectral sequence, the way we proved it in Lemma \ref{lemma-Leray}
is a spectral sequence of $\Gamma(Y, \mathcal{O}_Y)$-modules. However, it
is quite easy to see that it is in fact a spectral sequence of
$\Gamma(X, \mathcal{O}_X)$-modules. For example $f$ gives rise to
a morphism of ringed spaces
$f' :  (X, \mathcal{O}_X) \to (Y, f_*\mathcal{O}_X)$.
By Lemma \ref{lemma-modules-abelian} the terms $E_r^{p, q}$ of the
Leray spectral sequence for an $\mathcal{O}_X$-module $\mathcal{F}$
and $f$ are identical with those for $\mathcal{F}$ and $f'$
at least for $r \geq 2$. Namely, they both agree with the terms of the Leray
spectral sequence for $\mathcal{F}$ as an abelian sheaf.
And since $(f_*\mathcal{O}_X)(Y) = \mathcal{O}_X(X)$ we see the result.
It is often the case
that the Leray spectral sequence carries additional structure.
\end{remark}

\begin{lemma}
\label{lemma-apply-Leray}
Let $f : X \to Y$ be a morphism of ringed spaces.
Let $\mathcal{F}$ be an $\mathcal{O}_X$-module.
\begin{enumerate}
\item If $R^qf_*\mathcal{F} = 0$ for $q > 0$, then
$H^p(X, \mathcal{F}) = H^p(Y, f_*\mathcal{F})$ for all $p$.
\item If $H^p(Y, R^qf_*\mathcal{F}) = 0$ for all $q$ and $p > 0$, then
$H^q(X, \mathcal{F}) = H^0(Y, R^qf_*\mathcal{F})$ for all $q$.
\end{enumerate}
\end{lemma}

\begin{proof}
These are two simple conditions that force the Leray spectral sequence to
degenerate at $E_2$. You can also prove these facts directly (without using
the spectral sequence) which is a good exercise in cohomology of sheaves.
\end{proof}

\begin{lemma}
\label{lemma-higher-direct-images-compose}
\begin{slogan}
The total derived functor of a composition is the
composition of the total derived functors.
\end{slogan}
Let $f : X \to Y$ and $g : Y \to Z$ be morphisms of ringed spaces.
In this case $Rg_* \circ Rf_* = R(g \circ f)_*$ as functors
from $D^{+}(X) \to D^{+}(Z)$.
\end{lemma}

\begin{proof}
We are going to apply
Derived Categories, Lemma \ref{derived-lemma-compose-derived-functors}.
It is clear that $g_* \circ f_* = (g \circ f)_*$, see
Sheaves, Lemma \ref{sheaves-lemma-pushforward-composition}.
It remains to show that $f_*\mathcal{I}$ is $g_*$-acyclic.
This follows from Lemma \ref{lemma-pushforward-injective}
and the description of the
higher direct images $R^ig_*$ in
Lemma \ref{lemma-describe-higher-direct-images}.
\end{proof}

\begin{lemma}[Relative Leray spectral sequence]
\label{lemma-relative-Leray}
Let $f : X \to Y$ and $g : Y \to Z$ be morphisms of ringed spaces.
Let $\mathcal{F}$ be an $\mathcal{O}_X$-module.
There is a spectral sequence with
$$
E_2^{p, q} = R^pg_*(R^qf_*\mathcal{F})
$$
converging to $R^{p + q}(g \circ f)_*\mathcal{F}$.
This spectral sequence is functorial in $\mathcal{F}$, and there
is a version for bounded below complexes of $\mathcal{O}_X$-modules.
\end{lemma}

\begin{proof}
This is a Grothendieck spectral sequence for composition of functors
and follows from Lemma \ref{lemma-higher-direct-images-compose} and
Derived Categories, Lemma \ref{derived-lemma-grothendieck-spectral-sequence}.
\end{proof}














\section{Functoriality of cohomology}
\label{section-functoriality}

\begin{lemma}
\label{lemma-functoriality}
Let $f : X \to Y$ be a morphism of ringed spaces.
Let $\mathcal{G}^\bullet$, resp.\ $\mathcal{F}^\bullet$ be
a bounded below complex of $\mathcal{O}_Y$-modules,
resp.\ $\mathcal{O}_X$-modules. Let
$\varphi : \mathcal{G}^\bullet \to f_*\mathcal{F}^\bullet$
be a morphism of complexes. There is a canonical morphism
$$
\mathcal{G}^\bullet
\longrightarrow
Rf_*(\mathcal{F}^\bullet)
$$
in $D^{+}(Y)$. Moreover this construction is functorial in the triple
$(\mathcal{G}^\bullet, \mathcal{F}^\bullet, \varphi)$.
\end{lemma}

\begin{proof}
Choose an injective resolution $\mathcal{F}^\bullet \to \mathcal{I}^\bullet$.
By definition $Rf_*(\mathcal{F}^\bullet)$ is represented by
$f_*\mathcal{I}^\bullet$ in $K^{+}(\mathcal{O}_Y)$.
The composition
$$
\mathcal{G}^\bullet \to f_*\mathcal{F}^\bullet \to f_*\mathcal{I}^\bullet
$$
is a morphism in $K^{+}(Y)$ which turns
into the morphism of the lemma upon applying the
localization functor $j_Y : K^{+}(Y) \to D^{+}(Y)$.
\end{proof}

\noindent
Let $f : X \to Y$ be a morphism of ringed spaces.
Let $\mathcal{G}$ be an $\mathcal{O}_Y$-module and let
$\mathcal{F}$ be an $\mathcal{O}_X$-module. Recall that an
$f$-map $\varphi$ from $\mathcal{G}$ to $\mathcal{F}$ is a map
$\varphi : \mathcal{G} \to f_*\mathcal{F}$, or what is the same
thing, a map $\varphi : f^*\mathcal{G} \to \mathcal{F}$.
See Sheaves, Definition \ref{sheaves-definition-f-map}.
Such an $f$-map gives rise to a morphism of complexes
\begin{equation}
\label{equation-functorial-derived}
\varphi :
R\Gamma(Y, \mathcal{G})
\longrightarrow
R\Gamma(X, \mathcal{F})
\end{equation}
in $D^{+}(\mathcal{O}_Y(Y))$. Namely, we use the morphism
$\mathcal{G} \to Rf_*\mathcal{F}$ in $D^{+}(Y)$ of
Lemma \ref{lemma-functoriality}, and we apply $R\Gamma(Y, -)$.
By Lemma \ref{lemma-before-Leray} we see that
$R\Gamma(X, \mathcal{F}) = R\Gamma(Y, Rf_*\mathcal{F})$
and we get the displayed arrow. We spell this out completely in
Remark \ref{remark-explain-arrow} below.
In particular it gives
rise to maps on cohomology
\begin{equation}
\label{equation-functorial}
\varphi : H^i(Y, \mathcal{G}) \longrightarrow H^i(X, \mathcal{F}).
\end{equation}

\begin{remark}
\label{remark-explain-arrow}
Let $f : X \to Y$ be a morphism of ringed spaces.
Let $\mathcal{G}$ be an $\mathcal{O}_Y$-module.
Let $\mathcal{F}$ be an $\mathcal{O}_X$-module.
Let $\varphi$ be an $f$-map from $\mathcal{G}$ to $\mathcal{F}$.
Choose a resolution $\mathcal{F} \to \mathcal{I}^\bullet$
by a complex of injective $\mathcal{O}_X$-modules.
Choose resolutions $\mathcal{G} \to \mathcal{J}^\bullet$ and
$f_*\mathcal{I}^\bullet \to (\mathcal{J}')^\bullet$ by complexes
of injective $\mathcal{O}_Y$-modules. By
Derived Categories, Lemma \ref{derived-lemma-morphisms-lift}
there exists a map of complexes
$\beta$ such that the diagram
\begin{equation}
\label{equation-choice}
\xymatrix{
\mathcal{G} \ar[d] \ar[r] &
f_*\mathcal{F} \ar[r] &
f_*\mathcal{I}^\bullet \ar[d] \\
\mathcal{J}^\bullet \ar[rr]^\beta & &
(\mathcal{J}')^\bullet
}
\end{equation}
commutes. Applying global section functors we see
that we get a diagram
$$
\xymatrix{
 & & \Gamma(Y, f_*\mathcal{I}^\bullet) \ar[d]_{qis} \ar@{=}[r] &
\Gamma(X, \mathcal{I}^\bullet) \\
\Gamma(Y, \mathcal{J}^\bullet) \ar[rr]^\beta & &
\Gamma(Y, (\mathcal{J}')^\bullet) &
}
$$
The complex on the bottom left represents $R\Gamma(Y, \mathcal{G})$
and the complex on the top right represents $R\Gamma(X, \mathcal{F})$.
The vertical arrow is a quasi-isomorphism by
Lemma \ref{lemma-before-Leray} which becomes invertible after
applying the localization functor
$K^{+}(\mathcal{O}_Y(Y)) \to D^{+}(\mathcal{O}_Y(Y))$.
The arrow (\ref{equation-functorial-derived}) is given by the
composition of the horizontal map by the inverse of the vertical map.
\end{remark}





\section{Refinements and {\v C}ech cohomology}
\label{section-refinements-cech}

\noindent
Let $(X, \mathcal{O}_X)$ be a ringed space. Let
$\mathcal{U} : X = \bigcup_{i \in I} U_i$ and
$\mathcal{V} : X = \bigcup_{j \in J} V_j$ be open coverings.
Assume that $\mathcal{U}$ is a refinement of $\mathcal{V}$.
Choose a map $c : I \to J$ such that $U_i \subset V_{c(i)}$
for all $i \in I$. This induces a map of {\v C}ech complexes
$$
\gamma :
\check{\mathcal{C}}^\bullet(\mathcal{V}, \mathcal{F})
\longrightarrow
\check{\mathcal{C}}^\bullet(\mathcal{U}, \mathcal{F}),
\quad
(\xi_{j_0 \ldots j_p})
\longmapsto
(\xi_{c(i_0) \ldots c(i_p)}|_{U_{i_0 \ldots i_p}})
$$
functorial in the sheaf of $\mathcal{O}_X$-modules $\mathcal{F}$.
Suppose that $c' : I \to J$ is a second map such that
$U_i \subset V_{c'(i)}$ for all $i \in I$. Then the corresponding maps
$\gamma$ and $\gamma'$ are homotopic. Namely,
$\gamma' - \gamma = \text{d} \circ h + h \circ \text{d}$
with
$h : \check{\mathcal{C}}^{p + 1}(\mathcal{V}, \mathcal{F}) \to
\check{\mathcal{C}}^p(\mathcal{U}, \mathcal{F})$
given by the rule
$$
h(\alpha)_{i_0 \ldots i_p} =
\sum\nolimits_{a = 0}^{p}
(-1)^a
\alpha_{c(i_0)\ldots c(i_a) c'(i_a) \ldots c'(i_p)}
$$
We omit the computation showing this works; please see the discussion
following (\ref{equation-transformation}) for the proof in a more general
case. In particular, the map on {\v C}ech cohomology groups is independent
of the choice of $c$.

\begin{remark}
\label{remark-FAC-refinement-homotopy}
\begin{reference}
\cite[Chapter I, \S 3, no. 21, Proposition 3, pp. 214--215]{FAC}
\end{reference}
For the ordering $c$, $c'$ used in the formula for $h$, its boundary is
$\gamma' - \gamma$. This agrees with the total-complex calculation
following Equation (\ref{equation-transformation}). Thus the cochain map
depends on the refinement function, while its map on cohomology does not.
\end{remark}

\noindent
Moreover, it is clear that if
$\mathcal{W} : X = \bigcup_{k \in K} W_k$ is a third open covering
and $\mathcal{V}$ is a refinement of $\mathcal{W}$, then the composition
of the maps
$$
\check{\mathcal{C}}^\bullet(\mathcal{W}, \mathcal{F})
\longrightarrow
\check{\mathcal{C}}^\bullet(\mathcal{V}, \mathcal{F})
\longrightarrow
\check{\mathcal{C}}^\bullet(\mathcal{U}, \mathcal{F})
$$
associated to maps $I \to J$ and $J \to K$ is the map associated
to the composition $I \to K$.

\begin{remark}
\label{remark-FAC-two-covering-double-complex}
\begin{reference}
\cite[Chapter I, \S 4, no. 28, Proposition 2, pp. 220--221]{FAC}
\end{reference}
Let $\mathcal{U} : X = \bigcup_{i \in I} U_i$ and
$\mathcal{V} : X = \bigcup_{j \in J} V_j$ be open coverings and let
$\mathcal{F}$ be an abelian sheaf. The two-covering groups
$$
A^{p, q} =
\prod_{(i_0, \ldots, i_p), (j_0, \ldots, j_q)}
\mathcal{F}(U_{i_0 \ldots i_p} \cap V_{j_0 \ldots j_q})
$$
form a double complex in the convention of
Homology, Definition \ref{homology-definition-double-complex}:
the first and second differentials are the Cech differentials in the
$i$- and $j$-indices. They commute, and the differential on the associated
total complex is therefore $d_1 + (-1)^p d_2$ on $A^{p, q}$.

For a fixed $(j_0, \ldots, j_q)$, let
$\mathcal{U}_{j_0 \ldots j_q}$ be the covering of
$V_{j_0 \ldots j_q}$ by the opens
$U_i \cap V_{j_0 \ldots j_q}$. Taking cohomology in the first direction
gives
$$
H^p_I(A^{\bullet, \bullet})^q =
\prod_{(j_0, \ldots, j_q)}
\check{H}^p(\mathcal{U}_{j_0 \ldots j_q}, \mathcal{F}).
$$
In degree $p = 0$, the sheaf condition identifies the resulting edge
complex with $\check{\mathcal{C}}^\bullet(\mathcal{V}, \mathcal{F})$.
Interchanging the two coverings gives the analogous identification with
$\check{\mathcal{C}}^\bullet(\mathcal{U}, \mathcal{F})$.
\end{remark}

\begin{definition}
\label{definition-cech-cohomology-space}
\begin{reference}
\cite[Chapter I, \S 3, no. 22, Definition, p. 215]{FAC}
\end{reference}
Let $(X, \mathcal{O}_X)$ be a ringed space. Let $\mathcal{F}$ be an
$\mathcal{O}_X$-module. The groups
$$
\check{H}^p(X, \mathcal{F}) =
\colim_\mathcal{U} \check{H}^p(\mathcal{U}, \mathcal{F})
$$
where the colimit is over all open coverings of $X$ preordered by refinement,
are the {\it {\v C}ech cohomology groups} of $X$ with coefficients in
$\mathcal{F}$.
\end{definition}

\begin{remark}
\label{remark-FAC-cech-colimit-notation}
\begin{reference}
\cite[Chapter I, \S 3, no. 22 and footnote 2, p. 215]{FAC}
\cite[Chapter III, \S 1, no. 53, Proposition 4 and Remark,
pp. 245--246]{FAC}
\end{reference}
The source denotes this direct limit by $H^p(X, \mathcal{F})$. We use
$\check{H}^p(X, \mathcal{F})$ to distinguish it from derived-functor
cohomology. The colimit can be indexed by a set: replace every covering
family by the set of its open members, a subset of the topology of $X$.
Finite intersections give a common refinement of any two coverings.
This convention is still in force in no. 53, where the cofinal curve-covering
argument proves a statement about the direct-limit Cech groups. The
corresponding derived-functor vanishing theorem is Proposition
\ref{proposition-vanishing-Noetherian}; this comparison does not identify the
two theories in all degrees without an additional comparison theorem.
\end{remark}

\begin{lemma}
\label{lemma-cech-refinement-acyclic}
\begin{reference}
\cite[Chapter I, \S 4, no. 29, Propositions 3--5,
pp. 221--222]{FAC}
\end{reference}
Let $(X, \mathcal{O}_X)$ be a ringed space and let
$\mathcal{F}$ be an $\mathcal{O}_X$-module. Let
$\mathcal{U} : X = \bigcup_{i \in I} U_i$ and
$\mathcal{V} : X = \bigcup_{j \in J} V_j$ be open coverings such that
$\mathcal{V}$ refines $\mathcal{U}$. For a tuple
$(i_0, \ldots, i_p)$, denote by
$\mathcal{V}|_{U_{i_0 \ldots i_p}}$ the covering of
$U_{i_0 \ldots i_p}$ by the opens
$U_{i_0 \ldots i_p} \cap V_j$. If
$$
\check{H}^q(\mathcal{V}|_{U_{i_0 \ldots i_p}}, \mathcal{F}) = 0
$$
for all $q > 0$, $p \geq 0$, and $i_0, \ldots, i_p \in I$, then the
map
$$
\check{H}^n(\mathcal{U}, \mathcal{F})
\longrightarrow
\check{H}^n(\mathcal{V}, \mathcal{F})
$$
induced by refinement is an isomorphism for every $n \geq 0$.
\end{lemma}

\begin{proof}
Let $A^{\bullet, \bullet}$ be the double complex of
Remark \ref{remark-FAC-two-covering-double-complex}. The hypothesis and
Lemma \ref{lemma-cech-h0} show that the edge map
$$
\check{\mathcal{C}}^\bullet(\mathcal{U}, \mathcal{F})
\longrightarrow \text{Tot}(A^{\bullet, \bullet})
$$
is a quasi-isomorphism by
Homology, Lemma \ref{homology-lemma-double-complex-gives-resolution}.

Choose $c : J \to I$ such that $V_j \subset U_{c(j)}$, and denote the
induced refinement map of Cech complexes by $\gamma_c$. For every tuple
$(j_0, \ldots, j_q)$, the covering of $V_{j_0 \ldots j_q}$ by the
opens $V_{j_0 \ldots j_q} \cap U_i$ contains
$V_{j_0 \ldots j_q}$ itself, namely for $i = c(j_0)$. Hence
Lemma \ref{lemma-cech-trivial} and the second-variable version of the
same Homology lemma show that the other edge map
$$
\check{\mathcal{C}}^\bullet(\mathcal{V}, \mathcal{F})
\longrightarrow \text{Tot}(A^{\bullet, \bullet})
$$
is a quasi-isomorphism as well.

It remains to identify the isomorphism obtained from the two edge maps.
For $n > 0$ and an $n$-cocycle $f$ of
$\check{\mathcal{C}}^\bullet(\mathcal{U}, \mathcal{F})$, set
$$
g^p_{i_0 \ldots i_p,j_0 \ldots j_{n-p-1}} =
f_{i_0 \ldots i_p c(j_0) \ldots c(j_{n-p-1})}
\big|_{U_{i_0 \ldots i_p} \cap V_{j_0 \ldots j_{n-p-1}}}
$$
for $0 \leq p \leq n - 1$. Writing $D$ for the total differential, a
direct calculation using $d(f) = 0$ gives
$$
D(g^0 - g^1 + \ldots + (-1)^{n - 1}g^{n - 1})
= \iota_{\mathcal{V}}(\gamma_c(f)) - \iota_{\mathcal{U}}(f),
$$
where $\iota_{\mathcal{U}}$ and $\iota_{\mathcal{V}}$ denote the two edge
maps above. Thus the composite obtained from the edge maps is the usual refinement
map. For $n = 0$ this follows directly from the sheaf condition. The
result follows.
\end{proof}

\begin{lemma}
\label{lemma-cech-cofinal-refinements}
\begin{reference}
\cite[Chapter I, \S 4, no. 29, Theorem 1, p. 222]{FAC}
\end{reference}
Let $(X, \mathcal{O}_X)$ be a ringed space, let
$\mathcal{F}$ be an $\mathcal{O}_X$-module, and let
$\mathcal{U} : X = \bigcup_{i \in I} U_i$ be an open covering. Let
$\mathcal{B}$ be a collection of open coverings cofinal for refinement.
Assume that for every $\mathcal{V} \in \mathcal{B}$, every $p \geq 0$,
every $i_0, \ldots, i_p \in I$, and every $q > 0$, we have
$$
\check{H}^q(\mathcal{V}|_{U_{i_0 \ldots i_p}}, \mathcal{F}) = 0.
$$
Then the canonical map
$$
\check{H}^n(\mathcal{U}, \mathcal{F})
\longrightarrow \check{H}^n(X, \mathcal{F})
$$
is an isomorphism for every $n \geq 0$.
\end{lemma}

\begin{proof}
The subcollection of coverings $\mathcal{V} \in \mathcal{B}$ which refine
$\mathcal{U}$ is still cofinal. Indeed, given a covering $\mathcal{W}$,
first take the common refinement by the intersections of members of
$\mathcal{W}$ and $\mathcal{U}$ and then use cofinality of $\mathcal{B}$.
For every member of this subcollection, the map from
$\check{H}^n(\mathcal{U}, \mathcal{F})$ is an isomorphism by
Lemma \ref{lemma-cech-refinement-acyclic}. These isomorphisms are compatible
with further refinement. Taking the colimit proves the lemma.
\end{proof}

\begin{remark}
\label{remark-FAC-cech-cofinal-question}
\begin{reference}
\cite[Chapter I, \S 4, no. 29, Remarks 1 and 2, p. 222]{FAC}
\end{reference}
The source asks whether the pointwise vanishing hypothesis of
Lemma \ref{lemma-cech-cofinal-refinements} can be replaced by vanishing of
the corresponding groups only after passage to the cofinal system. We do
not assert that weakening here: the proof above uses the edge
quasi-isomorphism before taking the colimit. The usual Leray criterion for
comparison with derived-functor cohomology is
Lemma \ref{lemma-cech-spectral-sequence-application}; its hypothesis is
vanishing of derived cohomology on all finite intersections, and is not the
same condition.
\end{remark}

\begin{lemma}
\label{lemma-cech-h0-space}
\begin{reference}
\cite[Chapter I, \S 3, no. 22, Proposition 4, p. 215]{FAC}
\end{reference}
Let $(X, \mathcal{O}_X)$ be a ringed space. Let $\mathcal{F}$ be an
$\mathcal{O}_X$-module. The canonical map
$$
\Gamma(X, \mathcal{F}) \longrightarrow \check{H}^0(X, \mathcal{F})
$$
is an isomorphism.
\end{lemma}

\begin{proof}
For every open covering $\mathcal{U}$ of $X$, the canonical map
$\Gamma(X, \mathcal{F}) \to \check{H}^0(\mathcal{U}, \mathcal{F})$
is an isomorphism by Lemma \ref{lemma-cech-h0}. Under these isomorphisms
the transition maps are the identity on $\Gamma(X, \mathcal{F})$.
Taking the colimit proves the lemma.
\end{proof}

\begin{lemma}
\label{lemma-cech-cohomology-additive}
\begin{reference}
\cite[Chapter I, \S 3, no. 23, pp. 215--216]{FAC}
\end{reference}
Let $(X, \mathcal{A})$ be a ringed space. For every $p \geq 0$, the rule
$$
\mathcal{F} \longmapsto \check{H}^p(X, \mathcal{F})
$$
defines an additive functor
$$
\textit{Mod}(\mathcal{A})
\longrightarrow
\text{Mod}_{\Gamma(X, \mathcal{A})}.
$$
In particular, {\v C}ech cohomology commutes with finite direct sums.
For $p = 0$, the isomorphism of Lemma \ref{lemma-cech-h0-space} is
functorial and $\Gamma(X, \mathcal{A})$-linear.
\end{lemma}

\begin{proof}
A morphism $\mathcal{F} \to \mathcal{G}$ induces a termwise map of
{\v C}ech complexes for every open covering. These maps commute with the
differentials and with the maps induced by refinements. Their construction
commutes with addition, composition, and identity maps. Passing first to
cohomology and then to the colimit gives an additive functor.

Every term of the {\v C}ech complex of $\mathcal{F}$ is naturally a
$\Gamma(X, \mathcal{A})$-module by restriction of global sections of
$\mathcal{A}$. The differentials, refinement maps, and maps induced by
morphisms of $\mathcal{A}$-modules are linear. This proves the module and
linearity assertions. The final assertion follows from the construction
in the proof of Lemma \ref{lemma-cech-h0-space}.
\end{proof}

\medskip\noindent
It turns out that the maps $\gamma$ defined above are compatible with
the map to cohomology, in other words, the composition
$$
\check{H}^p(\mathcal{V}, \mathcal{F}) \to
\check{H}^p(\mathcal{U}, \mathcal{F})
\xrightarrow{\text{Lemma \ref{lemma-cech-cohomology}}}
H^p(X, \mathcal{F})
$$
is the canonical map from the first group to cohomology of
Lemma \ref{lemma-cech-cohomology}.
In the lemma below we will prove this in a slightly more general
setting. A consequence is that we obtain a well defined map
\begin{equation}
\label{equation-cech-to-cohomology}
\check{H}^p(X, \mathcal{F}) =
\colim_\mathcal{U} \check{H}^p(\mathcal{U}, \mathcal{F})
\longrightarrow
H^p(X, \mathcal{F})
\end{equation}
from {\v C}ech cohomology to cohomology.

\begin{lemma}
\label{lemma-functoriality-cech}
Let $f : X \to Y$ be a morphism of ringed spaces.
Let $\varphi : f^*\mathcal{G} \to \mathcal{F}$ be an $f$-map
from an $\mathcal{O}_Y$-module $\mathcal{G}$ to an
$\mathcal{O}_X$-module $\mathcal{F}$.
Let $\mathcal{U} : X = \bigcup_{i \in I} U_i$ and
$\mathcal{V} : Y = \bigcup_{j \in J} V_j$ be open coverings.
Assume that $\mathcal{U}$ is a refinement of
$f^{-1}\mathcal{V} : X = \bigcup_{j \in J} f^{-1}(V_j)$.
In this case there exists a commutative diagram
$$
\xymatrix{
\check{\mathcal{C}}^\bullet(\mathcal{U}, \mathcal{F}) \ar[r] &
R\Gamma(X, \mathcal{F}) \\
\check{\mathcal{C}}^\bullet(\mathcal{V}, \mathcal{G}) \ar[r]
\ar[u]^\gamma &
R\Gamma(Y, \mathcal{G}) \ar[u]
}
$$
in $D^{+}(\mathcal{O}_X(X))$ with horizontal arrows given by
Lemma \ref{lemma-cech-cohomology} and right vertical arrow by
(\ref{equation-functorial-derived}).
In particular we get commutative diagrams of cohomology groups
$$
\xymatrix{
\check{H}^p(\mathcal{U}, \mathcal{F}) \ar[r] &
H^p(X, \mathcal{F}) \\
\check{H}^p(\mathcal{V}, \mathcal{G}) \ar[r]
\ar[u]^\gamma &
H^p(Y, \mathcal{G}) \ar[u]
}
$$
where the right vertical arrow is (\ref{equation-functorial})
\end{lemma}

\begin{proof}
We first define the left vertical arrow. Namely, choose a map
$c : I \to J$ such that $U_i \subset f^{-1}(V_{c(i)})$ for all
$i \in I$. In degree $p$ we define the map by the rule
$$
\gamma(s)_{i_0 \ldots i_p} = \varphi(s)_{c(i_0) \ldots c(i_p)}
$$
This makes sense because $\varphi$ does indeed induce maps
$\mathcal{G}(V_{c(i_0) \ldots c(i_p)}) \to \mathcal{F}(U_{i_0 \ldots i_p})$
by assumption. It is also clear that this defines a morphism of complexes.
Choose injective resolutions
$\mathcal{F} \to \mathcal{I}^\bullet$ on $X$ and
$\mathcal{G} \to J^\bullet$ on $Y$. According to
the proof of Lemma \ref{lemma-cech-cohomology} we introduce the double
complexes $A^{\bullet, \bullet}$ and $B^{\bullet, \bullet}$
with terms
$$
B^{p, q} = \check{\mathcal{C}}^p(\mathcal{V}, \mathcal{J}^q)
\quad
\text{and}
\quad
A^{p, q} = \check{\mathcal{C}}^p(\mathcal{U}, \mathcal{I}^q).
$$
As in Remark \ref{remark-explain-arrow} above we also choose an
injective resolution
$f_*\mathcal{I} \to (\mathcal{J}')^\bullet$ on $Y$ and a morphism
of complexes $\beta : \mathcal{J} \to (\mathcal{J}')^\bullet$
making (\ref{equation-choice}) commutes. We introduce some more
double complexes, namely $(B')^{\bullet, \bullet}$ and
$(B'')^{\bullet, \bullet}$ with
$$
(B')^{p, q} = \check{\mathcal{C}}^p(\mathcal{V}, (\mathcal{J}')^q)
\quad
\text{and}
\quad
(B'')^{p, q} = \check{\mathcal{C}}^p(\mathcal{V}, f_*\mathcal{I}^q).
$$
Note that there is an $f$-map of complexes from
$f_*\mathcal{I}^\bullet$ to $\mathcal{I}^\bullet$. Hence
it is clear that the same rule as above defines a morphism
of double complexes
$$
\gamma : (B'')^{\bullet, \bullet} \longrightarrow A^{\bullet, \bullet}.
$$
Consider the diagram of complexes
$$
\xymatrix{
\check{\mathcal{C}}^\bullet(\mathcal{U}, \mathcal{F})
\ar[r] &
\text{Tot}(A^{\bullet, \bullet}) & & &
\Gamma(X, \mathcal{I}^\bullet) \ar[lll]^{qis}
\ar@{=}[ddl]\\
\check{\mathcal{C}}^\bullet(\mathcal{V}, \mathcal{G})
\ar[r] \ar[u]^\gamma &
\text{Tot}(B^{\bullet, \bullet}) \ar[r]^\beta &
\text{Tot}((B')^{\bullet, \bullet}) &
\text{Tot}((B'')^{\bullet, \bullet}) \ar[l] \ar[llu]_{s\gamma} \\
& \Gamma(Y, \mathcal{J}^\bullet) \ar[u]^{qis} \ar[r]^\beta &
\Gamma(Y, (\mathcal{J}')^\bullet) \ar[u] &
\Gamma(Y, f_*\mathcal{I}^\bullet) \ar[u] \ar[l]_{qis}
}
$$
The two horizontal arrows with targets $\text{Tot}(A^{\bullet, \bullet})$ and
$\text{Tot}(B^{\bullet, \bullet})$
are the ones explained in Lemma \ref{lemma-cech-cohomology}.
The left upper shape (a pentagon) is commutative simply
because (\ref{equation-choice}) is commutative.
The two lower squares are trivially commutative.
It is also immediate from the definitions that the
right upper shape (a square) is commutative.
The result of the lemma now follows from the definitions
and the fact that going around the diagram on the outer sides
from $\check{\mathcal{C}}^\bullet(\mathcal{V}, \mathcal{G})$
to $\Gamma(X, \mathcal{I}^\bullet)$ either on top or on bottom
is the same (where you have to invert any quasi-isomorphisms along the way).
\end{proof}





\section{Acyclicity criteria by coverings}
\label{section-covering-acyclicity-criteria}

\noindent
We first extend the notation for the \v Cech complex to the coverings by
subsets which occur below. Let $A$ be a topological space, let
$\mathcal R=(R_i)_{i\in I}$ be an indexed covering of $A$ by subsets, and
let $\mathcal F$ be an abelian sheaf on $A$. Put
$$
R_{i_0\ldots i_p}=R_{i_0}\cap\ldots\cap R_{i_p}
$$
and define
\begin{align*}
C^p(\mathcal R,\underline H^q(\mathcal F))
&=
\prod_{(i_0,\ldots,i_p)\in I^{p+1}}
H^q(R_{i_0\ldots i_p},
\mathcal F|_{R_{i_0\ldots i_p}}),\\
H^p(\mathcal R,\underline H^q(\mathcal F))
&=
H^p(C^\bullet(\mathcal R,\underline H^q(\mathcal F))).
\end{align*}
The differential is the alternating sum of the restriction maps, as in
Equation (\ref{equation-d-cech}). For $q=0$ we abbreviate these groups to
$C^p(\mathcal R,\mathcal F)$ and $H^p(\mathcal R,\mathcal F)$.
If $A\subset X$ and $\mathcal F$ is a sheaf on $X$, this notation always
means that the construction is applied to $\mathcal F|_A$.

If $L$ is a finite collection of coverings of $A$, then
$\mathcal R^L$ denotes their common intersection covering: its members are
the nonempty intersections obtained by choosing one member of every
covering in $L$. We use the convention that $\mathcal R^\emptyset$ is the
one-member covering $(A)$. The nerve of $\mathcal R$ has dimension at most
$d$ if no $d+2$ distinct members of $\mathcal R$ have nonempty
intersection.

\begin{lemma}[Leray--Cartan spectral sequence for closed coverings]
\label{lemma-cech-spectral-sequence-closed-covering}
\begin{reference}
\cite[Chapter III, \S 3.8, Theorem 3.8.1 and the following
remark]{Tohoku}
\end{reference}
Let $A$ be a topological space and let
$\mathcal R=(R_i)_{i\in I}$ be a covering. Assume one of the following.
\begin{enumerate}
\item Every $R_i$ is open in $A$.
\item The set $I$ is finite and every $R_i$ is closed in $A$.
\item The space $A$ is paracompact and $(R_i)_{i\in I}$ is a locally
finite covering by closed subsets.
\end{enumerate}
For every abelian sheaf $\mathcal F$ on $A$ there is a functorial first
quadrant spectral sequence
$$
E_1^{p,q}=
\prod_{(i_0,\ldots,i_p)\in I^{p+1}}
H^q(R_{i_0\ldots i_p},
\mathcal F|_{R_{i_0\ldots i_p}}),
\qquad
E_2^{p,q}=H^p(\mathcal R,\underline H^q(\mathcal F))
\Longrightarrow H^{p+q}(A,\mathcal F).
$$
Its bottom edge maps are canonical maps
$$
H^n(\mathcal R,\mathcal F)\longrightarrow H^n(A,\mathcal F),
$$
and its other edge maps are the restriction maps
$$
H^n(A,\mathcal F)\longrightarrow
H^0(\mathcal R,\underline H^n(\mathcal F)).
$$

These spectral sequences and both edge maps are natural for refinement of
coverings and for trace on a subset, whenever the source and target
coverings satisfy one of the three conditions above. A refinement function
$c:J\to I$ with $V_j\subset R_{c(j)}$ induces the usual restriction map,
independently of $c$ from the $E_2$-page onward, and induces the identity on
$H^*(A,\mathcal F)$. If $B\subset A$, trace induces the restriction map
$H^*(A,\mathcal F)\to H^*(B,\mathcal F|_B)$ on the abutments.

If $n>0$ and
$$
H^q(R_{i_0\ldots i_p},
\mathcal F|_{R_{i_0\ldots i_p}})=0
\quad\text{for }0<q<n,
$$
for all $p$ and all tuples of indices, then the edge maps give an exact
sequence
$$
H^n(\mathcal R,\mathcal F)\longrightarrow H^n(A,\mathcal F)
\longrightarrow H^0(\mathcal R,\underline H^n(\mathcal F)).
$$
\end{lemma}

\begin{proof}
In the open case this is Lemma \ref{lemma-cech-spectral-sequence}, using the
covering resolution of Lemma \ref{lemma-covering-resolution}. We give the
closed-cover argument because it is needed below.

For a tuple $\mathbf i=(i_0,\ldots,i_p)$, let
$j_{\mathbf i}:R_{\mathbf i}\to A$ be the closed inclusion. Define
$$
\mathcal K^p(\mathcal R,\mathcal F)=
\prod_{\mathbf i\in I^{p+1}}
(j_{\mathbf i})_*(\mathcal F|_{R_{\mathbf i}})
$$
with the alternating restriction differential. There is an augmentation
$$
\mathcal F\longrightarrow
\mathcal K^\bullet(\mathcal R,\mathcal F).
$$
This is a quasi-isomorphism. Indeed, in the finite case, and locally in the
locally finite case, stalks commute with the displayed products. At a point
$x$, the nonzero factors form the augmented cochain complex of the simplex
on the nonempty set
$$
I_x=\{i\in I\mid x\in R_i\}
$$
with coefficients in $\mathcal F_x$. Choosing one element of $I_x$ gives
the usual contracting homotopy.

We recall a convenient construction of the hypercohomology spectral
sequence of this resolution. Apply the exact Godement resolution functor of
Lemma \ref{lemma-godement-resolution} to every term
$\mathcal K^p(\mathcal R,\mathcal F)$. For each Godement degree its
exactness turns the augmented covering resolution above into a flasque
resolution of the corresponding Godement term of $\mathcal F$. Hence one
spectral sequence of the resulting double complex of global sections
collapses to the complex of global sections of the Godement resolution of
$\mathcal F$, and its total cohomology is $H^*(A,\mathcal F)$. The other
spectral sequence has
$$
E_1^{p,q}=H^q(A,\mathcal K^p(\mathcal R,\mathcal F)).
$$

Suppose first that $I$ is finite. Direct image by a closed inclusion is
exact, it sends injectives to injectives, and its global sections are the
global sections on the closed subset. Finite products commute with
cohomology. Consequently
$$
H^q(A,\mathcal K^p(\mathcal R,\mathcal F))=
\prod_{\mathbf i\in I^{p+1}}
H^q(R_{\mathbf i},\mathcal F|_{R_{\mathbf i}}).
$$

Suppose next that $A$ is paracompact and the closed covering is locally
finite. Choose a flasque resolution
$\mathcal F\to\mathcal I^\bullet$. Every $\mathcal I^r$ is soft by
Lemma \ref{lemma-flasque-soft-family-supports}, and its restriction to a
closed subset is soft. For fixed $p$, the family of closed subsets
$R_{\mathbf i}$, indexed by $I^{p+1}$, is locally finite. It follows
directly from the extension property that
$\mathcal K^p(\mathcal R,\mathcal I^r)$ is soft: a section over a closed
subset is a locally finite family of sections over its intersections with
the $R_{\mathbf i}$, and the members of this family can be extended one by
one. These sheaves are acyclic by Lemma
\ref{lemma-family-supports-flasque-soft-acyclic}.

The functor $\mathcal K^p(\mathcal R,-)$ is exact in this case, as is seen
on stalks where only finitely many factors occur. Thus
$\mathcal K^p(\mathcal R,\mathcal I^\bullet)$ computes the cohomology of
$\mathcal K^p(\mathcal R,\mathcal F)$. Each closed subspace
$R_{\mathbf i}$ is paracompact, and
$\mathcal I^\bullet|_{R_{\mathbf i}}$ is a soft resolution. Since products
of abelian groups are exact, we again obtain
$$
H^q(A,\mathcal K^p(\mathcal R,\mathcal F))=
\prod_{\mathbf i\in I^{p+1}}
H^q(R_{\mathbf i},\mathcal F|_{R_{\mathbf i}}).
$$
This proves the displayed spectral sequence.

The augmentation identifies the bottom and vertical edge maps with the
maps in the statement. In particular, the vertical edge sends a class to
the compatible family of its restrictions to the $R_i$. A refinement
function gives a morphism of the covering resolutions by restriction.
Two refinement functions give homotopic morphisms by the usual prism
homotopy, so the induced maps agree from $E_2$ onward. The augmentation
shows that the map on the common abutment is the identity. Restricting all
terms to $B$ gives the trace map and the restriction map on the abutments.
This also proves all asserted naturality.

Finally, under the last vanishing hypothesis the terms with
$0<q<n$ cannot contribute in total degree $n$. Thus the filtration in that
degree gives the asserted exact edge sequence.
\end{proof}

\begin{lemma}[Coverwise local killing]
\label{lemma-covering-acyclicity-local-killing}
\begin{reference}
\cite[Chapter III, \S 3.9, Theorem 3.9.1, pp. 179--180]{Tohoku}
\end{reference}
Let $X$ be a topological space and let $\mathcal S$ be a nonempty
collection of subsets of $X$. For every $A\in\mathcal S$, let
$\mathfrak R(A)$ be a nonempty collection of indexed coverings of $A$.
Assume the following.
\begin{enumerate}
\item The members of every $\mathcal R\in\mathfrak R(A)$ and all their
nonempty finite intersections belong to $\mathcal S$.
\item If $\mathcal R,\mathcal R'\in\mathfrak R(A)$ and $B$ is a member of
$\mathcal R$, then the trace $\mathcal R'|_B$ belongs to
$\mathfrak R(B)$.
\item One of the following three alternatives holds: all members of
$\mathcal S$ are open in $X$; all members are closed and all coverings in
the collections $\mathfrak R(A)$ are finite; or all members are closed,
$X$ is paracompact, and all those coverings are locally finite.
\end{enumerate}
Let $\mathcal F$ be an abelian sheaf on $X$ and let $n\geq0$. Suppose
\begin{enumerate}
\item[$A(n)$] $H^i(\mathcal R,\mathcal F)=0$ for
$1\leq i\leq n$, every $A\in\mathcal S$, and every
$\mathcal R\in\mathfrak R(A)$.
\item[$B(n-1)$] For $A\in\mathcal S$, $1\leq i\leq n-1$, and
$c\in H^i(A,\mathcal F|_A)$, there is a finite
$L\subset\mathfrak R(A)$ such that $c$ restricts to zero on every member
of $\mathcal R^L$.
\end{enumerate}
Then
$$
H^i(A,\mathcal F|_A)=0
\quad\text{for }A\in\mathcal S\text{ and }1\leq i\leq n-1.
$$
Moreover, a class $c\in H^n(A,\mathcal F|_A)$ is zero if and only if
there is a finite $L\subset\mathfrak R(A)$ such that $c$ restricts to
zero on every member of $\mathcal R^L$.
\end{lemma}

\begin{proof}
We argue by induction on $n$. The assertion for $n=0$ follows from the
convention $\mathcal R^\emptyset=(A)$, or directly by checking germs.
Suppose $n>0$. Applying the induction hypothesis in degree $n-1$ and using
$B(n-1)$ shows that
$$
H^i(A,\mathcal F|_A)=0
\quad\text{for }1\leq i\leq n-1
$$
and every $A\in\mathcal S$.

It remains to prove the characterization in degree $n$. The forward
implication follows by taking $L=\emptyset$. For the converse, suppose
$c\in H^n(A,\mathcal F|_A)$ is killed on the members of
$\mathcal R^L$. We use induction on $|L|$. The case $|L|=0$ is immediate.
If $L=\{\mathcal R\}$, all the nonzero intersections of members of
$\mathcal R$ belong to $\mathcal S$, so their cohomology vanishes in
degrees $1,\ldots,n-1$. Lemma
\ref{lemma-cech-spectral-sequence-closed-covering} therefore gives the exact
sequence
$$
H^n(\mathcal R,\mathcal F)\longrightarrow
H^n(A,\mathcal F|_A)\longrightarrow
H^0(\mathcal R,\underline H^n(\mathcal F)).
$$
The first group is zero by $A(n)$, and the image of $c$ in the last group
is its family of restrictions to the members of $\mathcal R$. Hence $c=0$.

Now let $L=\{\mathcal R^1,\ldots,\mathcal R^k\}$ with $k>1$. For a member
$B$ of $\mathcal R^1$, the trace coverings
$\mathcal R^2|_B,\ldots,\mathcal R^k|_B$ belong to $\mathfrak R(B)$.
Trace naturality in the preceding lemma says that $c|_B$ is killed on their
common intersection covering. The induction on $k$ gives $c|_B=0$ for
every member $B$ of $\mathcal R^1$. The case $|L|=1$ now gives $c=0$.
\end{proof}

\begin{lemma}[Equivalent covering acyclicity conditions]
\label{lemma-covering-acyclicity-conditions}
\begin{reference}
\cite[Chapter III, \S 3.9, Corollary 1 to Theorem 3.9.1,
p. 179]{Tohoku}
\end{reference}
In the situation of Lemma
\ref{lemma-covering-acyclicity-local-killing}, the following are equivalent.
\begin{enumerate}
\item $H^i(A,\mathcal F|_A)=0$ for every $A\in\mathcal S$ and
$1\leq i\leq n$.
\item Conditions $A(n)$ and $B(n)$ hold.
\end{enumerate}
\end{lemma}

\begin{proof}
Assume (1). Then $B(n)$ is immediate. In the spectral sequence of Lemma
\ref{lemma-cech-spectral-sequence-closed-covering}, every finite
intersection of members of a covering belongs to $\mathcal S$. Thus, in
total degrees at most $n$, only the row $q=0$ can be nonzero, and the bottom
edge maps identify
$$
H^i(\mathcal R,\mathcal F)=H^i(A,\mathcal F|_A)
\quad\text{for }0\leq i\leq n.
$$
This proves $A(n)$.

Conversely, apply Lemma \ref{lemma-covering-acyclicity-local-killing} using
$A(n)$ and $B(n-1)$. It gives the vanishing through degree $n-1$ and says
that a degree $n$ class is zero exactly when it is killed on a common
intersection covering. Condition $B(n)$ kills every such class.
\end{proof}

\begin{lemma}[Special coverings compute acyclicity]
\label{lemma-special-coverings-compute-acyclicity}
\begin{reference}
\cite[Chapter III, \S 3.9, Corollary 2 to Theorem 3.9.1,
p. 179]{Tohoku}
\end{reference}
In the situation of Lemma
\ref{lemma-covering-acyclicity-local-killing}, assume in addition that for
every $A\in\mathcal S$ and every open covering $\mathcal U$ of $A$, there
is a finite $L\subset\mathfrak R(A)$ such that $\mathcal R^L$ refines
$\mathcal U$. Then the following are equivalent.
\begin{enumerate}
\item $H^i(A,\mathcal F|_A)=0$ for every $A\in\mathcal S$ and
$1\leq i\leq n$.
\item $H^i(\mathcal R,\mathcal F)=0$ for every $A\in\mathcal S$,
$\mathcal R\in\mathfrak R(A)$, and $1\leq i\leq n$.
\end{enumerate}
\end{lemma}

\begin{proof}
Lemma \ref{lemma-kill-cohomology-class-on-covering}, applied on the
ringed space $(A,\underline{\mathbf Z}_A)$, gives an open covering on which
any positive-degree class becomes zero. By hypothesis a common intersection
covering $\mathcal R^L$ refines that open covering. Thus $B(m)$ holds for
every $m$. The result is now Lemma
\ref{lemma-covering-acyclicity-conditions}.
\end{proof}

\begin{lemma}[Bounded nerves and acyclicity]
\label{lemma-bounded-nerve-special-coverings-acyclic}
\begin{reference}
\cite[Chapter III, \S 3.9, Corollary 3 to Theorem 3.9.1,
p. 180]{Tohoku}
\end{reference}
In the situation of Lemma
\ref{lemma-special-coverings-compute-acyclicity}, suppose there is a
$d\geq0$ such that every $\mathcal R\in\mathfrak R(A)$ has nerve of
dimension at most $d$. If the equivalent conditions of that lemma hold for
$1\leq i\leq d$, then
$$
H^i(A,\mathcal F|_A)=0
\quad\text{for every }A\in\mathcal S\text{ and every }i>0.
$$
\end{lemma}

\begin{proof}
Choose a total ordering on the index set of $\mathcal R$. The usual
projection to alternating cochains and its simplicial homotopy inverse, as
in Lemma \ref{lemma-alternating-usual}, use only restriction maps and hence
apply to the subset covering considered here. The alternating complex is
zero in degrees greater than $d$. Consequently
$H^i(\mathcal R,\mathcal F)=0$ for $i>d$. By hypothesis it also vanishes
for $1\leq i\leq d$, so $A(m)$ holds for every $m$. The proof of Lemma
\ref{lemma-special-coverings-compute-acyclicity} shows that $B(m)$ holds for
every $m$. Applying Lemma \ref{lemma-covering-acyclicity-conditions} for
arbitrary $m$ proves the assertion.
\end{proof}


\begin{lemma}[Stars of finite closed coverings]
\label{lemma-finite-closed-covering-star-neighbourhood}
\begin{reference}
\cite[Chapter III, \S 3.9, p. 180]{Tohoku}
\end{reference}
Let $X$ be a topological space and let $\mathcal{R}$ be a finite covering
of $X$ by closed subsets. For $x \in X$ set
$$
E_x(\mathcal{R}) =
\bigcup_{\substack{A \in \mathcal{R} \\ x \in A}} A
\quad\text{and}\quad
O_x(\mathcal{R}) =
X \setminus
\bigcup_{\substack{A \in \mathcal{R} \\ x \notin A}} A.
$$
Then
$$
x \in O_x(\mathcal{R}) \subset E_x(\mathcal{R}),
$$
$O_x(\mathcal{R})$ is an open neighbourhood of $x$, and
$$
y \in O_x(\mathcal{R})
\quad\Longrightarrow\quad
E_y(\mathcal{R}) \subset E_x(\mathcal{R}).
$$
\end{lemma}

\begin{proof}
The union occurring in the definition of $O_x(\mathcal{R})$ is a finite
union of closed subsets. Thus $O_x(\mathcal{R})$ is open. None of those
closed subsets contains $x$, and hence $x \in O_x(\mathcal{R})$.

Let $y \in O_x(\mathcal{R})$. If $A \in \mathcal{R}$ contains $y$, then
$A$ contains $x$; otherwise $y$ would belong to the union whose complement
defines $O_x(\mathcal{R})$. Since $\mathcal{R}$ covers $X$, this first shows
that $y \in E_x(\mathcal{R})$. If $z \in E_y(\mathcal{R})$, then some
$A \in \mathcal{R}$ contains both $y$ and $z$. The preceding observation
gives $x \in A$, and hence $z \in E_x(\mathcal{R})$.
\end{proof}

\begin{lemma}[Separating finite closed coverings]
\label{lemma-separating-finite-closed-coverings-refine-open-covering}
\begin{reference}
\cite[Chapter III, \S 3.9, Proposition 3.9.2,
pp. 180--181]{Tohoku}
\end{reference}
Let $X$ be a quasi-compact topological space and let $\mathfrak R$ be a
nonempty collection of finite coverings of $X$ by closed subsets. Assume
that for every pair of distinct points $x, y \in X$ there is an
$\mathcal R \in \mathfrak R$ such that no member of $\mathcal R$ contains
both $x$ and $y$. Then for every open covering $\mathcal U$ of $X$ there
is a nonempty finite subset $L \subset \mathfrak R$ such that
$$
\mathcal R^L =
\left\{\bigcap_{\mathcal R \in L} A_{\mathcal R} \ \middle|\
A_{\mathcal R} \in \mathcal R,\
\bigcap_{\mathcal R \in L} A_{\mathcal R} \neq \emptyset\right\}
$$
is a covering of $X$ which refines $\mathcal U$.

Suppose moreover that $\mathcal S$ is a collection of closed subsets of
$X$ containing $X$, and that for every $A \in \mathcal S$ a collection
$\mathfrak R(A)$ of coverings of $A$ is given. If
$\mathfrak R = \mathfrak R(X)$ satisfies the hypothesis above and
$$
\mathcal R|_A =
\{A \cap B \mid B \in \mathcal R,\ A \cap B \neq \emptyset\}
\in \mathfrak R(A)
$$
for every $A \in \mathcal S$ and $\mathcal R \in \mathfrak R(X)$, then
finite common refinements of coverings in $\mathfrak R(A)$ are cofinal
among the open coverings of $A$, for every $A \in \mathcal S$.
\end{lemma}

\begin{proof}
If $X$ is empty, choose any $\mathcal R_1 \in \mathfrak R$; after discarding
empty members, $\mathcal R^{\{\mathcal R_1\}}$ is the empty covering of
$X$, and the assertion is immediate. Assume $X$ is nonempty. For every
$x \in X$, choose
$U_x \in \mathcal U$ containing $x$. Using the notation of
Lemma \ref{lemma-finite-closed-covering-star-neighbourhood}, the separation
hypothesis gives
$$
\bigcap_{\mathcal R \in \mathfrak R} E_x(\mathcal R) = \{x\}.
$$
Each $E_x(\mathcal R)$ is closed. Hence the opens
$X \setminus E_x(\mathcal R)$ cover the closed subset
$X \setminus U_x$, which is quasi-compact by Topology, Lemma
\ref{topology-lemma-closed-in-quasi-compact}. We can therefore choose
finitely many members of $\mathfrak R$ whose set $L_x$ satisfies
$$
\bigcap_{\mathcal R \in L_x} E_x(\mathcal R) \subset U_x.
$$
Adding an arbitrary member of $\mathfrak R$ if necessary, we may assume
$L_x$ is nonempty.

Let $\mathcal R_x$ be the finite closed covering whose nonempty members are
the intersections of one member of each $\mathcal R \in L_x$. Directly
from the definition,
$$
E_x(\mathcal R_x) =
\bigcap_{\mathcal R \in L_x} E_x(\mathcal R) \subset U_x.
$$
The preceding lemma shows that $O_x(\mathcal R_x)$ is an open neighbourhood
of $x$ contained in $E_x(\mathcal R_x)$. By quasi-compactness there are
$x_1, \ldots, x_m \in X$ such that the $O_{x_i}(\mathcal R_{x_i})$ cover
$X$.

Set $L = L_{x_1} \cup \ldots \cup L_{x_m}$ and let
$\mathcal V = \mathcal R^L$. Let $A \in \mathcal V$, choose a presentation
$$
A = \bigcap_{\mathcal R \in L} A_{\mathcal R},
\qquad A_{\mathcal R} \in \mathcal R,
$$
and choose $a \in A$. Choose $i$ such that
$a \in O_{x_i}(\mathcal R_{x_i})$. Then
$$
B = \bigcap_{\mathcal R \in L_{x_i}} A_{\mathcal R}
$$
is a member of $\mathcal R_{x_i}$ containing $a$, and $A \subset B$.
By the preceding lemma,
$$
A \subset B \subset E_a(\mathcal R_{x_i})
\subset E_{x_i}(\mathcal R_{x_i}) \subset U_{x_i}.
$$
Thus $\mathcal V$ refines $\mathcal U$.

For the final assertion, a closed subset $A$ of the quasi-compact space
$X$ is quasi-compact. The collection of trace coverings
$\mathcal R|_A$, for $\mathcal R \in \mathfrak R(X)$, is a nonempty
point-separating collection of finite closed coverings of $A$. Applying
the first assertion to $A$ gives the desired common refinement, and every
covering used belongs to $\mathfrak R(A)$ by assumption.
\end{proof}

\section{Cohomology on Hausdorff quasi-compact spaces}
\label{section-cohomology-LC}

\noindent
For such a space {\v C}ech cohomology agrees with cohomology.

\begin{lemma}
\label{lemma-cech-always}
\begin{reference}
\cite[Chapter III, Corollary 3.8.1(3), p. 176]{Tohoku}
\end{reference}
Let $X$ be a topological space. Let $\mathcal{F}$ be an abelian sheaf. Then
the map $\check{H}^1(X, \mathcal{F}) \to H^1(X, \mathcal{F})$ defined
in (\ref{equation-cech-to-cohomology}) is an isomorphism, and the map
$$
\check{H}^2(X, \mathcal{F}) \longrightarrow H^2(X, \mathcal{F})
$$
is injective.
\end{lemma}

\begin{proof}
Let $\mathcal{U}$ be an open covering of $X$.
By Lemma \ref{lemma-cech-spectral-sequence}
there is an exact sequence
$$
0 \to \check{H}^1(\mathcal{U}, \mathcal{F}) \to H^1(X, \mathcal{F})
\to \check{H}^0(\mathcal{U}, \underline{H}^1(\mathcal{F}))
\to \check{H}^2(\mathcal{U}, \mathcal{F}) \to H^2(X, \mathcal{F})
$$
Thus the map in degree $1$ is injective. To show surjectivity it suffices
to show that
any element of $\check{H}^0(\mathcal{U}, \underline{H}^1(\mathcal{F}))$
maps to zero after replacing $\mathcal{U}$ by a refinement.
This is immediate from the definitions and the fact that
$\underline{H}^1(\mathcal{F})$ is a presheaf of abelian groups
whose sheafification is zero by locality of cohomology, see
Lemma \ref{lemma-kill-cohomology-class-on-covering}.

The same argument proves injectivity in degree $2$. Indeed, a class in
$\check{H}^2(X, \mathcal{F})$ is represented on some covering
$\mathcal{U}$. If it maps to zero in $H^2(X, \mathcal{F})$, then by the
exact sequence above it comes from an element of
$\check{H}^0(\mathcal{U}, \underline{H}^1(\mathcal{F}))$. After replacing
$\mathcal{U}$ by a refinement, this element maps to zero by the preceding
paragraph. Refinement maps give morphisms of the filtered \v Cech double
complexes and hence morphisms of their spectral sequences by
Homology, Lemma
\ref{homology-lemma-spectral-sequence-filtered-complex-functorial}.
Hence the original class is zero in the colimit defining
$\check{H}^2(X, \mathcal{F})$.
\end{proof}

\begin{lemma}
\label{lemma-cech-cohomology-not-delta-functor}
\begin{reference}
\cite[Chapter III, Section 3.8, pp. 174 and 177--179]{Tohoku}
\end{reference}
There exist a finite topological space $X$ and a short exact sequence
$$
0 \to \mathcal{F} \to \mathcal{G} \to \mathcal{H} \to 0
$$
of abelian sheaves on $X$ such that
$$
\check{H}^1(X, \mathcal{G}) = 0, \qquad
\check{H}^1(X, \mathcal{H}) \cong \mathbf{Z}, \qquad
\check{H}^2(X, \mathcal{F}) = 0,
$$
whereas $H^2(X, \mathcal{F}) \cong \mathbf{Z}$. In particular, the
additive functors $\check{H}^p(X, -)$ do not in general form a
cohomological delta functor, and the map
$$
\check{H}^2(X, \mathcal{F}) \longrightarrow H^2(X, \mathcal{F})
$$
need not be surjective.
\end{lemma}

\begin{proof}
Let
$$
X = \{\eta, a, b, x_1, x_2\}
$$
be the finite space whose minimum open neighbourhoods of these five points
are, in the displayed order,
$$
\{\eta\},\quad \{\eta, a\},\quad \{\eta, b\},\quad
X \setminus \{x_2\},\quad X \setminus \{x_1\}.
$$
Every nonempty open contains $\eta$, so $X$ is irreducible. The closed
subsets
$$
Y_1 = \{a, x_1, x_2\}, \qquad
Y_2 = \{b, x_1, x_2\}
$$
are irreducible with generic points $a$ and $b$, respectively, and their
intersection $D = \{x_1, x_2\}$ is discrete. Put
$Y = Y_1 \cup Y_2 = X \setminus \{\eta\}$, and denote the closed
inclusions by
$$
i : Y \longrightarrow X, \qquad
i_r : Y_r \longrightarrow Y, \qquad
k : D \longrightarrow Y.
$$
Set
$$
\mathcal{G} = \underline{\mathbf{Z}}_X, \qquad
\mathcal{H} = i_*\underline{\mathbf{Z}}_Y, \qquad
\mathcal{F} = \ker(\mathcal{G} \longrightarrow \mathcal{H}).
$$
The resulting sequence $0 \to \mathcal{F} \to \mathcal{G} \to
\mathcal{H} \to 0$ is exact; this can be checked on stalks, or follows
from the canonical exact sequence for the closed complement $Y$.

Consider the covering $\mathcal{U} = \{U_1, U_2\}$, where
$$
U_1 = X \setminus \{x_2\}, \qquad
U_2 = X \setminus \{x_1\}.
$$
This covering refines every open covering of $X$: a member containing
$x_r$ contains its minimum open neighbourhood $U_r$. Hence its class is
final in the refinement preorder of
Definition \ref{definition-cech-cohomology-space}, and fixed-cover Cech
cohomology for $\mathcal{U}$ computes $\check{H}^q(X, -)$.

We have $U_1 \cap U_2 = \{\eta, a, b\}$. Every one of
$U_1$, $U_2$, and $U_1 \cap U_2$ meets $Y$, and restriction of a constant
integer to that intersection is injective. Thus $\mathcal{F}$ has no
sections on any of these opens, so its Cech complex for $\mathcal{U}$ is
zero. In particular, $\check{H}^2(X, \mathcal{F}) = 0$.

For $\mathcal{G}$ the alternating Cech complex in degrees zero and one is
$$
\mathbf{Z}^2 \longrightarrow \mathbf{Z}, \qquad
(m, n) \longmapsto n - m,
$$
and therefore $\check{H}^1(X, \mathcal{G}) = 0$. For $\mathcal{H}$ it is
$$
\mathbf{Z}^2 \longrightarrow \mathbf{Z}^2, \qquad
(m, n) \longmapsto (n - m, n - m),
$$
because $Y \cap (U_1 \cap U_2) = \{a, b\}$ is discrete. Consequently,
$\check{H}^1(X, \mathcal{H}) \cong \mathbf{Z}$. These alternating
computations also compute the usual Cech groups by
Lemma \ref{lemma-alternating-usual}.

It remains to compute derived-functor cohomology. On $Y$ there is a
stalkwise exact sequence
$$
0 \to \underline{\mathbf{Z}}_Y \to
(i_1)_*\underline{\mathbf{Z}}_{Y_1} \oplus
(i_2)_*\underline{\mathbf{Z}}_{Y_2} \to
k_*\underline{\mathbf{Z}}_D \to 0,
$$
where the last arrow is the difference of the two restrictions. The last
two terms are flasque: the constant sheaves on the irreducible spaces
$Y_1$, $Y_2$ are flasque, the constant sheaf on the discrete space $D$ is
flasque, and direct image along a closed inclusion preserves flasqueness.
Thus this is a flasque resolution of $\underline{\mathbf{Z}}_Y$. On global
sections its last arrow is
$$
\mathbf{Z}^2 \longrightarrow \mathbf{Z}^2, \qquad
(m, n) \longmapsto (m - n, m - n).
$$
It follows that
$$
H^1(Y, \underline{\mathbf{Z}}_Y) \cong \mathbf{Z}.
$$
By Lemma \ref{lemma-cohomology-and-closed-immersions}, this is
$H^1(X, \mathcal{H})$. Since $X$ is irreducible,
Lemma \ref{lemma-irreducible-constant-cohomology-zero} gives
$H^q(X, \mathcal{G}) = 0$ for $q > 0$. The long exact cohomology sequence
therefore gives
$$
H^2(X, \mathcal{F}) \cong H^1(X, \mathcal{H}) \cong \mathbf{Z}.
$$

If the functors $\check{H}^p(X, -)$ formed a cohomological delta functor,
the short exact sequence above would give an exact segment
$$
\check{H}^1(X, \mathcal{G}) \longrightarrow
\check{H}^1(X, \mathcal{H}) \longrightarrow
\check{H}^2(X, \mathcal{F}).
$$
The computed segment is $0 \to \mathbf{Z} \to 0$, which is not exact.
The comparison assertion follows from the two degree-two computations.
\end{proof}

\begin{remark}
\label{remark-FAC-cech-low-degree-exact-sequence}
\begin{reference}
\cite[Chapter I, \S 3, no. 24, Propositions 5 and 6 and
Corollaries 1 and 2, pp. 216--217]{FAC}
The affine-cover specialization and its consequences are
\cite[Chapter II, \S 3, no. 47, Theorem 5 and Corollaries 1 and 2,
pp. 239--240]{FAC}.
\end{reference}
Let
$$
0 \to \mathcal{F} \to \mathcal{G} \to \mathcal{H} \to 0
$$
be a short exact sequence of sheaves. For a fixed open covering
$\mathcal{U}$, the corresponding sequence of {\v C}ech complexes need
not be exact on the right. The source replaces
$\check{\mathcal{C}}^\bullet(\mathcal{U}, \mathcal{H})$ by the image
$\mathcal{C}_0^\bullet(\mathcal{U}, \mathcal{H})$ of
$\check{\mathcal{C}}^\bullet(\mathcal{U}, \mathcal{G})$. Write
$H_0^q(\mathcal{U}, \mathcal{H})$ for the cohomology of this image
complex and $H_0^q(X, \mathcal{H})$ for its colimit over refinements.
The resulting short exact sequences of complexes are compatible with
refinements.

In no. 47, vanishing for a coherent kernel on every finite intersection of
an affine-open covering makes the image complex equal to the ordinary Cech
complex of $\mathcal{H}$. The first corollary passes this equality to the
refinement colimit, and the second obtains the resulting long exact sequence.
This records the source's auxiliary construction without making the groups
$H_0^q$ a second modern cohomology theory.

The local lifting property for a surjective map of sheaves, see
Sheaves, Definition \ref{sheaves-definition-injective-surjective},
shows after taking the colimit over coverings that
$$
H_0^0(X, \mathcal{H}) \longrightarrow \check{H}^0(X, \mathcal{H})
$$
is an isomorphism and that
$$
H_0^1(X, \mathcal{H}) \longrightarrow \check{H}^1(X, \mathcal{H})
$$
is injective.
Thus the long exact sequence for the image complexes gives the exact
sequence
$$
0 \to \Gamma(X, \mathcal{F}) \to \Gamma(X, \mathcal{G}) \to
\Gamma(X, \mathcal{H}) \to \check{H}^1(X, \mathcal{F}) \to
\check{H}^1(X, \mathcal{G}) \to \check{H}^1(X, \mathcal{H}).
$$
By Lemmas \ref{lemma-cech-h0-space} and \ref{lemma-cech-always}, its terms
identify with degree zero and one derived-functor cohomology. In the current
framework the same low-degree exactness, and its continuation in every
degree, follow from the universal delta-functor construction discussed in
Section \ref{section-derived-functors}. In particular,
$H^1(X, \mathcal{F}) = 0$ implies that
$\Gamma(X, \mathcal{G}) \to \Gamma(X, \mathcal{H})$ is surjective.
\end{remark}

\begin{lemma}
\label{lemma-cech-paracompact}
\begin{reference}
\cite[Chapter I, \S 3, no. 25, Proposition 7 and its Corollary,
pp. 217--218]{FAC}
\end{reference}
Let $X$ be a paracompact topological space, see
Topology, Definition \ref{topology-definition-paracompact}. Let
$\mathcal{F}$ be an abelian sheaf on $X$. Then
the map $\check{H}^n(X, \mathcal{F}) \to H^n(X, \mathcal{F})$ defined
in (\ref{equation-cech-to-cohomology}) is an isomorphism for
all $n$.
\end{lemma}

\begin{proof}
We already know that $\check{H}^n(X, -) \to H^n(X, -)$
is an isomorphism of functors for $n = 0, 1$, see
Lemmas \ref{lemma-cech-h0-space} and \ref{lemma-cech-always}.
The functors $H^n(X, -)$ form a universal $\delta$-functor, see
Derived Categories, Lemma \ref{derived-lemma-higher-derived-functors}.
If we show that $\check{H}^n(X, -)$ forms a universal $\delta$-functor
and that $\check{H}^n(X, -) \to H^n(X, -)$ is compatible with boundary
maps, then the map will automatically be an isomorphism by uniqueness
of universal $\delta$-functors, see
Homology, Lemma \ref{homology-lemma-uniqueness-universal-delta-functor}.

\medskip\noindent
Let $0 \to \mathcal{F} \to \mathcal{G} \to \mathcal{H} \to 0$
be a short exact sequence of abelian sheaves on $X$.
Let $\mathcal{U} : X = \bigcup_{i \in I} U_i$ be an open covering.
This gives a complex of complexes
$$
0 \to \check{\mathcal{C}}^\bullet(\mathcal{U}, \mathcal{F}) \to
\check{\mathcal{C}}^\bullet(\mathcal{U}, \mathcal{G}) \to
\check{\mathcal{C}}^\bullet(\mathcal{U}, \mathcal{H}) \to 0
$$
which is in general not exact on the right. The sequence defines
the maps
$$
\check{H}^n(\mathcal{U}, \mathcal{F}) \to
\check{H}^n(\mathcal{U}, \mathcal{G}) \to
\check{H}^n(\mathcal{U}, \mathcal{H})
$$
but isn't good enough to define a boundary operator
$\delta : \check{H}^n(\mathcal{U}, \mathcal{H}) \to
\check{H}^{n + 1}(\mathcal{U}, \mathcal{F})$. Indeed
such a thing will not exist in general. However, given an
element $\overline{h} \in \check{H}^n(\mathcal{U}, \mathcal{H})$
which is the cohomology class of a cocycle
$h = (h_{i_0 \ldots i_n})$
we can choose open coverings
$$
U_{i_0 \ldots i_n} = \bigcup W_{i_0 \ldots i_n, k}
$$
such that $h_{i_0 \ldots i_n}|_{W_{i_0 \ldots i_n, k}}$
lifts to a section of $\mathcal{G}$ over $W_{i_0 \ldots i_n, k}$.
By Topology, Lemma \ref{topology-lemma-paracompact-refine-covering}
(this is where we use paracompactness)
we can choose an open covering $\mathcal{V} : X = \bigcup_{j \in J} V_j$
and $\alpha : J \to I$ such that $V_j \subset U_{\alpha(j)}$
(it is a refinement) and such that for all $j_0, \ldots, j_n \in J$
with $V_{j_0 \ldots j_n}$ nonempty there is a $k$ such that
$V_{j_0 \ldots j_n} \subset W_{\alpha(j_0) \ldots \alpha(j_n), k}$.
On an empty intersection the unique section lifts automatically.
We obtain maps of complexes
$$
\xymatrix{
0 \ar[r] &
\check{\mathcal{C}}^\bullet(\mathcal{U}, \mathcal{F}) \ar[d] \ar[r] &
\check{\mathcal{C}}^\bullet(\mathcal{U}, \mathcal{G}) \ar[d] \ar[r] &
\check{\mathcal{C}}^\bullet(\mathcal{U}, \mathcal{H}) \ar[d] \ar[r] &
0 \\
0 \ar[r] &
\check{\mathcal{C}}^\bullet(\mathcal{V}, \mathcal{F}) \ar[r] &
\check{\mathcal{C}}^\bullet(\mathcal{V}, \mathcal{G}) \ar[r] &
\check{\mathcal{C}}^\bullet(\mathcal{V}, \mathcal{H}) \ar[r] &
0
}
$$
In fact, the vertical arrows are the maps of complexes used
to define the transition maps between the {\v C}ech cohomology groups.
Our choice of refinement shows that we may choose
$$
g_{j_0 \ldots j_n} \in
\mathcal{G}(V_{j_0 \ldots j_n}),\quad
g_{j_0 \ldots j_n} \longmapsto
h_{\alpha(j_0) \ldots \alpha(j_n)}|_{V_{j_0 \ldots j_n}}
$$
The cochain $g = (g_{j_0 \ldots j_n})$ is not a cocycle
in general but we know that its {\v C}ech boundary $\text{d}(g)$
maps to zero in $\check{\mathcal{C}}^{n + 1}(\mathcal{V}, \mathcal{H})$
(by the commutative diagram above and the fact that $h$ is a cocycle).
Hence $\text{d}(g)$ is a cocycle in
$\check{\mathcal{C}}^\bullet(\mathcal{V}, \mathcal{F})$.
This allows us to define
$$
\delta(\overline{h}) = \text{class of }\text{d}(g)\text{ in }
\check{H}^{n + 1}(\mathcal{V}, \mathcal{F})
$$
Now, given an element $\xi \in \check{H}^n(X, \mathcal{H})$
we choose an open covering $\mathcal{U}$ and an element
$\overline{h} \in \check{H}^n(\mathcal{U}, \mathcal{H})$
mapping to $\xi$ in the colimit defining {\v C}ech cohomology.
Then we choose $\mathcal{V}$ and $g$ as above and set
$\delta(\xi)$ equal to the image of $\delta(\overline{h})$
in $\check{H}^{n + 1}(X, \mathcal{F})$.
At this point a lot of properties have to be checked, all of which
are straightforward. For example, we need to check that our construction
is independent of the choice of
$\mathcal{U}, \overline{h}, \mathcal{V}, \alpha : J \to I, g$.
The class of $\text{d}(g)$ is independent of the choice of the lifts
$g_{i_0 \ldots i_n}$ because the difference will be a coboundary.
Independence of $\alpha$ holds\footnote{This is an important
check because the nonuniqueness of $\alpha$ is the only thing preventing
us from taking the colimit of {\v C}ech complexes over all open
coverings of $X$ to get a short exact sequence of complexes computing
{\v C}ech cohomology.}
because a different choice
of $\alpha$ determines homotopic vertical maps of complexes
in the diagram above, see Section \ref{section-refinements-cech}.
For the other choices we use that given a finite collection
of open coverings of $X$ we can always find an open covering refining all
of them. We also need to check additivity which is shown in the same manner.
Finally, we need to check that the maps
$\check{H}^n(X, -) \to H^n(X, -)$ are compatible
with boundary maps. To do this we choose injective
resolutions
$$
\xymatrix{
0 \ar[r] &
\mathcal{F} \ar[r] \ar[d] &
\mathcal{G} \ar[r] \ar[d] &
\mathcal{H} \ar[r] \ar[d] &
0 \\
0 \ar[r] &
\mathcal{I}_1^\bullet \ar[r] &
\mathcal{I}_2^\bullet \ar[r] &
\mathcal{I}_3^\bullet \ar[r] &
0
}
$$
as in Derived Categories, Lemma \ref{derived-lemma-injective-resolution-ses}.
This will give a commutative diagram
$$
\xymatrix{
0 \ar[r] &
\check{\mathcal{C}}^\bullet(\mathcal{U}, \mathcal{F}) \ar[r] \ar[d] &
\check{\mathcal{C}}^\bullet(\mathcal{U}, \mathcal{G}) \ar[r] \ar[d] &
\check{\mathcal{C}}^\bullet(\mathcal{U}, \mathcal{H}) \ar[r] \ar[d] &
0 \\
0 \ar[r] &
\text{Tot}(\check{\mathcal{C}}^\bullet(\mathcal{U}, \mathcal{I}_1^\bullet))
\ar[r] &
\text{Tot}(\check{\mathcal{C}}^\bullet(\mathcal{U}, \mathcal{I}_2^\bullet))
\ar[r] &
\text{Tot}(\check{\mathcal{C}}^\bullet(\mathcal{U}, \mathcal{I}_3^\bullet))
\ar[r] &
0
}
$$
Here $\mathcal{U}$ is an open covering as above and
the vertical maps are those used to define the maps
$\check{H}^n(\mathcal{U}, -) \to H^n(X, -)$, see
Lemma \ref{lemma-cech-cohomology}.
The bottom complex is exact as the sequence of
complexes of injectives is termwise split exact.
Hence the boundary map in cohomology is computed
by the usual procedure for this lower exact sequence, see
Homology, Lemma \ref{homology-lemma-long-exact-sequence-cochain}.
The same will be true after passing to the refinement
$\mathcal{V}$ where the boundary map for {\v C}ech cohomology
was defined. Hence the boundary maps agree because they
use the same construction (whenever the first one is defined
on an element in {\v C}ech cohomology on a given open covering).
This finishes our discussion of the construction of
the structure of a $\delta$-functor on {\v C}ech cohomology
and why this structure is compatible with the given
$\delta$-functor structure on usual cohomology.

\medskip\noindent
Finally, we may apply Lemma \ref{lemma-injective-trivial-cech}
to see that higher {\v C}ech cohomology is trivial on injective
sheaves. Hence we see that {\v C}ech cohomology is a universal
$\delta$-functor by
Homology, Lemma \ref{homology-lemma-efface-implies-universal}.
\end{proof}

\begin{lemma}
\label{lemma-cech-Hausdorff-quasi-compact}
Let $X$ be a Hausdorff and quasi-compact topological space. Let
$\mathcal{F}$ be an abelian sheaf on $X$. Then
the map $\check{H}^n(X, \mathcal{F}) \to H^n(X, \mathcal{F})$ defined
in (\ref{equation-cech-to-cohomology}) is an isomorphism for all $n$.
\end{lemma}

\begin{proof}
Every open covering of $X$ has a finite subcover, which is a locally finite
refinement. Thus $X$ is paracompact and the result follows from
Lemma \ref{lemma-cech-paracompact}.
\end{proof}

\begin{lemma}
\label{lemma-cohomology-of-closed}
\begin{reference}
\cite[Expose V bis, 4.1.3]{SGA4}
\end{reference}
Let $X$ be a topological space. Let $Z \subset X$ be a quasi-compact subset
such that any two points of $Z$ have disjoint open neighbourhoods in $X$.
For every abelian sheaf $\mathcal{F}$ on $X$ the canonical
map
$$
\colim H^p(U, \mathcal{F})
\longrightarrow
H^p(Z, \mathcal{F}|_Z)
$$
where the colimit is over open neighbourhoods $U$ of $Z$ in $X$
is an isomorphism.
\end{lemma}

\begin{proof}
We first prove this for $p = 0$. Injectivity follows from
the definition of $\mathcal{F}|_Z$ and holds in general
(for any subset of any topological space $X$). Next, suppose that
$s \in H^0(Z, \mathcal{F}|_Z)$. Then we can find opens $U_i \subset X$
such that $Z \subset \bigcup U_i$ and such that $s|_{Z \cap U_i}$
comes from $s_i \in \mathcal{F}(U_i)$. It follows that
there exist opens $W_{ij} \subset U_i \cap U_j$ with
$W_{ij} \cap Z = U_i \cap U_j \cap Z$ such that
$s_i|_{W_{ij}} = s_j|_{W_{ij}}$. Applying
Topology, Lemma
\ref{topology-lemma-lift-covering-of-quasi-compact-hausdorff-subset}
we find opens $V_i$ of $X$ such that $V_i \subset U_i$ and
such that $V_i \cap V_j \subset W_{ij}$. Hence we see that
$s_i|_{V_i}$ glue to a section of $\mathcal{F}$ over the
open neighbourhood $\bigcup V_i$ of $Z$.

\medskip\noindent
To finish the proof, it suffices to show that if $\mathcal{I}$ is an
injective abelian sheaf on $X$, then $H^p(Z, \mathcal{I}|_Z) = 0$
for $p > 0$. This follows using short exact sequences and dimension
shifting; details omitted. Thus, suppose $\overline{\xi}$ is an element
of $H^p(Z, \mathcal{I}|_Z)$ for some $p > 0$.
By Lemma \ref{lemma-cech-Hausdorff-quasi-compact}
the element $\overline{\xi}$ comes from
$\check{H}^p(\mathcal{V}, \mathcal{I}|_Z)$
for some open covering $\mathcal{V} : Z = \bigcup V_i$ of $Z$.
Say $\overline{\xi}$ is the image of the class of a cocycle
$\xi = (\xi_{i_0 \ldots i_p})$ in
$\check{\mathcal{C}}^p(\mathcal{V}, \mathcal{I}|_Z)$.

\medskip\noindent
Let $\mathcal{I}' \subset \mathcal{I}|_Z$ be the subpresheaf
defined by the rule
$$
\mathcal{I}'(V) =
\{s \in \mathcal{I}|_Z(V) \mid
\exists (U, t),\ U \subset X\text{ open},
\ t \in \mathcal{I}(U),\ V = Z \cap U,\ s = t|_{Z \cap U} \}
$$
Then $\mathcal{I}|_Z$ is the sheafification of $\mathcal{I}'$.
Thus for every $(p + 1)$-tuple $i_0 \ldots i_p$ we can find an
open covering $V_{i_0 \ldots i_p} = \bigcup W_{i_0 \ldots i_p, k}$
such that $\xi_{i_0 \ldots i_p}|_{W_{i_0 \ldots i_p, k}}$ is
a section of $\mathcal{I}'$. Applying
Topology, Lemma \ref{topology-lemma-refine-covering}
we may after refining $\mathcal{V}$ assume that each
$\xi_{i_0 \ldots i_p}$ is a section of the presheaf $\mathcal{I}'$.

\medskip\noindent
Write $V_i = Z \cap U_i$ for some opens $U_i \subset X$.
Since $\mathcal{I}$ is flasque (Lemma \ref{lemma-injective-flasque})
and since $\xi_{i_0 \ldots i_p}$ is a section of $\mathcal{I}'$
for every $(p + 1)$-tuple $i_0 \ldots i_p$ we can choose
a section $s_{i_0 \ldots i_p} \in \mathcal{I}(U_{i_0 \ldots i_p})$
which restricts to $\xi_{i_0 \ldots i_p}$ on
$V_{i_0 \ldots i_p} = Z \cap U_{i_0 \ldots i_p}$.
(This appeal to injectives being flasque can be avoided by an
additional application of
Topology, Lemma
\ref{topology-lemma-lift-covering-of-quasi-compact-hausdorff-subset}.)
Let $s = (s_{i_0 \ldots i_p})$ be the corresponding cochain
for the open covering $U = \bigcup U_i$.
Since $\text{d}(\xi) = 0$ we see that the sections
$\text{d}(s)_{i_0 \ldots i_{p + 1}}$ restrict to zero
on $Z \cap U_{i_0 \ldots i_{p + 1}}$. Hence, by the initial
remarks of the proof, there exists open subsets
$W_{i_0 \ldots i_{p + 1}} \subset U_{i_0 \ldots i_{p + 1}}$
with $Z \cap W_{i_0 \ldots i_{p + 1}} = Z \cap U_{i_0 \ldots i_{p + 1}}$
such that $\text{d}(s)_{i_0 \ldots i_{p + 1}}|_{W_{i_0 \ldots i_{p + 1}}} = 0$.
By Topology, Lemma
\ref{topology-lemma-lift-covering-of-quasi-compact-hausdorff-subset}
we can find $U'_i \subset U_i$ such that $Z \subset \bigcup U'_i$
and such that $U'_{i_0 \ldots i_{p + 1}} \subset W_{i_0 \ldots i_{p + 1}}$.
Then $s' = (s'_{i_0 \ldots i_p})$ with
$s'_{i_0 \ldots i_p} = s_{i_0 \ldots i_p}|_{U'_{i_0 \ldots i_p}}$
is a cocycle for $\mathcal{I}$ for the open covering
$U' = \bigcup U'_i$ of an open neighbourhood of $Z$.
Since $\mathcal{I}$ has trivial higher {\v C}ech cohomology groups
(Lemma \ref{lemma-injective-trivial-cech})
we conclude that $s'$ is a coboundary. It follows that the image of
$\xi$ in the {\v C}ech complex for the open covering
$Z = \bigcup Z \cap U'_i$ is a coboundary and we are done.
\end{proof}






\section{The base change map}
\label{section-base-change-map}

\noindent
We will need to know how to construct the base change map in some cases.
Since we have not yet discussed derived pullback we only discuss
this in the case of a base change by a flat morphism of ringed spaces.
Before we state the result, let us discuss flat pullback on the derived
category. Namely, suppose that $g : X \to Y$ is a flat morphism of
ringed spaces. By Modules, Lemma \ref{modules-lemma-pullback-flat}
the functor $g^* : \textit{Mod}(\mathcal{O}_Y) \to
\textit{Mod}(\mathcal{O}_X)$ is exact. Hence it has a derived functor
$$
g^* : D^{+}(Y) \to D^{+}(X)
$$
which is computed by simply pulling back an representative of a given
object in $D^{+}(Y)$, see
Derived Categories, Lemma \ref{derived-lemma-right-derived-exact-functor}.
Hence as indicated we indicate this functor by $g^*$ rather than
$Lg^*$.

\begin{lemma}
\label{lemma-base-change-map-flat-case}
Let
$$
\xymatrix{
X' \ar[r]_{g'} \ar[d]_{f'} &
X \ar[d]^f \\
S' \ar[r]^g &
S
}
$$
be a commutative diagram of ringed spaces.
Let $\mathcal{F}^\bullet$ be a bounded below complex of
$\mathcal{O}_X$-modules.
Assume both $g$ and $g'$ are flat.
Then there exists a canonical base change map
$$
g^*Rf_*\mathcal{F}^\bullet
\longrightarrow
R(f')_*(g')^*\mathcal{F}^\bullet
$$
in $D^{+}(S')$.
\end{lemma}

\begin{proof}
Choose injective resolutions $\mathcal{F}^\bullet \to \mathcal{I}^\bullet$
and $(g')^*\mathcal{F}^\bullet \to \mathcal{J}^\bullet$.
By Lemma \ref{lemma-pushforward-injective-flat} we see that
$(g')_*\mathcal{J}^\bullet$ is a complex of injectives representing
$R(g')_*(g')^*\mathcal{F}^\bullet$. Hence by
Derived Categories, Lemmas \ref{derived-lemma-morphisms-lift}
and \ref{derived-lemma-morphisms-equal-up-to-homotopy}
the arrow $\beta$ in the diagram
$$
\xymatrix{
(g')_*(g')^*\mathcal{F}^\bullet \ar[r] &
(g')_*\mathcal{J}^\bullet \\
\mathcal{F}^\bullet \ar[u]^{adjunction} \ar[r] &
\mathcal{I}^\bullet \ar[u]_\beta
}
$$
exists and is unique up to homotopy.
Pushing down to $S$ we get
$$
f_*\beta :
f_*\mathcal{I}^\bullet
\longrightarrow
f_*(g')_*\mathcal{J}^\bullet
=
g_*(f')_*\mathcal{J}^\bullet
$$
By adjunction of $g^*$ and $g_*$ we get a map of complexes
$g^*f_*\mathcal{I}^\bullet \to (f')_*\mathcal{J}^\bullet$.
Note that this map is unique up to homotopy since the only
choice in the whole process was the choice of the map $\beta$
and everything was done on the level of complexes.
\end{proof}

\begin{remark}
\label{remark-correct-version-base-change-map}
The ``correct'' version of the base change map is the map
$$
Lg^* Rf_* \mathcal{F}^\bullet
\longrightarrow
R(f')_* L(g')^*\mathcal{F}^\bullet.
$$
The construction of this map involves
unbounded complexes, see Remark \ref{remark-base-change}.
\end{remark}





\section{Proper base change in topology}
\label{section-proper-base-change}

\noindent
In this section we prove a very general version of the proper base change
theorem in topology. It tells us that the stalks of the higher direct
images $R^pf_*$ can be computed on the fibre.

\begin{lemma}
\label{lemma-proper-base-change}
Let $f : (X, \mathcal{O}_X) \to (Y, \mathcal{O}_Y)$ be a morphism of
ringed spaces. Let $y \in Y$. Assume that
\begin{enumerate}
\item $f$ is closed,
\item $f$ is separated, and
\item $f^{-1}(y)$ is quasi-compact.
\end{enumerate}
Then for $E$ in $D^+(\mathcal{O}_X)$
we have $(Rf_*E)_y = R\Gamma(f^{-1}(y), E|_{f^{-1}(y)})$ in
$D^+(\mathcal{O}_{Y, y})$.
\end{lemma}

\begin{proof}
The base change map of Lemma \ref{lemma-base-change-map-flat-case}
gives a canonical map $(Rf_*E)_y \to R\Gamma(f^{-1}(y), E|_{f^{-1}(y)})$.
To prove this map is an isomorphism, we represent $E$ by a bounded
below complex of injectives $\mathcal{I}^\bullet$.
Set $Z = f^{-1}(\{y\})$. The assumptions of
Lemma \ref{lemma-cohomology-of-closed}
are satisfied, see Topology, Lemma \ref{topology-lemma-separated}.
Hence the restrictions
$\mathcal{I}^n|_Z$ are acyclic for $\Gamma(Z, -)$.
Thus $R\Gamma(Z, E|_Z)$ is represented by the
complex $\Gamma(Z, \mathcal{I}^\bullet|_Z)$, see
Derived Categories, Lemma \ref{derived-lemma-leray-acyclicity}.
In other words, we have to show the map
$$
\colim_V \mathcal{I}^\bullet(f^{-1}(V))
\longrightarrow
\Gamma(Z, \mathcal{I}^\bullet|_Z)
$$
is an isomorphism. Using Lemma \ref{lemma-cohomology-of-closed}
we see that it suffices to show that the collection of open neighbourhoods
$f^{-1}(V)$ of $Z = f^{-1}(\{y\})$
is cofinal in the system of all open neighbourhoods.
If $f^{-1}(\{y\}) \subset U$ is an open neighbourhood, then as $f$ is closed
the set $V = Y \setminus f(X \setminus U)$ is an open neighbourhood
of $y$ with $f^{-1}(V) \subset U$. This proves the lemma.
\end{proof}

\begin{theorem}[Proper base change]
\label{theorem-proper-base-change}
\begin{reference}
\cite[Expose V bis, 4.1.1]{SGA4}
\end{reference}
Consider a cartesian square of topological spaces
$$
\xymatrix{
X' = Y' \times_Y X \ar[d]_{f'} \ar[r]_-{g'} & X \ar[d]^f \\
Y' \ar[r]^g & Y
}
$$
Assume that $f$ is proper.
Let $E$ be an object of $D^+(X)$. Then the base change map
$$
g^{-1}Rf_*E \longrightarrow Rf'_*(g')^{-1}E
$$
of Lemma \ref{lemma-base-change-map-flat-case} is an isomorphism
in $D^+(Y')$.
\end{theorem}

\begin{proof}
Let $y' \in Y'$ be a point with image $y \in Y$. It suffices to show that
the base change map induces an isomorphism on stalks at $y'$.
As $f$ is proper it follows that $f'$ is proper, the
fibres of $f$ and $f'$ are quasi-compact and $f$ and $f'$ are closed, see
Topology, Theorem \ref{topology-theorem-characterize-proper} and
Lemma \ref{topology-lemma-base-change-separated}.
Thus we can apply Lemma \ref{lemma-proper-base-change} twice to see that
$$
(Rf'_*(g')^{-1}E)_{y'} = R\Gamma((f')^{-1}(y'), (g')^{-1}E|_{(f')^{-1}(y')})
$$
and
$$
(Rf_*E)_y = R\Gamma(f^{-1}(y), E|_{f^{-1}(y)})
$$
The induced map of fibres $(f')^{-1}(y') \to f^{-1}(y)$ is
a homeomorphism of topological spaces and the pull back of
$E|_{f^{-1}(y)}$ is $(g')^{-1}E|_{(f')^{-1}(y')}$. The
desired result follows.
\end{proof}

\begin{lemma}[Proper base change for sheaves of sets]
\label{lemma-proper-base-change-sheaves-of-sets}
Consider a cartesian square of topological spaces
$$
\xymatrix{
X' \ar[d]_{f'} \ar[r]_-{g'} & X \ar[d]^f \\
Y' \ar[r]^g & Y
}
$$
Assume that $f$ is proper. Then
$g^{-1}f_*\mathcal{F} = f'_*(g')^{-1}\mathcal{F}$
for any sheaf of sets $\mathcal{F}$ on $X$.
\end{lemma}

\begin{proof}
We argue exactly as in the proof of Theorem \ref{theorem-proper-base-change}
and we find it suffices to show
$(f_*\mathcal{F})_y = \Gamma(X_y, \mathcal{F}|_{X_y})$.
Then we argue as in Lemma \ref{lemma-proper-base-change}
to reduce this to the $p = 0$ case of Lemma \ref{lemma-cohomology-of-closed}
for sheaves of sets. The first part of the proof of
Lemma \ref{lemma-cohomology-of-closed}
works for sheaves of sets and this finishes the proof.
Some details omitted.
\end{proof}



\section{Cohomology and colimits}
\label{section-limits}

\noindent
Let $X$ be a ringed space. Let $(\mathcal{F}_i, \varphi_{ii'})$ be
a system of sheaves of $\mathcal{O}_X$-modules over the directed set $I$, see
Categories, Section \ref{categories-section-posets-limits}.
Since for each $i$ there is a canonical map
$\mathcal{F}_i \to \colim_i \mathcal{F}_i$ we get a
canonical map
$$
\colim_i H^p(X, \mathcal{F}_i)
\longrightarrow
H^p(X, \colim_i \mathcal{F}_i)
$$
for every $p \geq 0$. Of course there is a similar map for
every open $U \subset X$. These maps are in general not isomorphisms,
even for $p = 0$. In this section we generalize the results of
Sheaves, Lemma \ref{sheaves-lemma-directed-colimits-sections}.
See also
Modules, Lemma \ref{modules-lemma-finite-presentation-quasi-compact-colimit}
(in the special case $\mathcal{G} = \mathcal{O}_X$).

\begin{lemma}
\label{lemma-quasi-separated-cohomology-colimit}
Let $X$ be a ringed space. Assume that the underlying topological space
of $X$ has the following properties:
\begin{enumerate}
\item there exists a basis of quasi-compact open subsets, and
\item the intersection of any two quasi-compact opens is quasi-compact.
\end{enumerate}
Then for any directed system $(\mathcal{F}_i, \varphi_{ii'})$
of sheaves of $\mathcal{O}_X$-modules and for any quasi-compact open
$U \subset X$ the canonical map
$$
\colim_i H^q(U, \mathcal{F}_i)
\longrightarrow
H^q(U, \colim_i \mathcal{F}_i)
$$
is an isomorphism for every $q \geq 0$.
\end{lemma}

\begin{proof}
It is important in this proof to argue for all quasi-compact opens
$U \subset X$ at the same time.
The result is true for $q = 0$ and any quasi-compact open $U \subset X$ by
Sheaves, Lemma \ref{sheaves-lemma-directed-colimits-sections}
(combined with
Topology, Lemma \ref{topology-lemma-topology-quasi-separated-scheme}).
Assume that we have proved the result for all $q \leq q_0$ and let
us prove the result for $q = q_0 + 1$.

\medskip\noindent
By our conventions on directed systems the index set $I$ is directed,
and any system of $\mathcal{O}_X$-modules $(\mathcal{F}_i, \varphi_{ii'})$
over $I$ is directed.
By Injectives, Lemma \ref{injectives-lemma-sheaves-modules-space} the category
of $\mathcal{O}_X$-modules has functorial injective embeddings.
Thus for any system $(\mathcal{F}_i, \varphi_{ii'})$ there exists a
system $(\mathcal{I}_i, \varphi_{ii'})$ with each $\mathcal{I}_i$ an
injective $\mathcal{O}_X$-module and a morphism of systems given
by injective $\mathcal{O}_X$-module maps
$\mathcal{F}_i \to \mathcal{I}_i$. Denote $\mathcal{Q}_i$ the
cokernel so that we have short exact sequences
$$
0 \to
\mathcal{F}_i \to
\mathcal{I}_i \to
\mathcal{Q}_i \to 0.
$$
We claim that the sequence
$$
0 \to
\colim_i \mathcal{F}_i \to
\colim_i \mathcal{I}_i \to
\colim_i \mathcal{Q}_i \to 0.
$$
is also a short exact sequence of $\mathcal{O}_X$-modules.
We may check this on stalks. By
Sheaves, Sections \ref{sheaves-section-limits-presheaves}
and \ref{sheaves-section-limits-sheaves}
taking stalks commutes with colimits. Since a directed colimit
of short exact sequences of abelian groups is short exact
(see Algebra, Lemma \ref{algebra-lemma-directed-colimit-exact})
we deduce the result. We claim that
$H^q(U, \colim_i \mathcal{I}_i) = 0$ for all quasi-compact
open $U \subset X$ and all $q \geq 1$. Accepting this claim
for the moment consider the diagram
$$
\xymatrix{
\colim_i H^{q_0}(U, \mathcal{I}_i) \ar[d] \ar[r] &
\colim_i H^{q_0}(U, \mathcal{Q}_i) \ar[d] \ar[r] &
\colim_i H^{q_0 + 1}(U, \mathcal{F}_i) \ar[d] \ar[r] &
0 \ar[d] \\
H^{q_0}(U, \colim_i \mathcal{I}_i) \ar[r] &
H^{q_0}(U, \colim_i \mathcal{Q}_i) \ar[r] &
H^{q_0 + 1}(U, \colim_i \mathcal{F}_i) \ar[r] &
0
}
$$
The zero at the lower right corner comes from the claim and the
zero at the upper right corner comes from the fact that the sheaves
$\mathcal{I}_i$ are injective.
The top row is exact by an application of
Algebra, Lemma \ref{algebra-lemma-directed-colimit-exact}.
Hence by the snake lemma we deduce the
result for $q = q_0 + 1$.

\medskip\noindent
It remains to show that the claim is true. We will use
Lemma \ref{lemma-cech-vanish-basis}.
Let $\mathcal{B}$ be the collection of all quasi-compact open
subsets of $X$. This is a basis for the topology on $X$ by assumption.
Let $\text{Cov}$ be the collection of finite open coverings
$\mathcal{U} : U = \bigcup_{j = 1, \ldots, m} U_j$ with each
of $U$, $U_j$ quasi-compact open in $X$. By the result for $q = 0$
we see that for $\mathcal{U} \in \text{Cov}$ we have
$$
\check{\mathcal{C}}^\bullet(\mathcal{U}, \colim_i \mathcal{I}_i)
=
\colim_i \check{\mathcal{C}}^\bullet(\mathcal{U}, \mathcal{I}_i)
$$
because all the multiple intersections $U_{j_0 \ldots j_p}$
are quasi-compact. By Lemma \ref{lemma-injective-trivial-cech}
each of the complexes in the colimit of {\v C}ech complexes is
acyclic in degree $\geq 1$. Hence by
Algebra, Lemma \ref{algebra-lemma-directed-colimit-exact}
we see that also the {\v C}ech complex
$\check{\mathcal{C}}^\bullet(\mathcal{U}, \colim_i \mathcal{I}_i)$
is acyclic in degrees $\geq 1$. In other words we see that
$\check{H}^p(\mathcal{U},  \colim_i \mathcal{I}_i) = 0$
for all $p \geq 1$. Thus the assumptions of
Lemma \ref{lemma-cech-vanish-basis} are satisfied and the claim follows.
\end{proof}

\begin{lemma}
\label{lemma-higher-direct-image-colimit}
Let $f : X \to Y$ be a continuous map of topological spaces.
Let $(\mathcal{F}_i, \varphi_{ii'})$ be a system of
abelian sheaves on $X$. Set $\mathcal{F} = \colim \mathcal{F}_i$.
Let $p \geq 0$ be an integer. Assume the set of opens $V \subset Y$ such that
$H^p(f^{-1}(V), \mathcal{F}) = \colim H^p(f^{-1}(V), \mathcal{F}_i)$
is a basis for the topology on $Y$. Then
$R^pf_*\mathcal{F} = \colim R^pf_*\mathcal{F}_i$.
\end{lemma}

\begin{proof}
Recall that $R^pf_*\mathcal{F}$ is the sheafification of the presheaf
$\mathcal{G}$ sending $V$ to $H^p(f^{-1}(V), \mathcal{F})$, see
Lemma \ref{lemma-describe-higher-direct-images}. Similarly,
$R^pf_*\mathcal{F}_i$ is the sheafification of the presheaf
$\mathcal{G}_i$ sending $V$ to $H^p(f^{-1}(V), \mathcal{F}_i)$.
Recall that sheafification is the left adjoint to the inclusion
from sheaves to presheaves, see Sheaves, Section
\ref{sheaves-section-sheafification}. Hence sheafification commutes
with colimits, see Categories, Lemma \ref{categories-lemma-adjoint-exact}.
Hence it suffices to show that the map of presheaves (with colimit
in the category of presheaves)
$$
\colim \mathcal{G}_i \longrightarrow \mathcal{G}
$$
induces an isomorphism on sheafifications. For this it suffices to show
that the presheaves $\mathcal{G}$ and $\colim \mathcal{G}_i$
agree on a basis for the topology of $Y$. Namely, in this case the stalks
of their sheafifications, which can be computed directly from the
presheaf values on elements of the basis, agree.
The required agreement is exactly the assumption of the lemma.
\end{proof}

\noindent
Next we formulate the analogy of
Sheaves, Lemma \ref{sheaves-lemma-descend-opens}
for cohomology.
Let $X$ be a spectral space which is written as a cofiltered limit
of spectral spaces $X_i$ for a diagram with spectral transition morphisms
as in
Topology, Lemma \ref{topology-lemma-directed-inverse-limit-spectral-spaces}.
Assume given
\begin{enumerate}
\item an abelian sheaf $\mathcal{F}_i$ on $X_i$ for all
$i \in \Ob(\mathcal{I})$,
\item for $a : j \to i$ an $f_a$-map
$\varphi_a : \mathcal{F}_i \to \mathcal{F}_j$ of abelian sheaves (see
Sheaves, Definition \ref{sheaves-definition-f-map})
\end{enumerate}
such that $\varphi_c = \varphi_b \circ \varphi_a$
whenever $c = a \circ b$. Set $\mathcal{F} = \colim p_i^{-1}\mathcal{F}_i$
on $X$.

\begin{lemma}
\label{lemma-colimit}
In the situation discussed above.
Let $i \in \Ob(\mathcal{I})$ and let $U_i \subset X_i$ be quasi-compact open.
Then
$$
\colim_{a : j \to i} H^p(f_a^{-1}(U_i), \mathcal{F}_j) =
H^p(p_i^{-1}(U_i), \mathcal{F})
$$
for all $p \geq 0$. In particular we have
$H^p(X, \mathcal{F}) = \colim H^p(X_i, \mathcal{F}_i)$.
\end{lemma}

\begin{proof}
The case $p = 0$ is Sheaves, Lemma \ref{sheaves-lemma-descend-opens}.

\medskip\noindent
In this paragraph we show that we can find a map of systems
$(\gamma_i) : (\mathcal{F}_i, \varphi_a) \to (\mathcal{G}_i, \psi_a)$
with $\mathcal{G}_i$ an injective abelian sheaf and $\gamma_i$ injective.
For each $i$ we pick an injection $\mathcal{F}_i \to \mathcal{I}_i$
where $\mathcal{I}_i$ is an injective abelian sheaf on $X_i$.
Then we can consider the family of maps
$$
\gamma_i :
\mathcal{F}_i
\longrightarrow
\prod\nolimits_{b : k \to i} f_{b, *}\mathcal{I}_k = \mathcal{G}_i
$$
where the component maps are the maps adjoint to the maps
$f_b^{-1}\mathcal{F}_i \to \mathcal{F}_k \to \mathcal{I}_k$.
For $a : j \to i$ in $\mathcal{I}$ there is a canonical map
$$
\psi_a : f_a^{-1}\mathcal{G}_i \to \mathcal{G}_j
$$
whose components are the canonical maps
$f_b^{-1}f_{a \circ b, *}\mathcal{I}_k \to f_{b, *}\mathcal{I}_k$
for $b : k \to j$. Thus we find an injection
$\{\gamma_i\} : \{\mathcal{F}_i, \varphi_a) \to (\mathcal{G}_i, \psi_a)$
of systems of abelian sheaves. Note that $\mathcal{G}_i$ is an injective
sheaf of abelian groups on $X_i$, see
Lemma \ref{lemma-pushforward-injective-flat} and
Homology, Lemma \ref{homology-lemma-product-injectives}.
This finishes the construction.

\medskip\noindent
Arguing exactly as in the proof of
Lemma \ref{lemma-quasi-separated-cohomology-colimit}
we see that it suffices to prove that
$H^p(X, \colim f_i^{-1}\mathcal{G}_i) = 0$ for $p > 0$.

\medskip\noindent
Set $\mathcal{G} = \colim f_i^{-1}\mathcal{G}_i$.
To show vanishing of cohomology of $\mathcal{G}$ on every quasi-compact
open of $X$, it suffices to show that the {\v C}ech cohomology of
$\mathcal{G}$ for any open covering $\mathcal{U}$ of a quasi-compact open of
$X$ by finitely many quasi-compact opens is zero, see
Lemma \ref{lemma-cech-vanish-basis}.
Such a covering is the inverse by $p_i$ of such a covering $\mathcal{U}_i$
on the space $X_i$ for some $i$ by
Topology, Lemma \ref{topology-lemma-descend-opens}. We have
$$
\check{\mathcal{C}}^\bullet(\mathcal{U}, \mathcal{G}) =
\colim_{a : j \to i}
\check{\mathcal{C}}^\bullet(f_a^{-1}(\mathcal{U}_i), \mathcal{G}_j)
$$
by the case $p = 0$. The right hand side is a filtered colimit of
complexes each of which is acyclic in positive degrees by
Lemma \ref{lemma-injective-trivial-cech}. Thus we conclude by
Algebra, Lemma \ref{algebra-lemma-directed-colimit-exact}.
\end{proof}









\section{Vanishing on Noetherian topological spaces}
\label{section-vanishing-Noetherian}

\noindent
The aim is to prove a theorem of Grothendieck namely
Proposition \ref{proposition-vanishing-Noetherian}. See \cite{Tohoku}.

\begin{lemma}
\label{lemma-cohomology-and-closed-immersions}
\begin{reference}
\cite[Chapter I, \S 3, no. 26, Proposition 8, p. 219]{FAC}
\cite[Chapter II, \S 2, no. 39, p. 232]{FAC}
\end{reference}
Let $i : Z \to X$ be a closed immersion of topological spaces.
For any abelian sheaf $\mathcal{F}$ on $Z$ we have
$H^p(Z, \mathcal{F}) = H^p(X, i_*\mathcal{F})$.
\end{lemma}

\begin{proof}
This is true because $i_*$ is exact (see
Modules, Lemma \ref{modules-lemma-i-star-exact}),
and hence $R^pi_* = 0$ as a functor
(Derived Categories, Lemma \ref{derived-lemma-right-derived-exact-functor}).
Thus we may apply Lemma \ref{lemma-apply-Leray}.
\end{proof}

\begin{lemma}
\label{lemma-irreducible-constant-cohomology-zero}
Let $X$ be an irreducible topological space.
Then $H^p(X, \underline{A}) = 0$ for all $p > 0$
and any abelian group $A$.
\end{lemma}

\begin{proof}
Recall that $\underline{A}$ is the constant sheaf as defined
in Sheaves, Definition \ref{sheaves-definition-constant-sheaf}.
Since $X$ is irreducible, any nonempty open $U$ is
irreducible and a fortiori connected. Hence for $U \subset X$
nonempty open we have $\underline{A}(U) = A$.
We have $\underline{A}(\emptyset) = 0$. Thus $\underline{A}$
is a flasque abelian sheaf on $X$. The vanishing follows
from Lemma \ref{lemma-flasque-acyclic}.
\end{proof}

\begin{lemma}
\label{lemma-subsheaf-of-constant-sheaf}
\begin{reference}
\cite[Page 168]{Tohoku}.
\end{reference}
Let $X$ be a topological space such that the intersection of any
two quasi-compact opens is quasi-compact. Let
$\mathcal{F} \subset \underline{\mathbf{Z}}$
be a subsheaf generated by finitely many sections over quasi-compact opens.
Then there exists a finite filtration
$$
(0) = \mathcal{F}_0 \subset \mathcal{F}_1 \subset \ldots \subset
\mathcal{F}_n = \mathcal{F}
$$
by abelian subsheaves such that for each $0 < i \leq n$
there exists a short exact sequence
$$
0 \to j'_!\underline{\mathbf{Z}}_V \to j_!\underline{\mathbf{Z}}_U \to
\mathcal{F}_i/\mathcal{F}_{i - 1} \to 0
$$
with $j : U \to X$ and $j' : V \to X$ the inclusion of quasi-compact opens
into $X$.
\end{lemma}

\begin{proof}
Say $\mathcal{F}$ is generated by the sections $s_1, \ldots, s_t$ over the
quasi-compact opens $U_1, \ldots, U_t$. Since $U_i$ is quasi-compact and
$s_i$ a locally constant function to $\mathbf{Z}$ we may assume, after
possibly replacing $U_i$ by the parts of a finite decomposition into open
and closed subsets, that $s_i$ is a constant section.
Say $s_i = n_i$ with $n_i \in \mathbf{Z}$. Of course we can remove
$(U_i, n_i)$ from the list if $n_i = 0$. Flipping signs if necessary
we may also assume $n_i > 0$. Next, for any subset $I \subset \{1, \ldots, t\}$
we may add $\bigcap_{i \in I} U_i$ and $\gcd(n_i, i \in I)$ to the list.
After doing this we see that our list $(U_1, n_1), \ldots, (U_t, n_t)$
satisfies the following property:
For $x \in X$ set $I_x = \{i \in \{1, \ldots, t\} \mid x \in U_i\}$.
Then $\gcd(n_i, i \in I_x)$ is attained by $n_i$ for some $i \in I_x$.

\medskip\noindent
As our filtration we take $\mathcal{F}_0 = (0)$ and
$\mathcal{F}_n$ generated by the sections $n_i$ over $U_i$ for those
$i$ such that $n_i \leq n$. It is clear that
$\mathcal{F}_n = \mathcal{F}$ for $n \gg 0$. Moreover, the quotient
$\mathcal{F}_n/\mathcal{F}_{n - 1}$ is generated by the section
$n$ over $U = \bigcup_{n_i \leq n} U_i$ and the kernel of the map
$j_!\underline{\mathbf{Z}}_U \to \mathcal{F}_n/\mathcal{F}_{n - 1}$
is generated by the section $n$ over $V = \bigcup_{n_i \leq n - 1} U_i$.
Thus a short exact sequence as in the statement of the lemma.
\end{proof}

\begin{lemma}
\label{lemma-vanishing-generated-one-section}
\begin{reference}
This is a special case of \cite[Proposition 3.6.1]{Tohoku}.
\end{reference}
Let $X$ be a topological space. Let $d \geq 0$ be an integer. Assume
\begin{enumerate}
\item $X$ is quasi-compact,
\item the quasi-compact opens form a basis for $X$, and
\item the intersection of two quasi-compact opens is quasi-compact.
\item $H^p(X, j_!\underline{\mathbf{Z}}_U) = 0$ for all $p > d$
and any quasi-compact open $j : U \to X$.
\end{enumerate}
Then $H^p(X, \mathcal{F}) = 0$ for all $p > d$
and any abelian sheaf $\mathcal{F}$ on $X$.
\end{lemma}

\begin{proof}
Let $S = \coprod_{U \subset X} \mathcal{F}(U)$ where $U$ runs over the
quasi-compact opens of $X$.
For any finite subset $A = \{s_1, \ldots, s_n\} \subset S$,
let $\mathcal{F}_A$ be the subsheaf of $\mathcal{F}$ generated
by all $s_i$ (see
Modules, Definition \ref{modules-definition-generated-by-local-sections}).
Note that if $A \subset A'$, then $\mathcal{F}_A \subset \mathcal{F}_{A'}$.
Hence $\{\mathcal{F}_A\}$ forms a system over the
directed partially ordered set of finite subsets of $S$.
By Modules, Lemma \ref{modules-lemma-generated-by-local-sections-stalk}
it is clear that
$$
\colim_A \mathcal{F}_A = \mathcal{F}
$$
by looking at stalks. By
Lemma \ref{lemma-quasi-separated-cohomology-colimit} we have
$$
H^p(X, \mathcal{F}) =
\colim_A H^p(X, \mathcal{F}_A)
$$
Hence it suffices to prove the vanishing for the abelian sheaves
$\mathcal{F}_A$. In other words, it suffices to prove the
result when $\mathcal{F}$ is generated by finitely many local sections
over quasi-compact opens of $X$.

\medskip\noindent
Suppose that $\mathcal{F}$ is generated by the local sections
$s_1, \ldots, s_n$. Let $\mathcal{F}' \subset \mathcal{F}$
be the subsheaf generated by $s_1, \ldots, s_{n - 1}$.
Then we have a short exact sequence
$$
0 \to \mathcal{F}' \to \mathcal{F} \to \mathcal{F}/\mathcal{F}' \to 0
$$
From the long exact sequence of cohomology we see that it suffices
to prove the vanishing for the abelian sheaves $\mathcal{F}'$
and $\mathcal{F}/\mathcal{F}'$ which are generated by fewer than
$n$ local sections. Hence it suffices to prove the vanishing
for sheaves generated by at most one local section. These sheaves
are exactly the quotients of the sheaves $j_!\underline{\mathbf{Z}}_U$
where $U$ is a quasi-compact open of $X$.

\medskip\noindent
Assume now that we have a short exact sequence
$$
0 \to \mathcal{K} \to j_!\underline{\mathbf{Z}}_U \to \mathcal{F} \to 0
$$
with $U$ quasi-compact open in $X$.
It suffices to show that $H^q(X, \mathcal{K})$ is zero for $q \geq d + 1$.
As above we can write $\mathcal{K}$ as the filtered colimit of
subsheaves $\mathcal{K}'$ generated by finitely many sections over
quasi-compact opens. Then $\mathcal{F}$ is the filtered colimit of the
sheaves $j_!\underline{\mathbf{Z}}_U/\mathcal{K}'$. In this way we
reduce to the case that $\mathcal{K}$ is generated by finitely many
sections over quasi-compact opens. Note that $\mathcal{K}$
is a subsheaf of $\underline{\mathbf{Z}}_X$. Thus by
Lemma \ref{lemma-subsheaf-of-constant-sheaf} there exists a finite
filtration of $\mathcal{K}$ whose successive quotients $\mathcal{Q}$ fit
into a short exact sequence
$$
0 \to j''_!\underline{\mathbf{Z}}_W \to
j'_!\underline{\mathbf{Z}}_V \to \mathcal{Q} \to 0
$$
with $j'' : W \to X$ and $j' : V \to X$ the inclusions of quasi-compact opens.
Hence the vanishing of $H^p(X, \mathcal{Q})$ for $p > d$ follows
from our assumption (in the lemma) on the vanishing of the cohomology groups
of $j''_!\underline{\mathbf{Z}}_W$ and $j'_!\underline{\mathbf{Z}}_V$.
Returning to $\mathcal{K}$ this, via an induction argument using the
long exact cohomology sequence, implies the desired vanishing for it as well.
\end{proof}

\begin{example}
\label{example-datta}
Let $X = \mathbf{N}$ endowed with the topology whose opens are
$\emptyset$, $X$, and $U_n = \{i \mid i \leq n\}$ for $n \geq 1$.
An abelian sheaf $\mathcal{F}$ on $X$ is the same as an inverse
system of abelian groups $A_n = \mathcal{F}(U_n)$ and
$\Gamma(X, \mathcal{F}) = \lim A_n$. Since the inverse limit
functor is not an exact functor on the category of inverse systems,
we see that there is an abelian sheaf with nonzero $H^1$.
Finally, the reader can check that $H^p(X, j_!\mathbf{Z}_U) = 0$,
$p \geq 1$ if $j : U = U_n \to X$ is the inclusion. Thus we see
that $X$ is an example of a space satisfying conditions (2), (3), and (4) of
Lemma \ref{lemma-vanishing-generated-one-section} for $d = 0$
but not the conclusion.
\end{example}

\begin{lemma}
\label{lemma-subsheaf-irreducible}
Let $X$ be an irreducible topological space.
Let $\mathcal{H} \subset \underline{\mathbf{Z}}$ be
an abelian subsheaf of the constant sheaf.
Then there exists a nonempty open $U \subset X$ such
that $\mathcal{H}|_U = \underline{d\mathbf{Z}}_U$
for some $d \in \mathbf{Z}$.
\end{lemma}

\begin{proof}
Recall that $\underline{\mathbf{Z}}(V) = \mathbf{Z}$
for any nonempty open $V$ of $X$ (see proof of
Lemma \ref{lemma-irreducible-constant-cohomology-zero}).
If $\mathcal{H} = 0$, then the lemma holds with $d = 0$.
If $\mathcal{H} \not = 0$, then there exists a nonempty open
$U \subset X$ such that $\mathcal{H}(U) \not = 0$.
Say $\mathcal{H}(U) = n\mathbf{Z}$ for some $n \geq 1$.
Hence we see that
$\underline{n\mathbf{Z}}_U
\subset \mathcal{H}|_U \subset
\underline{\mathbf{Z}}_U$. If the first inclusion is strict we
can find a nonempty $U' \subset U$ and an integer $1 \leq n' < n$
such that
$\underline{n'\mathbf{Z}}_{U'}
\subset \mathcal{H}|_{U'} \subset
\underline{\mathbf{Z}}_{U'}$.
This process has to stop after a finite number of steps, and
hence we get the lemma.
\end{proof}

\begin{proposition}[Grothendieck]
\label{proposition-vanishing-Noetherian}
\begin{reference}
\cite[Theorem 3.6.5]{Tohoku}.
The projective-variety special case is
\cite[Chapter III, \S 1, no. 52, Proposition 3 and Corollary,
pp. 244--245]{FAC}.
Its closed-projective-variety application is
\cite[Chapter III, \S 3, no. 66, Theorem 1, p. 259]{FAC}.
The curve case and its historical question are
\cite[Chapter III, \S 1, no. 53, Proposition 4 and Remark, p. 246]{FAC}.
\end{reference}
\begin{history}
The FAC proposition proves vanishing above the dimension of the support for a
coherent algebraic sheaf on a classical variety locally closed in projective
space. Its proof constructs a bounded affine covering by homogeneous
principal opens and uses the alternating Cech complex. The proposition below
is the stronger topological theorem: it applies to every abelian sheaf on an
arbitrary Noetherian space. Applying it on the closed support and using closed
pushforward recovers the FAC bound without the quasi-projective hypothesis.
In the following number FAC gives a separate curve proof by a cofinal family
of covers whose higher alternating complexes are those of finite simplices,
then asks whether an analogous assertion holds in arbitrary dimension. Since
the source notation there denotes direct-limit Cech cohomology, the theorem
below supplies the affirmative derived-functor analogue rather than silently
identifying the two cohomology theories.

In no. 66 the source extends a coherent sheaf on a closed projective
subvariety by zero to projective space and combines the no. 52 bound with
cohomology invariance under closed pushforward. The proposition below makes
the resulting dimension bound independent of projectivity, coherence, and
the chosen ambient space.
\end{history}
Let $X$ be a Noetherian topological space.
If $\dim(X) \leq d$, then $H^p(X, \mathcal{F}) = 0$
for all $p > d$ and any abelian sheaf $\mathcal{F}$
on $X$.
\end{proposition}

\begin{proof}
We prove this lemma by induction on $d$.
So fix $d$ and assume the lemma holds for all
Noetherian topological spaces of dimension $< d$.

\medskip\noindent
Let $\mathcal{F}$ be an abelian sheaf on $X$.
Suppose $U \subset X$ is an open. Let $Z \subset X$
denote the closed complement.
Denote $j : U \to X$ and $i : Z \to X$ the inclusion maps.
Then there is a short exact sequence
$$
0 \to j_{!}j^*\mathcal{F} \to \mathcal{F} \to i_*i^*\mathcal{F} \to 0
$$
see Modules, Lemma \ref{modules-lemma-canonical-exact-sequence}.
Note that $j_!j^*\mathcal{F}$ is supported on
the topological closure $Z'$ of $U$, i.e., it is of
the form $i'_*\mathcal{F}'$ for some abelian sheaf $\mathcal{F}'$
on $Z'$, where $i' : Z' \to X$ is the inclusion.

\medskip\noindent
We can use this to reduce to the case where $X$ is irreducible.
Namely, according to
Topology, Lemma \ref{topology-lemma-Noetherian}
$X$ has finitely
many irreducible components. If $X$ has more than one irreducible
component, then let $Z \subset X$ be an irreducible component of $X$
and set $U = X \setminus Z$. By the above, and the long exact sequence
of cohomology, it suffices to prove the vanishing of
$H^p(X, i_*i^*\mathcal{F})$ and $H^p(X, i'_*\mathcal{F}')$ for $p > d$.
By Lemma \ref{lemma-cohomology-and-closed-immersions} it suffices to prove
$H^p(Z, i^*\mathcal{F})$ and $H^p(Z', \mathcal{F}')$ vanish for $p > d$.
Since $Z'$ and $Z$ have fewer irreducible components we indeed
reduce to the case of an irreducible $X$.

\medskip\noindent
If $d = 0$ and $X$ is irreducible, then $X$ is the only nonempty
open subset of $X$. Hence every sheaf is constant and higher cohomology
groups vanish (for example by
Lemma \ref{lemma-irreducible-constant-cohomology-zero}).

\medskip\noindent
Suppose $X$ is irreducible of dimension $d > 0$.
By Lemma \ref{lemma-vanishing-generated-one-section}
we reduce to the case where
$\mathcal{F} = j_!\underline{\mathbf{Z}}_U$ for some open $U \subset X$.
In this case we look at the short exact sequence
$$
0 \to j_!(\underline{\mathbf{Z}}_U) \to
\underline{\mathbf{Z}}_X \to i_*\underline{\mathbf{Z}}_Z \to 0
$$
where $Z = X \setminus U$.
By Lemma \ref{lemma-irreducible-constant-cohomology-zero}
we have the vanishing of $H^p(X, \underline{\mathbf{Z}}_X)$
for all $p \geq 1$. By induction we have
$H^p(X, i_*\underline{\mathbf{Z}}_Z) = H^p(Z, \underline{\mathbf{Z}}_Z) = 0$
for $p \geq d$. Hence we win by the long exact cohomology sequence.
\end{proof}





\section{Cohomology with support in a closed subset}
\label{section-cohomology-support}

\noindent
This section just discusses the bare minimum -- the discussion
will be continued in Section \ref{section-cohomology-support-bis}.

\medskip\noindent
Let $X$ be a topological space and let $Z \subset X$ be a closed subset.
Let $\mathcal{F}$ be an abelian sheaf on $X$. We let
$$
\Gamma_Z(X, \mathcal{F}) =
\{s \in \mathcal{F}(X) \mid \text{Supp}(s) \subset Z\}
$$
be the subset of sections whose support is contained in $Z$.
The support of a section is defined in
Modules, Definition \ref{modules-definition-support}.
Modules, Lemma \ref{modules-lemma-support-section-closed}
implies that $\Gamma_Z(X, \mathcal{F})$ is a subgroup of
$\Gamma(X, \mathcal{F})$. The same lemma guarantees that
the assignment $\mathcal{F} \mapsto \Gamma_Z(X, \mathcal{F})$
is a functor in $\mathcal{F}$.
This functor is left exact but not exact in general.

\medskip\noindent
Since the category of abelian sheaves has enough injectives
(Injectives, Lemma \ref{injectives-lemma-abelian-sheaves-space})
we we obtain a right derived functor
$$
R\Gamma_Z(X, -) : D^+(X) \longrightarrow D^+(\textit{Ab})
$$
by
Derived Categories, Lemma \ref{derived-lemma-enough-injectives-right-derived}.
The value of $R\Gamma_Z(X, -)$ on an object $K$ is computed by representing
$K$ by a bounded below complex $\mathcal{I}^\bullet$ of injective abelian
sheaves and taking $\Gamma_Z(X, \mathcal{I}^\bullet)$, see
Derived Categories, Lemma \ref{derived-lemma-injective-acyclic}.
The cohomology groups of an abelian sheaf $\mathcal{F}$
with support in $Z$ defined by
$H^q_Z(X, \mathcal{F}) = R^q\Gamma_Z(X, \mathcal{F})$.

\medskip\noindent
Let $\mathcal{I}$ be an injective abelian sheaf on $X$. Let
$U = X \setminus Z$. Then the restriction map
$\mathcal{I}(X) \to \mathcal{I}(U)$ is surjective
(Lemma \ref{lemma-injective-restriction-surjective})
with kernel $\Gamma_Z(X, \mathcal{I})$. It immediately follows that
for $K \in D^+(X)$ there is a distinguished triangle
$$
R\Gamma_Z(X, K) \to R\Gamma(X, K) \to R\Gamma(U, K) \to R\Gamma_Z(X, K)[1]
$$
in $D^+(\textit{Ab})$. As a consequence we obtain a long exact cohomology
sequence
$$
\ldots \to H^i_Z(X, K) \to H^i(X, K) \to H^i(U, K) \to
H^{i + 1}_Z(X, K) \to \ldots
$$
for any $K$ in $D^+(X)$.

\medskip\noindent
For an abelian sheaf $\mathcal{F}$ on $X$ we can consider the
{\it subsheaf of sections with support in $Z$}, denoted
$\mathcal{H}_Z(\mathcal{F})$, defined by the rule
$$
\mathcal{H}_Z(\mathcal{F})(U) =
\{s \in \mathcal{F}(U) \mid \text{Supp}(s) \subset U \cap Z\} =
\Gamma_{Z \cap U}(U, \mathcal{F}|_U)
$$
Using the equivalence of Modules, Lemma \ref{modules-lemma-i-star-exact}
we may view $\mathcal{H}_Z(\mathcal{F})$ as an abelian sheaf on $Z$, see
Modules, Remark \ref{modules-remark-sections-support-in-closed}.
Thus we obtain a functor
$$
\textit{Ab}(X) \longrightarrow \textit{Ab}(Z),\quad
\mathcal{F} \longmapsto
\mathcal{H}_Z(\mathcal{F})\text{ viewed as a sheaf on }Z
$$
This functor is left exact, but in general not exact. Exactly as above
we obtain a right derived functor
$$
R\mathcal{H}_Z : D^+(X) \longrightarrow D^+(Z)
$$
the derived functor. We set
$\mathcal{H}^q_Z(\mathcal{F}) = R^q\mathcal{H}_Z(\mathcal{F})$ so that
$\mathcal{H}^0_Z(\mathcal{F}) = \mathcal{H}_Z(\mathcal{F})$.

\medskip\noindent
Observe that we have
$\Gamma_Z(X, \mathcal{F}) = \Gamma(Z, \mathcal{H}_Z(\mathcal{F}))$
for any abelian sheaf $\mathcal{F}$. By
Lemma \ref{lemma-sections-with-support-acyclic} below
the functor $\mathcal{H}_Z$ transforms injective abelian sheaves
into sheaves right acyclic for $\Gamma(Z, -)$. Thus by
Derived Categories, Lemma \ref{derived-lemma-grothendieck-spectral-sequence}
we obtain a convergent Grothendieck spectral sequence
$$
E_2^{p, q} = H^p(Z, \mathcal{H}^q_Z(K)) \Rightarrow H^{p + q}_Z(X, K)
$$
functorial in $K$ in $D^+(X)$.

\begin{lemma}
\label{lemma-sections-with-support-acyclic}
Let $i : Z \to X$ be the inclusion of a closed subset.
Let $\mathcal{I}$ be an injective abelian sheaf on $X$.
Then $\mathcal{H}_Z(\mathcal{I})$ is an injective abelian sheaf on $Z$.
\end{lemma}

\begin{proof}
This follows from
Homology, Lemma \ref{homology-lemma-adjoint-preserve-injectives}
as $\mathcal{H}_Z(-)$ is right adjoint to the exact functor $i_*$.
See Modules, Lemmas \ref{modules-lemma-i-star-exact} and
\ref{modules-lemma-i-star-right-adjoint}.
\end{proof}




\section{Cohomology with supports in a family}
\label{section-cohomology-family-supports}

\begin{definition}
\label{definition-family-supports}
\begin{reference}
\cite[Chapter III, \S 3.2, pp. 156--157]{Tohoku}
\end{reference}
Let $X$ be a topological space. If $A \subset X$ is a subset and
$\mathcal{F}$ is a sheaf on $X$, we write
$\Gamma(A, \mathcal{F}) = \Gamma(A, \mathcal{F}|_A)$.
A {\it family of supports} on $X$ is a nonempty collection $\Phi$ of
closed subsets of $X$ such that
\begin{enumerate}
\item if $Z \in \Phi$ and $Z' \subset Z$ is closed, then $Z' \in \Phi$,
\item if $Z_1, Z_2 \in \Phi$, then $Z_1 \cup Z_2 \in \Phi$.
\end{enumerate}
For an abelian sheaf $\mathcal{F}$ on $X$ we set
$$
\Gamma_\Phi(X, \mathcal{F}) =
\{s \in \Gamma(X, \mathcal{F}) \mid \text{Supp}(s) \in \Phi\}.
$$
Equivalently,
$$
\Gamma_\Phi(X, \mathcal{F}) =
\colim_{Z \in \Phi} \Gamma_Z(X, \mathcal{F}).
$$
If $Z \subset X$ is closed, put
$$
\Phi_Z = \{Z' \subset X \mid Z' \text{ is closed and } Z' \subset Z\}.
$$
Then $\Gamma_{\Phi_Z}(X, -) = \Gamma_Z(X, -)$. If $\Phi$ is the family of all
closed subsets of $X$, then
$\Gamma_\Phi(X, -) = \Gamma(X, -)$.
\end{definition}

\begin{lemma}
\label{lemma-cohomology-family-supports}
\begin{reference}
\cite[Chapter III, \S 3.2, Proposition 3.2.1, p. 157]{Tohoku}
\end{reference}
Let $X$ be a topological space and let $\Phi$ be a family of supports on
$X$.
\begin{enumerate}
\item The functor
$$
\Gamma_\Phi(X, -) : \textit{Ab}(X) \longrightarrow \textit{Ab}
$$
is left exact.
\item Its right derived functor
$$
R\Gamma_\Phi(X, -) : D^+(X) \longrightarrow D^+(\textit{Ab})
$$
is everywhere defined. We set
$$
H^p_\Phi(X, \mathcal{F}) = R^p\Gamma_\Phi(X, \mathcal{F}).
$$
These groups vanish for $p < 0$, agree with
$\Gamma_\Phi(X, \mathcal{F})$ for $p = 0$, and for $p \geq 0$ form a
universal cohomological $\delta$-functor. They vanish for $p > 0$ when
$\mathcal{F}$ is injective. Conversely, these properties characterize
$H^p_\Phi(X, -)$ up to unique isomorphism of $\delta$-functors.
\item If $\mathcal{F} \to \mathcal{C}^\bullet$ is a nonnegative
resolution by sheaves right acyclic for $\Gamma_\Phi(X, -)$, then
$$
H^p_\Phi(X, \mathcal{F}) =
H^p(\Gamma_\Phi(X, \mathcal{C}^\bullet)).
$$
In particular, this applies to an injective resolution.
\end{enumerate}
\end{lemma}

\begin{proof}
Modules, Lemma \ref{modules-lemma-support-section-closed} shows that
$\Gamma_\Phi(X, \mathcal{F})$ is a subgroup of
$\Gamma(X, \mathcal{F})$ and that $\Gamma_\Phi(X, -)$ is an additive
functor.
Consider a short exact sequence
$0 \to \mathcal{F}' \to \mathcal{F} \to \mathcal{F}'' \to 0$.
Global sections are left exact. If a section of $\mathcal{F}$ supported
in a member of $\Phi$ comes from a section of $\mathcal{F}'$, then the
two sections have the same support because
$\mathcal{F}' \to \mathcal{F}$ is injective on stalks. Thus
$\Gamma_\Phi(X, -)$ is left exact.

The category $\textit{Ab}(X)$ has enough injectives by Injectives, Lemma
\ref{injectives-lemma-abelian-sheaves-space}. Hence the right derived
functor is everywhere defined by Derived Categories, Lemma
\ref{derived-lemma-enough-injectives-right-derived}. The assertions about
negative degrees and degree zero follow from Derived Categories, Lemma
\ref{derived-lemma-left-exact-higher-derived}. Injective objects are right
acyclic by Derived Categories, Lemma
\ref{derived-lemma-injective-acyclic}. The universal $\delta$-functor
assertion follows from Derived Categories, Lemma
\ref{derived-lemma-right-derived-delta-functor}. Finally, the assertion
about acyclic resolutions is Leray's acyclicity lemma, Derived Categories,
Lemma \ref{derived-lemma-leray-acyclicity}.

If another $\delta$-functor has the listed properties, then its positive
terms are effaceable by Homology, Lemma
\ref{homology-lemma-characterize-effaceable-functor}. It is universal by
Homology, Lemma
\ref{homology-lemma-efface-implies-universal}, and hence uniquely isomorphic
to the one above by Homology, Lemma
\ref{homology-lemma-uniqueness-universal-delta-functor}.
\end{proof}

\begin{lemma}
\label{lemma-family-supports-pullback}
\begin{reference}
\cite[Chapter III, \S 3.2, Proposition 3.2.2, pp. 157--158]{Tohoku}
\end{reference}
Let $f : X \to Y$ be a continuous map. Let $\Phi$ and $\Psi$ be families
of supports on $X$ and $Y$, respectively. Assume that
$$
B \in \Psi \quad \Longrightarrow \quad f^{-1}(B) \in \Phi.
$$
For every $p \in \mathbf{Z}$ there is a natural map
$$
H^p_\Psi(Y, \mathcal{F}) \longrightarrow
H^p_\Phi(X, f^{-1}\mathcal{F}).
$$
These maps form the unique morphism of cohomological $\delta$-functors
whose degree zero part is pullback of sections. They are transitive for
compositions of continuous maps and compatible families of supports.
\end{lemma}

\begin{proof}
The adjunction map $\mathcal{F} \to f_*f^{-1}\mathcal{F}$ induces a
natural map
$$
\Gamma(Y, \mathcal{F}) \longrightarrow
\Gamma(X, f^{-1}\mathcal{F}).
$$
By Sheaves, Lemma \ref{sheaves-lemma-pullback-abelian-stalk}, the support
of the pulled back section is the inverse image of the support of the
original section. Thus a section on the left supported on $B \in \Psi$
pulls back to a section supported on $f^{-1}(B) \in \Phi$. We obtain a
natural transformation
$$
\Gamma_\Psi(Y, -) \longrightarrow
\Gamma_\Phi(X, f^{-1}(-)).
$$
Pullback of abelian sheaves is exact by Modules, Lemma
\ref{modules-lemma-exactness-pushforward-pullback}. Therefore the functors
$\mathcal{F} \mapsto H^p_\Phi(X, f^{-1}\mathcal{F})$ form a
cohomological $\delta$-functor. The source $\delta$-functor is universal
by Lemma \ref{lemma-cohomology-family-supports}, so its degree zero map
extends uniquely to a morphism of $\delta$-functors. For two composable
continuous maps, the composite and the directly constructed morphism
extend the same pullback map in degree zero. Universality makes them equal,
which proves transitivity. In negative degrees both sides vanish and the
map is the unique map between zero groups.
\end{proof}

\begin{lemma}
\label{lemma-family-supports-restriction}
\begin{reference}
\cite[Chapter III, \S 3.2, p. 158]{Tohoku}
\end{reference}
Let $X$ be a topological space, let $Y \subset X$ be a subset, and let
$\Phi$ be a family of supports on $X$. Then
$$
\Phi|_Y = \{Z \cap Y \mid Z \in \Phi\}
$$
is a family of supports on $Y$, and there are natural restriction maps
$$
H^p_\Phi(X, \mathcal{F}) \longrightarrow
H^p_{\Phi|_Y}(Y, \mathcal{F}|_Y)
$$
for every $p \in \mathbf{Z}$.
\end{lemma}

\begin{proof}
The family axioms for $\Phi|_Y$ follow from those for $\Phi$ and the
description of the closed subsets of the subspace $Y$. Apply Lemma
\ref{lemma-family-supports-pullback} to the inclusion $Y \to X$.
\end{proof}

\begin{definition}
\label{definition-soft-family-supports}
\begin{reference}
\cite[Chapter III, \S 3.3, p. 159]{Tohoku}
\end{reference}
Let $X$ be a topological space.
\begin{enumerate}
\item A sheaf of sets $\mathcal{F}$ on $X$ is {\it soft} if the restriction
map
$$
\Gamma(X, \mathcal{F}) \longrightarrow \Gamma(A, \mathcal{F}|_A)
$$
is surjective for every closed subset $A \subset X$.
\item Let $\Phi$ be a family of supports on $X$. An abelian sheaf
$\mathcal{F}$ on $X$ is {\it $\Phi$-soft} if
$$
\Gamma_\Phi(X, \mathcal{F}) \longrightarrow
\Gamma(A, \mathcal{F}|_A)
$$
is surjective for every $A \in \Phi$.
\item A family of supports $\Phi$ is {\it paracompactifying} if every
$A \in \Phi$ is paracompact and admits a closed neighbourhood
$B \in \Phi$; here this means $A \subset B$ and $A$ is contained in the
interior of $B$.
\end{enumerate}
\end{definition}

\begin{definition}[Fine abelian sheaves]
\label{definition-fine-abelian-sheaf}
\begin{reference}
\cite[Chapter III, \S 3.3, p. 160]{Tohoku}
\end{reference}
Let $X$ be a topological space and let $\mathcal F$ be an abelian sheaf
on $X$. For an endomorphism $e:\mathcal F\to\mathcal F$, put
$$
\mathop{\rm Supp}(e)=
\overline{\{x\in X\mid e_x:\mathcal F_x\to\mathcal F_x
\text{ is nonzero}\}}.
$$
We say that $\mathcal F$ is \emph{fine} if, for every locally finite open
covering $X=\bigcup_{i\in I}U_i$, there are endomorphisms
$e_i:\mathcal F\to\mathcal F$ such that
$$
\mathop{\rm Supp}(e_i)\subset U_i
\quad\text{and}\quad
\sum_{i\in I}e_i=\mathrm{id}_{\mathcal F}.
$$
The sum is locally finite, since the supports are subordinate to the
locally finite covering.
\end{definition}

\begin{lemma}
\label{lemma-fine-phi-soft-and-tensor}
Let $X$ be a topological space.
\begin{enumerate}
\item If $\Phi$ is a paracompactifying family of supports on $X$, then
every fine abelian sheaf on $X$ is $\Phi$-soft.
\item If $\mathcal C$ is a fine abelian sheaf and $\mathcal F$ is any
abelian sheaf, then $\mathcal F\otimes_{\mathbf Z}\mathcal C$ is fine.
\end{enumerate}
\end{lemma}

\begin{proof}
Let $A\in\Phi$ and let
$s_A\in\Gamma(A,\mathcal F|_A)$. Choose a closed neighbourhood
$B\in\Phi$ of $A$. By
Lemma \ref{lemma-section-closed-paracompact-neighbourhood}, the section
$s_A$ extends to a section $s\in\mathcal F(V)$ for an open neighbourhood
$V$ of $A$. Shrinking $V$, we may assume $V\subset\mathop{\rm Int}(B)$.

Apply fineness to the two-element open covering
$$
X=V\cup(X\setminus A).
$$
Write the resulting endomorphisms as $e_V$ and $e_0$. The section
$e_V(s)$ over $V$ has support contained in the closed set
$\mathop{\rm Supp}(e_V)\subset V$. It therefore glues with the zero
section over $X\setminus\mathop{\rm Supp}(e_V)$ and extends to a section
$\widetilde s$ over $X$. On $A$ the endomorphism $e_0$ is zero, so
$e_V|_A=\mathrm{id}$ and $\widetilde s|_A=s_A$. Moreover,
$$
\mathop{\rm Supp}(\widetilde s)
\subset\mathop{\rm Supp}(e_V)\subset V\subset B.
$$
Thus $\widetilde s$ belongs to $\Gamma_\Phi(X,\mathcal F)$, which proves
(1).

For (2), let $e_i$ be a fine partition of the identity on $\mathcal C$
subordinate to a locally finite open covering $\{U_i\}$. The
endomorphisms
$$
\mathrm{id}_{\mathcal F}\otimes e_i:
\mathcal F\otimes_{\mathbf Z}\mathcal C
\longrightarrow
\mathcal F\otimes_{\mathbf Z}\mathcal C
$$
have support contained in $\mathop{\rm Supp}(e_i)\subset U_i$ and their
locally finite sum is the identity. Hence
$\mathcal F\otimes_{\mathbf Z}\mathcal C$ is fine.
\end{proof}

\begin{lemma}
\label{lemma-section-closed-paracompact-neighbourhood}
\begin{reference}
\cite[Chapter III, \S 3.3, p. 159]{Tohoku}
\end{reference}
Let $X$ be a topological space, let $A \subset X$ be closed, and let
$\mathcal{F}$ be a sheaf of sets on $X$. If $A$ has a paracompact
neighbourhood, then every section of $\mathcal{F}|_A$ over $A$ extends
to a section of $\mathcal{F}$ over an open neighbourhood of $A$ in $X$.
\end{lemma}

\begin{proof}
After replacing $X$ by a paracompact neighbourhood of $A$ and at the end
intersecting the open obtained with the interior of that neighbourhood,
we may assume that $X$ is paracompact. Let
$s \in \Gamma(A, \mathcal{F}|_A)$. There is an open covering
$A \subset \bigcup_{i \in I} U_i$ and sections
$s_i \in \mathcal{F}(U_i)$ which induce $s$ over $A \cap U_i$.
Add $U_0 = X \setminus A$ to obtain an open covering of $X$.

By Topology, Lemma \ref{topology-lemma-paracompact-shrinking}, there are
a locally finite open covering $X = \bigcup_{j \in J} V_j$, a map
$\alpha : J \to I \cup \{0\}$, and an open covering
$X = \bigcup_{j \in J} W_j$ such that
$$
\overline{W_j} \subset V_j \subset U_{\alpha(j)}.
$$
Let $J' \subset J$ be the set of $j$ such that $W_j$ meets $A$.
Then $\alpha(j) \neq 0$ for $j \in J'$, and the $W_j$, $j \in J'$,
cover $A$. For $j,k \in J'$, the sections
$s_{\alpha(j)}$ and $s_{\alpha(k)}$ agree over $A \cap V_j \cap V_k$.
Thus there is an open subset $E_{jk} \subset V_j \cap V_k$ containing
$A \cap V_j \cap V_k$ on which they agree.

The family $\{\overline{W_j}\}_{j \in J'}$ is locally finite. Hence so is
the family of closed subsets
$$
C_{jk} =
(\overline{W_j} \cap \overline{W_k}) \setminus E_{jk},
\qquad j,k \in J',
$$
and its union $C$ is closed. The set $C$ is disjoint from $A$. Therefore
$$
U = \left(\bigcup_{j \in J'} W_j\right) \setminus C
$$
is an open neighbourhood of $A$. The sections
$s_{\alpha(j)}|_{U \cap W_j}$ agree on pairwise intersections by the
construction of $C$. They glue to a section of $\mathcal{F}$ over $U$
which restricts to $s$ over $A$.
\end{proof}

\begin{lemma}
\label{lemma-flasque-soft-family-supports}
Let $X$ be a topological space.
\begin{enumerate}
\item If $X$ is paracompact, then every flasque sheaf of sets on $X$ is
soft.
\item If $\Phi$ is a paracompactifying family of supports on $X$, then
every flasque abelian sheaf on $X$ is $\Phi$-soft.
\item Every abelian sheaf on $X$ embeds in a flasque abelian sheaf. Hence,
if $\Phi$ is paracompactifying, every abelian sheaf embeds in a
$\Phi$-soft abelian sheaf.
\end{enumerate}
\end{lemma}

\begin{proof}
Let $\mathcal{F}$ be flasque. If $X$ is paracompact, then Lemma
\ref{lemma-section-closed-paracompact-neighbourhood} extends a section over
a closed subset $A$ to a section over an open neighbourhood of $A$.
Flasqueness extends the latter section over $X$. This proves (1).

Assume $\Phi$ is paracompactifying. Let $A \in \Phi$ and let
$s \in \Gamma(A, \mathcal{F}|_A)$. Choose a closed neighbourhood
$B \in \Phi$ of $A$. By definition $B$ is paracompact, so Lemma
\ref{lemma-section-closed-paracompact-neighbourhood} gives an open
$U \subset X$ containing $A$ and a section $t \in \mathcal{F}(U)$
inducing $s$. After shrinking $U$ we may assume $U \subset \mathop{\rm
Int}(B)$. Extend $t$ to a section $\widetilde t \in \mathcal{F}(X)$.
On the disjoint union $U \amalg (X \setminus B)$ define a section $r$
which is zero over $U$ and equal to $\widetilde t$ over $X \setminus B$.
Extend $r$ to a section $\widetilde r \in \mathcal{F}(X)$. Then
$\widetilde t - \widetilde r$ induces $s$ over $A$ and has support
contained in $B \in \Phi$. This proves (2).

Finally, the category of abelian sheaves on $X$ has enough injectives by
Injectives, Lemma \ref{injectives-lemma-abelian-sheaves-space}, and every
injective abelian sheaf is flasque by Lemma
\ref{lemma-injective-flasque}. This proves (3).
\end{proof}


\begin{lemma}
\label{lemma-glue-locally-finite-closed}
Let $X$ be a topological space, let $\{A_i\}_{i \in I}$ be a locally
finite closed covering of $X$, and let $\mathcal{F}$ be a sheaf of sets.
Sections
$$
s_i \in \Gamma(A_i, \mathcal{F}|_{A_i})
$$
which agree over every $A_i \cap A_j$ glue uniquely to a section of
$\mathcal{F}$ over $X$.
\end{lemma}

\begin{proof}
Use the description of sheaves by \'etale spaces from
Sheaves, Lemma
\ref{sheaves-lemma-sheaves-etale-spaces-equivalence}.
Thus the $s_i$ are continuous sections of a fixed \'etale space
$p : E \to X$. They define a unique set-theoretic section
$s : X \to E$. If $C \subset E$ is closed, then
$$
s^{-1}(C) = \bigcup_{i \in I}s_i^{-1}(C).
$$
Each set on the right is closed in $A_i$, hence in $X$, and these closed
sets form a locally finite family. Their union is therefore closed.
Thus $s$ is continuous, and hence is the desired section.
\end{proof}

\begin{lemma}
\label{lemma-extend-zero-from-closed}
Let $B \subset X$ be closed, let $\mathcal{F}$ be an abelian sheaf on
$X$, and let
$$
t \in \Gamma(B, \mathcal{F}|_B).
$$
If the support of $t$, computed on $B$, is contained in
$\mathop{\rm Int}(B)$, then $t$ extends by zero to a section
$\widetilde t \in \Gamma(X, \mathcal{F})$. Its support is the support
of $t$, viewed as a closed subset of $X$.
\end{lemma}

\begin{proof}
Let $S$ be the support of $t$. Since $B$ is closed, $S$ is closed in
$X$. The opens $\mathop{\rm Int}(B)$ and $X \setminus S$ cover $X$.
On the first use $t$, and on the second use the zero section. These
sections agree on the intersection because $t$ is zero off $S$.
They glue to $\widetilde t$. Its restriction to $B$ is $t$, and the
assertion about support follows stalkwise.
\end{proof}

\begin{lemma}
\label{lemma-soft-exact-paracompact}
Let $X$ be paracompact and let
$$
0 \longrightarrow \mathcal{K} \longrightarrow \mathcal{F}
\longrightarrow \mathcal{G} \longrightarrow 0
$$
be a short exact sequence of abelian sheaves on $X$.
\begin{enumerate}
\item If $\mathcal{K}$ is soft, then
$$
0 \longrightarrow \Gamma(X, \mathcal{K})
\longrightarrow \Gamma(X, \mathcal{F})
\longrightarrow \Gamma(X, \mathcal{G})
\longrightarrow 0
$$
is exact.
\item If $\mathcal{K}$ and $\mathcal{F}$ are soft, then
$\mathcal{G}$ is soft.
\end{enumerate}
\end{lemma}

\begin{proof}
Only surjectivity in (1) requires proof. Let
$s \in \Gamma(X, \mathcal{G})$. There is an open covering
$X = \bigcup U_i$ and sections $u_i \in \mathcal{F}(U_i)$ lifting
$s|_{U_i}$. By Topology, Lemma
\ref{topology-lemma-paracompact-shrinking}, there is a locally finite
closed covering $\{A_j\}_{j \in J}$ and a map $\alpha : J \to I$ such
that $A_j \subset U_{\alpha(j)}$. Namely, one may take
$A_j = \overline{W_j}$ in that lemma.

Well order $J$. We construct by transfinite recursion compatible
sections
$$
t_j \in \Gamma(A_j, \mathcal{F}|_{A_j})
$$
lifting $s|_{A_j}$. Suppose the sections have been constructed for the
predecessors of $j$. Their union
$$
C_j = \bigcup_{k < j} A_k
$$
is closed, since it is the union of a subfamily of a locally finite
family of closed subsets. By Lemma
\ref{lemma-glue-locally-finite-closed}, the preceding sections glue to a
section $t_{C_j}$ over $C_j$.

Put $D_j = A_j \cap C_j$ and initially take
$u_{\alpha(j)}|_{A_j}$. The difference
$$
u_{\alpha(j)}|_{D_j} - t_{C_j}|_{D_j}
$$
maps to zero in $\mathcal{G}|_{D_j}$, and hence comes from a section of
$\mathcal{K}|_{D_j}$. Since $D_j$ is closed and $\mathcal{K}$ is soft,
extend this section to a section of $\mathcal{K}$ over $X$ and subtract
its image from $u_{\alpha(j)}|_{A_j}$. The resulting section $t_j$
agrees with all preceding sections. At the end, Lemma
\ref{lemma-glue-locally-finite-closed} glues the $t_j$ to a global lift
of $s$. This proves (1).

For (2), let $A \subset X$ be closed and let
$s_A \in \Gamma(A, \mathcal{G}|_A)$. The closed subspace $A$ is
paracompact. Moreover, $\mathcal{K}|_A$ is soft: every closed subset of
$A$ is closed in $X$, and an extension over $X$ restricts to one over
$A$. Applying (1) on $A$ lifts $s_A$ to a section of
$\mathcal{F}|_A$. Softness of $\mathcal{F}$ extends that section over
$X$, and its image extends $s_A$. Thus $\mathcal{G}$ is soft.
\end{proof}

\begin{lemma}
\label{lemma-phi-soft-exact}
Let $\Phi$ be a paracompactifying family of supports on $X$, and let
$$
0 \longrightarrow \mathcal{K} \longrightarrow \mathcal{F}
\longrightarrow \mathcal{G} \longrightarrow 0
$$
be exact. If $\mathcal{K}$ and $\mathcal{F}$ are $\Phi$-soft, then
$\mathcal{G}$ is $\Phi$-soft and
$$
0 \longrightarrow \Gamma_\Phi(X, \mathcal{K})
\longrightarrow \Gamma_\Phi(X, \mathcal{F})
\longrightarrow \Gamma_\Phi(X, \mathcal{G})
\longrightarrow 0
$$
is exact.
\end{lemma}

\begin{proof}
For every $B \in \Phi$, the restrictions
$\mathcal{K}|_B$ and $\mathcal{F}|_B$ are soft. Indeed, if
$C \subset B$ is closed, then $C$ is closed in $X$ and belongs to
$\Phi$; a $\Phi$-supported extension over $X$ restricts to an extension
over $B$. Since $B$ is paracompact, Lemma
\ref{lemma-soft-exact-paracompact} shows that
$\mathcal{G}|_B$ is soft and that
$$
\Gamma(B, \mathcal{F}|_B) \longrightarrow
\Gamma(B, \mathcal{G}|_B)
$$
is surjective.

Let $s \in \Gamma_\Phi(X, \mathcal{G})$ and put
$A = \text{Supp}(s) \in \Phi$. Choose a closed neighbourhood
$B \in \Phi$ of $A$. Lift $s|_B$ to
$t \in \Gamma(B, \mathcal{F}|_B)$. Put
$$
D = B \setminus \mathop{\rm Int}(B).
$$
The image of $t|_D$ is zero because $D \cap A = \emptyset$. Hence
$t|_D$ comes from a section of $\mathcal{K}|_D$. Softness of
$\mathcal{K}|_B$ extends it to a section $r$ over $B$. Replacing $t$
by $t-r$, we obtain a lift of $s|_B$ which is zero over $D$. Its
support is therefore contained in $\mathop{\rm Int}(B)$, so Lemma
\ref{lemma-extend-zero-from-closed} extends it by zero to a global lift
of $s$. Its support is a closed subset of $B$, hence belongs to $\Phi$.
This proves the required surjectivity.

It remains to prove that $\mathcal{G}$ is $\Phi$-soft. Let
$A \in \Phi$ and let
$u \in \Gamma(A, \mathcal{G}|_A)$. Choose a closed neighbourhood
$B \in \Phi$ of $A$ and again put
$D = B \setminus \mathop{\rm Int}(B)$. On the disjoint closed subsets
$A,D \subset B$, take the sections $u$ and $0$. By Lemma
\ref{lemma-glue-locally-finite-closed}, they define a section over
$A \cup D$. Since $\mathcal{G}|_B$ is soft, this extends to a section
over $B$. It is zero over $D$, so Lemma
\ref{lemma-extend-zero-from-closed} extends it by zero to a section in
$\Gamma_\Phi(X, \mathcal{G})$ inducing $u$ over $A$.
\end{proof}

\begin{lemma}
\label{lemma-family-supports-acyclic-classes}
\begin{reference}
\cite[Chapter III, \S 3.3, Proposition 3.3.2, p. 159]{Tohoku}
\end{reference}
Let $\Phi$ be a family of supports on $X$. Each of the following choices
of a class $\mathcal{M}$ satisfies the hypotheses of Derived Categories,
Lemma \ref{derived-lemma-subcategory-right-acyclics-injectives} for the
functor $\Gamma_\Phi(X,-)$:
\begin{enumerate}
\item $\mathcal{M}$ is the class of flasque abelian sheaves;
\item $\Phi$ is paracompactifying and $\mathcal{M}$ is the class of
$\Phi$-soft abelian sheaves.
\end{enumerate}
In other words, every abelian sheaf embeds in an object of
$\mathcal{M}$, the class is closed under direct summands, and if
$$
0 \longrightarrow \mathcal{F}' \longrightarrow \mathcal{F}
\longrightarrow \mathcal{F}'' \longrightarrow 0
$$
is exact with $\mathcal{F}',\mathcal{F} \in \mathcal{M}$, then
$\mathcal{F}'' \in \mathcal{M}$ and applying $\Gamma_\Phi(X,-)$ leaves
the sequence exact.
\end{lemma}

\begin{proof}
The embedding assertion follows from Lemma
\ref{lemma-flasque-soft-family-supports}. Both classes are closed under
direct summands: extend a section after including it into the direct sum
and then project the resulting extension back to the summand. Projection
cannot enlarge support.

For flasque sheaves, apply Lemma \ref{lemma-flasque-acyclic} to the
ringed space $(X, \underline{\mathbf{Z}})$. Its proof shows that
$\mathcal{F}''$ is flasque and that sections over every open subset remain
exact. Let $s \in \Gamma_\Phi(X,\mathcal{F}'')$ and put
$A = \text{Supp}(s)$. Lift $s$ to
$t \in \Gamma(X,\mathcal{F})$. Over $U = X \setminus A$, the section
$t|_U$ comes from a section of $\mathcal{F}'|_U$. Extend the latter over
$X$ using flasqueness of $\mathcal{F}'$ and subtract it from $t$.
The resulting lift of $s$ is supported in $A \in \Phi$. Thus applying
$\Gamma_\Phi(X,-)$ is exact.

For $\Phi$-soft sheaves, the quotient and exactness assertions are
Lemma \ref{lemma-phi-soft-exact}.
\end{proof}

\begin{lemma}
\label{lemma-family-supports-flasque-soft-acyclic}
\begin{reference}
\cite[Chapter III, \S 3.3, Corollary to Proposition 3.3.2,
p. 159]{Tohoku}
\end{reference}
Let $\Phi$ be a family of supports on $X$.
\begin{enumerate}
\item If $\mathcal{F}$ is flasque, then
$H^p_\Phi(X,\mathcal{F}) = 0$ for $p > 0$.
\item If $\Phi$ is paracompactifying and $\mathcal{F}$ is
$\Phi$-soft, then $H^p_\Phi(X,\mathcal{F}) = 0$ for $p > 0$.
\item Consequently, a nonnegative resolution by flasque sheaves, or by
$\Phi$-soft sheaves when $\Phi$ is paracompactifying, computes
$H^p_\Phi(X,-)$ after applying $\Gamma_\Phi(X,-)$.
\end{enumerate}
\end{lemma}

\begin{proof}
The first two assertions follow from Lemma
\ref{lemma-family-supports-acyclic-classes} and Derived Categories,
Lemma \ref{derived-lemma-subcategory-right-acyclics-injectives}.
The last assertion follows from Lemma
\ref{lemma-cohomology-family-supports}.
\end{proof}


\begin{lemma}[Cartan resolution]
\label{lemma-cartan-resolution-fine-torsion-free}
\begin{reference}
\cite[Chapter III, \S 3.3, p. 160]{Tohoku}
\end{reference}
Let $X$ be a topological space, let $\Phi$ be a paracompactifying family
of supports on $X$, and suppose given an exact augmented complex
$$
0\longrightarrow\underline{\mathbf Z}\longrightarrow
\mathcal C^0\longrightarrow\mathcal C^1\longrightarrow\cdots
$$
of abelian sheaves on $X$ such that every $\mathcal C^n$ is fine and has
torsion-free stalks. For every abelian sheaf $\mathcal F$ on $X$, the
augmented complex
$$
0\longrightarrow\mathcal F\longrightarrow
\mathcal F\otimes_{\mathbf Z}\mathcal C^0\longrightarrow
\mathcal F\otimes_{\mathbf Z}\mathcal C^1\longrightarrow\cdots
$$
is an exact functorial resolution. Its terms are fine, hence
$\Phi$-soft and $\Gamma_\Phi(X,-)$-acyclic. Consequently there are
canonical functorial isomorphisms
$$
H^n_\Phi(X,\mathcal F)
\cong
H^n\bigl(\Gamma_\Phi
(X,\mathcal F\otimes_{\mathbf Z}\mathcal C^\bullet)\bigr),
\qquad n\geq0.
$$
Moreover, the assignment
$$
\mathcal F\longmapsto
\Gamma_\Phi(X,\mathcal F\otimes_{\mathbf Z}\mathcal C^\bullet)
$$
is an exact resolving functor for $\Gamma_\Phi(X,-)$. We call
$\mathcal F\otimes_{\mathbf Z}\mathcal C^\bullet$ the \emph{Cartan
resolution of $\mathcal F$ relative to $\mathcal C^\bullet$}. This
construction is conditional on the displayed fine torsion-free resolution;
it makes no assertion that such a resolution exists without the stated
hypothesis.
\end{lemma}

\begin{proof}
A torsion-free abelian group is flat over $\mathbf Z$ by More on Algebra,
Lemma \ref{more-algebra-lemma-dedekind-torsion-free-flat}. Flatness of a
sheaf of modules can be checked on stalks by Modules, Lemma
\ref{modules-lemma-flat-stalks-flat}. Hence every $\mathcal C^n$ is flat.

Let $\mathcal Z^0=\underline{\mathbf Z}$ and, for $n\geq0$, let
$\mathcal Z^{n+1}$ be the image of
$\mathcal C^n\to\mathcal C^{n+1}$. Exactness of the augmented complex
gives short exact sequences
$$
0\longrightarrow\mathcal Z^n\longrightarrow\mathcal C^n
\longrightarrow\mathcal Z^{n+1}\longrightarrow0.
$$
Every stalk of $\mathcal Z^{n+1}$ is a subgroup of the torsion-free group
$\mathcal C^{n+1}_x$, and is therefore torsion-free and flat over
$\mathbf Z$. Thus every $\mathcal Z^n$ is a flat abelian sheaf. Tensoring
the displayed short exact sequences with an arbitrary $\mathcal F$
preserves exactness. Since
$\mathcal F\otimes_{\mathbf Z}\underline{\mathbf Z}\cong\mathcal F$,
these tensor sequences splice to the asserted exact augmented resolution.
Flatness of each $\mathcal C^n$ also shows that this resolution is exact
as a functor of $\mathcal F$.

Lemma \ref{lemma-fine-phi-soft-and-tensor} shows that every term
$\mathcal F\otimes\mathcal C^n$ is fine and $\Phi$-soft. Such sheaves are
$\Gamma_\Phi(X,-)$-acyclic by
Lemma \ref{lemma-family-supports-flasque-soft-acyclic}. Therefore the
resolution computes supported cohomology, by the same lemma, and the
displayed isomorphisms are canonical and functorial. The exactness in
$\mathcal F$ gives, for every short exact sequence of sheaves, a termwise
short exact sequence of fine, hence $\Phi$-soft, resolution terms. Lemma
\ref{lemma-phi-soft-exact} shows that applying $\Gamma_\Phi(X,-)$ leaves
each of these sequences exact. Together with the augmentation in degree
zero and the acyclicity of the terms, this gives the final
resolving-functor assertion.
\end{proof}



\begin{lemma}
\label{lemma-extension-by-zero-locally-closed}
\begin{reference}
\cite[Chapter III, \S 3.5, p. 166]{Tohoku}
\end{reference}
Let $X$ be a topological space and let $k : Z \to X$ be the inclusion of
a locally closed subset. There is an exact functor
$$
k_! : \textit{Ab}(Z) \longrightarrow \textit{Ab}(X)
$$
characterized, up to unique isomorphism, by the following properties:
$k^{-1}k_!\mathcal{G}=\mathcal{G}$ and the restriction of
$k_!\mathcal{G}$ to $X \setminus Z$ is zero. More precisely, if
$$
Z \xrightarrow{j} C \xrightarrow{i} X
$$
is a factorization of $k$ with $j$ open and $i$ closed, then
$k_!\mathcal{G}=i_*j_!\mathcal{G}$.

These functors are transitive for inclusions of locally closed subsets.
For an abelian sheaf $\mathcal{F}$ on $X$ put
$$
\mathcal{F}_Z=k_!(\mathcal{F}|_Z).
$$
The functor $\mathcal{F}\mapsto\mathcal{F}_Z$ is exact. If $Y\subset X$
is closed, $U=X\setminus Y$, and $i:Y\to X$, $j:U\to X$ are the
inclusions, then there is a functorial short exact sequence
$$
0\longrightarrow \mathcal{F}_U\longrightarrow \mathcal{F}
\longrightarrow \mathcal{F}_Y\longrightarrow 0.
$$
\end{lemma}

\begin{proof}
A locally closed subset is open in a closed subset, so a factorization as
in the statement exists. The functor $j_!$ is exact by Modules, Lemma
\ref{modules-lemma-j-shriek-exact}, and $i_*$ is exact by Modules, Lemma
\ref{modules-lemma-i-star-exact}. Hence their composite is exact. The
description of its restrictions follows from Sheaves, Lemmas
\ref{sheaves-lemma-j-shriek-abelian} and
\ref{sheaves-lemma-stalks-closed-pushforward}.

Conversely, suppose that an abelian sheaf $\mathcal{H}$ on $X$ has the
prescribed restrictions. Since its support is contained in the closed set
$C$, Modules, Lemma \ref{modules-lemma-i-star-exact} identifies
$\mathcal{H}$ with $i_*(\mathcal{H}|_C)$. On $C$, the canonical exact
sequence of Modules, Lemma
\ref{modules-lemma-canonical-exact-sequence} identifies
$\mathcal{H}|_C$ with $j_!\mathcal{G}$. This also shows uniqueness and
independence of the factorization. Transitivity follows from the same
characterization. Restriction is exact, so the assertion about
$\mathcal{F}_Z$ follows. The final short exact sequence is precisely
Modules, Lemma \ref{modules-lemma-canonical-exact-sequence}.
\end{proof}

\begin{lemma}
\label{lemma-induced-family-supports}
\begin{reference}
\cite[Chapter III, \S 3.5, p. 166]{Tohoku}
\end{reference}
Let $\Phi$ be a family of supports on $X$ and let $Z\subset X$ be locally
closed. Then
$$
\Phi_Z=\{A\subset Z\mid A\in\Phi\}
$$
is a family of supports on $Z$. If $\Phi$ is paracompactifying, then so
is $\Phi_Z$. For every abelian sheaf $\mathcal{G}$ on $Z$, restriction
induces an isomorphism
$$
\Gamma_\Phi(X,k_!\mathcal{G})
\longrightarrow \Gamma_{\Phi_Z}(Z,\mathcal{G}),
$$
where $k:Z\to X$ is the inclusion.
\end{lemma}

\begin{proof}
Every member of $\Phi_Z$ is closed in $X$, hence closed in $Z$. If
$A\in\Phi_Z$ and $A'\subset A$ is closed in $Z$, then $A'$ is closed in
$X$ because $A$ is closed in $X$. Thus $A'\in\Phi_Z$. The other family
axioms are immediate.

Assume that $\Phi$ is paracompactifying and let $A\in\Phi_Z$. Choose a
closed neighbourhood $B\in\Phi$ of $A$ in $X$. Choose a closed subset
$C\subset X$ in which $Z$ is open. The space $B\cap C$ is paracompact.
Apply Topology, Lemma \ref{topology-lemma-paracompact-shrinking} to the
open covering
$$
(B\cap Z)\ \cup\ ((B\cap C)\setminus A)
$$
of $B\cap C$. Take the union $D$ of the closures of those members of the
shrinking which refine $B\cap Z$. Local finiteness shows that $D$ is
closed in $B\cap C$, hence in $X$. Moreover,
$$
A\subset \mathop{\rm Int}_Z(D),\qquad D\subset B\cap Z.
$$
Thus $D\in\Phi_Z$, and it is paracompact as a closed subset of $B$. This
proves that $\Phi_Z$ is paracompactifying.

Finally, a section of $k_!\mathcal{G}$ restricts to a section of
$\mathcal{G}$ with the same support. Conversely, a section of
$\mathcal{G}$ whose support belongs to $\Phi_Z$ extends by zero to a
section of $k_!\mathcal{G}$, because its support is closed in $X$. These
operations are inverse and preserve the indicated support conditions.
\end{proof}

\begin{lemma}
\label{lemma-open-extension-phi-soft}
Let $j:U\to X$ be the inclusion of an open subset and let $\Phi$ be a
paracompactifying family of supports on $X$. If an abelian sheaf
$\mathcal{G}$ on $U$ is $\Phi_U$-soft, then $j_!\mathcal{G}$ is
$\Phi$-soft.
\end{lemma}

\begin{proof}
Let $A\in\Phi$ and let
$$
s\in\Gamma(A,(j_!\mathcal{G})|_A).
$$
Its support $S$, computed on $A$, is closed in $X$, is contained in
$U$, and belongs to $\Phi_U$. By Lemma
\ref{lemma-induced-family-supports}, choose a closed neighbourhood
$D\in\Phi_U$ of $S$ in $U$.

The restriction $\mathcal{G}|_D$ is soft. Indeed, every closed subset of
$D$ belongs to $\Phi_U$, and a $\Phi_U$-supported extension over $U$
restricts to an extension over $D$. On the closed subsets
$$
A\cap D\quad\text{and}\quad D\setminus\mathop{\rm Int}_U(D)
$$
of $D$, prescribe respectively the section induced by $s$ and the zero
section. They agree on the intersection because
$S\subset\mathop{\rm Int}_U(D)$. Lemma
\ref{lemma-glue-locally-finite-closed} gives a section over their union,
and softness extends it over $D$. This extension is zero on
$D\setminus\mathop{\rm Int}_U(D)$. Lemma
\ref{lemma-extend-zero-from-closed}, applied on $U$, extends it by zero
to a section of $\mathcal{G}$ over $U$ supported in $D$. It therefore
defines a section of $j_!\mathcal{G}$ over $X$, supported in
$D\in\Phi$, which induces $s$ over $A$.
\end{proof}

\begin{lemma}
\label{lemma-family-supports-extension-by-zero}
\begin{reference}
\cite[Chapter III, \S 3.5, Theorem 3.5.1, pp. 166--167]{Tohoku}
\end{reference}
Let $\Phi$ be a family of supports on $X$, let $Z\subset X$ be locally
closed, and let $k:Z\to X$ be the inclusion. There are canonical maps
$$
H^p_{\Phi_Z}(Z,\mathcal{G})
\longrightarrow H^p_\Phi(X,k_!\mathcal{G}),
\qquad p\in\mathbf{Z},
$$
which form the unique morphism of cohomological $\delta$-functors whose
degree zero map is the isomorphism of Lemma
\ref{lemma-induced-family-supports}. These maps are isomorphisms in each
of the following cases:
\begin{enumerate}
\item $Z$ is closed;
\item $Z$ is open and $\Phi$ is paracompactifying.
\end{enumerate}
\end{lemma}

\begin{proof}
The functor $k_!$ is exact by Lemma
\ref{lemma-extension-by-zero-locally-closed}. Hence
$$
\mathcal{G}\longmapsto H^p_\Phi(X,k_!\mathcal{G})
$$
is a cohomological $\delta$-functor on $\textit{Ab}(Z)$. In degree zero
it agrees with $\Gamma_{\Phi_Z}(Z,-)$ by Lemma
\ref{lemma-induced-family-supports}. Universality in Lemma
\ref{lemma-cohomology-family-supports} therefore gives the displayed
maps.

Suppose first that $Z$ is closed. Embed $\mathcal{G}$ in an injective
abelian sheaf $\mathcal{I}$ on $Z$. The functor $k_*=k_!$ takes
injectives to injectives: it is right adjoint to the exact restriction
functor $k^{-1}$; see the discussion preceding Sheaves, Lemma
\ref{sheaves-lemma-stalks-closed-pushforward} and Modules, Lemma
\ref{modules-lemma-exactness-pushforward-pullback}. Thus the positive
terms of the target $\delta$-functor can be effaced.

Suppose next that $Z$ is open and $\Phi$ is paracompactifying. Then
$\Phi_Z$ is paracompactifying by Lemma
\ref{lemma-induced-family-supports}. Embed $\mathcal{G}$ in a
$\Phi_Z$-soft sheaf $\mathcal{S}$ using Lemma
\ref{lemma-flasque-soft-family-supports}. Exactness of $k_!$ gives an
embedding $k_!\mathcal{G}\to k_!\mathcal{S}$. Lemma
\ref{lemma-open-extension-phi-soft} says that $k_!\mathcal{S}$ is
$\Phi$-soft, so its positive $H^p_\Phi$ vanish by Lemma
\ref{lemma-family-supports-flasque-soft-acyclic}. Again the positive
terms of the target $\delta$-functor can be effaced.

In either case the target $\delta$-functor is universal by Homology,
Lemma \ref{homology-lemma-efface-implies-universal}. Applying
universality in both directions to the degree zero isomorphism and its
inverse shows that the displayed morphism is an isomorphism in every
degree.
\end{proof}

\begin{lemma}
\label{lemma-family-supports-closed-open-exact-sequence}
\begin{reference}
\cite[Chapter III, \S 3.5, p. 167]{Tohoku}
\end{reference}
Let $Y\subset X$ be closed, let $U=X\setminus Y$, and let $\Phi$ be a
family of supports on $X$. For every abelian sheaf $\mathcal{F}$ on $X$
there is a functorial long exact sequence
$$
\ldots\to H^p_\Phi(X,\mathcal{F}_U)\to
H^p_\Phi(X,\mathcal{F})\to
H^p_{\Phi_Y}(Y,\mathcal{F}|_Y)\to
H^{p+1}_\Phi(X,\mathcal{F}_U)\to\ldots.
$$
If $\Phi$ is paracompactifying, then the first and last terms may be
replaced by
$$
H^p_{\Phi_U}(U,\mathcal{F}|_U)
\quad\text{and}\quad
H^{p+1}_{\Phi_U}(U,\mathcal{F}|_U),
$$
respectively.
\end{lemma}

\begin{proof}
Apply the cohomological $\delta$-functor $H^p_\Phi(X,-)$ to the short
exact sequence of Lemma
\ref{lemma-extension-by-zero-locally-closed}. Identify the closed term,
and in the paracompactifying case the open terms, using Lemma
\ref{lemma-family-supports-extension-by-zero}.
\end{proof}


\begin{lemma}
\label{lemma-cohomological-functor-local-generators}
\begin{reference}
\cite[Chapter III, \S 3.6, Proposition 3.6.1, pp. 167--169]{Tohoku}
\end{reference}
Let $X$ be a topological space and let $\mathcal{A}$ be an abelian
category satisfying AB4. Let
$$
T^q:\textit{Ab}(X)\longrightarrow\mathcal{A},
\qquad q\in\mathbf{Z},
$$
be additive functors. Assume that every short exact sequence
$0\to\mathcal{F}'\to\mathcal{F}\to\mathcal{F}''\to0$ gives a
functorial long exact sequence
$$
\ldots\to T^q(\mathcal{F}')\to T^q(\mathcal{F})\to
T^q(\mathcal{F}'')\to T^{q+1}(\mathcal{F}')\to\ldots,
$$
and that every $T^q$ commutes with directed colimits. Fix
$p\in\mathbf{Z}$, and let
$\mathcal{C}\subset\mathcal{A}$ be a Serre subcategory stable under
arbitrary direct sums.

For an open subset $U\subset X$, write
$$
\mathbf{Z}_U=j_!\underline{\mathbf{Z}}_U,
$$
where $j:U\to X$ is the inclusion. If
$$
T^q(\mathbf{Z}_U)\in\mathcal{C}
$$
for every open $U\subset X$ and $q=p,p+1,p+2$, then
$$
T^p(\mathcal{F})\in\mathcal{C}
$$
for every abelian sheaf $\mathcal{F}$ on $X$.
\end{lemma}

\begin{proof}
We first record that $\mathcal{C}$ is stable under directed colimits.
Indeed, a directed colimit in the AB4 category $\mathcal{A}$ has the
standard presentation
$$
\bigoplus_{i\to j} M_i\longrightarrow
\bigoplus_i M_i\longrightarrow \colim_i M_i\longrightarrow 0.
$$
If all $M_i$ belong to $\mathcal{C}$, then both direct sums do, and the
cokernel does because $\mathcal{C}$ is a Serre subcategory.

Let $\mathcal{F}$ be an abelian sheaf on $X$. The set of all its local
sections generates $\mathcal{F}$ in the sense of Modules, Definition
\ref{modules-definition-generated-by-local-sections}. The subsheaves
generated by finite subsets of this set form a directed system with colimit
$\mathcal{F}$; this may be checked on stalks using Modules, Lemma
\ref{modules-lemma-generated-by-local-sections-stalk}. Since $T^p$ commutes
with directed colimits and $\mathcal{C}$ is stable under them, it is enough
to treat a sheaf generated by finitely many local sections.

We reduce further by induction on the number of generators. If
$\mathcal{F}'$ is generated by all but the last generator, then
$\mathcal{F}/\mathcal{F}'$ is generated by one local section. The long
exact sequence and the Serre property reduce the
assertion to $\mathcal{F}'$ and $\mathcal{F}/\mathcal{F}'$. Thus it is
enough to consider a sheaf generated by one section over an open subset
$U$. The adjunction for extension by zero gives an exact sequence
$$
0\longrightarrow\mathcal{R}\longrightarrow\mathbf{Z}_U
\longrightarrow\mathcal{F}\longrightarrow 0.
$$
The resulting exact sequence
$$
T^p(\mathbf{Z}_U)\longrightarrow T^p(\mathcal{F})
\longrightarrow T^{p+1}(\mathcal{R})
$$
shows that it remains to prove
$T^{p+1}(\mathcal{R})\in\mathcal{C}$.

The natural map $\mathbf{Z}_U\to\underline{\mathbf{Z}}_X$ is injective
on stalks, so we may view $\mathcal{R}$ as a subsheaf of the constant
sheaf $\underline{\mathbf{Z}}_X$. Every local section of $\mathcal{R}$
is locally a constant integer. Hence $\mathcal{R}$ is the directed colimit
of subsheaves generated by finitely many constant integer sections over
open subsets. Since $T^{p+1}$ commutes with directed colimits, we may assume
that $\mathcal{R}$ is generated by constant sections
$$
f_i=n_i\quad\text{over}\quad U_i,\qquad 1\leq i\leq k,
$$
where, after removing zero sections and changing signs, $n_i>0$.

For every nonempty subset $I\subset\{1,\ldots,k\}$, adjoin over
$\bigcap_{i\in I}U_i$ the constant section with value
$\gcd(n_i\mid i\in I)$. Bezout's identity shows that these new sections
belong to the subsheaf already generated, so this does not change
$\mathcal{R}$. After this finite enlargement, the generators have the
following property: for every $x\in X$, the greatest common divisor of the
$n_i$ with $x\in U_i$ is equal to one of these $n_i$. Reorder the generators
so that
$$
n_1\leq n_2\leq\ldots\leq n_k.
$$

Let $\mathcal{R}_m$ be the subsheaf generated by $f_1,\ldots,f_m$, put
$$
V_m=\bigcup_{i\leq m}U_i,
\qquad
Y_m=U_m\setminus V_{m-1},
$$
and set $\mathcal{R}_0=0$. We claim that
$$
\mathcal{R}_m/\mathcal{R}_{m-1}\cong\mathbf{Z}_{Y_m},
$$
where the right hand side is the locally closed extension by zero of the
constant sheaf on $Y_m$ from Lemma
\ref{lemma-extension-by-zero-locally-closed}.

It suffices to check the claim on stalks. The quotient vanishes away from
$V_m$. It also vanishes on $V_{m-1}$. Indeed, suppose
$x\in V_{m-1}\cap U_m$, and let $n_{i_0}$ be the greatest common divisor
of the $n_i$ for which $x\in U_i$. If $i_0<m$, then $n_{i_0}$ is an
earlier generator which divides $n_m$. If $i_0\geq m$, choose
$i<m$ with $x\in U_i$. Then
$$
n_{i_0}\leq n_i\leq n_m\leq n_{i_0},
$$
so $n_i=n_m$. In either case the germ of $f_m$ belongs to
$\mathcal{R}_{m-1,x}$. Over $Y_m$, no earlier generator is present and
the class of $f_m$ identifies the quotient with the constant integers.
This proves the claim.

The Serre property and induction on $m$ now reduce the assertion about
$T^{p+1}(\mathcal{R})$ to showing
$$
T^{p+1}(\mathbf{Z}_Y)\in\mathcal{C}
$$
for every locally closed subset $Y\subset X$. Write
$Y=V\setminus U$ with $U,V\subset X$ open, and put $W=U\cup V$. There is
a short exact sequence
$$
0\longrightarrow\mathbf{Z}_U\longrightarrow\mathbf{Z}_W
\longrightarrow\mathbf{Z}_Y\longrightarrow 0,
$$
as may be checked on stalks. Hence there is an exact sequence
$$
T^{p+1}(\mathbf{Z}_W)\longrightarrow
T^{p+1}(\mathbf{Z}_Y)\longrightarrow
T^{p+2}(\mathbf{Z}_U).
$$
The two outer terms belong to $\mathcal{C}$ by hypothesis, so the middle
term does as well. This completes the proof.
\end{proof}




\begin{lemma}
\label{lemma-compact-sections-colimit}
Let $A$ be a quasi-compact Hausdorff space and let
$(\mathcal{G}_i)_{i \in I}$ be a directed system of sheaves of abelian
groups on $A$. Then the canonical map
$$
\colim_i \Gamma(A, \mathcal{G}_i)
\longrightarrow
\Gamma(A, \colim_i \mathcal{G}_i)
$$
is an isomorphism.
\end{lemma}

\begin{proof}
Injectivity follows from Sheaves, Lemma
\ref{sheaves-lemma-directed-colimits-sections}, part (2). To prove
surjectivity, let $s$ be a section of the sheaf on the right. Choose a finite
open covering $A = \bigcup_{a = 1}^r U_a$, indices $i_a$, and sections
$s_a \in \mathcal{G}_{i_a}(U_a)$ which represent $s$ locally.

The space $A$ is paracompact, and hence Topology, Lemma
\ref{topology-lemma-paracompact-shrinking} gives, after taking a finite
subcover, open coverings
$$
A = \bigcup_{j = 1}^m W_j = \bigcup_{j = 1}^m V_j
$$
and a map $\alpha : \{1, \ldots, m\} \to \{1, \ldots, r\}$ such that
$$
\overline{W_j} \subset V_j \subset U_{\alpha(j)}.
$$
For every pair $j,j'$, the two corresponding sections have the same image
in the colimit near every point of the quasi-compact set
$\overline{W_j} \cap \overline{W_{j'}}$. Thus finitely many neighbourhoods
of this set and directedness give an index at which they agree on a
neighbourhood of the set. There are only finitely many pairs, so there is a
single $i \in I$, dominating all the $i_a$ and these agreement indices, at
which the resulting sections agree on every $W_j \cap W_{j'}$. They glue to
a section of $\mathcal{G}_i(A)$ which maps to $s$.
\end{proof}


\begin{lemma}
\label{lemma-colimit-soft-flasque}
\begin{reference}
\cite[Chapter III, \S 3.6, Lemma 3.6.4, p. 170]{Tohoku}
\end{reference}
Let $I$ be a directed set and let $(\mathcal{F}_i)_{i \in I}$ be a
directed system of abelian sheaves on a topological space $X$.
\begin{enumerate}
\item Suppose that $X$ is locally quasi-compact and Hausdorff, and let
$\Phi$ be the family of quasi-compact subsets of $X$. If every
$\mathcal{F}_i$ is $\Phi$-soft, then $\colim_i\mathcal{F}_i$ is
$\Phi$-soft.
\item Suppose that $X$ is Noetherian. If every $\mathcal{F}_i$ is flasque,
then $\colim_i\mathcal{F}_i$ is flasque.
\end{enumerate}
\end{lemma}

\begin{proof}
Assume the hypotheses of (1), set
$\mathcal{F} = \colim_i \mathcal{F}_i$, and let $A \in \Phi$. The subspace
$A$ is quasi-compact and Hausdorff. Restriction to $A$, being a pullback,
commutes with colimits. Lemma \ref{lemma-compact-sections-colimit} therefore
shows that every
$$
s \in \Gamma(A, \mathcal{F}|_A)
$$
comes from a section $s_i \in \Gamma(A, \mathcal{F}_i|_A)$ for some $i$.
As $\mathcal{F}_i$ is $\Phi$-soft, there is a section
$t_i \in \Gamma_\Phi(X, \mathcal{F}_i)$ restricting to $s_i$. The image
$t$ of $t_i$ in $\Gamma(X, \mathcal{F})$ has support contained in the
support of $t_i$, and hence has quasi-compact support. Thus
$t \in \Gamma_\Phi(X, \mathcal{F})$ and it restricts to $s$. This proves
(1).

Assume the hypotheses of (2), and again put
$\mathcal{F} = \colim_i \mathcal{F}_i$. Let $V \subset U \subset X$ be
open and let $s \in \mathcal{F}(V)$. Every subset of $X$ is Noetherian by
Topology, Lemma \ref{topology-lemma-Noetherian}, and hence every open subset
of $V$ is quasi-compact by Topology, Lemma
\ref{topology-lemma-Noetherian-quasi-compact}. Thus Sheaves, Lemma
\ref{sheaves-lemma-directed-colimits-sections}, part (4), shows that
$$
\colim_i \mathcal{F}_i(V) \longrightarrow \mathcal{F}(V)
$$
is an isomorphism. Choose $i$ and $s_i \in \mathcal{F}_i(V)$ mapping to
$s$. Since $\mathcal{F}_i$ is flasque, $s_i$ extends to a section over $U$.
Its image in $\mathcal{F}(U)$ extends $s$. Hence $\mathcal{F}$ is flasque.
\end{proof}


\begin{lemma}
\label{lemma-sections-colimit-paracompact-neighbourhoods}
Let $X$ be a topological space and let $Y \subset X$. Assume that for every
open neighbourhood $W$ of $Y$ there is a paracompact subset $P \subset W$
which is a neighbourhood of $Y$. Then:
\begin{enumerate}
\item for every sheaf of sets $\mathcal{F}$ on $X$ the canonical map
$$
\colim_{Y \subset U} \Gamma(U, \mathcal{F})
\longrightarrow \Gamma(Y, \mathcal{F}|_Y)
$$
is bijective, where the colimit is over the open neighbourhoods $U$ of $Y$,
ordered by reverse inclusion;
\item the space $Y$ is paracompact;
\item every closed subset $A \subset Y$ has the same property: every open
neighbourhood of $A$ in $X$ contains a paracompact neighbourhood of $A$.
\end{enumerate}
\end{lemma}

\begin{proof}
We first prove (1). Injectivity holds without the hypothesis on $Y$. Namely,
if sections over two open neighbourhoods have the same restriction to $Y$,
then their germs agree at every point of $Y$. They therefore agree over a
smaller open neighbourhood of $Y$.

Let $s \in \Gamma(Y, \mathcal{F}|_Y)$. There are open subsets $U_i \subset X$
covering $Y$ and sections $s_i \in \mathcal{F}(U_i)$ which induce $s$ over
$Y \cap U_i$. For every $i,i'$ there is an open
$E_{ii'} \subset U_i \cap U_{i'}$ containing
$Y \cap U_i \cap U_{i'}$ such that $s_i$ and $s_{i'}$ agree over
$E_{ii'}$. Choose a paracompact neighbourhood
$$
Y \subset P \subset \bigcup_i U_i.
$$
Apply Topology, Lemma \ref{topology-lemma-paracompact-shrinking} to the open
covering $\{P \cap U_i\}$ of $P$. We obtain locally finite open coverings
$P=\bigcup_{j\in J}V_j=\bigcup_{j\in J}W_j$ and a map $\alpha:J\to I$ such
that
$$
\overline{W_j}^{\,P}\subset V_j\subset P\cap U_{\alpha(j)}.
$$
The family of closed subsets
$$
C_{jj'}=
\left(\overline{W_j}^{\,P}\cap\overline{W_{j'}}^{\,P}\right)
\setminus\left(P\cap E_{\alpha(j)\alpha(j')}\right)
$$
is locally finite. Hence $C=\bigcup_{j,j'}C_{jj'}$ is closed in $P$.
Moreover, $C$ is disjoint from $Y$. Thus
$$
(P\setminus C)\cap\mathop{\rm Int}_X(P)
$$
is an open neighbourhood of $Y$ in $X$. On its intersections with the
$W_j$, the sections $s_{\alpha(j)}$ agree by construction, and hence they
glue. This proves surjectivity.

For (2), let $Y=\bigcup_i O_i$ be an open covering, with each $O_i$ open in
$Y$. Write $O_i=Y\cap U_i$ with $U_i\subset X$ open. Choose a paracompact
neighbourhood $P$ of $Y$ contained in $\bigcup_iU_i$. A locally finite open
refinement on $P$ of $\{P\cap U_i\}$ restricts to a locally finite open
refinement of the given covering on $Y$. Also, $Y$ is Hausdorff because it
is a subspace of $P$. Hence $Y$ is paracompact.

For (3), let $A\subset Y$ be closed and let $W\subset X$ be an open
neighbourhood of $A$. Write $A=Y\cap D$ with $D\subset X$ closed. The open
set $W\cup(X\setminus D)$ contains $Y$, so we can choose a paracompact
neighbourhood $P$ of $Y$ contained in it. Apply
Topology, Lemma \ref{topology-lemma-paracompact-shrinking} to the covering
$P=T_0\cup T_1$, where
$$
T_0=P\cap W,\qquad T_1=P\setminus D.
$$
Write the resulting data as
$$
P=\bigcup_{j\in J}V_j=\bigcup_{j\in J}W_j,\qquad
\overline{W_j}^{\,P}\subset V_j\subset T_{\alpha(j)}
$$
for a map $\alpha:J\to\{0,1\}$. Put
$$
J_0=\{j\in J\mid\alpha(j)=0\},\qquad
B=\bigcup_{j\in J_0}\overline{W_j}^{\,P}.
$$
Local finiteness shows that $B$ is closed in $P$, and hence paracompact,
and the displayed inclusions show that $B\subset W$. Every point of
$P\cap D$ belongs to some $W_j$ with $j\in J_0$. Thus
$\bigcup_{j\in J_0}W_j$ is an open subset of $P$ containing $A$.
Since $P$ is a neighbourhood of $Y$, this open subset of $P$ contains an
open neighbourhood of $A$ in $X$. Hence $B$ is a paracompact neighbourhood
of $A$ contained in $W$. This proves (3).
\end{proof}

\begin{lemma}
\label{lemma-cohomology-colimit-paracompact-neighbourhoods}
\begin{reference}
\cite[Chapter III, \S 3.10, Corollary 2, p. 182]{Tohoku}
\end{reference}
Let $X$ be a topological space and let $Y\subset X$. Assume that every open
neighbourhood of $Y$ contains a paracompact neighbourhood of $Y$. Then for
every abelian sheaf $\mathcal{F}$ on $X$ the canonical maps
$$
\colim_{Y\subset U}H^p(U,\mathcal{F}|_U)
\longrightarrow H^p(Y,\mathcal{F}|_Y)
$$
are isomorphisms for every $p\in\mathbf{Z}$. Here the colimit is over the
open neighbourhoods of $Y$, ordered by reverse inclusion. The isomorphisms
are functorial in $\mathcal{F}$ and compatible with the connecting maps.

If $Y$ is closed, it is enough that $Y$ have one paracompact neighbourhood.
\end{lemma}

\begin{proof}
Let $\mathcal{I}$ be an injective abelian sheaf on $X$. We claim that
$\mathcal{I}|_Y$ is soft. Let $A\subset Y$ be closed. By Lemma
\ref{lemma-sections-colimit-paracompact-neighbourhoods}, part (3), the
hypothesis of part (1) of that lemma holds for $A$. Thus a section of
$\mathcal{I}|_Y$ over $A$ is induced by a section of $\mathcal{I}$ over an
open neighbourhood of $A$ in $X$. Since $\mathcal{I}$ is flasque by Lemma
\ref{lemma-injective-flasque}, this section extends over $X$, proving the
claim.

The space $Y$ is paracompact by Lemma
\ref{lemma-sections-colimit-paracompact-neighbourhoods}, part (2). Hence the
family of all closed subsets of $Y$ is paracompactifying, and softness is
the same as softness for this family. Lemma
\ref{lemma-family-supports-flasque-soft-acyclic} therefore gives
$$
H^q(Y,\mathcal{I}|_Y)=0\quad\text{for }q>0.
$$

For every open neighbourhood $U$ of $Y$, consider the additive functor
$$
G_U(\mathcal{F})=\Gamma(U,\mathcal{F}|_U),
$$
and put $G(\mathcal{F})=\Gamma(Y,\mathcal{F}|_Y)$. Restriction to $U$
takes injectives to injectives by Lemma \ref{lemma-cohomology-of-open},
applied to the constant sheaves of rings. Consequently
$$
R^pG_U(\mathcal{F})=H^p(U,\mathcal{F}|_U).
$$
Restriction from $X$ to $Y$ is exact by Modules, Lemma
\ref{modules-lemma-exactness-pushforward-pullback}. The acyclicity proved
above therefore shows that an injective resolution on $X$, after restriction
to $Y$, computes cohomology on $Y$. Thus
$$
R^pG(\mathcal{F})=H^p(Y,\mathcal{F}|_Y).
$$
Part (1) of Lemma
\ref{lemma-sections-colimit-paracompact-neighbourhoods} says that
$G=\colim_U G_U$. The category of abelian sheaves on $X$ has enough
injectives by Injectives, Lemma
\ref{injectives-lemma-abelian-sheaves-space}, and abelian groups satisfy
AB5. Derived Categories, Lemma
\ref{derived-lemma-filtered-colimit-right-derived-functors} now gives the
asserted isomorphisms, functoriality, and compatibility with connecting
maps.

Finally, suppose that $Y$ is closed and has a paracompact neighbourhood
$P$. Given an open neighbourhood $W$ of $Y$, apply
Topology, Lemma \ref{topology-lemma-paracompact-shrinking} to
$$
P=(P\cap W)\cup(P\setminus Y).
$$
The construction in part (3), with $D=Y$, gives a paracompact
neighbourhood of $Y$ contained in $W$. Hence the hypothesis of the lemma
holds.
\end{proof}


\begin{lemma}
\label{lemma-cohomology-limit-relatively-compact-opens}
\begin{reference}
\cite[Chapter III, \S 3.10, (3.10.2) and Proposition 3.10.2,
pp. 182--183]{Tohoku}
\end{reference}
Let $X$ be a locally quasi-compact Hausdorff space and let
$\mathcal{F}$ be an abelian sheaf on $X$. Let $\mathcal{U}$ be the
directed set of open subsets $U\subset X$ such that $\overline{U}$ is
quasi-compact, ordered by inclusion. Restriction gives canonical maps
$$
\rho^p :
H^p(X,\mathcal{F})\longrightarrow
\lim_{U\in\mathcal{U}}H^p(U,\mathcal{F}|_U),
\qquad p\in\mathbf{Z}.
$$
These maps are isomorphisms for $p=0$, are functorial in $\mathcal{F}$,
and are compatible with the connecting maps.

Assume in addition that $X$ is countable at infinity, in other words that
$X$ is a countable union of quasi-compact subsets. Then:
\begin{enumerate}
\item The map $\rho^p$ is surjective for every $p\in\mathbf{Z}$.
\item Let $p>1$. The map $\rho^p$ is an isomorphism if, for every
$U\in\mathcal{U}$, there is a $V\in\mathcal{U}$ with $U\subset V$ such
that, for every $W\in\mathcal{U}$ with $V\subset W$, the images of
$$
H^{p-1}(W,\mathcal{F}|_W)\longrightarrow H^{p-1}(U,\mathcal{F}|_U)
$$
and
$$
H^{p-1}(V,\mathcal{F}|_V)\longrightarrow H^{p-1}(U,\mathcal{F}|_U)
$$
are equal.
\item The map $\rho^1$ is an isomorphism if the groups
$H^0(U,\mathcal{F}|_U)$, for $U\in\mathcal{U}$, are endowed with
complete metrizable topological group structures such that all
restriction maps are continuous, and if, for every $U\in\mathcal{U}$,
there is a $V\in\mathcal{U}$ with $U\subset V$ such that, for every
$W\in\mathcal{U}$ with $V\subset W$, the image of
$$
H^0(W,\mathcal{F}|_W)\longrightarrow H^0(U,\mathcal{F}|_U)
$$
is dense in the image of
$$
H^0(V,\mathcal{F}|_V)\longrightarrow H^0(U,\mathcal{F}|_U).
$$
\end{enumerate}
\end{lemma}

\begin{proof}
The set $\mathcal{U}$ is directed, since the closure of the union of two
of its members is the union of their quasi-compact closures. It also
covers $X$. Indeed, if $x\in X$, choose a quasi-compact neighbourhood
$E$ of $x$ and an open $U$ with $x\in U\subset E$. The set $E$ is
closed because $X$ is Hausdorff, and hence
$\overline{U}\subset E$ is quasi-compact. Compatibility of restriction
maps gives the maps $\rho^p$ and their compatibility with connecting
maps. A compatible family of sections over the members of
$\mathcal{U}$ glues uniquely over $X$, proving the assertion for $p=0$.

Suppose that $X$ is countable at infinity. We first choose a cofinal
sequence
$$
U_1\subset U_2\subset\ldots
$$
in $\mathcal{U}$. To see that this is possible, write
$X=\bigcup_{n\geq 1}K_n$ with $K_n$ quasi-compact, and replace $K_n$
by $K_1\cup\ldots\cup K_n$ so that the sequence is increasing. Every
quasi-compact subset $K\subset X$ is contained in a member of
$\mathcal{U}$: cover $K$ by finitely many open subsets having
quasi-compact closure and take their union. Starting with
$U_0=\emptyset$, we may therefore choose $U_n\in\mathcal{U}$ such that
$$
K_n\cup\overline{U_{n-1}}\subset U_n.
$$
Then $\bigcup U_n=X$. If $U\in\mathcal{U}$, quasi-compactness of
$\overline{U}$ and the fact that the $U_n$ form an increasing covering
show that $\overline{U}\subset U_n$ for some $n$. Thus the sequence is
cofinal. In particular, limits over $\mathcal{U}$ may be computed over
this sequence; this also follows from Categories, Lemma
\ref{categories-lemma-initial}, applied to the opposite index
categories.

Choose an injective resolution
$\mathcal{F}\to\mathcal{I}^\bullet$. Such resolutions exist by
Injectives, Lemma
\ref{injectives-lemma-abelian-sheaves-space}. Put
$$
C_n^\bullet=\Gamma(U_n,\mathcal{I}^\bullet|_{U_n}).
$$
Restriction to $U_n$ preserves injectives by Lemma
\ref{lemma-cohomology-of-open}, applied to the constant sheaves of
rings. Hence
$$
H^p(C_n^\bullet)=H^p(U_n,\mathcal{F}|_{U_n}).
$$
Every $\mathcal{I}^q$ is flasque by Lemma
\ref{lemma-injective-flasque}. Consequently the transition maps
$$
C_{n+1}^q\longrightarrow C_n^q
$$
are surjective. Thus every inverse system $(C_n^q)$ is Mittag-Leffler
and is right acyclic for $\lim$ by More on Algebra, Lemma
\ref{more-algebra-lemma-compute-Rlim}. Since the $U_n$ form an
increasing covering of $X$, the sheaf axiom gives
$$
\Gamma(X,\mathcal{I}^\bullet)=\lim_n C_n^\bullet.
$$
More on Algebra, Lemmas
\ref{more-algebra-lemma-compute-Rlim} and
\ref{more-algebra-lemma-break-long-exact-sequence} now give canonical
short exact sequences
\begin{equation}
\label{equation-cohomology-limit-relatively-compact-opens}
0\longrightarrow
R^1\lim_n H^{p-1}(U_n,\mathcal{F}|_{U_n})
\longrightarrow H^p(X,\mathcal{F})
\longrightarrow
\lim_n H^p(U_n,\mathcal{F}|_{U_n})
\longrightarrow 0.
\end{equation}
The right-hand map is $\rho^p$ under the cofinality identification.
This proves (1).

Assume the hypothesis of (2), and set
$$
A_n=H^{p-1}(U_n,\mathcal{F}|_{U_n}).
$$
Given $n$, choose $V$ as in (2), and then choose $m\geq n$ such that
$V\subset U_m$. For $k\geq m$, applying the hypothesis first to
$W=U_m$ and then to $W=U_k$ shows that
$$
\mathop{\rm Im}(A_k\longrightarrow A_n)
=
\mathop{\rm Im}(A_m\longrightarrow A_n).
$$
Thus $(A_n)$ is Mittag-Leffler in the sense of Homology, Definition
\ref{homology-definition-Mittag-Leffler}. More on Algebra, Lemma
\ref{more-algebra-lemma-compute-Rlim} gives
$R^1\lim A_n=0$, and (2) follows from
(\ref{equation-cohomology-limit-relatively-compact-opens}).

We prove the topological assertion needed for (3). Let
$$
A_1\longleftarrow A_2\longleftarrow A_3\longleftarrow\ldots
$$
be an inverse system of complete metrizable abelian topological groups
with continuous transition maps. Suppose that for every $n$ there is
an $m\geq n$ such that, for all $k\geq m$,
$$
\mathop{\rm Im}(A_k\longrightarrow A_n)
\quad\text{is dense in}\quad
\mathop{\rm Im}(A_m\longrightarrow A_n).
$$
We claim that $R^1\lim A_n=0$.

Choose a cofinal subsequence
$n_1<n_2<\ldots$ such that, for every $r$ and every
$k\geq n_{r+1}$, the image of $A_k$ in $A_{n_r}$ is dense in the
image of $A_{n_{r+1}}$. Passing to a cofinal subsequence gives an
isomorphic pro-object and does not change its derived limit by More on
Algebra, Lemma
\ref{more-algebra-lemma-pro-isomorphism-iso-Rlim}. After renumbering, we
may therefore assume that for every $r$ the image of $A_{r+2}$ in $A_r$
is dense in the image of $A_{r+1}$ in $A_r$.

Let $f_r:A_{r+1}\to A_r$ be the transition map and put
$$
B_r=\overline{f_r(A_{r+1})}\subset A_r.
$$
Each $B_r$ is a closed subgroup of a complete metrizable group, and
hence is complete and metrizable for the induced topology.
Continuity gives maps
$$
g_r:B_{r+1}\longrightarrow B_r.
$$
Their images are dense: they contain the image of $A_{r+2}$ in
$A_r$, which is dense in the image of $A_{r+1}$ and hence in $B_r$.
Moreover, the inclusions $B_r\to A_r$ define an isomorphism of
pro-objects. Indeed, the transition map
$A_{r+1}\to A_r$ factors through $B_r$, and these factorizations give
an inverse in the pro-category. Another application of More on
Algebra, Lemma \ref{more-algebra-lemma-pro-isomorphism-iso-Rlim} shows
that it is enough to prove $R^1\lim B_r=0$.

For every $r$, choose a complete metric $d_r$ inducing the topology
on $B_r$; no invariance property of these metrics will be used.
By More on Algebra, Lemma \ref{more-algebra-lemma-compute-Rlim}, it
suffices to prove that
$$
\prod_r B_r\longrightarrow\prod_r B_r,\qquad
(x_r)\longmapsto(x_r-g_r(x_{r+1}))
$$
is surjective. Fix $(a_r)\in\prod_rB_r$. We construct, for every
$N\geq1$, elements
$$
x_1^{(N)},\ldots,x_N^{(N)}
$$
such that
$$
x_r^{(N)}=a_r+g_r(x_{r+1}^{(N)})\quad(1\leq r<N)
$$
and, for $N\geq1$,
$$
d_r(x_r^{(N+1)},x_r^{(N)})<2^{-N}
\quad(1\leq r\leq N).
$$
For $N=1$, choose $x_1^{(1)}$ arbitrarily. Suppose the construction
has been made through $N$. For $z\in B_N$, define continuously and
recursively
$$
y_N(z)=z,\qquad
y_r(z)=a_r+g_r(y_{r+1}(z))\quad(r=N-1,\ldots,1).
$$
The induction hypothesis gives
$y_r(x_N^{(N)})=x_r^{(N)}$. Hence there is an open neighbourhood
$Q\subset B_N$ of $x_N^{(N)}$ such that
$$
d_r(y_r(z),x_r^{(N)})<2^{-N}
$$
for all $z\in Q$ and all $1\leq r\leq N$. Since $g_N(B_{N+1})$ is
dense in $B_N$, the translate $a_N+g_N(B_{N+1})$ is dense as well.
Choose $x_{N+1}^{(N+1)}\in B_{N+1}$ such that
$$
z=a_N+g_N(x_{N+1}^{(N+1)})\in Q,
$$
and set
$$
x_N^{(N+1)}=z,\qquad
x_r^{(N+1)}=y_r(z)\quad(r<N).
$$
This completes the induction.

For fixed $r$, the sequence $(x_r^{(N)})_{N\geq r}$ is Cauchy because
the successive distances are bounded by the summable series
$\sum_{N\geq r}2^{-N}$. Completeness gives a limit $x_r\in B_r$.
Passing to the limit in the displayed recursive identities, using
continuity of $g_r$, gives
$$
x_r=a_r+g_r(x_{r+1})
$$
for every $r$. Thus the displayed product map is surjective and
$R^1\lim A_n=0$, proving the claim.

Finally, set
$$
A_n=H^0(U_n,\mathcal{F}|_{U_n})
$$
with the topologies from (3). Given $n$, choose $V$ as in (3), and
choose $m\geq n$ with $V\subset U_m$. For $k\geq m$, the image of
$A_k$ in $A_n$ and the image of $A_m$ in $A_n$ are both dense in the
image of $H^0(V,\mathcal{F}|_V)$ in $A_n$. Since
$$
\mathop{\rm Im}(A_k\to A_n)
\subset
\mathop{\rm Im}(A_m\to A_n),
$$
the former is dense in the latter. Thus the density hypothesis of the
claim holds and gives $R^1\lim A_n=0$. Applying
(\ref{equation-cohomology-limit-relatively-compact-opens}) with $p=1$
proves (3).
\end{proof}


\begin{definition}
\label{definition-relative-family-supports}
\begin{reference}
\cite[Chapter III, \S 3.7, pp. 171--172]{Tohoku}
\end{reference}
Let $f : Y \to X$ be a continuous map and let $\Psi$ be a family of
supports on $Y$. For an open subset $U \subset X$, let $\Psi(U)$ be the
collection of closed subsets $A \subset f^{-1}(U)$ with the following
property: for every $x \in U$ there is an open neighbourhood $V \subset U$
of $x$ such that
$$
\overline{A \cap f^{-1}(V)} \in \Psi,
$$
where the closure is taken in $Y$.
\end{definition}

\begin{lemma}
\label{lemma-relative-direct-image-family-supports}
\begin{reference}
\cite[Chapter III, \S 3.7, pp. 171--172]{Tohoku}
\end{reference}
Let $f : Y \to X$ and $\Psi$ be as in Definition
\ref{definition-relative-family-supports}.
\begin{enumerate}
\item For every open $U \subset X$, the collection $\Psi(U)$ is a family
of supports on $f^{-1}(U)$.
\item For every abelian sheaf $\mathcal{F}$ on $Y$, the rule
$$
U \longmapsto
\Gamma_{\Psi(U)}(f^{-1}(U), \mathcal{F}|_{f^{-1}(U)})
$$
is a sheaf on $X$, denoted $f_\Psi\mathcal{F}$. This construction defines
an additive left exact functor
$$
f_\Psi : \textit{Ab}(Y) \longrightarrow \textit{Ab}(X).
$$
\item If $s \in \Gamma(X, f_\Psi\mathcal{F})$ corresponds to a section of
$\mathcal{F}$ on $Y$, then
$$
\text{Supp}_X(s) = \overline{f(\text{Supp}_Y(s))}.
$$
\item Let $\Phi$ be a family of supports on $X$. Then
$$
\Psi' = \{A \in \Psi(X) \mid \overline{f(A)} \in \Phi\}
$$
is a family of supports on $Y$, and there is a functorial equality
$$
\Gamma_\Phi(X, f_\Psi\mathcal{F})
= \Gamma_{\Psi'}(Y, \mathcal{F}).
$$
We say that $\Phi$ and $\Psi$ are adapted relative to $f$ if
$\Psi'=\Psi$.
\end{enumerate}
\end{lemma}

\begin{proof}
The empty set belongs to $\Psi(U)$. If $A \in \Psi(U)$ and $A' \subset A$
is closed in $f^{-1}(U)$, then $A' \in \Psi(U)$ by downward closure of
$\Psi$. If $A_1,A_2 \in \Psi(U)$, then at a point $x \in U$ intersect
neighbourhoods which work for $A_1$ and $A_2$. The closure belonging to
$A_1 \cup A_2$ is contained in the union of the two closures, which belongs
to $\Psi$. This proves (1).

The rule in (2) is a subpresheaf of $f_*\mathcal{F}$. Suppose that
$U=\bigcup U_i$ and compatible sections over the $U_i$ belong to the
displayed groups. They glue uniquely as sections of $f_*\mathcal{F}$. The
condition on the support is local on $U$: near $x \in U$, choose an $i$ with
$x \in U_i$ and use the condition for the section over $U_i$. Thus the
glued section belongs to the displayed group over $U$, proving the sheaf
condition. Additivity follows because the support of a sum is contained in
the union of the two supports. Restriction to the open subspace $f^{-1}(U)$
is exact by Modules, Lemma
\ref{modules-lemma-exactness-pushforward-pullback}. Left exactness now
follows from left exactness of
$\Gamma_{\Psi(U)}(f^{-1}(U),-)$ for every $U$, proved as in Lemma
\ref{lemma-cohomology-family-supports}.

Let $A=\text{Supp}_Y(s)$. A point $x \in X$ is outside
$\overline{f(A)}$ if and only if there is an open neighbourhood $V$ of $x$
such that $f^{-1}(V)$ is disjoint from $A$. This is equivalent to the
restriction of $s$ to $f^{-1}(V)$ being zero, and hence to the germ of $s$
in $f_\Psi\mathcal{F}$ at $x$ being zero. This proves (3).

The family axioms for $\Psi'$ follow from those for $\Psi(X)$ and $\Phi$,
using
$$
\overline{f(A_1 \cup A_2)}
= \overline{f(A_1)} \cup \overline{f(A_2)}.
$$
Finally, (3) says exactly that a global section on either side of the
equality in (4) belongs to that side if and only if it belongs to the
other side. This proves (4).
\end{proof}

\begin{lemma}
\label{lemma-relative-direct-image-acyclic}
\begin{reference}
\cite[Chapter III, \S 3.7, Lemma 3.7.1, pp. 172--173]{Tohoku}
\end{reference}
Let $f : Y \to X$, $\Psi$, $\Phi$, and $\Psi'$ be as in Lemma
\ref{lemma-relative-direct-image-family-supports}.
\begin{enumerate}
\item If $\Psi$ is the family of all closed subsets of $Y$, then
$f_\Psi=f_*$ and $f_\Psi$ sends injective abelian sheaves to injective
abelian sheaves.
\item If $\Phi$ is paracompactifying, then $f_\Psi$ sends flasque abelian
sheaves to $\Phi$-soft abelian sheaves.
\end{enumerate}
\end{lemma}

\begin{proof}
If $\Psi$ is the family of all closed subsets, then $\Psi(U)$ is the family
of all closed subsets of $f^{-1}(U)$. Hence $f_\Psi=f_*$. The final
assertion of (1) follows because $f_*$ is right adjoint to the exact functor
$f^{-1}$ by Modules, Lemma
\ref{modules-lemma-exactness-pushforward-pullback}, and hence preserves
injectives by Homology, Lemma
\ref{homology-lemma-adjoint-preserve-injectives}.

Assume $\Phi$ is paracompactifying and let $\mathcal{F}$ be flasque. Put
$\mathcal{G}=f_\Psi\mathcal{F}$. Let $B \in \Phi$ and let
$$
s \in \Gamma(B,\mathcal{G}|_B).
$$
Choose a closed paracompact neighbourhood $C \in \Phi$ of $B$. Lemma
\ref{lemma-section-closed-paracompact-neighbourhood} extends $s$ to a
section $t$ of $\mathcal{G}$ over an open neighbourhood $U$ of $B$.
After shrinking $U$ inside $\mathop{\rm Int}(C)$, apply Topology, Lemma
\ref{topology-lemma-paracompact-shrinking} on $C$ to the covering by
$U\cap C$ and $C\setminus B$. Let $W$ be the union of the members of the
resulting shrinking which are subordinate to $U\cap C$. The other members
are disjoint from $B$, so $B\subset W$, and local finiteness gives
$\overline{W}\subset U\cap C$, with closure computed in $C$. Replacing $W$
by $W\cap\mathop{\rm Int}(C)$ and taking its closure in $X$, we obtain a
closed neighbourhood $B_1$ of $B$ such that
$$
B_1 \subset U \cap C.
$$
In particular $B_1 \in \Phi$.

Regard $t$ as a section of $\mathcal{F}$ over $f^{-1}(U)$ and put
$A=\text{Supp}(t)$. The set
$$
D=A\cap f^{-1}(B_1)
$$
is closed in $Y$. On the open subset
$$
f^{-1}(\mathop{\rm Int}(B_1)) \cup (Y\setminus D)
$$
there is a section which equals $t$ on the first open and zero on the
second. Indeed, these sections agree on the intersection because $D$
contains the support of $t$ over $f^{-1}(\mathop{\rm Int}(B_1))$.
Flasqueness extends this section to an $h\in\Gamma(Y,\mathcal{F})$.

Its support $E$ is contained in $D$. We claim that $E\in\Psi(X)$. For a
point of $B_1$, use the defining local condition for $A\in\Psi(U)$ and
downward closure of $\Psi$; outside $B_1$, use a neighbourhood disjoint
from $B_1$. Moreover,
$$
\overline{f(E)}\subset B_1\in\Phi.
$$
Thus $E\in\Psi'$. By Lemma
\ref{lemma-relative-direct-image-family-supports}, the section $h$ defines
a section of $\mathcal{G}$ with support in $\Phi$. Over $B$ it induces
$s$. Hence $\mathcal{G}$ is $\Phi$-soft.
\end{proof}

\begin{lemma}[Relative Leray spectral sequence with supports]
\label{lemma-relative-Leray-family-supports}
\begin{reference}
\cite[Chapter III, \S 3.7, Theorem 3.7.3, p. 173]{Tohoku}
\end{reference}
Let $f : Y \to X$ be a continuous map, let $\Psi$ and $\Phi$ be families
of supports on $Y$ and $X$, and let $\Psi'$ be as in Lemma
\ref{lemma-relative-direct-image-family-supports}. Assume either that
$\Psi$ is the family of all closed subsets of $Y$ or that $\Phi$ is
paracompactifying. For every abelian sheaf $\mathcal{F}$ on $Y$ there is a
functorial first quadrant spectral sequence
$$
E_2^{p,q}=H^p_\Phi(X,R^qf_\Psi\mathcal{F})
\Rightarrow H^{p+q}_{\Psi'}(Y,\mathcal{F}).
$$
Moreover, for $q\geq0$, $R^qf_\Psi\mathcal{F}$ is the sheaf associated to
the presheaf
$$
U \longmapsto
H^q_{\Psi(U)}(f^{-1}(U),\mathcal{F}|_{f^{-1}(U)}).
$$
\end{lemma}

\begin{proof}
Lemma \ref{lemma-relative-direct-image-family-supports} gives an equality
of left exact functors
$$
\Gamma_{\Psi'}(Y,-)=\Gamma_\Phi(X,-)\circ f_\Psi.
$$
The hypotheses of the Grothendieck spectral sequence are satisfied. If
$\Psi$ is the family of all closed subsets, then $f_\Psi$ sends injectives
to injectives by Lemma \ref{lemma-relative-direct-image-acyclic}, and these
are right acyclic for $\Gamma_\Phi(X,-)$ by Lemma
\ref{lemma-cohomology-family-supports}. If $\Phi$ is paracompactifying, an
injective abelian sheaf on $Y$ is flasque by Lemma
\ref{lemma-injective-flasque}, its relative direct image is $\Phi$-soft by
Lemma \ref{lemma-relative-direct-image-acyclic}, and $\Phi$-soft sheaves are
right acyclic for $\Gamma_\Phi(X,-)$ by Lemma
\ref{lemma-family-supports-flasque-soft-acyclic}. Derived Categories, Lemma
\ref{derived-lemma-grothendieck-spectral-sequence} now gives the displayed
spectral sequence.

Apply Lemma \ref{lemma-derived-functor-sheafification} to the functor
$f_\Psi$. For every open $U\subset X$, Lemma
\ref{lemma-relative-direct-image-family-supports} identifies
$$
\Gamma(U,f_\Psi(-))
=\Gamma_{\Psi(U)}(f^{-1}(U),(-)|_{f^{-1}(U)}).
$$
Restriction to the open subspace $f^{-1}(U)$ preserves injective abelian
sheaves by Lemma \ref{lemma-cohomology-of-open}, applied to the constant
sheaves of rings. Hence the right derived functors of the latter are the
supported cohomology groups in the final display by Lemma
\ref{lemma-cohomology-family-supports}. This proves the presheaf
description.
\end{proof}


\section{Cohomology on spectral spaces}
\label{section-spectral}

\noindent
A key result on the cohomology of spectral spaces is Lemma \ref{lemma-colimit}
which loosely speaking says that cohomology commutes with cofiltered limits
in the category of spectral spaces as defined in
Topology, Definition \ref{topology-definition-spectral-space}.
This can be applied to give analogues of
Lemmas \ref{lemma-cohomology-of-closed} and \ref{lemma-proper-base-change}
as follows.

\begin{lemma}
\label{lemma-cohomology-of-neighbourhoods-of-closed}
Let $X$ be a spectral space. Let $\mathcal{F}$ be an abelian sheaf on $X$.
Let $E \subset X$ be a quasi-compact subset. Let $W \subset X$ be the set of
points of $X$ which specialize to a point of $E$.
\begin{enumerate}
\item $H^p(W, \mathcal{F}|_W) = \colim H^p(U, \mathcal{F})$
where the colimit is over quasi-compact open neighbourhoods of $E$,
\item $H^p(W \setminus E, \mathcal{F}|_{W \setminus E}) =
\colim H^p(U \setminus E, \mathcal{F}|_{U \setminus E})$
if $E$ is a constructible subset.
\end{enumerate}
\end{lemma}

\begin{proof}
From Topology, Lemma \ref{topology-lemma-make-spectral-space}
we see that $W = \lim U$ where the limit is over the quasi-compact
opens containing $E$. Each $U$ is a spectral space by
Topology, Lemma \ref{topology-lemma-spectral-sub}.
Thus we may apply Lemma \ref{lemma-colimit} to conclude that (1) holds.
The same proof works for part (2) except we use
Topology, Lemma \ref{topology-lemma-make-spectral-space-minus}.
\end{proof}

\begin{lemma}
\label{lemma-proper-base-change-spectral}
Let $f : X \to Y$ be a spectral map of spectral spaces. Let $y \in Y$.
Let $E \subset Y$ be the set of points specializing to $y$.
Let $\mathcal{F}$ be an abelian sheaf on $X$.
Then $(R^pf_*\mathcal{F})_y = H^p(f^{-1}(E), \mathcal{F}|_{f^{-1}(E)})$.
\end{lemma}

\begin{proof}
Observe that $E = \bigcap V$ where $V$ runs over the quasi-compact
open neighbourhoods of $y$ in $Y$. Hence $f^{-1}(E) = \bigcap f^{-1}(V)$.
This implies that $f^{-1}(E) = \lim f^{-1}(V)$ as topological spaces.
Since $f$ is spectral, each $f^{-1}(V)$ is a spectral space too
(Topology, Lemma \ref{topology-lemma-spectral-sub}).
We conclude that $f^{-1}(E)$ is a spectral space and that
$$
H^p(f^{-1}(E), \mathcal{F}|_{f^{-1}(E)}) =
\colim H^p(f^{-1}(V), \mathcal{F})
$$
by Lemma \ref{lemma-colimit}. On the other hand, the stalk of
$R^pf_*\mathcal{F}$ at $y$ is given by the colimit on the right.
\end{proof}

\begin{lemma}
\label{lemma-vanishing-for-profinite}
Let $X$ be a profinite topological space. Then $H^q(X, \mathcal{F}) = 0$
for all $q > 0$ and all abelian sheaves $\mathcal{F}$.
\end{lemma}

\begin{proof}
Any open covering of $X$ can be refined by a finite disjoint union
decomposition with open parts, see
Topology, Lemma \ref{topology-lemma-profinite-refine-open-covering}.
Hence if $\mathcal{F} \to \mathcal{G}$ is a surjection of abelian
sheaves on $X$, then $\mathcal{F}(X) \to \mathcal{G}(X)$ is surjective.
In other words, the global sections functor is an exact functor.
Therefore its higher derived functors are zero, see
Derived Categories, Lemma \ref{derived-lemma-right-derived-exact-functor}.
\end{proof}

\noindent
The following result on cohomological vanishing
improves Grothendieck's result
(Proposition \ref{proposition-vanishing-Noetherian})
and can be found in \cite{Scheiderer}.

\begin{proposition}
\label{proposition-cohomological-dimension-spectral}
\begin{reference}
Part (1) is the main theorem of \cite{Scheiderer}.
\end{reference}
Let $X$ be a spectral space of Krull dimension $d$.
Let $\mathcal{F}$ be an abelian sheaf on $X$.
\begin{enumerate}
\item $H^q(X, \mathcal{F}) = 0$ for $q > d$,
\item $H^d(X, \mathcal{F}) \to H^d(U, \mathcal{F})$ is surjective
for every quasi-compact open $U \subset X$,
\item $H^q_Z(X, \mathcal{F}) = 0$ for $q > d$ and any constructible
closed subset $Z \subset X$.
\end{enumerate}
\end{proposition}

\begin{proof}
We prove this result by induction on $d$.

\medskip\noindent
If $d = 0$, then $X$ is a profinite space, see
Topology, Lemma \ref{topology-lemma-characterize-profinite-spectral}.
Thus (1) holds by Lemma \ref{lemma-vanishing-for-profinite}.
If $U \subset X$ is quasi-compact open, then $U$ is
also closed as a quasi-compact subset of a Hausdorff space.
Hence $X = U \amalg (X \setminus U)$ as a topological space
and we see that (2) holds. Given $Z$ as in (3) we consider the
long exact sequence
$$
H^{q - 1}(X, \mathcal{F}) \to
H^{q - 1}(X \setminus Z, \mathcal{F}) \to
H^q_Z(X, \mathcal{F}) \to H^q(X, \mathcal{F})
$$
Since $X$ and $U = X \setminus Z$ are profinite (namely $U$ is quasi-compact
because $Z$ is constructible) and since
we have (2) and (1) we obtain the desired vanishing of the
cohomology groups with support in $Z$.

\medskip\noindent
Induction step. Assume $d \geq 1$ and assume
the proposition is valid for all spectral
spaces of dimension $< d$. We first prove part (2) for $X$.
Let $U$ be a quasi-compact open. Let $\xi \in H^d(U, \mathcal{F})$.
Set $Z = X \setminus U$. Let $W \subset X$ be the set of points
specializing to $Z$. By
Lemma \ref{lemma-cohomology-of-neighbourhoods-of-closed} we have
$$
H^d(W \setminus Z, \mathcal{F}|_{W \setminus Z}) =
\colim_{Z \subset V} H^d(V \setminus Z, \mathcal{F})
$$
where the colimit is over the quasi-compact open neighbourhoods $V$
of $Z$ in $X$.
By Topology, Lemma \ref{topology-lemma-make-spectral-space} we see that
$W \setminus Z$ is a spectral space.
Since every point of $W$ specializes to a point of $Z$, we see that
$W \setminus Z$ is a spectral space of Krull dimension $< d$.
By induction hypothesis we see that the image of $\xi$ in
$H^d(W \setminus Z, \mathcal{F}|_{W \setminus Z})$ is zero.
By the displayed formula, there exists a $Z \subset V \subset X$
quasi-compact open such that $\xi|_{V \setminus Z} = 0$.
Since $V \setminus Z = V \cap U$ we conclude by the Mayer-Vietoris
(Lemma \ref{lemma-mayer-vietoris}) for the covering $X = U \cup V$
that there exists a $\tilde \xi \in H^d(X, \mathcal{F})$ which restricts
to $\xi$ on $U$ and to zero on $V$. In other words, part (2) is true.

\medskip\noindent
Proof of part (1) assuming (2). Choose an injective resolution
$\mathcal{F} \to \mathcal{I}^\bullet$. Set
$$
\mathcal{G} = \Im(\mathcal{I}^{d - 1} \to \mathcal{I}^d) =
\Ker(\mathcal{I}^d \to \mathcal{I}^{d + 1})
$$
For $U \subset X$ quasi-compact open we have a map of exact sequences
as follows
$$
\xymatrix{
\mathcal{I}^{d - 1}(X) \ar[r] \ar[d] &
\mathcal{G}(X) \ar[r] \ar[d] &
H^d(X, \mathcal{F}) \ar[d] \ar[r] & 0 \\
\mathcal{I}^{d - 1}(U) \ar[r] &
\mathcal{G}(U) \ar[r] &
H^d(U, \mathcal{F}) \ar[r] & 0
}
$$
The sheaf $\mathcal{I}^{d - 1}$ is flasque by
Lemma \ref{lemma-injective-flasque} and the fact that $d \geq 1$.
By part (2) we see that the right vertical arrow is surjective.
We conclude by a diagram chase that the map
$\mathcal{G}(X) \to \mathcal{G}(U)$ is surjective.
By Lemma \ref{lemma-vanishing-ravi} we conclude that
$\check{H}^q(\mathcal{U}, \mathcal{G}) = 0$ for $q > 0$ and
any finite covering $\mathcal{U} : U = U_1 \cup \ldots \cup U_n$
of a quasi-compact open by quasi-compact opens. Applying
Lemma \ref{lemma-cech-vanish-basis} we find that $H^q(U, \mathcal{G}) = 0$
for all $q > 0$ and all quasi-compact opens $U$ of $X$.
By Leray's acyclicity lemma
(Derived Categories, Lemma \ref{derived-lemma-leray-acyclicity})
we conclude that
$$
H^q(X, \mathcal{F}) =
H^q\left(
\Gamma(X, \mathcal{I}^0) \to \ldots \to
\Gamma(X, \mathcal{I}^{d - 1}) \to \Gamma(X, \mathcal{G})
\right)
$$
In particular the cohomology group vanishes if $q > d$.

\medskip\noindent
Proof of (3).  Given $Z$ as in (3) we consider the long exact sequence
$$
H^{q - 1}(X, \mathcal{F}) \to
H^{q - 1}(X \setminus Z, \mathcal{F}) \to
H^q_Z(X, \mathcal{F}) \to H^q(X, \mathcal{F})
$$
Since $X$ and $U = X \setminus Z$ are spectral spaces
(Topology, Lemma \ref{topology-lemma-spectral-sub})
of dimension $\leq d$
and since we have (2) and (1) we obtain the desired vanishing.
\end{proof}













\section{The alternating {\v C}ech complex}
\label{section-alternating-cech}

\noindent
This section compares the {\v C}ech complex with the alternating {\v C}ech
complex and some related complexes.

\medskip\noindent
Let $X$ be a topological space. Let $\mathcal{U} : U = \bigcup_{i \in I} U_i$
be an open covering. For $p \geq 0$ set
$$
\check{\mathcal{C}}_{alt}^p(\mathcal{U}, \mathcal{F})
=
\left\{
\begin{matrix}
s \in  \check{\mathcal{C}}^p(\mathcal{U}, \mathcal{F})
\text{ such that }
s_{i_0 \ldots i_p} = 0 \text{ if } i_n = i_m \text{ for some } n \not = m\\
\text{ and }
s_{i_0\ldots i_n \ldots i_m \ldots i_p}
=
-s_{i_0\ldots i_m \ldots i_n \ldots i_p}
\text{ in any case.}
\end{matrix}
\right\}
$$
We omit the verification that the differential $d$ of
Equation (\ref{equation-d-cech}) maps
$\check{\mathcal{C}}^p_{alt}(\mathcal{U}, \mathcal{F})$ into
$\check{\mathcal{C}}^{p + 1}_{alt}(\mathcal{U}, \mathcal{F})$.

\begin{definition}
\label{definition-alternating-cech-complex}
\begin{reference}
\cite[Chapter I, \S 3, no. 18, p. 212]{FAC}
\end{reference}
Let $X$ be a topological space. Let $\mathcal{U} : U = \bigcup_{i \in I} U_i$
be an open covering. Let $\mathcal{F}$ be an abelian presheaf on $X$.
The complex $\check{\mathcal{C}}_{alt}^\bullet(\mathcal{U}, \mathcal{F})$
is the {\it alternating {\v C}ech complex} associated to $\mathcal{F}$ and the
open covering $\mathcal{U}$.
\end{definition}

\noindent
Hence there is a canonical morphism of complexes
$$
\check{\mathcal{C}}_{alt}^\bullet(\mathcal{U}, \mathcal{F})
\longrightarrow
\check{\mathcal{C}}^\bullet(\mathcal{U}, \mathcal{F})
$$
namely the inclusion of the alternating {\v C}ech complex into the
usual {\v C}ech complex.

\medskip\noindent
Suppose our covering $\mathcal{U} : U = \bigcup_{i \in I} U_i$ comes
equipped with a total ordering $<$ on $I$. In this case, set
$$
\check{\mathcal{C}}_{ord}^p(\mathcal{U}, \mathcal{F})
=
\prod\nolimits_{(i_0, \ldots, i_p) \in I^{p + 1}, i_0 < \ldots < i_p}
\mathcal{F}(U_{i_0\ldots i_p}).
$$
This is an abelian group. For
$s \in \check{\mathcal{C}}_{ord}^p(\mathcal{U}, \mathcal{F})$ we denote
$s_{i_0\ldots i_p}$ its value in $\mathcal{F}(U_{i_0\ldots i_p})$.
We define
$$
d : \check{\mathcal{C}}_{ord}^p(\mathcal{U}, \mathcal{F})
\longrightarrow
\check{\mathcal{C}}_{ord}^{p + 1}(\mathcal{U}, \mathcal{F})
$$
by the formula
$$
d(s)_{i_0\ldots i_{p + 1}}
=
\sum\nolimits_{j = 0}^{p + 1}
(-1)^j
s_{i_0\ldots \hat i_j \ldots i_{p + 1}}|_{U_{i_0\ldots i_{p + 1}}}
$$
for any $i_0 < \ldots < i_{p + 1}$. Note that this formula is identical
to Equation (\ref{equation-d-cech}).
It is straightforward to see that $d \circ d = 0$. In other words
$\check{\mathcal{C}}_{ord}^\bullet(\mathcal{U}, \mathcal{F})$ is a complex.

\begin{definition}
\label{definition-ordered-cech-complex}
Let $X$ be a topological space.
Let $\mathcal{U} : U = \bigcup_{i \in I} U_i$ be an open covering.
Assume given a total ordering on $I$.
Let $\mathcal{F}$ be an abelian presheaf on $X$.
The complex $\check{\mathcal{C}}_{ord}^\bullet(\mathcal{U}, \mathcal{F})$
is the {\it ordered {\v C}ech complex} associated to $\mathcal{F}$, the
open covering $\mathcal{U}$ and the given total ordering on $I$.
\end{definition}

\noindent
This complex is sometimes called the alternating {\v C}ech complex.
The reason is that there is an obvious comparison map between
the ordered {\v C}ech complex and the alternating {\v C}ech complex.
Namely, consider the map
$$
c :
\check{\mathcal{C}}_{ord}^\bullet(\mathcal{U}, \mathcal{F})
\longrightarrow
\check{\mathcal{C}}^\bullet(\mathcal{U}, \mathcal{F})
$$
given by the rule
$$
c(s)_{i_0\ldots i_p} =
\left\{
\begin{matrix}
0 &
\text{if} &
i_n = i_m \text{ for some } n \not = m\\
\text{sgn}(\sigma) s_{i_{\sigma(0)}\ldots i_{\sigma(p)}} &
\text{if} &
i_{\sigma(0)} < i_{\sigma(1)} < \ldots < i_{\sigma(p)}
\end{matrix}
\right.
$$
Here $\sigma$ denotes a permutation of $\{0, \ldots, p\}$ and
$\text{sgn}(\sigma)$ denotes its sign. The alternating and ordered
{\v C}ech complexes are often identified in the literature via the map
$c$. Namely we have the following easy lemma.

\begin{lemma}
\label{lemma-ordered-alternating}
Let $X$ be a topological space.
Let $\mathcal{U} : U = \bigcup_{i \in I} U_i$ be an open covering.
Assume $I$ comes equipped with a total ordering.
The map $c$ is a morphism of complexes. In fact it induces
an isomorphism
$$
c : \check{\mathcal{C}}_{ord}^\bullet(\mathcal{U}, \mathcal{F})
\to \check{\mathcal{C}}_{alt}^\bullet(\mathcal{U}, \mathcal{F})
$$
of complexes.
\end{lemma}

\begin{proof}
Omitted.
\end{proof}

\noindent
There is also a map
$$
\pi :
\check{\mathcal{C}}^\bullet(\mathcal{U}, \mathcal{F})
\longrightarrow
\check{\mathcal{C}}_{ord}^\bullet(\mathcal{U}, \mathcal{F})
$$
which is described by the rule
$$
\pi(s)_{i_0\ldots i_p} = s_{i_0\ldots i_p}
$$
whenever $i_0 < i_1 < \ldots < i_p$.

\begin{lemma}
\label{lemma-project-to-ordered}
Let $X$ be a topological space.
Let $\mathcal{U} : U = \bigcup_{i \in I} U_i$ be an open covering.
Assume $I$ comes equipped with a total ordering.
The map $\pi : \check{\mathcal{C}}^\bullet(\mathcal{U}, \mathcal{F})
\to \check{\mathcal{C}}_{ord}^\bullet(\mathcal{U}, \mathcal{F})$
is a morphism of complexes. It induces an isomorphism
$$
\pi : \check{\mathcal{C}}_{alt}^\bullet(\mathcal{U}, \mathcal{F})
\to \check{\mathcal{C}}_{ord}^\bullet(\mathcal{U}, \mathcal{F})
$$
of complexes which is a left inverse to the morphism $c$.
\end{lemma}

\begin{proof}
Omitted.
\end{proof}

\begin{remark}
\label{remark-compared-ordered-complexes}
This means that if we have two total orderings $<_1$ and $<_2$ on
the index set $I$, then we get an isomorphism of complexes
$\tau = \pi_2 \circ c_1 :
\check{\mathcal{C}}_{ord\text{-}1}(\mathcal{U}, \mathcal{F}) \to
\check{\mathcal{C}}_{ord\text{-}2}(\mathcal{U}, \mathcal{F})$.
It is clear that
$$
\tau(s)_{i_0 \ldots i_p} =
\text{sign}(\sigma) s_{i_{\sigma(0)} \ldots i_{\sigma(p)}}
$$
where $i_0 <_1 i_1 <_1 \ldots <_1 i_p$ and
$i_{\sigma(0)} <_2 i_{\sigma(1)} <_2 \ldots <_2 i_{\sigma(p)}$.
This is the sense in which the ordered {\v C}ech complex is independent
of the chosen total ordering.
\end{remark}

\begin{lemma}
\label{lemma-alternating-usual}
\begin{reference}
\cite[Chapter I, \S 3, no. 20, Proposition 2, pp. 213--214]{FAC}
\end{reference}
Let $X$ be a topological space.
Let $\mathcal{U} : U = \bigcup_{i \in I} U_i$ be an open covering.
Assume $I$ comes equipped with a total ordering.
The map $c \circ \pi$ is homotopic to the identity on
$\check{\mathcal{C}}^\bullet(\mathcal{U}, \mathcal{F})$.
In particular the inclusion map
$\check{\mathcal{C}}_{alt}^\bullet(\mathcal{U}, \mathcal{F}) \to
\check{\mathcal{C}}^\bullet(\mathcal{U}, \mathcal{F})$
is a homotopy equivalence.
\end{lemma}

\begin{proof}
For any multi-index $(i_0, \ldots, i_p) \in I^{p + 1}$ there exists
a unique permutation $\sigma : \{0, \ldots, p\} \to \{0, \ldots, p\}$
such that
$$
i_{\sigma(0)} \leq i_{\sigma(1)} \leq \ldots \leq i_{\sigma(p)}
\quad
\text{and}
\quad
\sigma(j) < \sigma(j + 1)
\quad
\text{if}
\quad
i_{\sigma(j)} = i_{\sigma(j + 1)}.
$$
We denote this permutation $\sigma = \sigma^{i_0 \ldots i_p}$.

\medskip\noindent
For any permutation $\sigma : \{0, \ldots, p\} \to \{0, \ldots, p\}$
and any $a$, $0 \leq a \leq p$ we denote $\sigma_a$
the unique permutation of $\{0, \ldots, p\}$ such that
$\sigma_a(j) = \sigma(j)$ for $0 \leq j < a$ and such that
$\sigma_a(a) < \sigma_a(a + 1) < \ldots < \sigma_a(p)$.
So if $p = 3$ and $\sigma$, $\tau$ are given by
$$
\begin{matrix}
\text{id} & 0 & 1 & 2 & 3 \\
\sigma & 3 & 2 & 1 & 0
\end{matrix}
\quad \text{and} \quad
\begin{matrix}
\text{id} & 0 & 1 & 2 & 3 \\
\tau & 3 & 0 & 2 & 1
\end{matrix}
$$
then we have
$$
\begin{matrix}
\text{id} & 0 & 1 & 2 & 3 \\
\sigma_0 & 0 & 1 & 2 & 3 \\
\sigma_1 & 3 & 0 & 1 & 2 \\
\sigma_2 & 3 & 2 & 0 & 1 \\
\sigma_3 & 3 & 2 & 1 & 0 \\
\end{matrix}
\quad \text{and} \quad
\begin{matrix}
\text{id} & 0 & 1 & 2 & 3 \\
\tau_0 & 0 & 1 & 2 & 3 \\
\tau_1 & 3 & 0 & 1 & 2 \\
\tau_2 & 3 & 0 & 1 & 2 \\
\tau_3 & 3 & 0 & 2 & 1 \\
\end{matrix}
$$
It is clear that always $\sigma_0 = \text{id}$ and $\sigma_p = \sigma$.

\medskip\noindent
Having introduced this notation we define for
$s \in \check{\mathcal{C}}^{p + 1}(\mathcal{U}, \mathcal{F})$
the element $h(s) \in \check{\mathcal{C}}^p(\mathcal{U}, \mathcal{F})$
to be the element with components
\begin{equation}
\label{equation-first-homotopy}
h(s)_{i_0\ldots i_p} =
\sum\nolimits_{0 \leq a \leq p}
(-1)^a \text{sign}(\sigma_a)
s_{i_{\sigma(0)} \ldots i_{\sigma(a)} i_{\sigma_a(a)} \ldots i_{\sigma_a(p)}}
\end{equation}
where $\sigma = \sigma^{i_0 \ldots i_p}$. The index
$i_{\sigma(a)}$ occurs twice in
$i_{\sigma(0)} \ldots i_{\sigma(a)} i_{\sigma_a(a)} \ldots i_{\sigma_a(p)}$
once in the first group of $a + 1$ indices and once in the second group
of $p - a + 1$ indices since $\sigma_a(j) = \sigma(a)$ for some
$j \geq a$ by definition of $\sigma_a$. Hence the sum makes sense since each
of the elements
$s_{i_{\sigma(0)} \ldots i_{\sigma(a)} i_{\sigma_a(a)} \ldots i_{\sigma_a(p)}}$
is defined over the open $U_{i_0 \ldots i_p}$.
Note also that for $a = 0$ we get $s_{i_0 \ldots i_p}$ and
for $a = p$ we get
$(-1)^p \text{sign}(\sigma) s_{i_{\sigma(0)} \ldots i_{\sigma(p)}}$.

\medskip\noindent
We claim that
$$
(dh + hd)(s)_{i_0 \ldots i_p} =
s_{i_0 \ldots i_p} -
\text{sign}(\sigma) s_{i_{\sigma(0)} \ldots i_{\sigma(p)}}
$$
where $\sigma = \sigma^{i_0 \ldots i_p}$. We omit the verification
of this claim. (There is a PARI/gp script called first-homotopy.gp
in the stacks-project subdirectory scripts which can be used to check
finitely many instances of this claim.
We wrote this script to make sure the signs are correct.)
Write
$$
\kappa :
\check{\mathcal{C}}^\bullet(\mathcal{U}, \mathcal{F})
\longrightarrow
\check{\mathcal{C}}^\bullet(\mathcal{U}, \mathcal{F})
$$
for the operator given by the rule
$$
\kappa(s)_{i_0 \ldots i_p} =
\text{sign}(\sigma^{i_0 \ldots i_p}) s_{i_{\sigma(0)} \ldots i_{\sigma(p)}}.
$$
The claim above implies that $\kappa$ is a morphism of complexes and that
$\kappa$ is homotopic to the identity map of the {\v C}ech complex.
This does not immediately imply the lemma since
the image of the operator $\kappa$ is not the alternating subcomplex.
Namely, the image of $\kappa$ is the ``semi-alternating'' complex
$\check{\mathcal{C}}_{semi\text{-}alt}^p(\mathcal{U}, \mathcal{F})$
where $s$ is a $p$-cochain of this complex if and only if
$$
s_{i_0 \ldots i_p} = \text{sign}(\sigma) s_{i_{\sigma(0)} \ldots i_{\sigma(p)}}
$$
for any $(i_0, \ldots, i_p) \in I^{p + 1}$ with
$\sigma = \sigma^{i_0 \ldots i_p}$.
We introduce yet another variant {\v C}ech complex, namely the semi-ordered
{\v C}ech complex defined by
$$
\check{\mathcal{C}}_{semi\text{-}ord}^p(\mathcal{U}, \mathcal{F})
=
\prod\nolimits_{i_0 \leq i_1 \leq \ldots \leq i_p}
\mathcal{F}(U_{i_0 \ldots i_p})
$$
It is easy to see that Equation (\ref{equation-d-cech}) also defines
a differential and hence that we get a complex. It is also clear
(analogous to Lemma \ref{lemma-project-to-ordered}) that the projection map
$$
\check{\mathcal{C}}_{semi\text{-}alt}^\bullet(\mathcal{U}, \mathcal{F})
\longrightarrow
\check{\mathcal{C}}_{semi\text{-}ord}^\bullet(\mathcal{U}, \mathcal{F})
$$
is an isomorphism of complexes.

\medskip\noindent
Hence the Lemma follows if we can show that the obvious inclusion map
$$
\check{\mathcal{C}}_{ord}^p(\mathcal{U}, \mathcal{F})
\longrightarrow
\check{\mathcal{C}}_{semi\text{-}ord}^p(\mathcal{U}, \mathcal{F})
$$
is a homotopy equivalence. To see this we use the homotopy
\begin{equation}
\label{equation-second-homotopy}
h(s)_{i_0 \ldots i_p} =
\left\{
\begin{matrix}
0 & \text{if} & i_0 < i_1 < \ldots < i_p \\
(-1)^a s_{i_0 \ldots i_{a - 1} i_a i_a i_{a + 1} \ldots i_p}
& \text{if} & i_0 < i_1 < \ldots < i_{a - 1} < i_a = i_{a + 1}
\end{matrix}
\right.
\end{equation}
We claim that
$$
(dh + hd)(s)_{i_0 \ldots i_p} =
\left\{
\begin{matrix}
0 & \text{if} & i_0 < i_1 < \ldots < i_p \\
s_{i_0 \ldots i_p}
& \text{else} &
\end{matrix}
\right.
$$
We omit the verification. (There is a PARI/gp script called second-homotopy.gp
in the stacks-project subdirectory scripts which can be used to check
finitely many instances of this claim.
We wrote this script to make sure the signs are correct.)
The claim clearly shows that the composition
$$
\check{\mathcal{C}}_{semi\text{-}ord}^\bullet(\mathcal{U}, \mathcal{F})
\longrightarrow
\check{\mathcal{C}}_{ord}^\bullet(\mathcal{U}, \mathcal{F})
\longrightarrow
\check{\mathcal{C}}_{semi\text{-}ord}^\bullet(\mathcal{U}, \mathcal{F})
$$
of the projection with the natural inclusion
is homotopic to the identity map as desired.
\end{proof}

\begin{lemma}
\label{lemma-cech-vanishing-cover-dimension}
\begin{reference}
\cite[Chapter I, \S 3, no. 20, Corollary, p. 214]{FAC}
\end{reference}
Let $X$ be a topological space. Let
$\mathcal{U} : X = \bigcup_{i \in I} U_i$ be an open covering. Let
$\mathcal{F}$ be an abelian sheaf on $X$. Let $n \geq 0$ be an integer.
Assume that
$$
U_{i_0 \ldots i_q} = \emptyset
$$
for every $q > n$ and every sequence of pairwise distinct elements
$i_0, \ldots, i_q$ of $I$. Then
$$
\check{H}^q(\mathcal{U}, \mathcal{F}) = 0
$$
for every $q > n$.
\end{lemma}

\begin{proof}
Choose a total ordering on $I$. Since $\mathcal{F}(\emptyset) = 0$, the
alternating \v{C}ech complex is zero in degrees greater than $n$.
Its inclusion into the usual \v{C}ech complex is a homotopy equivalence
by Lemma \ref{lemma-alternating-usual}. This proves the result.
\end{proof}

\begin{lemma}
\label{lemma-alternating-cech-trivial}
Let $X$ be a topological space. Let $\mathcal{F}$ be an abelian presheaf on $X$.
Let $\mathcal{U} : U = \bigcup_{i \in I} U_i$ be an open covering. If
$U_i = U$ for some $i \in I$, then the extended alternating {\v C}ech complex
$$
\mathcal{F}(U) \to \check{\mathcal{C}}_{alt}^\bullet(\mathcal{U}, \mathcal{F})
$$
obtained by putting $\mathcal{F}(U)$ in degree $-1$ with differential given by
the canonical map of $\mathcal{F}(U)$ into
$\check{\mathcal{C}}^0(\mathcal{U}, \mathcal{F})$
is homotopy equivalent to $0$. Similarly, for any total ordering on $I$
the extended ordered {\v C}ech complex
$$
\mathcal{F}(U) \to
\check{\mathcal{C}}_{ord}^\bullet(\mathcal{U}, \mathcal{F})
$$
is homotopy equivalent to $0$.
\end{lemma}

\begin{proof}[First proof]
Combine Lemmas \ref{lemma-cech-trivial} and \ref{lemma-alternating-usual}.
\end{proof}

\begin{proof}[Second proof]
Since the alternating and ordered {\v C}ech complexes are isomorphic
it suffices to prove this for the ordered one.
We will use standard notation: a cochain $s$ of degree $p$
in the extended ordered {\v C}ech complex has the form
$s = (s_{i_0 \ldots i_p})$ where $s_{i_0 \ldots i_p}$ is in
$\mathcal{F}(U_{i_0 \ldots i_p})$ and $i_0 < \ldots < i_p$.
With this notation we have
$$
d(x)_{i_0 \ldots i_{p + 1}} =
\sum\nolimits_j (-1)^j x_{i_0 \ldots \hat i_j \ldots i_p}
$$
Fix an index $i \in I$ with $U = U_i$.
As homotopy we use the maps
$$
h : \text{cochains of degree }p + 1 \to \text{cochains of degree }p
$$
given by the rule
$$
h(s)_{i_0 \ldots i_p} = 0 \text{ if } i \in \{i_0, \ldots, i_p\}
\text{ and }
h(s)_{i_0 \ldots i_p} =
(-1)^j s_{i_0 \ldots i_j i i_{j + 1} \ldots i_p} \text{ if not}
$$
Here $j$ is the unique index such that $i_j < i < i_{j + 1}$ in the
second case; also, since $U = U_i$ we have the equality
$$
\mathcal{F}(U_{i_0 \ldots i_p}) =
\mathcal{F}(U_{i_0 \ldots i_j i i_{j + 1} \ldots i_p})
$$
which we can use to make sense of thinking of
$(-1)^j s_{i_0 \ldots i_j i i_{j + 1} \ldots i_p}$
as an element of $\mathcal{F}(U_{i_0 \ldots i_p})$.
We will show by a computation that $d h + h d$ equals
the negative of the identity map which finishes the proof.
To do this fix $s$ a cochain of degree $p$ and let
$i_0 < \ldots < i_p$ be elements of $I$.

\medskip\noindent
Case I: $i \in \{i_0, \ldots, i_p\}$. Say $i = i_t$. Then we have
$h(d(s))_{i_0 \ldots i_p} = 0$. On the other hand we have
$$
d(h(s))_{i_0 \ldots i_p} =
\sum (-1)^j h(s)_{i_0 \ldots \hat i_j \ldots i_p} =
(-1)^t h(s)_{i_0 \ldots \hat i \ldots i_p} =
(-1)^t (-1)^{t - 1} s_{i_0 \ldots i_p}
$$
Thus $(dh + hd)(s)_{i_0 \ldots i_p} = -s_{i_0 \ldots i_p}$ as desired.

\medskip\noindent
Case II: $i \not \in \{i_0, \ldots, i_p\}$. Let $j$ be such that
$i_j < i < i_{j + 1}$. Then we see that
\begin{align*}
h(d(s))_{i_0 \ldots i_p}
& =
(-1)^j d(s)_{i_0 \ldots i_j i i_{j + 1} \ldots i_p} \\
& =
\sum\nolimits_{j' \leq j} (-1)^{j + j'}
s_{i_0 \ldots \hat i_{j'} \ldots i_j i i_{j + 1} \ldots i_p} -
s_{i_0 \ldots i_p} \\
&
+ \sum\nolimits_{j' > j} (-1)^{j + j' + 1}
s_{i_0 \ldots i_j i i_{j + 1} \ldots \hat i_{j'} \ldots i_p}
\end{align*}
On the other hand we have
\begin{align*}
d(h(s))_{i_0 \ldots i_p}
& =
\sum\nolimits_{j'} (-1)^{j'} h(s)_{i_0 \ldots \hat i_{j'} \ldots i_p} \\
& =
\sum\nolimits_{j' \leq j} (-1)^{j' + j - 1}
s_{i_0 \ldots \hat i_{j'} \ldots i_j i i_{j + 1} \ldots i_p} \\
& +
\sum\nolimits_{j' > j} (-1)^{j' + j}
s_{i_0 \ldots i_j i i_{j + 1} \ldots \hat i_{j'} \ldots i_p}
\end{align*}
Adding these up we obtain
$(dh + hd)(s)_{i_0 \ldots i_p} = - s_{i_0 \ldots i_p}$
as desired.
\end{proof}





\section{Alternative view of the {\v C}ech complex}
\label{section-locally-finite-cech}

\noindent
In this section we discuss an alternative way to establish the relationship
between the {\v C}ech complex and cohomology.

\begin{lemma}
\label{lemma-covering-resolution}
Let $X$ be a ringed space. Let $\mathcal{U} : X = \bigcup_{i \in I} U_i$
be an open covering of $X$. Let $\mathcal{F}$ be an $\mathcal{O}_X$-module.
Denote $\mathcal{F}_{i_0 \ldots i_p}$ the restriction of
$\mathcal{F}$ to $U_{i_0 \ldots i_p}$. There exists a complex
${\mathfrak C}^\bullet(\mathcal{U}, \mathcal{F})$
of $\mathcal{O}_X$-modules with
$$
{\mathfrak C}^p(\mathcal{U}, \mathcal{F}) =
\prod\nolimits_{i_0 \ldots i_p}
(j_{i_0 \ldots i_p})_* \mathcal{F}_{i_0 \ldots i_p}
$$
and differential
$d : {\mathfrak C}^p(\mathcal{U}, \mathcal{F})
\to {\mathfrak C}^{p + 1}(\mathcal{U}, \mathcal{F})$
as in Equation (\ref{equation-d-cech}). Moreover, there exists a canonical
map
$$
\mathcal{F} \to {\mathfrak C}^\bullet(\mathcal{U}, \mathcal{F})
$$
which is a quasi-isomorphism, i.e.,
${\mathfrak C}^\bullet(\mathcal{U}, \mathcal{F})$
is a resolution of $\mathcal{F}$.
\end{lemma}

\begin{proof}
We check
$$
0 \to \mathcal{F} \to \mathfrak{C}^0(\mathcal{U}, \mathcal{F}) \to
\mathfrak{C}^1(\mathcal{U}, \mathcal{F}) \to  \ldots
$$
is exact on stalks. Let $x \in X$ and choose $i_{\text{fix}} \in I$
such that $x \in U_{i_{\text{fix}}}$. Then define
$$
h : \mathfrak{C}^p(\mathcal{U}, \mathcal{F})_x
\to \mathfrak{C}^{p - 1}(\mathcal{U}, \mathcal{F})_x
$$
as follows: If $s \in \mathfrak{C}^p(\mathcal{U}, \mathcal{F})_x$, take
a representative
$$
\widetilde{s} \in
\mathfrak{C}^p(\mathcal{U}, \mathcal{F})(V) =
\prod\nolimits_{i_0 \ldots i_p}
\mathcal{F}(V \cap U_{i_0} \cap \ldots \cap U_{i_p})
$$
defined on some neighborhood $V$ of $x$, and set
$$
h(s)_{i_0 \ldots i_{p - 1}} =
\widetilde{s}_{i_{\text{fix}} i_0 \ldots i_{p - 1}, x}.
$$
By the same formula (for $p = 0$) we get a map
$\mathfrak{C}^{0}(\mathcal{U},\mathcal{F})_x \to \mathcal{F}_x$.
We compute formally as follows:
\begin{align*}
(dh + hd)(s)_{i_0 \ldots i_p}
& =
\sum\nolimits_{j = 0}^p
(-1)^j
h(s)_{i_0 \ldots \hat i_j \ldots i_p}
+
d(s)_{i_{\text{fix}} i_0 \ldots i_p}\\
& =
\sum\nolimits_{j = 0}^p
(-1)^j
s_{i_{\text{fix}} i_0 \ldots \hat i_j \ldots i_p}
+
s_{i_0 \ldots i_p}
+
\sum\nolimits_{j = 0}^p
(-1)^{j + 1}
s_{i_{\text{fix}} i_0 \ldots \hat i_j \ldots i_p} \\
& =
s_{i_0 \ldots i_p}
\end{align*}
This shows $h$ is a homotopy from the identity map of
the extended complex
$$
0 \to \mathcal{F}_x \to \mathfrak{C}^0(\mathcal{U}, \mathcal{F})_x
\to \mathfrak{C}^1(\mathcal{U}, \mathcal{F})_x \to \ldots
$$
to zero and we conclude.
\end{proof}

\noindent
With this lemma it is easy to reprove the {\v C}ech to cohomology spectral
sequence of Lemma \ref{lemma-cech-spectral-sequence}. Namely,
let $X$, $\mathcal{U}$, $\mathcal{F}$ as in
Lemma \ref{lemma-covering-resolution}
and let $\mathcal{F} \to \mathcal{I}^\bullet$ be an injective resolution.
Then we may consider the double complex
$$
A^{\bullet, \bullet} =
\Gamma(X, {\mathfrak C}^\bullet(\mathcal{U}, \mathcal{I}^\bullet)).
$$
By construction we have
$$
A^{p, q} = \prod\nolimits_{i_0 \ldots i_p} \mathcal{I}^q(U_{i_0 \ldots i_p})
$$
Consider the two spectral sequences of
Homology, Section \ref{homology-section-double-complex} associated
to this double complex, see especially
Homology, Lemma \ref{homology-lemma-ss-double-complex}.
For the spectral sequence $({}'E_r, {}'d_r)_{r \geq 0}$ we get
${}'E_2^{p, q} = \check{H}^p(\mathcal{U}, \underline{H}^q(\mathcal{F}))$
because taking products is exact
(Homology, Lemma \ref{homology-lemma-product-abelian-groups-exact}).
For the spectral sequence $({}''E_r, {}''d_r)_{r \geq 0}$ we get
${}''E_2^{p, q} = 0$ if $p > 0$ and ${}''E_2^{0, q} = H^q(X, \mathcal{F})$.
Namely, for fixed $q$ the complex of sheaves
${\mathfrak C}^\bullet(\mathcal{U}, \mathcal{I}^q)$
is a resolution (Lemma \ref{lemma-covering-resolution})
of the injective sheaf $\mathcal{I}^q$
by injective sheaves (by Lemmas \ref{lemma-cohomology-of-open} and
\ref{lemma-pushforward-injective-flat}
and
Homology, Lemma \ref{homology-lemma-product-injectives}).
Hence the cohomology of
$\Gamma(X, {\mathfrak C}^\bullet(\mathcal{U}, \mathcal{I}^q))$
is zero in positive degrees and equal to $\Gamma(X, \mathcal{I}^q)$
in degree $0$. Taking cohomology of the next differential
we get our claim about the spectral sequence $({}''E_r, {}''d_r)_{r \geq 0}$.
Whence the result since both spectral sequences converge to the
cohomology of the associated total complex of $A^{\bullet, \bullet}$.

\begin{definition}
\label{definition-covering-locally-finite}
Let $X$ be a topological space.
An open covering $X = \bigcup_{i \in I} U_i$ is said to be
{\it locally finite} if for every $x \in X$ there exists an open neighbourhood
$W$ of $x$ such that $\{i \in I \mid W \cap U_i \not = \emptyset\}$ is finite.
\end{definition}

\begin{remark}
\label{remark-locally-finite-sections}
Let $X = \bigcup_{i \in I} U_i$ be a locally finite open covering.
Denote $j_i : U_i \to X$ the inclusion map. Suppose that for each $i$
we are given an abelian sheaf $\mathcal{F}_i$ on $U_i$. Consider the
abelian sheaf $\mathcal{G} = \bigoplus_{i \in I} (j_i)_*\mathcal{F}_i$.
Then for $V \subset X$ open we actually have
$$
\Gamma(V, \mathcal{G}) = \prod\nolimits_{i \in I} \mathcal{F}_i(V \cap U_i).
$$
In other words we have
$$
\bigoplus\nolimits_{i \in I} (j_i)_*\mathcal{F}_i =
\prod\nolimits_{i \in I} (j_i)_*\mathcal{F}_i
$$
This seems strange until you realize that the direct sum of a collection
of sheaves is the sheafification of what you think it should be.
See discussion in Modules, Section \ref{modules-section-kernels}.
Thus we conclude that in this case the complex of
Lemma \ref{lemma-covering-resolution} has terms
$$
{\mathfrak C}^p(\mathcal{U}, \mathcal{F}) =
\bigoplus\nolimits_{i_0 \ldots i_p}
(j_{i_0 \ldots i_p})_* \mathcal{F}_{i_0 \ldots i_p}
$$
which is sometimes useful.
\end{remark}











\section{{\v C}ech cohomology of complexes}
\label{section-cech-cohomology-of-complexes}

\noindent
In general for sheaves of abelian groups
${\mathcal F}$ and ${\mathcal G}$ on $X$ there is a cup product map
$$
H^i(X, {\mathcal F}) \times H^j(X, {\mathcal G})
\longrightarrow
H^{i + j}(X, {\mathcal F} \otimes_{\mathbf Z} {\mathcal G}).
$$
In this section we define it using {\v C}ech cocycles by an explicit formula
for the cup product. If you are worried about the fact that cohomology may not
equal {\v C}ech cohomology, then you can use hypercoverings and still
use the cocycle notation. This also has the advantage that
it works to define the cup product for hypercohomology on any topos (insert
future reference here).

\medskip\noindent
Let ${\mathcal F}^\bullet$ be a bounded below complex of presheaves of abelian
groups on $X$. We can often compute $H^n(X, {\mathcal F}^\bullet)$
using {\v C}ech cocycles. Namely, let
${\mathcal U} : X = \bigcup_{i \in I} U_i$
be an open covering of $X$. Since the {\v C}ech complex
$\check{\mathcal{C}}^\bullet(\mathcal{U}, \mathcal{F})$
(Definition \ref{definition-cech-complex})
is functorial in the presheaf $\mathcal{F}$ we obtain a double complex
$\check{\mathcal{C}}^\bullet(\mathcal{U}, \mathcal{F}^\bullet)$.
The associated total complex to
$\check{\mathcal{C}}^\bullet({\mathcal U}, {\mathcal F}^\bullet)$
is the complex with degree $n$ term
$$
\text{Tot}^n(\check{\mathcal{C}}^\bullet({\mathcal U}, {\mathcal F}^\bullet))
=
\bigoplus\nolimits_{p + q = n}
\prod\nolimits_{i_0\ldots i_p} {\mathcal F}^q(U_{i_0\ldots i_p})
$$
see
Homology, Definition \ref{homology-definition-associated-simple-complex}.
A typical element in $\text{Tot}^n$ will be denoted
$\alpha = \{\alpha_{i_0\ldots i_p}\}$ where
$\alpha_{i_0 \ldots i_p} \in \mathcal{F}^q(U_{i_0\ldots i_p})$.
In other words the $\mathcal{F}$-degree of $\alpha_{i_0\ldots i_p}$ is
$q = n - p$. This notation requires us to be aware of the degree $\alpha$
lives in at all times. We indicate this situation by the formula
$\deg_{\mathcal F}(\alpha_{i_0\ldots i_p}) = q$.
According to our conventions in
Homology, Definition \ref{homology-definition-associated-simple-complex}
the differential of an element $\alpha$ of degree $n$ is given by
$$
d(\alpha)_{i_0\ldots i_{p + 1}}
=
\sum\nolimits_{j = 0}^{p + 1}
(-1)^j \alpha_{i_0 \ldots \hat i_j \ldots i_{p + 1}} +
(-1)^{p + 1}d_{{\mathcal F}}(\alpha_{i_0 \ldots i_{p + 1}})
$$
where $d_\mathcal{F}$ denotes the differential on the complex
$\mathcal{F}^\bullet$.
The expression $\alpha_{i_0 \ldots \hat i_j \ldots i_{p + 1}}$ means the
restriction of $\alpha_{i_0 \ldots \hat i_j \ldots i_{p + 1}}
\in {\mathcal F}(U_{i_0\ldots\hat i_j\ldots i_{p + 1}})$ to
$U_{i_0 \ldots i_{p + 1}}$.

\medskip\noindent
The construction of
$\text{Tot}(\check{\mathcal{C}}^\bullet({\mathcal U}, {\mathcal F}^\bullet))$
is functorial in ${\mathcal F}^\bullet$. As well there is a functorial
transformation
\begin{equation}
\label{equation-global-sections-to-cech}
\Gamma(X, {\mathcal F}^\bullet)
\longrightarrow
\text{Tot}(\check{\mathcal{C}}^\bullet({\mathcal U}, {\mathcal F}^\bullet))
\end{equation}
of complexes defined by the following rule: The section
$s\in \Gamma(X, {\mathcal F}^n)$
is mapped to the element $\alpha = \{\alpha_{i_0\ldots i_p}\}$
with $\alpha_{i_0} = s|_{U_{i_0}}$ and $\alpha_{i_0\ldots i_p} = 0$
for $p > 0$.

\medskip\noindent
Refinements. Let ${\mathcal V} = \{ V_j \}_{j\in J}$ be a
refinement of ${\mathcal U}$. This means there is a map $t: J \to I$
such that $V_j \subset U_{t(j)}$ for all $j\in J$. This gives
rise to a functorial transformation
\begin{equation}
\label{equation-transformation}
T_t :
\text{Tot}(\check{\mathcal{C}}^\bullet({\mathcal U}, {\mathcal F}^\bullet))
\longrightarrow
\text{Tot}(\check{\mathcal{C}}^\bullet({\mathcal V}, {\mathcal F}^\bullet)).
\end{equation}
defined by the rule
$$
T_t(\alpha)_{j_0\ldots j_p}
=
\alpha_{t(j_0)\ldots t(j_p)}|_{V_{j_0\ldots j_p}}.
$$
Given two maps $t, t' : J \to I$ as above the maps
$T_t$ and $T_{t'}$ constructed above are homotopic.
The homotopy is given by
$$
h(\alpha)_{j_0\ldots j_p}
=
\sum\nolimits_{a = 0}^{p}
(-1)^a
\alpha_{t(j_0)\ldots t(j_a) t'(j_a) \ldots t'(j_p)}
$$
for an element $\alpha$ of degree $n$. This works
because of the following computation, again with
$\alpha$ an element of degree $n$ (so $d(\alpha)$
has degree $n + 1$ and $h(\alpha)$ has degree $n - 1$):
\begin{align*}
(
d(h(\alpha)) + h(d(\alpha))
)_{j_0\ldots j_p}
= &
\sum\nolimits_{k = 0}^p
(-1)^k
h(\alpha)_{j_0 \ldots \hat j_k \ldots j_p}
+ \\
&
(-1)^p
d_{\mathcal F}(h(\alpha)_{j_0 \ldots j_p})
+ \\
&
\sum\nolimits_{a = 0}^p
(-1)^a
d(\alpha)_{t(j_0) \ldots t(j_a) t'(j_a) \ldots t'(j_p)}
\\
= &
\sum\nolimits_{k = 0}^p
\sum\nolimits_{a = 0}^{k - 1}
(-1)^{k + a}
\alpha_{t(j_0)\ldots t(j_a)t'(j_a)\ldots \hat{t'(j_k)}\ldots t'(j_p)}
+ \\
&
\sum\nolimits_{k = 0}^p
\sum\nolimits_{a = k + 1}^p
(-1)^{k + a - 1}
\alpha_{t(j_0)\ldots \hat{t(j_k)}\ldots t(j_a)t'(j_a)\ldots t'(j_p)}
+ \\
&
\sum\nolimits_{a = 0}^p
(-1)^{p + a}
d_{\mathcal F}(\alpha_{t(j_0)\ldots t(j_a) t'(j_a) \ldots t'(j_p)})
+ \\
&
\sum\nolimits_{a = 0}^p
\sum\nolimits_{k = 0}^a
(-1)^{a + k}
\alpha_{t(j_0)\ldots\hat{t(j_k)}\ldots t(j_a)t'(j_a)\ldots t'(j_p)}
+ \\
&
\sum\nolimits_{a = 0}^p
\sum\nolimits_{k = a}^p
(-1)^{a + k + 1}
\alpha_{t(j_0) \ldots t(j_a) t'(j_a) \ldots \hat{t'(j_k)} \ldots t'(j_p)}
+ \\
&
\sum\nolimits_{a = 0}^p
(-1)^{a + p + 1}
d_{\mathcal F}(\alpha_{t(j_0)\ldots t(j_a) t'(j_a) \ldots t'(j_p)})
\\
= &
\alpha_{t'(j_0)\ldots t'(j_p)} +
(-1)^{2p + 1}\alpha_{t(j_0)\ldots t(j_p)}
\\
= &
T_{t'}(\alpha)_{j_0\ldots j_p} - T_t(\alpha)_{j_0\ldots j_p}
\end{align*}
We leave it to the reader to verify the cancellations. (Note that the
terms having both $k$ and $a$ in the 1st, 2nd and 4th, 5th summands
cancel, except the ones where $a = k$ which only occur in the 4th and 5th
and these cancel against each other except for the two desired terms.)
It follows that the induced map
$$
H^n(T_t) :
H^n(
\text{Tot}(\check{\mathcal{C}}^\bullet({\mathcal U}, {\mathcal F}^\bullet))
)
\to
H^n(
\text{Tot}(\check{\mathcal{C}}^\bullet({\mathcal V}, {\mathcal F}^\bullet))
)
$$
is independent of the choice of $t$. We define
{\it {\v C}ech hypercohomology} as the colimit of the
{\v C}ech cohomology groups
over all refinements via the maps $H^\bullet(T_t)$.

\medskip\noindent
In the colimit (over all open coverings of $X$) the following lemma provides
a map of {\v C}ech hypercohomology into cohomology, which is often an
isomorphism and is always an isomorphism if we use hypercoverings.

\begin{lemma}
\label{lemma-cech-complex-complex}
Let $(X, \mathcal{O}_X)$ be a ringed space.
Let $\mathcal{U} : X = \bigcup_{i \in I} U_i$ be
an open covering. For a bounded below complex $\mathcal{F}^\bullet$
of $\mathcal{O}_X$-modules there is a canonical map
$$
\text{Tot}(\check{\mathcal{C}}^\bullet(\mathcal{U}, \mathcal{F}^\bullet))
\longrightarrow
R\Gamma(X, \mathcal{F}^\bullet)
$$
functorial in $\mathcal{F}^\bullet$ and compatible with
(\ref{equation-global-sections-to-cech}) and (\ref{equation-transformation}).
There is a spectral sequence $(E_r, d_r)_{r \geq 0}$ with
$$
E_2^{p, q} =
H^p(\text{Tot}(\check{\mathcal{C}}^\bullet(\mathcal{U},
\underline{H}^q(\mathcal{F}^\bullet)))
$$
converging to $H^{p + q}(X, \mathcal{F}^\bullet)$.
\end{lemma}

\begin{proof}
Let ${\mathcal I}^\bullet$ be a bounded below complex of injectives.
The map (\ref{equation-global-sections-to-cech}) for
$\mathcal{I}^\bullet$ is a map
$\Gamma(X, {\mathcal I}^\bullet) \to
\text{Tot}(\check{\mathcal{C}}^\bullet({\mathcal U}, {\mathcal I}^\bullet))$.
This is a quasi-isomorphism of complexes of abelian groups
as follows from
Homology, Lemma \ref{homology-lemma-double-complex-gives-resolution}
applied to the double complex
$\check{\mathcal{C}}^\bullet({\mathcal U}, {\mathcal I}^\bullet)$ using
Lemma \ref{lemma-injective-trivial-cech}.
Suppose ${\mathcal F}^\bullet \to {\mathcal I}^\bullet$ is a quasi-isomorphism
of ${\mathcal F}^\bullet$ into a bounded below complex of injectives.
Since $R\Gamma(X, {\mathcal F}^\bullet)$ is represented by the complex
$\Gamma(X, {\mathcal I}^\bullet)$ we obtain the map of the lemma
using
$$
\text{Tot}(\check{\mathcal{C}}^\bullet({\mathcal U}, {\mathcal F}^\bullet))
\longrightarrow
\text{Tot}(\check{\mathcal{C}}^\bullet({\mathcal U}, {\mathcal I}^\bullet)).
$$
We omit the verification of functoriality and compatibilities.
To construct the spectral sequence of the lemma, choose a Cartan-Eilenberg
resolution $\mathcal{F}^\bullet \to \mathcal{I}^{\bullet, \bullet}$, see
Derived Categories, Lemma \ref{derived-lemma-cartan-eilenberg}. In this
case $\mathcal{F}^\bullet \to \text{Tot}(\mathcal{I}^{\bullet, \bullet})$
is an injective resolution and hence
$$
\text{Tot}(\check{\mathcal{C}}^\bullet({\mathcal U},
\text{Tot}({\mathcal I}^{\bullet, \bullet})))
$$
computes $R\Gamma(X, \mathcal{F}^\bullet)$ as we've seen above.
By Homology, Remark \ref{homology-remark-triple-complex}
we can view this as the total complex associated to the
triple complex
$\check{\mathcal{C}}^\bullet({\mathcal U}, {\mathcal I}^{\bullet, \bullet})$
hence, using the same remark we can view it as the total complex
associate to the double complex $A^{\bullet, \bullet}$ with terms
$$
A^{n, m} =
\bigoplus\nolimits_{p + q = n}
\check{\mathcal{C}}^p({\mathcal U}, \mathcal{I}^{q, m})
$$
Since $\mathcal{I}^{q, \bullet}$ is an injective resolution of
$\mathcal{F}^q$ we can apply the first spectral sequence associated to
$A^{\bullet, \bullet}$
(Homology, Lemma \ref{homology-lemma-ss-double-complex})
to get a spectral sequence with
$$
E_1^{n, m} =
\bigoplus\nolimits_{p + q = n}
\check{\mathcal{C}}^p(\mathcal{U}, \underline{H}^m(\mathcal{F}^q))
$$
which is the $n$th term of the complex
$\text{Tot}(\check{\mathcal{C}}^\bullet(\mathcal{U},
\underline{H}^m(\mathcal{F}^\bullet))$. Hence we obtain
$E_2$ terms as described in the lemma. Convergence by
Homology, Lemma \ref{homology-lemma-first-quadrant-ss}.
\end{proof}

\begin{lemma}
\label{lemma-cech-complex-complex-computes}
Let $(X, \mathcal{O}_X)$ be a ringed space.
Let $\mathcal{U} : X = \bigcup_{i \in I} U_i$ be
an open covering. Let $\mathcal{F}^\bullet$ be a bounded below complex
of $\mathcal{O}_X$-modules. If $H^i(U_{i_0 \ldots i_p}, \mathcal{F}^q) = 0$
for all $i > 0$ and all $p, i_0, \ldots, i_p, q$, then the map
$
\text{Tot}(\check{\mathcal{C}}^\bullet(\mathcal{U}, \mathcal{F}^\bullet))
\to
R\Gamma(X, \mathcal{F}^\bullet)
$
of Lemma \ref{lemma-cech-complex-complex} is an isomorphism.
\end{lemma}

\begin{proof}
Immediate from the spectral sequence of Lemma \ref{lemma-cech-complex-complex}.
\end{proof}

\begin{remark}
\label{remark-shift-complex-cech-complex}
Let $(X, \mathcal{O}_X)$ be a ringed space. Let
$\mathcal{U} : X = \bigcup_{i \in I} U_i$ be
an open covering. Let $\mathcal{F}^\bullet$ be a bounded below complex
of $\mathcal{O}_X$-modules. Let $b$ be an integer.
We claim there is a commutative diagram
$$
\xymatrix{
\text{Tot}(\check{\mathcal{C}}^\bullet(\mathcal{U}, \mathcal{F}^\bullet))[b]
\ar[r] \ar[d]_\gamma &
R\Gamma(X, \mathcal{F}^\bullet)[b] \ar[d] \\
\text{Tot}(\check{\mathcal{C}}^\bullet(\mathcal{U}, \mathcal{F}^\bullet[b]))
\ar[r] &
R\Gamma(X, \mathcal{F}^\bullet[b])
}
$$
in the derived category where the map $\gamma$ is the map on complexes
constructed in Homology, Remark \ref{homology-remark-shift-double-complex}.
This makes sense because the double complex
$\check{\mathcal{C}}^\bullet(\mathcal{U}, \mathcal{F}^\bullet[b])$
is clearly the same as the double complex
$\check{\mathcal{C}}^\bullet(\mathcal{U}, \mathcal{F}^\bullet)[0, b]$
introduced in Homology, Remark \ref{homology-remark-shift-double-complex}.
To check that the diagram commutes, we may choose an injective resolution
$\mathcal{F}^\bullet \to \mathcal{I}^\bullet$ as in the proof of
Lemma \ref{lemma-cech-complex-complex}. Chasing diagrams, we see that
it suffices to check the diagram commutes when we replace $\mathcal{F}^\bullet$
by $\mathcal{I}^\bullet$. Then we consider the extended diagram
$$
\xymatrix{
\Gamma(X, \mathcal{I}^\bullet)[b] \ar[r] \ar[d] &
\text{Tot}(\check{\mathcal{C}}^\bullet(\mathcal{U}, \mathcal{I}^\bullet))[b]
\ar[r] \ar[d]_\gamma &
R\Gamma(X, \mathcal{I}^\bullet)[b] \ar[d] \\
\Gamma(X, \mathcal{I}^\bullet[b]) \ar[r] &
\text{Tot}(\check{\mathcal{C}}^\bullet(\mathcal{U}, \mathcal{I}^\bullet[b]))
\ar[r] &
R\Gamma(X, \mathcal{I}^\bullet[b])
}
$$
where the left horizontal arrows are (\ref{equation-global-sections-to-cech}).
Since in this case the horizontal arrows are isomorphisms in the derived
category (see proof of Lemma \ref{lemma-cech-complex-complex}) it
suffices to show that the left square commutes. This is true because
the map $\gamma$ uses the sign $1$ on the summands
$\check{\mathcal{C}}^0(\mathcal{U}, \mathcal{I}^{q + b})$, see
formula in Homology, Remark \ref{homology-remark-shift-double-complex}.
\end{remark}

\noindent
Let $X$ be a topological space, let $\mathcal{U} : X = \bigcup_{i \in I} U_i$
be an open covering, and let $\mathcal{F}^\bullet$ be a bounded below
complex of presheaves of abelian groups. Consider the map
$\tau :
\text{Tot}(\check{\mathcal{C}}^\bullet({\mathcal U}, {\mathcal F}^\bullet))
\to
\text{Tot}(\check{\mathcal{C}}^\bullet({\mathcal U}, {\mathcal F}^\bullet))$
defined by
$$
\tau(\alpha)_{i_0 \ldots i_p} = (-1)^{p(p + 1)/2} \alpha_{i_p \ldots i_0}.
$$
Then we have for an element $\alpha$ of degree $n$ that
\begin{align*}
& d(\tau(\alpha))_{i_0 \ldots i_{p + 1}} \\
& =
\sum\nolimits_{j = 0}^{p + 1}
(-1)^j
\tau(\alpha)_{i_0 \ldots \hat i_j \ldots i_{p + 1}}
+
(-1)^{p + 1}
d_{\mathcal F}(\tau(\alpha)_{i_0 \ldots i_{p + 1}})
\\
& =
\sum\nolimits_{j = 0}^{p + 1}
(-1)^{j + \frac{p(p + 1)}{2}}
\alpha_{i_{p + 1} \ldots \hat i_j \ldots i_0}
+
(-1)^{p + 1 + \frac{(p + 1)(p + 2)}{2}}
d_{\mathcal F}(\alpha_{i_{p + 1} \ldots i_0})
\end{align*}
On the other hand we have
\begin{align*}
& \tau(d(\alpha))_{i_0\ldots i_{p + 1}} \\
& =
(-1)^{\frac{(p + 1)(p + 2)}{2}} d(\alpha)_{i_{p + 1} \ldots i_0}
\\
& =
(-1)^{\frac{(p + 1)(p + 2)}{2}}
\left(
\sum\nolimits_{j = 0}^{p + 1}
(-1)^j
\alpha_{i_{p + 1}\ldots \hat i_{p + 1 - j} \ldots i_0}
+
(-1)^{p + 1}
d_{\mathcal F}(\alpha_{i_{p + 1}\ldots i_0})
\right)
\end{align*}
Thus we conclude that $d(\tau(\alpha)) = \tau(d(\alpha))$
because $p(p + 1)/2 \equiv (p + 1)(p + 2)/2 + p + 1 \bmod 2$. In other words
$\tau$ is an endomorphism of the complex
$\text{Tot}(\check{\mathcal{C}}^\bullet({\mathcal U}, {\mathcal F}^\bullet))$.
Note that the diagram
$$
\begin{matrix}
\Gamma(X, {\mathcal F}^\bullet) &
\longrightarrow &
\text{Tot}(\check{\mathcal{C}}^\bullet({\mathcal U}, {\mathcal F}^\bullet)) \\
\downarrow \text{id} & & \downarrow \tau \\
\Gamma(X, {\mathcal F}^\bullet) &
\longrightarrow &
\text{Tot}(\check{\mathcal{C}}^\bullet({\mathcal U}, {\mathcal F}^\bullet))
\end{matrix}
$$
commutes. In addition $\tau$ is clearly compatible with refinements.
This suggests that $\tau$ acts as the identity on {\v C}ech cohomology
(i.e., in the colimit -- provided {\v C}ech hypercohomology agrees with
hypercohomology, which is always the case if we use hypercoverings).
We claim that $\tau$ actually is homotopic to the identity on the
total {\v C}ech complex
$\text{Tot}(\check{\mathcal{C}}^\bullet({\mathcal U}, {\mathcal F}^\bullet))$.
To prove this, we use as homotopy
$$
h(\alpha)_{i_0 \ldots i_p}
=
\sum\nolimits_{a = 0}^p
\epsilon_p(a)
\alpha_{i_0 \ldots i_a i_p \ldots i_a}
\quad\text{with}\quad
\epsilon_p(a) = (-1)^{\frac{(p - a)(p - a - 1)}{2} + p}
$$
for $\alpha$ of degree $n$. As usual we omit writing
$|_{U_{i_0 \ldots i_p}}$. This works
because of the following computation, again with
$\alpha$ an element of degree $n$:
\begin{align*}
(d(h(\alpha)) + h(d(\alpha)))_{i_0 \ldots i_p}
= &
\sum\nolimits_{k = 0}^p
(-1)^k
h(\alpha)_{i_0 \ldots \hat i_k \ldots i_p}
+ \\
&
(-1)^p
d_{\mathcal F}(h(\alpha)_{i_0 \ldots i_p})
+ \\
&
\sum\nolimits_{a = 0}^p
\epsilon_p(a)
d(\alpha)_{i_0 \ldots i_a i_p \ldots i_a}
\\
= &
\sum\nolimits_{k = 0}^p
\sum\nolimits_{a = 0}^{k - 1}
(-1)^k \epsilon_{p - 1}(a)
\alpha_{i_0 \ldots i_a i_p \ldots \hat{i_k} \ldots i_a}
+ \\
&
\sum\nolimits_{k = 0}^p
\sum\nolimits_{a = k + 1}^p
(-1)^k \epsilon_{p - 1}(a - 1)
\alpha_{i_0 \ldots \hat{i_k} \ldots i_a i_p \ldots i_a}
+ \\
&
\sum\nolimits_{a = 0}^p
(-1)^p \epsilon_p(a)
d_{\mathcal F}(\alpha_{i_0 \ldots i_a i_p \ldots i_a})
+ \\
&
\sum\nolimits_{a = 0}^p
\sum\nolimits_{k = 0}^a
\epsilon_p(a) (-1)^k
\alpha_{i_0 \ldots \hat{i_k} \ldots i_a i_p \ldots i_a}
+ \\
&
\sum\nolimits_{a = 0}^p
\sum\nolimits_{k = a}^p
\epsilon_p(a) (-1)^{p + a + 1 - k}
\alpha_{i_0 \ldots i_a i_p \ldots \hat{i_k} \ldots i_a}
+ \\
&
\sum\nolimits_{a = 0}^p
\epsilon_p(a) (-1)^{p + 1}
d_{\mathcal F}(\alpha_{i_0 \ldots i_a i_p \ldots i_a})
\\
= &
\epsilon_p(0) \alpha_{i_p \ldots i_0} +
\epsilon_p(p) (-1)^{p + 1} \alpha_{i_0 \ldots i_p} \\
= &
(-1)^{\frac{p(p + 1)}{2}}\alpha_{i_p \ldots i_0}
- \alpha_{i_0 \ldots i_p}
\end{align*}
The cancellations follow because
$$
(-1)^k \epsilon_{p - 1}(a) + \epsilon_p(a)(-1)^{p + a + 1 - k} = 0
\quad\text{and}\quad
(-1)^k\epsilon_{p - 1}(a - 1) + \epsilon_p(a) (-1)^k = 0
$$
We leave it to the reader to verify the cancellations.

\medskip\noindent
Suppose we have two bounded below complexes of abelian sheaves
${\mathcal F}^\bullet$ and ${\mathcal G}^\bullet$. We define the complex
$\text{Tot}({\mathcal F}^\bullet\otimes_{\mathbf Z} {\mathcal G}^\bullet)$
to be to complex with terms
$\bigoplus_{p + q = n} {\mathcal F}^p \otimes {\mathcal G}^q$
and differential according to the rule
\begin{equation}
\label{equation-differential-tensor-product-complexes}
d(\alpha \otimes \beta) =
d(\alpha)\otimes \beta + (-1)^{\deg(\alpha)} \alpha \otimes d(\beta)
\end{equation}
when $\alpha$ and $\beta$ are homogeneous, see
Homology, Definition \ref{homology-definition-associated-simple-complex}.

\medskip\noindent
Suppose that $M^\bullet$ and $N^\bullet$ are two bounded below
complexes of abelian groups. Then if $m$, resp.\ $n$
is a cocycle for $M^\bullet$, resp.\ $N^\bullet$, it is immediate
that $m \otimes n$ is a cocycle for $\text{Tot}(M^\bullet\otimes N^\bullet)$.
Hence a cup product
$$
H^i(M^\bullet) \times H^j(N^\bullet)
\longrightarrow
H^{i + j}(Tot(M^\bullet\otimes N^\bullet)).
$$
This is discussed also in
More on Algebra, Section \ref{more-algebra-section-products-tor}.

\medskip\noindent
So the construction of the cup product in hypercohomology
of complexes rests on a construction of a map of complexes
\begin{equation}
\label{equation-needs-signs}
\text{Tot}\left(
\text{Tot}(\check{\mathcal{C}}^\bullet({\mathcal U}, {\mathcal F}^\bullet))
\otimes_{\mathbf Z}
\text{Tot}(\check{\mathcal{C}}^\bullet({\mathcal U}, {\mathcal G}^\bullet))
\right)
\longrightarrow
\text{Tot}(
\check{\mathcal{C}}^\bullet({\mathcal U},
\text{Tot}({\mathcal F}^\bullet\otimes {\mathcal G}^\bullet)
))
\end{equation}
This map is denoted $\cup$ and is given by the rule
$$
(\alpha \cup \beta)_{i_0 \ldots i_p}
=
\sum\nolimits_{r = 0}^p
\epsilon(n, m, p, r)
\alpha_{i_0 \ldots i_r} \otimes \beta_{i_r \ldots i_p}.
$$
where $\alpha$ has degree $n$ and $\beta$ has degree $m$
and with
$$
\epsilon(n, m, p, r) = (-1)^{(p + r)n + rp + r}.
$$
Note that $\epsilon(n, m, p, n) = 1$. Hence if
$\mathcal{F}^\bullet = \mathcal{F}[0]$ is the complex
consisting in a single abelian sheaf $\mathcal{F}$ placed in degree $0$,
then there no signs in the formula for $\cup$ (as in that case
$\alpha_{i_0 \ldots i_r} = 0$ unless $r = n$).
For an explanation of why there has to be a sign and how to
compute it see \cite[Exposee XVII]{SGA4} by Deligne.
To check (\ref{equation-needs-signs})
is a map of complexes we have to show that
$$
d(\alpha \cup \beta) =
d(\alpha) \cup \beta +
(-1)^{\deg(\alpha)} \alpha \cup d(\beta)
$$
by the definition of the differential on
$\text{Tot}(
\text{Tot}(\check{\mathcal{C}}^\bullet({\mathcal U}, {\mathcal F}^\bullet))
\otimes_{\mathbf Z}
\text{Tot}(\check{\mathcal{C}}^\bullet({\mathcal U}, {\mathcal G}^\bullet))
)$
as given in
Homology, Definition \ref{homology-definition-associated-simple-complex}.
We compute first
\begin{align*}
d(\alpha \cup \beta)_{i_0 \ldots i_{p + 1}}
= &
\sum\nolimits_{j = 0}^{p + 1}
(-1)^j
(\alpha \cup \beta)_{i_0 \ldots \hat i_j \ldots i_{p + 1}}
+
(-1)^{p + 1}
d_{{\mathcal F} \otimes {\mathcal G}}
((\alpha \cup \beta)_{i_0 \ldots i_{p + 1}})
\\
= &
\sum\nolimits_{j = 0}^{p + 1}
\sum\nolimits_{r = 0}^{j - 1}
(-1)^j \epsilon(n, m, p, r)
\alpha_{i_0 \ldots i_r} \otimes \beta_{i_r \ldots \hat i_j \ldots i_{p + 1}}
+ \\
&
\sum\nolimits_{j = 0}^{p + 1}
\sum\nolimits_{r = j + 1}^{p + 1}
(-1)^j \epsilon(n, m, p, r - 1)
\alpha_{i_0 \ldots \hat i_j \ldots i_r} \otimes \beta_{i_r \ldots i_{p + 1}}
+ \\
&
\sum\nolimits_{r = 0}^{p + 1}
(-1)^{p + 1} \epsilon(n, m, p + 1, r)
d_{{\mathcal F} \otimes {\mathcal G}}
(\alpha_{i_0 \ldots i_r} \otimes \beta_{i_r \ldots i_{p + 1}})
\end{align*}
and note that the summands in the last term equal
$$
(-1)^{p + 1} \epsilon(n, m, p + 1, r)
\left(
d_{\mathcal F}(\alpha_{i_0 \ldots i_r}) \otimes
\beta_{i_r \ldots i_{p + 1}} +
(-1)^{n - r}
\alpha_{i_0 \ldots i_r} \otimes d_{\mathcal G}(\beta_{i_r \ldots i_{p + 1}})
\right).
$$
because $\deg_\mathcal{F}(\alpha_{i_0 \ldots i_r}) = n - r$.
On the other hand
\begin{align*}
(d(\alpha) \cup \beta)_{i_0\ldots i_{p + 1}}
= &
\sum\nolimits_{r = 0}^{p + 1}
\epsilon(n + 1, m, p + 1, r)
d(\alpha)_{i_0\ldots i_r} \otimes \beta_{i_r\ldots i_{p + 1}}
\\
= &
\sum\nolimits_{r = 0}^{p + 1}
\sum\nolimits_{j = 0}^{r}
\epsilon(n + 1, m, p + 1, r) (-1)^j
\alpha_{i_0\ldots\hat{i_j}\ldots i_r} \otimes \beta_{i_r\ldots i_{p + 1}}
+ \\
&
\sum\nolimits_{r = 0}^{p + 1}
\epsilon(n + 1, m, p + 1, r) (-1)^r
d_{\mathcal F}(\alpha_{i_0 \ldots i_r}) \otimes \beta_{i_r\ldots i_{p + 1}}
\end{align*}
and
\begin{align*}
(\alpha \cup d(\beta))_{i_0\ldots i_{p + 1}}
= &
\sum\nolimits_{r = 0}^{p + 1}
\epsilon(n, m + 1, p + 1, r)
\alpha_{i_0 \ldots i_r} \otimes d(\beta)_{i_r \ldots i_{p + 1}}
\\
= &
\sum\nolimits_{r = 0}^{p + 1}
\sum\nolimits_{j = r}^{p + 1}
\epsilon(n, m + 1, p + 1, r) (-1)^{j - r}
\alpha_{i_0 \ldots i_r} \otimes \beta_{i_r \ldots \hat{i_j}\ldots i_{p + 1}}
+ \\
&
\sum\nolimits_{r = 0}^{p + 1}
\epsilon(n, m + 1, p + 1, r) (-1)^{p + 1 - r}
\alpha_{i_0 \ldots i_r} \otimes d_{\mathcal G}(\beta_{i_r \ldots i_{p + 1}})
\end{align*}
The desired equality holds if we have
\begin{align*}
(-1)^{p + 1} \epsilon(n, m, p + 1, r)
& =
\epsilon(n + 1, m, p + 1, r) (-1)^r \\
(-1)^{p + 1} \epsilon(n, m, p + 1, r) (-1)^{n - r}
& =
(-1)^n \epsilon(n, m + 1, p + 1, r) (-1)^{p + 1 - r} \\
\epsilon(n + 1, m, p + 1, r) (-1)^r
& =
(-1)^{1 + n} \epsilon(n, m + 1, p + 1, r - 1) \\
(-1)^j \epsilon(n, m, p, r)
& =
(-1)^n \epsilon(n, m + 1, p + 1, r) (-1)^{j - r} \\
(-1)^j \epsilon(n, m, p, r - 1)
& =
\epsilon(n + 1, m, p + 1, r) (-1)^j
\end{align*}
(The third equality is necessary to get the terms with $r = j$
from $d(\alpha) \cup \beta$ and $(-1)^n \alpha \cup d(\beta)$
to cancel each other.) We leave the verifications to the reader.
(Alternatively, check the script signs.gp in the scripts subdirectory
of the Stacks project.)

\medskip\noindent
Associativity of the cup product. Suppose that ${\mathcal F}^\bullet$,
${\mathcal G}^\bullet$ and ${\mathcal H}^\bullet$ are bounded below
complexes of abelian groups on $X$. The obvious map
(without the intervention of signs) is an isomorphism
of complexes
$$
\text{Tot}(
\text{Tot}({\mathcal F}^\bullet \otimes_{\mathbf Z} {\mathcal G}^\bullet)
\otimes_{\mathbf Z} {\mathcal H}^\bullet
)
\longrightarrow
\text{Tot}(
{\mathcal F}^\bullet \otimes_{\mathbf Z}
\text{Tot}({\mathcal G}^\bullet \otimes_{\mathbf Z} {\mathcal H}^\bullet)
).
$$
Another way to say this is that the triple complex
${\mathcal F}^\bullet \otimes_{\mathbf Z} {\mathcal G}^\bullet
\otimes_{\mathbf Z} {\mathcal H}^\bullet$ gives rise to a well defined
total complex with differential satisfying
$$
d(\alpha \otimes \beta \otimes \gamma) =
d(\alpha) \otimes \beta \otimes \gamma +
(-1)^{\deg(\alpha)} \alpha \otimes d(\beta) \otimes \gamma +
(-1)^{\deg(\alpha) + \deg(\beta)} \alpha \otimes \beta \otimes d(\gamma)
$$
for homogeneous elements. Using this map it is easy to verify that
$$
(\alpha \cup \beta) \cup \gamma = \alpha \cup ( \beta \cup \gamma)
$$
namely, if $\alpha$ has degree $a$, $\beta$ has degree $b$ and
$\gamma$ has degree $c$, then
\begin{align*}
((\alpha \cup \beta) \cup \gamma)_{i_0 \ldots i_p}
= &
\sum\nolimits_{r = 0}^p
\epsilon(a + b, c, p, r)
(\alpha \cup \beta)_{i_0 \ldots i_r} \otimes \gamma_{i_r \ldots i_p}
\\
= &
\sum\nolimits_{r = 0}^p
\sum\nolimits_{s = 0}^r
\epsilon(a + b, c, p, r) \epsilon(a, b, r, s)
\alpha_{i_0 \ldots i_s} \otimes
\beta_{i_s \ldots i_r} \otimes
\gamma_{i_r \ldots i_p}
\end{align*}
and
\begin{align*}
(\alpha \cup (\beta \cup \gamma)_{i_0\ldots i_p}
= &
\sum\nolimits_{s = 0}^p
\epsilon(a, b + c, p, s)
\alpha_{i_0 \ldots i_s} \otimes (\beta \cup \gamma)_{i_s \ldots i_p}
\\
= &
\sum\nolimits_{s = 0}^p
\sum\nolimits_{r = s}^p
\epsilon(a, b + c, p, s) \epsilon(b, c, p - s, r - s)
\alpha_{i_0 \ldots i_s} \otimes \beta_{i_s \ldots i_r} \otimes
\gamma_{i_r \ldots i_p}
\end{align*}
and a trivial mod $2$ calculation shows the signs match up.
(Alternatively, check the script signs.gp in the scripts subdirectory
of the Stacks project.)

\medskip\noindent
Finally, we indicate why the cup product preserves a graded commutative
structure, at least on a cohomological level. For this we use the operator
$\tau$ introduced above. Let ${\mathcal F}^\bullet$ be a bounded below
complexes of abelian groups, and assume we are given a graded commutative
multiplication
$$
\wedge^\bullet :
\text{Tot}({\mathcal F}^\bullet\otimes {\mathcal F}^\bullet)
\longrightarrow
{\mathcal F}^\bullet.
$$
This means the following: For $s$ a local section of
${\mathcal F}^a$, and $t$ a local section of ${\mathcal F}^b$
we have $s \wedge t$ a local section of ${\mathcal F}^{a + b}$.
Graded commutative means we have
$s \wedge t = (-1)^{ab} t \wedge s$. Since $\wedge$ is a map
of complexes we have
$d(s\wedge t) = d(s) \wedge t + (-1)^a s \wedge d(t)$.
The composition
$$
\text{Tot}(
\text{Tot}(\check{\mathcal{C}}^\bullet({\mathcal U}, {\mathcal F}^\bullet))
\otimes
\text{Tot}(\check{\mathcal{C}}^\bullet({\mathcal U}, {\mathcal F}^\bullet))
)
\to
\text{Tot}(
\check{\mathcal{C}}^\bullet({\mathcal U},
\text{Tot}({\mathcal F}^\bullet\otimes_{\mathbf Z}{\mathcal F}^\bullet))
)
\to
\text{Tot}(\check{\mathcal{C}}^\bullet({\mathcal U}, {\mathcal F}^\bullet))
$$
induces a cup product on cohomology
$$
H^n(
\text{Tot}(\check{\mathcal{C}}^\bullet({\mathcal U}, {\mathcal F}^\bullet))
)
\times
H^m(
\text{Tot}(\check{\mathcal{C}}^\bullet({\mathcal U}, {\mathcal F}^\bullet))
)
\longrightarrow
H^{n + m}(
\text{Tot}(\check{\mathcal{C}}^\bullet({\mathcal U}, {\mathcal F}^\bullet))
)
$$
and so in the colimit also a product on {\v C}ech cohomology
and therefore (using hypercoverings if needed) a product
in cohomology of ${\mathcal F}^\bullet$. We claim this product
(on cohomology) is graded commutative as well. To prove this
we first consider an element $\alpha$ of degree $n$ in
$\text{Tot}(\check{\mathcal{C}}^\bullet({\mathcal U}, {\mathcal F}^\bullet))$
and an element $\beta$ of degree $m$ in
$\text{Tot}(\check{\mathcal{C}}^\bullet({\mathcal U}, {\mathcal F}^\bullet))$
and we compute
\begin{align*}
\wedge^\bullet(\alpha \cup \beta)_{i_0 \ldots i_p}
= &
\sum\nolimits_{r = 0}^p
\epsilon(n, m, p, r)
\alpha_{i_0 \ldots i_r} \wedge \beta_{i_r \ldots i_p} \\
= &
\sum\nolimits_{r = 0}^p
\epsilon(n, m, p, r)
(-1)^{\deg(\alpha_{i_0 \ldots i_r})\deg(\beta_{i_r \ldots i_p})}
\beta_{i_r \ldots i_p} \wedge \alpha_{i_0 \ldots i_r}
\end{align*}
because $\wedge$ is graded commutative. On the other hand we have
\begin{align*}
\tau(\wedge^\bullet(\tau(\beta) \cup \tau(\alpha)))_{i_0 \ldots i_p}
= &
\chi(p)
\sum\nolimits_{r = 0}^p
\epsilon(m, n, p, r)
\tau(\beta)_{i_p \ldots i_{p - r}} \wedge \tau(\alpha)_{i_{p - r} \ldots i_0}
\\
= &
\chi(p)
\sum\nolimits_{r = 0}^p
\epsilon(m, n, p, r) \chi(r) \chi(p - r)
\beta_{i_{p - r} \ldots i_p} \wedge \alpha_{i_0 \ldots i_{p - r}}
\\
= &
\chi(p)
\sum\nolimits_{r = 0}^p
\epsilon(m, n, p, p - r) \chi(r) \chi(p - r)
\beta_{i_r \ldots i_p} \wedge \alpha_{i_0 \ldots i_r}
\end{align*}
where $\chi(t) = (-1)^{\frac{t(t + 1)}{2}}$. Since we proved earlier that
$\tau$ acts as the identity on cohomology we have to verify that
$$
\epsilon(n, m, p, r)
(-1)^{(n - r)(m - (p - r))}
=
(-1)^{nm} \chi(p)\epsilon(m, n, p, p - r) \chi(r) \chi(p - r)
$$
A trivial mod $2$ calculation shows these signs match up.
(Alternatively, check the script signs.gp in the scripts subdirectory
of the Stacks project.)

\medskip\noindent
Finally, we study the compatibility of cup product with boundary maps.
Suppose that
$$
0
\to
{\mathcal F}_1^\bullet
\to
{\mathcal F}_2^\bullet
\to
{\mathcal F}_3^\bullet
\to
0
\quad\text{and}\quad
0
\leftarrow
{\mathcal G}_1^\bullet
\leftarrow
{\mathcal G}_2^\bullet
\leftarrow
{\mathcal G}_3^\bullet
\leftarrow
0
$$
are short exact sequences of bounded below complexes of abelian
sheaves on $X$. Let ${\mathcal H}^\bullet$ be another bounded below
complex of abelian sheaves, and suppose we have maps of complexes
$$
\gamma_i :
\text{Tot}({\mathcal F}_i^\bullet \otimes_{\mathbf Z} {\mathcal G}_i^\bullet)
\longrightarrow
{\mathcal H}^\bullet
$$
which are compatible with the maps between the complexes, namely such that
the diagrams
$$
\xymatrix{
\text{Tot}({\mathcal F}_1^\bullet \otimes_{\mathbf Z} {\mathcal G}_1^\bullet)
\ar[d]_{\gamma_1}
&
\text{Tot}({\mathcal F}_1^\bullet \otimes_{\mathbf Z} {\mathcal G}_2^\bullet)
\ar[l] \ar[d]
\\
\mathcal{H}^\bullet &
\text{Tot}({\mathcal F}_2^\bullet \otimes_{\mathbf Z} {\mathcal G}_2^\bullet)
\ar[l]_-{\gamma_2}
}
$$
and
$$
\xymatrix{
\text{Tot}({\mathcal F}_2^\bullet \otimes_{\mathbf Z} {\mathcal G}_2^\bullet)
\ar[d]_{\gamma_2}
&
\text{Tot}({\mathcal F}_2^\bullet \otimes_{\mathbf Z} {\mathcal G}_3^\bullet)
\ar[l] \ar[d]
\\
\mathcal{H}^\bullet &
\text{Tot}({\mathcal F}_3^\bullet \otimes_{\mathbf Z} {\mathcal G}_3^\bullet)
\ar[l]_-{\gamma_3}
}
$$
are commutative.

\begin{lemma}
\label{lemma-compute-sign-cup-product-boundaries}
In the situation above, assume {\v C}ech cohomology agrees with cohomology
for the sheaves $\mathcal{F}_i^p$ and $\mathcal{G}_j^q$.
Let $a_3 \in H^n(X, \mathcal{F}_3^\bullet)$ and
$b_1 \in H^m(X, \mathcal{G}_1^\bullet)$. Then we have
$$
\gamma_1( \partial a_3 \cup b_1) =
(-1)^{n + 1} \gamma_3( a_3 \cup \partial b_1)
$$
in $H^{n + m + 1}(X, \mathcal{H}^\bullet)$ where $\partial$ indicates the
boundary map on cohomology associated to the short exact sequences of
complexes above.
\end{lemma}

\begin{proof}
We will use the following conventions and notation. We think of
${\mathcal F}_1^p$ as a subsheaf of ${\mathcal F}_2^p$ and we think of
${\mathcal G}_3^q$ as a subsheaf of ${\mathcal G}_2^q$. Hence if $s$ is
a local section of ${\mathcal F}_1^p$ we use $s$ to denote
the corresponding section of ${\mathcal F}_2^p$ as well. Similarly
for local sections of ${\mathcal G}_3^q$. Furthermore,
if $s$ is a local section of ${\mathcal F}_2^p$ then we denote
$\bar s$ its image in ${\mathcal F}_3^p$. Similarly for the
map ${\mathcal G}_2^q \to {\mathcal G}^q_1$. In particular if
$s$ is a local section of ${\mathcal F}_2^p$ and $\bar s = 0$
then $s$ is a local section of ${\mathcal F}_1^p$. The commutativity
of the diagrams above implies, for local sections $s$ of
${\mathcal F}_2^p$ and $t$ of ${\mathcal G}_3^q$ that
$\gamma_2(s \otimes t) = \gamma_3(\bar s \otimes t)$ as sections of
${\mathcal H}^{p + q}$.

\medskip\noindent
Let ${\mathcal U} : X =  \bigcup_{i \in I} U_i$
be an open covering of $X$. Suppose that $\alpha_3$,
resp.\ $\beta_1$ is a degree $n$, resp.\ $m$ cocycle of
$\text{Tot}(
\check{\mathcal{C}}^\bullet({\mathcal U}, {\mathcal F}_3^\bullet))$,
resp.\ $\text{Tot}(
\check{\mathcal{C}}^\bullet({\mathcal U}, {\mathcal G}_1^\bullet))$
representing $a_3$, resp.\ $b_1$. After refining $\mathcal{U}$ if necessary,
we can find cochains $\alpha_2$, resp.\ $\beta_2$ of
degree $n$, resp.\ $m$ in
$\text{Tot}(
\check{\mathcal{C}}^\bullet({\mathcal U}, {\mathcal F}_2^\bullet))$,
resp.\ $\text{Tot}(
\check{\mathcal{C}}^\bullet({\mathcal U}, {\mathcal G}_2^\bullet))$
mapping to $\alpha_3$, resp.\ $\beta_1$.
Then we see that
$$
\overline{d(\alpha_2)} = d(\bar \alpha_2) = 0
\quad\text{and}\quad
\overline{d(\beta_2)} = d(\bar \beta_2) = 0.
$$
This means that $\alpha_1 = d(\alpha_2)$ is a degree $n + 1$ cocycle in
$\text{Tot}(\check{\mathcal{C}}^\bullet({\mathcal U}, {\mathcal F}_1^\bullet))$
representing $\partial a_3$. Similarly, $\beta_3 = d(\beta_2)$ is
a degree $m + 1$ cocycle in
$\text{Tot}(\check{\mathcal{C}}^\bullet({\mathcal U}, {\mathcal G}_3^\bullet))$
representing $\partial b_1$.
Thus we may compute
\begin{align*}
d(\gamma_2(\alpha_2 \cup \beta_2))
& =
\gamma_2(d(\alpha_2 \cup \beta_2))
\\
& =
\gamma_2(d(\alpha_2) \cup \beta_2 + (-1)^n \alpha_2 \cup d(\beta_2) )
\\
& =
\gamma_2( \alpha_1 \cup \beta_2)  + (-1)^n \gamma_2( \alpha_2 \cup \beta_3)
\\
& =
\gamma_1(\alpha_1 \cup \beta_1) + (-1)^n \gamma_3(\alpha_3 \cup \beta_3)
\end{align*}
So this even tells us that the sign is $(-1)^{n + 1}$ as indicated
in the lemma\footnote{The sign depends on the convention for the
signs in the long exact sequence in cohomology associated to a triangle
in $D(X)$. The conventions in the Stacks project are (a) distinguished
triangles correspond to termwise split exact sequences and (b) the boundary
maps in the long exact sequence are given by the maps in the snake lemma
without the intervention of signs. See
Derived Categories, Section \ref{derived-section-homotopy-triangulated}.}.
\end{proof}

\begin{lemma}
\label{lemma-boundary-derivation-over-cup-product}
Let $X$ be a topological space. Let $\mathcal{O}' \to \mathcal{O}$ be a
surjection of sheaves of rings whose kernel $\mathcal{I} \subset \mathcal{O}'$
has square zero. Then $M = H^1(X, \mathcal{I})$ is a
$R = H^0(X, \mathcal{O})$-module and the boundary map
$\partial : R \to M$ associated to the short exact sequence
$$
0 \to \mathcal{I} \to \mathcal{O}' \to \mathcal{O} \to 0
$$
is a derivation (Algebra, Definition \ref{algebra-definition-derivation}).
\end{lemma}

\begin{proof}
The map $\mathcal{O}' \to \SheafHom(\mathcal{I}, \mathcal{I})$
factors through $\mathcal{O}$ as $\mathcal{I} \cdot \mathcal{I} = 0$
by assumption. Hence $\mathcal{I}$ is a sheaf of $\mathcal{O}$-modules
and this defines the $R$-module structure on $M$.
The boundary map is additive hence it suffices to prove
the Leibniz rule. Let $f \in R$. Choose an open covering
$\mathcal{U} : X = \bigcup U_i$ such that there exist
$f_i \in \mathcal{O}'(U_i)$ lifting $f|_{U_i} \in \mathcal{O}(U_i)$.
Observe that $f_i - f_j$ is an element of $\mathcal{I}(U_i \cap U_j)$.
Then $\partial(f)$ corresponds to the {\v C}ech cohomology class of
the $1$-cocycle $\alpha$ with $\alpha_{i_0i_1} = f_{i_0} - f_{i_1}$.
(Observe that by Lemma \ref{lemma-cech-h1} the first {\v C}ech cohomology
group with respect to $\mathcal{U}$ is a submodule of $M$.)
Next, let $g \in R$ be a second element and assume (after possibly
refining the open covering) that $g_i \in \mathcal{O}'(U_i)$ lifts
$g|_{U_i} \in \mathcal{O}(U_i)$. Then we see that
$\partial(g)$ is given by the cocycle $\beta$ with
$\beta_{i_0i_1} = g_{i_0} - g_{i_1}$. Since $f_ig_i \in \mathcal{O}'(U_i)$
lifts $fg|_{U_i}$ we see that
$\partial(fg)$ is given by the cocycle $\gamma$ with
$$
\gamma_{i_0i_1} = f_{i_0}g_{i_0} - f_{i_1}g_{i_1} =
(f_{i_0} - f_{i_1})g_{i_0} + f_{i_1}(g_{i_0} - g_{i_1}) =
\alpha_{i_0i_1}g + f\beta_{i_0i_1}
$$
by our definition of the $\mathcal{O}$-module structure on $\mathcal{I}$.
This proves the Leibniz rule and the proof is complete.
\end{proof}









\section{Flat resolutions}
\label{section-flat}

\noindent
A reference for the material in this section is \cite{Spaltenstein}.
Let $(X, \mathcal{O}_X)$ be a ringed space. By
Modules, Lemma \ref{modules-lemma-module-quotient-flat}
any $\mathcal{O}_X$-module is a quotient of a flat $\mathcal{O}_X$-module.
By
Derived Categories, Lemma \ref{derived-lemma-subcategory-left-resolution}
any bounded above complex of $\mathcal{O}_X$-modules has a left
resolution by a bounded above complex of flat $\mathcal{O}_X$-modules.
However, for unbounded complexes, it turns out that flat resolutions
aren't good enough.

\begin{lemma}
\label{lemma-derived-tor-exact}
Let $(X, \mathcal{O}_X)$ be a ringed space.
Let $\mathcal{G}^\bullet$ be a complex of $\mathcal{O}_X$-modules.
The functors
$$
K(\textit{Mod}(\mathcal{O}_X))
\longrightarrow
K(\textit{Mod}(\mathcal{O}_X)),
\quad
\mathcal{F}^\bullet \longmapsto
\text{Tot}(\mathcal{G}^\bullet \otimes_{\mathcal{O}_X} \mathcal{F}^\bullet)
$$
and
$$
K(\textit{Mod}(\mathcal{O}_X))
\longrightarrow
K(\textit{Mod}(\mathcal{O}_X)),
\quad
\mathcal{F}^\bullet \longmapsto
\text{Tot}(\mathcal{F}^\bullet \otimes_{\mathcal{O}_X} \mathcal{G}^\bullet)
$$
are exact functors of triangulated categories.
\end{lemma}

\begin{proof}
This follows from Derived Categories, Remark
\ref{derived-remark-double-complex-as-tensor-product-of}.
\end{proof}

\begin{definition}
\label{definition-K-flat}
Let $(X, \mathcal{O}_X)$ be a ringed space.
A complex $\mathcal{K}^\bullet$ of $\mathcal{O}_X$-modules is
called {\it K-flat} if for every acyclic complex $\mathcal{F}^\bullet$
of $\mathcal{O}_X$-modules the complex
$$
\text{Tot}(\mathcal{F}^\bullet \otimes_{\mathcal{O}_X} \mathcal{K}^\bullet)
$$
is acyclic.
\end{definition}

\begin{lemma}
\label{lemma-K-flat-quasi-isomorphism}
Let $(X, \mathcal{O}_X)$ be a ringed space.
Let $\mathcal{K}^\bullet$ be a K-flat complex.
Then the functor
$$
K(\textit{Mod}(\mathcal{O}_X))
\longrightarrow
K(\textit{Mod}(\mathcal{O}_X)), \quad
\mathcal{F}^\bullet
\longmapsto
\text{Tot}(\mathcal{F}^\bullet \otimes_{\mathcal{O}_X} \mathcal{K}^\bullet)
$$
transforms quasi-isomorphisms into quasi-isomorphisms.
\end{lemma}

\begin{proof}
Follows from
Lemma \ref{lemma-derived-tor-exact}
and the fact that quasi-isomorphisms are characterized by having
acyclic cones.
\end{proof}

\begin{lemma}
\label{lemma-check-K-flat-stalks}
Let $(X, \mathcal{O}_X)$ be a ringed space. Let $\mathcal{K}^\bullet$
be a complex of $\mathcal{O}_X$-modules. Then $\mathcal{K}^\bullet$
is K-flat if and only if for all $x \in X$ the complex
$\mathcal{K}_x^\bullet$ of $\mathcal{O}_{X, x}$-modules is K-flat
(More on Algebra, Definition \ref{more-algebra-definition-K-flat}).
\end{lemma}

\begin{proof}
If $\mathcal{K}_x^\bullet$ is K-flat for all $x \in X$ then we see
that $\mathcal{K}^\bullet$ is K-flat because $\otimes$ and
direct sums commute with taking stalks and because we can check exactness
at stalks, see
Modules, Lemma \ref{modules-lemma-abelian}.
Conversely, assume $\mathcal{K}^\bullet$ is K-flat. Pick $x \in X$
and let
$M^\bullet$ be an acyclic complex of $\mathcal{O}_{X, x}$-modules.
Then $i_{x, *}M^\bullet$ is an acyclic complex of $\mathcal{O}_X$-modules.
Thus $\text{Tot}(i_{x, *}M^\bullet \otimes_{\mathcal{O}_X} \mathcal{K}^\bullet)$
is acyclic. Taking stalks at $x$ shows that
$\text{Tot}(M^\bullet \otimes_{\mathcal{O}_{X, x}} \mathcal{K}_x^\bullet)$
is acyclic.
\end{proof}

\begin{lemma}
\label{lemma-tensor-product-K-flat}
Let $(X, \mathcal{O}_X)$ be a ringed space.
If $\mathcal{K}^\bullet$, $\mathcal{L}^\bullet$ are K-flat complexes
of $\mathcal{O}_X$-modules, then
$\text{Tot}(\mathcal{K}^\bullet \otimes_{\mathcal{O}_X} \mathcal{L}^\bullet)$
is a K-flat complex of $\mathcal{O}_X$-modules.
\end{lemma}

\begin{proof}
Follows from the isomorphism
$$
\text{Tot}(\mathcal{M}^\bullet \otimes_{\mathcal{O}_X}
\text{Tot}(\mathcal{K}^\bullet \otimes_{\mathcal{O}_X} \mathcal{L}^\bullet))
=
\text{Tot}(\text{Tot}(\mathcal{M}^\bullet \otimes_{\mathcal{O}_X}
\mathcal{K}^\bullet) \otimes_{\mathcal{O}_X} \mathcal{L}^\bullet)
$$
and the definition.
\end{proof}

\begin{lemma}
\label{lemma-K-flat-two-out-of-three}
Let $(X, \mathcal{O}_X)$ be a ringed space.
Let $(\mathcal{K}_1^\bullet, \mathcal{K}_2^\bullet, \mathcal{K}_3^\bullet)$
be a distinguished triangle in $K(\textit{Mod}(\mathcal{O}_X))$.
If two out of three of $\mathcal{K}_i^\bullet$ are K-flat, so is the third.
\end{lemma}

\begin{proof}
Follows from
Lemma \ref{lemma-derived-tor-exact}
and the fact that in a distinguished triangle in
$K(\textit{Mod}(\mathcal{O}_X))$
if two out of three are acyclic, so is the third.
\end{proof}

\begin{lemma}
\label{lemma-K-flat-two-out-of-three-ses}
Let $(X, \mathcal{O}_X)$ be a ringed space. Let
$0 \to \mathcal{K}_1^\bullet \to \mathcal{K}_2^\bullet \to
\mathcal{K}_3^\bullet \to 0$ be a short exact sequence of complexes
such that the terms of $\mathcal{K}_3^\bullet$ are flat $\mathcal{O}_X$-modules.
If two out of three of $\mathcal{K}_i^\bullet$ are K-flat, so is the third.
\end{lemma}

\begin{proof}
By Modules, Lemma \ref{modules-lemma-flat-tor-zero}
for every complex $\mathcal{L}^\bullet$
we obtain a short exact sequence
$$
0 \to
\text{Tot}(\mathcal{L}^\bullet \otimes_{\mathcal{O}_X} \mathcal{K}_1^\bullet)
\to
\text{Tot}(\mathcal{L}^\bullet \otimes_{\mathcal{O}_X} \mathcal{K}_1^\bullet)
\to
\text{Tot}(\mathcal{L}^\bullet \otimes_{\mathcal{O}_X} \mathcal{K}_1^\bullet)
\to 0
$$
of complexes. Hence the lemma follows from the long exact sequence of
cohomology sheaves and the definition of K-flat complexes.
\end{proof}

\begin{lemma}
\label{lemma-pullback-K-flat}
Let $f : (X, \mathcal{O}_X) \to (Y, \mathcal{O}_Y)$ be a morphism of
ringed spaces. The pullback of a K-flat complex of $\mathcal{O}_Y$-modules
is a K-flat complex of $\mathcal{O}_X$-modules.
\end{lemma}

\begin{proof}
We can check this on stalks, see
Lemma \ref{lemma-check-K-flat-stalks}.
Hence this follows from
Sheaves, Lemma \ref{sheaves-lemma-stalk-pullback-modules}
and
More on Algebra, Lemma \ref{more-algebra-lemma-base-change-K-flat}.
\end{proof}

\begin{lemma}
\label{lemma-bounded-flat-K-flat}
Let $(X, \mathcal{O}_X)$ be a ringed space. A bounded above complex
of flat $\mathcal{O}_X$-modules is K-flat.
\end{lemma}

\begin{proof}
We can check this on stalks, see
Lemma \ref{lemma-check-K-flat-stalks}.
Thus this lemma follows from
Modules, Lemma \ref{modules-lemma-flat-stalks-flat}
and
More on Algebra, Lemma \ref{more-algebra-lemma-derived-tor-quasi-isomorphism}.
\end{proof}

\noindent
In the following lemma by a colimit of a system of complexes we mean
the termwise colimit.

\begin{lemma}
\label{lemma-colimit-K-flat}
Let $(X, \mathcal{O}_X)$ be a ringed space.
Let $\mathcal{K}_1^\bullet \to \mathcal{K}_2^\bullet \to \ldots$
be a system of K-flat complexes.
Then $\colim_i \mathcal{K}_i^\bullet$ is K-flat.
\end{lemma}

\begin{proof}
Because we are taking termwise colimits it is clear that
$$
\colim_i \text{Tot}(
\mathcal{F}^\bullet \otimes_{\mathcal{O}_X} \mathcal{K}_i^\bullet)
=
\text{Tot}(\mathcal{F}^\bullet \otimes_{\mathcal{O}_X}
\colim_i \mathcal{K}_i^\bullet)
$$
Hence the lemma follows from the fact that filtered colimits are
exact.
\end{proof}

\begin{lemma}
\label{lemma-resolution-by-direct-sums-extensions-by-zero}
Let $(X, \mathcal{O}_X)$ be a ringed space.
For any complex $\mathcal{G}^\bullet$ of $\mathcal{O}_X$-modules
there exists a commutative diagram of complexes of $\mathcal{O}_X$-modules
$$
\xymatrix{
\mathcal{K}_1^\bullet \ar[d] \ar[r] &
\mathcal{K}_2^\bullet \ar[d] \ar[r] & \ldots \\
\tau_{\leq 1}\mathcal{G}^\bullet \ar[r] &
\tau_{\leq 2}\mathcal{G}^\bullet \ar[r] & \ldots
}
$$
with the following properties: (1) the vertical arrows are quasi-isomorphisms
and termwise surjective,
(2) each $\mathcal{K}_n^\bullet$ is a bounded above complex whose terms
are direct sums of $\mathcal{O}_X$-modules of the form
$j_{U!}\mathcal{O}_U$, and
(3) the maps $\mathcal{K}_n^\bullet \to \mathcal{K}_{n + 1}^\bullet$ are
termwise split injections whose cokernels are direct sums of
$\mathcal{O}_X$-modules of the form $j_{U!}\mathcal{O}_U$. Moreover, the map
$\colim \mathcal{K}_n^\bullet \to \mathcal{G}^\bullet$ is a quasi-isomorphism.
\end{lemma}

\begin{proof}
The existence of the diagram and properties (1), (2), (3) follows immediately
from
Modules, Lemma \ref{modules-lemma-module-quotient-flat}
and
Derived Categories, Lemma \ref{derived-lemma-special-direct-system}.
The induced map
$\colim \mathcal{K}_n^\bullet \to \mathcal{G}^\bullet$
is a quasi-isomorphism because filtered colimits are exact.
\end{proof}

\begin{lemma}
\label{lemma-K-flat-resolution}
Let $(X, \mathcal{O}_X)$ be a ringed space.
For any complex $\mathcal{G}^\bullet$ there exists a $K$-flat complex
$\mathcal{K}^\bullet$ whose terms are flat $\mathcal{O}_X$-modules
and a quasi-isomorphism $\mathcal{K}^\bullet \to \mathcal{G}^\bullet$
which is termwise surjective.
\end{lemma}

\begin{proof}
Choose a diagram as in
Lemma \ref{lemma-resolution-by-direct-sums-extensions-by-zero}.
Each complex $\mathcal{K}_n^\bullet$ is a bounded
above complex of flat modules, see
Modules, Lemma \ref{modules-lemma-j-shriek-flat}.
Hence $\mathcal{K}_n^\bullet$ is K-flat by
Lemma \ref{lemma-bounded-flat-K-flat}.
Thus $\colim \mathcal{K}_n^\bullet$ is K-flat by
Lemma \ref{lemma-colimit-K-flat}.
The induced map
$\colim \mathcal{K}_n^\bullet \to \mathcal{G}^\bullet$
is a quasi-isomorphism and termwise surjective by construction.
Property (3) of Lemma \ref{lemma-resolution-by-direct-sums-extensions-by-zero}
shows that $\colim \mathcal{K}_n^m$ is a direct sum of
flat modules and hence flat which proves the final assertion.
\end{proof}

\begin{lemma}
\label{lemma-derived-tor-quasi-isomorphism-other-side}
Let $(X, \mathcal{O}_X)$ be a ringed space. Let
$\alpha : \mathcal{P}^\bullet \to \mathcal{Q}^\bullet$ be a
quasi-isomorphism of K-flat complexes of $\mathcal{O}_X$-modules.
For every complex $\mathcal{F}^\bullet$ of $\mathcal{O}_X$-modules
the induced map
$$
\text{Tot}(\text{id}_{\mathcal{F}^\bullet} \otimes \alpha) :
\text{Tot}(\mathcal{F}^\bullet \otimes_{\mathcal{O}_X} \mathcal{P}^\bullet)
\longrightarrow
\text{Tot}(\mathcal{F}^\bullet \otimes_{\mathcal{O}_X} \mathcal{Q}^\bullet)
$$
is a quasi-isomorphism.
\end{lemma}

\begin{proof}
Choose a quasi-isomorphism $\mathcal{K}^\bullet \to \mathcal{F}^\bullet$
with $\mathcal{K}^\bullet$ a K-flat complex, see
Lemma \ref{lemma-K-flat-resolution}.
Consider the commutative diagram
$$
\xymatrix{
\text{Tot}(\mathcal{K}^\bullet
\otimes_{\mathcal{O}_X} \mathcal{P}^\bullet) \ar[r] \ar[d] &
\text{Tot}(\mathcal{K}^\bullet
\otimes_{\mathcal{O}_X} \mathcal{Q}^\bullet) \ar[d] \\
\text{Tot}(\mathcal{F}^\bullet
\otimes_{\mathcal{O}_X} \mathcal{P}^\bullet) \ar[r] &
\text{Tot}(\mathcal{F}^\bullet
\otimes_{\mathcal{O}_X} \mathcal{Q}^\bullet)
}
$$
The result follows as by
Lemma \ref{lemma-K-flat-quasi-isomorphism}
the vertical arrows and the top horizontal arrow are quasi-isomorphisms.
\end{proof}

\noindent
Let $(X, \mathcal{O}_X)$ be a ringed space.
Let $\mathcal{F}^\bullet$ be an object of $D(\mathcal{O}_X)$.
Choose a K-flat resolution $\mathcal{K}^\bullet \to \mathcal{F}^\bullet$, see
Lemma \ref{lemma-K-flat-resolution}.
By
Lemma \ref{lemma-derived-tor-exact}
we obtain an exact functor of triangulated categories
$$
K(\mathcal{O}_X)
\longrightarrow
K(\mathcal{O}_X),
\quad
\mathcal{G}^\bullet
\longmapsto
\text{Tot}(\mathcal{G}^\bullet \otimes_{\mathcal{O}_X} \mathcal{K}^\bullet)
$$
By
Lemma \ref{lemma-K-flat-quasi-isomorphism}
this functor induces a functor
$D(\mathcal{O}_X) \to D(\mathcal{O}_X)$ simply because
$D(\mathcal{O}_X)$ is the localization of $K(\mathcal{O}_X)$
at quasi-isomorphisms. As the category of $K$-flat resolutions
of $\mathcal{F}^\bullet$ is cofiltered and as we have
Lemma \ref{lemma-derived-tor-quasi-isomorphism-other-side}
the resulting functor (up to isomorphism)
does not depend on the choice of $\mathcal{K}^\bullet$.

\begin{definition}
\label{definition-derived-tor}
Let $(X, \mathcal{O}_X)$ be a ringed space.
Let $\mathcal{F}^\bullet$ be an object of $D(\mathcal{O}_X)$.
The {\it derived tensor product}
$$
- \otimes_{\mathcal{O}_X}^{\mathbf{L}} \mathcal{F}^\bullet :
D(\mathcal{O}_X)
\longrightarrow
D(\mathcal{O}_X)
$$
is the exact functor of triangulated categories described above.
\end{definition}

\noindent
It is clear from our explicit constructions that
there is a canonical isomorphism
$$
\mathcal{F}^\bullet \otimes_{\mathcal{O}_X}^{\mathbf{L}} \mathcal{G}^\bullet
\cong
\mathcal{G}^\bullet \otimes_{\mathcal{O}_X}^{\mathbf{L}} \mathcal{F}^\bullet
$$
for $\mathcal{G}^\bullet$ and $\mathcal{F}^\bullet$ in $D(\mathcal{O}_X)$.
Here we use sign rules as given in More on Algebra, Section
\ref{more-algebra-section-sign-rules}. Hence when we write
$\mathcal{F}^\bullet \otimes_{\mathcal{O}_X}^{\mathbf{L}} \mathcal{G}^\bullet$
we will usually be agnostic about which variable we are using to
define the derived tensor product with.

\begin{definition}
\label{definition-tor}
Let $(X, \mathcal{O}_X)$ be a ringed space.
Let $\mathcal{F}$, $\mathcal{G}$ be $\mathcal{O}_X$-modules.
The {\it Tor}'s of $\mathcal{F}$ and $\mathcal{G}$ are define by
the formula
$$
\text{Tor}_p^{\mathcal{O}_X}(\mathcal{F}, \mathcal{G}) =
H^{-p}(\mathcal{F} \otimes_{\mathcal{O}_X}^\mathbf{L} \mathcal{G})
$$
with derived tensor product as defined above.
\end{definition}

\noindent
This definition implies that for every short exact sequence
of $\mathcal{O}_X$-modules
$0 \to \mathcal{F}_1 \to \mathcal{F}_2 \to \mathcal{F}_3 \to 0$
we have a long exact cohomology sequence
$$
\xymatrix{
\mathcal{F}_1 \otimes_{\mathcal{O}_X} \mathcal{G} \ar[r] &
\mathcal{F}_2 \otimes_{\mathcal{O}_X} \mathcal{G} \ar[r] &
\mathcal{F}_3 \otimes_{\mathcal{O}_X} \mathcal{G} \ar[r] & 0 \\
\text{Tor}_1^{\mathcal{O}_X}(\mathcal{F}_1, \mathcal{G}) \ar[r] &
\text{Tor}_1^{\mathcal{O}_X}(\mathcal{F}_2, \mathcal{G}) \ar[r] &
\text{Tor}_1^{\mathcal{O}_X}(\mathcal{F}_3, \mathcal{G}) \ar[ull]
}
$$
for every $\mathcal{O}_X$-module $\mathcal{G}$. This will be called
the long exact sequence of $\text{Tor}$ associated to the situation.

\begin{lemma}
\label{lemma-flat-tor-zero}
\begin{slogan}
Tor measures the deviation of flatness.
\end{slogan}
Let $(X, \mathcal{O}_X)$ be a ringed space.
Let $\mathcal{F}$ be an $\mathcal{O}_X$-module.
The following are equivalent
\begin{enumerate}
\item $\mathcal{F}$ is a flat $\mathcal{O}_X$-module, and
\item $\text{Tor}_1^{\mathcal{O}_X}(\mathcal{F}, \mathcal{G}) = 0$
for every $\mathcal{O}_X$-module $\mathcal{G}$.
\end{enumerate}
\end{lemma}

\begin{proof}
If $\mathcal{F}$ is flat, then $\mathcal{F} \otimes_{\mathcal{O}_X} -$
is an exact functor and the satellites vanish. Conversely assume (2)
holds. Then if $\mathcal{G} \to \mathcal{H}$ is injective with cokernel
$\mathcal{Q}$, the long exact sequence of $\text{Tor}$ shows that
the kernel of
$\mathcal{F} \otimes_{\mathcal{O}_X} \mathcal{G} \to
\mathcal{F} \otimes_{\mathcal{O}_X} \mathcal{H}$
is a quotient of
$\text{Tor}_1^{\mathcal{O}_X}(\mathcal{F}, \mathcal{Q})$
which is zero by assumption. Hence $\mathcal{F}$ is flat.
\end{proof}

\begin{lemma}
\label{lemma-factor-through-K-flat}
Let $(X, \mathcal{O}_X)$ be a ringed space.
Let $a : \mathcal{K}^\bullet \to \mathcal{L}^\bullet$ be a map of complexes
of $\mathcal{O}_X$-modules. If $\mathcal{K}^\bullet$ is K-flat, then
there exist a complex $\mathcal{N}^\bullet$ and maps of complexes
$b : \mathcal{K}^\bullet \to \mathcal{N}^\bullet$
and $c : \mathcal{N}^\bullet \to \mathcal{L}^\bullet$ such that
\begin{enumerate}
\item $\mathcal{N}^\bullet$ is K-flat,
\item $c$ is a quasi-isomorphism,
\item $a$ is homotopic to $c \circ b$.
\end{enumerate}
If the terms of $\mathcal{K}^\bullet$ are flat, then we may choose
$\mathcal{N}^\bullet$, $b$, and $c$
such that the same is true for $\mathcal{N}^\bullet$.
\end{lemma}

\begin{proof}
We will use that the homotopy category $K(\textit{Mod}(\mathcal{O}_X))$
is a triangulated category, see Derived Categories, Proposition
\ref{derived-proposition-homotopy-category-triangulated}.
Choose a distinguished triangle
$\mathcal{K}^\bullet \to \mathcal{L}^\bullet \to
\mathcal{C}^\bullet \to \mathcal{K}^\bullet[1]$.
Choose a quasi-isomorphism $\mathcal{M}^\bullet \to \mathcal{C}^\bullet$ with
$\mathcal{M}^\bullet$ K-flat with flat terms, see
Lemma \ref{lemma-K-flat-resolution}.
By the axioms of triangulated categories,
we may fit the composition
$\mathcal{M}^\bullet \to \mathcal{C}^\bullet \to \mathcal{K}^\bullet[1]$
into a distinguished triangle
$\mathcal{K}^\bullet \to \mathcal{N}^\bullet \to
\mathcal{M}^\bullet \to \mathcal{K}^\bullet[1]$.
By Lemma \ref{lemma-K-flat-two-out-of-three} we see that
$\mathcal{N}^\bullet$ is K-flat.
Again using the axioms of triangulated categories,
we can choose a map $\mathcal{N}^\bullet \to \mathcal{L}^\bullet$ fitting into
the following morphism of distinguished triangles
$$
\xymatrix{
\mathcal{K}^\bullet \ar[r] \ar[d] &
\mathcal{N}^\bullet \ar[r] \ar[d] &
\mathcal{M}^\bullet \ar[r] \ar[d] &
\mathcal{K}^\bullet[1] \ar[d] \\
\mathcal{K}^\bullet \ar[r] &
\mathcal{L}^\bullet \ar[r] &
\mathcal{C}^\bullet \ar[r] &
\mathcal{K}^\bullet[1]
}
$$
Since two out of three of the arrows are quasi-isomorphisms, so is
the third arrow $\mathcal{N}^\bullet \to \mathcal{L}^\bullet$
by the long exact sequences
of cohomology associated to these distinguished triangles
(or you can look at the image of this diagram in $D(\mathcal{O}_X)$ and use
Derived Categories, Lemma \ref{derived-lemma-third-isomorphism-triangle}
if you like). This finishes the proof of (1), (2), and (3).
To prove the final assertion, we may choose $\mathcal{N}^\bullet$
such that $\mathcal{N}^n \cong \mathcal{M}^n \oplus \mathcal{K}^n$, see
Derived Categories, Lemma
\ref{derived-lemma-improve-distinguished-triangle-homotopy}.
Hence we get the desired flatness
if the terms of $\mathcal{K}^\bullet$ are flat.
\end{proof}











\section{Derived pullback}
\label{section-derived-pullback}

\noindent
Let $f : (X, \mathcal{O}_X) \to (Y, \mathcal{O}_Y)$
be a morphism of ringed spaces. We can use K-flat resolutions to define
a derived pullback functor
$$
Lf^* : D(\mathcal{O}_Y) \to D(\mathcal{O}_X)
$$
Namely, for every complex of $\mathcal{O}_Y$-modules $\mathcal{G}^\bullet$
we can choose a K-flat resolution
$\mathcal{K}^\bullet \to \mathcal{G}^\bullet$ and set
$Lf^*\mathcal{G}^\bullet = f^*\mathcal{K}^\bullet$.
You can use
Lemmas \ref{lemma-pullback-K-flat},
\ref{lemma-K-flat-resolution}, and
\ref{lemma-derived-tor-quasi-isomorphism-other-side}
to see that this is well defined. However, to cross all the t's and dot all
the i's it is perhaps more convenient to use some general theory.

\begin{lemma}
\label{lemma-derived-base-change}
The construction above is independent of choices and defines an exact
functor of triangulated categories
$Lf^* : D(\mathcal{O}_Y) \to D(\mathcal{O}_X)$.
\end{lemma}

\begin{proof}
To see this we use the general theory developed in
Derived Categories, Section \ref{derived-section-derived-functors}.
Set $\mathcal{D} = K(\mathcal{O}_Y)$ and $\mathcal{D}' = D(\mathcal{O}_X)$.
Let us write $F : \mathcal{D} \to \mathcal{D}'$ the exact functor
of triangulated categories defined by the rule
$F(\mathcal{G}^\bullet) = f^*\mathcal{G}^\bullet$.
We let $S$ be the set of quasi-isomorphisms in
$\mathcal{D} = K(\mathcal{O}_Y)$.
This gives a situation as in
Derived Categories, Situation \ref{derived-situation-derived-functor}
so that
Derived Categories, Definition
\ref{derived-definition-right-derived-functor-defined}
applies. We claim that $LF$ is everywhere defined.
This follows from
Derived Categories, Lemma \ref{derived-lemma-find-existence-computes}
with $\mathcal{P} \subset \Ob(\mathcal{D})$ the collection
of $K$-flat complexes: (1) follows from
Lemma \ref{lemma-K-flat-resolution}
and to see (2) we have to show that for a quasi-isomorphism
$\mathcal{K}_1^\bullet  \to \mathcal{K}_2^\bullet$ between
K-flat complexes of $\mathcal{O}_Y$-modules the map
$f^*\mathcal{K}_1^\bullet  \to f^*\mathcal{K}_2^\bullet$ is a
quasi-isomorphism. To see this write this as
$$
f^{-1}\mathcal{K}_1^\bullet \otimes_{f^{-1}\mathcal{O}_Y} \mathcal{O}_X
\longrightarrow
f^{-1}\mathcal{K}_2^\bullet \otimes_{f^{-1}\mathcal{O}_Y} \mathcal{O}_X
$$
The functor $f^{-1}$ is exact, hence the map
$f^{-1}\mathcal{K}_1^\bullet  \to f^{-1}\mathcal{K}_2^\bullet$ is a
quasi-isomorphism. By
Lemma \ref{lemma-pullback-K-flat}
applied to the morphism $(X, f^{-1}\mathcal{O}_Y) \to (Y, \mathcal{O}_Y)$
the complexes $f^{-1}\mathcal{K}_1^\bullet$ and $f^{-1}\mathcal{K}_2^\bullet$
are K-flat complexes of $f^{-1}\mathcal{O}_Y$-modules. Hence
Lemma \ref{lemma-derived-tor-quasi-isomorphism-other-side}
guarantees that the displayed map is a quasi-isomorphism.
Thus we obtain a derived functor
$$
LF :
D(\mathcal{O}_Y) = S^{-1}\mathcal{D}
\longrightarrow
\mathcal{D}' = D(\mathcal{O}_X)
$$
see
Derived Categories, Equation (\ref{derived-equation-everywhere}).
Finally,
Derived Categories, Lemma \ref{derived-lemma-find-existence-computes}
also guarantees that
$LF(\mathcal{K}^\bullet) = F(\mathcal{K}^\bullet) = f^*\mathcal{K}^\bullet$
when $\mathcal{K}^\bullet$ is K-flat, i.e., $Lf^* = LF$ is
indeed computed in the way described above.
\end{proof}

\begin{lemma}
\label{lemma-derived-pullback-composition}
Let $f : X \to Y$ and $g : Y \to Z$ be morphisms of ringed spaces.
Then $Lf^* \circ Lg^* = L(g \circ f)^*$ as functors
$D(\mathcal{O}_Z) \to D(\mathcal{O}_X)$.
\end{lemma}

\begin{proof}
Let $E$ be an object of $D(\mathcal{O}_Z)$.
By construction $Lg^*E$ is computed by choosing a K-flat complex
$\mathcal{K}^\bullet$ representing $E$ on $Z$ and
setting $Lg^*E = g^*\mathcal{K}^\bullet$.
By Lemma \ref{lemma-pullback-K-flat} we see that $g^*\mathcal{K}^\bullet$
is K-flat on $Y$. Then $Lf^*Lg^*E$ is given by
$f^*g^*\mathcal{K}^\bullet = (g \circ f)^*\mathcal{K}^\bullet$
which also represents $L(g \circ f)^*E$.
\end{proof}

\begin{lemma}
\label{lemma-pullback-tensor-product}
Let $f : (X, \mathcal{O}_X) \to (Y, \mathcal{O}_Y)$
be a morphism of ringed spaces. There is a canonical bifunctorial
isomorphism
$$
Lf^*(
\mathcal{F}^\bullet \otimes_{\mathcal{O}_Y}^{\mathbf{L}} \mathcal{G}^\bullet
) =
Lf^*\mathcal{F}^\bullet
\otimes_{\mathcal{O}_X}^{\mathbf{L}}
Lf^*\mathcal{G}^\bullet
$$
for $\mathcal{F}^\bullet, \mathcal{G}^\bullet \in \Ob(D(\mathcal{O}_Y))$.
\end{lemma}

\begin{proof}
We may assume that $\mathcal{F}^\bullet$ and $\mathcal{G}^\bullet$
are K-flat complexes. In this case
$\mathcal{F}^\bullet \otimes_{\mathcal{O}_Y}^{\mathbf{L}} \mathcal{G}^\bullet$
is just the total complex associated to the double complex
$\mathcal{F}^\bullet \otimes_{\mathcal{O}_Y} \mathcal{G}^\bullet$.
By
Lemma \ref{lemma-tensor-product-K-flat}
$\text{Tot}(\mathcal{F}^\bullet \otimes_{\mathcal{O}_Y} \mathcal{G}^\bullet)$
is K-flat also. Hence the isomorphism of the lemma comes from the
isomorphism
$$
\text{Tot}(f^*\mathcal{F}^\bullet \otimes_{\mathcal{O}_X}
f^*\mathcal{G}^\bullet)
\longrightarrow
f^*\text{Tot}(\mathcal{F}^\bullet \otimes_{\mathcal{O}_Y} \mathcal{G}^\bullet)
$$
whose constituents are the isomorphisms
$f^*\mathcal{F}^p \otimes_{\mathcal{O}_X} f^*\mathcal{G}^q \to
f^*(\mathcal{F}^p \otimes_{\mathcal{O}_Y} \mathcal{G}^q)$ of
Modules, Lemma \ref{modules-lemma-tensor-product-pullback}.
\end{proof}

\begin{lemma}
\label{lemma-variant-derived-pullback}
Let $f : (X, \mathcal{O}_X) \to (Y, \mathcal{O}_Y)$
be a morphism of ringed spaces. There is a canonical bifunctorial
isomorphism
$$
\mathcal{F}^\bullet
\otimes_{\mathcal{O}_X}^{\mathbf{L}}
Lf^*\mathcal{G}^\bullet
=
\mathcal{F}^\bullet
\otimes_{f^{-1}\mathcal{O}_Y}^{\mathbf{L}}
f^{-1}\mathcal{G}^\bullet
$$
for $\mathcal{F}^\bullet$ in $D(\mathcal{O}_X)$ and
$\mathcal{G}^\bullet$ in $D(\mathcal{O}_Y)$.
\end{lemma}

\begin{proof}
Let $\mathcal{F}$ be an $\mathcal{O}_X$-module and let $\mathcal{G}$
be an $\mathcal{O}_Y$-module. Then
$\mathcal{F} \otimes_{\mathcal{O}_X} f^*\mathcal{G} =
\mathcal{F} \otimes_{f^{-1}\mathcal{O}_Y} f^{-1}\mathcal{G}$
because
$f^*\mathcal{G} =
\mathcal{O}_X \otimes_{f^{-1}\mathcal{O}_Y} f^{-1}\mathcal{G}$.
The lemma follows from this and the definitions.
\end{proof}

\begin{lemma}
\label{lemma-tensor-pull-compatibility}
Let $f : (X, \mathcal{O}_X) \to (Y, \mathcal{O}_Y)$ be a morphism
of ringed spaces. Let $\mathcal{K}^\bullet$ and $\mathcal{M}^\bullet$
be complexes of $\mathcal{O}_Y$-modules. The diagram
$$
\xymatrix{
Lf^*(\mathcal{K}^\bullet
\otimes_{\mathcal{O}_Y}^\mathbf{L}
\mathcal{M}^\bullet) \ar[r] \ar[d] &
Lf^*\text{Tot}(\mathcal{K}^\bullet
\otimes_{\mathcal{O}_Y}
\mathcal{M}^\bullet) \ar[d] \\
Lf^*\mathcal{K}^\bullet \otimes_{\mathcal{O}_X}^\mathbf{L}
Lf^*\mathcal{M}^\bullet \ar[d] &
f^*\text{Tot}(\mathcal{K}^\bullet
\otimes_{\mathcal{O}_Y}
\mathcal{M}^\bullet) \ar[d] \\
f^*\mathcal{K}^\bullet \otimes_{\mathcal{O}_X}^\mathbf{L}
f^*\mathcal{M}^\bullet \ar[r] &
\text{Tot}(f^*\mathcal{K}^\bullet \otimes_{\mathcal{O}_X}
f^*\mathcal{M}^\bullet)
}
$$
commutes.
\end{lemma}

\begin{proof}
We will use the existence of K-flat resolutions as in
Lemma \ref{lemma-pullback-K-flat}. If we choose such
resolutions $\mathcal{P}^\bullet \to \mathcal{K}^\bullet$
and $\mathcal{Q}^\bullet \to \mathcal{M}^\bullet$, then
we see that
$$
\xymatrix{
Lf^*\text{Tot}(\mathcal{P}^\bullet
\otimes_{\mathcal{O}_Y}
\mathcal{Q}^\bullet) \ar[r] \ar[d] &
Lf^*\text{Tot}(\mathcal{K}^\bullet
\otimes_{\mathcal{O}_Y}
\mathcal{M}^\bullet) \ar[d] \\
f^*\text{Tot}(\mathcal{P}^\bullet
\otimes_{\mathcal{O}_Y}
\mathcal{Q}^\bullet) \ar[d] \ar[r] &
f^*\text{Tot}(\mathcal{K}^\bullet
\otimes_{\mathcal{O}_Y}
\mathcal{M}^\bullet) \ar[d] \\
\text{Tot}(f^*\mathcal{P}^\bullet \otimes_{\mathcal{O}_X}
f^*\mathcal{Q}^\bullet) \ar[r] &
\text{Tot}(f^*\mathcal{K}^\bullet \otimes_{\mathcal{O}_X}
f^*\mathcal{M}^\bullet)
}
$$
commutes. However, now the left hand side of the diagram
is the left hand side of the diagram by our choice of
$\mathcal{P}^\bullet$ and $\mathcal{Q}^\bullet$ and
Lemma \ref{lemma-tensor-product-K-flat}.
\end{proof}









\section{Cohomology of unbounded complexes}
\label{section-unbounded}

\noindent
Let $(X, \mathcal{O}_X)$ be a ringed space.
The category $\textit{Mod}(\mathcal{O}_X)$ is a Grothendieck
abelian category: it has all colimits,
filtered colimits are exact, and it has a generator, namely
$$
\bigoplus\nolimits_{U \subset X\text{ open}} j_{U!}\mathcal{O}_U,
$$
see Modules, Section \ref{modules-section-kernels} and
Lemmas \ref{modules-lemma-j-shriek-flat} and
\ref{modules-lemma-module-quotient-flat}.
By Injectives, Theorem
\ref{injectives-theorem-K-injective-embedding-grothendieck}
for every complex $\mathcal{F}^\bullet$ of $\mathcal{O}_X$-modules
there exists an injective quasi-isomorphism
$\mathcal{F}^\bullet \to \mathcal{I}^\bullet$ to a K-injective complex
of $\mathcal{O}_X$-modules all of whose terms are injective
$\mathcal{O}_X$-modules and moreover this embedding can be
chosen functorial in the complex $\mathcal{F}^\bullet$.
It follows from
Derived Categories, Lemma \ref{derived-lemma-enough-K-injectives-implies}
that
\begin{enumerate}
\item any exact functor $F : K(\textit{Mod}(\mathcal{O}_X)) \to \mathcal{D}$
into a trianguated category $\mathcal{D}$ has a right derived functor
$RF : D(\mathcal{O}_X) \to \mathcal{D}$,
\item for any additive functor
$F : \textit{Mod}(\mathcal{O}_X) \to \mathcal{A}$
into an abelian category $\mathcal{A}$ we consider the exact functor
$F : K(\textit{Mod}(\mathcal{O}_X)) \to K(\mathcal{A})$ induced by $F$
and we obtain a right derived functor
$RF : D(\mathcal{O}_X) \to D(\mathcal{A})$.
\end{enumerate}
By construction we have $RF(\mathcal{F}^\bullet) = F(\mathcal{I}^\bullet)$
where $\mathcal{F}^\bullet \to \mathcal{I}^\bullet$ is as above.

\medskip\noindent
Here are some examples of the above:
\begin{enumerate}
\item The functor $\Gamma(X, -) : \textit{Mod}(\mathcal{O}_X) \to
\text{Mod}_{\Gamma(X, \mathcal{O}_X)}$ gives rise to
$$
R\Gamma(X, -) : D(\mathcal{O}_X) \to D(\Gamma(X, \mathcal{O}_X))
$$
We shall use the notation $H^i(X, K) = H^i(R\Gamma(X, K))$ for cohomology.
\item For an open $U \subset X$ we consider the functor
$\Gamma(U, -) : \textit{Mod}(\mathcal{O}_X) \to
\text{Mod}_{\Gamma(U, \mathcal{O}_X)}$. This gives rise to
$$
R\Gamma(U, -) : D(\mathcal{O}_X) \to D(\Gamma(U, \mathcal{O}_X))
$$
We shall use the notation $H^i(U, K) = H^i(R\Gamma(U, K))$ for cohomology.
\item For a morphism of ringed spaces
$f : (X, \mathcal{O}_X) \to (Y, \mathcal{O}_Y)$ we consider the functor
$f_* : \textit{Mod}(\mathcal{O}_X) \to \textit{Mod}(\mathcal{O}_Y)$
which gives rise to the total direct image
$$
Rf_* : D(\mathcal{O}_X) \longrightarrow D(\mathcal{O}_Y)
$$
on unbounded derived categories.
\end{enumerate}

\begin{lemma}
\label{lemma-adjoint}
Let $f : (X, \mathcal{O}_X) \to (Y, \mathcal{O}_Y)$ be a morphism of
ringed spaces. The functor $Rf_*$ defined above and the functor $Lf^*$
defined in Lemma \ref{lemma-derived-base-change} are adjoint:
$$
\Hom_{D(\mathcal{O}_X)}(Lf^*\mathcal{G}^\bullet, \mathcal{F}^\bullet)
=
\Hom_{D(\mathcal{O}_Y)}(\mathcal{G}^\bullet, Rf_*\mathcal{F}^\bullet)
$$
bifunctorially in $\mathcal{F}^\bullet \in \Ob(D(\mathcal{O}_X))$ and
$\mathcal{G}^\bullet \in \Ob(D(\mathcal{O}_Y))$.
\end{lemma}

\begin{proof}
This follows formally from the fact that $Rf_*$ and $Lf^*$ exist, see
Derived Categories, Lemma \ref{derived-lemma-derived-adjoint-functors}.
\end{proof}

\begin{lemma}
\label{lemma-derived-pushforward-composition}
Let $f : X \to Y$ and $g : Y \to Z$ be morphisms of ringed spaces.
Then $Rg_* \circ Rf_* = R(g \circ f)_*$ as functors
$D(\mathcal{O}_X) \to D(\mathcal{O}_Z)$.
\end{lemma}

\begin{proof}
By Lemma \ref{lemma-adjoint} we see that $Rg_* \circ Rf_*$
is adjoint to $Lf^* \circ Lg^*$. We have
$Lf^* \circ Lg^* = L(g \circ f)^*$ by
Lemma \ref{lemma-derived-pullback-composition}
and hence by
uniqueness of adjoint functors we have $Rg_* \circ Rf_* = R(g \circ f)_*$.
\end{proof}

\begin{remark}
\label{remark-base-change}
The construction of unbounded derived functor $Lf^*$ and $Rf_*$
allows one to construct the base change map in full generality.
Namely, suppose that
$$
\xymatrix{
X' \ar[r]_{g'} \ar[d]_{f'} &
X \ar[d]^f \\
S' \ar[r]^g &
S
}
$$
is a commutative diagram of ringed spaces. Let $K$ be an object of
$D(\mathcal{O}_X)$. Then there exists a canonical base change
map
$$
Lg^*Rf_*K \longrightarrow R(f')_*L(g')^*K
$$
in $D(\mathcal{O}_{S'})$. Namely, this map is adjoint to a map
$L(f')^*Lg^*Rf_*K \to L(g')^*K$
Since $L(f')^*Lg^* = L(g')^*Lf^*$ we see this is the same as a map
$L(g')^*Lf^*Rf_*K \to L(g')^*K$
which we can take to be $L(g')^*$ of the adjunction map
$Lf^*Rf_*K \to K$.
\end{remark}

\begin{remark}
\label{remark-compose-base-change}
Consider a commutative diagram
$$
\xymatrix{
X' \ar[r]_k \ar[d]_{f'} & X \ar[d]^f \\
Y' \ar[r]^l \ar[d]_{g'} & Y \ar[d]^g \\
Z' \ar[r]^m & Z
}
$$
of ringed spaces. Then the base change maps of
Remark \ref{remark-base-change}
for the two squares compose to give the base
change map for the outer rectangle. More precisely,
the composition
\begin{align*}
Lm^* \circ R(g \circ f)_*
& =
Lm^* \circ Rg_* \circ Rf_* \\
& \to Rg'_* \circ Ll^* \circ Rf_* \\
& \to Rg'_* \circ Rf'_* \circ Lk^* \\
& = R(g' \circ f')_* \circ Lk^*
\end{align*}
is the base change map for the rectangle. We omit the verification.
\end{remark}

\begin{remark}
\label{remark-compose-base-change-horizontal}
Consider a commutative diagram
$$
\xymatrix{
X'' \ar[r]_{g'} \ar[d]_{f''} & X' \ar[r]_g \ar[d]_{f'} & X \ar[d]^f \\
Y'' \ar[r]^{h'} & Y' \ar[r]^h & Y
}
$$
of ringed spaces. Then the base change maps of
Remark \ref{remark-base-change}
for the two squares compose to give the base
change map for the outer rectangle. More precisely,
the composition
\begin{align*}
L(h \circ h')^* \circ Rf_*
& =
L(h')^* \circ Lh_* \circ Rf_* \\
& \to L(h')^* \circ Rf'_* \circ Lg^* \\
& \to Rf''_* \circ L(g')^* \circ Lg^* \\
& = Rf''_* \circ L(g \circ g')^*
\end{align*}
is the base change map for the rectangle. We omit the verification.
\end{remark}

\begin{lemma}
\label{lemma-adjoints-push-pull-compatibility}
Let $f : (X, \mathcal{O}_X) \to (Y, \mathcal{O}_Y)$
be a morphism of ringed spaces. Let $\mathcal{K}^\bullet$
be a complex of $\mathcal{O}_X$-modules.
The diagram
$$
\xymatrix{
Lf^*f_*\mathcal{K}^\bullet \ar[r] \ar[d] &
f^*f_*\mathcal{K}^\bullet \ar[d] \\
Lf^*Rf_*\mathcal{K}^\bullet \ar[r] &
\mathcal{K}^\bullet
}
$$
coming from $Lf^* \to f^*$ on complexes, $f_* \to Rf_*$ on complexes,
and adjunction $Lf^* \circ Rf_* \to \text{id}$
commutes in $D(\mathcal{O}_X)$.
\end{lemma}

\begin{proof}
We will use the existence of K-flat resolutions and
K-injective resolutions, see Lemma \ref{lemma-pullback-K-flat}
and the discussion above. Choose a quasi-isomorphism
$\mathcal{K}^\bullet \to \mathcal{I}^\bullet$ where $\mathcal{I}^\bullet$
is K-injective as a complex of $\mathcal{O}_X$-modules.
Choose a quasi-isomorphism $\mathcal{Q}^\bullet \to f_*\mathcal{I}^\bullet$
where $\mathcal{Q}^\bullet$ is K-flat as a complex of
$\mathcal{O}_Y$-modules. We can choose a K-flat complex of
$\mathcal{O}_Y$-modules $\mathcal{P}^\bullet$
and a diagram of morphisms of complexes
$$
\xymatrix{
\mathcal{P}^\bullet \ar[r] \ar[d] &
f_*\mathcal{K}^\bullet \ar[d] \\
\mathcal{Q}^\bullet \ar[r] & f_*\mathcal{I}^\bullet
}
$$
commutative up to homotopy where the top horizontal arrow
is a quasi-isomorphism. Namely, we can first choose such a
diagram for some complex $\mathcal{P}^\bullet$ because
the quasi-isomorphisms form a multiplicative system in
the homotopy category of complexes and then we can replace
$\mathcal{P}^\bullet$ by a K-flat complex.
Taking pullbacks we obtain a diagram of morphisms of complexes
$$
\xymatrix{
f^*\mathcal{P}^\bullet \ar[r] \ar[d] &
f^*f_*\mathcal{K}^\bullet \ar[d] \ar[r] &
\mathcal{K}^\bullet \ar[d] \\
f^*\mathcal{Q}^\bullet \ar[r] &
f^*f_*\mathcal{I}^\bullet \ar[r] &
\mathcal{I}^\bullet
}
$$
commutative up to homotopy. The outer rectangle witnesses the
truth of the statement in the lemma.
\end{proof}

\begin{remark}
\label{remark-cup-product}
Let $f : (X, \mathcal{O}_X) \to (Y, \mathcal{O}_Y)$ be a morphism of
ringed spaces. The adjointness of $Lf^*$ and $Rf_*$ allows us to construct
a relative cup product
$$
Rf_*K \otimes_{\mathcal{O}_Y}^\mathbf{L} Rf_*L
\longrightarrow
Rf_*(K \otimes_{\mathcal{O}_X}^\mathbf{L} L)
$$
in $D(\mathcal{O}_Y)$ for all $K, L$ in $D(\mathcal{O}_X)$.
Namely, this map is adjoint to a map
$Lf^*(Rf_*K \otimes_{\mathcal{O}_Y}^\mathbf{L} Rf_*L) \to
K \otimes_{\mathcal{O}_X}^\mathbf{L} L$ for which we can take the
composition of the isomorphism
$Lf^*(Rf_*K \otimes_{\mathcal{O}_Y}^\mathbf{L} Rf_*L) =
Lf^*Rf_*K \otimes_{\mathcal{O}_X}^\mathbf{L} Lf^*Rf_*L$
(Lemma \ref{lemma-pullback-tensor-product})
with the map
$Lf^*Rf_*K \otimes_{\mathcal{O}_X}^\mathbf{L} Lf^*Rf_*L
\to K \otimes_{\mathcal{O}_X}^\mathbf{L} L$
coming from the counit $Lf^* \circ Rf_* \to \text{id}$.
\end{remark}






\section{Cohomology of filtered complexes}
\label{section-cohomology-filtered-object}

\noindent
Filtered complexes of sheaves frequently come up in a natural fashion
when studying cohomology of algebraic varieties, for example the de Rham
complex comes with its Hodge filtration. In this sectionwe use the very general
Injectives, Lemma \ref{injectives-lemma-K-injective-embedding-filtration}
to find construct spectral sequences on cohomology and we relate these to
previously constructed spectral sequences.

\begin{lemma}
\label{lemma-spectral-sequence-filtered-object}
Let $(X, \mathcal{O}_X)$ be a ringed space. Let $\mathcal{F}^\bullet$ be a
filtered complex of $\mathcal{O}_X$-modules. There exists a canonical
spectral sequence $(E_r, \text{d}_r)_{r \geq 1}$ of bigraded
$\Gamma(X, \mathcal{O}_X)$-modules with $d_r$ of bidegree $(r, -r + 1)$ and
$$
E_1^{p, q} = H^{p + q}(X, \text{gr}^p\mathcal{F}^\bullet)
$$
If for every $n$ we have
$$
H^n(X, F^p\mathcal{F}^\bullet) = 0\text{ for }p \gg 0
\quad\text{and}\quad
H^n(X, F^p\mathcal{F}^\bullet) = H^n(X, \mathcal{F}^\bullet)\text{ for }p \ll 0
$$
then the spectral sequence is bounded and converges to
$H^*(X, \mathcal{F}^\bullet)$.
\end{lemma}

\begin{proof}
(For a proof in case the complex is a bounded below complex
of modules with finite filtrations, see the remark below.)
Choose an map of filtered complexes
$j : \mathcal{F}^\bullet \to \mathcal{J}^\bullet$ as in
Injectives, Lemma
\ref{injectives-lemma-K-injective-embedding-filtration}.
The spectral sequence is the spectral sequence of
Homology, Section \ref{homology-section-filtered-complex}
associated to the filtered complex
$$
\Gamma(X, \mathcal{J}^\bullet)
\quad\text{with}\quad
F^p\Gamma(X, \mathcal{J}^\bullet) = \Gamma(X, F^p\mathcal{J}^\bullet)
$$
Since cohomology is computed by evaluating on K-injective representatives
we see that the $E_1$ page is as stated in the lemma.
The convergence and boundedness under the stated conditions
follows from Homology, Lemma \ref{homology-lemma-ss-converges-trivial}.
\end{proof}

\begin{remark}
\label{remark-spectral-sequence-filtered-object}
Let $(X, \mathcal{O}_X)$ be a ringed space. Let $\mathcal{F}^\bullet$ be a
filtered complex of $\mathcal{O}_X$-modules. If $\mathcal{F}^\bullet$
is bounded from below and for each $n$ the filtration on $\mathcal{F}^n$
is finite, then there is a construction of the spectral sequence
in Lemma \ref{lemma-spectral-sequence-filtered-object}
avoiding Injectives, Lemma
\ref{injectives-lemma-K-injective-embedding-filtration}.
Namely, by
Derived Categories, Lemma
\ref{derived-lemma-right-resolution-by-filtered-injectives}
there is a filtered quasi-isomorphism
$i : \mathcal{F}^\bullet \to \mathcal{I}^\bullet$
of filtered complexes
with $\mathcal{I}^\bullet$ bounded below,
the filtration on $\mathcal{I}^n$ is finite for all $n$,
and with each $\text{gr}^p\mathcal{I}^n$ an
injective $\mathcal{O}_X$-module.
Then we take the spectral sequence associated to
$$
\Gamma(X, \mathcal{I}^\bullet)
\quad\text{with}\quad
F^p\Gamma(X, \mathcal{I}^\bullet) = \Gamma(X, F^p\mathcal{I}^\bullet)
$$
Since cohomology can be computed by evaluating on
bounded below complexes of injectives
we see that the $E_1$ page is as stated in the lemma.
The convergence and boundedness under the stated conditions
follows from
Homology, Lemma \ref{homology-lemma-biregular-ss-converges}.
In fact, this is a special case of the spectral sequence
in Derived Categories, Lemma \ref{derived-lemma-ss-filtered-derived}.
\end{remark}

\begin{example}
\label{example-spectral-sequence}
Let $(X, \mathcal{O}_X)$ be a ringed space. Let $\mathcal{F}^\bullet$ be a
complex of $\mathcal{O}_X$-modules. We can apply
Lemma \ref{lemma-spectral-sequence-filtered-object}
with $F^p\mathcal{F}^\bullet = \tau_{\leq -p}\mathcal{F}^\bullet$.
(If $\mathcal{F}^\bullet$ is bounded below we can use
Remark \ref{remark-spectral-sequence-filtered-object}.)
Then we get a spectral sequence
$$
E_1^{p, q} = H^{p + q}(X, H^{-p}(\mathcal{F}^\bullet)[p]) =
H^{2p + q}(X, H^{-p}(\mathcal{F}^\bullet))
$$
After renumbering $p = -j$ and $q = i + 2j$ we find that for any
$K \in D(\mathcal{O}_X)$ there is a spectral sequence
$(E'_r, d'_r)_{r \geq 2}$ of bigraded modules with
$d'_r$ of bidegree $(r, -r + 1)$, with
$$
(E'_2)^{i, j} = H^i(X, H^j(K))
$$
If $K$ is bounded below (for example), then this spectral sequence
is bounded and converges to $H^{i + j}(X, K)$.
In the bounded below case this spectral sequence is an example
of the second spectral sequence of
Derived Categories, Lemma \ref{derived-lemma-two-ss-complex-functor}
(constructed using Cartan-Eilenberg resolutions).
\end{example}

\begin{example}
\label{example-spectral-sequence-bis}
Let $(X, \mathcal{O}_X)$ be a ringed space. Let $\mathcal{F}^\bullet$ be a
complex of $\mathcal{O}_X$-modules. We can apply
Lemma \ref{lemma-spectral-sequence-filtered-object}
with $F^p\mathcal{F}^\bullet = \sigma_{\geq p}\mathcal{F}^\bullet$.
Then we get a spectral sequence
$$
E_1^{p, q} = H^{p + q}(X, \mathcal{F}^p[-p]) = H^q(X, \mathcal{F}^p)
$$
If $\mathcal{F}^\bullet$ is bounded below, then
\begin{enumerate}
\item we can use
Remark \ref{remark-spectral-sequence-filtered-object}
to construct this spectral sequence,
\item the spectral sequence is bounded and converges to
$H^{i + j}(X, \mathcal{F}^\bullet)$, and
\item the spectral sequence is equal to the first spectral sequence of
Derived Categories, Lemma \ref{derived-lemma-two-ss-complex-functor}
(constructed using Cartan-Eilenberg resolutions).
\end{enumerate}
\end{example}

\begin{lemma}
\label{lemma-relative-spectral-sequence-filtered-object}
Let $f : (X, \mathcal{O}_X) \to (Y, \mathcal{O}_Y)$ be a morphism of
ringed spaces. Let $\mathcal{F}^\bullet$ be a filtered complex of
$\mathcal{O}_X$-modules. There exists a canonical spectral sequence
$(E_r, \text{d}_r)_{r \geq 1}$ of bigraded
$\mathcal{O}_Y$-modules with $d_r$ of bidegree $(r, -r + 1)$ and
$$
E_1^{p, q} = R^{p + q}f_*\text{gr}^p\mathcal{F}^\bullet
$$
If for every $n$ we have
$$
R^nf_*F^p\mathcal{F}^\bullet = 0 \text{ for }p \gg 0
\quad\text{and}\quad
R^nf_*F^p\mathcal{F}^\bullet = R^nf_*\mathcal{F}^\bullet \text{ for }p \ll 0
$$
then the spectral sequence is bounded and converges to
$Rf_*\mathcal{F}^\bullet$.
\end{lemma}

\begin{proof}
The proof is exactly the same as the proof of
Lemma \ref{lemma-spectral-sequence-filtered-object}.
\end{proof}





\section{Godement resolution}
\label{section-godement}

\noindent
A reference is \cite{Godement}.

\medskip\noindent
Let $(X, \mathcal{O}_X)$ be a ringed space. Denote $X_{disc}$ the discrete
topological space with the same points as $X$. Denote $f : X_{disc} \to X$
the obvious continuous map. Set $\mathcal{O}_{X_{disc}} = f^{-1}\mathcal{O}_X$.
Then $f : (X_{disc}, \mathcal{O}_{X_{disc}}) \to (X, \mathcal{O}_X)$
is a flat morphism of ringed spaces. We can apply the
{\it dual} of the material in
Simplicial, Section \ref{simplicial-section-standard} to the adjoint pair of
functors $f^*, f_*$ on sheaves of modules. Thus we obtain an augmented
cosimplicial object
$$
\xymatrix{
\text{id} \ar[r] &
f_*f^* \ar@<1ex>[r] \ar@<-1ex>[r] &
f_*f^*f_*f^* \ar@<0ex>[l] \ar@<-2ex>[r] \ar@<0ex>[r] \ar@<2ex>[r] &
f_*f^*f_*f^*f_*f^* \ar@<1ex>[l] \ar@<-1ex>[l]
}
$$
in the category of functors from $\textit{Mod}(\mathcal{O}_X)$
to itself, see Simplicial, Lemma \ref{simplicial-lemma-standard-simplicial}.
Moreover, the augmentation
$$
\xymatrix{
f^* \ar[r] &
f^*f_*f^* \ar@<1ex>[r] \ar@<-1ex>[r] &
f^*f_*f^*f_*f^* \ar@<0ex>[l] \ar@<-2ex>[r] \ar@<0ex>[r] \ar@<2ex>[r] &
f^*f_*f^*f_*f^*f_*f^* \ar@<1ex>[l] \ar@<-1ex>[l]
}
$$
is a homotopy equivalence, see
Simplicial, Lemma \ref{simplicial-lemma-standard-simplicial-homotopy}.

\begin{lemma}
\label{lemma-godement-resolution}
Let $(X, \mathcal{O}_X)$ be a ringed space. For every sheaf of
$\mathcal{O}_X$-modules $\mathcal{F}$ there is a resolution
$$
0 \to
\mathcal{F} \to
f_*f^*\mathcal{F} \to
f_*f^*f_*f^*\mathcal{F} \to
f_*f^*f_*f^*f_*f^*\mathcal{F} \to \ldots
$$
functorial in $\mathcal{F}$ such that each term
$f_*f^* \ldots f_*f^*\mathcal{F}$ is a flasque
$\mathcal{O}_X$-module and such that for all $x \in X$ the
map
$$
\mathcal{F}_x[0] \to \Big(
(f_*f^*\mathcal{F})_x \to
(f_*f^*f_*f^*\mathcal{F})_x \to
(f_*f^*f_*f^*f_*f^*\mathcal{F})_x \to \ldots
\Big)
$$
is a homotopy equivalence in the category of complexes
of $\mathcal{O}_{X, x}$-modules.
Moreover, the assignment of the nonnegative resolution complex to
$\mathcal{F}$ is an exact functor from $\textit{Mod}(\mathcal{O}_X)$ to
$\text{CoCh}_{\geq 0}(\textit{Mod}(\mathcal{O}_X))$.
\end{lemma}

\begin{proof}
The complex $f_*f^*\mathcal{F} \to  f_*f^*f_*f^*\mathcal{F} \to
f_*f^*f_*f^*f_*f^*\mathcal{F} \to \ldots$ is the complex associated
to the cosimplicial object with terms
$f_*f^*\mathcal{F}, f_*f^*f_*f^*\mathcal{F},
f_*f^*f_*f^*f_*f^*\mathcal{F}, \ldots$ described above, see
Simplicial, Section \ref{simplicial-section-dold-kan-cosimplicial}.
The augmentation gives rise to the map $\mathcal{F} \to f_*f^*\mathcal{F}$
as indicated. For any abelian sheaf $\mathcal{H}$ on $X_{disc}$ the
pushforward $f_*\mathcal{H}$ is flasque because $X_{disc}$
is a discrete space and the pushforward of a flasque sheaf is flasque.
Hence the terms of the complex are flasque $\mathcal{O}_X$-modules.

\medskip\noindent
If $x \in X_{disc} = X$ is a point, then $(f^*\mathcal{G})_x = \mathcal{G}_x$
for any $\mathcal{O}_X$-module $\mathcal{G}$. Hence $f^*$ is an exact functor
and a complex of $\mathcal{O}_X$-modules
$\mathcal{G}_1 \to \mathcal{G}_2 \to \mathcal{G}_3$
is exact if and only if
$f^*\mathcal{G}_1 \to f^*\mathcal{G}_2 \to f^*\mathcal{G}_3$
is exact (see Modules, Lemma \ref{modules-lemma-abelian}).
The result mentioned in the introduction to this section
proves the pullback by $f^*$ gives a homotopy equivalence from
the constant cosimplicial object $f^*\mathcal{F}$ to the
cosimplicial object with terms
$f_*f^*\mathcal{F}, f_*f^*f_*f^*\mathcal{F},
f_*f^*f_*f^*f_*f^*\mathcal{F}, \ldots$.
By Simplicial, Lemma \ref{simplicial-lemma-homotopy-equivalence-s-Q}
we obtain that
$$
f^*\mathcal{F}[0] \to \Big(
f^*f_*f^*\mathcal{F} \to
f^*f_*f^*f_*f^*\mathcal{F} \to
f^*f_*f^*f_*f^*f_*f^*\mathcal{F} \to \ldots
\Big)
$$
is a homotopy equivalence. This immediately implies the two remaining
statements of the lemma.

It remains to prove exactness of the resolution functor. We have already
seen that $f^*$ is exact. If $\mathcal{H}$ is an
$\mathcal{O}_{X_{disc}}$-module, then
$$
(f_*\mathcal{H})(U) = \prod_{x \in U} \mathcal{H}_x.
$$
Kernels and surjections of modules are preserved by products. Hence $f_*$
is exact. It follows that every iterate $f_*f^*\ldots f_*f^*$ is exact.
The differentials of the resolution are natural transformations between
these iterates, so the resulting functor to complexes is exact degree by
degree.
\end{proof}

\begin{lemma}
\label{lemma-family-supports-colimit}
\begin{reference}
\cite[Chapter III, \S 3.6, Proposition 3.6.3, p. 170]{Tohoku}
\end{reference}
Let $I$ be a directed set, let $(\mathcal{F}_i)_{i \in I}$ be a directed
system of abelian sheaves on a topological space $X$, and let $\Phi$ be a
family of supports. Assume one of the following:
\begin{enumerate}
\item $X$ is locally quasi-compact and Hausdorff, and $\Phi$ is the family
of quasi-compact subsets of $X$;
\item $X$ is Noetherian, and $\Phi$ is the family of all closed subsets of
$X$.
\end{enumerate}
Then the canonical map
$$
\colim_i H^p_\Phi(X, \mathcal{F}_i)
\longrightarrow
H^p_\Phi(X, \colim_i \mathcal{F}_i)
$$
is an isomorphism for every $p \in \mathbf{Z}$.
\end{lemma}

\begin{proof}
Put $\mathcal{F} = \colim_i \mathcal{F}_i$. We first show that
$\Gamma_\Phi(X, -)$ commutes with directed colimits.

Assume (1). For injectivity, suppose that a compactly supported section
$s_i$ maps to zero, and let $K$ be its support. Apply Lemma
\ref{lemma-compact-sections-colimit} to the restrictions to the
quasi-compact Hausdorff space $K$. At some later stage the restriction of
$s_i$ to $K$ is zero. Its image was already zero away from $K$, so it is
zero globally at that stage.

For surjectivity, let $s \in \Gamma_\Phi(X, \mathcal{F})$ and let $K$ be
its support. Local quasi-compactness, the Hausdorff property, and
quasi-compactness of $K$ give a quasi-compact closed neighbourhood $B$ of
$K$. Thus $K \subset \mathop{\rm Int}(B)$. By Lemma
\ref{lemma-compact-sections-colimit}, the restriction $s|_B$ comes from
some
$$
t_i \in \Gamma(B, \mathcal{F}_i|_B).
$$
Its image is zero on the quasi-compact closed subset
$D = B \setminus \mathop{\rm Int}(B)$. After increasing $i$, injectivity
in Lemma \ref{lemma-compact-sections-colimit}, applied to $D$, shows that
$t_i|_D=0$. Hence the support of $t_i$ on $B$ is contained in
$\mathop{\rm Int}(B)$. Lemma \ref{lemma-extend-zero-from-closed} extends
$t_i$ by zero to a section of $\mathcal{F}_i$ on $X$ with quasi-compact
support. Its image is $s$.

Assume (2). Every open subset of $X$, and every intersection of two such
opens, is Noetherian and hence quasi-compact by Topology, Lemmas
\ref{topology-lemma-Noetherian} and
\ref{topology-lemma-Noetherian-quasi-compact}. Sheaves, Lemma
\ref{sheaves-lemma-directed-colimits-sections}, part (4), therefore shows
that global sections commute with directed colimits. Since every global
section has closed support, $\Gamma_\Phi(X, -)=\Gamma(X, -)$ in this case.

Write
$$
\mathcal{F}_i \longrightarrow \mathcal{G}_i^\bullet
$$
for the functorial nonnegative flasque resolution of Lemma
\ref{lemma-godement-resolution}, applied with the constant sheaf of rings
$\underline{\mathbf{Z}}_X$. Filtered colimits are exact in the category of
abelian sheaves by Modules, Lemma
\ref{modules-lemma-limits-colimits}. Consequently
$$
\mathcal{F} \longrightarrow
\mathcal{G}^\bullet = \colim_i \mathcal{G}_i^\bullet
$$
is again a resolution.

In case (1), the family $\Phi$ is paracompactifying. Indeed, its members
are quasi-compact Hausdorff and hence paracompact, and local
quasi-compactness gives every member a quasi-compact closed neighbourhood.
Thus every $\mathcal{G}_i^n$ is $\Phi$-soft by Lemma
\ref{lemma-flasque-soft-family-supports}, and every $\mathcal{G}^n$ is
$\Phi$-soft by Lemma \ref{lemma-colimit-soft-flasque}. In case (2), every
$\mathcal{G}^n$ is flasque by the same lemma. Lemma
\ref{lemma-family-supports-flasque-soft-acyclic} therefore shows that these
resolutions compute supported cohomology.

The degree zero isomorphisms proved above apply term by term to these
resolutions. Since filtered colimits of abelian groups are exact, they
commute with cohomology of complexes. We obtain
$$
\begin{aligned}
H^p_\Phi(X, \mathcal{F})
&= H^p(\Gamma_\Phi(X, \mathcal{G}^\bullet)) \\
&= H^p(\colim_i \Gamma_\Phi(X, \mathcal{G}_i^\bullet)) \\
&= \colim_i H^p(\Gamma_\Phi(X, \mathcal{G}_i^\bullet)) \\
&= \colim_i H^p_\Phi(X, \mathcal{F}_i).
\end{aligned}
$$
\end{proof}

\begin{lemma}
\label{lemma-family-supports-colimit-extension-by-zero}
\begin{reference}
\cite[Chapter III, \S 3.10, Corollary 1, p. 182]{Tohoku}
\end{reference}
Let $X$ be a topological space, let $\Phi$ be a paracompactifying family
of supports on $X$, and let $\mathcal{F}$ be an abelian sheaf on $X$.
Let $\mathcal{U}_\Phi$ be the directed set of open subsets $U\subset X$
such that $\overline{U}\in\Phi$, ordered by inclusion. Then the canonical
maps
$$
\colim_{U\in\mathcal{U}_\Phi} H^p(X,\mathcal{F}_U)
\longrightarrow H^p_\Phi(X,\mathcal{F})
$$
are isomorphisms for every $p\in\mathbf{Z}$. Here
$\mathcal{F}_U=j_!(\mathcal{F}|_U)$ for the inclusion $j:U\to X$.
These isomorphisms are functorial in $\mathcal{F}$ and compatible with the
connecting maps.
\end{lemma}

\begin{proof}
The set $\mathcal{U}_\Phi$ is nonempty since
$\emptyset\in\mathcal{U}_\Phi$, and it is directed: if
$U,V\in\mathcal{U}_\Phi$, then
$$
\overline{U\cup V}=\overline{U}\cup\overline{V}\in\Phi.
$$
For $U\in\mathcal{U}_\Phi$, consider the additive functor
$$
G_U(\mathcal{F})=\Gamma(X,\mathcal{F}_U).
$$
The inclusions $\mathcal{F}_U\to\mathcal{F}_V$ for $U\subset V$ make these
functors into a directed system. There is a functorial equality
$$
\colim_{U\in\mathcal{U}_\Phi}G_U(\mathcal{F})
=\Gamma_\Phi(X,\mathcal{F}).
$$
Indeed, a section in $G_U(\mathcal{F})$ has support contained in $U$, and
therefore in $\overline{U}\in\Phi$. Conversely, if a section has support
$A\in\Phi$, choose a closed neighbourhood $B\in\Phi$ of $A$. For
$U=\mathop{\rm Int}(B)$ we have $A\subset U$ and
$\overline{U}\subset B$, so the section comes from
$\Gamma(X,\mathcal{F}_U)$.

We next identify the derived functors of $G_U$. The functor
$\mathcal{F}\mapsto\mathcal{F}_U$ is exact by Lemma
\ref{lemma-extension-by-zero-locally-closed}. Let $\mathcal{I}$ be an
injective abelian sheaf on $X$. Then $\mathcal{I}|_U$ is injective by
Lemma \ref{lemma-cohomology-of-open}, applied to the constant sheaves of
rings. Put $B=\overline{U}$, let $i:B\to X$ be the inclusion, and write
$\mathcal{I}_U=i_*\mathcal{H}$; this is possible because
$\mathcal{I}_U$ is supported on $B$. Since $B\in\Phi$, the induced family
$\Phi_B$ is the family of all closed subsets of $B$. Lemmas
\ref{lemma-cohomology-and-closed-immersions} and
\ref{lemma-family-supports-extension-by-zero}, first for the closed subset
$B$ and then for the open subset $U$, give, for $q>0$,
$$
\begin{aligned}
H^q(X,\mathcal{I}_U)
&=H^q(B,\mathcal{H})
=H^q_{\Phi_B}(B,\mathcal{H}) \\
&=H^q_\Phi(X,\mathcal{I}_U)
=H^q_{\Phi_U}(U,\mathcal{I}|_U)=0.
\end{aligned}
$$
The last equality is Lemma \ref{lemma-cohomology-family-supports}, since
$\mathcal{I}|_U$ is injective. Thus $\mathcal{I}_U$ is acyclic for global
sections. If $\mathcal{F}\to\mathcal{I}^\bullet$ is an injective
resolution, exactness of $(-)_U$ makes
$\mathcal{F}_U\to\mathcal{I}^\bullet_U$ a resolution by sheaves acyclic
for global sections. It follows that
$$
R^pG_U(\mathcal{F})=H^p(X,\mathcal{F}_U).
$$

The category of abelian sheaves on $X$ has enough injectives by Injectives,
Lemma \ref{injectives-lemma-abelian-sheaves-space}, and the category of
abelian groups satisfies AB5. Derived Categories, Lemma
\ref{derived-lemma-filtered-colimit-right-derived-functors}, applied to the
functors $G_U$, now gives
$$
\colim_{U\in\mathcal{U}_\Phi}H^p(X,\mathcal{F}_U)
=R^p\Gamma_\Phi(X,\mathcal{F})
=H^p_\Phi(X,\mathcal{F}).
$$
The same lemma gives functoriality and compatibility with the connecting
maps.
\end{proof}


\begin{definition}
\label{definition-cohomological-dimension}
\begin{reference}
\cite[Chapter III, \S 3.6, p. 170]{Tohoku}
\end{reference}
Let $X$ be a topological space and let $n \geq 0$. We say that $X$ has
{\it cohomological dimension at most $n$} if
$$
H^q(X, \mathcal{F})=0
$$
for every abelian sheaf $\mathcal{F}$ on $X$ and every $q>n$.
\end{definition}

\begin{lemma}
\label{lemma-cohomological-dimension-open-generators}
\begin{reference}
\cite[Chapter III, \S 3.6, Corollary to Proposition 3.6.3,
p. 170]{Tohoku}
\end{reference}
Let $X$ be either a quasi-compact Hausdorff space or a Noetherian space,
and let $n \geq 0$. Then $X$ has cohomological dimension at most $n$ if
and only if
$$
H^q(X, \mathbf{Z}_U)=0
$$
for every open subset $U \subset X$ and every $q>n$, where
$\mathbf{Z}_U=j_!\underline{\mathbf{Z}}_U$ for the inclusion
$j:U\to X$.
\end{lemma}

\begin{proof}
One implication is immediate. Conversely, fix $q>n$. The functors
$$
T^r=H^r(X,-),\qquad r\in\mathbf{Z},
$$
form an integer-indexed cohomological functor from abelian sheaves to
abelian groups. They commute with directed colimits by Lemma
\ref{lemma-family-supports-colimit}: in the quasi-compact Hausdorff case
compactly supported cohomology is ordinary cohomology, and in the
Noetherian case use the family of all closed subsets. Apply Lemma
\ref{lemma-cohomological-functor-local-generators} with $p=q$ and with the
zero Serre subcategory of abelian groups. Its three required degrees are
$q,q+1,q+2$, all strictly larger than $n$. The assumed vanishing therefore
gives $H^q(X,\mathcal{F})=0$ for every abelian sheaf $\mathcal{F}$.
\end{proof}


\begin{lemma}
\label{lemma-godement-resolution-bounded-below}
Let $(X, \mathcal{O}_X)$ be a ringed space. Let
$\mathcal{F}^\bullet$ be a bounded below complex of
$\mathcal{O}_X$-modules. There exists a quasi-isomorphism
$\mathcal{F}^\bullet \to \mathcal{G}^\bullet$
where $\mathcal{G}^\bullet$ be a bounded below complex of flasque
$\mathcal{O}_X$-modules and for all $x \in X$ the
map $\mathcal{F}^\bullet_x \to \mathcal{G}^\bullet_x$
is a homotopy equivalence in the category of complexes
of $\mathcal{O}_{X, x}$-modules.
\end{lemma}

\begin{proof}
Let $\mathcal{A}$ be the category of complexes of $\mathcal{O}_X$-modules
and let $\mathcal{B}$ be the category of complexes of $\mathcal{O}_X$-modules.
Then we can apply the discussion above to the adjoint functors
$f^*$ and $f_*$ between $\mathcal{A}$ and $\mathcal{B}$.
Arguing exactly as in the proof of
Lemma \ref{lemma-godement-resolution}
we get a resolution
$$
0 \to
\mathcal{F}^\bullet \to
f_*f^*\mathcal{F}^\bullet \to
f_*f^*f_*f^*\mathcal{F}^\bullet \to
f_*f^*f_*f^*f_*f^*\mathcal{F}^\bullet \to \ldots
$$
in the abelian category $\mathcal{A}$ such that each term of each
$f_*f^*\ldots f_*f^*\mathcal{F}^\bullet$ is a flasque
$\mathcal{O}_X$-module and such that for all $x \in X$ the
map
$$
\mathcal{F}^\bullet_x[0] \to \Big(
(f_*f^*\mathcal{F}^\bullet)_x \to
(f_*f^*f_*f^*\mathcal{F}^\bullet)_x \to
(f_*f^*f_*f^*f_*f^*\mathcal{F}^\bullet)_x \to \ldots
\Big)
$$
is a homotopy equivalence in the category of complexes of complexes
of $\mathcal{O}_{X, x}$-modules. Since a complex of complexes is the
same thing as a double complex, we can consider the induced map
$$
\mathcal{F}^\bullet \to
\mathcal{G}^\bullet =
\text{Tot}(
f_*f^*\mathcal{F}^\bullet \to
f_*f^*f_*f^*\mathcal{F}^\bullet \to
f_*f^*f_*f^*f_*f^*\mathcal{F}^\bullet \to \ldots
)
$$
Since the complex $\mathcal{F}^\bullet$ is bounded below, the
same is true for $\mathcal{G}^\bullet$ and in fact each term
of $\mathcal{G}^\bullet$ is a finite direct sum of
terms of the complexes $f_*f^*\ldots f_*f^*\mathcal{F}^\bullet$
and hence is flasque. The final assertion of the lemma
now follows from
Homology, Lemma \ref{homology-lemma-homotopy-complex-complexes}.
Since this in particular shows that
$\mathcal{F}^\bullet \to \mathcal{G}^\bullet$
is a quasi-isomorphism, the proof is complete.
\end{proof}














\section{Cup product}
\label{section-cup-product}

\noindent
Let $(X, \mathcal{O}_X)$ be a ringed space. Let $K, M$ be objects
of $D(\mathcal{O}_X)$. Set $A = \Gamma(X, \mathcal{O}_X)$.
The (global) cup product in this setting is a map
$$
\mu :
R\Gamma(X, K) \otimes_A^\mathbf{L} R\Gamma(X, M)
\longrightarrow
R\Gamma(X, K \otimes_{\mathcal{O}_X}^\mathbf{L} M)
$$
in $D(A)$. We define it as the relative cup product for the
morphism of ringed spaces $f : (X, \mathcal{O}_X) \to (pt, A)$
as in Remark \ref{remark-cup-product} via $D(pt, A) = D(A)$.
This map in particular defines pairings
$$
\cup :
H^i(X, K) \times H^j(X, M)
\longrightarrow
H^{i + j}(X, K \otimes_{\mathcal{O}_X}^\mathbf{L} M)
$$
Namely, given $\xi \in H^i(X, K) = H^i(R\Gamma(X, K))$ and
$\eta \in H^j(X, M) = H^j(R\Gamma(X, M))$ we can
first ``tensor'' them to get an element $\xi \otimes \eta$ in
$H^{i + j}(R\Gamma(X, K) \otimes_A^\mathbf{L} R\Gamma(X, M))$, see
More on Algebra, Section \ref{more-algebra-section-products-tor}.
Then we can apply $\mu$ to get the desired element
$\xi \cup \eta = \mu(\xi \otimes \eta)$
of $H^{i + j}(X, K \otimes_{\mathcal{O}_X}^\mathbf{L} M)$.

\medskip\noindent
Here is another way to think of the cup product of $\xi$ and $\eta$.
Namely, we can write
$$
H^i(R\Gamma(X, K)) = \Hom_{D(\mathcal{O}_X)}(\mathcal{O}_X[-i], K)
\quad\text{and}\quad
H^j(R\Gamma(X, M)) = \Hom_{D(\mathcal{O}_X)}(\mathcal{O}_X[-j], M)
$$
because $\Hom(\mathcal{O}_X, -) = \Gamma(X, -)$.
Thus $\xi$ and $\eta$ are the ``same'' thing as maps
$$
\tilde \xi : \mathcal{O}_X[-i] \to K
\quad\text{and}\quad
\tilde \eta : \mathcal{O}_X[-j] \to M
$$
Combining this with the functoriality of the derived tensor product
we obtain
$$
\mathcal{O}_X[-i - j] =
\mathcal{O}_X[-i] \otimes_{\mathcal{O}_X}^\mathbf{L} \mathcal{O}_X[-j]
\xrightarrow{\tilde \xi \otimes \tilde \eta}
K \otimes_{\mathcal{O}_X}^\mathbf{L} M
$$
which by the same token as above is an element of
$H^{i + j}(X, K \otimes_{\mathcal{O}_X}^\mathbf{L} M)$.

\begin{lemma}
\label{lemma-second-cup-equals-first}
This construction gives the cup product.
\end{lemma}

\begin{proof}
With $f : (X, \mathcal{O}_X) \to (pt, A)$ as above we have
$Rf_*(-) = R\Gamma(X, -)$ and our map $\mu$ is adjoint to the map
$$
Lf^*(Rf_*K \otimes_A^\mathbf{L} Rf_*M) =
Lf^*Rf_*K \otimes_{\mathcal{O}_X}^\mathbf{L} Lf^*Rf_*M
\xrightarrow{\epsilon_K \otimes \epsilon_M}
K \otimes_{\mathcal{O}_X}^\mathbf{L} M
$$
where $\epsilon$ is the counit of the adjunction between
$Lf^*$ and $Rf_*$.
If we think of $\xi$ and $\eta$ as maps $\xi : A[-i] \to R\Gamma(X, K)$
and $\eta : A[-j] \to R\Gamma(X, M)$, then
the tensor $\xi \otimes \eta$ corresponds to the map\footnote{There
is a sign hidden here, namely, the equality is defined by
the composition
$$
A[-i - j] \to (A \otimes_A^\mathbf{L} A)[-i - j] \to
A[-i] \otimes_A^\mathbf{L} A[-j]
$$
where in the second step we use the identification of
More on Algebra, Item (\ref{more-algebra-item-shift-tensor})
which uses a sign in principle.
Except, in this case the sign is $+1$ by our convention and even if it wasn't
$+1$ it wouldn't matter since we used the same sign
in the identification
$\mathcal{O}_X[-i - j] =
\mathcal{O}_X[-i] \otimes_{\mathcal{O}_X}^\mathbf{L} \mathcal{O}_X[-j]$.}
$$
A[-i - j] = A[-i] \otimes_A^\mathbf{L} A[-j]
\xrightarrow{\xi \otimes \eta}
R\Gamma(X, K) \otimes_A^\mathbf{L} R\Gamma(X, M)
$$
By definition the cup product $\xi \cup \eta$ is the map
$A[-i - j] \to R\Gamma(X, K \otimes_{\mathcal{O}_X}^\mathbf{L} M)$
which is adjoint to
$$
(\epsilon_K \otimes \epsilon_M) \circ Lf^*(\xi \otimes \eta) =
(\epsilon_K \circ Lf^*\xi) \otimes (\epsilon_M \circ Lf^*\eta)
$$
However, it is easy to see that
$\epsilon_K \circ Lf^*\xi = \tilde \xi$ and
$\epsilon_M \circ Lf^*\eta = \tilde \eta$.
We conclude that $\widetilde{\xi \cup \eta} = \tilde \xi \otimes \tilde \eta$
which means we have the desired agreement.
\end{proof}

\begin{remark}
\label{remark-cup-with-element-map-total-cohomology}
Let $(X, \mathcal{O}_X)$ be a ringed space. Let $K, M$ be objects
of $D(\mathcal{O}_X)$. Set $A = \Gamma(X, \mathcal{O}_X)$.
Given $\xi \in H^i(X, K)$ we get an associated map
$$
\xi = ``\xi \cup -'' :
R\Gamma(X, M)[-i]
\to
R\Gamma(X, K \otimes_{\mathcal{O}_X}^\mathbf{L} M)
$$
by representing $\xi$ as a map $\xi : A[-i] \to R\Gamma(X, K)$ as in the
proof of Lemma \ref{lemma-second-cup-equals-first}
and then using the composition
$$
R\Gamma(X, M)[-i] = A[-i] \otimes_A^\mathbf{L} R\Gamma(X, M)
\xrightarrow{\xi \otimes 1}
R\Gamma(X, K) \otimes_A^\mathbf{L} R\Gamma(X, M)
\to
R\Gamma(X, K \otimes_{\mathcal{O}_X}^\mathbf{L} M)
$$
where the second arrow is the global cup product $\mu$ above.
On cohomology this recovers the cup product by $\xi$ as is clear
from Lemma \ref{lemma-second-cup-equals-first} and its proof.
\end{remark}

\noindent
Let us formulate and prove a natural compatibility of the
relative cup product. Namely, suppose that we have a morphism
$f : (X, \mathcal{O}_X) \to (Y, \mathcal{O}_Y)$ of ringed spaces.
Let $\mathcal{K}^\bullet$ and $\mathcal{M}^\bullet$
be complexes of $\mathcal{O}_X$-modules.
There is a naive cup product
$$
\text{Tot}(
f_*\mathcal{K}^\bullet
\otimes_{\mathcal{O}_Y}
f_*\mathcal{M}^\bullet)
\longrightarrow
f_*\text{Tot}(\mathcal{K}^\bullet
\otimes_{\mathcal{O}_X}
\mathcal{M}^\bullet)
$$
We claim that this is related to the relative cup product.

\begin{lemma}
\label{lemma-cup-compatible-with-naive}
In the situation above the following diagram commutes
$$
\xymatrix{
f_*\mathcal{K}^\bullet
\otimes_{\mathcal{O}_Y}^\mathbf{L}
f_*\mathcal{M}^\bullet \ar[r] \ar[d]
&
Rf_*\mathcal{K}^\bullet
\otimes_{\mathcal{O}_Y}^\mathbf{L}
Rf_*\mathcal{M}^\bullet \ar[d]^{\text{Remark \ref{remark-cup-product}}} \\
\text{Tot}(
f_*\mathcal{K}^\bullet
\otimes_{\mathcal{O}_Y}
f_*\mathcal{M}^\bullet) \ar[d]_{\text{naive cup product}} &
Rf_*(\mathcal{K}^\bullet
\otimes_{\mathcal{O}_X}^\mathbf{L}
\mathcal{M}^\bullet) \ar[d] \\
f_*\text{Tot}(\mathcal{K}^\bullet
\otimes_{\mathcal{O}_X}
\mathcal{M}^\bullet) \ar[r] &
Rf_*\text{Tot}(\mathcal{K}^\bullet
\otimes_{\mathcal{O}_X}
\mathcal{M}^\bullet)
}
$$
\end{lemma}

\begin{proof}
By the construction in Remark \ref{remark-cup-product} we see that
going around the diagram clockwise the map
$$
f_*\mathcal{K}^\bullet
\otimes_{\mathcal{O}_Y}^\mathbf{L}
f_*\mathcal{M}^\bullet
\longrightarrow
Rf_*\text{Tot}(\mathcal{K}^\bullet
\otimes_{\mathcal{O}_X}
\mathcal{M}^\bullet)
$$
is adjoint to the map
\begin{align*}
Lf^*(f_*\mathcal{K}^\bullet
\otimes_{\mathcal{O}_Y}^\mathbf{L}
f_*\mathcal{M}^\bullet)
& =
Lf^*f_*\mathcal{K}^\bullet
\otimes_{\mathcal{O}_Y}^\mathbf{L}
Lf^*f_*\mathcal{M}^\bullet \\
& \to
Lf^*Rf_*\mathcal{K}^\bullet
\otimes_{\mathcal{O}_Y}^\mathbf{L}
Lf^*Rf_*\mathcal{M}^\bullet \\
& \to
\mathcal{K}^\bullet
\otimes_{\mathcal{O}_Y}^\mathbf{L}
\mathcal{M}^\bullet \\
& \to
\text{Tot}(\mathcal{K}^\bullet
\otimes_{\mathcal{O}_X}
\mathcal{M}^\bullet)
\end{align*}
By Lemma \ref{lemma-adjoints-push-pull-compatibility} this is also equal to
\begin{align*}
Lf^*(f_*\mathcal{K}^\bullet
\otimes_{\mathcal{O}_Y}^\mathbf{L}
f_*\mathcal{M}^\bullet)
& =
Lf^*f_*\mathcal{K}^\bullet
\otimes_{\mathcal{O}_Y}^\mathbf{L}
Lf^*f_*\mathcal{M}^\bullet \\
& \to
f^*f_*\mathcal{K}^\bullet
\otimes_{\mathcal{O}_Y}^\mathbf{L}
f^*f_*\mathcal{M}^\bullet \\
& \to
\mathcal{K}^\bullet
\otimes_{\mathcal{O}_Y}^\mathbf{L}
\mathcal{M}^\bullet \\
& \to
\text{Tot}(\mathcal{K}^\bullet
\otimes_{\mathcal{O}_X}
\mathcal{M}^\bullet)
\end{align*}
Going around anti-clockwise we obtain the map adjoint to the map
\begin{align*}
Lf^*(f_*\mathcal{K}^\bullet
\otimes_{\mathcal{O}_Y}^\mathbf{L}
f_*\mathcal{M}^\bullet)
& \to
Lf^*\text{Tot}(
f_*\mathcal{K}^\bullet
\otimes_{\mathcal{O}_Y}
f_*\mathcal{M}^\bullet) \\
& \to
Lf^*f_*\text{Tot}(\mathcal{K}^\bullet
\otimes_{\mathcal{O}_X}
\mathcal{M}^\bullet) \\
& \to
Lf^*Rf_*\text{Tot}(\mathcal{K}^\bullet
\otimes_{\mathcal{O}_X}
\mathcal{M}^\bullet) \\
& \to
\text{Tot}(\mathcal{K}^\bullet
\otimes_{\mathcal{O}_X}
\mathcal{M}^\bullet)
\end{align*}
By Lemma \ref{lemma-adjoints-push-pull-compatibility} this is also equal to
\begin{align*}
Lf^*(f_*\mathcal{K}^\bullet
\otimes_{\mathcal{O}_Y}^\mathbf{L}
f_*\mathcal{M}^\bullet)
& \to
Lf^*\text{Tot}(
f_*\mathcal{K}^\bullet
\otimes_{\mathcal{O}_Y}
f_*\mathcal{M}^\bullet) \\
& \to
Lf^*f_*\text{Tot}(\mathcal{K}^\bullet
\otimes_{\mathcal{O}_X}
\mathcal{M}^\bullet) \\
& \to
f^*f_*\text{Tot}(\mathcal{K}^\bullet
\otimes_{\mathcal{O}_X}
\mathcal{M}^\bullet) \\
& \to
\text{Tot}(\mathcal{K}^\bullet
\otimes_{\mathcal{O}_X}
\mathcal{M}^\bullet)
\end{align*}
Now the proof is finished by a contemplation of the diagram
$$
\xymatrix{
Lf^*(f_*\mathcal{K}^\bullet
\otimes_{\mathcal{O}_Y}^\mathbf{L}
f_*\mathcal{M}^\bullet) \ar[d] \ar[rr] & &
Lf^*f_*\mathcal{K}^\bullet \otimes_{\mathcal{O}_X}^\mathbf{L}
Lf^*f_*\mathcal{M}^\bullet \ar[d] \\
Lf^*\text{Tot}(
f_*\mathcal{K}^\bullet
\otimes_{\mathcal{O}_Y}
f_*\mathcal{M}^\bullet) \ar[d]_{naive} \ar[r] &
f^*\text{Tot}(
f_*\mathcal{K}^\bullet
\otimes_{\mathcal{O}_Y}
f_*\mathcal{M}^\bullet) \ar[ldd]^{naive} \ar[dd] &
f^*f_*\mathcal{K}^\bullet \otimes_{\mathcal{O}_X}^\mathbf{L}
f^*f_*\mathcal{M}^\bullet \ar[dd] \ar[ldd] \\
Lf^*f_*\text{Tot}(\mathcal{K}^\bullet
\otimes_{\mathcal{O}_X}
\mathcal{M}^\bullet) \ar[d] \\
f^*f_*\text{Tot}(\mathcal{K}^\bullet \otimes_{\mathcal{O}_X}
\mathcal{M}^\bullet) \ar[rd] &
\text{Tot}(f^*f_*\mathcal{K}^\bullet \otimes_{\mathcal{O}_X}
f^*f_*\mathcal{M}^\bullet) \ar[d] &
\mathcal{K}^\bullet \otimes_{\mathcal{O}_X}^\mathbf{L}
\mathcal{M}^\bullet \ar[ld] \\
& \text{Tot}(\mathcal{K}^\bullet
\otimes_{\mathcal{O}_X}
\mathcal{M}^\bullet)
}
$$
All of the polygons in this diagram commute. The top one commutes
by Lemma \ref{lemma-tensor-pull-compatibility}.
The square with the two naive cup products commutes because
$Lf^* \to f^*$ is functorial in the complex of modules.
Similarly with the square involving the two maps
$\mathcal{A}^\bullet \otimes^\mathbf{L} \mathcal{B}^\bullet \to
\text{Tot}(\mathcal{A}^\bullet \otimes \mathcal{B}^\bullet)$.
Finally, the commutativity of the remaining square
is true on the level of complexes and may be viewed as the
definition of the naive cup product (by the adjointness
of $f^*$ and $f_*$). The proof is finished because
going around the diagram on the outside are the two maps
given above.
\end{proof}

\noindent
Let $(X, \mathcal{O}_X)$ be a ring space. Let $\mathcal{K}^\bullet$ and
$\mathcal{M}^\bullet$ be complexes of $\mathcal{O}_X$-modules.
Then we have a ``naive'' cup product
$$
\mu' :
\text{Tot}(
\Gamma(X, \mathcal{K}^\bullet) \otimes_A \Gamma(X, \mathcal{M}^\bullet))
\longrightarrow
\Gamma(X, \text{Tot}(
\mathcal{K}^\bullet \otimes_{\mathcal{O}_X} \mathcal{M}^\bullet))
$$
By Lemma \ref{lemma-cup-compatible-with-naive}
applied to the morphism $(X, \mathcal{O}_X) \to (pt, A)$
this naive cup product is related to the cup product $\mu$
defined in the first paragraph of this section by the
following commutative diagram
$$
\xymatrix{
\Gamma(X, \mathcal{K}^\bullet)
\otimes_A^\mathbf{L}
\Gamma(X, \mathcal{M}^\bullet) \ar[d] \ar[r] &
R\Gamma(X, \mathcal{K}^\bullet)
\otimes_A^\mathbf{L}
R\Gamma(X, \mathcal{M}^\bullet) \ar[d]^-\mu \\
\text{Tot}(\Gamma(X, \mathcal{K}^\bullet)
\otimes_A
\Gamma(X, \mathcal{M}^\bullet)) \ar[d]_-{\mu'} &
R\Gamma(X, \mathcal{K}^\bullet
\otimes_{\mathcal{O}_X}^\mathbf{L}
\mathcal{M}^\bullet) \ar[d] \\
\Gamma(X, \text{Tot}(\mathcal{K}^\bullet
\otimes_{\mathcal{O}_X}
\mathcal{M}^\bullet)) \ar[r] &
R\Gamma(X, \text{Tot}(\mathcal{K}^\bullet
\otimes_{\mathcal{O}_X}
\mathcal{M}^\bullet))
}
$$
in $D(A)$. On cohomology we obtain the commutative diagram
$$
\xymatrix{
H^i(\Gamma(X, \mathcal{K}^\bullet)) \times
H^j(\Gamma(X, \mathcal{M}^\bullet)) \ar[d] \ar[r] &
H^{i + j}(X,
\text{Tot}(\mathcal{K}^\bullet \otimes_{\mathcal{O}_X} \mathcal{M}^\bullet)) \\
H^i(X, \mathcal{K}^\bullet) \times
H^j(X, \mathcal{M}^\bullet) \ar[r]^\cup &
H^{i + j}(X, \mathcal{K}^\bullet \otimes_{\mathcal{O}_X}^\mathbf{L}
\mathcal{M}^\bullet) \ar[u]
}
$$
relating the naive cup product with the actual cuproduct.

\begin{lemma}
\label{lemma-diagrams-commute}
Let $(X, \mathcal{O}_X)$ be a ringed space. Let
$\mathcal{K}^\bullet$ and $\mathcal{M}^\bullet$
be bounded below complexes of $\mathcal{O}_X$-modules.
Let $\mathcal{U} : X = \bigcup_{i \in I} U_i$ be an open covering
Then
$$
\xymatrix{
\text{Tot}(\check{\mathcal{C}}^\bullet(\mathcal{U}, \mathcal{K}^\bullet))
\otimes_A^\mathbf{L}
\text{Tot}(\check{\mathcal{C}}^\bullet(\mathcal{U}, \mathcal{M}^\bullet))
\ar[d] \ar[r] &
R\Gamma(X, \mathcal{K}^\bullet)
\otimes_A^\mathbf{L}
R\Gamma(X, \mathcal{M}^\bullet) \ar[d]^\mu \\
\text{Tot}(
\text{Tot}(\check{\mathcal{C}}^\bullet(\mathcal{U}, \mathcal{K}^\bullet))
\otimes_A
\text{Tot}(\check{\mathcal{C}}^\bullet(\mathcal{U}, \mathcal{M}^\bullet)))
\ar[d]^{(\ref{equation-needs-signs})} &
R\Gamma(X,
\mathcal{K}^\bullet \otimes_{\mathcal{O}_X}^\mathbf{L} \mathcal{M}^\bullet)
\ar[d] \\
\text{Tot}(
\check{\mathcal{C}}^\bullet({\mathcal U},
\text{Tot}(\mathcal{K}^\bullet \otimes_{\mathcal{O}_X} \mathcal{M}^\bullet)
)) \ar[r] &
R\Gamma(X,
\text{Tot}(\mathcal{K}^\bullet \otimes_{\mathcal{O}_X} \mathcal{M}^\bullet))
}
$$
where the horizontal arrows are the ones in
Lemma \ref{lemma-cech-complex-complex}
commutes in $D(A)$.
\end{lemma}

\begin{proof}
Choose quasi-isomorphisms of complexes
$a : \mathcal{K}^\bullet \to \mathcal{K}_1^\bullet$ and
$b : \mathcal{M}^\bullet \to \mathcal{M}_1^\bullet$
as in Lemma \ref{lemma-godement-resolution-bounded-below}.
Since the maps $a$ and $b$ on stalks are homotopy equivalences
we see that the induced map
$$
\text{Tot}(\mathcal{K}^\bullet \otimes_{\mathcal{O}_X} \mathcal{M}^\bullet)
\to
\text{Tot}(\mathcal{K}_1^\bullet \otimes_{\mathcal{O}_X} \mathcal{M}_1^\bullet)
$$
is a homotopy equivalence on stalks too (More on Algebra, Lemma
\ref{more-algebra-lemma-derived-tor-homotopy}) and hence a quasi-isomorphism.
Thus the targets
$$
R\Gamma(X,
\text{Tot}(\mathcal{K}^\bullet
\otimes_{\mathcal{O}_X} \mathcal{M}^\bullet)) =
R\Gamma(X,
\text{Tot}(\mathcal{K}_1^\bullet
\otimes_{\mathcal{O}_X} \mathcal{M}_1^\bullet))
$$
of the two diagrams are the same in $D(A)$. It follows that it suffices
to prove the diagram commutes for $\mathcal{K}$ and $\mathcal{M}$
replaced by $\mathcal{K}_1$ and $\mathcal{M}_1$. This reduces us to
the case discussed in the next paragraph.

\medskip\noindent
Assume $\mathcal{K}^\bullet$ and $\mathcal{M}^\bullet$ are bounded
below complexes of flasque $\mathcal{O}_X$-modules and
consider the diagram relating the cup product with the cup product
(\ref{equation-needs-signs}) on {\v C}ech complexes.
Then we can consider the commutative diagram
$$
\xymatrix{
\Gamma(X, \mathcal{K}^\bullet)
\otimes_A^\mathbf{L}
\Gamma(X, \mathcal{M}^\bullet) \ar[d] \ar[r] &
\text{Tot}(\check{\mathcal{C}}^\bullet(\mathcal{U}, \mathcal{K}^\bullet))
\otimes_A^\mathbf{L}
\text{Tot}(\check{\mathcal{C}}^\bullet(\mathcal{U}, \mathcal{M}^\bullet))
\ar[d] \\
\text{Tot}(\Gamma(X, \mathcal{K}^\bullet)
\otimes_A
\Gamma(X, \mathcal{M}^\bullet)) \ar[d] \ar[r] &
\text{Tot}(
\text{Tot}(\check{\mathcal{C}}^\bullet(\mathcal{U}, \mathcal{K}^\bullet))
\otimes_A
\text{Tot}(\check{\mathcal{C}}^\bullet(\mathcal{U}, \mathcal{M}^\bullet)))
\ar[d]^{(\ref{equation-needs-signs})} \\
\Gamma(X, \text{Tot}(\mathcal{K}^\bullet
\otimes_{\mathcal{O}_X}
\mathcal{M}^\bullet)) \ar[r] &
\text{Tot}(
\check{\mathcal{C}}^\bullet({\mathcal U},
\text{Tot}(\mathcal{K}^\bullet \otimes_{\mathcal{O}_X} \mathcal{M}^\bullet)
))
}
$$
In this diagram the horizontal arrows are isomorphisms in $D(A)$ because
for a bounded below complex of flasque modules such as $\mathcal{K}^\bullet$
we have
$$
\Gamma(X, \mathcal{K}^\bullet) =
\text{Tot}(\check{\mathcal{C}}^\bullet(\mathcal{U}, \mathcal{K}^\bullet)) =
R\Gamma(X, \mathcal{K}^\bullet)
$$
in $D(A)$. This follows from
Lemma \ref{lemma-flasque-acyclic},
Derived Categories, Lemma \ref{derived-lemma-leray-acyclicity}, and
Lemma \ref{lemma-cech-complex-complex-computes}.
Hence the commutativity of the diagram of the lemma involving
(\ref{equation-needs-signs}) follows from the already proven
commutativity of Lemma \ref{lemma-cup-compatible-with-naive}
where $f$ is the morphism to a point (see discussion
following Lemma \ref{lemma-cup-compatible-with-naive}).
\end{proof}

\begin{lemma}
\label{lemma-cup-product-associative}
Let $f : (X, \mathcal{O}_X) \to (Y, \mathcal{O}_Y)$
be a morphism of ringed spaces. The relative cup product of
Remark \ref{remark-cup-product} is associative in the sense that
the diagram
$$
\xymatrix{
Rf_*K \otimes_{\mathcal{O}_Y}^\mathbf{L}
Rf_*L \otimes_{\mathcal{O}_Y}^\mathbf{L}
Rf_*M \ar[r] \ar[d] &
Rf_*(K \otimes_{\mathcal{O}_X}^\mathbf{L} L)
\otimes_{\mathcal{O}_Y}^\mathbf{L} Rf_*M \ar[d] \\
Rf_*K \otimes_{\mathcal{O}_Y}^\mathbf{L}
Rf_*(L \otimes_{\mathcal{O}_X}^\mathbf{L} M) \ar[r] &
Rf_*(K \otimes_{\mathcal{O}_X}^\mathbf{L}
L \otimes_{\mathcal{O}_X}^\mathbf{L} M)
}
$$
is commutative in $D(\mathcal{O}_Y)$ for all $K, L, M$ in $D(\mathcal{O}_X)$.
\end{lemma}

\begin{proof}
Going around either side we obtain the map adjoint to the obvious map
\begin{align*}
Lf^*(Rf_*K \otimes_{\mathcal{O}_Y}^\mathbf{L}
Rf_*L \otimes_{\mathcal{O}_Y}^\mathbf{L}
Rf_*M) & =
Lf^*(Rf_*K) \otimes_{\mathcal{O}_X}^\mathbf{L}
Lf^*(Rf_*L) \otimes_{\mathcal{O}_X}^\mathbf{L}
Lf^*(Rf_*M) \\
& \to
K \otimes_{\mathcal{O}_X}^\mathbf{L}
L \otimes_{\mathcal{O}_X}^\mathbf{L} M
\end{align*}
in $D(\mathcal{O}_X)$.
\end{proof}

\begin{lemma}
\label{lemma-cup-product-commutative}
Let $f : (X, \mathcal{O}_X) \to (Y, \mathcal{O}_Y)$
be a morphism of ringed spaces. The relative cup product of
Remark \ref{remark-cup-product} is commutative in the sense that
the diagram
$$
\xymatrix{
Rf_*K \otimes_{\mathcal{O}_Y}^\mathbf{L} Rf_*L \ar[r] \ar[d]_\psi &
Rf_*(K \otimes_{\mathcal{O}_X}^\mathbf{L} L) \ar[d]^{Rf_*\psi} \\
Rf_*L \otimes_{\mathcal{O}_Y}^\mathbf{L} Rf_*K \ar[r] &
Rf_*(L \otimes_{\mathcal{O}_X}^\mathbf{L} K)
}
$$
is commutative in $D(\mathcal{O}_Y)$ for all $K, L$ in $D(\mathcal{O}_X)$.
Here $\psi$ is the commutativity constraint on the derived category
(Lemma \ref{lemma-symmetric-monoidal-derived}).
\end{lemma}

\begin{proof}
Omitted.
\end{proof}

\begin{lemma}
\label{lemma-compose-cup-product}
Let $f : (X, \mathcal{O}_X) \to (Y, \mathcal{O}_Y)$ and
$g : (Y, \mathcal{O}_Y) \to (Z, \mathcal{O}_Z)$
be morphisms of ringed spaces. The relative cup product of
Remark \ref{remark-cup-product} is compatible with compositions
in the sense that the diagram
$$
\xymatrix{
R(g \circ f)_*K \otimes_{\mathcal{O}_Z}^\mathbf{L} R(g \circ f)_*L
\ar@{=}[rr] \ar[d] & &
Rg_*Rf_*K \otimes_{\mathcal{O}_Z}^\mathbf{L} Rg_*Rf_*L \ar[d] \\
R(g \circ f)_*(K \otimes_{\mathcal{O}_X}^\mathbf{L} L) \ar@{=}[r] &
Rg_*Rf_*(K \otimes_{\mathcal{O}_X}^\mathbf{L} L) &
Rg_*(Rf_*K \otimes_{\mathcal{O}_Y}^\mathbf{L}  Rf_*L) \ar[l]
}
$$
is commutative in $D(\mathcal{O}_Z)$ for all $K, L$ in $D(\mathcal{O}_X)$.
\end{lemma}

\begin{proof}
This is true because going around the diagram either way we obtain the map
adjoint to the map
\begin{align*}
& L(g \circ f)^*\left(R(g \circ f)_*K
\otimes_{\mathcal{O}_Z}^\mathbf{L}
R(g \circ f)_*L\right) \\
& =
L(g \circ f)^*R(g \circ f)_*K
\otimes_{\mathcal{O}_X}^\mathbf{L}
L(g \circ f)^*R(g \circ f)_*L) \\
& \to
K \otimes_{\mathcal{O}_X}^\mathbf{L} L
\end{align*}
in $D(\mathcal{O}_X)$. To see this one uses that the composition
of the counits like so
$$
L(g \circ f)^*R(g \circ f)_* =
Lf^* Lg^* Rg_* Rf_*  \to
Lf^* Rf_* \to \text{id}
$$
is the counit for $L(g \circ f)^*$ and $R(g \circ f)_*$. See
Categories, Lemma \ref{categories-lemma-compose-counits}.
\end{proof}

\begin{lemma}
\label{lemma-base-change-cup-product}
Consider a commutative square
$$
\xymatrix{
(X', \mathcal{O}_{X'}) \ar[r]_{g'} \ar[d]_{f'} &
(X, \mathcal{O}_X) \ar[d]^f \\
(Y', \mathcal{O}_{Y'}) \ar[r]^g & (Y, \mathcal{O}_Y)
}
$$
of ringed spaces. Let $K, L$ in $D(\mathcal{O}_X)$.
The relative cup product
is compatible with the square in the sense that the diagram
$$
\xymatrix{
Lg^*(Rf_*K \otimes_{\mathcal{O}_Y}^\mathbf{L} Rf_*L)
\ar[r] \ar@{=}[d] &
Lg^*(Rf_*(K \otimes_{\mathcal{O}_X}^\mathbf{L} L)) \ar[d] \\
Lg^*Rf_*K \otimes_{\mathcal{O}_{Y'}}^\mathbf{L} Lg^*Rf_*L \ar[d] &
R(f')_*L(g')^*(K \otimes_{\mathcal{O}_X}^\mathbf{L} L) \ar@{=}[d] \\
R(f')_*(L(g')^*K \otimes_{\mathcal{O}_{Y'}} R(f')_*(L(g')^*L \ar[r] &
R(f')_*(L(g')^*K \otimes_{\mathcal{O}_{X'}}^\mathbf{L} L(g')^*L)
}
$$
is commutative in $D(\mathcal{O}_{Y'})$. The horizontal arrows are given
by the relative cup product (Remark \ref{remark-cup-product})
and the vertical arrows are given
by the base change map (Remark \ref{remark-base-change})
and Lemma \ref{lemma-pullback-tensor-product}.
\end{lemma}

\begin{proof}
Omitted.
\end{proof}














\section{Some properties of K-injective complexes}
\label{section-properties-K-injective}

\noindent
Let $(X, \mathcal{O}_X)$ be a ringed space. Let $U \subset X$ be an
open subset. Denote $j : (U, \mathcal{O}_U) \to (X, \mathcal{O}_X)$
the corresponding open immersion. The pullback functor $j^*$ is exact
as it is just the restriction functor. Thus derived pullback $Lj^*$ is
computed on any complex by simply restricting the complex. We often
simply denote the corresponding functor
$$
D(\mathcal{O}_X) \to D(\mathcal{O}_U),
\quad
E \mapsto j^*E = E|_U
$$
Similarly, extension by zero
$j_! : \textit{Mod}(\mathcal{O}_U) \to \textit{Mod}(\mathcal{O}_X)$
(see Sheaves, Section \ref{sheaves-section-open-immersions})
is an exact functor (Modules, Lemma \ref{modules-lemma-j-shriek-exact}).
Thus it induces a functor
$$
j_! : D(\mathcal{O}_U) \to D(\mathcal{O}_X),\quad
F \mapsto j_!F
$$
by simply applying $j_!$ to any complex representing the object $F$.

\begin{lemma}
\label{lemma-restrict-K-injective-to-open}
Let $X$ be a ringed space. Let $U \subset X$ be an open subspace.
The restriction of a K-injective complex of $\mathcal{O}_X$-modules
to $U$ is a K-injective complex of $\mathcal{O}_U$-modules.
\end{lemma}

\begin{proof}
Follows from
Derived Categories, Lemma \ref{derived-lemma-adjoint-preserve-K-injectives}
and the fact that the restriction functor has the
exact left adjoint $j_!$.
For the construction of $j_!$ see
Sheaves, Section \ref{sheaves-section-open-immersions}
and for exactness see Modules, Lemma \ref{modules-lemma-j-shriek-exact}.
\end{proof}

\begin{lemma}
\label{lemma-unbounded-cohomology-of-open}
Let $X$ be a ringed space. Let $U \subset X$ be an open subspace.
For $K$ in $D(\mathcal{O}_X)$ we have
$H^p(U, K) = H^p(U, K|_U)$.
\end{lemma}

\begin{proof}
Let $\mathcal{I}^\bullet$ be a K-injective complex of $\mathcal{O}_X$-modules
representing $K$. Then
$$
H^q(U, K) = H^q(\Gamma(U, \mathcal{I}^\bullet)) =
H^q(\Gamma(U, \mathcal{I}^\bullet|_U))
$$
by construction of cohomology. By Lemma \ref{lemma-restrict-K-injective-to-open}
the complex $\mathcal{I}^\bullet|_U$ is a K-injective complex
representing $K|_U$ and the lemma follows.
\end{proof}

\begin{lemma}
\label{lemma-sheafification-cohomology}
Let $(X, \mathcal{O}_X)$ be a ringed space. Let $K$ be an object of
$D(\mathcal{O}_X)$. The sheafification of
$$
U \mapsto H^q(U, K) = H^q(U, K|_U)
$$
is the $q$th cohomology sheaf $H^q(K)$ of $K$.
\end{lemma}

\begin{proof}
The equality $H^q(U, K) = H^q(U, K|_U)$ holds by
Lemma \ref{lemma-unbounded-cohomology-of-open}.
Choose a K-injective complex $\mathcal{I}^\bullet$ representing $K$.
Then
$$
H^q(U, K) =
\frac{\Ker(\mathcal{I}^q(U) \to \mathcal{I}^{q + 1}(U))}
{\Im(\mathcal{I}^{q - 1}(U) \to \mathcal{I}^q(U))}.
$$
by our construction of cohomology. Since
$H^q(K) = \Ker(\mathcal{I}^q \to \mathcal{I}^{q + 1})/
\Im(\mathcal{I}^{q - 1} \to \mathcal{I}^q)$ the result is clear.
\end{proof}

\begin{lemma}
\label{lemma-restrict-direct-image-open}
Let $f : (X, \mathcal{O}_X) \to (Y, \mathcal{O}_Y)$ be a morphism of ringed
spaces. Given an open subspace $V \subset Y$, set $U = f^{-1}(V)$ and denote
$g : U \to V$ the induced morphism. Then
$(Rf_*E)|_V = Rg_*(E|_U)$ for $E$ in $D(\mathcal{O}_X)$.
\end{lemma}

\begin{proof}
Represent $E$ by a K-injective complex $\mathcal{I}^\bullet$ of
$\mathcal{O}_X$-modules. Then $Rf_*(E) = f_*\mathcal{I}^\bullet$
and $Rg_*(E|_U) = g_*(\mathcal{I}^\bullet|_U)$ by
Lemma \ref{lemma-restrict-K-injective-to-open}.
Since it is clear that $(f_*\mathcal{F})|_V = g_*(\mathcal{F}|_U)$
for any sheaf $\mathcal{F}$ on $X$ the result follows.
\end{proof}

\begin{lemma}
\label{lemma-Leray-unbounded}
Let $f : X \to Y$ be a morphism of ringed spaces.
Then $R\Gamma(Y, -) \circ Rf_* = R\Gamma(X, -)$ as functors
$D(\mathcal{O}_X) \to D(\Gamma(Y, \mathcal{O}_Y))$.
More generally for $V \subset Y$ open and $U = f^{-1}(V)$
we have $R\Gamma(U, -) = R\Gamma(V, -) \circ Rf_*$.
\end{lemma}

\begin{proof}
Let $Z$ be the ringed space consisting of a singleton
space with $\Gamma(Z, \mathcal{O}_Z) = \Gamma(Y, \mathcal{O}_Y)$.
There is a canonical morphism $Y \to Z$ of ringed spaces
inducing the identification on global sections of structure sheaves.
Then $D(\mathcal{O}_Z) = D(\Gamma(Y, \mathcal{O}_Y))$.
Hence the assertion $R\Gamma(Y, -) \circ Rf_* = R\Gamma(X, -)$
follows from Lemma \ref{lemma-derived-pushforward-composition}
applied to $X \to Y \to Z$.

\medskip\noindent
The second (more general) statement follows from the first statement
after applying Lemma \ref{lemma-restrict-direct-image-open}.
\end{proof}

\begin{lemma}
\label{lemma-unbounded-describe-higher-direct-images}
Let $f : (X, \mathcal{O}_X) \to (Y, \mathcal{O}_Y)$ be a morphism of ringed
spaces. Let $K$ be in $D(\mathcal{O}_X)$. Then $H^i(Rf_*K)$ is the sheaf
associated to the presheaf
$$
V \mapsto H^i(f^{-1}(V), K) = H^i(V, Rf_*K)
$$
\end{lemma}

\begin{proof}
The equality $H^i(f^{-1}(V), K) = H^i(V, Rf_*K)$ follows upon taking
cohomology from the second statement in
Lemma \ref{lemma-Leray-unbounded}. Then the statement on sheafification
follows from Lemma \ref{lemma-sheafification-cohomology}.
\end{proof}

\begin{lemma}
\label{lemma-modules-abelian-unbounded}
Let $X$ be a ringed space. Let $K$ be an object of $D(\mathcal{O}_X)$
and denote $K_{ab}$ its image in $D(\underline{\mathbf{Z}}_X)$.
\begin{enumerate}
\item For any open $U \subset X$ there is a canonical map
$R\Gamma(U, K) \to R\Gamma(U, K_{ab})$
which is an isomorphism in $D(\textit{Ab})$.
\item Let $f : X \to Y$ be a morphism of ringed spaces.
There is a canonical map $Rf_*K \to Rf_*(K_{ab})$ which
is an isomorphism in $D(\underline{\mathbf{Z}}_Y)$.
\end{enumerate}
\end{lemma}

\begin{proof}
The map is constructed as follows. Choose a K-injective complex
$\mathcal{I}^\bullet$ representing $K$. Choose a quasi-isomorpism
$\mathcal{I}^\bullet \to \mathcal{J}^\bullet$ where $\mathcal{J}^\bullet$
is a K-injective complex of abelian groups. Then the map in (1)
is given by $\Gamma(U, \mathcal{I}^\bullet) \to \Gamma(U, \mathcal{J}^\bullet)$
and the map in (2) is given by
$f_*\mathcal{I}^\bullet \to f_*\mathcal{J}^\bullet$.
To show that these maps are isomorphisms, it suffices to prove
they induce isomorphisms on cohomology groups and cohomology sheaves.
By Lemmas \ref{lemma-unbounded-cohomology-of-open} and
\ref{lemma-unbounded-describe-higher-direct-images}
it suffices to show that the map
$$
H^0(X, K) \longrightarrow H^0(X, K_{ab})
$$
is an isomorphism. Observe that
$$
H^0(X, K) = \Hom_{D(\mathcal{O}_X)}(\mathcal{O}_X, K)
$$
and similarly for the other group. Choose any complex $\mathcal{K}^\bullet$
of $\mathcal{O}_X$-modules representing $K$. By construction of the
derived category as a localization we have
$$
\Hom_{D(\mathcal{O}_X)}(\mathcal{O}_X, K) =
\colim_{s : \mathcal{F}^\bullet \to \mathcal{O}_X}
\Hom_{K(\mathcal{O}_X)}(\mathcal{F}^\bullet, \mathcal{K}^\bullet)
$$
where the colimit is over quasi-isomorphisms $s$ of complexes of
$\mathcal{O}_X$-modules. Similarly, we have
$$
\Hom_{D(\underline{\mathbf{Z}}_X)}(\underline{\mathbf{Z}}_X, K) =
\colim_{s : \mathcal{G}^\bullet \to \underline{\mathbf{Z}}_X}
\Hom_{K(\underline{\mathbf{Z}}_X)}(\mathcal{G}^\bullet, \mathcal{K}^\bullet)
$$
Next, we observe that the quasi-isomorphisms
$s : \mathcal{G}^\bullet \to \underline{\mathbf{Z}}_X$
with $\mathcal{G}^\bullet$ bounded above complex of flat
$\underline{\mathbf{Z}}_X$-modules is cofinal in the system.
(This follows from Modules, Lemma \ref{modules-lemma-module-quotient-flat} and
Derived Categories, Lemma \ref{derived-lemma-subcategory-left-resolution};
see discussion in Section \ref{section-flat}.)
Hence we can construct an inverse to the map
$H^0(X, K) \longrightarrow H^0(X, K_{ab})$
by representing an element $\xi \in H^0(X, K_{ab})$ by a pair
$$
(s : \mathcal{G}^\bullet \to \underline{\mathbf{Z}}_X,
a : \mathcal{G}^\bullet \to \mathcal{K}^\bullet)
$$
with $\mathcal{G}^\bullet$ a bounded above complex of flat
$\underline{\mathbf{Z}}_X$-modules and sending this to
$$
(\mathcal{G}^\bullet \otimes_{\underline{\mathbf{Z}}_X} \mathcal{O}_X
\to \mathcal{O}_X,
\mathcal{G}^\bullet  \otimes_{\underline{\mathbf{Z}}_X} \mathcal{O}_X
\to \mathcal{K}^\bullet)
$$
The only thing to note here is that the first arrow
is a quasi-isomorphism by
Lemmas \ref{lemma-derived-tor-quasi-isomorphism-other-side} and
\ref{lemma-bounded-flat-K-flat}.
We omit the detailed verification that this construction
is indeed an inverse.
\end{proof}

\begin{lemma}
\label{lemma-adjoint-lower-shriek-restrict}
Let $(X, \mathcal{O}_X)$ be a ringed space. Let $U \subset X$ be an
open subset. Denote $j : (U, \mathcal{O}_U) \to (X, \mathcal{O}_X)$
the corresponding open immersion. The restriction functor
$D(\mathcal{O}_X) \to D(\mathcal{O}_U)$ is a right adjoint to
extension by zero $j_! : D(\mathcal{O}_U) \to D(\mathcal{O}_X)$.
\end{lemma}

\begin{proof}
This follows formally from the fact that $j_!$ and $j^*$ are adjoint and
exact (and hence $Lj_! = j_!$ and $Rj^* = j^*$ exist), see
Derived Categories, Lemma \ref{derived-lemma-derived-adjoint-functors}.
\end{proof}

\begin{lemma}
\label{lemma-K-injective-flat}
Let $f : X \to Y$ be a flat morphism of ringed spaces.
If $\mathcal{I}^\bullet$ is a K-injective complex of $\mathcal{O}_X$-modules,
then $f_*\mathcal{I}^\bullet$ is K-injective as a complex of
$\mathcal{O}_Y$-modules.
\end{lemma}

\begin{proof}
This is true because
$$
\Hom_{K(\mathcal{O}_Y)}(\mathcal{F}^\bullet, f_*\mathcal{I}^\bullet)
=
\Hom_{K(\mathcal{O}_X)}(f^*\mathcal{F}^\bullet, \mathcal{I}^\bullet)
$$
by
Sheaves, Lemma
\ref{sheaves-lemma-adjoint-pullback-pushforward-modules}
and the fact that $f^*$ is exact as $f$ is assumed to be flat.
\end{proof}





\section{Unbounded Mayer-Vietoris}
\label{section-unbounded-mayer-vietoris}

\noindent
There is a Mayer-Vietoris sequence for unbounded cohomology as well.

\begin{lemma}
\label{lemma-exact-sequence-lower-shriek}
Let $(X, \mathcal{O}_X)$ be a ringed space.
Let $X = U \cup V$ be the union of two open subspaces.
For any object $E$ of $D(\mathcal{O}_X)$ we have a distinguished
triangle
$$
j_{U \cap V!}E|_{U \cap V} \to
j_{U!}E|_U \oplus j_{V!}E|_V \to E \to
j_{U \cap V!}E|_{U \cap V}[1]
$$
in $D(\mathcal{O}_X)$.
\end{lemma}

\begin{proof}
We have seen in Section \ref{section-properties-K-injective}
that the restriction functors and the extension
by zero functors are computed by just applying the functors to
any complex. Let $\mathcal{E}^\bullet$ be a complex of $\mathcal{O}_X$-modules
representing $E$. The distinguished triangle of the lemma is the
distinguished triangle associated (by
Derived Categories, Section
\ref{derived-section-canonical-delta-functor} and especially
Lemma \ref{derived-lemma-derived-canonical-delta-functor})
to the short exact sequence of complexes of $\mathcal{O}_X$-modules
$$
0 \to j_{U \cap V!}\mathcal{E}^\bullet|_{U \cap V} \to
j_{U!}\mathcal{E}^\bullet|_U \oplus j_{V!}\mathcal{E}^\bullet|_V
\to \mathcal{E}^\bullet \to 0
$$
To see this sequence is exact one checks on stalks using
Sheaves, Lemma \ref{sheaves-lemma-j-shriek-modules}
(computation omitted).
\end{proof}

\begin{lemma}
\label{lemma-exact-sequence-j-star}
Let $(X, \mathcal{O}_X)$ be a ringed space.
Let $X = U \cup V$ be the union of two open subspaces.
For any object $E$ of $D(\mathcal{O}_X)$ we have a distinguished
triangle
$$
E \to
Rj_{U, *}E|_U \oplus Rj_{V, *}E|_V \to
Rj_{U \cap V, *}E|_{U \cap V} \to
E[1]
$$
in $D(\mathcal{O}_X)$.
\end{lemma}

\begin{proof}
Choose a K-injective complex $\mathcal{I}^\bullet$ representing $E$
whose terms $\mathcal{I}^n$ are injective objects of
$\textit{Mod}(\mathcal{O}_X)$, see Injectives, Theorem
\ref{injectives-theorem-K-injective-embedding-grothendieck}.
We have seen that $\mathcal{I}^\bullet|U$ is a K-injective complex
as well (Lemma \ref{lemma-restrict-K-injective-to-open}). Hence
$Rj_{U, *}E|_U$ is represented by $j_{U, *}\mathcal{I}^\bullet|_U$.
Similarly for $V$ and $U \cap V$. Hence the distinguished triangle
of the lemma is the distinguished triangle associated (by
Derived Categories, Section
\ref{derived-section-canonical-delta-functor} and especially
Lemma \ref{derived-lemma-derived-canonical-delta-functor})
to the short exact sequence of complexes
$$
0 \to
\mathcal{I}^\bullet \to
j_{U, *}\mathcal{I}^\bullet|_U \oplus j_{V, *}\mathcal{I}^\bullet|_V \to
j_{U \cap V, *}\mathcal{I}^\bullet|_{U \cap V} \to
0.
$$
This sequence is exact because for any $W \subset X$ open
and any $n$ the sequence
$$
0 \to
\mathcal{I}^n(W) \to
\mathcal{I}^n(W \cap U) \oplus \mathcal{I}^n(W \cap V) \to
\mathcal{I}^n(W \cap U \cap V) \to
0
$$
is exact (see proof of Lemma \ref{lemma-mayer-vietoris}).
\end{proof}

\begin{lemma}
\label{lemma-mayer-vietoris-hom}
Let $(X, \mathcal{O}_X)$ be a ringed space. Let $X = U \cup V$ be
the union of two open subspaces of $X$.
For objects $E$, $F$ of $D(\mathcal{O}_X)$ we have a
Mayer-Vietoris sequence
$$
\xymatrix{
& \ldots \ar[r] & \Ext^{-1}(E_{U \cap V}, F_{U \cap V}) \ar[lld] \\
\Hom(E, F) \ar[r] &
\Hom(E_U, F_U) \oplus
\Hom(E_V, F_V) \ar[r] &
\Hom(E_{U \cap V}, F_{U \cap V})
}
$$
where the subscripts denote restrictions to the relevant opens
and the $\Hom$'s and $\Ext$'s are taken in the relevant
derived categories.
\end{lemma}

\begin{proof}
Use the distinguished triangle of
Lemma \ref{lemma-exact-sequence-lower-shriek}
to obtain a long exact sequence of $\Hom$'s
(from Derived Categories, Lemma \ref{derived-lemma-representable-homological})
and use that
$$
\Hom_{D(\mathcal{O}_X)}(j_{U!}E|_U, F) =
\Hom_{D(\mathcal{O}_U)}(E|_U, F|_U)
$$
by Lemma \ref{lemma-adjoint-lower-shriek-restrict}.
\end{proof}

\begin{lemma}
\label{lemma-unbounded-mayer-vietoris}
Let $(X, \mathcal{O}_X)$ be a ringed space. Suppose that
$X = U \cup V$ is a union of two open subsets. For an object $E$
of $D(\mathcal{O}_X)$ we have a distinguished triangle
$$
R\Gamma(X, E) \to R\Gamma(U, E) \oplus R\Gamma(V, E) \to
R\Gamma(U \cap V, E) \to R\Gamma(X, E)[1]
$$
and in particular a long exact cohomology sequence
$$
\ldots \to
H^n(X, E) \to
H^n(U, E) \oplus H^0(V, E) \to
H^n(U \cap V, E) \to
H^{n + 1}(X, E) \to \ldots
$$
The construction of the distinguished triangle and the
long exact sequence is functorial in $E$.
\end{lemma}

\begin{proof}
Choose a K-injective complex $\mathcal{I}^\bullet$
representing $E$. We may assume $\mathcal{I}^n$ is an injective
object of $\textit{Mod}(\mathcal{O}_X)$ for all $n$, see
Injectives, Theorem
\ref{injectives-theorem-K-injective-embedding-grothendieck}.
Then $R\Gamma(X, E)$ is computed by $\Gamma(X, \mathcal{I}^\bullet)$.
Similarly for $U$, $V$, and $U \cap V$ by
Lemma \ref{lemma-restrict-K-injective-to-open}.
Hence the distinguished triangle of the lemma is the distinguished
triangle associated (by
Derived Categories, Section
\ref{derived-section-canonical-delta-functor} and especially
Lemma \ref{derived-lemma-derived-canonical-delta-functor})
to the short exact sequence of complexes
$$
0 \to
\mathcal{I}^\bullet(X) \to
\mathcal{I}^\bullet(U) \oplus \mathcal{I}^\bullet(V) \to
\mathcal{I}^\bullet(U \cap V) \to
0.
$$
We have seen this is a short exact sequence in the proof of
Lemma \ref{lemma-mayer-vietoris}.
The final statement follows from the functoriality of the construction
in Injectives, Theorem
\ref{injectives-theorem-K-injective-embedding-grothendieck}.
\end{proof}

\begin{lemma}
\label{lemma-unbounded-relative-mayer-vietoris}
Let $f : X \to Y$ be a morphism of ringed spaces.
Suppose that $X = U \cup V$ is a union of two open subsets.
Denote $a = f|_U : U \to Y$, $b = f|_V : V \to Y$, and
$c = f|_{U \cap V} : U \cap V \to Y$.
For every object $E$ of $D(\mathcal{O}_X)$ there exists a
distinguished triangle
$$
Rf_*E \to
Ra_*(E|_U) \oplus Rb_*(E|_V) \to
Rc_*(E|_{U \cap V}) \to
Rf_*E[1]
$$
This triangle is functorial in $E$.
\end{lemma}

\begin{proof}
Choose a K-injective complex $\mathcal{I}^\bullet$
representing $E$. We may assume $\mathcal{I}^n$ is an injective
object of $\textit{Mod}(\mathcal{O}_X)$ for all $n$, see
Injectives, Theorem
\ref{injectives-theorem-K-injective-embedding-grothendieck}.
Then $Rf_*E$ is computed by $f_*\mathcal{I}^\bullet$.
Similarly for $U$, $V$, and $U \cap V$ by
Lemma \ref{lemma-restrict-K-injective-to-open}.
Hence the distinguished triangle of the lemma is the distinguished
triangle associated (by
Derived Categories, Section
\ref{derived-section-canonical-delta-functor} and especially
Lemma \ref{derived-lemma-derived-canonical-delta-functor})
to the short exact sequence of complexes
$$
0 \to
f_*\mathcal{I}^\bullet \to
a_*\mathcal{I}^\bullet|_U \oplus b_*\mathcal{I}^\bullet|_V \to
c_*\mathcal{I}^\bullet|_{U \cap V} \to
0.
$$
This is a short exact sequence of complexes by
Lemma \ref{lemma-relative-mayer-vietoris}
and the fact that $R^1f_*\mathcal{I} = 0$
for an injective object $\mathcal{I}$ of $\textit{Mod}(\mathcal{O}_X)$.
The final statement follows from the functoriality of the construction
in Injectives, Theorem
\ref{injectives-theorem-K-injective-embedding-grothendieck}.
\end{proof}

\begin{lemma}
\label{lemma-pushforward-restriction}
Let $(X, \mathcal{O}_X)$ be a ringed space. Let $j : U \to X$ be an
open subspace. Let $T \subset X$ be a closed subset contained in $U$.
\begin{enumerate}
\item If $E$ is an object of $D(\mathcal{O}_X)$ whose cohomology sheaves
are supported on $T$, then $E \to Rj_*(E|_U)$ is an isomorphism.
\item If $F$ is an object of $D(\mathcal{O}_U)$ whose cohomology sheaves
are supported on $T$, then $j_!F \to Rj_*F$ is an isomorphism.
\end{enumerate}
\end{lemma}

\begin{proof}
Let $V = X \setminus T$ and $W = U \cap V$. Note that $X = U \cup V$ is an
open covering of $X$. Denote $j_W : W \to V$ the open immersion.
Let $E$ be an object of $D(\mathcal{O}_X)$ whose cohomology sheaves are
supported on $T$. By
Lemma \ref{lemma-restrict-direct-image-open} we have
$(Rj_*E|_U)|_V = Rj_{W, *}(E|_W) = 0$ because $E|_W = 0$ by our assumption.
On the other hand, $Rj_*(E|_U)|_U = E|_U$. Thus (1) is clear.
Let $F$ be an object of $D(\mathcal{O}_U)$ whose cohomology sheaves
are supported on $T$. By
Lemma \ref{lemma-restrict-direct-image-open} we have
$(Rj_*F)|_V = Rj_{W, *}(F|_W) = 0$ because $F|_W = 0$ by our assumption.
We also have $(j_!F)|_V = j_{W!}(F|_W) = 0$ (the first equality is immediate
from the definition of extension by zero). Since both
$(Rj_*F)|_U = F$ and $(j_!F)|_U = F$ we see that (2) holds.
\end{proof}

\begin{lemma}
\label{lemma-mayer-vietoris-cup}
Let $(X, \mathcal{O}_X)$ be a ringed space. Set $A = \Gamma(X, \mathcal{O}_X)$.
Suppose that $X = U \cup V$ is a union of two open subsets. For objects
$K$ and $M$ of $D(\mathcal{O}_X)$ we have a map of distinguished triangles
$$
\xymatrix{
R\Gamma(X, K) \otimes_A^\mathbf{L} R\Gamma(X, M) \ar[r] \ar[d] &
R\Gamma(X, K \otimes_{\mathcal{O}_X}^\mathbf{L} M) \ar[d] \\
R\Gamma(X, K) \otimes_A^\mathbf{L}
(R\Gamma(U, M) \oplus R\Gamma(V, M)) \ar[r] \ar[d] &
R\Gamma(U, K \otimes_{\mathcal{O}_X}^\mathbf{L} M)
\oplus R\Gamma(V, K \otimes_{\mathcal{O}_X}^\mathbf{L} M)) \ar[d] \\
R\Gamma(X, K) \otimes_A^\mathbf{L} R\Gamma(U \cap V, M) \ar[r] \ar[d] &
R\Gamma(U \cap V, K \otimes_{\mathcal{O}_X}^\mathbf{L} M) \ar[d] \\
R\Gamma(X, K) \otimes_A^\mathbf{L} R\Gamma(X, M)[1] \ar[r] &
R\Gamma(X, K \otimes_{\mathcal{O}_X}^\mathbf{L} M)[1]
}
$$
where
\begin{enumerate}
\item the horizontal arrows are given by cup product,
\item on the right hand side we have the distinguished triangle
of Lemma \ref{lemma-unbounded-mayer-vietoris} for
$K \otimes_{\mathcal{O}_X}^\mathbf{L} M$, and
\item on the left hand side we have the exact functor
$R\Gamma(X, K) \otimes_A^\mathbf{L} - $ applied to the
distinguished triangle of Lemma \ref{lemma-unbounded-mayer-vietoris} for $M$.
\end{enumerate}
\end{lemma}

\begin{proof}
Choose a K-flat complex $T^\bullet$ of flat $A$-modules representing
$R\Gamma(X, K)$, see
More on Algebra, Lemma \ref{more-algebra-lemma-K-flat-resolution}.
Denote $T^\bullet \otimes_A \mathcal{O}_X$ the pullback of $T^\bullet$
by the morphism of ringed spaces $(X, \mathcal{O}_X) \to (pt, A)$.
There is a natural adjunction map
$\epsilon : T^\bullet \otimes_A \mathcal{O}_X \to K$ in $D(\mathcal{O}_X)$.
Observe that $T^\bullet \otimes_A \mathcal{O}_X$ is a K-flat
complex of $\mathcal{O}_X$-modules with flat terms, see
Lemma \ref{lemma-pullback-K-flat} and
Modules, Lemma \ref{modules-lemma-pullback-flat}.
By Lemma \ref{lemma-factor-through-K-flat} we can find a morphism of complexes
$$
T^\bullet \otimes_A \mathcal{O}_X \longrightarrow \mathcal{K}^\bullet
$$
of $\mathcal{O}_X$-modules representing $\epsilon$
such that $\mathcal{K}^\bullet$ is a
K-flat complex with flat terms. Namely, by the construction of
$D(\mathcal{O}_X)$ we can first represent $\epsilon$ by some map of complexes
$e : T^\bullet \otimes_A \mathcal{O}_X \to \mathcal{L}^\bullet$
of $\mathcal{O}_X$-modules representing $\epsilon$
and then we can apply the lemma to $e$. Choose a K-injective
complex $\mathcal{I}^\bullet$ whose terms are injective $\mathcal{O}_X$-modules
representing $M$. Finally, choose a quasi-isomorphism
$$
\text{Tot}(\mathcal{K}^\bullet \otimes_\mathcal{O} \mathcal{I}^\bullet)
\longrightarrow
\mathcal{J}^\bullet
$$
into a K-injective complex whose terms are injective $\mathcal{O}_X$-modules.
Observe that source and target of this arrow represent
$K \otimes_{\mathcal{O}_X}^\mathbf{L} M$ in $D(\mathcal{O}_X)$.
At this point, for any open $W \subset X$ we obtain a map of complexes
$$
\text{Tot}(T^\bullet \otimes_A \mathcal{I}^\bullet(W))
\to
\text{Tot}(\mathcal{K}^\bullet(W) \otimes_A \mathcal{I}^\bullet(W))
\to
\mathcal{J}^\bullet(W)
$$
of $A$-modules whose composition represents the map
$$
R\Gamma(X, K) \otimes_A^\mathbf{L} R\Gamma(W, M)
\longrightarrow
R\Gamma(W, K \otimes_{\mathcal{O}_X}^\mathbf{L} M)
$$
in $D(A)$. Clearly, these maps are compatible with restriction mappings.
OK, so now we can consider the following commutative(!) diagram
of complexes of $A$-modules
$$
\xymatrix{
0 \ar[d] & 0 \ar[d] \\
\text{Tot}(T^\bullet \otimes_A \mathcal{I}^\bullet(X)) \ar[d] \ar[r] &
\mathcal{J}^\bullet(X) \ar[d] \\
\text{Tot}(T^\bullet \otimes_A
(\mathcal{I}^\bullet(U) \oplus \mathcal{I}^\bullet(V)) \ar[d] \ar[r] &
\mathcal{J}^\bullet(U) \oplus \mathcal{J}^\bullet(V) \ar[d] \\
\text{Tot}(T^\bullet \otimes_A \mathcal{I}^\bullet(U \cap V)) \ar[r] \ar[d] &
\mathcal{J}^\bullet(U \cap V) \ar[d] \\
0 & 0
}
$$
By the proof of Lemma \ref{lemma-mayer-vietoris} the columns are
exact sequences of complexes of $A$-modules (this also uses that
$\text{Tot}(T^\bullet \otimes_A -)$ transforms short exact sequences
of complexes of $A$-modules into short exact sequences as the terms
of $T^\bullet$ are flat $A$-modules). Since the distinguished triangles
of Lemma \ref{lemma-unbounded-mayer-vietoris}
are the distinguished triangles associated to these
short exact sequences of complexes, the desired result follows from
the functoriality of ``taking the associated distinguished triangle''
discussed in
Derived Categories, Section \ref{derived-section-canonical-delta-functor}.
\end{proof}












\section{Cohomology with support in a closed subset, II}
\label{section-cohomology-support-bis}

\noindent
We continue the discussion started in
Section \ref{section-cohomology-support}.

\medskip\noindent
Let $(X, \mathcal{O}_X)$ be a ringed space. Let $Z \subset X$ be
a closed subset. In this situation we can consider the functor
$\textit{Mod}(\mathcal{O}_X) \to \textit{Mod}(\mathcal{O}_X(X))$
given by $\mathcal{F} \mapsto \Gamma_Z(X, \mathcal{F})$. See
Modules, Definition \ref{modules-definition-support}
and
Modules, Lemma \ref{modules-lemma-support-section-closed}.
Using K-injective resolutions, see Section \ref{section-unbounded},
we obtain the right derived functor
$$
R\Gamma_Z(X, - ) : D(\mathcal{O}_X) \to D(\mathcal{O}_X(X))
$$
Given an object $K$ in $D(\mathcal{O}_X)$ we denote
$H^q_Z(X, K) = H^q(R\Gamma_Z(X, K))$ the cohomology module with
support in $Z$. We will see later
(Lemma \ref{lemma-sections-support-abelian-unbounded}) that this
agrees with the construction in Section \ref{section-cohomology-support}.

\medskip\noindent
For an $\mathcal{O}_X$-module $\mathcal{F}$ we can consider the
{\it subsheaf of sections with support in $Z$}, denoted
$\mathcal{H}_Z(\mathcal{F})$, defined by the rule
$$
\mathcal{H}_Z(\mathcal{F})(U) =
\{s \in \mathcal{F}(U) \mid \text{Supp}(s) \subset U \cap Z\} =
\Gamma_{Z \cap U}(U, \mathcal{F}|_U)
$$
As discussed in
Modules, Remark \ref{modules-remark-sections-support-in-closed-modules}
we may view $\mathcal{H}_Z(\mathcal{F})$ as an $\mathcal{O}_X|_Z$-module
on $Z$ and we obtain a functor
$$
\textit{Mod}(\mathcal{O}_X) \longrightarrow \textit{Mod}(\mathcal{O}_X|_Z),
\quad
\mathcal{F} \longmapsto \mathcal{H}_Z(\mathcal{F})
\text{ viewed as an }\mathcal{O}_X|_Z\text{-module on }Z
$$
This functor is left exact, but in general not exact. Exactly as above
we obtain a right derived functor
$$
R\mathcal{H}_Z : D(\mathcal{O}_X) \longrightarrow D(\mathcal{O}_X|_Z)
$$
We set $\mathcal{H}^q_Z(K) = H^q(R\mathcal{H}_Z(K))$ so that
$\mathcal{H}^0_Z(\mathcal{F}) = \mathcal{H}_Z(\mathcal{F})$
for any sheaf of $\mathcal{O}_X$-modules $\mathcal{F}$.

\begin{lemma}
\label{lemma-cohomology-with-support-sheaf-on-support}
Let $(X, \mathcal{O}_X)$ be a ringed space. Let $i : Z \to X$ be the
inclusion of a closed subset.
\begin{enumerate}
\item $R\mathcal{H}_Z : D(\mathcal{O}_X) \to D(\mathcal{O}_X|_Z)$
is right adjoint to $i_* : D(\mathcal{O}_X|_Z) \to D(\mathcal{O}_X)$.
\item For $K$ in $D(\mathcal{O}_X|_Z)$ we have $R\mathcal{H}_Z(i_*K) = K$.
\item Let $\mathcal{G}$ be a sheaf of
$\mathcal{O}_X|_Z$-modules on $Z$. Then
$\mathcal{H}^p_Z(i_*\mathcal{G}) = 0$ for $p > 0$.
\end{enumerate}
\end{lemma}

\begin{proof}
The functor $i_*$ is exact, so $i_* = Ri_* = Li_*$. Hence part (1)
of the lemma follows from
Modules, Lemma \ref{modules-lemma-adjoint-section-with-support}
and
Derived Categories, Lemma \ref{derived-lemma-derived-adjoint-functors}.
Let $K$ be as in (2). We can represent $K$ by a K-injective complex
$\mathcal{I}^\bullet$ of $\mathcal{O}_X|_Z$-modules. By
Lemma \ref{lemma-K-injective-flat}
the complex $i_*\mathcal{I}^\bullet$, which represents $i_*K$,
is a K-injective complex of $\mathcal{O}_X$-modules. Thus
$R\mathcal{H}_Z(i_*K)$ is computed by
$\mathcal{H}_Z(i_*\mathcal{I}^\bullet) = \mathcal{I}^\bullet$
which proves (2). Part (3) is a special case of (2).
\end{proof}

\noindent
Let $(X, \mathcal{O}_X)$ be a ringed space and let $Z \subset X$
be a closed subset. The category of $\mathcal{O}_X$-modules whose
support is contained in $Z$ is a Serre subcategory of the
category of all $\mathcal{O}_X$-modules, see
Homology, Definition \ref{homology-definition-serre-subcategory}
and
Modules, Lemma \ref{modules-lemma-support-section-closed}.
We denote $D_Z(\mathcal{O}_X)$
the strictly full saturated triangulated subcategory of
$D(\mathcal{O}_X)$ consisting of complexes whose cohomology sheaves
are supported on $Z$, see
Derived Categories, Section \ref{derived-section-triangulated-sub}.

\begin{lemma}
\label{lemma-complexes-with-support-on-closed}
Let $(X, \mathcal{O}_X)$ be a ringed space. Let $i : Z \to X$ be the
inclusion of a closed subset.
\begin{enumerate}
\item For $K$ in $D(\mathcal{O}_X|_Z)$ we have $i_*K$ in $D_Z(\mathcal{O}_X)$.
\item The functor $i_* : D(\mathcal{O}_X|_Z) \to D_Z(\mathcal{O}_X)$
is an equivalence with quasi-inverse
$i^{-1}|_{D_Z(\mathcal{O}_X)} = R\mathcal{H}_Z|_{D_Z(\mathcal{O}_X)}$.
\item The functor
$i_* \circ R\mathcal{H}_Z : D(\mathcal{O}_X) \to D_Z(\mathcal{O}_X)$
is right adjoint to the inclusion functor
$D_Z(\mathcal{O}_X) \to D(\mathcal{O}_X)$.
\end{enumerate}
\end{lemma}

\begin{proof}
Part (1) is immediate from the definitions. Part (3) is a formal
consequence of part (2) and
Lemma \ref{lemma-cohomology-with-support-sheaf-on-support}.
In the rest of the proof we prove part (2).

\medskip\noindent
Let us think of $i$ as the morphism of ringed spaces
$i : (Z, \mathcal{O}_X|_Z) \to (X, \mathcal{O}_X)$.
Recall that $i^*$ and $i_*$ is an adjoint pair of functors.
Since $i$ is a closed immersion, $i_*$ is exact.
Since $i^{-1}\mathcal{O}_X = \mathcal{O}_X|_Z$ is the structure
sheaf of $(Z, \mathcal{O}_X|_Z)$ we see that $i^* = i^{-1}$
is exact and we see that that $i^*i_* = i^{-1}i_*$
is isomorphic to the identify functor. See
Modules, Lemmas \ref{modules-lemma-exactness-pushforward-pullback} and
\ref{modules-lemma-i-star-exact}. Thus
$i_* : D(\mathcal{O}_X|_Z) \to D_Z(\mathcal{O}_X)$
is fully faithful and $i^{-1}$ determines
a left inverse. On the other hand, suppose that $K$ is an object of
$D_Z(\mathcal{O}_X)$ and consider the adjunction map
$K \to i_*i^{-1}K$.
Using exactness of $i_*$ and $i^{-1}$ this induces the adjunction maps
$H^n(K) \to i_*i^{-1}H^n(K)$ on cohomology sheaves. Since these cohomology
sheaves are supported on $Z$ we see these adjunction maps are isomorphisms
and we conclude that $i_* : D(\mathcal{O}_X|_Z) \to D_Z(\mathcal{O}_X)$
is an equivalence.

\medskip\noindent
To finish the proof it suffices to show that $R\mathcal{H}_Z(K) = i^{-1}K$ if
$K$ is an object of $D_Z(\mathcal{O}_X)$. To do this we can use that
$K = i_*i^{-1}K$ as we've just proved this is the case. Then
Lemma \ref{lemma-cohomology-with-support-sheaf-on-support}
tells us what we want.
\end{proof}

\begin{lemma}
\label{lemma-sections-with-support-K-injective}
Let $(X, \mathcal{O}_X)$ be a ringed space. Let $i : Z \to X$
be the inclusion of a closed subset. If $\mathcal{I}^\bullet$ is a K-injective
complex of $\mathcal{O}_X$-modules, then
$\mathcal{H}_Z(\mathcal{I}^\bullet)$ is K-injective complex of
$\mathcal{O}_X|_Z$-modules.
\end{lemma}

\begin{proof}
Since $i_* : \textit{Mod}(\mathcal{O}_X|_Z) \to \textit{Mod}(\mathcal{O}_X)$
is exact and left adjoint to $\mathcal{H}_Z$
(Modules, Lemma \ref{modules-lemma-adjoint-section-with-support})
this follows from
Derived Categories, Lemma \ref{derived-lemma-adjoint-preserve-K-injectives}.
\end{proof}

\begin{lemma}
\label{lemma-local-to-global-sections-with-support}
Let $(X, \mathcal{O}_X)$ be a ringed space. Let $i : Z \to X$ be the
inclusion of a closed subset. Then
$R\Gamma(Z, - ) \circ R\mathcal{H}_Z = R\Gamma_Z(X, - )$
as functors $D(\mathcal{O}_X) \to D(\mathcal{O}_X(X))$.
\end{lemma}

\begin{proof}
Follows from the construction of right derived functors using
K-injective resolutions, Lemma \ref{lemma-sections-with-support-K-injective},
and the fact that $\Gamma_Z(X, -) = \Gamma(Z, -) \circ \mathcal{H}_Z$.
\end{proof}

\begin{lemma}
\label{lemma-triangle-sections-with-support}
Let $(X, \mathcal{O}_X)$ be a ringed space. Let $i : Z \to X$ be the
inclusion of a closed subset. Let $U = X \setminus Z$.
There is a distinguished triangle
$$
R\Gamma_Z(X, K) \to R\Gamma(X, K) \to R\Gamma(U, K) \to
R\Gamma_Z(X, K)[1]
$$
in $D(\mathcal{O}_X(X))$ functorial for $K$ in $D(\mathcal{O}_X)$.
\end{lemma}

\begin{proof}
Choose a K-injective complex $\mathcal{I}^\bullet$ all of whose terms
are injective $\mathcal{O}_X$-modules representing $K$. See
Section \ref{section-unbounded}. Recall that $\mathcal{I}^\bullet|_U$
is a K-injective complex of $\mathcal{O}_U$-modules, see
Lemma \ref{lemma-restrict-K-injective-to-open}. Hence each
of the derived functors in the distinguished triangle is gotten
by applying the underlying functor to $\mathcal{I}^\bullet$.
Hence we find that it suffices to prove that
for an injective $\mathcal{O}_X$-module $\mathcal{I}$ we have
a short exact sequence
$$
0 \to \Gamma_Z(X, \mathcal{I}) \to \Gamma(X, \mathcal{I})
\to \Gamma(U, \mathcal{I}) \to 0
$$
This follows from Lemma \ref{lemma-injective-restriction-surjective}
and the definitions.
\end{proof}

\begin{lemma}
\label{lemma-triangle-sections-with-support-sheaves}
Let $(X, \mathcal{O}_X)$ be a ringed space. Let $i : Z \to X$ be the
inclusion of a closed subset. Denote $j : U = X \setminus Z \to X$
the inclusion of the complement. There is a distinguished triangle
$$
i_*R\mathcal{H}_Z(K) \to K \to Rj_*(K|_U) \to
i_*R\mathcal{H}_Z(K)[1]
$$
in $D(\mathcal{O}_X)$ functorial for $K$ in $D(\mathcal{O}_X)$.
\end{lemma}

\begin{proof}
Choose a K-injective complex $\mathcal{I}^\bullet$ all of whose terms
are injective $\mathcal{O}_X$-modules representing $K$. See
Section \ref{section-unbounded}. Recall that $\mathcal{I}^\bullet|_U$
is a K-injective complex of $\mathcal{O}_U$-modules, see
Lemma \ref{lemma-restrict-K-injective-to-open}. Hence each
of the derived functors in the distinguished triangle is gotten
by applying the underlying functor to $\mathcal{I}^\bullet$.
Hence it suffices to prove that
for an injective $\mathcal{O}_X$-module $\mathcal{I}$ we have
a short exact sequence
$$
0 \to i_*\mathcal{H}_Z(\mathcal{I}) \to \mathcal{I}
\to j_*(\mathcal{I}|_U) \to 0
$$
This follows from Lemma \ref{lemma-injective-restriction-surjective}
and the definitions.
\end{proof}

\begin{lemma}
\label{lemma-sections-support-in-closed-disjoint-open}
Let $(X, \mathcal{O}_X)$ be a ringed space. Let $Z \subset X$
be a closed subset. Let $j : U \to X$ be the inclusion of
an open subset with $U \cap Z = \emptyset$. Then
$R\mathcal{H}_Z(Rj_*K) = 0$ for all $K$ in $D(\mathcal{O}_U)$.
\end{lemma}

\begin{proof}
Choose a K-injective complex $\mathcal{I}^\bullet$ of $\mathcal{O}_U$-modules
representing $K$. Then $j_*\mathcal{I}^\bullet$ represents $Rj_*K$. By
Lemma \ref{lemma-K-injective-flat} the complex $j_*\mathcal{I}^\bullet$ is a
K-injective complex of $\mathcal{O}_X$-modules. Hence
$\mathcal{H}_Z(j_*\mathcal{I}^\bullet)$ represents $R\mathcal{H}_Z(Rj_*K)$.
Thus it suffices to show that $\mathcal{H}_Z(j_*\mathcal{G}) = 0$
for any abelian sheaf $\mathcal{G}$ on $U$. Thus we have to show that
a section $s$ of $j_*\mathcal{G}$ over some open $W$ which is supported
on $W \cap Z$ is zero. The support condition means that
$s|_{W \setminus W \cap Z} = 0$. Since $j_*\mathcal{G}(W) =
\mathcal{G}(U \cap W) = j_*\mathcal{G}(W \setminus W \cap Z)$
this implies that $s$ is zero as desired.
\end{proof}

\begin{lemma}
\label{lemma-sections-support-abelian-unbounded}
Let $(X, \mathcal{O}_X)$ be a ringed space. Let $Z \subset X$
be a closed subset. Let $K$ be an object of $D(\mathcal{O}_X)$
and denote $K_{ab}$ its image in $D(\underline{\mathbf{Z}}_X)$.
\begin{enumerate}
\item There is a canonical map $R\Gamma_Z(X, K) \to R\Gamma_Z(X, K_{ab})$
which is an isomorphism in $D(\textit{Ab})$.
\item There is a canonical map
$R\mathcal{H}_Z(K) \to R\mathcal{H}_Z(K_{ab})$
which is an isomorphism in $D(\underline{\mathbf{Z}}_Z)$.
\end{enumerate}
\end{lemma}

\begin{proof}
Proof of (1). The map is constructed as follows. Choose a K-injective complex
of $\mathcal{O}_X$-modules $\mathcal{I}^\bullet$ representing $K$.
Choose a quasi-isomorpism
$\mathcal{I}^\bullet \to \mathcal{J}^\bullet$ where $\mathcal{J}^\bullet$
is a K-injective complex of abelian groups. Then the map in (1)
is given by
$$
\Gamma_Z(X, \mathcal{I}^\bullet) \to \Gamma_Z(X, \mathcal{J}^\bullet)
$$
determined by the fact that $\Gamma_Z$ is a functor on abelian sheaves.
An easy check shows that the resulting map combined with the canonical
maps of Lemma \ref{lemma-modules-abelian-unbounded}
fit into a morphism of distinguished triangles
$$
\xymatrix{
R\Gamma_Z(X, K) \ar[r] \ar[d] &
R\Gamma(X, K) \ar[r] \ar[d] &
R\Gamma(U, K) \ar[d] \\
R\Gamma_Z(X, K_{ab}) \ar[r] &
R\Gamma(X, K_{ab}) \ar[r] &
R\Gamma(U, K_{ab})
}
$$
of Lemma \ref{lemma-triangle-sections-with-support}.
Since two of the three arrows are isomorphisms by the lemma cited,
we conclude by Derived Categories, Lemma
\ref{derived-lemma-third-isomorphism-triangle}.

\medskip\noindent
The proof of (2) is omitted. Hint: use the same argument with
Lemma \ref{lemma-triangle-sections-with-support-sheaves}
for the distinguished triangle.
\end{proof}

\begin{remark}
\label{remark-support-cup-product}
Let $(X, \mathcal{O}_X)$ be a ringed space. Let $i : Z \to X$
be the inclusion of a closed subset. Given $K$ and $M$ in
$D(\mathcal{O}_X)$ there is a canonical map
$$
K|_Z \otimes_{\mathcal{O}_X|_Z}^\mathbf{L} R\mathcal{H}_Z(M)
\longrightarrow
R\mathcal{H}_Z(K \otimes_{\mathcal{O}_X}^\mathbf{L} M)
$$
in $D(\mathcal{O}_X|_Z)$. Here $K|_Z = i^{-1}K$ is the restriction of
$K$ to $Z$ viewed as an object of $D(\mathcal{O}_X|_Z)$. By adjointness
of $i_*$ and $R\mathcal{H}_Z$ of
Lemma \ref{lemma-cohomology-with-support-sheaf-on-support} to construct this
map it suffices to produce a canonical map
$$
i_*\left(K|_Z \otimes_{\mathcal{O}_X|_Z}^\mathbf{L} R\mathcal{H}_Z(M)\right)
\longrightarrow
K \otimes_{\mathcal{O}_X}^\mathbf{L} M
$$
To construct this map, we choose a K-injective complex $\mathcal{I}^\bullet$
of $\mathcal{O}_X$-modules representing $M$ and a K-flat complex
$\mathcal{K}^\bullet$ of $\mathcal{O}_X$-modules representing $K$.
Observe that $\mathcal{K}^\bullet|_Z$ is a K-flat complex of
$\mathcal{O}_X|_Z$-modules representing $K|_Z$, see
Lemma \ref{lemma-pullback-K-flat}. Hence we need to produce a map
of complexes
$$
i_*\text{Tot}\left(
\mathcal{K}^\bullet|_Z \otimes_{\mathcal{O}_X|_Z}
\mathcal{H}_Z(\mathcal{I}^\bullet)\right)
\longrightarrow
\text{Tot}(\mathcal{K}^\bullet \otimes_{\mathcal{O}_X} \mathcal{I}^\bullet)
$$
of $\mathcal{O}_X$-modules. For this it suffices to produce maps
$$
i_*(\mathcal{K}^a|_Z \otimes_{\mathcal{O}_X|_Z}
\mathcal{H}_Z(\mathcal{I}^b))
\longrightarrow
\mathcal{K}^a \otimes_{\mathcal{O}_X} \mathcal{I}^b
$$
Looking at stalks (for example), we see that the left hand side of this
formula is equal to
$\mathcal{K}^a \otimes_{\mathcal{O}_X} i_*\mathcal{H}_Z(\mathcal{I}^b)$
and we can use the inclusion
$\mathcal{H}_Z(\mathcal{I}^b) \to \mathcal{I}^b$ to get our map.
\end{remark}

\begin{remark}
\label{remark-support-cup-product-global}
With notation as in Remark \ref{remark-support-cup-product}
we obtain a canonical cup product
\begin{align*}
H^a(X, K) \times H^b_Z(X, M)
& =
H^a(X, K) \times H^b(Z, R\mathcal{H}_Z(M)) \\
& \to
H^a(Z, K|_Z) \times H^b(Z, R\mathcal{H}_Z(M)) \\
& \to
H^{a + b}(Z, K|_Z \otimes_{\mathcal{O}_X|_Z}^\mathbf{L} R\mathcal{H}_Z(M)) \\
& \to
H^{a + b}(Z, R\mathcal{H}_Z(K \otimes_{\mathcal{O}_X}^\mathbf{L} M)) \\
& =
H^{a + b}_Z(X, K \otimes_{\mathcal{O}_X}^\mathbf{L} M)
\end{align*}
Here the equal signs are given by
Lemma \ref{lemma-local-to-global-sections-with-support},
the first arrow is restriction to $Z$, the second
arrow is the cup product (Section \ref{section-cup-product}),
and the third arrow is the map from Remark \ref{remark-support-cup-product}.
\end{remark}

\begin{lemma}
\label{lemma-support-cup-product}
With notation as in Remark \ref{remark-support-cup-product} the diagram
$$
\xymatrix{
H^i(X, K) \times H^j_Z(X, M) \ar[r] \ar[d] &
H^{i + j}_Z(X, K \otimes_{\mathcal{O}_X}^\mathbf{L} M) \ar[d] \\
H^i(X, K) \times H^j(X, M) \ar[r] &
H^{i + j}(X, K \otimes_{\mathcal{O}_X}^\mathbf{L} M)
}
$$
commutes where the top horizontal arrow is the cup product of
Remark \ref{remark-support-cup-product-global}.
\end{lemma}

\begin{proof}
Omitted.
\end{proof}

\begin{remark}
\label{remark-support-functorial}
Let $f : (X', \mathcal{O}_{X'}) \to (X, \mathcal{O}_X)$ be a morphism
of ringed spaces. Let $Z \subset X$ be a closed subset and $Z' = f^{-1}(Z)$.
Denote $f|_{Z'} : (Z', \mathcal{O}_{X'}|_{Z'}) \to (Z, \mathcal{O}_X|Z)$ be the
induced morphism of ringed spaces. For any $K$ in $D(\mathcal{O}_X)$ there
is a canonical map
$$
L(f|_{Z'})^*R\mathcal{H}_Z(K) \longrightarrow R\mathcal{H}_{Z'}(Lf^*K)
$$
in $D(\mathcal{O}_{X'}|_{Z'})$. Denote $i : Z \to X$ and $i' : Z' \to X'$
the inclusion maps. By
Lemma \ref{lemma-complexes-with-support-on-closed} part (2)
applied to $i'$ it is the same thing to give a map
$$
i'_* L(f|_{Z'})^* R\mathcal{H}_Z(K)
\longrightarrow
i'_*R\mathcal{H}_{Z'}(Lf^*K)
$$
in $D_{Z'}(\mathcal{O}_{X'})$. The map of functors
$Lf^* \circ i_* \to i'_* \circ L(f|_{Z'})^*$ of
Remark \ref{remark-base-change} is an isomorphism in this case
(follows by checking what happens on stalks using that $i_*$ and $i'_*$
are exact and that $\mathcal{O}_{Z, z} = \mathcal{O}_{X, z}$
and similarly for $Z'$). Hence it suffices to construct a the top
horizontal arrow in the following diagram
$$
\xymatrix{
Lf^* i_* R\mathcal{H}_Z(K) \ar[rr] \ar[rd] & &
i'_* R\mathcal{H}_{Z'}(Lf^*K) \ar[ld] \\
& Lf^*K
}
$$
The complex $Lf^* i_* R\mathcal{H}_Z(K)$ is supported on $Z'$. The south-east
arrow comes from the adjunction mapping $i_*R\mathcal{H}_Z(K) \to K$
(Lemma \ref{lemma-cohomology-with-support-sheaf-on-support}). Since the
adjunction mapping $i'_* R\mathcal{H}_{Z'}(Lf^*K) \to Lf^*K$ is universal by
Lemma \ref{lemma-complexes-with-support-on-closed} part (3), we find that
the south-east arrow factors uniquely over the south-west arrow and
we obtain the desired arrow.
\end{remark}

\begin{lemma}
\label{lemma-support-functorial}
With notation and assumptions as in Remark \ref{remark-support-functorial}
the diagram
$$
\xymatrix{
H^p_Z(X, K) \ar[r] \ar[d] & H^p_{Z'}(X', Lf^*K) \ar[d] \\
H^p(X, K) \ar[r] & H^p(X', Lf^*K)
}
$$
commutes. Here the top horizontal arrow comes from the identifications
$H^p_Z(X, K) = H^p(Z, R\mathcal{H}_Z(K))$ and
$H^p_{Z'}(X', Lf^*K) = H^p(Z', R\mathcal{H}_{Z'}(K'))$,
the pullback map
$H^p(Z, R\mathcal{H}_Z(K)) \to H^p(Z', L(f|_{Z'})^*R\mathcal{H}_Z(K))$,
and the map constructed in Remark \ref{remark-support-functorial}.
\end{lemma}

\begin{proof}
Omitted. Hints:
Using that $H^p(Z, R\mathcal{H}_Z(K)) = H^p(X, i_*R\mathcal{H}_Z(K))$
and similarly for $R\mathcal{H}_{Z'}(Lf^*K)$ this follows from
the functoriality of the pullback maps and the commutative diagram
used to define the map of Remark \ref{remark-support-functorial}.
\end{proof}






\section{Inverse systems and cohomology, I}
\label{section-inverse-systems}

\noindent
Let $A$ be a ring and let $I \subset A$ be an ideal.
We prove some results on inverse systems of sheaves of $A/I^n$-modules.

\begin{lemma}
\label{lemma-ML-general}
Let $I$ be an ideal of a ring $A$. Let $X$ be a topological space.
Let
$$
\ldots \to \mathcal{F}_3 \to \mathcal{F}_2 \to \mathcal{F}_1
$$
be an inverse system of sheaves of $A$-modules on $X$
such that $\mathcal{F}_n = \mathcal{F}_{n + 1}/I^n\mathcal{F}_{n + 1}$.
Let $p \geq 0$. Assume
$$
\bigoplus\nolimits_{n \geq 0} H^{p + 1}(X, I^n\mathcal{F}_{n + 1})
$$
satisfies the ascending chain condition as a graded
$\bigoplus_{n \geq 0} I^n/I^{n + 1}$-module.
Then the inverse system $M_n = H^p(X, \mathcal{F}_n)$ satisfies the
Mittag-Leffler condition\footnote{In fact, there exists
a $c \geq 0$ such that $\Im(M_n \to M_{n - c})$ is the stable image
for all $n \geq c$.}.
\end{lemma}

\begin{proof}
Set $N_n = H^{p + 1}(X, I^n\mathcal{F}_{n + 1})$ and let
$\delta_n : M_n \to N_n$ be the boundary map on cohomology
coming from the short exact sequence
$0 \to I^n\mathcal{F}_{n + 1} \to \mathcal{F}_{n + 1} \to \mathcal{F}_n \to 0$.
Then $\bigoplus \Im(\delta_n) \subset \bigoplus N_n$ is a graded submodule.
Namely, if $s \in M_n$ and $f \in I^m$, then we have a commutative diagram
$$
\xymatrix{
0 \ar[r] &
I^n\mathcal{F}_{n + 1} \ar[d]_f \ar[r] &
\mathcal{F}_{n + 1} \ar[d]_f \ar[r] &
\mathcal{F}_n \ar[d]_f \ar[r] & 0 \\
0 \ar[r] &
I^{n + m}\mathcal{F}_{n + m + 1} \ar[r] &
\mathcal{F}_{n + m + 1} \ar[r] &
\mathcal{F}_{n + m} \ar[r] & 0
}
$$
The middle vertical map is given by lifting a local section of
$\mathcal{F}_{n + 1}$ to a section of $\mathcal{F}_{n + m + 1}$
and then multiplying by $f$; similarly for the other vertical arrows.
We conclude that $\delta_{n + m}(fs) = f \delta_n(s)$.
By assumption we can find $s_j \in M_{n_j}$, $j = 1, \ldots, N$
such that $\delta_{n_j}(s_j)$
generate $\bigoplus \Im(\delta_n)$ as a graded module. Let $n > c = \max(n_j)$.
Let $s \in M_n$. Then we can find $f_j \in I^{n - n_j}$ such that
$\delta_n(s) = \sum f_j \delta_{n_j}(s_j)$. We conclude that
$\delta(s - \sum f_j s_j) = 0$, i.e., we can find $s' \in M_{n + 1}$
mapping to $s - \sum f_js_j$ in $M_n$. It follows that
$$
\Im(M_{n + 1} \to M_{n - c}) = \Im(M_n \to M_{n - c})
$$
Namely, the elements $f_js_j$ map to zero in $M_{n - c}$.
This proves the lemma.
\end{proof}

\begin{lemma}
\label{lemma-ML-general-better}
Let $I$ be an ideal of a ring $A$. Let $X$ be a topological space. Let
$$
\ldots \to \mathcal{F}_3 \to \mathcal{F}_2 \to \mathcal{F}_1
$$
be an inverse system of $A$-modules on $X$
such that $\mathcal{F}_n = \mathcal{F}_{n + 1}/I^n\mathcal{F}_{n + 1}$.
Let $p \geq 0$. Given $n$ define
$$
N_n =
\bigcap\nolimits_{m \geq n}
\Im\left(
H^{p + 1}(X, I^n\mathcal{F}_{m + 1}) \to H^{p + 1}(X, I^n\mathcal{F}_{n + 1})
\right)
$$
If $\bigoplus N_n$ satisfies the ascending chain condition as a graded
$\bigoplus_{n \geq 0} I^n/I^{n + 1}$-module, then the inverse system
$M_n = H^p(X, \mathcal{F}_n)$ satisfies the Mittag-Leffler
condition\footnote{In fact, there exists
a $c \geq 0$ such that $\Im(M_n \to M_{n - c})$ is the stable image
for all $n \geq c$.}.
\end{lemma}

\begin{proof}
The proof is exactly the same as the proof of Lemma \ref{lemma-ML-general}.
In fact, the result will follow from the arguments given there
as soon as we show that
$\bigoplus N_n$ is a graded $\bigoplus_{n \geq 0} I^n/I^{n + 1}$-submodule
of $\bigoplus H^{p + 1}(X, I^n\mathcal{F}_{n + 1})$
and that the boundary maps
$\delta_n : M_n \to H^{p + 1}(X, I^n\mathcal{F}_{n + 1})$
have image contained in $N_n$.

\medskip\noindent
Suppose that $\xi \in N_n$ and $f \in I^k$.
Choose $m \gg n + k$. Choose
$\xi' \in H^{p + 1}(X, I^n\mathcal{F}_{m + 1})$ lifting
$\xi$. We consider the diagram
$$
\xymatrix{
0 \ar[r] &
I^n\mathcal{F}_{m + 1} \ar[d]_f \ar[r] &
\mathcal{F}_{m + 1} \ar[d]_f \ar[r] &
\mathcal{F}_n \ar[d]_f \ar[r] & 0 \\
0 \ar[r] &
I^{n + k}\mathcal{F}_{m + 1} \ar[r] &
\mathcal{F}_{m + 1} \ar[r] &
\mathcal{F}_{n + k} \ar[r] & 0
}
$$
constructed as in the proof of Lemma \ref{lemma-ML-general}.
We get an induced map on cohomology and we see that
$f \xi' \in H^{p + 1}(X, I^{n + k}\mathcal{F}_{m + 1})$
maps to $f \xi$. Since this is true for all $m \gg n + k$
we see that $f\xi$ is in $N_{n + k}$ as desired.

\medskip\noindent
To see the boundary maps $\delta_n$ have image contained in $N_n$
we consider the diagrams
$$
\xymatrix{
0 \ar[r] &
I^n\mathcal{F}_{m + 1} \ar[d] \ar[r] &
\mathcal{F}_{m + 1} \ar[d] \ar[r] &
\mathcal{F}_n \ar[d] \ar[r] & 0 \\
0 \ar[r] &
I^n\mathcal{F}_{n + 1} \ar[r] &
\mathcal{F}_{n + 1} \ar[r] &
\mathcal{F}_n \ar[r] & 0
}
$$
for $m \geq n$. Looking at the induced maps on cohomology we conclude.
\end{proof}

\begin{lemma}
\label{lemma-topology-I-adic-general}
Let $I$ be an ideal of a ring $A$. Let $X$ be a topological space. Let
$$
\ldots \to \mathcal{F}_3 \to \mathcal{F}_2 \to \mathcal{F}_1
$$
be an inverse system of sheaves of $A$-modules on $X$ such that
$\mathcal{F}_n = \mathcal{F}_{n + 1}/I^n\mathcal{F}_{n + 1}$.
Let $p \geq 0$. Assume
$$
\bigoplus\nolimits_{n \geq 0} H^p(X, I^n\mathcal{F}_{n + 1})
$$
satisfies the ascending chain condition as a graded
$\bigoplus_{n \geq 0} I^n/I^{n + 1}$-module.
Then the limit topology on $M = \lim H^p(X, \mathcal{F}_n)$
is the $I$-adic topology.
\end{lemma}

\begin{proof}
Set $F^n = \Ker(M \to H^p(X, \mathcal{F}_n))$ for $n \geq 1$ and $F^0 = M$.
Observe that $I F^n \subset F^{n + 1}$. In particular $I^n M \subset F^n$.
Hence the $I$-adic topology is finer than the limit topology. For
the converse, we will show that given $n$
there exists an $m \geq n$ such that $F^m \subset I^nM$\footnote{In fact,
there exist a $c \geq 0$ such that $F^{n + c} \subset I^nM$ for all $n$.}.
We have injective maps
$$
F^n/F^{n + 1} \longrightarrow H^p(X, \mathcal{F}_{n + 1})
$$
whose image is contained in the image of
$H^p(X, I^n\mathcal{F}_{n + 1}) \to H^p(X, \mathcal{F}_{n + 1})$.
Denote
$$
E_n \subset H^p(X, I^n\mathcal{F}_{n + 1})
$$
the inverse image of $F^n/F^{n + 1}$. Then $\bigoplus E_n$ is
a graded $\bigoplus I^n/I^{n + 1}$-submodule of
$\bigoplus H^p(X, I^n\mathcal{F}_{n + 1})$ and
$\bigoplus E_n \to \bigoplus F^n/F^{n + 1}$ is a homomorphism of graded
modules; details omitted. By assumption $\bigoplus E_n$ is generated by
finitely many homogeneous elements over $\bigoplus I^n/I^{n + 1}$.
Since $E_n \to F^n/F^{n + 1}$ is surjective, we see that
the same thing is true of $\bigoplus F^n/F^{n + 1}$.
Hence we can find $r$ and $c_1, \ldots, c_r \geq 0$ and
$a_i \in F^{c_i}$ whose images in $\bigoplus F^n/F^{n + 1}$ generate.
Set $c = \max(c_i)$.

\medskip\noindent
For $n \geq c$ we claim that $I F^n = F^{n + 1}$. The claim shows that
$F^{n + c} = I^nF^c \subset I^nM$ as desired. To prove the claim
suppose $a \in F^{n + 1}$. The image of
$a$ in $F^{n + 1}/F^{n + 2}$ is a linear combination
of our $a_i$. Therefore $a  - \sum f_i a_i \in F^{n + 2}$
for some $f_i \in I^{n + 1 - c_i}$. Since
$I^{n + 1 - c_i} = I \cdot I^{n - c_i}$ as $n \geq c_i$ we can write
$f_i = \sum g_{i, j} h_{i, j}$ with $g_{i, j} \in I$
and $h_{i, j}a_i \in F^n$. Thus we see that
$F^{n + 1} = F^{n + 2} + IF^n$.
A simple induction argument gives $F^{n + 1} = F^{n + e} + IF^n$
for all $e > 0$. It follows that $IF^n$ is dense in $F^{n + 1}$.
Choose generators $k_1, \ldots, k_r$ of $I$ and consider
the continuous map
$$
u : (F^n)^{\oplus r} \longrightarrow F^{n + 1},\quad
(x_1, \ldots, x_r) \mapsto \sum k_i x_i
$$
(in the limit topology).
By the above the image of $(F^m)^{\oplus r}$ under $u$ is dense in
$F^{m + 1}$ for all $m \geq n$. By the open mapping lemma
(More on Algebra, Lemma \ref{more-algebra-lemma-open-mapping}) we find
that $u$ is open. Hence $u$ is surjective. Hence $IF^n = F^{n + 1}$
for $n \geq c$. This concludes the proof.
\end{proof}

\begin{lemma}
\label{lemma-topology-I-adic-general-better}
Let $I$ be an ideal of a ring $A$. Let $X$ be a topological space. Let
$$
\ldots \to \mathcal{F}_3 \to \mathcal{F}_2 \to \mathcal{F}_1
$$
be an inverse system of sheaves of $A$-modules on $X$ such that
$\mathcal{F}_n = \mathcal{F}_{n + 1}/I^n\mathcal{F}_{n + 1}$.
Let $p \geq 0$. Given $n$ define
$$
N_n =
\bigcap\nolimits_{m \geq n}
\Im\left(
H^p(X, I^n\mathcal{F}_{m + 1}) \to H^p(X, I^n\mathcal{F}_{n + 1})
\right)
$$
If $\bigoplus N_n$ satisfies the ascending chain condition as a graded
$\bigoplus_{n \geq 0} I^n/I^{n + 1}$-module, then
the limit topology on $M = \lim H^p(X, \mathcal{F}_n)$
is the $I$-adic topology.
\end{lemma}

\begin{proof}
The proof is exactly the same as the proof of
Lemma \ref{lemma-topology-I-adic-general}.
In fact, the result will follow from the arguments given there
as soon as we show that
$\bigoplus N_n$ is a graded $\bigoplus_{n \geq 0} I^n/I^{n + 1}$-submodule
of $\bigoplus H^{p + 1}(X, I^n\mathcal{F}_{n + 1})$
and that $F^n/F^{n + 1} \subset H^p(X, \mathcal{F}_{n + 1})$
is contained in the image of $N_n \to H^p(X, \mathcal{F}_{n + 1})$.
In the proof of Lemma \ref{lemma-ML-general-better}
we have seen the statement on the module structure.

\medskip\noindent
Let $t \in F^n$. Choose an element $s \in H^p(X, I^n\mathcal{F}_{n + 1})$
which maps to the image of $t$ in $H^p(X, \mathcal{F}_{n + 1})$. We have
to show that $s$ is in $N_n$. Now $F^n$ is the kernel of the map from
$M \to H^p(X, \mathcal{F}_n)$ hence for all $m \geq n$ we can map $t$
to an element $t_m \in H^p(X, \mathcal{F}_{m + 1})$ which maps to zero
in $H^p(X, \mathcal{F}_n)$. Consider the cohomology sequence
$$
H^{p - 1}(X, \mathcal{F}_n) \to
H^p(X, I^n\mathcal{F}_{m + 1}) \to
H^p(X, \mathcal{F}_{m + 1}) \to
H^p(X, \mathcal{F}_n)
$$
coming from the short exact sequence
$0 \to I^n\mathcal{F}_{m + 1} \to \mathcal{F}_{m + 1} \to \mathcal{F}_n \to 0$.
We can choose $s_m \in H^p(X, I^n\mathcal{F}_{m + 1})$ mapping to $t_m$.
Comparing the sequence above with the one for $m = n$ we see that
$s_m$ maps to $s$ up to an element in the image of
$H^{p - 1}(X, \mathcal{F}_n) \to H^p(X, I^n\mathcal{F}_{n + 1})$.
However, this map factors through the map
$H^p(X, I^n\mathcal{F}_{m + 1}) \to H^p(X, I^n\mathcal{F}_{n + 1})$
and we see that $s$ is in the image as desired.
\end{proof}





\section{Inverse systems and cohomology, II}
\label{section-inverse-systems-bis}

\noindent
This section continues the discussion in
Section \ref{section-inverse-systems}
in the setting where the ideal is principal.

\begin{lemma}
\label{lemma-equivalent-f-good}
Let $(X, \mathcal{O}_X)$ be a ringed space.
Let $f \in \Gamma(X, \mathcal{O}_X)$. Let
$$
\ldots \to \mathcal{F}_3 \to \mathcal{F}_2 \to \mathcal{F}_1
$$
be inverse system of $\mathcal{O}_X$-modules. Consider the conditions
\begin{enumerate}
\item for all $n \geq 1$ the map
$f : \mathcal{F}_{n + 1} \to \mathcal{F}_{n + 1}$ factors
through $\mathcal{F}_{n + 1} \to \mathcal{F}_n$ to give a
short exact sequence
$0 \to \mathcal{F}_n \to \mathcal{F}_{n + 1} \to \mathcal{F}_1 \to 0$,
\item for all $n \geq 1$ the map
$f^n : \mathcal{F}_{n + 1} \to \mathcal{F}_{n + 1}$
factors through $\mathcal{F}_{n + 1} \to \mathcal{F}_1$
to give a short exact sequence
$0 \to \mathcal{F}_1 \to \mathcal{F}_{n + 1} \to \mathcal{F}_n \to 0$
\item there exists an $\mathcal{O}_X$-module $\mathcal{G}$
which is $f$-divisible such that $\mathcal{F}_n = \mathcal{G}[f^n]$, and
\item there exists an $\mathcal{O}_X$-module $\mathcal{F}$
which is $f$-torsion free such that
$\mathcal{F}_n = \mathcal{F}/f^n\mathcal{F}$.
\end{enumerate}
Then (4) $\Rightarrow$ (3) $\Leftrightarrow$ (2) $\Leftrightarrow$ (1).
\end{lemma}

\begin{proof}
We omit the proof of the equivalence of (1) and (2). We omit the proof
that (3) implies (1). Given $\mathcal{F}_n$ as in (1) to prove (3) we set
$\mathcal{G} = \colim \mathcal{F}_n$ where the maps
$\mathcal{F}_1 \to \mathcal{F}_2 \to \mathcal{F}_3 \to \ldots$
are as in (1). The map $f : \mathcal{G} \to \mathcal{G}$
is surjective as the image of $\mathcal{F}_{n + 1} \subset \mathcal{G}$
is $\mathcal{F}_n \subset \mathcal{G}$ by the short exact sequence (1).
Thus $\mathcal{G}$ is an $f$-divisible $\mathcal{O}_X$-module
with $\mathcal{F}_n = \mathcal{G}[f^n]$.

\medskip\noindent
Assume given $\mathcal{F}$ as in (4). The map
$\mathcal{F}/f^{n + 1}\mathcal{F} \to \mathcal{F}/f^n\mathcal{F}$
is always surjective with kernel the image of the map
$\mathcal{F}/f\mathcal{F} \to \mathcal{F}/f^{n + 1}\mathcal{F}$
induced by multiplication with $f^n$. To verify (2)
it suffices to see that the kernel of
$f^n : \mathcal{F} \to \mathcal{F}/f^{n + 1}\mathcal{F}$
is $f\mathcal{F}$. To see this it suffices to show that given
sections $s, t$ of $\mathcal{F}$ over an open $U \subset X$
with $f^ns = f^{n + 1}t$ we have $s = ft$. This is clear because
$f : \mathcal{F} \to \mathcal{F}$ is injective as $\mathcal{F}$
is $f$-torsion free.
\end{proof}

\begin{lemma}
\label{lemma-topology-I-adic-f}
Suppose $X$, $f$, $(\mathcal{F}_n)$ satisfy condition (1) of
Lemma \ref{lemma-equivalent-f-good}. Let $p \geq 0$ and set
$H^p = \lim H^p(X, \mathcal{F}_n)$.
Then $f^cH^p$ is the kernel of $H^p \to H^p(X, \mathcal{F}_c)$ for all
$c \geq 1$. Thus the limit topology on $H^p$ is the $f$-adic topology.
\end{lemma}

\begin{proof}
Let $c \geq 1$. It is clear that $f^c H^p$ maps to zero in
$H^p(X, \mathcal{F}_c)$. If $\xi = (\xi_n) \in H^p$ is small in the
limit topology, then $\xi_c = 0$, and hence $\xi_n$
maps to zero in $H^p(X, \mathcal{F}_c)$ for $n \geq c$.
Consider the inverse system of short exact sequences
$$
0 \to \mathcal{F}_{n - c} \xrightarrow{f^c} \mathcal{F}_n \to
\mathcal{F}_c \to 0
$$
and the corresponding inverse system of long exact cohomology sequences
$$
H^{p - 1}(X, \mathcal{F}_c) \to
H^p(X, \mathcal{F}_{n - c}) \to
H^p(X, \mathcal{F}_n) \to
H^p(X, \mathcal{F}_c)
$$
Since the term $H^{p - 1}(X, \mathcal{F}_c)$ is independent of
$n$ we can choose a compatible sequence of elements
$\xi'_n \in H^p(X, \mathcal{F}_{n - c})$
lifting $\xi_n$. Setting $\xi' = (\xi'_n)$ we see that $\xi = f^c \xi'$
as desired.
\end{proof}

\begin{lemma}
\label{lemma-limit-finite}
Let $A$ be a Noetherian ring complete with respect to a principal ideal $(f)$.
Let $X$ be a topological space. Let
$$
\ldots \to \mathcal{F}_3 \to \mathcal{F}_2 \to \mathcal{F}_1
$$
be an inverse system of sheaves of $A$-modules. Assume
\begin{enumerate}
\item $\Gamma(X, \mathcal{F}_1)$ is a finite $A$-module,
\item $X$, $f$, $(\mathcal{F}_n)$ satisfy condition (1) of
Lemma \ref{lemma-equivalent-f-good}.
\end{enumerate}
Then
$$
M = \lim \Gamma(X, \mathcal{F}_n)
$$
is a finite $A$-module, $f$ is a nonzerodivisor on $M$, and
$M/fM$ is the image of $M$ in $\Gamma(X, \mathcal{F}_1)$.
\end{lemma}

\begin{proof}
By Lemma \ref{lemma-topology-I-adic-f} we have
$M/fM \subset H^0(X, \mathcal{F}_1)$. From (1) and the Noetherian
property of $A$ we get that $M/fM$ is a finite $A$-module.
Observe that $\bigcap f^nM = 0$ as $f^nM$ maps to zero in
$H^0(X, \mathcal{F}_n)$. By
Algebra, Lemma \ref{algebra-lemma-finite-over-complete-ring}
we conclude that $M$ is finite over $A$.
Finally, suppose $s = (s_n) \in M = \lim H^0(X, \mathcal{F}_n)$
satisfies $fs = 0$. Then $s_{n + 1}$ is in the kernel of
$\mathcal{F}_{n + 1} \to \mathcal{F}_n$ by condition (1) of
Lemma \ref{lemma-equivalent-f-good}.
Hence $s_n = 0$. Since $n$ was arbitrary, we see $s = 0$.
Thus $f$ is a nonzerodivisor on $M$.
\end{proof}

\begin{lemma}
\label{lemma-ML}
Let $A$ be a ring. Let $f \in A$. Let $X$ be a topological space. Let
$$
\ldots \to \mathcal{F}_3 \to \mathcal{F}_2 \to \mathcal{F}_1
$$
be an inverse system of sheaves of $A$-modules. Let $p \geq 0$. Assume
\begin{enumerate}
\item either $H^{p + 1}(X, \mathcal{F}_1)$ is an $A$-module of finite length
or $A$ is Noetherian and $H^{p + 1}(X, \mathcal{F}_1)$ is a finite $A$-module,
\item $X$, $f$, $(\mathcal{F}_n)$ satisfy condition (1) of
Lemma \ref{lemma-equivalent-f-good}.
\end{enumerate}
Then the inverse system $M_n = H^p(X, \mathcal{F}_n)$ satisfies the
Mittag-Leffler condition.
\end{lemma}

\begin{proof}
Set $I = (f)$. We will use the criterion of Lemma \ref{lemma-ML-general}.
Observe that $f^n : \mathcal{F}_1 \to I^n\mathcal{F}_{n + 1}$
is an isomorphism for all $n \geq 0$.
Thus it suffices to show that
$$
\bigoplus\nolimits_{n \geq 1} H^{p + 1}(X, \mathcal{F}_1) \cdot f^{n + 1}
$$
is a graded $S = \bigoplus_{n \geq 0} A/(f) \cdot f^n$-module satisfying the
ascending chain condition. If $A$ is not Noetherian, then
$H^{p + 1}(X, \mathcal{F}_1)$ has finite length and the result holds.
If $A$ is Noetherian, then $S$ is a Noetherian ring and the result
holds as the module is finite over $S$ by the assumed finiteness
of $H^{p + 1}(X, \mathcal{F}_1)$. Some details omitted.
\end{proof}

\begin{lemma}
\label{lemma-ML-better}
Let $A$ be a ring. Let $f \in A$. Let $X$ be a topological space. Let
$$
\ldots \to \mathcal{F}_3 \to \mathcal{F}_2 \to \mathcal{F}_1
$$
be an inverse system of sheaves of $A$-modules. Let $p \geq 0$. Assume
\begin{enumerate}
\item either there is an $m \geq 1$ such that the image of
$H^{p + 1}(X, \mathcal{F}_m) \to H^{p + 1}(X, \mathcal{F}_1)$
is an $A$-module of finite length or $A$ is Noetherian
and the intersection of the images of
$H^{p + 1}(X, \mathcal{F}_m) \to H^{p + 1}(X, \mathcal{F}_1)$
is a finite $A$-module,
\item $X$, $f$, $(\mathcal{F}_n)$ satisfy condition (1) of
Lemma \ref{lemma-equivalent-f-good}.
\end{enumerate}
Then the inverse system $M_n = H^p(X, \mathcal{F}_n)$ satisfies the
Mittag-Leffler condition.
\end{lemma}

\begin{proof}
Set $I = (f)$. We will use the criterion of Lemma \ref{lemma-ML-general-better}
involving the modules $N_n$. For $m \geq n$ we have
$I^n\mathcal{F}_{m + 1} = \mathcal{F}_{m + 1 - n}$. Thus we see that
$$
N_n = \bigcap\nolimits_{m \geq 1} \Im\left(
H^{p + 1}(X, \mathcal{F}_m) \to H^{p + 1}(X, \mathcal{F}_1)
\right)
$$
is independent of $n$ and
$\bigoplus N_n = \bigoplus N_1 \cdot f^{n + 1}$.
Thus we conclude exactly as in the proof of Lemma \ref{lemma-ML}.
\end{proof}

\begin{remark}
\label{remark-compare-ML}
Let $(X, \mathcal{O}_X)$ be a ringed space.
Let $f \in \Gamma(X, \mathcal{O}_X)$.
Let $\mathcal{F}$ be $\mathcal{O}_X$-module.
If $\mathcal{F}$ is $f$-torsion free, then for every $p \geq 0$
we have a short exact sequence of inverse systems
$$
0 \to \{H^p(X, \mathcal{F})/f^nH^p(X, \mathcal{F})\}
\to \{H^p(X, \mathcal{F}/f^n\mathcal{F})\}
\to \{H^{p + 1}(X, \mathcal{F})[f^n]\} \to 0
$$
Since the first inverse system has the Mittag-Leffler condition (ML)
we learn three things from this:
\begin{enumerate}
\item There is a short exact sequence
$$
0 \to \widehat{H^p(X, \mathcal{F})} \to
\lim H^p(X, \mathcal{F}/f^n\mathcal{F}) \to
T_f(H^{p + 1}(X, \mathcal{F})) \to 0
$$
where $\widehat{\ }$ denotes the usual $f$-adic completion and
$T_f( - )$ denotes the $f$-adic Tate module from
More on Algebra, Example
\ref{more-algebra-example-spectral-sequence-principal}.
\item We have
$R^1\lim H^p(X, \mathcal{F}/f^n\mathcal{F}) =
R^1\lim H^{p + 1}(X, \mathcal{F})[f^n]$.
\item The system $\{H^{p + 1}(X, \mathcal{F})[f^n]\}$ is ML if and only if
$\{H^p(X, \mathcal{F}/f^n\mathcal{F})\}$ is ML.
\end{enumerate}
See Homology, Lemma \ref{homology-lemma-Mittag-Leffler} and
More on Algebra, Lemmas \ref{more-algebra-lemma-six-term-Rlim} and
\ref{more-algebra-lemma-Mittag-Leffler}.
\end{remark}










\section{Derived limits}
\label{section-derived-limits}

\noindent
Let $(X, \mathcal{O}_X)$ be a ringed space. Since the triangulated category
$D(\mathcal{O}_X)$ has products
(Injectives, Lemma \ref{injectives-lemma-derived-products})
it follows that $D(\mathcal{O}_X)$ has derived limits, see
Derived Categories, Definition \ref{derived-definition-derived-limit}.
If $(K_n)$ is an inverse system in $D(\mathcal{O}_X)$ then we
denote $R\lim K_n$ the derived limit.

\begin{lemma}
\label{lemma-RGamma-commutes-with-Rlim}
Let $(X, \mathcal{O}_X)$ be a ringed space. For $U \subset X$ open the
functor $R\Gamma(U, -)$ commutes with $R\lim$. Moreover, there are
short exact sequences
$$
0 \to
R^1\lim H^{m - 1}(U, K_n) \to H^m(U, R\lim K_n) \to
\lim H^m(U, K_n) \to 0
$$
for any inverse system $(K_n)$ in $D(\mathcal{O}_X)$ and any $m \in \mathbf{Z}$.
\end{lemma}

\begin{proof}
The first statement follows from
Injectives, Lemma \ref{injectives-lemma-RF-commutes-with-Rlim}.
Then we may apply
More on Algebra, Remark \ref{more-algebra-remark-compare-derived-limit}
to $R\lim R\Gamma(U, K_n) = R\Gamma(U, R\lim K_n)$ to get the short
exact sequences.
\end{proof}

\begin{lemma}
\label{lemma-Rf-commutes-with-Rlim}
Let $f : (X, \mathcal{O}_X) \to (Y, \mathcal{O}_Y)$ be a morphism of ringed
spaces. Then $Rf_*$ commutes with $R\lim$, i.e., $Rf_*$ commutes with
derived limits.
\end{lemma}

\begin{proof}
Let $(K_n)$ be an inverse system in $D(\mathcal{O}_X)$. Consider the defining
distinguished triangle
$$
R\lim K_n \to \prod K_n \to \prod K_n
$$
in $D(\mathcal{O}_X)$. Applying the exact functor $Rf_*$ we obtain
the distinguished triangle
$$
Rf_*(R\lim K_n) \to Rf_*\left(\prod K_n\right) \to Rf_*\left(\prod K_n\right)
$$
in $D(\mathcal{O}_Y)$. Thus we see that it suffices to prove that
$Rf_*$ commutes with products in the derived category (which are not just
given by products of complexes, see
Injectives, Lemma \ref{injectives-lemma-derived-products}).
However, since $Rf_*$ is a right adjoint by Lemma \ref{lemma-adjoint}
this follows formally (see
Categories, Lemma \ref{categories-lemma-adjoint-exact}).
Caution: Note that we cannot apply
Categories, Lemma \ref{categories-lemma-adjoint-exact}
directly as $R\lim K_n$ is not a limit in $D(\mathcal{O}_X)$.
\end{proof}

\begin{remark}
\label{remark-discuss-derived-limit}
Let $(X, \mathcal{O}_X)$ be a ringed space. Let $(K_n)$ be an inverse
system in $D(\mathcal{O}_X)$. Set $K = R\lim K_n$. For each $n$ and $m$
let $\mathcal{H}^m_n = H^m(K_n)$ be the $m$th cohomology sheaf of
$K_n$ and similarly set $\mathcal{H}^m = H^m(K)$. Let us denote
$\underline{\mathcal{H}}^m_n$ the presheaf
$$
U \longmapsto \underline{\mathcal{H}}^m_n(U) = H^m(U, K_n)
$$
Similarly we set $\underline{\mathcal{H}}^m(U) = H^m(U, K)$.
By Lemma \ref{lemma-sheafification-cohomology} we see that
$\mathcal{H}^m_n$ is the sheafification of
$\underline{\mathcal{H}}^m_n$ and $\mathcal{H}^m$ is the
sheafification of $\underline{\mathcal{H}}^m$.
Here is a diagram
$$
\xymatrix{
K \ar@{=}[d] &
\underline{\mathcal{H}}^m \ar[d] \ar[r] &
\mathcal{H}^m \ar[d] \\
R\lim K_n &
\lim \underline{\mathcal{H}}^m_n \ar[r] &
\lim \mathcal{H}^m_n
}
$$
In general it may not be the case that
$\lim \mathcal{H}^m_n$ is the sheafification of
$\lim \underline{\mathcal{H}}^m_n$.
If $U \subset X$ is an open, then we have short exact
sequences
\begin{equation}
\label{equation-ses-Rlim-over-U}
0 \to
R^1\lim \underline{\mathcal{H}}^{m - 1}_n(U) \to
\underline{\mathcal{H}}^m(U) \to
\lim \underline{\mathcal{H}}^m_n(U) \to 0
\end{equation}
by Lemma \ref{lemma-RGamma-commutes-with-Rlim}.
\end{remark}

\noindent
The following lemma applies to an inverse system of quasi-coherent
modules with surjective transition maps on a scheme.

\begin{lemma}
\label{lemma-inverse-limit-is-derived-limit}
Let $(X, \mathcal{O}_X)$ be a ringed space. Let $(\mathcal{F}_n)$ be an
inverse system of $\mathcal{O}_X$-modules. Let $\mathcal{B}$ be a set
of opens of $X$. Assume
\begin{enumerate}
\item every open of $X$ has a covering whose members are elements of
$\mathcal{B}$,
\item $H^p(U, \mathcal{F}_n) = 0$ for $p > 0$ and $U \in \mathcal{B}$,
\item the inverse system $\mathcal{F}_n(U)$ has vanishing $R^1\lim$
for $U \in \mathcal{B}$.
\end{enumerate}
Then $R\lim \mathcal{F}_n = \lim \mathcal{F}_n$ and we have
$H^p(U, \lim \mathcal{F}_n) = 0$ for $p > 0$ and $U \in \mathcal{B}$.
\end{lemma}

\begin{proof}
Set $K_n = \mathcal{F}_n$ and $K = R\lim \mathcal{F}_n$. Using the notation
of Remark \ref{remark-discuss-derived-limit} and assumption (2) we see that for
$U \in \mathcal{B}$ we have $\underline{\mathcal{H}}_n^m(U) = 0$
when $m \not = 0$ and $\underline{\mathcal{H}}_n^0(U) = \mathcal{F}_n(U)$.
From Equation (\ref{equation-ses-Rlim-over-U}) and assumption (3)
we see that $\underline{\mathcal{H}}^m(U) = 0$
when $m \not = 0$ and equal to $\lim \mathcal{F}_n(U)$
when $m = 0$. Sheafifying using (1) we find that
$\mathcal{H}^m = 0$ when $m \not = 0$ and equal to
$\lim \mathcal{F}_n$ when $m = 0$. Hence $K = \lim \mathcal{F}_n$.
Since $H^m(U, K) = \underline{\mathcal{H}}^m(U) = 0$ for $m > 0$
(see above) we see that the second assertion holds.
\end{proof}

\begin{lemma}
\label{lemma-cohomology-derived-limit-injective}
Let $(X, \mathcal{O}_X)$ be a ringed space. Let $(K_n)$ be an
inverse system in $D(\mathcal{O}_X)$. Let $x \in X$ and $m \in \mathbf{Z}$.
Assume there exist an integer $n(x)$ and a fundamental system $\mathfrak{U}_x$
of open neighbourhoods of $x$ such that for $U \in \mathfrak{U}_x$
\begin{enumerate}
\item $R^1\lim H^{m - 1}(U, K_n) = 0$, and
\item $H^m(U, K_n) \to H^m(U, K_{n(x)})$ is injective
for $n \geq n(x)$.
\end{enumerate}
Then the map on stalks $H^m(R\lim K_n)_x \to H^m(K_{n(x)})_x$ is injective.
\end{lemma}

\begin{proof}
Let $\gamma$ be an element of $H^m(R\lim K_n)_x$ which maps to zero
in $H^m(K_{n(x)})_x$. Since $H^m(R\lim K_n)$ is the sheafification
of $U \mapsto H^m(U, R\lim K_n)$
(by Lemma \ref{lemma-sheafification-cohomology})
we can choose $U \in \mathfrak{U}_x$
and an element $\tilde \gamma \in H^m(U, R\lim K_n)$ mapping to $\gamma$.
Then $\tilde\gamma$ maps to $\tilde\gamma_{n(x)} \in H^m(U, K_{n(x)})$.
Using that $H^m(K_{n(x)})$ is the sheafification of
$U \mapsto H^m(U, K_{n(x)})$
(by Lemma \ref{lemma-sheafification-cohomology} again)
we see that after shrinking $U$ we may assume that $\tilde\gamma_{n(x)} = 0$.
For this $U$ we consider the short exact sequence
$$
0 \to
R^1\lim H^{m - 1}(U, K_n) \to H^m(U, R\lim K_n) \to
\lim H^m(U, K_n) \to 0
$$
of Lemma \ref{lemma-RGamma-commutes-with-Rlim}.
By assumption (1) the group on the left is zero and by
assumption (2) the group on the right maps injectively
into $H^m(U, K_{n(x)})$. We conclude $\tilde\gamma = 0$
and hence $\gamma = 0$ as desired.
\end{proof}

\begin{lemma}
\label{lemma-is-limit-per-point}
Let $(X, \mathcal{O}_X)$ be a ringed space. Let $E \in D(\mathcal{O}_X)$.
Assume that for every $x \in X$ there exist
a function $p(x, -) : \mathbf{Z} \to \mathbf{Z}$ and
a fundamental system $\mathfrak{U}_x$ of open neighbourhoods of $x$
such that
$$
H^p(U, H^{m - p}(E)) = 0 \text{ for }
U \in \mathfrak{U}_x \text{ and } p > p(x, m)
$$
Then the map $E \to R\lim \tau_{\geq -n} E$ of
Derived Categories, Remark
\ref{derived-remark-map-into-derived-limit-truncations}
is an isomorphism in $D(\mathcal{O}_X)$.
\end{lemma}

\begin{proof}
Set $K_n = \tau_{\geq -n}E$ and $K = R\lim K_n$.
The canonical map $E \to K$
comes from the canonical maps $E \to K_n = \tau_{\geq -n}E$.
We have to show that $E \to K$ induces an isomorphism
$H^m(E) \to H^m(K)$ of cohomology sheaves. In the rest of the
proof we fix $m$. If $n \geq -m$, then
the map $E \to \tau_{\geq -n}E = K_n$ induces an isomorphism
$H^m(E) \to H^m(K_n)$.
To finish the proof it suffices to show that for every $x \in X$
there exists an integer $n(x) \geq -m$ such that the map
$H^m(K)_x \to H^m(K_{n(x)})_x$ is injective. Namely, then
the composition
$$
H^m(E)_x \to H^m(K)_x \to H^m(K_{n(x)})_x
$$
is a bijection and the second arrow is injective, hence the
first arrow is bijective. Set
$$
n(x) = 1 + \max\{-m, p(x, m - 1) - m, -1 + p(x, m) - m, -2 + p(x, m + 1) - m\}.
$$
so that in any case $n(x) \geq -m$. Claim: the maps
$$
H^{m - 1}(U, K_{n + 1}) \to H^{m - 1}(U, K_n)
\quad\text{and}\quad
H^m(U, K_{n + 1}) \to H^m(U, K_n)
$$
are isomorphisms for $n \geq n(x)$ and $U \in \mathfrak{U}_x$.
The claim implies conditions
(1) and (2) of Lemma \ref{lemma-cohomology-derived-limit-injective}
are satisfied and hence implies the desired injectivity.
Recall (Derived Categories, Remark
\ref{derived-remark-truncation-distinguished-triangle})
that we have distinguished triangles
$$
H^{-n - 1}(E)[n + 1] \to
K_{n + 1} \to K_n \to H^{-n - 1}(E)[n + 2]
$$
Looking at the associated long exact cohomology sequence the claim follows if
$$
H^{m + n}(U, H^{-n - 1}(E)),\quad
H^{m + n + 1}(U, H^{-n - 1}(E)),\quad
H^{m + n + 2}(U, H^{-n - 1}(E))
$$
are zero for $n \geq n(x)$ and $U \in \mathfrak{U}_x$.
This follows from our choice of $n(x)$
and the assumption in the lemma.
\end{proof}

\begin{lemma}
\label{lemma-is-limit-spaltenstein}
\begin{reference}
\cite[Proposition 3.13]{Spaltenstein}
\end{reference}
Let $(X, \mathcal{O}_X)$ be a ringed space. Let $E \in D(\mathcal{O}_X)$.
Assume that for every $x \in X$ there exist an integer $d_x \geq 0$ and
a fundamental system $\mathfrak{U}_x$ of open neighbourhoods of $x$
such that
$$
H^p(U, H^q(E)) = 0 \text{ for }
U \in \mathfrak{U}_x,\ p > d_x, \text{ and }q < 0
$$
Then the map $E \to R\lim \tau_{\geq -n} E$ of
Derived Categories, Remark
\ref{derived-remark-map-into-derived-limit-truncations}
is an isomorphism in $D(\mathcal{O}_X)$.
\end{lemma}

\begin{proof}
This follows from Lemma \ref{lemma-is-limit-per-point}
with $p(x, m) = d_x + \max(0, m)$.
\end{proof}

\begin{lemma}
\label{lemma-is-limit}
Let $(X, \mathcal{O}_X)$ be a ringed space. Let $E \in D(\mathcal{O}_X)$.
Assume there exist a function $p(-) : \mathbf{Z} \to \mathbf{Z}$
and a set $\mathcal{B}$ of opens of $X$ such that
\begin{enumerate}
\item every open in $X$ has a covering whose members are
elements of $\mathcal{B}$, and
\item $H^p(U, H^{m - p}(E)) = 0$ for $p > p(m)$ and $U \in \mathcal{B}$.
\end{enumerate}
Then the map $E \to R\lim \tau_{\geq -n} E$ of
Derived Categories, Remark
\ref{derived-remark-map-into-derived-limit-truncations}
is an isomorphism in $D(\mathcal{O}_X)$.
\end{lemma}

\begin{proof}
Apply Lemma \ref{lemma-is-limit-per-point}
with $p(x, m) = p(m)$ and
$\mathfrak{U}_x = \{U \in \mathcal{B} \mid x \in U\}$.
\end{proof}

\begin{lemma}
\label{lemma-is-limit-dimension}
Let $(X, \mathcal{O}_X)$ be a ringed space. Let $E \in D(\mathcal{O}_X)$.
Assume there exist an integer $d \geq 0$ and a basis $\mathcal{B}$ for the
topology of $X$ such that
$$
H^p(U, H^q(E)) = 0 \text{ for }
U \in \mathcal{B},\ p > d, \text{ and }q < 0
$$
Then the map $E \to R\lim \tau_{\geq -n} E$ of
Derived Categories, Remark
\ref{derived-remark-map-into-derived-limit-truncations}
is an isomorphism in $D(\mathcal{O}_X)$.
\end{lemma}

\begin{proof}
Apply Lemma \ref{lemma-is-limit-spaltenstein} with $d_x = d$
and $\mathfrak{U}_x = \{U \in \mathcal{B} \mid x \in U\}$.
\end{proof}

\noindent
The lemmas above can be used to compute cohomology
in certain situations.

\begin{lemma}
\label{lemma-cohomology-over-U-trivial}
Let $(X, \mathcal{O}_X)$ be a ringed space. Let $K$
be an object of $D(\mathcal{O}_X)$.
Let $\mathcal{B}$ be a set of opens of $X$. Assume
\begin{enumerate}
\item every open of $X$ has a covering whose members are
elements of $\mathcal{B}$,
\item $H^p(U, H^q(K)) = 0$ for all $p > 0$, $q \in \mathbf{Z}$, and
$U \in \mathcal{B}$.
\end{enumerate}
Then $H^q(U, K) = H^0(U, H^q(K))$ for $q \in \mathbf{Z}$
and $U \in \mathcal{B}$.
\end{lemma}

\begin{proof}
Observe that $K = R\lim \tau_{\geq -n} K$ by
Lemma \ref{lemma-is-limit-dimension} with $d = 0$.
Let $U \in \mathcal{B}$. By Equation (\ref{equation-ses-Rlim-over-U})
we get a short exact sequence
$$
0 \to R^1\lim H^{q - 1}(U, \tau_{\geq -n}K) \to
H^q(U, K) \to \lim H^q(U, \tau_{\geq -n}K) \to 0
$$
Condition (2) implies
$H^q(U, \tau_{\geq -n} K) = H^0(U, H^q(\tau_{\geq -n} K))$
for all $q$ by using the spectral sequence of
Example \ref{example-spectral-sequence}.
The spectral sequence converges because $\tau_{\geq -n}K$ is bounded
below. If $n > -q$ then we have $H^q(\tau_{\geq -n}K) = H^q(K)$.
Thus the systems on the left and the right of the displayed
short exact sequence are eventually constant with values
$H^0(U, H^{q - 1}(K))$ and $H^0(U, H^q(K))$. The lemma follows.
\end{proof}

\noindent
Here is another case where we can describe the derived limit.

\begin{lemma}
\label{lemma-derived-limit-suitable-system}
Let $(X, \mathcal{O}_X)$ be a ringed space. Let $(K_n)$
be an inverse system of objects of $D(\mathcal{O}_X)$.
Let $\mathcal{B}$ be a set of opens of $X$. Assume
\begin{enumerate}
\item every open of $X$ has a covering whose members are
elements of $\mathcal{B}$,
\item for all $U \in \mathcal{B}$ and all $q \in \mathbf{Z}$ we have
\begin{enumerate}
\item $H^p(U, H^q(K_n)) = 0$ for $p > 0$,
\item the inverse system $H^0(U, H^q(K_n))$ has vanishing $R^1\lim$.
\end{enumerate}
\end{enumerate}
Then $H^q(R\lim K_n) = \lim H^q(K_n)$ for $q \in \mathbf{Z}$.
\end{lemma}

\begin{proof}
Set $K = R\lim K_n$. Let $U \in \mathcal{B}$.
By Lemma \ref{lemma-cohomology-over-U-trivial} and (2)(a)
we have $H^q(U, K_n) = H^0(U, H^q(K_n))$. By
Lemma \ref{lemma-RGamma-commutes-with-Rlim} and (2)(b) we have
$H^q(U, K) = \lim H^0(U, H^q(K_n))$. Thus $H^q(U, K)$ is the inverse limit
the sections of the sheaves $H^q(K_n)$ over $U$. Since
$\lim H^q(K_n)$ is a sheaf we find using assumption (1) that $H^q(K)$,
which is the sheafification of the presheaf $U \mapsto H^q(U, K)$,
is equal to $\lim H^q(K_n)$. This proves the lemma.
\end{proof}








\section{Producing K-injective resolutions}
\label{section-K-injective}


\noindent
Let $(X, \mathcal{O}_X)$ be a ringed space.
Let $\mathcal{F}^\bullet$ be a complex of $\mathcal{O}_X$-modules.
The category $\textit{Mod}(\mathcal{O}_X)$ has enough injectives, hence
we can use
Derived Categories, Lemma \ref{derived-lemma-special-inverse-system}
produce a diagram
$$
\xymatrix{
\ldots \ar[r] &
\tau_{\geq -2}\mathcal{F}^\bullet \ar[r] \ar[d] &
\tau_{\geq -1}\mathcal{F}^\bullet \ar[d] \\
\ldots \ar[r] & \mathcal{I}_2^\bullet \ar[r] & \mathcal{I}_1^\bullet
}
$$
in the category of complexes of $\mathcal{O}_X$-modules such that
\begin{enumerate}
\item the vertical arrows are quasi-isomorphisms,
\item $\mathcal{I}_n^\bullet$ is a bounded below complex of injectives,
\item the arrows $\mathcal{I}_{n + 1}^\bullet \to \mathcal{I}_n^\bullet$
are termwise split surjections.
\end{enumerate}
The category of $\mathcal{O}_X$-modules has limits (they are computed
on the level of presheaves), hence we can form the termwise limit
$\mathcal{I}^\bullet = \lim_n \mathcal{I}_n^\bullet$. By
Derived Categories, Lemmas
\ref{derived-lemma-bounded-below-injectives-K-injective} and
\ref{derived-lemma-limit-K-injectives}
this is a K-injective complex. In general the canonical map
\begin{equation}
\label{equation-into-candidate-K-injective}
\mathcal{F}^\bullet \to \mathcal{I}^\bullet
\end{equation}
may not be a quasi-isomorphism. In the following lemma we describe some
conditions under which it is.

\begin{lemma}
\label{lemma-K-injective}
In the situation described above.
Denote $\mathcal{H}^m = H^m(\mathcal{F}^\bullet)$ the $m$th cohomology sheaf.
Let $\mathcal{B}$ be a set of open subsets of $X$.
Let $d \in \mathbf{N}$.
Assume
\begin{enumerate}
\item every open in $X$ has a covering whose members are
elements of $\mathcal{B}$,
\item for every $U \in \mathcal{B}$ we have $H^p(U, \mathcal{H}^q) = 0$
for $p > d$ and $q < 0$\footnote{It suffices if
$\forall m$, $\exists p(m)$, $H^p(U. \mathcal{H}^{m - p}) = 0$ for
$p > p(m)$, see Lemma \ref{lemma-is-limit}.}.
\end{enumerate}
Then (\ref{equation-into-candidate-K-injective}) is a quasi-isomorphism.
\end{lemma}

\begin{proof}
By Derived Categories, Lemma \ref{derived-lemma-difficulty-K-injectives}
it suffices to show that the map
$\mathcal{F}^\bullet \to R\lim \tau_{\geq -n} \mathcal{F}^\bullet$
is an isomorphism. This is Lemma \ref{lemma-is-limit-dimension}.
\end{proof}

\noindent
Here is a technical lemma about the cohomology sheaves of the inverse
limit of a system of complexes of sheaves. In some sense this lemma
is the wrong thing to try to prove as one should take derived
limits and not actual inverse limits.

\begin{lemma}
\label{lemma-inverse-limit-complexes}
Let $(X, \mathcal{O}_X)$ be a ringed space. Let $(\mathcal{F}_n^\bullet)$
be an inverse system of complexes of $\mathcal{O}_X$-modules.
Let $m \in \mathbf{Z}$. Assume there exist a set $\mathcal{B}$
of open subsets of $X$ and an integer $n_0$ such that
\begin{enumerate}
\item every open in $X$ has a covering whose members are
elements of $\mathcal{B}$,
\item for every $U \in \mathcal{B}$
\begin{enumerate}
\item the systems of abelian groups
$\mathcal{F}_n^{m - 2}(U)$ and $\mathcal{F}_n^{m - 1}(U)$
have vanishing $R^1\lim$ (for example these have the Mittag-Leffler
condition),
\item the system of abelian groups $H^{m - 1}(\mathcal{F}_n^\bullet(U))$
has vanishing $R^1\lim$ (for example it has the Mittag-Leffler condition), and
\item we have
$H^m(\mathcal{F}_n^\bullet(U)) = H^m(\mathcal{F}_{n_0}^\bullet(U))$
for all $n \geq n_0$.
\end{enumerate}
\end{enumerate}
Then the maps
$H^m(\mathcal{F}^\bullet) \to \lim H^m(\mathcal{F}_n^\bullet) \to
H^m(\mathcal{F}_{n_0}^\bullet)$
are isomorphisms of sheaves where
$\mathcal{F}^\bullet = \lim \mathcal{F}_n^\bullet$ is the termwise
inverse limit.
\end{lemma}

\begin{proof}
Let $U \in \mathcal{B}$. Note that $H^m(\mathcal{F}^\bullet(U))$ is the
cohomology of
$$
\lim_n \mathcal{F}_n^{m - 2}(U) \to
\lim_n \mathcal{F}_n^{m - 1}(U) \to
\lim_n \mathcal{F}_n^m(U) \to
\lim_n \mathcal{F}_n^{m + 1}(U)
$$
in the third spot from the left. By assumptions (2)(a) and (2)(b)
we may apply
More on Algebra, Lemma \ref{more-algebra-lemma-apply-Mittag-Leffler-again}
to conclude that
$$
H^m(\mathcal{F}^\bullet(U)) = \lim H^m(\mathcal{F}_n^\bullet(U))
$$
By assumption (2)(c) we conclude
$$
H^m(\mathcal{F}^\bullet(U)) = H^m(\mathcal{F}_n^\bullet(U))
$$
for all $n \geq n_0$. By assumption (1) we conclude that the sheafification of
$U \mapsto H^m(\mathcal{F}^\bullet(U))$ is equal to the sheafification
of $U \mapsto H^m(\mathcal{F}_n^\bullet(U))$ for all $n \geq n_0$.
Thus the inverse system of sheaves $H^m(\mathcal{F}_n^\bullet)$ is
constant for $n \geq n_0$ with value $H^m(\mathcal{F}^\bullet)$ which
proves the lemma.
\end{proof}












\section{Inverse systems and cohomology, III}
\label{section-inverse-systems-tri}

\noindent
This section continues the discussion in
Section \ref{section-inverse-systems-bis}
using derived limits.

\begin{lemma}
\label{lemma-formal-functions-principal}
Let $(X, \mathcal{O}_X)$ be a ringed space. Let $A \to \Gamma(X, \mathcal{O}_X)$
be a ring map and let $f \in A$. Let $E$ be an object of $D(\mathcal{O}_X)$.
Denote
$$
E_n = E \otimes_{\mathcal{O}_X} (\mathcal{O}_X \xrightarrow{f^n} \mathcal{O}_X)
$$
and set $E^\wedge = R\lim E_n$. For $p \in \mathbf{Z}$
is a canonical commutative diagram
$$
\xymatrix{
& 0 & 0 \\
0 \ar[r] &
\widehat{H^p(X, E)} \ar[r] \ar[u] &
\lim H^p(X, E_n) \ar[r] \ar[u] &
T_f(H^{p + 1}(X, E)) \ar[r] &
0 \\
0 \ar[r] &
H^0(H^p(X, E)^\wedge) \ar[r] \ar[u] &
H^p(X, E^\wedge) \ar[r] \ar[u] &
T_f(H^{p + 1}(X, E)) \ar[r] \ar@{=}[u] &
0 \\
&
R^1\lim H^p(X, E)[f^n] \ar[u] \ar[r]^\cong &
R^1\lim H^{p - 1}(X, E_n) \ar[u] \\
& 0 \ar[u] & 0 \ar[u]
}
$$
with exact rows and columns
where $\widehat{H^p(X, E)} = \lim H^p(X, E)/f^n H^p(X, E)$ is the usual
$f$-adic completion, $H^p(X, E)^\wedge$ is the derived $f$-adic completion,
and $T_f(H^{p + 1}(X, E))$ is the $f$-adic Tate module, see
More on Algebra, Example
\ref{more-algebra-example-spectral-sequence-principal}.
Finally, we have $H^p(X, E^\wedge) = H^p(R\Gamma(X, E)^\wedge)$.
\end{lemma}

\begin{proof}
Observe that $R\Gamma(X, E^\wedge) = R\lim R\Gamma(X, E_n)$ by
Lemma \ref{lemma-Rf-commutes-with-Rlim}.
On the other hand, we have
$$
R\Gamma(X, E_n) =
R\Gamma(X, E) \otimes_A^\mathbf{L} (A \xrightarrow{f^n} A)
$$
(details omitted). We find that $R\Gamma(X, E^\wedge)$ is the derived
$f$-adic completion $R\Gamma(X, E)^\wedge$. Whence the diagram by
More on Algebra, Lemma
\ref{more-algebra-lemma-compare-completion-derived-completion}.
\end{proof}

\begin{lemma}
\label{lemma-stabilizes}
Let $\mathcal{A}$ be an abelian category. Let $f : M \to M$ be a morphism
of $\mathcal{A}$. If $M[f^n] = \Ker(f^n : M \to M)$ stabilizes, then the
inverse systems
$$
(M \xrightarrow{f^n} M)
\quad\text{and}\quad
\Coker(f^n : M \to M)
$$
are pro-isomorphic in $D(\mathcal{A})$.
\end{lemma}

\begin{proof}
There is clearly a map from the first inverse system to the second. Suppose that
$M[f^c] = M[f^{c + 1}] = M[f^{c + 2}] = \ldots$.
Then we can define an arrow of inverse systems in $D(\mathcal{A})$
in the other direction by the diagrams
$$
\xymatrix{
M/M[f^c] \ar[r]_-{f^{n + c}} \ar[d]_{f^c} &
M \ar[d]^1 \\
M \ar[r]^{f^n} & M
}
$$
Since the top horizontal arrow is injective the complex
in the top row is quasi-isomorphic to $\Coker(f^{n + c} : M \to M)$.
Some details omitted.
\end{proof}

\begin{example}
\label{example-formal-functions-principal}
Let $(X, \mathcal{O}_X)$ be a ringed space. Let $A \to \Gamma(X, \mathcal{O}_X)$
be a ring map and let $f \in A$. Let $\mathcal{F}$ be an $\mathcal{O}_X$-module.
Assume there is a $c$ such that $\mathcal{F}[f^c] = \mathcal{F}[f^n]$
for all $n \geq c$. We are going to apply
Lemma \ref{lemma-formal-functions-principal} with
$E = \mathcal{F}$. By Lemma \ref{lemma-stabilizes} we see that
the inverse system $(E_n)$ is pro-isomorphic to the inverse system
$(\mathcal{F}/f^n\mathcal{F})$. We conclude that for $p \in \mathbf{Z}$
we obtain a commutative diagram
$$
\xymatrix{
& 0 & 0 \\
0 \ar[r] &
\widehat{H^p(X, \mathcal{F})} \ar[r] \ar[u] &
\lim H^p(X, \mathcal{F}/f^n\mathcal{F}) \ar[r] \ar[u] &
T_f(H^{p + 1}(X, \mathcal{F})) \ar[r] &
0 \\
0 \ar[r] &
H^0(H^p(X, \mathcal{F})^\wedge) \ar[r] \ar[u] &
H^p(R\Gamma(X, \mathcal{F})^\wedge) \ar[r] \ar[u] &
T_f(H^{p + 1}(X, \mathcal{F})) \ar[r] \ar@{=}[u] &
0 \\
&
R^1\lim H^p(X, \mathcal{F})[f^n] \ar[u] \ar[r]^\cong &
R^1\lim H^{p - 1}(X, \mathcal{F}/f^n\mathcal{F}) \ar[u] \\
& 0 \ar[u] & 0 \ar[u]
}
$$
with exact rows and columns where
$\widehat{H^p(X, \mathcal{F})} =
\lim H^p(X, \mathcal{F})/f^n H^p(X, \mathcal{F})$
is the usual $f$-adic completion
and $M^\wedge$ denotes derived $f$-adic completion
for $M$ in $D(A)$.
\end{example}
















\section{{\v C}ech cohomology of unbounded complexes}
\label{section-cech-cohomology-of-unbounded-complexes}

\noindent
The construction of Section \ref{section-cech-cohomology-of-complexes}
isn't the ``correct'' one for unbounded complexes. The problem is that
in the Stacks project we use direct sums in the totalization of a
double complex and we would have to replace this by a product. Instead
of doing so in this section we assume the covering is finite and
we use the alternating {\v C}ech complex.

\medskip\noindent
Let $(X, \mathcal{O}_X)$ be a ringed space.
Let ${\mathcal F}^\bullet$ be a complex of presheaves of
$\mathcal{O}_X$-modules. Let ${\mathcal U} : X = \bigcup_{i \in I} U_i$
be a {\bf finite} open covering of $X$. Since the alternating
{\v C}ech complex
$\check{\mathcal{C}}_{alt}^\bullet(\mathcal{U}, \mathcal{F})$
(Section \ref{section-alternating-cech})
is functorial in the presheaf $\mathcal{F}$ we obtain a double complex
$\check{\mathcal{C}}^\bullet_{alt}(\mathcal{U}, \mathcal{F}^\bullet)$.
In this section we work with the associated total complex.
The construction of
$\text{Tot}(\check{\mathcal{C}}^\bullet_{alt}({\mathcal U},
{\mathcal F}^\bullet))$
is functorial in ${\mathcal F}^\bullet$. As well there is a functorial
transformation
\begin{equation}
\label{equation-global-sections-to-alternating-cech}
\Gamma(X, {\mathcal F}^\bullet)
\longrightarrow
\text{Tot}(\check{\mathcal{C}}^\bullet_{alt}({\mathcal U},
{\mathcal F}^\bullet))
\end{equation}
of complexes defined by the following rule: The section
$s\in \Gamma(X, {\mathcal F}^n)$
is mapped to the element $\alpha = \{\alpha_{i_0\ldots i_p}\}$
with $\alpha_{i_0} = s|_{U_{i_0}}$ and $\alpha_{i_0\ldots i_p} = 0$
for $p > 0$.

\begin{lemma}
\label{lemma-alternating-cech-complex-complex}
Let $(X, \mathcal{O}_X)$ be a ringed space.
Let $\mathcal{U} : X = \bigcup_{i \in I} U_i$ be
a finite open covering. For a complex $\mathcal{F}^\bullet$
of $\mathcal{O}_X$-modules there is a canonical map
$$
\text{Tot}(\check{\mathcal{C}}^\bullet_{alt}(\mathcal{U}, \mathcal{F}^\bullet))
\longrightarrow
R\Gamma(X, \mathcal{F}^\bullet)
$$
functorial in $\mathcal{F}^\bullet$ and compatible with
(\ref{equation-global-sections-to-alternating-cech}).
\end{lemma}

\begin{proof}
Let ${\mathcal I}^\bullet$ be a K-injective complex whose terms
are injective $\mathcal{O}_X$-modules.
The map (\ref{equation-global-sections-to-alternating-cech}) for
$\mathcal{I}^\bullet$ is a map
$\Gamma(X, {\mathcal I}^\bullet) \to
\text{Tot}(\check{\mathcal{C}}^\bullet_{alt}({\mathcal U},
{\mathcal I}^\bullet))$.
This is a quasi-isomorphism of complexes of abelian groups
as follows from
Homology, Lemma \ref{homology-lemma-double-complex-gives-resolution}
applied to the double complex
$\check{\mathcal{C}}^\bullet_{alt}({\mathcal U},
{\mathcal I}^\bullet)$ using
Lemmas \ref{lemma-injective-trivial-cech} and \ref{lemma-alternating-usual}.
Suppose ${\mathcal F}^\bullet \to {\mathcal I}^\bullet$ is a quasi-isomorphism
of ${\mathcal F}^\bullet$ into a K-injective complex whose terms
are injectives (Injectives, Theorem
\ref{injectives-theorem-K-injective-embedding-grothendieck}).
Since $R\Gamma(X, {\mathcal F}^\bullet)$ is represented by the complex
$\Gamma(X, {\mathcal I}^\bullet)$ we obtain the map of the lemma
using
$$
\text{Tot}(\check{\mathcal{C}}^\bullet_{alt}({\mathcal U},
{\mathcal F}^\bullet))
\longrightarrow
\text{Tot}(\check{\mathcal{C}}^\bullet_{alt}({\mathcal U},
{\mathcal I}^\bullet)).
$$
We omit the verification of functoriality and compatibilities.
\end{proof}

\begin{lemma}
\label{lemma-alternating-cech-complex-complex-ss}
Let $(X, \mathcal{O}_X)$ be a ringed space. Let
$\mathcal{U} : X = \bigcup_{i \in I} U_i$ be a finite open covering. Let
$\mathcal{F}^\bullet$ be a complex of $\mathcal{O}_X$-modules.
Let $\mathcal{B}$ be a set of open subsets of $X$. Assume
\begin{enumerate}
\item every open in $X$ has a covering whose members are
elements of $\mathcal{B}$,
\item we have $U_{i_0\ldots i_p} \in \mathcal{B}$ for all
$i_0, \ldots, i_p \in I$,
\item for every $U \in \mathcal{B}$ and $p > 0$ we have
\begin{enumerate}
\item $H^p(U, \mathcal{F}^q) = 0$,
\item $H^p(U, \Coker(\mathcal{F}^{q - 1} \to \mathcal{F}^q)) = 0$, and
\item $H^p(U, H^q(\mathcal{F})) = 0$.
\end{enumerate}
\end{enumerate}
Then the map
$$
\text{Tot}(\check{\mathcal{C}}^\bullet_{alt}(\mathcal{U}, \mathcal{F}^\bullet))
\longrightarrow
R\Gamma(X, \mathcal{F}^\bullet)
$$
of Lemma \ref{lemma-alternating-cech-complex-complex}
is an isomorphism in $D(\textit{Ab})$.
\end{lemma}

\begin{proof}
First assume $\mathcal{F}^\bullet$ is bounded below. In this case the map
$$
\text{Tot}(\check{\mathcal{C}}^\bullet_{alt}(\mathcal{U}, \mathcal{F}^\bullet))
\longrightarrow
\text{Tot}(\check{\mathcal{C}}^\bullet(\mathcal{U}, \mathcal{F}^\bullet))
$$
is a quasi-isomorphism by Lemma \ref{lemma-alternating-usual}.
Namely, the map of double complexes
$\check{\mathcal{C}}^\bullet_{alt}(\mathcal{U}, \mathcal{F}^\bullet) \to
\check{\mathcal{C}}^\bullet(\mathcal{U}, \mathcal{F}^\bullet)$
induces an isomorphism between the first pages of the second spectral sequences
associated to these complexes
(by Homology, Lemma \ref{homology-lemma-ss-double-complex})
and these spectral sequences converge
(Homology, Lemma \ref{homology-lemma-first-quadrant-ss}).
Thus the conclusion in this case by
Lemma \ref{lemma-cech-complex-complex-computes} and assumption (3)(a).

\medskip\noindent
In general, by assumption (3)(c) we may choose a resolution
$\mathcal{F}^\bullet \to \mathcal{I}^\bullet = \lim \mathcal{I}_n^\bullet$
as in Lemma \ref{lemma-K-injective}.
Then the map of the lemma becomes
$$
\lim_n
\text{Tot}(\check{\mathcal{C}}^\bullet_{alt}(\mathcal{U},
\tau_{\geq -n}\mathcal{F}^\bullet))
\longrightarrow
\Gamma(X, \mathcal{I}^\bullet) =
\lim_n \Gamma(X, \mathcal{I}_n^\bullet)
$$
Here the arrow is in the derived category, but the equality on the
right holds on the level of complexes.
Note that (3)(b) shows that $\tau_{\geq -n}\mathcal{F}^\bullet$
is a bounded below complex satisfying the hypothesis of the lemma.
Thus the case of bounded below complexes shows each of the maps
$$
\text{Tot}(\check{\mathcal{C}}^\bullet_{alt}(\mathcal{U},
\tau_{\geq -n}\mathcal{F}^\bullet))
\longrightarrow
\Gamma(X, \mathcal{I}_n^\bullet)
$$
is a quasi-isomorphism. The cohomologies of the complexes on the left
hand side in given degree are eventually
constant (as the alternating {\v C}ech complex is finite).
Hence the same is true on the right hand side.
Thus the cohomology of the limit on the right hand side is
this constant value by
Homology, Lemma \ref{homology-lemma-apply-Mittag-Leffler-again}
(or the stronger More on Algebra, Lemma
\ref{more-algebra-lemma-apply-Mittag-Leffler-again})
and we win.
\end{proof}






\section{Hom complexes}
\label{section-hom-complexes}

\noindent
Let $(X, \mathcal{O}_X)$ be a ringed space. Let
$\mathcal{L}^\bullet$ and $\mathcal{M}^\bullet$ be two complexes
of $\mathcal{O}_X$-modules. We construct a complex
of $\mathcal{O}_X$-modules
$\SheafHom^\bullet(\mathcal{L}^\bullet, \mathcal{M}^\bullet)$.
Namely, for each $n$ we set
$$
\SheafHom^n(\mathcal{L}^\bullet, \mathcal{M}^\bullet) =
\prod\nolimits_{n = p + q}
\SheafHom_{\mathcal{O}_X}(\mathcal{L}^{-q}, \mathcal{M}^p)
$$
It is a good idea to think of $\SheafHom^n$ as the
sheaf of $\mathcal{O}_X$-modules of all $\mathcal{O}_X$-linear
maps from $\mathcal{L}^\bullet$ to $\mathcal{M}^\bullet$
(viewed as graded $\mathcal{O}_X$-modules) which are homogeneous
of degree $n$. In this terminology, we define the differential by the rule
$$
\text{d}(f) =
\text{d}_\mathcal{M} \circ f - (-1)^n f \circ \text{d}_\mathcal{L}
$$
for
$f \in \SheafHom^n_{\mathcal{O}_X}(\mathcal{L}^\bullet, \mathcal{M}^\bullet)$.
We omit the verification that $\text{d}^2 = 0$.
This construction is a special case of
Differential Graded Algebra, Example \ref{dga-example-category-complexes}.
It follows immediately from the construction that we have
\begin{equation}
\label{equation-cohomology-hom-complex}
H^n(\Gamma(U, \SheafHom^\bullet(\mathcal{L}^\bullet, \mathcal{M}^\bullet))) =
\Hom_{K(\mathcal{O}_U)}(\mathcal{L}^\bullet, \mathcal{M}^\bullet[n])
\end{equation}
for all $n \in \mathbf{Z}$ and every open $U \subset X$.

\begin{lemma}
\label{lemma-compose}
Let $(X, \mathcal{O}_X)$ be a ringed space.
Given complexes $\mathcal{K}^\bullet, \mathcal{L}^\bullet, \mathcal{M}^\bullet$
of $\mathcal{O}_X$-modules there is an isomorphism
$$
\SheafHom^\bullet(\mathcal{K}^\bullet,
\SheafHom^\bullet(\mathcal{L}^\bullet, \mathcal{M}^\bullet))
=
\SheafHom^\bullet(\text{Tot}(\mathcal{K}^\bullet \otimes_{\mathcal{O}_X}
\mathcal{L}^\bullet), \mathcal{M}^\bullet)
$$
of complexes of $\mathcal{O}_X$-modules functorial in
$\mathcal{K}^\bullet, \mathcal{L}^\bullet, \mathcal{M}^\bullet$.
\end{lemma}

\begin{proof}
Omitted. Hint: This is proved in exactly the same way as
More on Algebra, Lemma \ref{more-algebra-lemma-compose}.
\end{proof}

\begin{lemma}
\label{lemma-composition}
Let $(X, \mathcal{O}_X)$ be a ringed space. Given complexes
$\mathcal{K}^\bullet, \mathcal{L}^\bullet, \mathcal{M}^\bullet$
of $\mathcal{O}_X$-modules there is a canonical morphism
$$
\text{Tot}\left(
\SheafHom^\bullet(\mathcal{L}^\bullet, \mathcal{M}^\bullet)
\otimes_{\mathcal{O}_X}
\SheafHom^\bullet(\mathcal{K}^\bullet, \mathcal{L}^\bullet)
\right)
\longrightarrow
\SheafHom^\bullet(\mathcal{K}^\bullet, \mathcal{M}^\bullet)
$$
of complexes of $\mathcal{O}_X$-modules.
\end{lemma}

\begin{proof}
Omitted. Hint: This is proved in exactly the same way as
More on Algebra, Lemma \ref{more-algebra-lemma-composition}.
\end{proof}

\begin{lemma}
\label{lemma-diagonal-better}
Let $(X, \mathcal{O}_X)$ be a ringed space. Given complexes
$\mathcal{K}^\bullet, \mathcal{L}^\bullet, \mathcal{M}^\bullet$
of $\mathcal{O}_X$-modules there is a canonical morphism
$$
\text{Tot}\left(
\mathcal{K}^\bullet \otimes_{\mathcal{O}_X}
\SheafHom^\bullet(\mathcal{M}^\bullet, \mathcal{L}^\bullet)
\right)
\longrightarrow
\SheafHom^\bullet(\mathcal{M}^\bullet,
\text{Tot}(\mathcal{K}^\bullet \otimes_{\mathcal{O}_X} \mathcal{L}^\bullet))
$$
of complexes of $\mathcal{O}_X$-modules functorial in all three complexes.
\end{lemma}

\begin{proof}
Omitted. Hint: This is proved in exactly the same way as
More on Algebra, Lemma \ref{more-algebra-lemma-diagonal-better}.
\end{proof}

\begin{lemma}
\label{lemma-diagonal}
Let $(X, \mathcal{O}_X)$ be a ringed space. Given complexes
$\mathcal{K}^\bullet, \mathcal{L}^\bullet$
of $\mathcal{O}_X$-modules there is a canonical morphism
$$
\mathcal{K}^\bullet
\longrightarrow
\SheafHom^\bullet(\mathcal{L}^\bullet,
\text{Tot}(\mathcal{K}^\bullet \otimes_{\mathcal{O}_X} \mathcal{L}^\bullet))
$$
of complexes of $\mathcal{O}_X$-modules functorial in both complexes.
\end{lemma}

\begin{proof}
Omitted. Hint: This is proved in exactly the same way as
More on Algebra, Lemma \ref{more-algebra-lemma-diagonal}.
\end{proof}

\begin{lemma}
\label{lemma-evaluate-and-more}
Let $(X, \mathcal{O}_X)$ be a ringed space. Given complexes
$\mathcal{K}^\bullet, \mathcal{L}^\bullet, \mathcal{M}^\bullet$
of $\mathcal{O}_X$-modules there is a canonical morphism
$$
\text{Tot}(\SheafHom^\bullet(\mathcal{L}^\bullet,
\mathcal{M}^\bullet) \otimes_{\mathcal{O}_X} \mathcal{K}^\bullet)
\longrightarrow
\SheafHom^\bullet(\SheafHom^\bullet(\mathcal{K}^\bullet,
\mathcal{L}^\bullet), \mathcal{M}^\bullet)
$$
of complexes of $\mathcal{O}_X$-modules functorial in all three complexes.
\end{lemma}

\begin{proof}
Omitted. Hint: This is proved in exactly the same way as
More on Algebra, Lemma \ref{more-algebra-lemma-evaluate-and-more}.
\end{proof}

\begin{lemma}
\label{lemma-RHom-into-K-injective}
Let $(X, \mathcal{O}_X)$ be a ringed space. Let $L$ and $M$
be objects of $D(\mathcal{O}_X)$. Let $\mathcal{I}^\bullet$
be a K-injective complex of $\mathcal{O}_X$-modules representing $M$. Let
$\mathcal{L}^\bullet$ be a complex of $\mathcal{O}_X$-modules
representing $L$. Then
$$
H^0(\Gamma(U, \SheafHom^\bullet(\mathcal{L}^\bullet, \mathcal{I}^\bullet))) =
\Hom_{D(\mathcal{O}_U)}(L|_U, M|_U)
$$
for all $U \subset X$ open.
\end{lemma}

\begin{proof}
We have
\begin{align*}
H^0(\Gamma(U, \SheafHom^\bullet(\mathcal{L}^\bullet, \mathcal{I}^\bullet)))
& =
\Hom_{K(\mathcal{O}_U)}(\mathcal{L}^\bullet|_U, \mathcal{I}^\bullet|_U) \\
& =
\Hom_{D(\mathcal{O}_U)}(L|_U, M|_U)
\end{align*}
The first equality is (\ref{equation-cohomology-hom-complex}).
The second equality is true because $\mathcal{I}^\bullet|_U$
is K-injective by Lemma \ref{lemma-restrict-K-injective-to-open}.
\end{proof}

\begin{lemma}
\label{lemma-RHom-well-defined}
Let $(X, \mathcal{O}_X)$ be a ringed space. Let
$(\mathcal{I}')^\bullet \to \mathcal{I}^\bullet$
be a quasi-isomorphism of K-injective complexes of $\mathcal{O}_X$-modules.
Let $(\mathcal{L}')^\bullet \to \mathcal{L}^\bullet$
be a quasi-isomorphism of complexes of $\mathcal{O}_X$-modules.
Then
$$
\SheafHom^\bullet(\mathcal{L}^\bullet, (\mathcal{I}')^\bullet)
\longrightarrow
\SheafHom^\bullet((\mathcal{L}')^\bullet, \mathcal{I}^\bullet)
$$
is a quasi-isomorphism.
\end{lemma}

\begin{proof}
Let $M$ be the object of $D(\mathcal{O}_X)$ represented by
$\mathcal{I}^\bullet$ and $(\mathcal{I}')^\bullet$.
Let $L$ be the object of $D(\mathcal{O}_X)$ represented by
$\mathcal{L}^\bullet$ and $(\mathcal{L}')^\bullet$.
By Lemma \ref{lemma-RHom-into-K-injective}
we see that the sheaves
$$
H^0(\SheafHom^\bullet(\mathcal{L}^\bullet, (\mathcal{I}')^\bullet))
\quad\text{and}\quad
H^0(\SheafHom^\bullet((\mathcal{L}')^\bullet, \mathcal{I}^\bullet))
$$
are both equal to the sheaf associated to the presheaf
$$
U \longmapsto \Hom_{D(\mathcal{O}_U)}(L|_U, M|_U)
$$
Thus the map is a quasi-isomorphism.
\end{proof}

\begin{lemma}
\label{lemma-RHom-from-K-flat-into-K-injective}
Let $(X, \mathcal{O}_X)$ be a ringed space. Let $\mathcal{I}^\bullet$
be a K-injective complex of $\mathcal{O}_X$-modules. Let
$\mathcal{L}^\bullet$ be a K-flat complex of $\mathcal{O}_X$-modules.
Then $\SheafHom^\bullet(\mathcal{L}^\bullet, \mathcal{I}^\bullet)$
is a K-injective complex of $\mathcal{O}_X$-modules.
\end{lemma}

\begin{proof}
Namely, if $\mathcal{K}^\bullet$ is an acyclic complex of
$\mathcal{O}_X$-modules, then
\begin{align*}
\Hom_{K(\mathcal{O}_X)}(\mathcal{K}^\bullet,
\SheafHom^\bullet(\mathcal{L}^\bullet, \mathcal{I}^\bullet))
& =
H^0(\Gamma(X,
\SheafHom^\bullet(\mathcal{K}^\bullet,
\SheafHom^\bullet(\mathcal{L}^\bullet, \mathcal{I}^\bullet)))) \\
& =
H^0(\Gamma(X, \SheafHom^\bullet(\text{Tot}(
\mathcal{K}^\bullet \otimes_{\mathcal{O}_X} \mathcal{L}^\bullet),
\mathcal{I}^\bullet))) \\
& =
\Hom_{K(\mathcal{O}_X)}(
\text{Tot}(\mathcal{K}^\bullet \otimes_{\mathcal{O}_X} \mathcal{L}^\bullet),
\mathcal{I}^\bullet) \\
& =
0
\end{align*}
The first equality by (\ref{equation-cohomology-hom-complex}).
The second equality by Lemma \ref{lemma-compose}.
The third equality by (\ref{equation-cohomology-hom-complex}).
The final equality because
$\text{Tot}(\mathcal{K}^\bullet \otimes_{\mathcal{O}_X} \mathcal{L}^\bullet)$
is acyclic because $\mathcal{L}^\bullet$ is K-flat
(Definition \ref{definition-K-flat}) and because $\mathcal{I}^\bullet$
is K-injective.
\end{proof}








\section{Internal hom in the derived category}
\label{section-internal-hom}

\noindent
Let $(X, \mathcal{O}_X)$ be a ringed space. Let $L, M$ be objects
of $D(\mathcal{O}_X)$. We would like to construct an object
$R\SheafHom(L, M)$ of $D(\mathcal{O}_X)$ such that for every third
object $K$ of $D(\mathcal{O}_X)$ there exists a canonical bijection
\begin{equation}
\label{equation-internal-hom}
\Hom_{D(\mathcal{O}_X)}(K, R\SheafHom(L, M))
=
\Hom_{D(\mathcal{O}_X)}(K \otimes_{\mathcal{O}_X}^\mathbf{L} L, M)
\end{equation}
Observe that this formula defines $R\SheafHom(L, M)$ up to unique
isomorphism by the Yoneda lemma
(Categories, Lemma \ref{categories-lemma-yoneda}).

\medskip\noindent
To construct such an object, choose a K-injective complex
$\mathcal{I}^\bullet$ representing $M$ and any complex of
$\mathcal{O}_X$-modules $\mathcal{L}^\bullet$ representing $L$.
Then we set
$$
R\SheafHom(L, M) = \SheafHom^\bullet(\mathcal{L}^\bullet, \mathcal{I}^\bullet)
$$
where the right hand side is the complex of $\mathcal{O}_X$-modules
constructed in Section \ref{section-hom-complexes}.
This is well defined by Lemma \ref{lemma-RHom-well-defined}.
We get a functor
$$
D(\mathcal{O}_X)^{opp} \times D(\mathcal{O}_X) \longrightarrow D(\mathcal{O}_X),
\quad
(K, L) \longmapsto R\SheafHom(K, L)
$$
As a prelude to proving (\ref{equation-internal-hom})
we compute the cohomology groups of $R\SheafHom(K, L)$.

\begin{lemma}
\label{lemma-section-RHom-over-U}
Let $(X, \mathcal{O}_X)$ be a ringed space. Let $L, M$ be objects
of $D(\mathcal{O}_X)$. For every open $U$ we have
$$
H^0(U, R\SheafHom(L, M)) =
\Hom_{D(\mathcal{O}_U)}(L|_U, M|_U)
$$
and in particular $H^0(X, R\SheafHom(L, M)) = \Hom_{D(\mathcal{O}_X)}(L, M)$.
\end{lemma}

\begin{proof}
Choose a K-injective complex $\mathcal{I}^\bullet$ of
$\mathcal{O}_X$-modules representing $M$ and a K-flat complex
$\mathcal{L}^\bullet$ representing $L$. Then
$\SheafHom^\bullet(\mathcal{L}^\bullet, \mathcal{I}^\bullet)$
is K-injective by Lemma \ref{lemma-RHom-from-K-flat-into-K-injective}.
Hence we can compute cohomology over $U$ by simply taking sections over $U$
and the result follows from Lemma \ref{lemma-RHom-into-K-injective}.
\end{proof}

\begin{lemma}
\label{lemma-internal-hom}
Let $(X, \mathcal{O}_X)$ be a ringed space. Let $K, L, M$ be objects
of $D(\mathcal{O}_X)$. With the construction as described above
there is a canonical isomorphism
$$
R\SheafHom(K, R\SheafHom(L, M)) =
R\SheafHom(K \otimes_{\mathcal{O}_X}^\mathbf{L} L, M)
$$
in $D(\mathcal{O}_X)$ functorial in $K, L, M$
which recovers (\ref{equation-internal-hom}) by taking $H^0(X, -)$.
\end{lemma}

\begin{proof}
Choose a K-injective complex $\mathcal{I}^\bullet$ representing
$M$ and a K-flat complex of $\mathcal{O}_X$-modules $\mathcal{L}^\bullet$
representing $L$. Let $\mathcal{K}^\bullet$ be any complex of
$\mathcal{O}_X$-modules representing $K$. Then
we have
$$
\SheafHom^\bullet(\mathcal{K}^\bullet,
\SheafHom^\bullet(\mathcal{L}^\bullet, \mathcal{I}^\bullet))
=
\SheafHom^\bullet(
\text{Tot}(\mathcal{K}^\bullet \otimes_{\mathcal{O}_X} \mathcal{L}^\bullet),
\mathcal{I}^\bullet)
$$
by Lemma \ref{lemma-compose}.
Note that the left hand side represents
$R\SheafHom(K, R\SheafHom(L, M))$ (use
Lemma \ref{lemma-RHom-from-K-flat-into-K-injective})
and that the right hand side represents
$R\SheafHom(K \otimes_{\mathcal{O}_X}^\mathbf{L} L, M)$.
This proves the displayed formula of the lemma.
Taking global sections and using Lemma \ref{lemma-section-RHom-over-U}
we obtain (\ref{equation-internal-hom}).
\end{proof}

\begin{lemma}
\label{lemma-restriction-RHom-to-U}
Let $(X, \mathcal{O}_X)$ be a ringed space. Let $K, L$ be objects
of $D(\mathcal{O}_X)$. The construction of $R\SheafHom(K, L)$
commutes with restrictions to opens, i.e.,
for every open $U$ we have
$R\SheafHom(K|_U, L|_U) = R\SheafHom(K, L)|_U$.
\end{lemma}

\begin{proof}
This is clear from the construction and
Lemma \ref{lemma-restrict-K-injective-to-open}.
\end{proof}

\begin{lemma}
\label{lemma-RHom-triangulated}
Let $(X, \mathcal{O}_X)$ be a ringed space. The bifunctor $R\SheafHom(- , -)$
transforms distinguished triangles into distinguished triangles in both
variables.
\end{lemma}

\begin{proof}
This follows from the observation that the assignment
$$
(\mathcal{L}^\bullet, \mathcal{M}^\bullet) \longmapsto
\SheafHom^\bullet(\mathcal{L}^\bullet, \mathcal{M}^\bullet)
$$
transforms a termwise split short exact sequences of complexes in either
variable into a termwise split short exact sequence. Details omitted.
\end{proof}

\begin{lemma}
\label{lemma-internal-hom-composition}
Let $(X, \mathcal{O}_X)$ be a ringed space. Given $K, L, M$ in
$D(\mathcal{O}_X)$ there is a canonical morphism
$$
R\SheafHom(L, M) \otimes_{\mathcal{O}_X}^\mathbf{L} R\SheafHom(K, L)
\longrightarrow R\SheafHom(K, M)
$$
in $D(\mathcal{O}_X)$ functorial in $K, L, M$.
\end{lemma}

\begin{proof}
Choose a K-injective complex $\mathcal{I}^\bullet$ representing $M$,
a K-injective complex $\mathcal{J}^\bullet$ representing $L$, and
any complex of $\mathcal{O}_X$-modules $\mathcal{K}^\bullet$ representing $K$.
By Lemma \ref{lemma-composition} there is a map of complexes
$$
\text{Tot}\left(
\SheafHom^\bullet(\mathcal{J}^\bullet, \mathcal{I}^\bullet)
\otimes_{\mathcal{O}_X}
\SheafHom^\bullet(\mathcal{K}^\bullet, \mathcal{J}^\bullet)
\right)
\longrightarrow
\SheafHom^\bullet(\mathcal{K}^\bullet, \mathcal{I}^\bullet)
$$
The complexes of $\mathcal{O}_X$-modules
$\SheafHom^\bullet(\mathcal{J}^\bullet, \mathcal{I}^\bullet)$,
$\SheafHom^\bullet(\mathcal{K}^\bullet, \mathcal{J}^\bullet)$, and
$\SheafHom^\bullet(\mathcal{K}^\bullet, \mathcal{I}^\bullet)$
represent $R\SheafHom(L, M)$, $R\SheafHom(K, L)$, and $R\SheafHom(K, M)$.
If we choose a K-flat complex $\mathcal{H}^\bullet$ and a quasi-isomorphism
$\mathcal{H}^\bullet \to
\SheafHom^\bullet(\mathcal{K}^\bullet, \mathcal{J}^\bullet)$,
then there is a map
$$
\text{Tot}\left(
\SheafHom^\bullet(\mathcal{J}^\bullet, \mathcal{I}^\bullet)
\otimes_{\mathcal{O}_X} \mathcal{H}^\bullet
\right)
\longrightarrow
\text{Tot}\left(
\SheafHom^\bullet(\mathcal{J}^\bullet, \mathcal{I}^\bullet)
\otimes_{\mathcal{O}_X}
\SheafHom^\bullet(\mathcal{K}^\bullet, \mathcal{J}^\bullet)
\right)
$$
whose source represents
$R\SheafHom(L, M) \otimes_{\mathcal{O}_X}^\mathbf{L} R\SheafHom(K, L)$.
Composing the two displayed arrows gives the desired map. We omit the
proof that the construction is functorial.
\end{proof}

\begin{lemma}
\label{lemma-internal-hom-diagonal-better}
Let $(X, \mathcal{O}_X)$ be a ringed space. Given $K, L, M$
in $D(\mathcal{O}_X)$ there is a canonical morphism
$$
K \otimes_{\mathcal{O}_X}^\mathbf{L} R\SheafHom(M, L)
\longrightarrow
R\SheafHom(M, K \otimes_{\mathcal{O}_X}^\mathbf{L} L)
$$
in $D(\mathcal{O}_X)$ functorial in $K, L, M$.
\end{lemma}

\begin{proof}
Choose a K-flat complex $\mathcal{K}^\bullet$ representing $K$,
and a K-injective complex $\mathcal{I}^\bullet$ representing $L$, and
choose any complex of $\mathcal{O}_X$-modules $\mathcal{M}^\bullet$
representing $M$. Choose a quasi-isomorphism
$\text{Tot}(\mathcal{K}^\bullet \otimes_{\mathcal{O}_X} \mathcal{I}^\bullet)
\to \mathcal{J}^\bullet$
where $\mathcal{J}^\bullet$ is K-injective. Then we use the map
$$
\text{Tot}\left(
\mathcal{K}^\bullet \otimes_{\mathcal{O}_X}
\SheafHom^\bullet(\mathcal{M}^\bullet, \mathcal{I}^\bullet)
\right)
\to
\SheafHom^\bullet(\mathcal{M}^\bullet,
\text{Tot}(\mathcal{K}^\bullet \otimes_{\mathcal{O}_X} \mathcal{I}^\bullet))
\to
\SheafHom^\bullet(\mathcal{M}^\bullet, \mathcal{J}^\bullet)
$$
where the first map is the map from Lemma \ref{lemma-diagonal-better}.
\end{proof}

\begin{lemma}
\label{lemma-internal-hom-diagonal}
Let $(X, \mathcal{O}_X)$ be a ringed space. Given $K, L$ in $D(\mathcal{O}_X)$
there is a canonical morphism
$$
K \longrightarrow R\SheafHom(L, K \otimes_{\mathcal{O}_X}^\mathbf{L} L)
$$
in $D(\mathcal{O}_X)$ functorial in both $K$ and $L$.
\end{lemma}

\begin{proof}
Choose a K-flat complex $\mathcal{K}^\bullet$ representing $K$
and any complex of $\mathcal{O}_X$-modules $\mathcal{L}^\bullet$
representing $L$. Choose a K-injective complex $\mathcal{J}^\bullet$
and a quasi-isomorphism
$\text{Tot}(\mathcal{K}^\bullet \otimes_{\mathcal{O}_X} \mathcal{L}^\bullet)
\to \mathcal{J}^\bullet$. Then we use
$$
\mathcal{K}^\bullet \to
\SheafHom^\bullet(\mathcal{L}^\bullet,
\text{Tot}(\mathcal{K}^\bullet \otimes_{\mathcal{O}_X} \mathcal{L}^\bullet))
\to
\SheafHom^\bullet(\mathcal{L}^\bullet, \mathcal{J}^\bullet)
$$
where the first map comes from Lemma \ref{lemma-diagonal}.
\end{proof}

\begin{lemma}
\label{lemma-dual}
Let $(X, \mathcal{O}_X)$ be a ringed space. Let $L$ be an
object of $D(\mathcal{O}_X)$. Set $L^\vee = R\SheafHom(L, \mathcal{O}_X)$.
For $M$ in $D(\mathcal{O}_X)$ there is a canonical map
\begin{equation}
\label{equation-eval}
M \otimes^\mathbf{L}_{\mathcal{O}_X} L^\vee
\longrightarrow
R\SheafHom(L, M)
\end{equation}
which induces a canonical map
$$
H^0(X, M \otimes^\mathbf{L}_{\mathcal{O}_X} L^\vee)
\longrightarrow
\Hom_{D(\mathcal{O}_X)}(L, M)
$$
functorial in $M$ in $D(\mathcal{O}_X)$.
\end{lemma}

\begin{proof}
The map (\ref{equation-eval}) is a special case of
Lemma \ref{lemma-internal-hom-composition}
using the identification $M = R\SheafHom(\mathcal{O}_X, M)$.
\end{proof}

\begin{lemma}
\label{lemma-internal-hom-evaluate}
Let $(X, \mathcal{O}_X)$ be a ringed space. Let $K, L, M$ be objects of
$D(\mathcal{O}_X)$. There is a canonical morphism
$$
R\SheafHom(L, M) \otimes_{\mathcal{O}_X}^\mathbf{L} K
\longrightarrow
R\SheafHom(R\SheafHom(K, L), M)
$$
in $D(\mathcal{O}_X)$ functorial in $K, L, M$.
\end{lemma}

\begin{proof}
Choose
a K-injective complex $\mathcal{I}^\bullet$ representing $M$,
a K-injective complex $\mathcal{J}^\bullet$ representing $L$, and
a K-flat complex $\mathcal{K}^\bullet$ representing $K$.
The map is defined using the map
$$
\text{Tot}(\SheafHom^\bullet(\mathcal{J}^\bullet,
\mathcal{I}^\bullet) \otimes_{\mathcal{O}_X} \mathcal{K}^\bullet)
\longrightarrow
\SheafHom^\bullet(\SheafHom^\bullet(\mathcal{K}^\bullet,
\mathcal{J}^\bullet), \mathcal{I}^\bullet)
$$
of Lemma \ref{lemma-evaluate-and-more}. By our particular
choice of complexes the left hand side represents
$R\SheafHom(L, M) \otimes_{\mathcal{O}_X}^\mathbf{L} K$
and the right hand side represents
$R\SheafHom(R\SheafHom(K, L), M)$. We omit the proof that
this is functorial in all three objects of $D(\mathcal{O}_X)$.
\end{proof}

\begin{remark}
\label{remark-tensor-internal-hom}
Let $(X, \mathcal{O}_X)$ be a ringed space. For $K, K', M, M'$ in
$D(\mathcal{O}_X)$ there is a canonical map
$$
R\SheafHom(K, K') \otimes_{\mathcal{O}_X}^\mathbf{L}
R\SheafHom(M, M')
\longrightarrow
R\SheafHom(K \otimes_{\mathcal{O}_X}^\mathbf{L} M,
K' \otimes_{\mathcal{O}_X}^\mathbf{L} M')
$$
Namely, by (\ref{equation-internal-hom}) is the same thing as a map
$$
R\SheafHom(K, K') \otimes_{\mathcal{O}_X}^\mathbf{L}
R\SheafHom(M, M') \otimes_{\mathcal{O}_X}^\mathbf{L}
K \otimes_{\mathcal{O}_X}^\mathbf{L} M
\longrightarrow
K' \otimes_{\mathcal{O}_X}^\mathbf{L} M'
$$
For this we can first flip the middle two factors
(with sign rules as in More on Algebra, Section
\ref{more-algebra-section-sign-rules})
and use the maps
$$
R\SheafHom(K, K') \otimes_{\mathcal{O}_X}^\mathbf{L} K \to K'
\quad\text{and}\quad
R\SheafHom(M, M') \otimes_{\mathcal{O}_X}^\mathbf{L} M \to M'
$$
from Lemma \ref{lemma-internal-hom-composition} when thinking
of $K = R\SheafHom(\mathcal{O}_X, K)$ and similarly for
$K'$, $M$, and $M'$.
\end{remark}

\begin{remark}
\label{remark-projection-formula-for-internal-hom}
Let $f : X \to Y$ be a morphism of ringed spaces.
Let $K, L$ be objects of $D(\mathcal{O}_X)$. We claim there is a canonical map
$$
Rf_*R\SheafHom(L, K) \longrightarrow R\SheafHom(Rf_*L, Rf_*K)
$$
Namely, by (\ref{equation-internal-hom}) this is the same thing
as a map
$Rf_*R\SheafHom(L, K) \otimes_{\mathcal{O}_Y}^\mathbf{L} Rf_*L \to Rf_*K$.
For this we can use the composition
$$
Rf_*R\SheafHom(L, K) \otimes_{\mathcal{O}_Y}^\mathbf{L} Rf_*L \to
Rf_*(R\SheafHom(L, K) \otimes_{\mathcal{O}_X}^\mathbf{L} L) \to
Rf_*K
$$
where the first arrow is the relative cup product
(Remark \ref{remark-cup-product}) and the second arrow is $Rf_*$ applied
to the canonical map
$R\SheafHom(L, K) \otimes_{\mathcal{O}_X}^\mathbf{L} L \to K$
coming from Lemma \ref{lemma-internal-hom-composition}
(with $\mathcal{O}_X$ in one of the spots).
\end{remark}

\begin{remark}
\label{remark-relative-cup-and-composition}
Let $h : X \to Y$ be a morphism of ringed spaces.
Let $K, M$ be objects of $D(\mathcal{O}_Y)$.
The diagram
$$
\xymatrix{
Rf_*R\SheafHom_{\mathcal{O}_X}(K, M)
\otimes_{\mathcal{O}_Y}^\mathbf{L} Rf_*K
\ar[r] \ar[d] &
Rf_*\left(R\SheafHom_{\mathcal{O}_X}(K, M)
\otimes_{\mathcal{O}_X}^\mathbf{L} K\right)
\ar[d] \\
R\SheafHom_{\mathcal{O}_Y}(Rf_*K, Rf_*M) \otimes_{\mathcal{O}_Y}^\mathbf{L}
Rf_*K \ar[r] &
Rf_*M
}
$$
is commutative. Here the left vertical arrow comes from
Remark \ref{remark-projection-formula-for-internal-hom}.
The top horizontal arrow is Remark \ref{remark-cup-product}.
The other two arrows are instances of the map in
Lemma \ref{lemma-internal-hom-composition} (with one of the entries
replaced with $\mathcal{O}_X$ or $\mathcal{O}_Y$).
\end{remark}

\begin{remark}
\label{remark-prepare-fancy-base-change}
Let $h : X \to Y$ be a morphism of ringed spaces.
Let $K, L$ be objects of $D(\mathcal{O}_Y)$. We claim there is a
canonical map
$$
Lh^*R\SheafHom(K, L) \longrightarrow R\SheafHom(Lh^*K, Lh^*L)
$$
in $D(\mathcal{O}_X)$. Namely, by (\ref{equation-internal-hom})
proved in Lemma \ref{lemma-internal-hom}
such a map is the same thing as a map
$$
Lh^*R\SheafHom(K, L) \otimes^\mathbf{L} Lh^*K \longrightarrow Lh^*L
$$
The source of this arrow is $Lh^*(\SheafHom(K, L) \otimes^\mathbf{L} K)$
by Lemma \ref{lemma-pullback-tensor-product}
hence it suffices to construct a canonical map
$$
R\SheafHom(K, L) \otimes^\mathbf{L} K \longrightarrow L.
$$
For this we take the arrow corresponding to
$$
\text{id} :
R\SheafHom(K, L)
\longrightarrow
R\SheafHom(K, L)
$$
via (\ref{equation-internal-hom}).
\end{remark}

\begin{remark}
\label{remark-fancy-base-change}
Suppose that
$$
\xymatrix{
X' \ar[r]_h \ar[d]_{f'} &
X \ar[d]^f \\
S' \ar[r]^g &
S
}
$$
is a commutative diagram of ringed spaces. Let $K, L$ be objects
of $D(\mathcal{O}_X)$. We claim there exists a canonical base change
map
$$
Lg^*Rf_*R\SheafHom(K, L)
\longrightarrow
R(f')_*R\SheafHom(Lh^*K, Lh^*L)
$$
in $D(\mathcal{O}_{S'})$. Namely, we take the map adjoint to
the composition
\begin{align*}
L(f')^*Lg^*Rf_*R\SheafHom(K, L)
& =
Lh^*Lf^*Rf_*R\SheafHom(K, L) \\
& \to
Lh^*R\SheafHom(K, L) \\
& \to
R\SheafHom(Lh^*K, Lh^*L)
\end{align*}
where the first arrow uses the adjunction mapping
$Lf^*Rf_* \to \text{id}$ and the second arrow is the canonical map
constructed in Remark \ref{remark-prepare-fancy-base-change}.
\end{remark}





\section{Ext sheaves}
\label{section-ext}

\noindent
Let $(X, \mathcal{O}_X)$ be a ringed space. Let $K, L \in D(\mathcal{O}_X)$.
Using the construction of the internal hom in the derived category we
obtain a well defined sheaves of $\mathcal{O}_X$-modules
$$
\SheafExt^n(K, L) = H^n(R\SheafHom(K, L))
$$
by taking the $n$th cohomology sheaf of the object $R\SheafHom(K, L)$
of $D(\mathcal{O}_X)$. We will sometimes write
$\SheafExt^n_{\mathcal{O}_X}(K, L)$ for this object.
By Lemma \ref{lemma-section-RHom-over-U}
we see that this $\SheafExt^n$-sheaf
is the sheafification of the rule
$$
U \longmapsto \Ext^n_{D(\mathcal{O}_U)}(K|_U, L|_U)
$$
By Example \ref{example-spectral-sequence} there is always a spectral
sequence
$$
E_2^{p, q} = H^p(X, \SheafExt^q(K, L))
$$
converging to $\Ext^{p + q}_{D(\mathcal{O}_X)}(K, L)$
in favorable situations (for example if $L$ is bounded below and
$K$ is bounded above).

\begin{lemma}[Centre actions on Ext for left modules]
\label{lemma-ext-center-actions-left-modules}
\begin{reference}
\cite[Chapter IV, \S 4.2, p. 188]{Tohoku}
\end{reference}
\begin{history}
The printed source makes the global Ext groups modules over
$Z(\Gamma(X,\mathcal A))$. For a sheaf of noncommutative rings, a globally
central section need not remain central after restriction, so that action is
not defined in general. The canonical scalar ring is instead
$\Gamma(X,\mathcal Z_{\mathcal A})$, where
$\mathcal Z_{\mathcal A}$ is the sheaf-theoretic centre from Modules, Lemma
\ref{modules-lemma-sheaf-hom-left-modules}. The two rings agree when
$\mathcal A$ is commutative.
\end{history}
Let $X$ be a topological space, let $\mathcal A$ be a sheaf of not
necessarily commutative rings with $1$, and let $\mathcal F,\mathcal G$ be
sheaves of unital left $\mathcal A$-modules. Write
$\mathcal Z_{\mathcal A}$ for the sheaf-theoretic centre of $\mathcal A$.
For every $p\geq0$, the global Ext group
\[
\Ext^p_{\mathcal A}(X;\mathcal F,\mathcal G)
=
\Ext^p_{\textit{Mod}(\mathcal A)}(\mathcal F,\mathcal G)
\]
is canonically a $\Gamma(X,\mathcal Z_{\mathcal A})$-module, and the
right-derived sheaf
\[
\SheafExt^p_{\mathcal A}(\mathcal F,\mathcal G)
=
R^p\SheafHom_{\mathcal A}(\mathcal F,-)(\mathcal G)
\]
is canonically a sheaf of $\mathcal Z_{\mathcal A}$-modules. These actions
are functorial in both module arguments and agree in degree zero with the
actions of Modules, Lemma
\ref{modules-lemma-sheaf-hom-left-modules}.

If the canonical inclusion
\[
\Gamma(X,\mathcal Z_{\mathcal A})
\longrightarrow Z(\Gamma(X,\mathcal A))
\]
is an equality---in particular, if $\mathcal A$ is commutative---then the
global Ext groups are modules over $Z(\Gamma(X,\mathcal A))$.
\end{lemma}

\begin{proof}
The category of sheaves of unital left $\mathcal A$-modules is abelian and
has enough injectives; the proof of Injectives, Lemma
\ref{injectives-lemma-sheaves-modules-space} does not use commutativity.
Choose an injective resolution
\[
\mathcal G[0]\longrightarrow\mathcal I^\bullet.
\]
For every $q$, Modules, Lemma
\ref{modules-lemma-sheaf-hom-left-modules} makes
$\SheafHom_{\mathcal A}(\mathcal F,\mathcal I^q)$ a sheaf of
$\mathcal Z_{\mathcal A}$-modules. Since the differentials of
$\mathcal I^\bullet$ are $\mathcal A$-linear, the induced differentials on
these Hom sheaves are $\mathcal Z_{\mathcal A}$-linear. Their cohomology
sheaves are therefore $\mathcal Z_{\mathcal A}$-modules, and they compute
the displayed right-derived sheaves.

Taking global sections gives a complex
\[
\Hom_{\mathcal A}(\mathcal F,\mathcal I^\bullet)
=
\Gamma\bigl(X,
\SheafHom_{\mathcal A}(\mathcal F,\mathcal I^\bullet)\bigr)
\]
of $\Gamma(X,\mathcal Z_{\mathcal A})$-modules. Its degree-$p$ cohomology is
$\Ext^p_{\textit{Mod}(\mathcal A)}(\mathcal F,\mathcal G)$, which gives the
global action. The construction is natural in $\mathcal G$ by functoriality
of derived functors, and natural in $\mathcal F$ by precomposition. Degree
zero is exactly the action constructed in the cited Modules lemma. The final
assertion follows by transporting scalars through the displayed equality.
\end{proof}


\begin{lemma}
\label{lemma-constant-sheaf-ext-double-complex}
\begin{reference}
\cite[Chapter IV, \S 4.3, Proposition 4.3.2, p. 191]{Tohoku}
\end{reference}
Let $X$ be a topological space, let $R$ be a ring, let $M$ be an
$R$-module, and let $\mathcal{A}$ be an $\underline{R}$-module on $X$.
Choose a free resolution
$$
\ldots \longrightarrow L_2 \longrightarrow L_1 \longrightarrow
L_0 \longrightarrow M \longrightarrow 0
$$
and let $\mathcal{A} \to \mathcal{C}^\bullet(\mathcal{A})$ be the
Godement resolution of Lemma \ref{lemma-godement-resolution}. There is a
first quadrant double complex
$$
K^{p, q} = \Hom_R(L_p, \Gamma(X, \mathcal{C}^q(\mathcal{A}))),
\qquad p, q \geq 0.
$$
One of its spectral sequences has
$$
{}'E_2^{p, q} = \Ext_R^p(M, H^q(X, \mathcal{A})).
$$
If every $L_p$ is finite free, then the other spectral sequence has
$$
{}''E_2^{p, q} =
H^p(X, \SheafExt^q_{\underline{R}}(\underline{M}, \mathcal{A})).
$$
In this case it is the local-to-global Ext spectral sequence described
above.
Both spectral sequences are bounded, converge to
$$
\Ext^{p + q}_{D(\underline{R})}(\underline{M}, \mathcal{A}),
$$
and induce finite filtrations on these groups. From the $E_2$-page onward
they are functorial in $\mathcal{A}$.
\end{lemma}

\begin{proof}
Write $L^{-p} = L_p$ and regard $L^\bullet$ as a bounded above complex
of projective $R$-modules. The complex
$\Gamma(X, \mathcal{C}^\bullet(\mathcal{A}))$ represents
$R\Gamma(X, \mathcal{A})$: the Godement resolution is a resolution by
flasque modules, and flasque modules are acyclic by
Lemma \ref{lemma-flasque-acyclic}. Hence the total complex of
$K^{\bullet, \bullet}$ represents
$$
R\Hom_R(M, R\Gamma(X, \mathcal{A})).
$$
Indeed, this follows by using $L^\bullet$ and
Derived Categories, Lemma
\ref{derived-lemma-morphisms-from-projective-complex}.

The constant sheaf functor is exact by Sites, Modules, Lemma
\ref{sites-modules-lemma-constant-exact} and is left adjoint to global
sections by Sites, Example \ref{sites-example-constant-sheaf}. Derived
Categories, Lemma \ref{derived-lemma-pre-derived-adjoint-functors}
therefore gives canonical identifications
$$
\begin{aligned}
H^n(\text{Tot}(K^{\bullet, \bullet}))
& = \Hom_{D(R)}(M, R\Gamma(X, \mathcal{A})[n]) \\
& = \Hom_{D(\underline{R})}(\underline{M}, \mathcal{A}[n]) \\
& = \Ext^n_{D(\underline{R})}(\underline{M}, \mathcal{A}).
\end{aligned}
$$

For fixed $p$, freeness of $L_p$ gives
$$
H^q(\Hom_R(L_p, \Gamma(X, \mathcal{C}^\bullet(\mathcal{A}))))
= \Hom_R(L_p, H^q(X, \mathcal{A})).
$$
Taking cohomology in the $L_\bullet$-direction gives
${}'E_2^{p, q} = \Ext_R^p(M, H^q(X, \mathcal{A}))$.

Assume now that every $L_p$ is finite free. Consider the nonnegative
complex of sheaves
$$
\mathbf{H}(\mathcal{F}) =
\SheafHom^\bullet_{\underline{R}}(\underline{L}_\bullet, \mathcal{F}).
$$
This is a resolving functor for
$\SheafHom_{\underline{R}}(\underline{M}, -)$ in the sense of
Derived Categories, Definition
\ref{derived-definition-resolving-functor}. Namely, its terms are exact
functors because the $L_p$ are finite free, and the augmentation induces
an isomorphism in degree zero. If $\mathcal{I}$ is injective, then
$\mathcal{I}|_U$ is injective for every open $U \subset X$ by
Lemma \ref{lemma-cohomology-of-open}, applied to the constant sheaves of
rings. Applying
$\Hom_{\underline{R}|_U}(\underline{L}_\bullet|_U, \mathcal{I}|_U)$
to the free resolution shows that
$H^p(\mathbf{H}(\mathcal{I})) = 0$ for $p > 0$. Thus Derived
Categories, Lemma
\ref{derived-lemma-resolving-functor-computes-derived} gives
$$
H^p(\mathbf{H}(\mathcal{A})) =
\SheafExt^p_{\underline{R}}(\underline{M}, \mathcal{A}).
$$

The Godement construction is exact and additive by
Lemma \ref{lemma-godement-resolution}, and hence it commutes with finite
direct sums. Moreover, every functor
$\mathcal{F} \mapsto \Gamma(X, \mathcal{C}^q(\mathcal{F}))$ is exact:
the Godement construction takes a short exact sequence to a termwise
short exact sequence, and its first term is flasque, so the assertion
follows from Lemma \ref{lemma-flasque-acyclic}. Consequently, finite
freeness gives natural identifications
$$
\begin{aligned}
K^{p, q}
& = \Gamma(X,
\mathcal{C}^q(\SheafHom_{\underline{R}}
(\underline{L}_p, \mathcal{A})))
\end{aligned}
$$
compatible with both differentials. Taking cohomology first in the
$L_\bullet$-direction and then in the Godement direction gives
$$
{}''E_1^{p, q} =
\Gamma(X, \mathcal{C}^p(
\SheafExt^q_{\underline{R}}(\underline{M}, \mathcal{A})))
$$
and hence
$$
{}''E_2^{p, q} =
H^p(X, \SheafExt^q_{\underline{R}}(\underline{M}, \mathcal{A})).
$$

The double complex is first quadrant, so boundedness, convergence, and
the finite filtrations follow from Homology, Lemma
\ref{homology-lemma-first-quadrant-ss}. Functoriality follows from the
functoriality of the Godement resolution and of the two filtrations.
\end{proof}


\begin{lemma}
\label{lemma-constant-sheaf-ext-semisimple-quotient}
\begin{reference}
\cite[Chapter IV, \S 4.3, Corollary 1 to Proposition 4.3.2,
p. 192]{Tohoku}
\end{reference}
Let $X$ be a topological space, let $R$ be a ring, let $I \subset R$ be
an ideal, and put $S = R/I$. Assume that every short exact sequence of
$S$-modules splits. Let $M$ be an $R$-module and let $\mathcal{A}$ be an
$\underline{S}$-module on $X$, viewed as an $\underline{R}$-module. Then,
for every $n \geq 0$, there is a canonical isomorphism
$$
\Ext^n_{D(\underline{R})}(\underline{M}, \mathcal{A}) =
\bigoplus_{p + q = n}
\Hom_S(\text{Tor}_p^R(S, M), H^q(X, \mathcal{A})).
$$
More precisely, both spectral sequences of the double complex in
Lemma \ref{lemma-constant-sheaf-ext-double-complex} degenerate at the
$E_2$-page and give this decomposition of their common abutment.
\end{lemma}

\begin{proof}
Choose a free resolution $L_\bullet \to M$ and let
$\mathcal{A} \to \mathcal{C}^\bullet(\mathcal{A})$ be its Godement
resolution in $\underline{S}$-modules. Its underlying complex is the
Godement resolution after restriction of scalars to $\underline{R}$.
Set
$$
Q^{-p} = S \otimes_R L_p
\quad\text{and}\quad
C^q = \Gamma(X, \mathcal{C}^q(\mathcal{A})).
$$
Since every $C^q$ is an $S$-module, tensor-Hom adjunction gives
isomorphisms
$$
\Hom_R(L_p, C^q) = \Hom_S(S \otimes_R L_p, C^q)
$$
compatible with both differentials. Thus the total complex occurring
in Lemma \ref{lemma-constant-sheaf-ext-double-complex} is
$$
\Hom_S^\bullet(Q^\bullet, C^\bullet).
$$
We have
$$
H^{-p}(Q^\bullet) = \text{Tor}_p^R(S, M)
\quad\text{and}\quad
H^q(C^\bullet) = H^q(X, \mathcal{A}).
$$

We use the splitting hypothesis. If $D^\bullet$ is any complex of
$S$-modules, choose splittings of the inclusions of boundaries into
cycles and of cycles into terms. These give decompositions
$$
D^i = B^i(D) \oplus H^i(D) \oplus B^{i + 1}(D)
$$
under which the differential maps the last summand isomorphically onto
the first summand in degree $i + 1$ and is zero on the other summands.
Consequently, $D^\bullet$ is the direct sum of its cohomology, placed in
the corresponding degrees with zero differential, and a contractible
complex.

Apply this observation first to $Q^\bullet$ and then to $C^\bullet$.
It shows that the two spectral sequences of
$\Hom_S^\bullet(Q^\bullet, C^\bullet)$ degenerate at $E_2$. It also shows
that the canonical map which sends the class of a degree $n$ cocycle to
the maps it induces on cohomology,
$$
H^n(\Hom_S^\bullet(Q^\bullet, C^\bullet))
\longrightarrow
\bigoplus_{p + q = n}
\Hom_S(H^{-p}(Q^\bullet), H^q(C^\bullet)),
$$
is an isomorphism. The map is independent of the chosen splittings;
the splittings are used only to prove that the canonical map is
bijective. For fixed $n$ the displayed sum is finite. Combining this
with the abutment identification of
Lemma \ref{lemma-constant-sheaf-ext-double-complex} proves the formula.
\end{proof}


\begin{lemma}
\label{lemma-constant-sheaf-ext-quotient-spectral-sequence}
\begin{reference}
\cite[Chapter IV, \S 4.3, p. 192]{Tohoku}
\end{reference}
Let $X$ be a topological space, let $R$ be a ring, let $I \subset R$ be
an ideal, and put $S = R/I$. Let $M$ be an $R$-module and let
$\mathcal{A}$ be an $\underline{S}$-module on $X$, viewed as an
$\underline{R}$-module. Choose a free resolution $L_\bullet \to M$, put
$$
Q^{-a} = S \otimes_R L_a,
$$
and let $C^\bullet$ be the complex of global sections of the Godement
resolution of $\mathcal{A}$. For $q \geq 0$, let
$\mathcal{E}_q^\bullet(Q, C)$ be the total complex whose degree $n$ term is
$$
\mathcal{E}_q^n(Q, C) =
\bigoplus_{a + b = n}\Ext_S^q(Q^{-a}, C^b).
$$
There are two relative $S$ hyper-Ext spectral sequences. The first has
$$
{}^{\mathrm I}E_2^{p, q} = H^p(\mathcal{E}_q^\bullet(Q, C)).
$$
It degenerates: its terms with $q > 0$ vanish, while
$\mathcal{E}_0^\bullet(Q, C) = \Hom_S^\bullet(Q^\bullet, C^\bullet)$.
Thus its abutment is canonically
$$
H^*(\Hom_S^\bullet(Q^\bullet, C^\bullet)) =
\Ext^*_{D(\underline{R})}(\underline{M}, \mathcal{A}).
$$
The second is a first quadrant spectral sequence with
$$
{}^{\mathrm{II}}E_2^{p, q} =
\bigoplus_{a + b = q}
\Ext_S^p(\text{Tor}_a^R(S, M), H^b(X, \mathcal{A}))
$$
converging to
$$
\Ext^{p + q}_{D(\underline{R})}(\underline{M}, \mathcal{A}).
$$
It is bounded, induces a finite filtration on each group in the
abutment, and both spectral sequences are additive and functorial in
$\mathcal{A}$ from the $E_2$-page onward.
\end{lemma}

\begin{proof}
Choose a free resolution $L_\bullet \to M$ and let
$\mathcal{A} \to \mathcal{C}^\bullet(\mathcal{A})$ be its Godement
resolution in $\underline{S}$-modules, as in
Lemma \ref{lemma-constant-sheaf-ext-semisimple-quotient}. Put
$$
Q^{-a} = S \otimes_R L_a
\quad\text{and}\quad
C^b = \Gamma(X, \mathcal{C}^b(\mathcal{A})).
$$
Thus $Q^\bullet$ is a bounded above complex of free $S$-modules,
$C^\bullet$ is bounded below, and
$$
H^{-a}(Q^\bullet) = \text{Tor}_a^R(S, M),
\qquad
H^b(C^\bullet) = H^b(X, \mathcal{A}).
$$
Tensor-Hom adjunction identifies
$\Hom_S^\bullet(Q^\bullet, C^\bullet)$ with the total complex occurring
in Lemma \ref{lemma-constant-sheaf-ext-double-complex}.

Choose a Cartan--Eilenberg resolution
$$
C^\bullet \longrightarrow J^{\bullet, \bullet}
$$
in the category of $S$-modules; it exists by Derived Categories,
Lemma \ref{derived-lemma-cartan-eilenberg}. For $p, q \geq 0$ set
$$
D^{p, q} =
\bigoplus_{a + b = q}
\Hom_S(Q^{-a}, J^{b, p}),
$$
with the standard total signs in the $a$- and $b$-directions. This is a
first quadrant double complex.

For every $a$ and $b$, the complex $J^{b, \bullet}$ is an injective
resolution of $C^b$. The spectral sequence of $D$ obtained by first taking
cohomology in this resolution direction has, after interchanging its two
indices,
$$
{}^{\mathrm I}E_1^{p, q} =
\bigoplus_{a + b = p}\Ext_S^q(Q^{-a}, C^b)
$$
and hence the page
$$
{}^{\mathrm I}E_2^{p, q} = H^p(\mathcal{E}_q^\bullet(Q, C)).
$$
Since $Q^{-a}$ is free, $\Ext_S^q(Q^{-a}, C^b)$ vanishes for $q > 0$.
For $q = 0$ the complex $\mathcal{E}_0^\bullet(Q, C)$ is exactly
$\Hom_S^\bullet(Q^\bullet, C^\bullet)$. This proves the assertions about
the first relative hyper-Ext spectral sequence. Equivalently, the augmented
complex
$$
\Hom_S(Q^{-a}, C^b)
\longrightarrow
\Hom_S(Q^{-a}, J^{b, \bullet})
$$
is exact away from degree zero. After taking the finite sums with
$a + b = q$, the variant in the second variable of Homology, Lemma
\ref{homology-lemma-double-complex-gives-resolution} gives a
quasi-isomorphism
$$
\Hom_S^\bullet(Q^\bullet, C^\bullet)
\longrightarrow
\text{Tot}(D^{\bullet, \bullet}).
$$
Consequently, the total cohomology of $D$ is the abutment identified in
Lemma \ref{lemma-constant-sheaf-ext-double-complex}.

The other filtration is the second relative $S$ hyper-Ext spectral
sequence. It first takes cohomology in the $q$-direction. Fix $p$.
The defining properties of a
Cartan--Eilenberg resolution imply that the cycles, boundaries, and
cohomology modules of the complex $J^{\bullet, p}$ are injective.
Hence its differentials split. The canonical map obtained by sending a
cocycle in the Hom complex to the maps it induces on cohomology is
therefore an isomorphism
$$
H^q(D^{p, \bullet})
\longrightarrow
\bigoplus_{a + b = q}
\Hom_S\bigl(
H^{-a}(Q^\bullet), H_J^b(J^{\bullet, p})
\bigr).
$$
Indeed, after choosing the splittings, this follows by decomposing
$J^{\bullet, p}$ into its cohomology and a contractible complex and
using the injectivity of the modules $H_J^b(J^{\bullet, p})$. The map
itself is independent of these splittings.

As $p$ varies, the complex
$$
H_J^b(J^{\bullet, \bullet})
$$
is an injective resolution of $H^b(C^\bullet)$ by the definition of a
Cartan--Eilenberg resolution. Taking cohomology in the $p$-direction
therefore gives
$$
{}^{\mathrm{II}}E_2^{p, q} =
\bigoplus_{a + b = q}
\Ext_S^p(H^{-a}(Q^\bullet), H^b(C^\bullet))
\quad\Longrightarrow\quad
H^{p + q}(\text{Tot}(D^{\bullet, \bullet})).
$$
Substituting the two cohomology identifications above gives the stated
$E_2$-page.

Finally, $D$ is first quadrant. Boundedness, convergence, and the finite
filtrations follow from Homology, Lemma
\ref{homology-lemma-first-quadrant-ss}. The usual comparison maps
between Cartan--Eilenberg resolutions preserve both filtrations and give
additivity and functoriality from the $E_2$-page onward, as in Derived
Categories, Remark
\ref{derived-remark-functorial-ss}.
\end{proof}


\begin{lemma}
\label{lemma-constant-sheaf-ext-hereditary-quotient}
\begin{reference}
\cite[Chapter IV, \S 4.3, p. 192]{Tohoku}
\end{reference}
In the situation of
Lemma \ref{lemma-constant-sheaf-ext-quotient-spectral-sequence}, assume
that $S$ has global dimension at most $1$; this is the hereditary
hypothesis in the cited source. Then, for every $n \geq 0$, there is a
canonical short exact sequence
$$
\begin{aligned}
0 \longrightarrow{}&
\bigoplus_{a + b = n - 1}
\Ext_S^1(\text{Tor}_a^R(S, M), H^b(X, \mathcal{A})) \\
\longrightarrow{}&
\Ext^n_{D(\underline{R})}(\underline{M}, \mathcal{A}) \\
\longrightarrow{}&
\bigoplus_{a + b = n}
\Hom_S(\text{Tor}_a^R(S, M), H^b(X, \mathcal{A}))
\longrightarrow 0.
\end{aligned}
$$
Here the first direct sum is zero when $n = 0$. These sequences are
additive and functorial in $\mathcal{A}$.
\end{lemma}

\begin{proof}
Since $S$ has global dimension at most $1$, Algebra, Lemma
\ref{algebra-lemma-projective-dimension-ext} shows that
$\Ext_S^p(N, N') = 0$ for all $S$-modules $N$, $N'$ and all $p > 1$.
Thus the spectral sequence of
Lemma \ref{lemma-constant-sheaf-ext-quotient-spectral-sequence} has only
the columns $p = 0, 1$. Every differential from the $E_2$-page onward
therefore has zero source or zero target, and $E_2 = E_\infty$.

The finite filtration on the abutment in degree $n$ has two possibly
nonzero graded pieces and gives a natural short exact sequence
$$
0 \longrightarrow E_\infty^{1, n - 1}
\longrightarrow
\Ext^n_{D(\underline{R})}(\underline{M}, \mathcal{A})
\longrightarrow E_\infty^{0, n} \longrightarrow 0.
$$
Substituting the $E_2$-page of the cited spectral sequence gives the
displayed sequence.
\end{proof}

\begin{lemma}
\label{lemma-constant-sheaf-ext-stalk-spectral-sequences}
\begin{reference}
\cite[Chapter IV, \S 4.3, Remark 2, p. 193]{Tohoku}
\end{reference}
Let $X$ be a topological space, let $R$ be a ring, let $M$ be an
$R$-module, let $\mathcal{A}$ be an $\underline{R}$-module on $X$, and let
$x \in X$. Denote by $\mathcal{N}_x$ the filtered set of open
neighbourhoods of $x$, ordered by reverse inclusion. There is a functorial
first quadrant spectral sequence
\begin{equation}
\label{equation-constant-sheaf-ext-stalk-first-spectral-sequence}
{}^{\mathrm{II}}E_2^{p, q} =
\colim_{U \in \mathcal{N}_x}
\Ext_R^p(M, H^q(U, \mathcal{A}))
\quad\Longrightarrow\quad
\SheafExt^{p + q}_{\underline{R}}(\underline{M}, \mathcal{A})_x.
\end{equation}
It induces a finite filtration on the abutment in every total degree.

Let $I \subset R$ be a two-sided ideal, put $S = R/I$, and assume that
$\mathcal{A}$ is an $\underline{S}$-module, viewed as an
$\underline{R}$-module. There is also a functorial first quadrant spectral
sequence
\begin{equation}
\label{equation-constant-sheaf-ext-stalk-quotient-spectral-sequence}
{}^{\mathrm{III}}E_2^{p, q} =
\bigoplus_{a + b = q}
\colim_{U \in \mathcal{N}_x}
\Ext_S^p(\text{Tor}_a^R(S, M), H^b(U, \mathcal{A}))
\quad\Longrightarrow\quad
\SheafExt^{p + q}_{\underline{R}}(\underline{M}, \mathcal{A})_x.
\end{equation}
It too induces a finite filtration in every total degree. Both spectral
sequences are additive and functorial in $\mathcal{A}$ from the $E_2$-page
onward.
\end{lemma}

\begin{proof}
Choose once and for all a free resolution $L_\bullet \to M$ and the
Godement resolution
$\mathcal{A} \to \mathcal{C}^\bullet(\mathcal{A})$ of
Lemma \ref{lemma-godement-resolution}. The construction of the Godement
resolution commutes with restriction to open subspaces: this follows
directly from its construction by the discrete space having the same
points. Thus, for $U \in \mathcal{N}_x$, the complex
$$
C_U^\bullet = \Gamma(U, \mathcal{C}^\bullet(\mathcal{A}))
$$
computes $R\Gamma(U, \mathcal{A})$, and an inclusion $V \subset U$ gives
a restriction map $C_U^\bullet \to C_V^\bullet$. These maps commute
strictly with composition.

Apply Lemma \ref{lemma-constant-sheaf-ext-double-complex} on every $U$,
using the fixed resolution $L_\bullet$. The restriction maps just described
give a directed system, from the $E_2$-page onward, of first quadrant
spectral sequences with
$$
E_2^{p, q}(U) = \Ext_R^p(M, H^q(U, \mathcal{A}))
$$
and abutment
$$
H_U^n =
\Ext^n_{D(\underline{R}|_U)}
(\underline{M}|_U, \mathcal{A}|_U).
$$
For the quotient construction, apply
Lemma \ref{lemma-constant-sheaf-ext-quotient-spectral-sequence} on every
$U$. Its functoriality with respect to the maps
$C_U^\bullet \to C_V^\bullet$ gives a directed system from the $E_2$-page
onward. Its $E_2$-page is
$$
\bigoplus_{a + b = q}
\Ext_S^p(\text{Tor}_a^R(S, M), H^b(U, \mathcal{A})).
$$
The canonical comparison maps between the Cartan--Eilenberg resolutions
used in that lemma are functorial from the $E_2$-page onward, so the maps
for restrictions are compatible with identities and composition.

We explain carefully the passage to the filtered colimit. For either
directed system write $E_r(U)$ for its $r$th page and set
$$
\overline E_r^{p, q} =
\colim_{U \in \mathcal{N}_x} E_r^{p, q}(U),
\qquad
\overline d_r = \colim_{U \in \mathcal{N}_x} d_r(U).
$$
Filtered colimits of abelian groups are exact. Consequently they commute
with kernels, images, and cokernels, and hence
$$
H(\overline E_r, \overline d_r) = \overline E_{r + 1}.
$$
Thus the colimit pages form spectral sequences with the displayed
$E_2$-pages. In the quotient spectral sequence the sum indexed by
$a + b = q$ is finite for fixed $q$, so it commutes with the filtered
colimit and gives exactly the second displayed $E_2$-page.

For completeness, convergence also passes to the colimit. In total degree
$n$, each spectral sequence over $U$ has a finite filtration
$$
H_U^n = F_U^0H_U^n \supset F_U^1H_U^n \supset \ldots
\supset F_U^{n + 1}H_U^n = 0
$$
whose graded pieces are $E_\infty^{p, n - p}(U)$. Restriction maps preserve
these filtrations. Exactness of filtered colimits gives exact sequences
$$
0 \longrightarrow \colim_U F_U^{p + 1}H_U^n
\longrightarrow \colim_U F_U^pH_U^n
\longrightarrow \colim_U E_\infty^{p, n - p}(U)
\longrightarrow 0.
$$
Hence $F^p = \colim_U F_U^pH_U^n$ is a finite filtration of
$\colim_U H_U^n$. For a fixed bidegree $(p, q)$, a first quadrant spectral
sequence stabilizes after a page depending only on $p$ and $q$, and not on
$U$. Consequently
$\colim_U E_\infty^{p, q}(U) = \overline E_\infty^{p, q}$, so the graded
pieces of this filtration are the stable page of the colimit spectral
sequence.

Finally, by Lemma \ref{lemma-section-RHom-over-U}, the sheaf
$\SheafExt^n_{\underline{R}}(\underline{M}, \mathcal{A})$ is the
sheafification of the presheaf $U \mapsto H_U^n$. Sheafification does not
change stalks, and therefore
\begin{equation}
\label{equation-constant-sheaf-ext-stalk-colimit}
\SheafExt^n_{\underline{R}}(\underline{M}, \mathcal{A})_x
= \colim_{U \in \mathcal{N}_x}
\Ext^n_{D(\underline{R}|_U)}
(\underline{M}|_U, \mathcal{A}|_U).
\end{equation}
This identifies the colimit abutments and their finite filtrations with
those asserted in the lemma. Additivity and functoriality pass to filtered
colimits.
\end{proof}


\begin{lemma}
\label{lemma-constant-sheaf-ext-stalk-comparison}
\begin{reference}
\cite[Chapter IV, \S 4.3, Remark 2, p. 193]{Tohoku}
\end{reference}
In the situation of
Lemma \ref{lemma-constant-sheaf-ext-stalk-spectral-sequences}, there are
canonical maps
\begin{equation}
\label{equation-constant-sheaf-ext-stalk-comparison}
\SheafExt^n_{\underline{R}}(\underline{M}, \mathcal{A})_x
\longrightarrow \Ext_R^n(M, \mathcal{A}_x),
\qquad n \geq 0.
\end{equation}
They are isomorphisms if either one of the following sets of hypotheses
holds.
\begin{enumerate}
\item For every $p, q \geq 0$, the canonical map
$$
\colim_{U \in \mathcal{N}_x}
\Ext_R^p(M, H^q(U, \mathcal{A}))
\longrightarrow
\Ext_R^p\left(M,
\colim_{U \in \mathcal{N}_x}H^q(U, \mathcal{A})\right)
$$
is an isomorphism.
\item There is a two-sided ideal $I \subset R$ such that, with $S = R/I$,
the sheaf $\mathcal{A}$ is an $\underline{S}$-module, and for every
$a, b, p \geq 0$ the canonical map
$$
\begin{aligned}
&
\colim_{U \in \mathcal{N}_x}
\Ext_S^p(\text{Tor}_a^R(S, M), H^b(U, \mathcal{A}))
\\
&\qquad\longrightarrow
\Ext_S^p\left(\text{Tor}_a^R(S, M),
\colim_{U \in \mathcal{N}_x}H^b(U, \mathcal{A})\right)
\end{aligned}
$$
is an isomorphism.
\end{enumerate}
\end{lemma}

\begin{proof}
The exact stalk functor sends $\underline{M}|_U$ to $M$ and
$\mathcal{A}|_U$ to $\mathcal{A}_x$. It therefore induces maps
$$
\Ext^n_{D(\underline{R}|_U)}
(\underline{M}|_U, \mathcal{A}|_U)
\longrightarrow \Ext_R^n(M, \mathcal{A}_x)
$$
compatible with restriction in $U$. Their colimit, under the identification
of Equation (\ref{equation-constant-sheaf-ext-stalk-colimit}), is the map
in Equation (\ref{equation-constant-sheaf-ext-stalk-comparison}).

We first record the cohomology groups appearing after passage to the stalk:
$$
\colim_{U \in \mathcal{N}_x} H^q(U, \mathcal{A}) =
\begin{cases}
\mathcal{A}_x & q = 0, \\
0 & q > 0.
\end{cases}
$$
For $q = 0$ this is the definition of the stalk. If $q > 0$, then every
class over a neighbourhood $U$ vanishes after restriction to an open
covering of $U$ by
Lemma \ref{lemma-kill-cohomology-class-on-covering}. Choosing a member of
that covering which contains $x$ shows that the germ of the class is zero.

The fixed Godement resolution gives strict maps
$C_U^\bullet \to \mathcal{C}^\bullet(\mathcal{A})_x$. By
Lemma \ref{lemma-godement-resolution}, the augmentation
$\mathcal{A}_x[0] \to \mathcal{C}^\bullet(\mathcal{A})_x$ is a homotopy
equivalence. Thus the double complexes formed with the stalk complex compute
$\Ext_R^*(M, \mathcal{A}_x)$, and the maps from the neighbourhood complexes
give a comparison from the spectral sequence of Equation
(\ref{equation-constant-sheaf-ext-stalk-first-spectral-sequence}) to this
stalk spectral sequence. Its $E_2$-page is concentrated in the row $q = 0$
and that row is $\Ext_R^p(M, \mathcal{A}_x)$. Under hypothesis (1), the map
on $E_2$-pages is an isomorphism. Therefore it is an isomorphism on every
later page and on the graded pieces of the finite abutment filtrations.
Descending induction on those finite filtrations proves that the map on
abutments is an isomorphism. By construction that map is the canonical map
in Equation (\ref{equation-constant-sheaf-ext-stalk-comparison}).

Under hypothesis (2), use instead the spectral sequence of Equation
(\ref{equation-constant-sheaf-ext-stalk-quotient-spectral-sequence}). The
quotient hyper-Ext construction applied to the stalk complex has
$$
E_2^{p, q} =
\Ext_S^p(\text{Tor}_q^R(S, M), \mathcal{A}_x)
\quad\Longrightarrow\quad
\Ext_R^{p + q}(M, \mathcal{A}_x),
$$
because the stalk complex has cohomology $\mathcal{A}_x$ in degree zero and
zero in higher degrees. The displayed colimit calculation and hypothesis
(2) identify the two $E_2$-pages.
The same finite-filtration argument proves that the canonical map on
abutments is an isomorphism.
\end{proof}


\begin{lemma}
\label{lemma-constant-sheaf-ext-stalk-noetherian-quotient}
\begin{reference}
\cite[Chapter IV, \S 4.3, Remark 2(b), p. 193]{Tohoku}
\end{reference}
Let $X$ be a topological space, let $R$ be a ring, let $I \subset R$ be a
two-sided ideal, and put $S = R/I$. Let $M$ be an $R$-module, let
$\mathcal{A}$ be an $\underline{S}$-module on $X$, and let $x \in X$.
Assume that $S$ is left Noetherian and that
$\text{Tor}_a^R(S, M)$ is a finite $S$-module for every $a \geq 0$. Then
the canonical maps
$$
\SheafExt^n_{\underline{R}}(\underline{M}, \mathcal{A})_x
\longrightarrow \Ext_R^n(M, \mathcal{A}_x)
$$
are isomorphisms for all $n \geq 0$.
\end{lemma}

\begin{proof}
Put $T_a = \text{Tor}_a^R(S, M)$. Since $S$ is left Noetherian and $T_a$
is finite, $T_a$ has a resolution by finite free $S$-modules. Indeed, choose
a surjection from a finite free module onto $T_a$; its kernel is finite by
the left Noetherian hypothesis, and iteration gives the desired resolution.

Let $(N_i)$ be any filtered system of $S$-modules and let $P_\bullet$ be
such a resolution of $T_a$. Since every $P_j$ is finite free, the canonical
maps
$$
\colim_i \Hom_S(P_j, N_i)
\longrightarrow \Hom_S(P_j, \colim_i N_i)
$$
are isomorphisms. Filtered colimits of abelian groups are exact, so taking
cohomology of the Hom complexes gives isomorphisms
$$
\colim_i \Ext_S^p(T_a, N_i)
\longrightarrow \Ext_S^p(T_a, \colim_i N_i)
$$
for every $p \geq 0$. Apply this with
$N_U = H^b(U, \mathcal{A})$, where $\mathcal{N}_x$ is the filtered set of
open neighbourhoods of $x$ ordered by reverse inclusion. Hypothesis (2) of
Lemma \ref{lemma-constant-sheaf-ext-stalk-comparison} holds for every
$a, b, p$, and the result follows from that lemma.
\end{proof}




\section{Cohomology with group actions}
\label{section-cohomology-group-actions}

\begin{definition}
\label{definition-sheaf-group-invariants}
\begin{reference}
\cite[Chapter IV, \S 4.4, pp. 193--194]{Tohoku}
\end{reference}
Let $X$ be a topological space, let $k$ be a ring, and let $G$ be a group.
A \emph{sheaf of $k$-modules with a $G$-action} is an
$\underline{k}$-module $\mathcal{A}$ together with an action of $G$ by
$\underline{k}$-linear automorphisms. We denote the category of these
objects by $\textit{Mod}_{k, G}(X)$. For
$\mathcal{A} \in \textit{Mod}_{k, G}(X)$, its \emph{invariant subsheaf}
$\mathcal{A}^G$ is given by
$$
\mathcal{A}^G(U) = \mathcal{A}(U)^G
$$
for every open $U \subset X$.
For every $x \in X$ the canonical map
$$
(\mathcal{A}^G)_x \longrightarrow (\mathcal{A}_x)^G
$$
is injective, and it is an isomorphism if $G$ is finitely generated.
Indeed, injectivity follows from equality of germs, and for a finite set of
generators a fixed germ has a representative which is invariant after
intersecting finitely many neighbourhoods of $x$.
\end{definition}

\begin{lemma}
\label{lemma-sheaves-group-action-injectives}
The category $\textit{Mod}_{k, G}(X)$ is abelian and has enough injectives.
The forgetful functor to $\textit{Mod}(\underline{k})$ and the functors
$$
(-)^G : \textit{Mod}_{k, G}(X) \longrightarrow
\textit{Mod}(\underline{k})
$$
and
$$
\Gamma(X, -) : \textit{Mod}_{k, G}(X) \longrightarrow
\text{Mod}_{k, G}
$$
send injective objects to injective objects. Here $G$ is given the discrete
topology in the notation on the right.
\end{lemma}

\begin{proof}
Kernels and cokernels in $\textit{Mod}_{k, G}(X)$ are computed on the
underlying $\underline{k}$-modules, so this category is abelian. Let
$$
u : \textit{Mod}_{k, G}(X) \longrightarrow
\textit{Mod}(\underline{k})
$$
be the forgetful functor. It is exact and has a right adjoint
$$
\operatorname{Coind}(\mathcal{F}) = \prod_{g \in G}\mathcal{F}.
$$
Writing a section of this product as $(s_g)_{g \in G}$, the action is
given by
$$
(h \cdot (s_g))_g = s_{gh}.
$$
The adjunction sends a $G$-equivariant map
$\mathcal{A} \to \operatorname{Coind}(\mathcal{F})$ to its component at
the identity. Its inverse sends a map $\varphi : u\mathcal{A} \to
\mathcal{F}$ to the map whose $g$th component is
$a \mapsto \varphi(g \cdot a)$.

Choose an injection $\iota : u\mathcal{A} \to \mathcal{I}$ into an injective
$\underline{k}$-module. The resulting map
$$
\mathcal{A} \longrightarrow \operatorname{Coind}(\mathcal{I}),
\qquad
a \longmapsto (g \mapsto \iota(g \cdot a)),
$$
is injective because its component at the identity is injective.
The object $\operatorname{Coind}(\mathcal{I})$ is injective by Homology,
Lemma \ref{homology-lemma-adjoint-preserve-injectives}. Thus the category has
enough injectives.

The forgetful functor $u$ is right adjoint to the induction functor
$$
\mathcal{F} \longmapsto \bigoplus_{g \in G}\mathcal{F}.
$$
Here $G$ acts by left translation on the summands. This induction functor
is exact, and hence $u$ preserves injectives. The
invariant-subsheaf functor is right adjoint to the exact functor which equips
a $\underline{k}$-module with the trivial $G$-action. It therefore preserves
injectives by the same lemma. Finally, $\Gamma(X, -)$ is right adjoint to the
constant-sheaf functor from $k\text{-}G$-modules to
$\textit{Mod}_{k, G}(X)$. This left adjoint is exact by Sites, Modules,
Lemma \ref{sites-modules-lemma-constant-exact}. Hence global sections also
preserve injectives.
\end{proof}

\begin{definition}
\label{definition-sheaf-group-operator-cohomology}
\begin{reference}
\cite[Chapter IV, \S 4.4, p. 194]{Tohoku}
\end{reference}
Let $X$, $k$, $G$, and $\mathcal{A}$ be as in Definition
\ref{definition-sheaf-group-invariants}. For $q \geq 0$ we set
$$
\mathcal{H}^q(G, \mathcal{A}) = R^q((-)^G)(\mathcal{A}).
$$
These are the \emph{sheaf-valued group cohomology functors}. We also set
$$
H^n(X; G, \mathcal{A}) =
R^n\bigl(\mathcal{A} \longmapsto \Gamma(X, \mathcal{A})^G\bigr)
(\mathcal{A})
$$
for $n \geq 0$ and call these the \emph{cohomology groups of $X$ with
operators $G$ and coefficients in $\mathcal{A}$}. In negative degrees both
functors are zero. Notice that
$$
\Gamma(X, \mathcal{A}^G) = \Gamma(X, \mathcal{A})^G.
$$
\end{definition}

\begin{lemma}
\label{lemma-sheaf-group-operator-spectral-sequences}
\begin{reference}
\cite[Chapter IV, \S 4.4, Theorem 4.4.1, p. 194]{Tohoku}
\end{reference}
\begin{history}
The cited source further says that both spectral sequences below degenerate
when $G$ acts trivially and $k$ is a field. The proof below gives the
canonical direct sum formula through the group-resolution--Godement complex.
We do not assert degeneration of the first spectral sequence without
additional finiteness hypotheses.
\end{history}
Let $X$ be a topological space, let $k$ be a ring, let $G$ be a group, and
let $\mathcal{A}$ be a sheaf of $k$-modules with a $G$-action. There are
functorial first quadrant spectral sequences
$$
{}^{\mathrm I}E_2^{p, q} =
H^p(X, \mathcal{H}^q(G, \mathcal{A}))
$$
and
$$
{}^{\mathrm{II}}E_2^{p, q} =
H^p(G, H^q(X, \mathcal{A}))
$$
converging to $H^{p + q}(X; G, \mathcal{A})$ and inducing finite
filtrations on the abutment. Here group cohomology is as in \'Etale
Cohomology, Definition
\ref{etale-cohomology-definition-galois-cohomology}, with $G$ discrete.
For a free resolution $P_\bullet \to k$ of the trivial $k\text{-}G$-module,
we write
$$
H_p(G, k) = H_p(k \otimes_{k[G]} P_\bullet),
$$
where $k$ is a right $k[G]$-module via the augmentation.

Assume the action of $G$ on $\mathcal{A}$ is trivial. If $k$ is a field,
then for every $n \geq 0$ there are canonical isomorphisms
$$
H^n(X; G, \mathcal{A}) =
\bigoplus_{p + q = n}
\Hom_k(H_p(G, k), H^q(X, \mathcal{A})).
$$
If $k = \mathbf{Z}$, then there are canonical short exact sequences
$$
\begin{aligned}
0 \longrightarrow{}&
\bigoplus_{p + q = n - 1}\quad
\Ext_{\mathbf{Z}}^1(H_p(G, \mathbf{Z}), H^q(X, \mathcal{A})) \\
\longrightarrow{}& H^n(X; G, \mathcal{A}) \\
\longrightarrow{}&
\bigoplus_{p + q = n}\quad
\Hom_{\mathbf{Z}}(H_p(G, \mathbf{Z}), H^q(X, \mathcal{A}))
\longrightarrow 0
\end{aligned}
$$
for $n \geq 0$, where the first direct sum is zero for $n = 0$.
Both the isomorphisms and the short exact sequences are additive and
functorial in $\mathcal{A}$.
\end{lemma}

\begin{proof}
The intermediate categories of sheaves of $k$-modules and of
$k\text{-}G$-modules have enough injectives by Injectives, Lemmas
\ref{injectives-lemma-sheaves-modules-space} and
\ref{injectives-lemma-G-modules}.
Apply Derived Categories, Lemma
\ref{derived-lemma-grothendieck-spectral-sequence} first to the invariant
subsheaf functor followed by global sections. The hypotheses hold by Lemma
\ref{lemma-sheaves-group-action-injectives}, and this gives the first
spectral sequence. Apply the same lemma to global sections with values in
$k\text{-}G$-modules followed by the invariant-module functor. Since the
forgetful functor preserves injectives, the derived functors of equivariant
global sections have underlying modules $H^q(X, \mathcal{A})$. This gives
the second spectral sequence; its group-cohomology notation agrees with
\'Etale Cohomology, Lemma
\ref{etale-cohomology-lemma-modules-abelian}. The two composites agree by the last formula
in Definition \ref{definition-sheaf-group-operator-cohomology}, and their
common derived functor is the one defining $H^*(X; G, -)$. Boundedness,
convergence, finite filtrations, additivity, and functoriality are part of
the cited lemma.

We prove the last two assertions. Choose a resolution
$$
\ldots \longrightarrow P_2 \longrightarrow P_1 \longrightarrow
P_0 \longrightarrow k \longrightarrow 0
$$
of the trivial $k\text{-}G$-module by free $k\text{-}G$-modules, and let
$\mathcal{A} \to \mathcal{C}^\bullet(\mathcal{A})$ be the Godement
resolution. Put
$$
Q_p = k \otimes_{k[G]} P_p,
\qquad
C^q = \Gamma(X, \mathcal{C}^q(\mathcal{A})).
$$
Here $k$ is the right $k[G]$-module given by the augmentation, and we write
$Q^{-p} = Q_p$ when we regard $Q$ as a cochain complex.
The Godement resolution is functorial, so the $G$-action on $C^\bullet$ is
trivial. The total complex of
$$
\Hom_{k, G}(P_p, C^q) = \Hom_k(Q_p, C^q)
$$
computes $H^*(X; G, \mathcal{A})$. Indeed, by the preceding identification
of equivariant derived global sections with ordinary sheaf cohomology, and
because the Godement terms are flasque, $C^\bullet$ represents
$R\Gamma(X, \mathcal{A})$ in the equivariant derived category. The bounded above
complex $P^\bullet$, with $P^{-p} = P_p$, computes derived invariants by
Derived Categories, Lemma
\ref{derived-lemma-morphisms-from-projective-complex}. The comparison with
the derived functor of the composite follows from the second application of
Derived Categories, Lemma
\ref{derived-lemma-grothendieck-spectral-sequence} above.

The homology of $Q_\bullet$ is $H_*(G, k)$, while the cohomology of
$C^\bullet$ is $H^*(X, \mathcal{A})$. If $k$ is a field, every complex of
$k$-vector spaces is the direct sum of its cohomology and a contractible
complex. Applying this first to $Q_\bullet$ and then to $C^\bullet$, as in
the proof of Lemma
\ref{lemma-constant-sheaf-ext-semisimple-quotient}, proves the canonical
direct sum formula.

Let $k = \mathbf{Z}$. Apply the Cartan--Eilenberg construction in the proof
of Lemma \ref{lemma-constant-sheaf-ext-quotient-spectral-sequence} to the
complexes $Q^\bullet$ and $C^\bullet$. It gives a functorial first quadrant
spectral sequence
$$
E_2^{r, s} =
\bigoplus_{a + b = s}
\Ext_{\mathbf{Z}}^r(H_a(G, \mathbf{Z}), H^b(X, \mathcal{A}))
$$
converging to $H^{r + s}(X; G, \mathcal{A})$. Every abelian group has
projective dimension at most $1$: choose a surjection from a free abelian
group and use that its kernel is free. Hence only the columns $r = 0, 1$
can be nonzero. As in the proof of Lemma
\ref{lemma-constant-sheaf-ext-hereditary-quotient}, the resulting two-step
filtration is exactly the displayed short exact sequence.
\end{proof}

\begin{lemma}
\label{lemma-sheaf-group-cohomology-restriction}
\begin{reference}
\cite[Chapter IV, \S 4.4, discussion preceding Theorem 4.4.1,
p. 194]{Tohoku}
\end{reference}
Let $X$ be a topological space, let $k$ be a ring, let $G$ be a group,
and let $\mathcal{A}$ be a sheaf of $k$-modules with a $G$-action. If
$U \subset X$ is open, then there are canonical isomorphisms
$$
\mathcal{H}^q_X(G, \mathcal{A})|_U
\longrightarrow
\mathcal{H}^q_U(G, \mathcal{A}|_U)
$$
for every $q \geq 0$.
For every $x \in X$ there are canonical maps
$$
(\mathcal{H}^q_X(G, \mathcal{A}))_x
\longrightarrow H^q(G, \mathcal{A}_x),
$$
where the group cohomology on the right is computed in
$\text{Mod}_{k, G}$. If $G$ is finite, these maps are isomorphisms.
\end{lemma}

\begin{proof}
Let $j : U \to X$ be the inclusion. Ordinary extension by zero, equipped
with the induced $G$-action, gives a functor
$$
j_! : \textit{Mod}_{k, G}(U) \longrightarrow
\textit{Mod}_{k, G}(X).
$$
It is exact by the stalk description in Sheaves, Lemma
\ref{sheaves-lemma-j-shriek-modules}, and it is left adjoint to
restriction. Hence restriction preserves injectives by Homology, Lemma
\ref{homology-lemma-adjoint-preserve-injectives}. Restricting an injective
resolution of $\mathcal{A}$ therefore gives an injective resolution of
$\mathcal{A}|_U$. Since restriction is exact and commutes with taking the
invariant subsheaf, taking cohomology gives the displayed isomorphisms.

Fix $x \in X$. The functors
$$
F^q(\mathcal{A}) = (\mathcal{H}^q_X(G, \mathcal{A}))_x
$$
form an effaceable cohomological delta functor: taking stalks is exact,
and an injection of $\mathcal{A}$ into an injective object kills every
$F^q$ for $q > 0$. Hence $F^*$ is universal by Homology, Lemma
\ref{homology-lemma-efface-implies-universal}. The canonical map
$(\mathcal{A}^G)_x \to (\mathcal{A}_x)^G$ of Definition
\ref{definition-sheaf-group-invariants} therefore extends uniquely to the
displayed maps.

Assume now that $G$ is finite. The delta functor
$$
E^q(\mathcal{A}) = H^q(G, \mathcal{A}_x)
$$
is effaceable as well. Indeed, let $u\mathcal{A}$ denote the underlying
sheaf and use the canonical injection
$$
\mathcal{A} \longrightarrow \operatorname{Coind}(u\mathcal{A})
= \prod_{g \in G}u\mathcal{A}
$$
from the proof of Lemma
\ref{lemma-sheaves-group-action-injectives}. Since the product is finite,
its stalk is $\operatorname{Coind}(\mathcal{A}_x)$. For every $k$-module
$M$, the $k\text{-}G$-module $\operatorname{Coind}(M)$ is acyclic for
invariants. Namely, apply coinduction to an injective resolution
$M \to J^\bullet$; finite products are exact, coinduction sends injectives
to injectives by Homology, Lemma
\ref{homology-lemma-adjoint-preserve-injectives}, and
$\operatorname{Coind}(J^\bullet)^G = J^\bullet$. Thus
$E^q(\operatorname{Coind}(u\mathcal{A})) = 0$ for $q > 0$.

It follows again from Homology, Lemma
\ref{homology-lemma-efface-implies-universal} that $E^*$ is universal.
For finite $G$ the degree zero map is an isomorphism by Definition
\ref{definition-sheaf-group-invariants}. Its inverse extends to a morphism
of universal delta functors, and uniqueness shows that the two extensions
are inverse in every degree.
\end{proof}

\begin{lemma}
\label{lemma-sheaf-group-cohomology-change-coefficients}
\begin{reference}
\cite[Chapter IV, \S 4.4, Remark 1, pp. 194--195]{Tohoku}
\end{reference}
Let $X$ be a topological space, let $k$ be a ring, let $G$ be a group, and
let $\mathcal{A}$ be a sheaf of $k$-modules with a $G$-action. Write the
subscripts $k$ and $\mathbf{Z}$ to indicate whether the derived functors are
computed in $\textit{Mod}_{k, G}(X)$ or, after forgetting the $k$-module
structure, in $\textit{Mod}_{\mathbf{Z}, G}(X)$. Then there are canonical
functorial isomorphisms
$$
\mathcal{H}^q_k(G, \mathcal{A})
\longrightarrow
\mathcal{H}^q_{\mathbf{Z}}(G, \mathcal{A})
$$
of underlying abelian sheaves and
$$
H^n_k(X; G, \mathcal{A})
\longrightarrow
H^n_{\mathbf{Z}}(X; G, \mathcal{A})
$$
of abelian groups for all $q, n \geq 0$.
\end{lemma}

\begin{proof}
Let $v : \textit{Mod}_{k, G}(X) \to
\textit{Mod}_{\mathbf{Z}, G}(X)$ be restriction of scalars, and let
$\mathcal{I}$ be injective in $\textit{Mod}_{k, G}(X)$. Its underlying
$k$-module is injective by Lemma
\ref{lemma-sheaves-group-action-injectives}, and hence it is flasque.

We first show that $\Gamma(U, \mathcal{I})$ is an injective
$k\text{-}G$-module for every open $U \subset X$. If $j : U \to X$ is the
inclusion, then
$$
M \longmapsto j_!\underline{M}_U
$$
is an exact left adjoint to $\Gamma(U, -)$ on the equivariant categories.
Here exactness follows from exactness of the constant sheaf functor and the
stalk description of $j_!$. Thus the claim follows from Homology, Lemma
\ref{homology-lemma-adjoint-preserve-injectives}.

By \'Etale Cohomology, Lemma
\ref{etale-cohomology-lemma-modules-abelian}, we have
$$
H^p_{\mathbf{Z}}(G, \Gamma(U, \mathcal{I})) = 0
$$
for $p > 0$. Since $v\mathcal{I}$ is flasque, we also have
$H^q(U, v\mathcal{I}) = 0$ for $q > 0$. The second spectral sequence of
Lemma \ref{lemma-sheaf-group-operator-spectral-sequences}, applied over
$\mathbf{Z}$ on $U$, therefore gives
$$
H^m_{\mathbf{Z}}(U; G, v\mathcal{I}) = 0
$$
for every $m > 0$ and every open $U \subset X$.

Apply Lemma \ref{lemma-derived-functor-sheafification} to the invariant
subsheaf functor on $\textit{Mod}_{\mathbf{Z}, G}(X)$. It identifies
$\mathcal{H}^m_{\mathbf{Z}}(G, v\mathcal{I})$ with the sheaf associated to
the presheaf
$$
U \longmapsto H^m_{\mathbf{Z}}(U; G, v\mathcal{I}),
$$
where we use Lemma \ref{lemma-sheaf-group-cohomology-restriction} to
identify the local derived functor. This presheaf is zero for $m > 0$ by
the preceding paragraph. Taking $U = X$ also
shows directly that $v\mathcal{I}$ is acyclic for invariant global sections.

Now apply $v$ to an injective resolution
$\mathcal{A} \to \mathcal{I}^\bullet$ in
$\textit{Mod}_{k, G}(X)$. The resulting resolution is termwise acyclic for
both $\mathbf{Z}$-coefficient functors. Leray's acyclicity lemma, Derived
Categories, Lemma \ref{derived-lemma-leray-acyclicity}, shows that it
computes their right derived functors. Term by term, the invariant subsheaf
and the invariant global sections are unchanged by $v$. The two cohomology
complexes are therefore identical after forgetting the $k$-module
structure, which gives the canonical functorial isomorphisms in the
statement.
\end{proof}


\begin{definition}
\label{definition-orbit-space-equivariant-cohomology}
\begin{reference}
\cite[Chapter V, \S 5.2, pp. 199--200]{Tohoku}
\end{reference}
Let $G$ act on a topological space $X$, let $Y = X/G$, and let
$q : X \to Y$ be the orbit map. For an abelian $G$-sheaf $\mathcal{A}$
on $X$ set
$$
\Gamma_X^G(\mathcal{A}) = \Gamma(X, \mathcal{A})^G.
$$
Thus
$$
\Gamma_X^G = \Gamma(Y, -) \circ q_*^G.
$$
For $n \geq 0$ we define
$$
H^n(X; G, \mathcal{A}) = R^n\Gamma_X^G(\mathcal{A})
$$
and
$$
\mathcal{H}^n(G, \mathcal{A}) = R^nq_*^G(\mathcal{A}).
$$
In negative degrees both functors are zero. The source category is
$\textit{Mod}_G(\underline{\mathbf{Z}}_X)$. It has enough injectives by
Injectives, Lemma
\ref{injectives-lemma-equivariant-modules-enough-injectives}, and $q_*^G$
is additive and left exact by Modules, Lemma
\ref{modules-lemma-invariant-direct-image-G-O-modules}. Consequently these
are universal cohomological delta functors. If $G$ acts trivially on $X$,
the definitions agree with Definition
\ref{definition-sheaf-group-operator-cohomology} for $k = \mathbf{Z}$.
\end{definition}

\begin{lemma}
\label{lemma-localize-higher-invariant-direct-image}
\begin{reference}
\cite[Chapter V, \S 5.2, Proposition 5.2.1, p. 200]{Tohoku}
\end{reference}
In the situation of Definition
\ref{definition-orbit-space-equivariant-cohomology}, let $U \subset Y$ be
open, put $X_U = q^{-1}(U)$, and let $q_U : X_U \to U$ be the restriction
of $q$. Then there are canonical isomorphisms
$$
\mathcal{H}^n(G, \mathcal{A})|_U
=
R^nq_{U,*}^G(\mathcal{A}|_{X_U})
$$
for every $n \geq 0$. Moreover, $\mathcal{H}^n(G, \mathcal{A})$ is the
sheaf associated to the presheaf
$$
U \longmapsto H^n(X_U; G, \mathcal{A}|_{X_U})
$$
on $Y$.
\end{lemma}

\begin{proof}
The open subset $X_U$ is $G$-stable. Equivariant extension by zero gives an
exact functor from abelian $G$-sheaves on $X_U$ to abelian $G$-sheaves on
$X$. It is left adjoint to restriction: this is the ordinary adjunction of
Sheaves, Lemma \ref{sheaves-lemma-j-shriek-modules}, with the action maps
carried along. Hence restriction to $X_U$ preserves injectives by Homology,
Lemma \ref{homology-lemma-adjoint-preserve-injectives}.

Restrict an injective resolution of $\mathcal{A}$. Since invariant direct
image commutes with the displayed restrictions already in degree zero, the
resulting complexes give the first isomorphism. For the second assertion,
apply Lemma \ref{lemma-derived-functor-sheafification} to $q_*^G$ and use
the first assertion to identify the local derived functors.
\end{proof}

\begin{lemma}
\label{lemma-higher-invariant-direct-image-slice}
\begin{reference}
\cite[Chapter V, \S 5.2, Corollary to Proposition 5.2.1,
p. 200]{Tohoku}
\end{reference}
Let $G$ act on $X$, let $H \subset G$ be a subgroup, and let $V \subset X$
be an open $H$-stable subset. Assume that
$$
Z = \bigcup_{g \in G} gV
    = \coprod_{gH \in G/H} gV.
$$
Let $q : X \to X/G$ be the orbit map, put $U = q(V)$, and let
$p : V \to U = V/H$ be the quotient map. If $\mathcal{A}$ is an abelian
$G$-sheaf on $X$, then there are canonical isomorphisms
$$
\mathcal{H}^n(G, \mathcal{A})|_U
\longrightarrow
R^np_*^H(\mathcal{A}|_V)
$$
for every $n \geq 0$.
\end{lemma}

\begin{proof}
The subset $Z=q^{-1}(U)$ is open and $G$-stable. By the preceding lemma we
may replace $X$ by $Z$. The induced bijection $V/H \to U$ is a
homeomorphism: $q|_V$ is open, and two points of $V$ lie in the same
$G$-orbit exactly when they lie in the same $H$-orbit. Restriction to $V$
gives an exact equivalence from
abelian $G$-sheaves on $Z$ to abelian $H$-sheaves on $V$. A quasi-inverse
is obtained by putting the translate of a given sheaf on each open and
closed component $gV$ and using the $G$-action to transport between the
components. Equivalently, on espaces \'etal\'es this is the construction
$$
E \longmapsto G \mathbin{\times^H} E.
$$
The disjointness assumption shows directly that restriction followed by
this construction, in either order, is naturally isomorphic to the identity.
In particular the equivalence preserves injectives.

For every open $W \subset U$, restriction of sections gives a functorial
isomorphism
$$
\mathcal{A}(q^{-1}(W))^G
\longrightarrow
\mathcal{A}(V \cap q^{-1}(W))^H.
$$
Indeed, an invariant section on the left is determined by its restriction
to $V$, and an $H$-invariant section on the right extends uniquely to the
other components by transport. Thus the equivalence identifies $q_*^G$
with $p_*^H$. Applying it to an injective resolution proves the result.
\end{proof}

\begin{lemma}
\label{lemma-higher-invariant-direct-image-stalk}
\begin{reference}
\cite[Chapter V, \S 5.3, Theorem 5.3.1, p. 204]{Tohoku}
\end{reference}
Let $G$ act discontinuously on $X$, let $q : X \to Y = X/G$ be the orbit
map, and let $\mathcal{A}$ be an abelian $G$-sheaf on $X$. Let $y \in Y$
and $x \in q^{-1}(y)$. There are canonical functorial isomorphisms
$$
(R^nq_*^G\mathcal{A})_y
\longrightarrow
H^n(G_x, \mathcal{A}_x)
$$
for every $n \geq 0$. The group cohomology on the right is computed with
$G_x$ viewed as a discrete group.
\end{lemma}

\begin{proof}
Put $H=G_x$. Choose an open $H$-invariant neighbourhood $V$ of $x$ as in
Sheaves, Lemma
\ref{sheaves-lemma-discontinuous-group-action-neighbourhoods}. Then
$$
q^{-1}(q(V)) = \coprod_{gH \in G/H}gV.
$$
By Lemmas \ref{lemma-localize-higher-invariant-direct-image} and
\ref{lemma-higher-invariant-direct-image-slice}, it is enough to prove the
statement for the finite group $H$ acting on $V$, with $x$ fixed by $H$.
Write $p : V \to V/H$ and continue to denote the image of $x$ by $y$.

On the category of abelian $H$-sheaves on $V$, consider the cohomological
delta functors
$$
F^n(\mathcal{B}) = (R^np_*^H\mathcal{B})_y
\quad\text{and}\quad
E^n(\mathcal{B}) = H^n(H, \mathcal{B}_x).
$$
Exactness of the stalk functor makes $E^*$ a cohomological delta functor.
The delta functor $F^*$ is effaceable in positive degrees: embed
$\mathcal{B}$ in an injective object, use the vanishing of positive right
derived functors on injectives, and then take the stalk. Hence it is
universal by Homology, Lemma
\ref{homology-lemma-efface-implies-universal}.

We show that $E^*$ is effaceable as well. Let $v\mathcal{B}$ be the
underlying abelian sheaf. The adjoint of the identity of $v\mathcal{B}$ in
the coinduction adjunction of Injectives, Lemma
\ref{injectives-lemma-equivariant-modules-enough-injectives}, is a
monomorphism
$$
\mathcal{B}
\longrightarrow
\operatorname{Coind}_H(v\mathcal{B}).
$$
Its component in the factor indexed by the identity is the identity. Since
$H$ is finite and fixes $x$, stalks commute with the product defining
coinduction and
$$
(\operatorname{Coind}_H(v\mathcal{B}))_x
\cong \prod_{h \in H}(T_hv\mathcal{B})_x
\cong \prod_{h \in H}\mathcal{B}_x
=
\operatorname{Coind}_H(\mathcal{B}_x).
$$
Every coinduced $H$-module is acyclic for invariants. Indeed, apply
coinduction to an injective resolution of the underlying abelian group;
coinduction is exact, takes injectives to injectives, and its invariant
module is the original module. Thus the displayed monomorphism effaces
$E^n$ for every $n>0$, and $E^*$ is universal by the same Homology lemma.

In degree zero, Sheaves, Lemma
\ref{sheaves-lemma-invariant-direct-image-stalk} gives the canonical
isomorphism
$$
F^0(\mathcal{B})=(p_*^H\mathcal{B})_y
\longrightarrow
\mathcal{B}_x^H=E^0(\mathcal{B}).
$$
Universality extends it uniquely to a morphism $F^*\to E^*$, while the
inverse in degree zero extends uniquely in the other direction. The two
composites have identity degree-zero component, so uniqueness makes them
identities in every degree. The comparison is independent of the chosen
slice: after passage to a smaller common $H$-invariant slice, both
constructions are the unique morphism of universal delta functors extending
the same degree-zero counit map. This proves the canonical functorial
isomorphisms.
\end{proof}


\begin{lemma}[Equivariant quotient spectral sequences]
\label{lemma-orbit-space-equivariant-spectral-sequences}
\begin{reference}
\cite[Chapter V, \S 5.2, Theorem 5.2.1, pp. 200--201]{Tohoku}
\end{reference}
In the situation of Definition
\ref{definition-orbit-space-equivariant-cohomology}, there are additive
functorial first quadrant spectral sequences
\begin{equation}
\label{equation-orbit-space-equivariant-spectral-sequence-I}
{}^{\mathrm I}E_2^{p, q}
=
H^p(Y, \mathcal{H}^q(G, \mathcal{A}))
\quad\Longrightarrow\quad
H^{p + q}(X; G, \mathcal{A})
\end{equation}
and
\begin{equation}
\label{equation-orbit-space-equivariant-spectral-sequence-II}
{}^{\mathrm{II}}E_2^{p, q}
=
H^p(G, H^q(X, \mathcal{A}))
\quad\Longrightarrow\quad
H^{p + q}(X; G, \mathcal{A}).
\end{equation}
Here the ordinary cohomology groups $H^q(X, \mathcal{A})$ carry their
natural $G$-module structures, and the group cohomology in the second
spectral sequence is computed with $G$ discrete. Both spectral sequences
induce finite filtrations on their abutments.
\end{lemma}

\begin{proof}
The defining functor has the two factorizations
$$
\Gamma_X^G
=
\Gamma(Y, -) \circ q_*^G
=
(-)^G \circ \Gamma(X, -).
$$

For the first factorization, the source category has enough injectives by
Injectives, Lemma
\ref{injectives-lemma-equivariant-modules-enough-injectives}, and the
category of abelian sheaves on $Y$ has enough injectives by Injectives,
Lemma \ref{injectives-lemma-abelian-sheaves-space}. If $\mathcal{I}$ is
injective in $\textit{Mod}_G(\underline{\mathbf{Z}}_X)$, then
$q_*^G\mathcal{I}$ is flasque by Injectives, Lemma
\ref{injectives-lemma-equivariant-injective-sections-invariants}, and hence
is right acyclic for global sections by Lemma
\ref{lemma-flasque-acyclic}. Derived Categories, Lemma
\ref{derived-lemma-grothendieck-spectral-sequence} therefore gives
(\ref{equation-orbit-space-equivariant-spectral-sequence-I}).

For the second factorization, the category of $G$-modules has enough
injectives by Injectives, Lemma \ref{injectives-lemma-G-modules}. The same
equivariant-injectives lemma shows that
$\Gamma(X, \mathcal{I})$ is an injective $\mathbf{Z}[G]$-module whenever
$\mathcal{I}$ is an injective abelian $G$-sheaf. It is consequently right
acyclic for invariants, so the Grothendieck spectral sequence applies again.

It remains to identify the derived functors of
$$
\Gamma(X, -) :
\textit{Mod}_G(\underline{\mathbf{Z}}_X)
\longrightarrow \text{Mod}_G.
$$
The forgetful functor from equivariant to ordinary abelian sheaves is right
adjoint to the induction functor of Modules, Lemma
\ref{modules-lemma-induction-G-O-modules}. Induction is exact: it is a direct
sum of exact translates, and exactness may be checked on stalks. Hence the
forgetful functor preserves injectives by Homology, Lemma
\ref{homology-lemma-adjoint-preserve-injectives}. An equivariant injective
resolution is therefore also an ordinary injective resolution after
forgetting the action. Its global-section cohomology is
$H^q(X, \mathcal{A})$, with the action induced by the equivariant resolution.
The derived functors of invariants are the group cohomology functors of
\'Etale Cohomology, Definition
\ref{etale-cohomology-definition-galois-cohomology}. This identifies the
second $E_2$-page. Functoriality, convergence, and the finite filtrations are
part of the cited Grothendieck spectral-sequence lemma.
\end{proof}

\begin{lemma}
\label{lemma-orbit-space-equivariant-spectral-sequence-edges}
\begin{reference}
\cite[Chapter V, \S 5.2, (5.2.6)--(5.2.9), p. 201]{Tohoku}
\end{reference}
The spectral sequences of Lemma
\ref{lemma-orbit-space-equivariant-spectral-sequences} have edge
homomorphisms
\begin{equation}
\label{equation-orbit-space-equivariant-edge-I}
H^n(Y, q_*^G\mathcal{A})
\longrightarrow
H^n(X; G, \mathcal{A})
\longrightarrow
\Gamma(Y, \mathcal{H}^n(G, \mathcal{A}))
\end{equation}
and
\begin{equation}
\label{equation-orbit-space-equivariant-edge-II}
H^n(G, \Gamma(X, \mathcal{A}))
\longrightarrow
H^n(X; G, \mathcal{A})
\longrightarrow
H^n(X, \mathcal{A})^G.
\end{equation}
Under Lemma \ref{lemma-localize-higher-invariant-direct-image}, the second
map in
(\ref{equation-orbit-space-equivariant-edge-I}) sends a class to the section
locally represented, over $U \subset Y$, by its restriction to
$q^{-1}(U)$.

After composing the first map in
(\ref{equation-orbit-space-equivariant-edge-I}) with the second map in
(\ref{equation-orbit-space-equivariant-edge-II}) and then including
invariants into $H^n(X, \mathcal{A})$, one obtains the ordinary pullback map
associated to the counit
$$
q^{-1}q_*^G\mathcal{A} \longrightarrow \mathcal{A}
$$
of Sheaves, Lemma
\ref{sheaves-lemma-inverse-invariant-direct-image}.

In low degrees the edge maps fit into functorial exact sequences
$$
\begin{aligned}
0 \longrightarrow H^1(Y, q_*^G\mathcal{A})
&\longrightarrow H^1(X; G, \mathcal{A})
\longrightarrow \Gamma(Y, \mathcal{H}^1(G, \mathcal{A})) \\
&\longrightarrow H^2(Y, q_*^G\mathcal{A})
\longrightarrow H^2(X; G, \mathcal{A})
\end{aligned}
$$
and
$$
\begin{aligned}
0 \longrightarrow H^1(G, \Gamma(X, \mathcal{A}))
&\longrightarrow H^1(X; G, \mathcal{A})
\longrightarrow H^1(X, \mathcal{A})^G \\
&\longrightarrow H^2(G, \Gamma(X, \mathcal{A}))
\longrightarrow H^2(X; G, \mathcal{A}).
\end{aligned}
$$
\end{lemma}

\begin{proof}
A first quadrant cohomological spectral sequence converging with a finite
filtration has canonical edge maps
$$
E_2^{n, 0} \longrightarrow H^n \longrightarrow E_2^{0, n}
$$
and the low-degree exact sequence
$$
0 \longrightarrow E_2^{1, 0} \longrightarrow H^1
\longrightarrow E_2^{0, 1} \longrightarrow E_2^{2, 0}
\longrightarrow H^2.
$$
Substitution of the two $E_2$-pages proves all the displayed formulas.

Tracing the Grothendieck edge maps through the first factorization identifies
the right edge with restriction to saturated open subsets followed by
sheafification, as described in Lemma
\ref{lemma-localize-higher-invariant-direct-image}. Tracing the first edge
of the first factorization and the right edge of the second factorization
shows that their composite is induced in degree zero by invariant pullback
of sections. Naturality of the Grothendieck construction identifies its
higher-degree maps with the pullback associated to the displayed counit.
\end{proof}

\begin{lemma}[Third equivariant Leray spectral sequence]
\label{lemma-orbit-space-equivariant-third-spectral-sequence}
\begin{reference}
\cite[Chapter V, \S 5.2, (5.2.10), p. 201]{Tohoku}
\end{reference}
In the situation of Definition
\ref{definition-orbit-space-equivariant-cohomology}, regard $Y$ as a
$G$-space with trivial action. There is an additive functorial first quadrant
spectral sequence
$$
E_2^{p,q} = H^p(Y;G,R^qq_*\mathcal{A})
\quad\Longrightarrow\quad
H^{p+q}(X;G,\mathcal{A}).
$$
Here $R^qq_*\mathcal{A}$ is the ordinary higher direct image endowed with its
canonical $G$-action. The spectral sequence induces finite filtrations on the
abutment and its bottom edge is the functorial map
$$
H^n(Y;G,q_*\mathcal{A})
\longrightarrow
H^n(X;G,\mathcal{A}).
$$
\end{lemma}

\begin{proof}
On the equivariant sheaf categories there is a factorization
$$
\Gamma_X^G = \Gamma_Y^G \circ q_*.
$$
Both categories have enough injectives by Injectives, Lemma
\ref{injectives-lemma-equivariant-modules-enough-injectives}. Moreover,
$q_*$ takes injective abelian $G$-sheaves on $X$ to injective abelian
$G$-sheaves on $Y$ by Injectives, Lemma
\ref{injectives-lemma-pushforward-equivariant-injectives}. Thus the
Grothendieck spectral sequence of Derived Categories, Lemma
\ref{derived-lemma-grothendieck-spectral-sequence} applies.

By Injectives, Lemma
\ref{injectives-lemma-equivariant-injective-underlying}, an equivariant
injective resolution is an injective resolution after forgetting the action.
Consequently the inner derived functors are the ordinary higher direct images,
with the action induced by the equivariant resolution. The outer derived
functors are the groups $H^p(Y;G,-)$ of Definition
\ref{definition-sheaf-group-operator-cohomology}. This gives the displayed
$E_2$-page and abutment. The cited spectral-sequence lemma supplies
additivity, functoriality, convergence, and the finite filtrations; its bottom
edge is the displayed map.
\end{proof}

\begin{lemma}
\label{lemma-orbit-space-equivariant-pushforward-acyclic}
\begin{reference}
\cite[Chapter V, \S 5.2, Proposition 5.2.2, pp. 201--202]{Tohoku}
\end{reference}
In the situation of Definition
\ref{definition-orbit-space-equivariant-cohomology}, assume that
$$
R^qq_*\mathcal{A}=0
\quad\text{for all }q>0.
$$
Then the bottom-edge maps of Lemma
\ref{lemma-orbit-space-equivariant-third-spectral-sequence} are canonical
functorial isomorphisms
$$
H^n(Y;G,q_*\mathcal{A})
\longrightarrow
H^n(X;G,\mathcal{A})
$$
for every $n\geq 0$. There are also canonical functorial isomorphisms of
abelian sheaves on $Y$
$$
\mathcal{H}^n(G,\mathcal{A})
\longrightarrow
\mathcal{H}^n(G,q_*\mathcal{A}),
$$
where the expression on the right is sheaf-valued group cohomology on the
trivial $G$-space $Y$. Finally, the two spectral sequences of Lemma
\ref{lemma-orbit-space-equivariant-spectral-sequences} for $\mathcal{A}$ are
canonically identified, from the $E_2$-page onward, with the corresponding two
spectral sequences for $q_*\mathcal{A}$ on $Y$.
\end{lemma}

\begin{proof}
Choose an injective resolution
$\mathcal{A}\to\mathcal{I}^\bullet$ in the category of abelian $G$-sheaves
on $X$. Every $q_*\mathcal{I}^r$ is injective in the equivariant category on
$Y$ by Injectives, Lemma
\ref{injectives-lemma-pushforward-equivariant-injectives}. After forgetting
the action, $\mathcal{I}^\bullet$ is an ordinary injective resolution by
Injectives, Lemma \ref{injectives-lemma-equivariant-injective-underlying}.
The hypothesis therefore says exactly that
$$
q_*\mathcal{A}\longrightarrow q_*\mathcal{I}^\bullet
$$
is an injective resolution in the equivariant category on $Y$.

The third spectral sequence is concentrated in the row $q=0$, and hence its
bottom edges give the first isomorphisms. Term by term we have
$$
q_*^G\mathcal{I}^r=(q_*\mathcal{I}^r)^G.
$$
The complex on the left computes
$\mathcal{H}^*(G,\mathcal{A})$, while the complex on the right computes
$\mathcal{H}^*(G,q_*\mathcal{A})$. This proves the sheaf isomorphisms.

For the first quotient spectral sequence, the complexes obtained by taking
invariant direct images are thus identical. For the second, the identity
$$
\Gamma(X,\mathcal{I}^r)=\Gamma(Y,q_*\mathcal{I}^r)
$$
is compatible with the $G$-actions in every degree. Hence the two pairs of
filtered double complexes agree after using the displayed pushed-forward
resolution. The usual comparison theorem gives the asserted canonical
identifications from the $E_2$-page onward.
\end{proof}

\begin{example}
\label{example-orbit-space-equivariant-five-term-comparison}
\begin{reference}
\cite[Chapter V, \S 5.2, (5.2.12) and the following example,
p. 203]{Tohoku}
\end{reference}
In the situation of Definition
\ref{definition-orbit-space-equivariant-cohomology}, let
$$
u : H^1(Y, q_*^G\mathcal{A}) \longrightarrow H^1(X, \mathcal{A})^G
$$
be the composite of the left edge map of
(\ref{equation-orbit-space-equivariant-edge-I}) with the right edge map of
(\ref{equation-orbit-space-equivariant-edge-II}), and let
$$
u' : H^1(G, \Gamma(X, \mathcal{A}))
\longrightarrow \Gamma(Y, \mathcal{H}^1(G, \mathcal{A}))
$$
be the composite of the left edge map of
(\ref{equation-orbit-space-equivariant-edge-II}) with the right edge map of
(\ref{equation-orbit-space-equivariant-edge-I}). Set
$$
K = \Ker\bigl(H^1(X, \mathcal{A})^G
\longrightarrow H^2(G, \Gamma(X, \mathcal{A}))\bigr)
$$
and
$$
K' = \Ker\bigl(\Gamma(Y, \mathcal{H}^1(G, \mathcal{A}))
\longrightarrow H^2(Y, q_*^G\mathcal{A})\bigr).
$$
There is a canonical isomorphism $\Ker(u)\cong\Ker(u')$ and canonical
exact sequences
$$
\begin{aligned}
0 \longrightarrow \Ker(u) \longrightarrow H^1(Y, q_*^G\mathcal{A})
&\longrightarrow K \\
&\longrightarrow \Coker(u')
\longrightarrow H^2(Y, q_*^G\mathcal{A})
\longrightarrow H^2(X; G, \mathcal{A})
\end{aligned}
$$
and
$$
\begin{aligned}
0 \longrightarrow \Ker(u')
\longrightarrow H^1(G, \Gamma(X, \mathcal{A}))
&\longrightarrow K' \\
&\longrightarrow \Coker(u)
\longrightarrow H^2(G, \Gamma(X, \mathcal{A})) \\
&\longrightarrow H^2(X; G, \mathcal{A}).
\end{aligned}
$$
Thus $\Im(u)=\Ker(K\to\Coker(u'))$ as subobjects of $K$.
\end{example}

\begin{proof}
Apply Homology, Lemma
\ref{homology-lemma-common-middle-five-term-comparison} to the two exact
sequences in Lemma
\ref{lemma-orbit-space-equivariant-spectral-sequence-edges}, taking the
second one as the unprimed sequence and the first one as the primed
sequence. In the notation of that lemma, $u$ and $u'$ are exactly the two
maps displayed above, and its two canonical exact sequences are the ones
displayed here.
\end{proof}

\begin{lemma}
\label{lemma-orbit-space-equivariant-cohomology-degeneration}
\begin{reference}
\cite[Chapter V, \S 5.2, Proposition 5.2.3, p. 202]{Tohoku}
\end{reference}
In the situation of Definition
\ref{definition-orbit-space-equivariant-cohomology}, assume that
$$
\mathcal{H}^q(G, \mathcal{A}) = 0
\quad\text{for all }q > 0.
$$
Then the edge maps give canonical functorial isomorphisms
$$
H^n(Y, q_*^G\mathcal{A})
\longrightarrow
H^n(X; G, \mathcal{A})
$$
for every $n \geq 0$. Consequently the second spectral sequence of Lemma
\ref{lemma-orbit-space-equivariant-spectral-sequences} may be written
$$
H^p(G, H^q(X, \mathcal{A}))
\quad\Longrightarrow\quad
H^{p + q}(Y, q_*^G\mathcal{A}).
$$
\end{lemma}

\begin{proof}
Under the hypothesis, the first spectral sequence of Lemma
\ref{lemma-orbit-space-equivariant-spectral-sequences} is concentrated in
the row $q=0$. It therefore degenerates and its edge maps give the displayed
isomorphisms. Substituting this identification of the abutment in the second
spectral sequence gives the final assertion.
\end{proof}

\begin{lemma}
\label{lemma-orbit-space-equivariant-cohomology-global-acyclic}
\begin{reference}
\cite[Chapter V, \S 5.2, Proposition 5.2.4, p. 202]{Tohoku}
\end{reference}
In the situation of Definition
\ref{definition-orbit-space-equivariant-cohomology}, assume that
$$
H^q(X, \mathcal{A})=0
\quad\text{for all }q>0.
$$
Then the edge maps give canonical functorial isomorphisms
$$
H^n(G, \Gamma(X, \mathcal{A}))
\longrightarrow
H^n(X; G, \mathcal{A})
$$
for every $n \geq 0$. Consequently the first spectral sequence of Lemma
\ref{lemma-orbit-space-equivariant-spectral-sequences} may be written
$$
H^p(Y, \mathcal{H}^q(G, \mathcal{A}))
\quad\Longrightarrow\quad
H^{p + q}(G, \Gamma(X, \mathcal{A})).
$$
\end{lemma}

\begin{proof}
Under the hypothesis, the second spectral sequence of Lemma
\ref{lemma-orbit-space-equivariant-spectral-sequences} is concentrated in
the row $q=0$. It therefore degenerates and its edge maps give the displayed
isomorphisms. Substituting this identification of the abutment in the first
spectral sequence gives the final assertion.
\end{proof}

\begin{lemma}
\label{lemma-finite-group-cohomology-invertible-order}
Let $G$ be a finite group of order $m$ and let $M$ be a $G$-module.
\begin{enumerate}
\item The group $H^p(G, M)$ is annihilated by $m$ for every $p>0$.
\item If multiplication by $m$ is an automorphism of $M$, then
$H^p(G, M)=0$ for every $p>0$.
\end{enumerate}
\end{lemma}

\begin{proof}
Let
$$
C = \operatorname{Coind}_{\{1\}}^G(M) = \prod_{g \in G} M,
$$
where in the coinduction term $M$ denotes the underlying abelian group.
Writing an element of $C$ as $(c_g)_{g \in G}$, the action is
$(h c)_g = c_{gh}$. There is a natural monomorphism
$$
M \longrightarrow C, \qquad x \longmapsto (g x)_{g \in G}.
$$
Put $N=C/M$. The coinduced module $C$ is acyclic for invariants. Indeed,
coinduction is exact and is right adjoint to the exact forgetful functor.
Hence it carries injective abelian groups to injective $G$-modules by
Homology, Lemma
\ref{homology-lemma-adjoint-preserve-injectives}, and taking invariants of a
coinduced module recovers the original abelian group.

If $\overline{c} \in N^G$ and $c \in C$ lifts it, then
$$
c' = \sum_{g \in G}gc
$$
is invariant and maps to $m\overline{c}$. Thus the cokernel of
$C^G\to N^G$ is annihilated by $m$. The long exact sequence in group
cohomology shows that $H^1(G,M)$ is this cokernel and gives
$$
H^p(G,M) \cong H^{p-1}(G,N) \quad\text{for }p>1.
$$
Induction on $p$ proves (1). If multiplication by $m$ is an automorphism of
$M$, then it induces an automorphism of $H^p(G,M)$. This automorphism is zero
for $p>0$ by (1), and hence $H^p(G,M)=0$. This proves (2).
\end{proof}

\begin{lemma}
\label{lemma-orbit-space-equivariant-cohomology-stabilizer-orders}
\begin{reference}
\cite[Chapter V, \S 5.3, Corollaries 1 and 3, pp. 204--205]{Tohoku}
\end{reference}
Let $G$ act discontinuously on $X$, let $q : X \to Y=X/G$ be the orbit
map, and let $\mathcal{A}$ be an abelian $G$-sheaf on $X$. Assume that for
every $x \in X$ multiplication by $|G_x|$ is an automorphism of
$\mathcal{A}_x$. Then there are canonical functorial isomorphisms
$$
H^n(Y, q_*^G\mathcal{A})
\longrightarrow
H^n(X; G, \mathcal{A})
$$
for every $n \geq 0$,
and an additive functorial first quadrant spectral sequence
$$
H^p(G, H^q(X, \mathcal{A}))
\quad\Longrightarrow\quad
H^{p + q}(Y, q_*^G\mathcal{A}).
$$
In particular, if the action is free, these conclusions hold for every
abelian $G$-sheaf on $X$.
\end{lemma}

\begin{proof}
Apply Lemma \ref{lemma-finite-group-cohomology-invertible-order}, part (2),
to $H=G_x$ and $M=\mathcal{A}_x$. Lemma
\ref{lemma-higher-invariant-direct-image-stalk} gives
$$
\mathcal{H}^q(G, \mathcal{A})_y
=
H^q(G_x, \mathcal{A}_x)
=0
$$
for $q>0$, where $y=q(x)$. Hence
$\mathcal{H}^q(G, \mathcal{A})=0$ for $q>0$, and Lemma
\ref{lemma-orbit-space-equivariant-cohomology-degeneration} proves the two
conclusions. If the action is free, then $|G_x|=1$ for every $x$, so the
hypothesis is automatic.
\end{proof}

\begin{lemma}
\label{lemma-orbit-space-equivariant-cohomology-finite-group}
\begin{reference}
\cite[Chapter V, \S 5.2, Corollary to Proposition 5.2.3,
p. 202]{Tohoku}
\end{reference}
Let $G$ be a finite group of order $m$ acting on a Hausdorff space $X$,
let $q : X \to Y=X/G$ be the orbit map, and let $\mathcal{A}$ be an
abelian $G$-sheaf on $X$. Assume that multiplication by $m$ is an
automorphism of $\mathcal{A}$. Then the edge maps give canonical functorial
isomorphisms
$$
H^n(Y, q_*^G\mathcal{A})
\longrightarrow
H^n(X; G, \mathcal{A})
\longrightarrow
H^n(X, \mathcal{A})^G
$$
for every $n \geq 0$. The hypothesis on multiplication holds, for example,
if $\mathcal{A}$ is a sheaf of vector spaces over a field whose
characteristic does not divide $m$.
\end{lemma}

\begin{proof}
By Sheaves, Lemma
\ref{sheaves-lemma-discontinuous-group-action-neighbourhoods}, the action
satisfies condition (D). For $x \in X$, put $d=|G_x|$ and $e=m/d$.
Multiplication by $d$ on $\mathcal{A}_x$ is an automorphism: its inverse is
the composite of multiplication by $e$ with the inverse of multiplication
by $m$. Hence Lemma
\ref{lemma-orbit-space-equivariant-cohomology-stabilizer-orders} gives the
first displayed isomorphism.

For every $q \geq 0$, multiplication by $m$ on
$H^q(X, \mathcal{A})$ is an automorphism by functoriality. Lemma
\ref{lemma-finite-group-cohomology-invertible-order}, part (2), therefore
gives
$$
H^p(G, H^q(X, \mathcal{A})) = 0 \quad\text{for }p>0.
$$
The second spectral sequence of Lemma
\ref{lemma-orbit-space-equivariant-spectral-sequences} is consequently
concentrated in the column $p=0$. It degenerates, and Lemma
\ref{lemma-orbit-space-equivariant-spectral-sequence-edges} identifies its
right edge with the second displayed isomorphism.
\end{proof}

\begin{lemma}
\label{lemma-prime-support-torsion-serre-subcategory}
Let $P$ be a set of prime numbers. The full subcategory
$\mathcal{C}_P\subset\textit{Ab}$ consisting of the abelian groups in which
every element is annihilated by a positive integer all of whose prime
divisors belong to $P$ is a Serre subcategory.
\end{lemma}

\begin{proof}
The zero group belongs to $\mathcal{C}_P$, and the condition is preserved by
subgroups and quotients. Suppose that
$$
0\longrightarrow A\longrightarrow B\longrightarrow C\longrightarrow0
$$
is exact with $A,C\in\mathcal{C}_P$. Given $b\in B$, choose a positive
integer $r$, with all prime divisors in $P$, which annihilates the image of
$b$ in $C$. Then $rb$ comes from $A$, so there is another such integer $s$
with $srb=0$. Thus $B\in\mathcal{C}_P$, and the subcategory is closed under
extensions. This is the Serre condition of Homology, Definition
\ref{homology-definition-serre-subcategory}.
\end{proof}

\begin{lemma}
\label{lemma-orbit-space-equivariant-cohomology-modulo-stabilizer-primes}
\begin{reference}
\cite[Chapter V, \S 5.3, Corollary 2, pp. 204--205]{Tohoku}
\end{reference}
\begin{history}
In the proof, the cited source writes
$H^p(X,\mathcal{H}^q(G,\mathcal{A}))$. The first spectral sequence of
Theorem 5.2.1 lives on the orbit space $Y$, so the correct group is the one
with $Y$ used below.
\end{history}
Let $G$ act discontinuously on $X$, let $q:X\to Y=X/G$ be the orbit map,
and let $\mathcal{A}$ be an abelian $G$-sheaf on $X$. Assume that the orders
$|G_x|$ have a least common multiple $n$. Let $P$ be the set of prime
divisors of $n$, let $\mathcal{C}_P$ be the Serre subcategory of Lemma
\ref{lemma-prime-support-torsion-serre-subcategory}, and write
$$
\pi:\textit{Ab}\longrightarrow\textit{Ab}/\mathcal{C}_P
$$
for the exact quotient functor of Homology, Lemma
\ref{homology-lemma-serre-subcategory-is-kernel}. Then the edge maps induce
canonical functorial isomorphisms
$$
\pi\bigl(H^r(Y,q_*^G\mathcal{A})\bigr)
\longrightarrow
\pi\bigl(H^r(X;G,\mathcal{A})\bigr)
$$
for every $r\geq0$. Moreover, applying $\pi$ to the second equivariant
quotient spectral sequence gives an additive functorial first quadrant
spectral sequence in $\textit{Ab}/\mathcal{C}_P$
$$
\pi\bigl(H^p(G,H^q(X,\mathcal{A}))\bigr)
\quad\Longrightarrow\quad
\pi\bigl(H^{p+q}(Y,q_*^G\mathcal{A})\bigr)
$$
with finite filtrations on its abutments.
\end{lemma}

\begin{proof}
Let $y\in Y$ and choose $x\in X$ above $y$. For $q>0$, Lemma
\ref{lemma-finite-group-cohomology-invertible-order}, part (1), and Lemma
\ref{lemma-higher-invariant-direct-image-stalk} give
$$
|G_x|\,\mathcal{H}^q(G,\mathcal{A})_y
=|G_x|\,H^q(G_x,\mathcal{A}_x)=0.
$$
Since $|G_x|$ divides $n$, multiplication by $n$ is zero on every stalk and
hence on $\mathcal{H}^q(G,\mathcal{A})$. It follows by functoriality that
every group
$$
{}^{\mathrm I}E_2^{p,q}
=H^p(Y,\mathcal{H}^q(G,\mathcal{A})),\qquad q>0,
$$
is annihilated by $n$ and therefore belongs to $\mathcal{C}_P$.
Since $\mathcal{C}_P$ is a Serre subcategory, every
${}^{\mathrm I}E_s^{p,q}$ with $q>0$ belongs to $\mathcal{C}_P$ for
$s\geq2$.

Consider the left edge map of the first spectral sequence of Lemma
\ref{lemma-orbit-space-equivariant-spectral-sequences}. Its kernel has a
finite filtration by subquotients of terms with positive second index: these
are the images of the differentials
$E_s^{r-s,s-1}\to E_s^{r,0}$ for $s\geq2$. Its cokernel has the finite
abutment filtration with graded pieces $E_\infty^{p,q}$ for $q>0$. Hence
both kernel and cokernel belong to $\mathcal{C}_P$. The exact functor $\pi$
therefore turns the edge map into the displayed isomorphism.

Finally, apply $\pi$ to every page and differential of the second spectral
sequence of Lemma
\ref{lemma-orbit-space-equivariant-spectral-sequences}. Exactness of $\pi$
shows that it commutes with passage from one page to the next, and it carries
the finite abutment filtration to a finite filtration with the corresponding
graded pieces. This gives a spectral sequence with abutment
$\pi(H^{p+q}(X;G,\mathcal{A}))$. The preceding edge isomorphism identifies
that abutment canonically with
$\pi(H^{p+q}(Y,q_*^G\mathcal{A}))$, as asserted.
\end{proof}


\begin{lemma}
\label{lemma-equivariant-support-excision}
Let $G$ act on a topological space $X$, let $Z \subset X$ be a closed
$G$-stable subset, and put $W = X \setminus Z$. Let $i : Z \to X$ and
$j : W \to X$ be the inclusions. If $\mathcal{F}$ is an abelian $G$-sheaf
on $X$, then
$$
\mathcal{F}_W = j_!(\mathcal{F}|_W)
\quad\text{and}\quad
\mathcal{F}_Z = i_*(\mathcal{F}|_Z)
$$
have natural structures of abelian $G$-sheaves, and there is a functorial
short exact sequence
$$
0 \longrightarrow \mathcal{F}_W \longrightarrow \mathcal{F}
\longrightarrow \mathcal{F}_Z \longrightarrow 0.
$$
For every $n \in \mathbf{Z}$ there is a canonical functorial isomorphism
$$
H^n(X;G,\mathcal{F}_Z)
\longrightarrow
H^n(Z;G,\mathcal{F}|_Z).
$$
Here the action of $G$ on $Z$ is the restricted action.
\end{lemma}

\begin{proof}
The functoriality of extension by zero with respect to the translations by
$G$ gives the asserted actions. The short exact sequence is the sequence of
Lemma \ref{lemma-extension-by-zero-locally-closed}, with its equivariant
structure retained.

Equivariant direct image along $i$ is exact, because this can be checked on
the underlying sheaves and Modules, Lemma
\ref{modules-lemma-i-star-exact} applies. It is right adjoint to the exact
equivariant restriction functor. Hence it preserves injectives by Homology,
Lemma \ref{homology-lemma-adjoint-preserve-injectives}. Finally,
$$
\Gamma_X^G(i_*\mathcal{G})=\Gamma_Z^G(\mathcal{G})
$$
for every abelian $G$-sheaf $\mathcal{G}$ on $Z$. Apply this equality to an
injective resolution on $Z$ to obtain the displayed isomorphisms.
\end{proof}

\begin{lemma}[Equivariant support filtration]
\label{lemma-orbit-space-equivariant-cohomology-support-filtration}
\begin{reference}
\cite[Chapter V, \S 5.4, (5.4.1), (5.4.2), and Proposition 5.4.1,
pp. 205--206]{Tohoku}
\end{reference}
Let $G$ act discontinuously on $X$, let $q : X \to Y=X/G$ be the orbit
map, and let $\mathcal{A}$ be an abelian $G$-sheaf on $X$. Let
$Y_0 \subset Y$ be closed and assume that
$$
\mathcal{H}^r(G,\mathcal{A})|_{Y \setminus Y_0}=0
\quad\text{for every }r>0.
$$
Put
$$
X_0=q^{-1}(Y_0),\qquad V=Y\setminus Y_0,\qquad
U=X\setminus X_0=q^{-1}(V),\qquad
\mathcal{B}=q_*^G\mathcal{A}.
$$
Use subscripts to denote extension by zero from these locally closed
subsets. There are canonical short exact sequences
\begin{equation}
\label{equation-orbit-space-equivariant-support-short-exact}
\begin{gathered}
0\longrightarrow\mathcal{A}_U\longrightarrow\mathcal{A}
\longrightarrow\mathcal{A}_{X_0}\longrightarrow0,\\
0\longrightarrow\mathcal{B}_V\longrightarrow\mathcal{B}
\longrightarrow\mathcal{B}_{Y_0}\longrightarrow0,
\end{gathered}
\end{equation}
and canonical isomorphisms
$$
q_*^G\mathcal{A}_U\cong\mathcal{B}_V,
\qquad
q_*^G\mathcal{A}_{X_0}\cong\mathcal{B}_{Y_0}.
$$

For $r\in\mathbf{Z}$ set
\begin{align*}
L^r&=H^r(Y,\mathcal{B}_V), &
M^r&=H^r(Y,\mathcal{B}), &
N^r&=H^r(Y_0,\mathcal{B}|_{Y_0}),\\
\widetilde L^r&=H^r(X;G,\mathcal{A}_U), &
\widetilde M^r&=H^r(X;G,\mathcal{A}), &
\widetilde N^r&=H^r(X_0;G,\mathcal{A}|_{X_0}).
\end{align*}
The groups in negative degrees are zero. The long exact sequences associated
to (\ref{equation-orbit-space-equivariant-support-short-exact}) fit into a
canonical commutative diagram with exact rows
\begin{equation}
\label{equation-orbit-space-equivariant-support-diagram}
\xymatrix@C=1.4em{
\cdots \ar[r] &
L^r \ar[r]\ar[d]^{\lambda^r} &
M^r \ar[r]\ar[d]^{\alpha^r} &
N^r \ar[r]\ar[d]^{\rho^r} &
L^{r+1} \ar[r]\ar[d]^{\lambda^{r+1}} & \cdots\\
\cdots \ar[r] &
\widetilde L^r \ar[r] &
\widetilde M^r \ar[r] &
\widetilde N^r \ar[r] &
\widetilde L^{r+1} \ar[r] & \cdots.
}
\end{equation}
Every $\lambda^r$ is an isomorphism.

Set
$$
Q^r=\Coker(\rho^r)
=\widetilde N^r/\Im(\rho^r),
$$
and let $\beta^r:\widetilde M^r\to Q^r$ be induced by the lower row of
(\ref{equation-orbit-space-equivariant-support-diagram}). There are
canonical homomorphisms $\partial^r:Q^{r-1}\to M^r$ which give a complex
\begin{equation}
\label{equation-orbit-space-equivariant-support-complex}
\cdots\longrightarrow
\widetilde M^{n-1}\xrightarrow{\beta^{n-1}}Q^{n-1}
\xrightarrow{\partial^n}M^n
\xrightarrow{\alpha^n}\widetilde M^n
\xrightarrow{\beta^n}Q^n\longrightarrow\cdots.
\end{equation}
This complex is exact except possibly at $M^n$, and there is a canonical
isomorphism
$$
\Ker(\alpha^n)/\Im(\partial^n)
\longrightarrow
\Ker(\rho^n).
$$

Consequently, for every $n\geq0$ there is a three-step filtration
$$
0\subset\Im(\partial^n)\subset\Ker(\alpha^n)\subset M^n
$$
whose successive quotients are canonically isomorphic to
$$
\Coker(\beta^{n-1}),\qquad
\Ker(\rho^n),\qquad
\Ker(\beta^n),
$$
respectively. The first one has the explicit description
\begin{align*}
\Coker(\beta^{n-1})
={}&H^{n-1}(X_0;G,\mathcal{A}|_{X_0})\Big/
\Big(
\Im\bigl(H^{n-1}(X;G,\mathcal{A})\to
H^{n-1}(X_0;G,\mathcal{A}|_{X_0})\bigr)\\
&\qquad+\Im\bigl(H^{n-1}(Y_0,\mathcal{B}|_{Y_0})\to
H^{n-1}(X_0;G,\mathcal{A}|_{X_0})\bigr)\Big).
\end{align*}
\end{lemma}

\begin{proof}
The subset $X_0$ is closed and $G$-stable, and $U$ is open and $G$-stable.
The first sequence in
(\ref{equation-orbit-space-equivariant-support-short-exact}) follows from
Lemma \ref{lemma-equivariant-support-excision}. The second follows from
Lemma \ref{lemma-extension-by-zero-locally-closed} on $Y$.

Apply the left exact functor $q_*^G$ to the first sequence. The resulting
monomorphism
$$
q_*^G\mathcal{A}_U\longrightarrow\mathcal{B}
$$
identifies its source with the subobject $\mathcal{B}_V$. Indeed, if
$y=q(x)$, then Sheaves, Lemma
\ref{sheaves-lemma-invariant-direct-image-stalk} identifies its map on
stalks with
$$
(\mathcal{A}_U)_x^{G_x}\longrightarrow\mathcal{A}_x^{G_x}.
$$
This is an isomorphism for $y\in V$ and has zero source for $y\in Y_0$,
which is exactly the stalk description of
$\mathcal{B}_V\to\mathcal{B}$. The induced monomorphism
$$
\mathcal{B}_{Y_0}=\mathcal{B}/\mathcal{B}_V
\longrightarrow q_*^G\mathcal{A}_{X_0}
$$
is also an isomorphism on stalks: over $Y_0$ its map is the identity on
$\mathcal{A}_x^{G_x}$, and over $V$ both sides vanish. This proves the two
isomorphisms of invariant direct images.

The left edge maps of Lemma
\ref{lemma-orbit-space-equivariant-spectral-sequence-edges}, applied to the
three terms of the first short exact sequence, give the vertical arrows in
(\ref{equation-orbit-space-equivariant-support-diagram}). They commute with
the connecting maps: in the derived category the left edge is induced by
$$
q_*^G\mathcal{F}[0]\longrightarrow Rq_*^G\mathcal{F}
$$
after applying $R\Gamma(Y,-)$, and hence is a morphism of cohomological
delta functors. Lemma \ref{lemma-cohomology-and-closed-immersions}
identifies the upper closed term with $N^r$, while Lemma
\ref{lemma-equivariant-support-excision} identifies the lower closed term
with $\widetilde N^r$. This proves the diagram.

We next prove that the maps $\lambda^r$ are isomorphisms. For $s>0$,
Lemma \ref{lemma-localize-higher-invariant-direct-image} and the support
hypothesis give
$$
\mathcal{H}^s(G,\mathcal{A}_U)|_V
=\mathcal{H}^s(G,\mathcal{A})|_V=0.
$$
If $y\in Y_0$ and $x\in q^{-1}(y)$, then
$(\mathcal{A}_U)_x=0$, and Lemma
\ref{lemma-higher-invariant-direct-image-stalk} gives
$$
\mathcal{H}^s(G,\mathcal{A}_U)_y
=H^s(G_x,(\mathcal{A}_U)_x)=0.
$$
Thus $\mathcal{H}^s(G,\mathcal{A}_U)=0$ for all $s>0$. Lemma
\ref{lemma-orbit-space-equivariant-cohomology-degeneration}, together with
$q_*^G\mathcal{A}_U=\mathcal{B}_V$, now proves that every $\lambda^r$ is an
isomorphism.

It remains to construct and analyse
(\ref{equation-orbit-space-equivariant-support-complex}). Write the maps in
the upper exact row of
(\ref{equation-orbit-space-equivariant-support-diagram}) as
$$
L^r\xrightarrow{a^r}M^r\xrightarrow{b^r}N^r
\xrightarrow{c^r}L^{r+1},
$$
and write $\widetilde a^r,\widetilde b^r,\widetilde c^r$ for the
corresponding maps in the lower row. The composite
$$
\widetilde N^{r-1}\xrightarrow{\widetilde c^{r-1}}
\widetilde L^r\xrightarrow{(\lambda^r)^{-1}}L^r
\xrightarrow{a^r}M^r
$$
kills $\Im(\rho^{r-1})$ by commutativity and exactness of the upper row.
It therefore induces $\partial^r:Q^{r-1}\to M^r$. Commutativity and
exactness immediately show that each product of two consecutive arrows in
(\ref{equation-orbit-space-equivariant-support-complex}) is zero.

We give the remaining diagram chase. If
$x\in\widetilde M^r$ has $\beta^r(x)=0$, then
$\widetilde b^r(x)=\rho^r(z)$ for some $z\in N^r$. We have
$$
\lambda^{r+1}c^r(z)
=\widetilde c^r\rho^r(z)
=\widetilde c^r\widetilde b^r(x)=0.
$$
Thus $c^r(z)=0$, so $z=b^r(y)$ for some $y\in M^r$. Exactness of the
lower row gives
$$
x-\alpha^r(y)=\widetilde a^r(t)
$$
for some $t\in\widetilde L^r$. Write $t=\lambda^r(s)$. Commutativity now
gives $x=\alpha^r(y+a^r(s))$. Hence the complex is exact at
$\widetilde M^r$.

Suppose that $z\in\widetilde N^r$ represents a class in
$\Ker(\partial^{r+1})$. Then
$$
a^{r+1}(\lambda^{r+1})^{-1}\widetilde c^r(z)=0.
$$
There is therefore a $y\in N^r$ with
$c^r(y)=(\lambda^{r+1})^{-1}\widetilde c^r(z)$. It follows that
$\widetilde c^r(z-\rho^r(y))=0$, so
$z-\rho^r(y)=\widetilde b^r(x)$ for some $x\in\widetilde M^r$. Thus the
class of $z$ is $\beta^r(x)$, proving exactness at $Q^r$.

Finally, $b^n$ sends $\Ker(\alpha^n)$ into $\Ker(\rho^n)$. This map is
surjective. Indeed, if $y\in\Ker(\rho^n)$, then
$\lambda^{n+1}c^n(y)=0$, so $y=b^n(x)$ for some $x\in M^n$. Since
$\widetilde b^n\alpha^n(x)=0$, we may write
$\alpha^n(x)=\widetilde a^n(t)$. Write $t=\lambda^n(s)$ and replace $x$ by
$x-a^n(s)$; it still maps to $y$ and now belongs to
$\Ker(\alpha^n)$. The kernel of
$$
b^n:\Ker(\alpha^n)\longrightarrow\Ker(\rho^n)
$$
is $\Im(\partial^n)$. To see this, if $x=a^n(s)$ belongs to
$\Ker(\alpha^n)$, then $\lambda^n(s)$ comes from some
$z\in\widetilde N^{n-1}$, and the class of $z$ maps to $x$ under
$\partial^n$. The converse follows directly from the definition of
$\partial^n$. This proves the assertion about the defect and completes the
exactness analysis.

Exactness at $Q^{n-1}$ identifies $\Im(\partial^n)$ with
$\Coker(\beta^{n-1})$. The defect isomorphism gives the middle quotient,
and exactness at $\widetilde M^n$ identifies
$M^n/\Ker(\alpha^n)$ with $\Ker(\beta^n)$. Finally, expanding the
definitions of $Q^{n-1}$ and $\beta^{n-1}$ gives the displayed formula for
$\Coker(\beta^{n-1})$.
\end{proof}


\begin{lemma}
\label{lemma-orbit-space-equivariant-cohomology-trivial-fixed-part}
\begin{reference}
\cite[Chapter V, \S 5.4, Corollary 5.4.1, p. 206]{Tohoku}
\end{reference}
In the situation of Lemma
\ref{lemma-orbit-space-equivariant-cohomology-support-filtration}, assume
that $\mathcal{A}$ is a $G$-sheaf of $k$-vector spaces for a field $k$,
that $G$ acts trivially on $X_0$, and that the restricted equivariant
structure on
$$
\mathcal{A}_0=\mathcal{A}|_{X_0}
$$
is the trivial $G$-action. Then
(\ref{equation-orbit-space-equivariant-support-complex}) is exact. For every
$n\geq 0$ there are canonical isomorphisms
$$
\begin{aligned}
H^n(X_0;G,\mathcal{A}_0)
&\cong
\bigoplus_{p+q=n}
\Hom_k\bigl(H_p(G,k),H^q(X_0,\mathcal{A}_0)\bigr)\\
&\cong
\bigoplus_{p+q=n}
H^p(G,k)\otimes_k H^q(X_0,\mathcal{A}_0).
\end{aligned}
$$
\end{lemma}

\begin{proof}
If $X_0$ is empty, all the groups in the displayed formula vanish. The map
$\rho^n$ of Lemma
\ref{lemma-orbit-space-equivariant-cohomology-support-filtration} is then a
map from zero to zero, so that lemma also gives exactness. Assume that
$X_0$ is nonempty and choose $x\in X_0$. Since $G$ acts trivially on
$X_0$, we have $G_x=G$. Thus $G$ is finite by Sheaves, Definition
\ref{sheaves-definition-discontinuous-group-action}.

The restriction
$$
q_0:X_0\longrightarrow Y_0
$$
of the orbit map is a homeomorphism. Indeed, it is a continuous bijection.
If $W\subset X_0$ is open, write $W=X_0\cap W'$ with $W'$ open in $X$.
Since $X_0$ is saturated and the orbit map is open, we have
$q_0(W)=Y_0\cap q(W')$, which is open in $Y_0$. The invariant-direct-image
isomorphisms in the preceding lemma and the trivial action on
$\mathcal{A}_0$ identify
$$
\mathcal{B}|_{Y_0}\cong q_{0,*}^G\mathcal{A}_0
\cong q_{0,*}\mathcal{A}_0.
$$

Under these identifications, $\rho^n$ is the left edge map
$$
H^n(Y_0,q_{0,*}^G\mathcal{A}_0)
\longrightarrow H^n(X_0;G,\mathcal{A}_0).
$$
By Lemma \ref{lemma-orbit-space-equivariant-spectral-sequence-edges}, its
composite with the right edge map to
$H^n(X_0,\mathcal{A}_0)^G$ and the inclusion of invariants is ordinary
pullback by $q_0$. The action on the latter cohomology group is trivial,
and pullback by the homeomorphism $q_0$ is an isomorphism. Hence every
$\rho^n$ is split injective. The preceding lemma identifies the only
possible defect of exactness of
(\ref{equation-orbit-space-equivariant-support-complex}) with
$\Ker(\rho^n)$, proving exactness.

It remains to prove the displayed formula. Lemma
\ref{lemma-sheaf-group-cohomology-change-coefficients} allows us to compute
the equivariant cohomology in the category of $G$-sheaves of
$k$-vector spaces. Since the actions on $X_0$ and $\mathcal{A}_0$ are
trivial, Lemma \ref{lemma-sheaf-group-operator-spectral-sequences} gives
the first displayed isomorphism.

Choose the standard bar resolution of the trivial $k[G]$-module $k$.
Since $G$ is finite, every term after tensoring over $k[G]$ with $k$ is a
finite-dimensional $k$-vector space. The universal coefficient
isomorphism over the field $k$ gives
$$
H^p(G,k)\cong\Hom_k(H_p(G,k),k),
$$
and evaluation gives the canonical isomorphism
$$
\Hom_k(H_p(G,k),k)\otimes_k E
\longrightarrow\Hom_k(H_p(G,k),E)
$$
for every $k$-vector space $E$. Applying this with
$E=H^q(X_0,\mathcal{A}_0)$ proves the second isomorphism.
\end{proof}

\begin{definition}
\label{definition-equivariant-cech-cohomology}
\begin{reference}
\cite[Chapter V, \S 5.5, p. 208]{Tohoku}
\end{reference}
Let $G$ act on a topological space $X$. A \emph{$G$-covering} of $X$
is a covering $\mathcal{U}=(U_i)_{i\in I}$ together with an action of
$G$ on $I$ such that
$$
gU_i=U_{gi}
$$
for every $g\in G$ and $i\in I$. It is called open, respectively closed,
if every $U_i$ is open, respectively closed.
The $G$-covering is called \emph{fixed-point-free} if the action of $G$
on $I$ is free, that is, if $gi\neq i$ for every $i\in I$ and every
$g\neq 1$.

Let $\mathcal{U}$ be an open $G$-covering and let $\mathcal{P}$ be an
abelian $G$-presheaf on $X$. Thus there are transport maps
$$
\rho_g(U):\mathcal{P}(U)\longrightarrow\mathcal{P}(gU)
$$
satisfying the analogues of the identities in Sheaves, Definition
\ref{sheaves-definition-G-sheaf}. The ordinary, nonalternating
\v Cech complex $\check{\mathcal{C}}^\bullet(\mathcal{U},\mathcal{P})$
is a complex of $G$-modules for the action
$$
(g\cdot s)_{i_0\ldots i_p}
=\rho_g(U_{g^{-1}i_0\ldots g^{-1}i_p})
\left(s_{g^{-1}i_0\ldots g^{-1}i_p}\right).
$$
We define the \emph{equivariant \v Cech cohomology} by
$$
\check H^n(\mathcal{U};G,\mathcal{P})
=H^n\left(R\Gamma_G(
\check{\mathcal{C}}^\bullet(\mathcal{U},\mathcal{P}))\right),
$$
where $\Gamma_G(M)=M^G$ on the category of $G$-modules.

More explicitly, let $P_\bullet\to\mathbf Z$ be a projective resolution
of the trivial $\mathbf Z[G]$-module. Then
$\check H^n(\mathcal{U};G,\mathcal{P})$ is the degree-$n$ cohomology of
the total complex of the first quadrant double complex
$$
C^{a,b}(\mathcal{U};G,\mathcal{P})
=\Hom_{\mathbf Z[G]}\left(
P_a,\check{\mathcal{C}}^b(\mathcal{U},\mathcal{P})\right).
$$
For the standard bar resolution this is the double complex of group
cochains with values in the \v Cech complex used by the source.
\end{definition}

\begin{lemma}[Equivariant \v Cech spectral sequence]
\label{lemma-equivariant-cech-spectral-sequence}
\begin{reference}
\cite[Chapter V, \S 5.5, Theorem 5.5.3, pp. 208--209]{Tohoku}
\end{reference}
Let $G$ act on a topological space $X$, let
$\mathcal{U}=(U_i)_{i\in I}$ be an open $G$-covering, and let
$\mathcal{A}$ be an abelian $G$-sheaf on $X$. For $q\geq 0$, let
$\underline H^q(\mathcal{A})$ be the abelian $G$-presheaf
$$
V\longmapsto H^q(V,\mathcal{A}),
$$
with transport induced by the action on $X$ and $\mathcal{A}$. There is
a bounded first quadrant spectral sequence, additive and functorial in
$\mathcal{A}$ from the $E_2$-page onward, with
\begin{equation}
\label{equation-equivariant-cech-spectral-sequence}
E_2^{p,q}(\mathcal{A})
=\check H^p(\mathcal{U};G,\underline H^q(\mathcal{A}))
\end{equation}
converging to $H^{p+q}(X;G,\mathcal{A})$ and inducing a finite filtration
on every group in the abutment.

The bottom edge maps are canonical homomorphisms
\begin{equation}
\label{equation-equivariant-cech-edge}
\check H^n(\mathcal{U};G,\mathcal{A})
\longrightarrow H^n(X;G,\mathcal{A}).
\end{equation}
They are isomorphisms for $n=0$ and injective for $n=1$.
The associated five term exact sequence begins
$$
\begin{aligned}
0 \longrightarrow
&\check H^1(\mathcal{U};G,\mathcal{A})
\longrightarrow H^1(X;G,\mathcal{A})
\longrightarrow
\check H^0(\mathcal{U};G,\underline H^1(\mathcal{A})) \\
\longrightarrow
&\check H^2(\mathcal{U};G,\mathcal{A})
\longrightarrow H^2(X;G,\mathcal{A}).
\end{aligned}
$$
\end{lemma}

\begin{proof}
Let
$$
\mathbf K(\mathcal{A})
=\check{\mathcal{C}}^\bullet(\mathcal{U},\mathcal{A}),
$$
regarded by Definition
\ref{definition-equivariant-cech-cohomology} as a nonnegative complex of
$G$-modules. This is an additive functor from abelian $G$-sheaves on $X$
to $\text{CoCh}_{\geq 0}(\text{Mod}_G)$. The sheaf condition and Lemma
\ref{lemma-cech-h0} give a natural isomorphism
$$
\Gamma(X,\mathcal{A})\longrightarrow H^0(\mathbf K(\mathcal{A}))
$$
of $G$-modules.

The category of abelian $G$-sheaves on $X$ and the category of
$G$-modules have enough injectives by Injectives, Lemmas
\ref{injectives-lemma-equivariant-modules-enough-injectives} and
\ref{injectives-lemma-G-modules}, respectively.

Let $\mathcal{I}$ be an injective abelian $G$-sheaf. By Injectives,
Lemma \ref{injectives-lemma-equivariant-injective-underlying}, its
underlying abelian sheaf is injective. Lemma
\ref{lemma-injective-trivial-cech} therefore gives
$$
H^n(\mathbf K(\mathcal{I}))=0\quad\text{for }n>0.
$$
Moreover, Injectives, Lemma
\ref{injectives-lemma-equivariant-injective-sections-invariants} says
that $\Gamma(X,\mathcal{I})$ is an injective $\mathbf Z[G]$-module. It is
therefore right acyclic for the invariant-module functor.

Apply Derived Categories, Lemma
\ref{derived-lemma-complex-valued-functor-spectral-sequence} to global
sections with values in $G$-modules, the invariant-module functor, and
$\mathbf K$. Its abutment is the right derived functor of invariant global
sections, hence is $H^*(X;G,\mathcal{A})$ by Definition
\ref{definition-orbit-space-equivariant-cohomology}.

It remains to identify the derived complexes of $\mathbf K$. Choose an
injective resolution $\mathcal{A}\to\mathcal{I}^\bullet$ in the category
of abelian $G$-sheaves. For every $r\geq 0$, exactness of products of
abelian groups gives
$$
\begin{aligned}
(R^q\mathbf K(\mathcal{A}))^r
&=H^q\left(
\prod_{(i_0,\ldots,i_r)\in I^{r+1}}
\mathcal{I}^\bullet(U_{i_0\ldots i_r})\right)\\
&=\prod_{(i_0,\ldots,i_r)\in I^{r+1}}
H^q(U_{i_0\ldots i_r},\mathcal{A}).
\end{aligned}
$$
These identifications commute with the \v Cech differentials and the
$G$-actions. Consequently there is a canonical isomorphism of complexes
$$
R^q\mathbf K(\mathcal{A})
\cong
\check{\mathcal{C}}^\bullet(
\mathcal{U},\underline H^q(\mathcal{A})).
$$
The hypercohomology identification in Definition
\ref{definition-equivariant-cech-cohomology} now gives
(\ref{equation-equivariant-cech-spectral-sequence}).

The maps (\ref{equation-equivariant-cech-edge}) are the bottom edge maps
of this first quadrant spectral sequence. In total degrees zero and one
there are no nonzero differentials which can alter the terms in bidegrees
$(0,0)$ and $(1,0)$, which proves the last assertion.
\end{proof}

\begin{definition}[Equivariant \v Cech complex for a closed covering]
\label{definition-equivariant-cech-closed-covering}
\begin{reference}
\cite[Chapter V, \S 5.5, pp. 208--209]{Tohoku}
\end{reference}
Let $G$ act on a topological space $X$, and let
$\mathcal U=(U_i)_{i\in I}$ be a closed $G$-covering.  Thus every $U_i$ is
closed, $gU_i=U_{gi}$, and the $U_i$ cover $X$.  Put
$U_{i_0\ldots i_p}=U_{i_0}\cap\cdots\cap U_{i_p}$.  If $\mathcal P$ is an
abelian $G$-equivariant coefficient system on these finite intersections,
write
$$
C^p_{\mathrm{cl}}(\mathcal U,\mathcal P)=
\prod_{(i_0,\ldots,i_p)\in I^{p+1}}
\mathcal P(U_{i_0\ldots i_p}).
$$
Here $\mathcal P(U_{\mathbf i})$ means
$\Gamma(U_{\mathbf i},\mathcal P|_{U_{\mathbf i}})$ when $\mathcal P$ is a
sheaf; the coefficient-system transport maps carry this group to the group
on $gU_{\mathbf i}$.  The alternating restriction maps make this a complex
of $G$-modules.  We define
$$
\check H^n_{\mathrm{cl}}(\mathcal U;G,\mathcal P)=
H^n\left(R\Gamma_G\bigl(
C^\bullet_{\mathrm{cl}}(\mathcal U,\mathcal P)\bigr)\right).
$$
For a $G$-sheaf $\mathcal A$, the coefficient system used below is
$$
\underline H^q_{\mathrm{cl}}(\mathcal A)(U_{i_0\ldots i_p})
 =H^q(U_{i_0\ldots i_p},\mathcal A|_{U_{i_0\ldots i_p}}),
$$
with transport induced by the action.  This is the same derived-invariants
construction as Definition \ref{definition-equivariant-cech-cohomology};
the subscript only records that the covering members need not be open.
\end{definition}

\begin{lemma}[Equivariant \v Cech spectral sequence for closed coverings]
\label{lemma-equivariant-cech-spectral-sequence-closed-covering}
\begin{reference}
\cite[Chapter V, \S 5.5, Theorem 5.5.3 and Corollary 2,
pp. 208--209]{Tohoku}
\end{reference}
Let $G$ act on a topological space $X$, let
$\mathcal U=(U_i)_{i\in I}$ be a closed $G$-covering, and let $\mathcal A$ be
an abelian $G$-sheaf on $X$.  Assume either that $I$ is finite, or that $X$
is paracompact and $\mathcal U$ is locally finite.  There is a first-quadrant
spectral sequence, additive and functorial in $\mathcal A$ from the
$E_2$-page onward,
$$
E_2^{p,q}=\check H^p_{\mathrm{cl}}
 (\mathcal U;G,\underline H^q_{\mathrm{cl}}(\mathcal A))
 \Longrightarrow H^{p+q}(X;G,\mathcal A),
$$
with a finite filtration in each total degree.  Its bottom edge maps are
canonical functorial homomorphisms
$$
\check H^n_{\mathrm{cl}}(\mathcal U;G,\mathcal A)
 \longrightarrow H^n(X;G,\mathcal A).
$$
They are isomorphisms for $n=0$ and injective for $n=1$, and the associated
five-term exact sequence begins
$$
\begin{aligned}
0\longrightarrow{}&\check H^1_{\mathrm{cl}}(\mathcal U;G,\mathcal A)
\longrightarrow H^1(X;G,\mathcal A)
\longrightarrow \check H^0_{\mathrm{cl}}
 (\mathcal U;G,\underline H^1_{\mathrm{cl}}(\mathcal A))\\
\longrightarrow{}&\check H^2_{\mathrm{cl}}(\mathcal U;G,\mathcal A)
\longrightarrow H^2(X;G,\mathcal A).
\end{aligned}
$$
\end{lemma}

\begin{proof}
For a tuple $\mathbf i=(i_0,\ldots,i_r)$ let
$j_{\mathbf i}:U_{\mathbf i}\hookrightarrow X$ be the closed inclusion and
put
$$
\mathcal K^r(\mathcal U,\mathcal A)=
\prod_{\mathbf i\in I^{r+1}}
(j_{\mathbf i})_*(\mathcal A|_{U_{\mathbf i}}),
$$
with the alternating restriction differential.  The augmentation
$$
\mathcal A[0]\longrightarrow
\mathcal K^\bullet(\mathcal U,\mathcal A)
$$
is equivariant, where $G$ acts by transport and permutation of the factors.
It is a quasi-isomorphism.  Indeed, the stalk argument in the proof of Lemma
\ref{lemma-cech-spectral-sequence-closed-covering} applies unchanged: in the
finite case, and locally in the locally finite case, the nonzero factors are
the augmented cochain complex of a nonempty simplex.

Let $F$ be global sections from abelian $G$-sheaves to $G$-modules and let
$G_0$ be the invariant-module functor.  If $\mathcal I$ is an injective
abelian $G$-sheaf, then $F(\mathcal I)$ is right acyclic for $G_0$ by
Injectives, Lemma
\ref{injectives-lemma-equivariant-injective-sections-invariants}.  Apply
Derived Categories, Lemma
\ref{derived-lemma-hyperderived-composite-functor-spectral-sequence-resolution}
to the displayed equivariant resolution.  Its $E_2$-page is
$$
H^p\left(RG_0\bigl(
(R^qF)(\mathcal K^\bullet(\mathcal U,\mathcal A))\bigr)\right)
$$
and it converges to $R^{p+q}(G_0\circ F)(\mathcal A)$.

It remains to identify the complex inside $RG_0$.  The cohomology
calculation in the proof of Lemma
\ref{lemma-cech-spectral-sequence-closed-covering} is functorial and
equivariant.  Injectives, Lemma
\ref{injectives-lemma-equivariant-injective-underlying} identifies the
underlying groups of $R^qF$ with ordinary sheaf cohomology, carrying the
induced $G$-action.  In the finite case the closed-cover calculation uses
exact direct image and finite products for closed inclusions; in the locally
finite paracompact case it uses the soft resolution and locally finite
products.  It gives
$$
(R^qF)(\mathcal K^\bullet(\mathcal U,\mathcal A))
\cong
C^\bullet_{\mathrm{cl}}(
\mathcal U,\underline H^q_{\mathrm{cl}}(\mathcal A))
$$
as complexes of $G$-modules.  Definition
\ref{definition-equivariant-cech-closed-covering} now identifies the
$E_2$-page with the one in the statement, while Definition
\ref{definition-orbit-space-equivariant-cohomology} identifies the abutment
with $H^{p+q}(X;G,\mathcal A)$.  The general edge and convergence statements
give the finite filtrations, the bottom edge maps, and the five-term exact
sequence asserted above.
\end{proof}

\begin{lemma}
\label{lemma-equivariant-cech-spectral-sequence-closed-acyclic-cover}
\begin{reference}
\cite[Chapter V, \S 5.5, Corollaries 1 and 2 to Theorem 5.5.3,
pp. 208--209]{Tohoku}
\end{reference}
In the situation of Lemma
\ref{lemma-equivariant-cech-spectral-sequence-closed-covering}, assume in
addition that
$$
H^q(U_{i_0\ldots i_p},\mathcal A|_{U_{i_0\ldots i_p}})=0
\quad\text{for every }q>0\text{ and every }(i_0,\ldots,i_p).
$$
Then all bottom edge maps
$$
\check H^n_{\mathrm{cl}}(\mathcal U;G,\mathcal A)
\longrightarrow H^n(X;G,\mathcal A)
$$
are isomorphisms.
\end{lemma}

\begin{proof}
The $E_2$-page of Lemma
\ref{lemma-equivariant-cech-spectral-sequence-closed-covering} is
concentrated in the row $q=0$.  The spectral sequence therefore degenerates,
and its bottom edge maps are isomorphisms in every degree.
\end{proof}

\begin{lemma}
\label{lemma-equivariant-cech-spectral-sequence-acyclic-cover}
\begin{reference}
\cite[Chapter V, \S 5.5, Corollary 1 to Theorem 5.5.3, p. 209]{Tohoku}
\end{reference}
In the situation of Lemma
\ref{lemma-equivariant-cech-spectral-sequence}, assume that
$$
H^q(U_{i_0\ldots i_p},\mathcal{A})=0
$$
for every $q>0$, every $p\geq 0$, and every
$i_0,\ldots,i_p\in I$. Then the edge maps
(\ref{equation-equivariant-cech-edge}) are isomorphisms for every
$n\geq 0$.
\end{lemma}

\begin{proof}
Under the stated hypothesis the $E_2$-page of
(\ref{equation-equivariant-cech-spectral-sequence}) is concentrated in
the row $q=0$. Hence the spectral sequence degenerates and its bottom
edge maps are isomorphisms.
\end{proof}


\begin{lemma}
\label{lemma-equivariant-cech-free-index}
\begin{reference}
\cite[Chapter V, \S 5.5, Proposition 5.5.4,
pp. 209--210]{Tohoku}
\end{reference}
Let $I$ be a free $G$-set. For every $r\geq 0$ and
$\mathbf{i}=(i_0,\ldots,i_r)\in I^{r+1}$, let $M_{\mathbf{i}}$ be an
abelian group. Suppose given isomorphisms
$$
\rho_g:M_{\mathbf{i}}\longrightarrow M_{g\mathbf{i}}
$$
satisfying $\rho_1=1$ and $\rho_g\rho_h=\rho_{gh}$. Put
$$
C^r=\prod_{\mathbf{i}\in I^{r+1}}M_{\mathbf{i}}
$$
with $G$-action
$$
(g\cdot s)_{\mathbf{i}}
=\rho_g(s_{g^{-1}\mathbf{i}}).
$$
For a set $R_r$ of representatives of $I^{r+1}/G$, there is an
isomorphism of $G$-modules
\begin{equation}
\label{equation-equivariant-cech-free-index-product}
C^r\cong
\prod_{\mathbf{i}\in R_r}
\operatorname{Coind}_{\{1\}}^G(M_{\mathbf{i}})
\cong
\operatorname{Coind}_{\{1\}}^G
\left(\prod_{\mathbf{i}\in R_r}M_{\mathbf{i}}\right).
\end{equation}
In particular, every $C^r$ is right acyclic for the invariant-module
functor.

Suppose that $G$-equivariant differentials make $C^\bullet$ a
nonnegative complex of $G$-modules. Then the canonical maps
$$
H^n((C^\bullet)^G)
\longrightarrow H^n(R\Gamma_G(C^\bullet))
$$
are isomorphisms for every $n\geq 0$. Consequently, if $\mathcal{U}$ is
any fixed-point-free $G$-covering for which a nonalternating \v Cech
coefficient complex is defined, the same conclusion applies to that
complex. In particular, if $\mathcal{U}$ is open and $\mathcal{P}$ is an
abelian $G$-presheaf, then
\begin{equation}
\label{equation-equivariant-cech-free-cover}
\check H^n(\mathcal{U};G,\mathcal{P})
\cong
H^n\left(
\check{\mathcal{C}}^\bullet(\mathcal{U},\mathcal{P})^G
\right).
\end{equation}
\end{lemma}

\begin{proof}
The diagonal action of $G$ on $I^{r+1}$ is free. For
$\mathbf{i}\in R_r$, the map from the product over its orbit to the
coinduced module is
$$
(s_{g\mathbf{i}})_{g\in G}
\longmapsto
(\rho_g(s_{g^{-1}\mathbf{i}}))_{g\in G}.
$$
It is an isomorphism, and it is $G$-equivariant for the convention
$$
(h m)_g=m_{gh}
$$
on the coinduced module. Taking the product over $R_r$, and commuting
the two products, proves
(\ref{equation-equivariant-cech-free-index-product}).

For an abelian group $M$, the coinduced module
$\operatorname{Coind}_{\{1\}}^G(M)$ is right acyclic for invariants.
Indeed, coinduction is exact and is right adjoint to the exact forgetful
functor. It therefore carries an injective resolution of $M$ to an
injective resolution of the coinduced module, and taking invariants of
that resolution recovers the original resolution of $M$.

Thus every term of $C^\bullet$ is right acyclic for invariants. Derived
Categories, Lemma \ref{derived-lemma-leray-acyclicity} identifies the
termwise invariant complex with $R\Gamma_G(C^\bullet)$, proving the
first assertion about cohomology. For the final assertion take
$$
M_{\mathbf{i}}
=\mathcal{P}(U_{i_0\ldots i_r})
$$
and use Definition \ref{definition-equivariant-cech-cohomology}.
\end{proof}


\begin{lemma}
\label{lemma-equivariant-cech-nonhomogeneous-cochains}
\begin{reference}
\cite[Chapter V, \S 5.5, (5.5.9)--(5.5.11),
pp. 210--211]{Tohoku}
\end{reference}
Let $\mathcal{U}=(U_i)_{i\in I}$ be a fixed-point-free $G$-covering
equipped with an abelian $G$-equivariant coefficient system
$\mathcal{P}$ on its finite intersections, so that its nonalternating
\v Cech complex is defined. This includes the case where $\mathcal{U}$ is
open and $\mathcal{P}$ is an abelian $G$-presheaf. Choose a set $J$ of
representatives for $I/G$ and identify
$$
I=G\times J,
\qquad
U_{g j}=gU_j.
$$
Here and below we identify $j\in J$ with $(1,j)\in G\times J$.
An invariant $n$-cochain
$$
f=(f_{g_0j_0,\ldots,g_nj_n})
$$
is uniquely determined by the nonhomogeneous cochain
\begin{equation}
\label{equation-equivariant-cech-nonhomogeneous-restriction}
F_{j_0,g_1j_1,\ldots,g_nj_n}
=f_{j_0,g_1j_1,\ldots,g_nj_n},
\end{equation}
whose indicated component belongs to
$$
\mathcal{P}(U_{j_0}\cap g_1U_{j_1}\cap\cdots\cap g_nU_{j_n}).
$$
The inverse correspondence is
\begin{equation}
\label{equation-equivariant-cech-nonhomogeneous-reconstruction}
f_{g_0j_0,\ldots,g_nj_n}
=g_0\cdot
F_{j_0,g_0^{-1}g_1j_1,\ldots,g_0^{-1}g_nj_n}.
\end{equation}
Under this correspondence the \v Cech differential is
\begin{equation}
\label{equation-equivariant-cech-nonhomogeneous-differential}
\begin{aligned}
(dF)_{j_0,g_1j_1,\ldots,g_{n+1}j_{n+1}}
={}&g_1\cdot
F_{j_1,g_1^{-1}g_2j_2,\ldots,
g_1^{-1}g_{n+1}j_{n+1}} \\
&+\sum_{\alpha=1}^{n+1}(-1)^\alpha
F_{j_0,g_1j_1,\ldots,
\widehat{g_\alpha j_\alpha},\ldots,
g_{n+1}j_{n+1}}.
\end{aligned}
\end{equation}
Every term on the right is implicitly restricted to the intersection
indexed by the left side.

In the single-orbit case, write $D=U_j$. Then the nonhomogeneous
$n$-cochains are families
$$
F_{g_1,\ldots,g_n}
\in
\mathcal{P}(D\cap g_1D\cap\cdots\cap g_nD),
$$
and the differential becomes
\begin{equation}
\label{equation-equivariant-cech-nonhomogeneous-single-orbit}
\begin{aligned}
(dF)_{g_1,\ldots,g_{n+1}}
={}&g_1\cdot
F_{g_1^{-1}g_2,\ldots,g_1^{-1}g_{n+1}} \\
&+\sum_{\alpha=1}^{n+1}(-1)^\alpha
F_{g_1,\ldots,\widehat{g_\alpha},\ldots,g_{n+1}}.
\end{aligned}
\end{equation}
If $\mathcal{P}(\emptyset)=0$, as is the case when $\mathcal{P}$ is a
sheaf, then this complex is determined by the set
$$
L=\{g\in G\mid D\cap gD\neq\emptyset\},
$$
the partially defined products $g^{-1}g'$ in $L$, the groups on the
finite intersections of the translates indexed by $L$, their restriction
maps, and the corresponding transport maps.
\end{lemma}

\begin{proof}
Equation
(\ref{equation-equivariant-cech-nonhomogeneous-restriction}) is restriction
to the components whose first group coordinate is $1$. Invariance gives
(\ref{equation-equivariant-cech-nonhomogeneous-reconstruction}), and the
latter formula also shows that every nonhomogeneous cochain determines a
unique invariant cochain. Substitute
(\ref{equation-equivariant-cech-nonhomogeneous-reconstruction}) in the
ordinary nonalternating \v Cech differential and separate the term which
omits the first index. This gives
(\ref{equation-equivariant-cech-nonhomogeneous-differential}). The
single-orbit formula is its specialization. If an indicated total
intersection is nonempty, then every $g_i$ and every $g_i^{-1}g_j$ used by
the differential belongs to $L$. Components on empty intersections are
zero by the final hypothesis, which proves the last assertion.
\end{proof}


\begin{lemma}
\label{lemma-equivariant-cech-refinements}
\begin{reference}
\cite[Chapter V, \S 5.5, (5.5.12), pp. 211--212]{Tohoku}
\end{reference}
\begin{history}
The cited source replaces only the coarser covering by a covering whose
index set is free and then assumes that an equivariant refinement function
from the finer covering exists. This need not be the case: a $G$-set with a
fixed point admits no $G$-map to a free $G$-set. The proof below replaces
both coverings. The two replacement maps are equivariant
quasi-isomorphisms, which gives the intended canonical refinement map.
\end{history}
Let $G$ act on a topological space $X$, let
$$
\mathcal{U}=(U_i)_{i\in I}
\quad\text{and}\quad
\mathcal{V}=(V_j)_{j\in J}
$$
be open $G$-coverings, and let $\mathcal{P}$ be an abelian
$G$-presheaf. If $\mathcal{V}$ refines $\mathcal{U}$ as an ordinary
open covering, then there are canonical homomorphisms
\begin{equation}
\label{equation-equivariant-cech-refinement}
\check H^n(\mathcal{U};G,\mathcal{P})
\longrightarrow
\check H^n(\mathcal{V};G,\mathcal{P}),
\qquad n\geq 0.
\end{equation}
They are functorial in $\mathcal{P}$ and transitive under further
refinement.

If the refinement is represented by a $G$-equivariant map
$c:J\to I$ such that $V_j\subset U_{c(j)}$, then
(\ref{equation-equivariant-cech-refinement}) is induced by the usual
refinement map of \v Cech complexes. In general, put
$$
I^{\mathrm{fr}}=G\times I,
\qquad
U^{\mathrm{fr}}_{(g,i)}=gU_i,
$$
with $h(g,i)=(hg,i)$, and define $\mathcal{V}^{\mathrm{fr}}$ in the
same way. The equivariant maps
$$
\psi_{\mathcal U}:I^{\mathrm{fr}}\longrightarrow I,
\quad (g,i)\longmapsto gi
$$
and $\psi_{\mathcal V}$ induce isomorphisms on equivariant \v Cech
cohomology. Given any refinement function $c:J\to I$, the canonical map
is the composite
$$
\check H^n(\mathcal U;G,\mathcal P)
\xrightarrow{\sim}
\check H^n(\mathcal U^{\mathrm{fr}};G,\mathcal P)
\longrightarrow
\check H^n(\mathcal V^{\mathrm{fr}};G,\mathcal P)
\xleftarrow{\sim}
\check H^n(\mathcal V;G,\mathcal P),
$$
where the middle arrow is induced by
$$
c^{\mathrm{fr}}:G\times J\longrightarrow G\times I,
\qquad (g,j)\longmapsto(g,c(j)),
$$
and the last isomorphism is inverted in forming the composite.
\end{lemma}

\begin{proof}
Suppose first that $c:J\to I$ is equivariant. The ordinary refinement
map
$$
\check{\mathcal C}^\bullet(\mathcal U,\mathcal P)
\longrightarrow
\check{\mathcal C}^\bullet(\mathcal V,\mathcal P)
$$
is then a map of complexes of $G$-modules. If $c'$ is a second
equivariant refinement function, the usual prism homotopy of Section
\ref{section-refinements-cech} is $G$-equivariant: every index appearing
in its formula is obtained from $c$ and $c'$, and restriction commutes
with transport. Thus the two maps agree in the derived category of
$G$-modules and induce the same map after $R\Gamma_G$.

For an arbitrary $\mathcal U$, the map $\psi_{\mathcal U}$ is an
equivariant refinement function. The nonequivariant map
$$
s:I\longrightarrow G\times I,
\qquad i\longmapsto(1,i)
$$
is a refinement function in the opposite direction. We have
$\psi_{\mathcal U}s=1_I$, while $s\psi_{\mathcal U}$ and
$1_{G\times I}$ induce homotopic ordinary refinement maps. Consequently
$$
\check{\mathcal C}^\bullet(\mathcal U,\mathcal P)
\longrightarrow
\check{\mathcal C}^\bullet(\mathcal U^{\mathrm{fr}},\mathcal P)
$$
is a quasi-isomorphism. It is equivariant, even though the auxiliary
map $s$ need not be. Hence it becomes an isomorphism after applying
$R\Gamma_G$. The same holds for $\mathcal V$.

Now choose an ordinary refinement function $c:J\to I$. The map
$c^{\mathrm{fr}}$ in the statement is equivariant and satisfies
$$
V^{\mathrm{fr}}_{(g,j)}=gV_j
\subset gU_{c(j)}=U^{\mathrm{fr}}_{c^{\mathrm{fr}}(g,j)}.
$$
This proves the existence of the displayed composite. Two choices of
$c$ give equivariantly homotopic middle maps, so the composite is
canonical. If $c$ was equivariant to begin with, then
$$
\psi_{\mathcal U}c^{\mathrm{fr}}=c\psi_{\mathcal V}.
$$
Thus the strictly commutative square of refinement maps shows that the
displayed composite is the usual map induced by $c$. Given a third
refinement, the free replacement of the middle
covering occurs once with each orientation and the two isomorphisms
cancel. The composite of the two equivariant refinement functions is
the one obtained from the composite ordinary refinement function. This
proves transitivity. Functoriality in $\mathcal P$ is immediate from the
construction.
\end{proof}

\begin{definition}
\label{definition-equivariant-cech-cohomology-space}
\begin{reference}
\cite[Chapter V, \S 5.5, (5.5.13), p. 212]{Tohoku}
\end{reference}
Let $G$ act on a topological space $X$ and let $\mathcal P$ be an
abelian $G$-presheaf. We define the equivariant \v Cech cohomology of
$X$ by
$$
\check H^n(X;G,\mathcal P)
=\colim_{\mathcal U}\check H^n(\mathcal U;G,\mathcal P),
$$
where $\mathcal U$ runs through the open $G$-coverings of $X$, ordered
by ordinary refinement, and the transition maps are those of Lemma
\ref{lemma-equivariant-cech-refinements}.
\end{definition}

\begin{lemma}
\label{lemma-equivariant-cech-free-coverings-cofinal}
The preorder of open $G$-coverings is filtered, and its subpreorder of
fixed-point-free open $G$-coverings is cofinal. Consequently the colimit
in Definition \ref{definition-equivariant-cech-cohomology-space} may be
taken over the fixed-point-free open $G$-coverings.
\end{lemma}

\begin{proof}
The intersection covering of two open $G$-coverings, indexed by the
product with its diagonal $G$-action, is a common refinement. Thus the
coverings form a filtered preorder. The construction and transitivity in
Lemma \ref{lemma-equivariant-cech-refinements} make the displayed colimit
well defined. There is no size issue. Indeed, the $G$-stable subsets
$S\subset \operatorname{Open}(X)$ which cover $X$, indexed by their
distinct open members, form a set of coverings cofinal among all open
$G$-coverings. Their free replacements, indexed by $G\times S$, form a
set cofinal among the fixed-point-free coverings. Finally,
$\mathcal U^{\mathrm{fr}}$ is a fixed-point-free refinement of every
$\mathcal U$, and the corresponding transition map is an isomorphism.
Hence the fixed-point-free coverings are cofinal for this colimit.
\end{proof}

\begin{lemma}[Global equivariant \v Cech spectral sequence]
\label{lemma-global-equivariant-cech-spectral-sequence}
\begin{reference}
\cite[Chapter V, \S 5.5, (5.5.14)--(5.5.15), p. 212]{Tohoku}
\end{reference}
Let $G$ act on a topological space $X$ and let $\mathcal A$ be an
abelian $G$-sheaf. There is an additive functorial first quadrant
spectral sequence
\begin{equation}
\label{equation-global-equivariant-cech-spectral-sequence}
E_2^{p,q}(\mathcal A)
=\check H^p(X;G,\underline H^q(\mathcal A))
\Longrightarrow H^{p+q}(X;G,\mathcal A),
\end{equation}
which induces a finite filtration on every group in the abutment. Its
bottom edge maps are canonical functorial homomorphisms
\begin{equation}
\label{equation-global-equivariant-cech-edge}
\check H^n(X;G,\mathcal A)
\longrightarrow H^n(X;G,\mathcal A).
\end{equation}
\end{lemma}

\begin{proof}
We use the cofinal filtered system of fixed-point-free open
$G$-coverings from Lemma
\ref{lemma-equivariant-cech-free-coverings-cofinal}. If
$\mathcal V$ refines $\mathcal U$ in this system, an equivariant
refinement function exists: choose it on one representative of every
free $G$-orbit of the index set of $\mathcal V$ and extend it
equivariantly. The resulting map of \v Cech complexes is compatible with
the augmentations from global sections. Applying it to an equivariant
injective resolution of $\mathcal A$ gives a morphism between the double
complexes used in the proof of Lemma
\ref{lemma-equivariant-cech-spectral-sequence}. The equivariant prism
homotopy proves independence of the refinement function. We therefore
obtain a filtered system of the fixed-cover spectral sequences whose
maps on the common abutment $H^*(X;G,\mathcal A)$ are identities.

Take the filtered colimit page by page. Filtered colimits of abelian
groups are exact, so they commute with the cohomology which passes from
one page to the next. The $E_2$-page is the one in
(\ref{equation-global-equivariant-cech-spectral-sequence}) by Definition
\ref{definition-equivariant-cech-cohomology-space}. In every total degree
the fixed-cover spectral sequences have uniformly finite filtrations. If
$F^p_{\mathcal U}H^n$ denotes the filtration for $\mathcal U$, put
$$
F^pH^n=\colim_{\mathcal U}F^p_{\mathcal U}H^n.
$$
Exactness of filtered colimits identifies its graded pieces with the
colimits of the graded pieces, and the first quadrant condition gives
$F^0H^n=H^n(X;G,\mathcal A)$ and $F^{n+1}H^n=0$. Thus this is a finite
filtration of the same abutment. The
colimit of the bottom edge maps is
(\ref{equation-global-equivariant-cech-edge}). Additivity and
functoriality in $\mathcal A$ follow from the corresponding fixed-cover
properties.
\end{proof}

\begin{lemma}
\label{lemma-equivariant-cech-zero-associated-sheaf-degree-zero}
Let $G$ act on a topological space $X$ and let $\mathcal P$ be an
abelian $G$-presheaf whose associated sheaf is zero. Then
$$
\check H^0(X;G,\mathcal P)=0.
$$
\end{lemma}

\begin{proof}
By Lemmas \ref{lemma-equivariant-cech-free-coverings-cofinal} and
\ref{lemma-equivariant-cech-free-index}, a class in
$\check H^0(X;G,\mathcal P)$ is represented on a fixed-point-free open
$G$-covering $\mathcal U=(U_i)_{i\in I}$ by an invariant ordinary \v Cech
$0$-cocycle $(s_i)_{i\in I}$. Let $J$ be the set of pairs $(i,V)$ where
$V\subset U_i$ is nonempty open and $s_i|_V=0$. The action
$$
g(i,V)=(gi,gV)
$$
makes $(V)_{(i,V)\in J}$ a fixed-point-free open $G$-covering refining
$\mathcal U$: the vanishing of the associated sheaf says that these opens
cover every $U_i$, and invariance of $(s_i)$ makes $J$ stable under $G$.
The pulled back cocycle is zero, which proves the lemma.
\end{proof}


\begin{lemma}
\label{lemma-equivariant-cech-zero-associated-sheaf}
\begin{reference}
\cite[Chapter V, \S 5.5, Lemma 5.5.5, pp. 212--213]{Tohoku}
\end{reference}
\begin{history}
Our nonalternating \v Cech complex has a factor
$\mathcal P(\emptyset)$ for every tuple with empty intersection. Thus the
normalization $\mathcal P(\emptyset)=0$ in the statement is necessary in
our convention; the cited source defines presheaves only on nonempty open
subsets, so this is the canonical extension of its convention. For
example, if $X$ is a discrete two point space,
$G=\{1\}$, $\mathcal P(\emptyset)=M\neq 0$, and
$\mathcal P(U)=0$ for $U\neq\emptyset$, then the associated sheaf is zero
but the two cross components give a nonzero class in degree one.

The cited proof also passes from vanishing on neighbourhoods of finitely
many points to vanishing on their union. That passage is not valid for a
nonseparated presheaf. The proof below instead constructs an equivariant
refinement whose individual members lie in neighbourhoods on which the
relevant sections are zero.
\end{history}
Let $G$ act discontinuously on a Hausdorff space $X$, and let
$q:X\to Y=X/G$ be the orbit map. Assume that $Y$ is paracompact. Let
$\mathcal P$ be an abelian $G$-presheaf on $X$ such that
$$
\mathcal P(\emptyset)=0
$$
and whose associated sheaf is zero. Then
$$
\check H^n(X;G,\mathcal P)=0
$$
for every $n>0$.
\end{lemma}

\begin{proof}
Let $\xi\in\check H^n(X;G,\mathcal P)$. By Lemma
\ref{lemma-equivariant-cech-free-coverings-cofinal}, it is represented on
a fixed-point-free open $G$-covering
$$
\mathcal U=(U_i)_{i\in I}.
$$
Lemma \ref{lemma-equivariant-cech-free-index} allows us to represent it by
an invariant cocycle
$$
f\in\check{\mathcal C}^n(\mathcal U,\mathcal P)^G.
$$

We first construct a pair of equivariant coverings adapted to the
discontinuous action. For every $x\in X$, choose $i(x)\in I$ with
$x\in U_{i(x)}$ and put $H_x=G_x$. Lemma
\ref{sheaves-lemma-discontinuous-group-action-neighbourhoods} and the
finiteness of $H_x$ give an open $H_x$-invariant neighbourhood $S_x$ of
$x$ such that
$$
S_x\subset U_{i(x)}
\quad\text{and}\quad
gS_x\cap S_x=\emptyset\quad\text{for }g\notin H_x.
$$
Indeed, intersect a neighbourhood supplied by that lemma with the finitely
many opens $hU_{i(x)}$, $h\in H_x$.

The opens $q(S_x)$ cover $Y$. By Topology, Lemma
\ref{topology-lemma-paracompact-shrinking}, there are a locally finite
open covering $(V_a)_{a\in A}$ of $Y$, an open covering
$(W_a)_{a\in A}$ of $Y$, and choices $x(a)\in X$ such that
$$
\overline{W_a}\subset V_a\subset q(S_{x(a)}).
$$
Write $S_a=S_{x(a)}$, $H_a=H_{x(a)}$, and $i_a=i(x(a))$. For
$(g,a)\in G\times A$ set
$$
\mathcal V_{g,a}=gS_a\cap q^{-1}(V_a),
\qquad
\mathcal W_{g,a}=gS_a\cap q^{-1}(W_a).
$$
With the action on the first factor, $\mathcal V$ and $\mathcal W$ are
fixed-point-free open $G$-coverings. The map
$$
G\times A\longrightarrow I,
\qquad (g,a)\longmapsto gi_a
$$
is an equivariant refinement function from $\mathcal V$ to $\mathcal U$.
Pull $f$ back to $\mathcal V$ and denote the resulting invariant cocycle
again by $f$.

We record two properties of these coverings. First,
$$
\overline{\mathcal W_{g,a}}\subset\mathcal V_{g,a}.
$$
Indeed,
$$
q^{-1}(q(S_a))=\coprod_{G/H_a}gS_a,
$$
so every distinct translate $gS_a$ is both open and closed in this inverse
image; now use $\overline{W_a}\subset V_a\subset q(S_a)$. Second, the
indexed family $(\mathcal W_{g,a})_{(g,a)\in G\times A}$ is locally finite.
To see this at $z\in X$, first use local finiteness of
$(\overline{W_a})$ near $q(z)$. For each of the finitely many remaining
$a$, either $q(z)\notin\overline{W_a}$, in which case a neighbourhood of
$z$ avoids all the corresponding members, or $q(z)\in V_a$, in which
case a neighbourhood of $z$ lies in one translate of $S_a$ and meets only
the finitely many duplicate indices coming from $H_a$.

Choose a set $R\subset X$ of representatives for $X/G$. For every
$r\in R$, choose an index $\tau(r)\in G\times A$ such that
$r\in\mathcal W_{\tau(r)}$. The closure inclusion and local finiteness just
proved give an open neighbourhood $M_r$ of $r$ such that
\begin{equation}
\label{equation-equivariant-cech-zero-shrinking}
M_r\subset\mathcal W_{\tau(r)}
\quad\text{and}\quad
M_r\cap\mathcal W_j\neq\emptyset
\ \Longrightarrow\ M_r\subset\mathcal V_j.
\end{equation}
This is the same elementary shrinking used in the proof of Topology,
Lemma \ref{topology-lemma-paracompact-refine-covering}: if
$r\in\overline{\mathcal W_j}$ shrink inside $\mathcal V_j$, and otherwise
shrink to avoid $\mathcal W_j$; only finitely many conditions occur.

Only finitely many members of $\mathcal V$ contain $r$. Hence there are
only finitely many $(n+1)$-tuples $j_0,\ldots,j_n$ with
$$
r\in\mathcal V_{j_0\ldots j_n}.
$$
Shrink $M_r$ inside every $\mathcal V_j$ which contains $r$.
For each such tuple the germ at $r$ of the section $f_{j_0\ldots j_n}$
is zero, because the associated sheaf of $\mathcal P$ is zero. Shrinking
$M_r$ finitely many more times, we may therefore assume that
\begin{equation}
\label{equation-equivariant-cech-zero-local}
f_{j_0\ldots j_n}|_{M_r}=0
\quad\text{whenever}\quad
r\in\mathcal V_{j_0\ldots j_n}.
\end{equation}
Finally put
$$
N_r=\bigcap_{h\in G_r}hM_r.
$$
This is an open $G_r$-invariant neighbourhood of $r$. It is contained in
$\mathcal W_{\tau(r)}$, and hence, by invariance, in
$\mathcal W_{h\tau(r)}$ for every $h\in G_r$. It also retains
(\ref{equation-equivariant-cech-zero-shrinking}) and
(\ref{equation-equivariant-cech-zero-local}).

Define a fixed-point-free open $G$-covering $\mathcal N$, indexed by
$G\times R$, by
$$
\mathcal N_{g,r}=gN_r.
$$
The equivariant map
$$
c:G\times R\longrightarrow G\times A,
\qquad c(g,r)=g\tau(r)
$$
is a refinement function from $\mathcal N$ to $\mathcal V$.
We claim that $c^*f=0$ as a cochain. Consider a tuple
$(g_0,r_0),\ldots,(g_n,r_n)$ with nonempty intersection and put
$d_a=g_0^{-1}g_a$. Then
$$
N_{r_0}\cap d_1N_{r_1}\cap\cdots\cap d_nN_{r_n}\neq\emptyset.
$$
Since
$d_aN_{r_a}\subset\mathcal W_{d_a\tau(r_a)}$,
(\ref{equation-equivariant-cech-zero-shrinking}) for $N_{r_0}$ gives
$$
N_{r_0}\subset
\mathcal V_{\tau(r_0)}\cap
\mathcal V_{d_1\tau(r_1)}\cap\cdots\cap
\mathcal V_{d_n\tau(r_n)}.
$$
Equation (\ref{equation-equivariant-cech-zero-local}) says that the
corresponding component of $f$ is zero on $N_{r_0}$, and invariance of
$f$ gives the same conclusion before applying $g_0^{-1}$. Thus the
component of $c^*f$ on the original intersection is zero. Components on
empty intersections are zero by $\mathcal P(\emptyset)=0$, proving the
claim.

The class $\xi$ consequently becomes zero after the refinement
$\mathcal N\to\mathcal V\to\mathcal U$. It is therefore zero in the
colimit of Definition
\ref{definition-equivariant-cech-cohomology-space}.
\end{proof}

\begin{lemma}[Equivariant \v Cech comparison]
\label{lemma-equivariant-cech-comparison-discontinuous}
\begin{reference}
\cite[Chapter V, \S 5.5, Theorem 5.5.6, p. 213]{Tohoku}
\end{reference}
\begin{history}
The cited statement uses the phrase ``discrete group of homeomorphisms'';
its proof invokes Lemma 5.5.5, whose action hypothesis is condition (D) of
Sheaves, Definition \ref{sheaves-definition-discontinuous-group-action}.
We record this proof-supported formulation. The Hausdorff hypothesis of
Lemma 5.5.5 follows here from condition (D) and paracompactness of the
orbit space, as observed below.
\end{history}
Let $G$ act discontinuously on a topological space $X$, let
$q:X\to Y=X/G$ be the orbit map, and assume that $Y$ is paracompact. For
every abelian $G$-sheaf $\mathcal A$ on $X$, the canonical maps
\begin{equation}
\label{equation-equivariant-cech-comparison-discontinuous}
\check H^n(X;G,\mathcal A)
\longrightarrow H^n(X;G,\mathcal A)
\end{equation}
are isomorphisms for every $n\geq 0$.
\end{lemma}

\begin{proof}
The space $Y$ is Hausdorff by the definition of paracompactness. We first
observe that $X$ is Hausdorff as well. Two points in distinct $G$-orbits
are separated by the inverse images of disjoint neighbourhoods in $Y$.
If $x' = gx\neq x$, then $g\notin G_x$, and Sheaves, Lemma
\ref{sheaves-lemma-discontinuous-group-action-neighbourhoods} gives an
open neighbourhood $V$ of $x$ such that $V\cap gV=\emptyset$. Thus $V$
and $gV$ separate $x$ and $x'$.

For $r>0$, put
$$
\mathcal P_r=\underline H^r(\mathcal A),
\qquad
\mathcal P_r(U)=H^r(U,\mathcal A).
$$
We have $\mathcal P_r(\emptyset)=H^r(\emptyset,\mathcal A)=0$.
Regarding $\mathcal A$ as a sheaf of $\underline{\mathbf Z}_X$-modules,
Lemma \ref{lemma-kill-cohomology-class-on-covering} shows that the sheaf
associated to $\mathcal P_r$ is zero. Hence Lemma
\ref{lemma-equivariant-cech-zero-associated-sheaf} gives
\begin{equation}
\label{equation-equivariant-cech-comparison-positive}
\check H^p(X;G,\mathcal P_r)=0
\quad\text{for }p>0,\ r>0.
\end{equation}

For $p=0$, Lemma
\ref{lemma-equivariant-cech-zero-associated-sheaf-degree-zero} gives
\begin{equation}
\label{equation-equivariant-cech-comparison-zero}
\check H^0(X;G,\mathcal P_r)=0
\quad\text{for }r>0.
\end{equation}

Apply the global spectral sequence of Lemma
\ref{lemma-global-equivariant-cech-spectral-sequence}. Equations
(\ref{equation-equivariant-cech-comparison-positive}) and
(\ref{equation-equivariant-cech-comparison-zero}) show that its
$E_2$-page is zero in every row $q>0$. The row $q=0$ is
$$
E_2^{p,0}=\check H^p(X;G,\mathcal A),
$$
because $\underline H^0(\mathcal A)=\mathcal A$. The spectral sequence
therefore degenerates, and its bottom edge maps are exactly
(\ref{equation-equivariant-cech-comparison-discontinuous}).
\end{proof}

\begin{lemma}[Low-degree equivariant \v Cech comparison]
\label{lemma-equivariant-cech-comparison-low-degree}
\begin{reference}
\cite[Chapter V, \S 5.5, Remarks II, p. 213]{Tohoku}
\end{reference}
\begin{history}
The official print has ${E_2^{0,n}=0}$, which is the relevant vanishing
for the indexing in (5.5.15). The frozen intake layer transcribed the
first index as $1$; we retain the authority's index and prove the
vanishing directly below.
\end{history}
Let $G$ act on a topological space $X$ and let $\mathcal A$ be an abelian
$G$-sheaf. With no further hypotheses, the canonical edge maps give a
functorial isomorphism
$$
\check H^1(X;G,\mathcal A)
\longrightarrow H^1(X;G,\mathcal A)
$$
and a functorial injection
$$
\check H^2(X;G,\mathcal A)
\longrightarrow H^2(X;G,\mathcal A).
$$
\end{lemma}

\begin{proof}
For $r>0$, let $\mathcal P_r$ be the presheaf
$U\mapsto H^r(U,\mathcal A)$. Regarding $\mathcal A$ as a sheaf of
$\underline{\mathbf Z}_X$-modules, Lemma
\ref{lemma-kill-cohomology-class-on-covering} shows that the sheaf
associated to $\mathcal P_r$ is zero. Hence Lemma
\ref{lemma-equivariant-cech-zero-associated-sheaf-degree-zero} gives
$$
E_2^{0,r}=\check H^0(X;G,\mathcal P_r)=0
\quad\text{for every }r>0
$$
in the spectral sequence of Lemma
\ref{lemma-global-equivariant-cech-spectral-sequence}. Its five-term
exact sequence begins
$$
0\longrightarrow E_2^{1,0}\longrightarrow H^1(X;G,\mathcal A)
\longrightarrow E_2^{0,1}\longrightarrow E_2^{2,0}
\longrightarrow H^2(X;G,\mathcal A).
$$
Here $E_2^{p,0}=\check H^p(X;G,\mathcal A)$ because
$\underline H^0(\mathcal A)=\mathcal A$. Thus the outer edge maps
displayed here are the canonical maps in the statement. Since
$E_2^{0,1}=0$, the first is an isomorphism and the second is injective.
\end{proof}


\begin{lemma}[Global equivariant \v Cech edge range]
\label{lemma-global-equivariant-cech-edge-range}
\begin{reference}
\cite[Chapter III, \S 3.8, paragraph following
Corollary 3.8.1(3), p. 176]{Tohoku}
\end{reference}
\begin{history}
The cited source states the case where $G$ is trivial. The formulation below
uses the global equivariant spectral sequence and contains that statement.
\end{history}
Let $G$ act on a topological space $X$, let $\mathcal A$ be an abelian
$G$-sheaf, and let $n \geq 1$. Assume that
$$
\check H^p(X;G,\underline H^q(\mathcal A))=0
$$
for every $p \geq 0$ and every $q$ with $0 < q < n$. Then the canonical
edge maps
$$
\check H^i(X;G,\mathcal A)
\longrightarrow
H^i(X;G,\mathcal A)
$$
are isomorphisms for $0 \leq i \leq n$ and are injective for $i=n+1$.
\end{lemma}

\begin{proof}
Apply the spectral sequence of Lemma
\ref{lemma-global-equivariant-cech-spectral-sequence}. For $q>0$, regard
$\mathcal A$ as a sheaf of $\underline{\mathbf Z}_X$-modules. Lemma
\ref{lemma-kill-cohomology-class-on-covering} shows that the sheaf
associated to the presheaf $\underline H^q(\mathcal A)$ is zero. Hence
Lemma \ref{lemma-equivariant-cech-zero-associated-sheaf-degree-zero}
gives
$$
E_2^{0,q}=\check H^0(X;G,\underline H^q(\mathcal A))=0
\quad\text{for every }q>0.
$$

For $r\geq 2$, a differential leaving $E_r^{i,0}$ has target
$E_r^{i+r,1-r}=0$. A differential entering it has source
$E_r^{i-r,r-1}$. Suppose $0\leq i\leq n+1$ and this source lies in the
first quadrant. If $r\leq n$, then $0<r-1<n$, so the source is zero by
the hypothesis. The only remaining possibility is
$i=r=n+1$, whose source is $E_{n+1}^{0,n}=0$ by the preceding
paragraph. Consequently
$$
E_2^{i,0}=E_\infty^{i,0}
\quad\text{for }0\leq i\leq n+1.
$$

Let $0\leq i\leq n$. Every other graded piece in total degree $i$ has
the form $E_\infty^{p,q}$ with $q>0$. It vanishes by the hypothesis
unless $p=0$ and $q=i=n$, and that remaining term vanishes by
$E_2^{0,n}=0$. Thus the finite filtration on
$H^i(X;G,\mathcal A)$ is concentrated in $E_\infty^{i,0}$, and the
bottom edge map is an isomorphism. In total degree $n+1$, the stable
term $E_\infty^{n+1,0}$ is the last filtration step and therefore
injects into the abutment. This proves the final assertion.
\end{proof}


\begin{definition}
\label{definition-equivariant-torsor}
Let $G$ act on a topological space $X$, and let $\mathcal A$ be a
$G$-sheaf of possibly noncommutative groups. A
{\it $G$-equivariant $\mathcal A$-torsor} is an $\mathcal A$-torsor
$\mathcal F$ in the sense of Definition \ref{definition-torsor},
together with the structure of a $G$-sheaf, such that
$$
\rho_g^{\mathcal F}(a \cdot s) =
\rho_g^{\mathcal A}(a) \cdot \rho_g^{\mathcal F}(s)
$$
for every $g \in G$ and all local sections $a$ and $s$ for which this
formula is defined. A morphism of $G$-equivariant $\mathcal A$-torsors
is both $\mathcal A$-equivariant and $G$-equivariant.
\end{definition}

\begin{definition}
\label{definition-nonabelian-equivariant-cech-h1-cover}
Let $G$ act on a topological space $X$, and let $\mathcal A$ be a
$G$-sheaf of possibly noncommutative groups. Let
$\mathcal U = (U_i)_{i \in I}$ be a fixed-point-free open $G$-covering
of $X$. An {\it equivariant $1$-cocycle} with values in $\mathcal A$ is
a family
$$
a_{ij} \in \mathcal A(U_i \cap U_j), \qquad i, j \in I,
$$
such that, after restriction to the indicated intersections,
$$
a_{ik} = a_{jk}a_{ij}
\quad\text{and}\quad
\rho_g^{\mathcal A}(a_{ij}) = a_{gi,gj}.
$$
Two equivariant $1$-cocycles $a$ and $a'$ are {\it equivalent} if there
is a family
$$
b_i \in \mathcal A(U_i),
\qquad
\rho_g^{\mathcal A}(b_i) = b_{gi},
$$
such that
$$
a'_{ij} = b_j a_{ij}b_i^{-1}.
$$
We denote the resulting pointed set by
$\check H^1_{\mathrm{nab}}(\mathcal U;G,\mathcal A)$. Its distinguished
point is the class of the identity cocycle.
\end{definition}

\begin{lemma}
\label{lemma-nonabelian-equivariant-cech-refinements}
Let $G$ act on a topological space $X$, and let $\mathcal A$ be a
$G$-sheaf of possibly noncommutative groups. Let
$\mathcal U=(U_i)_{i\in I}$ and $\mathcal V=(V_j)_{j\in J}$ be
fixed-point-free open $G$-coverings of $X$. If $\mathcal V$ refines
$\mathcal U$ as an ordinary open covering, then there is a canonical map
of pointed sets
$$
\check H^1_{\mathrm{nab}}(\mathcal U;G,\mathcal A)
\longrightarrow
\check H^1_{\mathrm{nab}}(\mathcal V;G,\mathcal A).
$$
These maps are transitive under further refinement.
\end{lemma}

\begin{proof}
There is a $G$-equivariant refinement function $c:J\to I$. Indeed,
choose $c(j)$ with $V_j\subset U_{c(j)}$ for one representative of
every $G$-orbit of $J$ and extend equivariantly; this is possible because
the action on $J$ is free. Pullback along $c$ is given by
$$
(c^*a)_{jk}=
a_{c(j)c(k)}|_{V_j\cap V_k}.
$$
It preserves cocycles, equivalence, and the distinguished point.

If $c'$ is another equivariant refinement function, then
$$
b_j=a_{c(j)c'(j)}|_{V_j}
$$
is an invariant family. The cocycle identities give
$$
(c'^*a)_{jk}=b_k(c^*a)_{jk}b_j^{-1}.
$$
Thus $c$ and $c'$ induce the same map on equivalence classes.
Composing equivariant refinement functions shows transitivity.
\end{proof}

\begin{definition}
\label{definition-nonabelian-equivariant-cech-h1}
Let $G$ act on a topological space $X$, and let $\mathcal A$ be a
$G$-sheaf of possibly noncommutative groups. Set
$$
\check H^1_{\mathrm{nab}}(X;G,\mathcal A)=
\colim_{\mathcal U}
\check H^1_{\mathrm{nab}}(\mathcal U;G,\mathcal A),
$$
where $\mathcal U$ runs through the fixed-point-free open
$G$-coverings and the transition maps are those of Lemma
\ref{lemma-nonabelian-equivariant-cech-refinements}. The intersection
covering of two such coverings is again fixed-point-free, and Lemma
\ref{lemma-equivariant-cech-free-coverings-cofinal} supplies a set-sized
cofinal system, so this is a well-defined pointed set.
\end{definition}

\begin{lemma}
\label{lemma-nonabelian-equivariant-cech-h1-torsors}
\begin{reference}
\cite[Chapter V, \S 5.5, Remarks II, p. 213]{Tohoku}
\end{reference}
\begin{history}
The cited source calls the objects below $G$-$A$-bundles with structural
sheaf $A$. We use the modern term equivariant torsor. The final sentences
of the cited paragraph observe that earlier sheaf-theoretic developments
can be transported to these objects and propose a broader noncommutative
homological algebra. The result below proves the precise degree-one
classification, not those programmatic proposals.
\end{history}
Let $G$ act on a topological space $X$, and let $\mathcal A$ be a
$G$-sheaf of possibly noncommutative groups. There is a canonical
bijection of pointed sets between
$\check H^1_{\mathrm{nab}}(X;G,\mathcal A)$ and the set of isomorphism
classes of $G$-equivariant $\mathcal A$-torsors.

If $\mathcal A$ is abelian, there are canonical bijections
$$
\left\{
\begin{matrix}
\text{isomorphism classes of} \\
G\text{-equivariant }\mathcal A\text{-torsors}
\end{matrix}
\right\}
\longleftrightarrow
\check H^1(X;G,\mathcal A)
\longleftrightarrow
H^1(X;G,\mathcal A).
$$
\end{lemma}

\begin{proof}
Let $\mathcal F$ be a $G$-equivariant $\mathcal A$-torsor. Choose an
ordinary open covering $(V_j)_{j\in J}$ and sections
$s_j\in\mathcal F(V_j)$. Put
$$
U_{(g,j)}=gV_j,
\qquad
s_{(g,j)}=\rho_g^{\mathcal F}(s_j),
$$
and let $G$ act on the index set $G\times J$ by left multiplication on
the first factor. This is a fixed-point-free open $G$-covering with an
equivariant family of local sections. There are unique sections
$$
a_{ij}\in\mathcal A(U_i\cap U_j)
$$
such that $a_{ij}\cdot s_i=s_j$. Simple transitivity and associativity
give
$$
a_{jk}a_{ij}\cdot s_i=a_{jk}\cdot s_j=s_k=a_{ik}\cdot s_i,
$$
and hence $a_{ik}=a_{jk}a_{ij}$. Compatibility of the two actions gives
$\rho_g^{\mathcal A}(a_{ij})=a_{gi,gj}$. Changing the equivariant local
sections to $b_i\cdot s_i$ changes the cocycle to
$b_ja_{ij}b_i^{-1}$. Common refinements show that the resulting class
does not depend on any of the choices, and an equivariant isomorphism of
torsors does not change it.

Conversely, let $a=(a_{ij})$ be an equivariant cocycle on $\mathcal U$.
On every $U_i$ take the trivial left $\mathcal A|_{U_i}$-torsor. On an
overlap glue the underlying sheaves of sets by the
$\mathcal A$-equivariant maps
$$
\varphi_{ij}(h)=h a_{ij}^{-1}.
$$
The cocycle condition gives
$$
\varphi_{jk}(\varphi_{ij}(h))
=h a_{ij}^{-1}a_{jk}^{-1}
=h(a_{jk}a_{ij})^{-1}
=h a_{ik}^{-1}.
$$
Thus Sheaves, Lemma \ref{sheaves-lemma-glue-sheaves} produces an
$\mathcal A$-torsor $\mathcal F_a$. The maps
$$
h\longmapsto\rho_g^{\mathcal A}(h)
$$
from the copy over $U_i$ to the copy over $U_{gi}$ are compatible with
the glueing maps by invariance of $a$. They therefore give
$\mathcal F_a$ the structure of a $G$-sheaf compatible with the left
$\mathcal A$-action.

If $a'_{ij}=b_ja_{ij}b_i^{-1}$, then the local maps
$$
h\longmapsto h b_i^{-1}
$$
glue to a $G$-equivariant $\mathcal A$-torsor isomorphism
$\mathcal F_a\to\mathcal F_{a'}$. Pullback along a refinement likewise
does not change the isomorphism class of the glued torsor. Starting with
a cocycle, the local sections corresponding to the identity recover
that cocycle. Starting with a torsor, the maps
$$
\mathcal A|_{U_i}\longrightarrow\mathcal F|_{U_i},
\qquad h\longmapsto h\cdot s_i,
$$
glue to an equivariant isomorphism with the torsor constructed from its
cocycle. The two constructions are therefore mutually inverse, and the
identity cocycle corresponds to the trivial equivariant torsor
$\mathcal A$.

Finally suppose that $\mathcal A$ is abelian and write its operation
additively. For every fixed-point-free covering the conditions above
become
$$
a_{jk}-a_{ik}+a_{ij}=0
\quad\text{and}\quad
a'_{ij}=a_{ij}+b_j-b_i.
$$
Consequently
$$
\check H^1_{\mathrm{nab}}(\mathcal U;G,\mathcal A)
\cong
H^1\left(\check{\mathcal C}^\bullet(
\mathcal U,\mathcal A)^G\right).
$$
Lemma \ref{lemma-equivariant-cech-free-index} identifies the latter
group with $\check H^1(\mathcal U;G,\mathcal A)$. The identifications
commute with refinement maps. Lemma
\ref{lemma-equivariant-cech-free-coverings-cofinal} gives the first
bijection in the statement after passing to the colimit, and Lemma
\ref{lemma-equivariant-cech-comparison-low-degree} gives the second.
\end{proof}


\begin{lemma}
\label{lemma-equivariant-internal-hom}
\begin{reference}
\cite[Chapter V, \S 5.6, pp. 213--214 and 216]{Tohoku}
\end{reference}
Let $(X, \mathcal O)$ be a $G$-ringed space and let $\mathcal A$ and
$\mathcal B$ be $G$-$\mathcal O$-modules. The internal hom
$$
\mathcal H=\SheafHom_{\mathcal O}(\mathcal A,\mathcal B)
$$
has a canonical structure of a $G$-$\mathcal O$-module. Explicitly, if
$g\in G$, $U\subset X$ is open, and $\varphi\in\mathcal H(U)$, then
$\rho_g^{\mathcal H}(U)(\varphi)$ is the morphism over $gU$ whose component
over an open $V\subset gU$ is
$$
\mathcal A(V)
\xrightarrow{\rho_{g^{-1}}^{\mathcal A}(V)}
\mathcal A(g^{-1}V)
\xrightarrow{\varphi_{g^{-1}V}}
\mathcal B(g^{-1}V)
\xrightarrow{\rho_g^{\mathcal B}(g^{-1}V)}
\mathcal B(V).
$$
In pointwise notation this is the conjugation formula
$$
(g\varphi)(a)=g\varphi(g^{-1}a).
$$
For every $G$-stable open $U\subset X$ set
$$
\Hom_{\mathcal O,G}(U;\mathcal A,\mathcal B)=
\Hom_{\textit{Mod}_G(\mathcal O|_U)}
(\mathcal A|_U,\mathcal B|_U).
$$
Then there are canonical equalities
$$
\Hom_{\mathcal O,G}(U;\mathcal A,\mathcal B)=
\Gamma(U,\mathcal H)^G.
$$
In particular,
$$
\begin{aligned}
\Hom_{\textit{Mod}_G(\mathcal O)}(\mathcal A,\mathcal B)
&=\Hom_{\mathcal O,G}(X;\mathcal A,\mathcal B) \\
&=\Gamma_X^G(\mathcal H)
=\Gamma(X,\mathcal H)^G.
\end{aligned}
$$
\end{lemma}

\begin{proof}
The displayed composite is $\mathcal O(V)$-linear by the semilinearity of
the transports on $\mathcal A$ and $\mathcal B$. Compatibility with
restriction maps is immediate. The identities
$$
\rho_{gh}=\rho_g\circ\rho_h
\quad\text{and}\quad
\rho_{g^{-1}}\circ\rho_g=\text{id}
$$
for the transports on $\mathcal A$ and $\mathcal B$ show that the displayed
rule defines a $G$-action on $\mathcal H$. If $a\in\mathcal O(U)$, then
semilinearity on $\mathcal B$ gives
$$
\rho_g^{\mathcal H}(U)(a\varphi)=
\rho_g^{\mathcal O}(U)(a)\rho_g^{\mathcal H}(U)(\varphi).
$$
Thus $\mathcal H$ is a $G$-$\mathcal O$-module in the sense of Modules,
Definition \ref{modules-definition-G-O-module}.

Suppose that $U$ is $G$-stable. By the definition of internal hom, a section
$\varphi\in\Gamma(U,\mathcal H)$ is an $\mathcal O|_U$-linear morphism
$\mathcal A|_U\to\mathcal B|_U$. The section is fixed by $g$ exactly when
$$
\rho_g^{\mathcal B}(W)\circ\varphi_W=
\varphi_{gW}\circ\rho_g^{\mathcal A}(W)
$$
for every open $W\subset U$. This is precisely the condition that
$\varphi$ be $G$-equivariant. This proves the first equality. Taking
$U=X$ and using Definition
\ref{definition-orbit-space-equivariant-cohomology} proves the remaining
equalities.
\end{proof}


\begin{lemma}
\label{lemma-equivariant-hom-functors}
\begin{reference}
\cite[Chapter V, \S 5.6, p. 214]{Tohoku}
\end{reference}
Let $(X, \mathcal O)$ be a $G$-ringed space, let
$q : X \to Y = X/G$ be the orbit map, and put
$$
\mathcal O' = q_*^G\mathcal O.
$$
Fix a $G$-$\mathcal O$-module $\mathcal A$. For a
$G$-$\mathcal O$-module $\mathcal B$ put
$$
\begin{aligned}
\mathbf h_{\mathcal O,\mathcal A}(\mathcal B)
&=\SheafHom_{\mathcal O}(\mathcal A,\mathcal B), \\
h_{\mathcal O,\mathcal A}(\mathcal B)
&=\Hom_{\mathcal O}(\mathcal A,\mathcal B)
=\Gamma\left(X,
\SheafHom_{\mathcal O}(\mathcal A,\mathcal B)\right), \\
\mathbf h_{\mathcal O,G,\mathcal A}(\mathcal B)
&=q_*^G\SheafHom_{\mathcal O}(\mathcal A,\mathcal B), \\
h_{\mathcal O,G,\mathcal A}(\mathcal B)
&=\Hom_{\textit{Mod}_G(\mathcal O)}(\mathcal A,\mathcal B)
=\Hom_{\mathcal O,G}(X;\mathcal A,\mathcal B).
\end{aligned}
$$
These assignments define additive left exact functors, with values,
respectively, in $G$-$\mathcal O$-modules, abelian $G$-modules,
$\mathcal O'$-modules, and abelian groups. There are equalities of functors
$$
\mathbf h_{\mathcal O,G,\mathcal A}
=q_*^G\circ\mathbf h_{\mathcal O,\mathcal A}
$$
and, after forgetting the indicated module structures,
$$
\begin{aligned}
h_{\mathcal O,G,\mathcal A}
&=\Gamma(Y,-)\circ\mathbf h_{\mathcal O,G,\mathcal A} \\
&=(-)^G\circ h_{\mathcal O,\mathcal A} \\
&=\Gamma_X^G\circ\mathbf h_{\mathcal O,\mathcal A}.
\end{aligned}
$$
\end{lemma}

\begin{proof}
The first functor is well defined by Lemma
\ref{lemma-equivariant-internal-hom}. It is additive and left exact by
Modules, Lemma \ref{modules-lemma-internal-hom-exact}: the conjugation
action is functorial in $\mathcal B$, and kernels in
$\textit{Mod}_G(\mathcal O)$ are computed on the underlying
$\mathcal O$-modules by Modules, Lemma
\ref{modules-lemma-AB5-AB3-star-generator-G-O-modules}. The other three
functors are obtained from it by applying global sections, invariant direct
image, and invariants after global sections. These are additive left exact
functors; for invariant direct image on modules, see Modules, Lemma
\ref{modules-lemma-invariant-direct-image-G-O-modules}.

The first equality of functors is a definition. If
$\mathcal H=\SheafHom_{\mathcal O}(\mathcal A,\mathcal B)$, then
$$
\Gamma(Y,q_*^G\mathcal H)
=\Gamma(X,\mathcal H)^G
=\Hom_{\textit{Mod}_G(\mathcal O)}(\mathcal A,\mathcal B).
$$
The first equality follows from the sectionwise definition of invariant
direct image, and the second is Lemma
\ref{lemma-equivariant-internal-hom}. This gives all three expressions for
$h_{\mathcal O,G,\mathcal A}$; the last uses Definition
\ref{definition-orbit-space-equivariant-cohomology}.
\end{proof}


\begin{definition}
\label{definition-equivariant-ext}
\begin{reference}
\cite[Chapter V, \S 5.6, p. 214]{Tohoku}
\end{reference}
\begin{history}
The printed formula (5.6.7) includes $X$ among the arguments of the
sheaf-valued Ext functor. The corrected source removes it: unlike the global
Ext group in (5.6.6), this object is a sheaf on $Y$, as is also clear from
the restriction formula (5.6.8).
\end{history}
In the situation of Lemma \ref{lemma-equivariant-hom-functors}, for
$n \in \mathbf Z$ and $G$-$\mathcal O$-modules $\mathcal A,\mathcal B$ set
$$
\Ext^n_{\mathcal O,G}(X;\mathcal A,\mathcal B)
=
R^nh_{\mathcal O,G,\mathcal A}(\mathcal B)
$$
and
$$
\SheafExt^n_{\mathcal O,G}(\mathcal A,\mathcal B)
=
R^n\mathbf h_{\mathcal O,G,\mathcal A}(\mathcal B).
$$
Thus the first expression is an abelian group and the second is an
$\mathcal O'$-module on $Y$, where $\mathcal O'=q_*^G\mathcal O$.
These right derived functors exist because $\textit{Mod}_G(\mathcal O)$
has enough injectives by Injectives, Lemma
\ref{injectives-lemma-equivariant-modules-enough-injectives}.
\end{definition}

\begin{lemma}
\label{lemma-equivariant-ext-properties}
\begin{reference}
\cite[Chapter V, \S 5.6, p. 214]{Tohoku}
\end{reference}
The assignments of Definition \ref{definition-equivariant-ext} are additive
bifunctors
$$
\textit{Mod}_G(\mathcal O)^{opp}\times
\textit{Mod}_G(\mathcal O)
\longrightarrow
\textit{Ab}
$$
and
$$
\textit{Mod}_G(\mathcal O)^{opp}\times
\textit{Mod}_G(\mathcal O)
\longrightarrow
\textit{Mod}(\mathcal O'),
$$
respectively, and they vanish in negative degrees. For fixed $\mathcal A$,
the functors in degrees $n\geq 0$ form universal cohomological delta functors
in $\mathcal B$, with
$$
\Ext^0_{\mathcal O,G}(X;\mathcal A,\mathcal B)
=
\Hom_{\textit{Mod}_G(\mathcal O)}(\mathcal A,\mathcal B)
=
\Hom_{\mathcal O,G}(X;\mathcal A,\mathcal B)
$$
and
$$
\SheafExt^0_{\mathcal O,G}(\mathcal A,\mathcal B)
=
q_*^G\SheafHom_{\mathcal O}(\mathcal A,\mathcal B)
=
\mathbf h_{\mathcal O,G,\mathcal A}(\mathcal B).
$$
There are canonical isomorphisms
$$
\Ext^n_{\mathcal O,G}(X;\mathcal A,\mathcal B)
\longrightarrow
\Ext^n_{\textit{Mod}_G(\mathcal O)}(\mathcal A,\mathcal B)
$$
for every $n\geq 0$. For fixed $\mathcal B$, both families in degrees
$n\geq 0$ are contravariant cohomological delta functors in $\mathcal A$.
\end{lemma}

\begin{proof}
For fixed $\mathcal A$, Lemma \ref{lemma-equivariant-hom-functors} gives
additive left exact functors, and the source category has enough injectives.
Derived Categories, Lemmas \ref{derived-lemma-left-exact-higher-derived} and
\ref{derived-lemma-right-derived-delta-functor} give the assertions about
negative degrees, degree zero, and universal cohomological delta functors in
$\mathcal B$.

Let $\mathcal B[0]\to\mathcal I^\bullet$ be an injective resolution in
$\textit{Mod}_G(\mathcal O)$. Then
$$
\Ext^n_{\mathcal O,G}(X;\mathcal A,\mathcal B)
=
H^n\left(
\Hom_{\textit{Mod}_G(\mathcal O)}
(\mathcal A,\mathcal I^\bullet)
\right).
$$
Derived Categories, Lemma \ref{derived-lemma-compute-ext-resolutions}
identifies the right hand side with
$$
\Ext^n_{\textit{Mod}_G(\mathcal O)}(\mathcal A,\mathcal B),
$$
as asserted.

It remains to check exactness in the first variable when the second variable
is injective. For the group-valued functor this is the defining property of
an injective object. For the sheaf-valued functor, let $\mathcal I$ be
injective, let $U\subset Y$ be open, and put $W=q^{-1}(U)$. Since $W$ is
$G$-stable, extension by zero along $W\subset X$ carries a natural
equivariant structure and is left adjoint to restriction by Sheaves, Lemma
\ref{sheaves-lemma-j-shriek-modules}. It is exact on underlying abelian
sheaves by
Modules, Lemma \ref{modules-lemma-j-shriek-exact}. Hence
$\mathcal I|_W$ is injective by Homology, Lemma
\ref{homology-lemma-adjoint-preserve-injectives}. Lemma
\ref{lemma-equivariant-internal-hom} gives
$$
\Gamma\left(
U,\mathbf h_{\mathcal O,G,\mathcal A}(\mathcal I)
\right)
=
\Hom_{\textit{Mod}_G(\mathcal O|_W)}
(\mathcal A|_W,\mathcal I|_W).
$$
Consequently $\mathbf h_{\mathcal O,G,-}(\mathcal I)$ is exact: a short
exact sequence in its argument gives an exact sequence on sections over
every open $U\subset Y$. Applying Derived Categories, Lemma
\ref{derived-lemma-partial-right-derived-bifunctor} to the group-valued and
sheaf-valued bifunctors proves bifunctoriality and supplies the contravariant
cohomological delta functors in $\mathcal A$.
\end{proof}

\begin{lemma}
\label{lemma-localize-equivariant-ext}
\begin{reference}
\cite[Chapter V, \S 5.6, formula (5.6.8), p. 214]{Tohoku}
\end{reference}
Let $U\subset Y$ be open, put $X_U=q^{-1}(U)$, and write
$\mathcal O_U=\mathcal O|_{X_U}$ and $q_U:X_U\to U$. Under the canonical
identification
$$
\mathcal O'|_U=q_{U,*}^G\mathcal O_U,
$$
there are canonical isomorphisms of $\mathcal O'|_U$-modules
$$
\SheafExt^n_{\mathcal O,G}(\mathcal A,\mathcal B)|_U
\longrightarrow
\SheafExt^n_{\mathcal O_U,G}
(\mathcal A|_{X_U},\mathcal B|_{X_U})
$$
for every $n\geq 0$.
\end{lemma}

\begin{proof}
Equivariant extension by zero along the open $G$-stable subspace
$X_U\subset X$ is an exact left adjoint to restriction, as in the proof of
Lemma \ref{lemma-equivariant-ext-properties}. Hence restriction to $X_U$
preserves injectives.

Choose an injective resolution $\mathcal B[0]\to\mathcal I^\bullet$ in
$\textit{Mod}_G(\mathcal O)$. Its restriction is an injective resolution of
$\mathcal B|_{X_U}$ in $\textit{Mod}_G(\mathcal O_U)$. The sectionwise
description in Lemma \ref{lemma-equivariant-internal-hom} gives canonical
isomorphisms
$$
\left(
q_*^G\SheafHom_{\mathcal O}(\mathcal A,\mathcal I^r)
\right)|_U
\longrightarrow
q_{U,*}^G\SheafHom_{\mathcal O_U}
(\mathcal A|_{X_U},\mathcal I^r|_{X_U})
$$
in every degree $r$. Restriction is exact, so taking cohomology proves the
result.
\end{proof}

\begin{lemma}
\label{lemma-equivariant-ext-slice}
\begin{reference}
\cite[Chapter V, \S 5.6, following formula (5.6.8), p. 214]{Tohoku}
\end{reference}
Let $H\subset G$ be a subgroup and let $V\subset X$ be an open $H$-stable
subset such that
$$
Z=\bigcup_{g\in G}gV
=
\coprod_{gH\in G/H}gV.
$$
Put $U=q(V)$ and let $p:V\to U=V/H$ be the quotient map. There is a
canonical identification
$$
p_*^H(\mathcal O|_V)=\mathcal O'|_U,
$$
and for every $n\geq 0$ there are canonical isomorphisms of
$\mathcal O'|_U$-modules
$$
\SheafExt^n_{\mathcal O,G}(\mathcal A,\mathcal B)|_U
\longrightarrow
\SheafExt^n_{\mathcal O|_V,H}
(\mathcal A|_V,\mathcal B|_V).
$$

In particular, if the underlying action of $G$ on $X$ is discontinuous,
then for every $x\in X$, with $H=G_x$, there are neighbourhoods
$V\subset X$ of $x$ and $U=q(V)\subset Y$ as above. The group $H$ is
finite, and on $U$ the equivariant Ext sheaves are computed by the finite
group $G_x$ acting on $V$.
\end{lemma}

\begin{proof}
The subset $Z=q^{-1}(U)$ is open and $G$-stable. By Lemma
\ref{lemma-localize-equivariant-ext} we may replace $X$ by $Z$. The induced
map $V/H\to U$ is a homeomorphism, as in the proof of Lemma
\ref{lemma-higher-invariant-direct-image-slice}.

Restriction to $V$ gives an exact equivalence
$$
\textit{Mod}_G(\mathcal O|_Z)
\longrightarrow
\textit{Mod}_H(\mathcal O|_V).
$$
A quasi-inverse is obtained by placing the translates of an
$H$-$\mathcal O|_V$-module on the open and closed components $gV$ and
transporting the module structures by the $G$-action. On espaces
\'etal\'es this is the construction
$$
E\longmapsto G\mathbin{\times^H}E.
$$
The disjointness hypothesis shows that the two composites are naturally
isomorphic to the identity. The construction respects the semilinear module
structures, so this is an equivalence of the displayed module categories.
In particular, it preserves injective objects.

For $W\subset U$ open, restriction gives functorial isomorphisms
$$
\begin{aligned}
&\Hom_{\textit{Mod}_G(\mathcal O|_{q^{-1}(W)})}
(\mathcal A|_{q^{-1}(W)},\mathcal B|_{q^{-1}(W)})\\
&\qquad\longrightarrow
\Hom_{\textit{Mod}_H(\mathcal O|_{p^{-1}(W)})}
(\mathcal A|_{p^{-1}(W)},\mathcal B|_{p^{-1}(W)}).
\end{aligned}
$$
By Lemma \ref{lemma-equivariant-internal-hom}, these isomorphisms identify
the degree-zero sheaf-valued Hom functors on $U$. Applied to the ring itself,
the same sectionwise identification gives
$p_*^H(\mathcal O|_V)=\mathcal O'|_U$.

Choose an injective resolution after replacing $X$ by $Z$. Restriction to
$V$ turns it into an injective resolution in
$\textit{Mod}_H(\mathcal O|_V)$, and the preceding degree-zero
identification gives an isomorphism between the resulting complexes of
sheaves on $U$. Taking cohomology proves the displayed Ext isomorphisms.

For a discontinuous action, take $H=G_x$ and choose $V$ using Sheaves,
Lemma \ref{sheaves-lemma-discontinuous-group-action-neighbourhoods}. That
lemma gives the required disjoint decomposition, and $H=G_x$ is finite by
definition.
\end{proof}


\begin{lemma}[Equivariant Hom and injective stalks]
\label{lemma-equivariant-hom-injective-stalk}
\begin{reference}
\cite[Chapter V, \S 5.6, formula (5.6.9), pp. 214--215]{Tohoku}
\end{reference}
\begin{history}
The parenthetical description following (5.6.9) in the printed source says
that $\mathcal U_x$ is generated by $\mathcal R_x$ and
$\mathbf Z[G]$. The definition on page 197 uses $G_x$, and only $G_x$
acts on a stalk at $x$. Thus the corrected source has
$\mathbf Z[G_x]$, as in Injectives, Definition
\ref{injectives-definition-stabilizer-skew-group-ring}.
\end{history}
For this lemma and the next one only, rings are not required to be
commutative. Let $G$ act discontinuously on $X$, and let $\mathcal R$ be a
$G$-sheaf of associative rings with $1$ whose restriction maps and
$G$-action preserve $1$. All modules are unital left modules. Let
$$
q:X\longrightarrow Y=X/G,
\qquad x\in X,
\qquad y=q(x),
\qquad H=G_x,
$$
and let $\mathcal U_x$ be the skew group ring of $\mathcal R_x$ and $H$.
For $G$-$\mathcal R$-modules $\mathcal A$ and $\mathcal B$, write
$\mathcal H\!om_{\mathcal R}(\mathcal A,\mathcal B)$ for the sheaf of
abelian groups of left $\mathcal R$-linear maps. It has the conjugation
$G$-action, but no module structure on this internal Hom is asserted.
\begin{enumerate}
\item If $\mathcal A$ is locally of finite presentation as a left
$\mathcal R$-module on a neighbourhood of $x$, then the natural map
$$
\left(q_*^G\mathcal H\!om_{\mathcal R}
(\mathcal A,\mathcal B)\right)_y
\longrightarrow
\mathop{\rm Hom}\nolimits_{\mathcal U_x}
(\mathcal A_x,\mathcal B_x)
$$
is an isomorphism of abelian groups.
\item Assume that on some neighbourhood $N$ of $x$ the following left
coherence condition holds: for every open $T\subset N$ and every morphism
$$
(\mathcal R|_T)^{\oplus p}\longrightarrow
(\mathcal R|_T)^{\oplus r}
$$
of left $\mathcal R|_T$-modules, with $p,r$ finite, its image is locally of
finite presentation. If $\mathcal R_x$ is left Noetherian and $\mathcal I$
is injective in $\textit{Mod}_G(\mathcal R)$, then $\mathcal I_x$ is an
injective left $\mathcal U_x$-module.
\end{enumerate}
The hypothesis on images in (2) holds if the left regular
$\mathcal R|_N$-module is coherent. Moreover, a coherent left
$\mathcal R$-module is locally of finite presentation. Hence both parts
apply under the hypotheses of the source: $\mathcal R$ is coherent with
left Noetherian stalks and $\mathcal A$ is coherent as an underlying left
$\mathcal R$-module.
\end{lemma}

\begin{proof}
Put
$$
\mathcal H=\mathcal H\!om_{\mathcal R}(\mathcal A,\mathcal B).
$$
The conjugation construction in Lemma
\ref{lemma-equivariant-internal-hom} uses no commutativity when $\mathcal H$
is regarded only as an abelian sheaf. Sheaves, Lemma
\ref{sheaves-lemma-invariant-direct-image-stalk} gives
$$
(q_*^G\mathcal H)_y=\mathcal H_x^H.
$$
Choose, after shrinking about $x$, a finite presentation
$$
(\mathcal R|_N)^{\oplus r}\longrightarrow
(\mathcal R|_N)^{\oplus p}\longrightarrow
\mathcal A|_N\longrightarrow 0.
$$
As sheaves of abelian groups,
$\mathcal H|_N$ is the kernel of the corresponding map
$\mathcal B|_N^{\oplus p}\to\mathcal B|_N^{\oplus r}$. Exactness of stalks
therefore gives a natural isomorphism
$$
\mathcal H_x\longrightarrow
\mathop{\rm Hom}\nolimits_{\mathcal R_x}
(\mathcal A_x,\mathcal B_x).
$$
It is $H$-equivariant. A left $\mathcal R_x$-linear map between the two
stalks is fixed by conjugation precisely when it commutes with $H$, and this
is precisely $\mathcal U_x$-linearity. Taking $H$-fixed parts proves (1).

We prove (2). First, $\mathcal U_x$ is left Noetherian. Indeed, as a left
$\mathcal R_x$-module it is the finite direct sum
$$
\mathcal U_x=\bigoplus_{h\in H}\mathcal R_x[h].
$$
It is therefore a Noetherian left $\mathcal R_x$-module. Every left ideal
of $\mathcal U_x$ is an $\mathcal R_x$-submodule, so it has finitely many
$\mathcal R_x$-generators; the same elements generate it as a left
$\mathcal U_x$-ideal.

We use the left-handed Baer criterion proved in the proof of Injectives,
Lemma \ref{injectives-lemma-injective-left-module-stalks}. Let
$J\subset\mathcal U_x$ be a left ideal and let
$$
\varphi:J\longrightarrow\mathcal I_x
$$
be $\mathcal U_x$-linear. Choose left generators
$a_1,\ldots,a_n$ of $J$. By Sheaves, Lemma
\ref{sheaves-lemma-discontinuous-group-action-neighbourhoods}, after
shrinking inside $N$ there is an open $H$-invariant neighbourhood $V$ of
$x$ such that
$$
gV\cap V=\emptyset\quad\text{for }g\notin H.
$$
For an open $T\subset X$, put
$$
\mathcal L(T)=
\mathop{\rm Ind}\nolimits_G(j_{T!}\mathcal R|_T),
$$
as in Modules, Lemma \ref{modules-lemma-generator-G-O-modules}. Its
construction and adjunction use only left module structures. If $T\subset
V$ is an $H$-invariant neighbourhood of $x$, then only the summands indexed
by $H$ meet $x$, and there is a canonical isomorphism of left
$\mathcal U_x$-modules
$$
\mathcal L(T)_x
=
\bigoplus_{h\in H}\mathcal R_x[h]
=
\mathcal U_x.
$$
More explicitly, an element $m$ in the summand indexed by $h$ maps to
$\rho^{\mathcal R}_{h,x}(m)[h]$. Under this identification the $H$-action
on the stalk is left multiplication on the symbols $[h]$, with the
prescribed semilinear action on the coefficients.

Using $\mathcal L(V)_x=\mathcal U_x$, represent all the $a_i$ by sections
$\widetilde a_i\in\mathcal L(V)(W)$ on a common $H$-invariant open
neighbourhood $W\subset V$ of $x$. By the adjunction defining
$\mathcal L(W)$, these sections give a $G$-equivariant morphism
$$
f:\mathcal L(W)^{\oplus n}\longrightarrow\mathcal L(V).
$$
Let $\mathcal F$ be its image. On the stalk at $x$, the map $f$ is
$$
\mathcal U_x^{\oplus n}\longrightarrow\mathcal U_x,
\qquad
(u_1,\ldots,u_n)\longmapsto\sum u_i a_i.
$$
Indeed, on the $i$th copy the adjunction sends $1$ to $a_i$, and
$\mathcal U_x$-linearity forces $u_i$ to map to $u_i a_i$; thus the
orientation is right multiplication by the chosen generator, as required
for the left ideal generated by the $a_i$.
Consequently $\mathcal F_x=J$.

On $W$, the condition on $V$ gives isomorphisms of underlying left modules
$$
\mathcal L(W)|_W\cong(\mathcal R|_W)^{\oplus |H|},
\qquad
\mathcal L(V)|_W\cong(\mathcal R|_W)^{\oplus |H|}.
$$
Thus $\mathcal F|_W$ is the image of a morphism between finite free left
$\mathcal R|_W$-modules. The left coherence hypothesis, followed by one
more $H$-invariant shrinking about $x$, makes $\mathcal F$ locally of
finite presentation near $x$.

Part (1), applied to $\mathcal F$ and $\mathcal I$, shows that $\varphi$ is
represented over some open neighbourhood $Q\subset q(W)$ of $y$ by a
$G$-equivariant left $\mathcal R$-linear map
$$
v:\mathcal F|_{q^{-1}(Q)}\longrightarrow
\mathcal I|_{q^{-1}(Q)}.
$$
Put $Z=q^{-1}(Q)$ and let $j:Z\to X$ be the open immersion. Equivariant
extension by zero is exact and is left adjoint to restriction; the proofs
of Sheaves, Lemma \ref{sheaves-lemma-j-shriek-modules} and Modules, Lemma
\ref{modules-lemma-j-shriek-exact} again use only the left action. Hence
$$
j_!(\mathcal F|_Z)\longrightarrow j_!(\mathcal L(V)|_Z)
$$
is a monomorphism, and $v$ gives a map from its source to $\mathcal I$.
Injectivity of $\mathcal I$ extends this map to the target. On the stalk at
$x$ the resulting $\mathcal U_x$-linear map
$$
\mathcal U_x=\mathcal L(V)_x\longrightarrow\mathcal I_x
$$
extends $\varphi$. The left Baer criterion proves that $\mathcal I_x$ is
injective.

For the final assertion in the statement, finite free left modules are
coherent when the left regular module is coherent. The image of a morphism
from a finite free module is of finite type; as a finite type submodule of
the coherent finite free target it is locally of finite presentation.
Likewise, a coherent $\mathcal A$ is locally of finite presentation. These
observations use only left module structures.
\end{proof}


\begin{lemma}[Equivariant Ext at a stalk]
\label{lemma-equivariant-ext-stalk}
\begin{reference}
\cite[Chapter V, \S 5.6, formula (5.6.10), p. 215]{Tohoku}
\end{reference}
In the situation of Lemma \ref{lemma-equivariant-hom-injective-stalk},
assume the left coherence and left Noetherian hypotheses in part (2), and
assume that $\mathcal A$ is locally of finite presentation near $x$. Define
an abelian sheaf on $Y$ by
$$
\mathcal E^n_{\mathcal R,G}(\mathcal A,\mathcal B)
=
R^n\left(
\mathcal C\longmapsto
q_*^G\mathcal H\!om_{\mathcal R}(\mathcal A,\mathcal C)
\right)(\mathcal B).
$$
There are canonical isomorphisms
$$
\mathcal E^n_{\mathcal R,G}(\mathcal A,\mathcal B)_y
\longrightarrow
\mathop{\rm Ext}\nolimits^n_{\mathcal U_x}
(\mathcal A_x,\mathcal B_x)
$$
for every $n\geq 0$. If $\mathcal R$ is commutative, then
$\mathcal E^n_{\mathcal R,G}(\mathcal A,\mathcal B)$ is the underlying
abelian sheaf of
$\SheafExt^n_{\mathcal R,G}(\mathcal A,\mathcal B)$ from Definition
\ref{definition-equivariant-ext}.
\end{lemma}

\begin{proof}
The category $\textit{Mod}_G(\mathcal R)$ has enough injectives: the proof
of Injectives, Lemma
\ref{injectives-lemma-equivariant-modules-enough-injectives} uses only
unital left module structures. Choose an injective resolution
$$
\mathcal B[0]\longrightarrow\mathcal I^\bullet.
$$
Part (2) of Lemma \ref{lemma-equivariant-hom-injective-stalk} shows that
every $\mathcal I^r_x$ is injective over $\mathcal U_x$. Exactness of
stalks therefore makes
$$
\mathcal B_x[0]\longrightarrow\mathcal I^\bullet_x
$$
an injective resolution of $\mathcal B_x$.

Stalks of abelian sheaves are exact. Hence the left hand side in the
statement is the $n$th cohomology of the complex
$$
\left(
q_*^G\mathcal H\!om_{\mathcal R}
(\mathcal A,\mathcal I^\bullet)
\right)_y.
$$
Part (1) of the preceding lemma gives a canonical termwise isomorphism of
this complex with
$$
\mathop{\rm Hom}\nolimits_{\mathcal U_x}
(\mathcal A_x,\mathcal I^\bullet_x).
$$
Its $n$th cohomology is
$\mathop{\rm Ext}\nolimits^n_{\mathcal U_x}
(\mathcal A_x,\mathcal B_x)$, proving the comparison.

When $\mathcal R$ is commutative, the functor inside the displayed right
derived functor is the underlying abelian-sheaf-valued functor of
$\mathbf h_{\mathcal R,G,\mathcal A}$ in Lemma
\ref{lemma-equivariant-hom-functors}. Forgetting the module structure on
the target is exact, and both derived functors are computed by the same
injective resolution. This proves the final assertion.
\end{proof}


\begin{lemma}
\label{lemma-skew-group-hom-injective-acyclic}
\begin{reference}
\cite[Chapter V, \S 5.6, Corollary and proof, pp. 215--216]{Tohoku}
\end{reference}
\begin{history}
The cited proof compares two actions on a coinduced module. With the
left-action convention used here, the equivariant map in the direction
claimed there has component $f_g\mapsto g^{-1}f_g$, rather than
$f_g\mapsto gf_g$. We use the corrected orientation below.
\end{history}
Let $R$ be a ring endowed with an action of a group $G$. Let $A$ and $B$
be left $R$-modules endowed with compatible semilinear $G$-actions, and
assume that $B$ is injective in the abelian category of such modules. Then
the $G$-module
$$
H=\Hom_R(A,B),
\qquad
(h\varphi)(a)=h\varphi(h^{-1}a),
$$
is right acyclic for invariants. Thus
$$
H^n(G,H)=0
$$
for every $n>0$. We do not assert that $H$ is an injective $G$-module.
\end{lemma}

\begin{proof}
Let
$$
B'=\prod_{g\in G}B.
$$
Write an element as $f=(f_g)_{g\in G}$ and define an $R$-module structure
and a $G$-action by
$$
(r f)_g=g(r)f_g,
\qquad
(h f)_g=f_{gh}.
$$
The group-action identities show that this makes $B'$ an $R$-module with
compatible semilinear $G$-action. The map
$$
\eta:B\longrightarrow B',
\qquad
\eta(b)_g=gb,
$$
is $R$-linear and $G$-equivariant. It is a monomorphism because its
component at $g=1$ is the identity. Since $B$ is injective, $\eta$ splits
in the equivariant category.

We compare $\Hom_R(A,B')$ with a coinduced $G$-module. First define
$$
\Phi:\Hom_R(A,B')\longrightarrow\prod_{g\in G}H,
\qquad
\Phi(u)_g(a)=g^{-1}u(a)_g.
$$
Each $\Phi(u)_g$ is $R$-linear. The map is bijective; its inverse sends a
family $(v_g)$ to the homomorphism given by
$$
u(a)_g=gv_g(a).
$$
Transporting the conjugation action on $\Hom_R(A,B')$ through $\Phi$
gives
$$
(h\star v)_g=h\cdot v_{gh},
$$
where the action on the right is the conjugation action on $H$.

Give
$$
\operatorname{Coind}_{\{1\}}^G(H)=\prod_{g\in G}H
$$
the standard action
$$
(hv)_g=v_{gh};
$$
here the given $G$-action on $H$ is forgotten in forming the coinduction.
The map
$$
\Psi:
\operatorname{Coind}_{\{1\}}^G(H)
\longrightarrow
\left(\prod_{g\in G}H,\star\right),
\qquad
\Psi(v)_g=g^{-1}\cdot v_g,
$$
is $G$-equivariant, since
$$
\Psi(hv)_g
=g^{-1}\cdot v_{gh}
=h\cdot\bigl((gh)^{-1}\cdot v_{gh}\bigr)
=(h\star\Psi(v))_g.
$$
Consequently $\Psi^{-1}\Phi$ is a $G$-isomorphism
$$
\Hom_R(A,B')
\longrightarrow
\operatorname{Coind}_{\{1\}}^G(H).
$$
Directly, its $g$th component is
$$
u\longmapsto\left(a\longmapsto u(g^{-1}a)_g\right).
$$

Every module coinduced from the trivial subgroup is right acyclic for
invariants. Indeed, apply coinduction to an injective resolution of its
underlying abelian group. Coinduction is exact, is right adjoint to the
exact forgetful functor, and the invariant module of a coinduced module is
the original abelian group.

It follows that $\Hom_R(A,B')$ is right acyclic for invariants. The
splitting of $\eta$ makes $\Hom_R(A,B)$ a $G$-equivariant direct summand
of $\Hom_R(A,B')$. Group cohomology respects split direct summands, which
proves the lemma.
\end{proof}


\begin{lemma}
\label{lemma-equivariant-product-internal-hom}
\begin{reference}
\cite[Chapter V, \S 5.6, proof of Lemma 5.6.1, p. 215]{Tohoku}
\end{reference}
\begin{history}
The symbols $I_x$ and $H_x$ below denote the point factors defining
pointwise product sheaves. They are not asserted to be the stalks of those
product sheaves; the two notions need not agree when $x$ is not isolated.
The displayed formula in the cited source also repeats
$(H_x)^{G_x}$ over every point of a single orbit. An invariant family on
an orbit is instead determined by one value at an orbit representative,
giving the corrected factor $H_{\xi(y)}^{G_{\xi(y)}}$ used below.
\end{history}
Let $(X,\mathcal O)$ be a $G$-ringed space, let $q:X\to Y=X/G$, and
choose representatives $\xi(y)\in y$ as in Injectives, Definition
\ref{injectives-definition-equivariant-product-module}. Put
$$
G_y=G_{\xi(y)},
\qquad
\mathcal U_y=\mathcal U_{\xi(y)}.
$$
Let $(I_y)_{y\in Y}$ be a family of $\mathcal U_y$-modules, let $I_x$
be the transported factors of that definition, and put
$$
\mathcal B=\mathcal P(I).
$$
For a $G$-$\mathcal O$-module $\mathcal A$, set
$$
H_x=\Hom_{\mathcal O_x}(\mathcal A_x,I_x).
$$
These modules have compatible semilinear transports
$$
H_x\longrightarrow H_{gx},
\qquad
\varphi\longmapsto
\left(a\longmapsto g\varphi(g^{-1}a)\right).
$$
Write
$$
H_y=H_{\xi(y)}
=\Hom_{\mathcal O_{\xi(y)}}(\mathcal A_{\xi(y)},I_y).
$$
There is a canonical isomorphism of $G$-$\mathcal O$-modules
$$
\SheafHom_{\mathcal O}(\mathcal A,\mathcal P(I))
\longrightarrow
\mathcal P(H).
$$
For every open $U\subset Y$, this induces an isomorphism
$$
\left(
q_*^G\SheafHom_{\mathcal O}(\mathcal A,\mathcal P(I))
\right)(U)
\longrightarrow
\prod_{y\in U}H_y^{G_y}.
$$
In particular,
$q_*^G\SheafHom_{\mathcal O}(\mathcal A,\mathcal P(I))$
is flasque. There is also a canonical isomorphism of $G$-modules
$$
\Hom_{\mathcal O}(\mathcal A,\mathcal P(I))
\longrightarrow
\prod_{y\in Y}\operatorname{Coind}_{G_y}^G(H_y).
$$
If every $H_y$ is right acyclic for $G_y$-invariants, then the $G$-module
on the left is right acyclic for $G$-invariants.
\end{lemma}

\begin{proof}
For an open $W\subset X$, the pointwise product description and the
stalk--skyscraper adjunction give
$$
\begin{aligned}
\Gamma\left(
W,\SheafHom_{\mathcal O}(\mathcal A,\mathcal P(I))
\right)
&=
\Hom_{\mathcal O|_W}
(\mathcal A|_W,\mathcal P(I)|_W) \\
&=
\prod_{x\in W}
\Hom_{\mathcal O_x}(\mathcal A_x,I_x) \\
&=
\mathcal P(H)(W).
\end{aligned}
$$
These identifications commute with restrictions. The displayed transport
formula shows that they also commute with the $G$-actions, proving the
first assertion.

For $U\subset Y$, an invariant family
$$
(s_x)_{x\in q^{-1}(U)}\in\prod_{x\in q^{-1}(U)}H_x
$$
is uniquely determined by its values at the representatives $\xi(y)$,
and such a value may be chosen precisely in $H_y^{G_y}$. This proves the
invariant-direct-image formula. Its restriction maps are projections
between products, so the resulting sheaf is flasque.

Similarly,
$$
\Gamma(X,\mathcal P(H))
=
\prod_{y\in Y}\prod_{x\in y}H_x.
$$
For each orbit, transport and inversion of cosets identify
$$
\prod_{x\in y}H_x
=
\operatorname{Coind}_{G_y}^G(H_y),
$$
as in Injectives, Definition
\ref{injectives-definition-equivariant-product-module}. This gives the
global formula.

Coinduction from $G_y$ is exact because $\mathbf Z[G]$ is free as a left
$\mathbf Z[G_y]$-module, and it is right adjoint to the exact restriction
functor. Applying it to an injective $G_y$-module resolution gives
canonical isomorphisms
$$
H^n\left(G,\operatorname{Coind}_{G_y}^G(M)\right)
=
H^n(G_y,M).
$$
Finally, group cohomology commutes with arbitrary products. Indeed, for a
projective resolution $P_\bullet\to\mathbf Z$ of the trivial
$\mathbf Z[G]$-module,
$$
\Hom_{\mathbf Z[G]}
\left(P_\bullet,\prod_iM_i\right)
=
\prod_i\Hom_{\mathbf Z[G]}(P_\bullet,M_i),
$$
and products of abelian groups are exact. The final assertion follows.
\end{proof}


\begin{lemma}
\label{lemma-equivariant-hom-injective-acyclic}
\begin{reference}
\cite[Chapter V, \S 5.6, Lemma 5.6.1 and proof, p. 215]{Tohoku}
\end{reference}
Let $(X,\mathcal O)$ be a $G$-ringed space, let $q:X\to Y=X/G$, and let
$\mathcal A,\mathcal B$ be $G$-$\mathcal O$-modules. Assume that
$\mathcal B$ is injective in $\textit{Mod}_G(\mathcal O)$ and put
$$
\mathcal H=
\SheafHom_{\mathcal O}(\mathcal A,\mathcal B).
$$
Then:
\begin{enumerate}
\item the $\mathcal O'=q_*^G\mathcal O$-module
$$
q_*^G\mathcal H
=
\mathbf h_{\mathcal O,G,\mathcal A}(\mathcal B)
$$
is flasque on $Y$;
\item the $G$-module
$$
\Gamma(X,\mathcal H)
=
\Hom_{\mathcal O}(\mathcal A,\mathcal B)
$$
is right acyclic for invariants;
\item the abelian $G$-sheaf $\mathcal H$ is right acyclic for
$\Gamma_X^G$, i.e.,
$$
H^n(X;G,\mathcal H)=0
$$
for every $n>0$.
\end{enumerate}
In fact, the underlying abelian sheaf of $\mathcal H$ is flasque. None of
these assertions says that $\Hom_{\mathcal O}(\mathcal A,\mathcal B)$ is
an injective $G$-module.
\end{lemma}

\begin{proof}
By Injectives, Lemma
\ref{injectives-lemma-equivariant-product-injectives}, there is a
monomorphism
$$
\mathcal B\longrightarrow\mathcal P(I)
$$
where every $I_y$ is an injective $\mathcal U_y$-module. Since
$\mathcal B$ is injective, this monomorphism splits equivariantly.

First consider $\mathcal P(I)$. For every $y\in Y$, the modules
$\mathcal A_{\xi(y)}$ and $I_y$ are modules over the one-point
$G_y$-ringed space with ring $\mathcal O_{\xi(y)}$. Equivalently, they
are modules over the skew group ring $\mathcal U_y$ of Injectives,
Definition \ref{injectives-definition-stabilizer-skew-group-ring}. Since
$I_y$ is injective, Lemma
\ref{lemma-skew-group-hom-injective-acyclic} shows that
$$
H_y=
\Hom_{\mathcal O_{\xi(y)}}(\mathcal A_{\xi(y)},I_y)
$$
is right acyclic for $G_y$-invariants.

Lemma \ref{lemma-equivariant-product-internal-hom} now shows that
$$
q_*^G\SheafHom_{\mathcal O}(\mathcal A,\mathcal P(I))
$$
is flasque and that
$$
\Hom_{\mathcal O}(\mathcal A,\mathcal P(I))
$$
is right acyclic for $G$-invariants. It also identifies the underlying
internal Hom with the pointwise product sheaf $\mathcal P(H)$, which is
flasque because all its restriction maps are projections.

Applying $\SheafHom_{\mathcal O}(\mathcal A,-)$ to the splitting of
$\mathcal B\to\mathcal P(I)$ makes $\mathcal H$ a $G$-equivariant direct
summand of
$$
\SheafHom_{\mathcal O}(\mathcal A,\mathcal P(I)).
$$
Invariant direct image and global sections preserve this splitting. A
direct summand of a flasque sheaf is flasque, and group cohomology
preserves split direct summands. This proves (1), (2), and the additional
flasqueness assertion.

Finally, since $\mathcal H$ is flasque, Lemma
\ref{lemma-orbit-space-equivariant-cohomology-global-acyclic} gives
canonical isomorphisms
$$
H^n(G,\Gamma(X,\mathcal H))
\longrightarrow
H^n(X;G,\mathcal H).
$$
The left hand side vanishes for $n>0$ by (2), proving (3).
\end{proof}


\begin{lemma}
\label{lemma-equivariant-higher-invariant-direct-image-modules}
\begin{reference}
\cite[Chapter V, \S 5.6, proof of Lemma 5.6.1 and the following
discussion, p. 215]{Tohoku}
\end{reference}
Let $(X,\mathcal O)$ be a $G$-ringed space, let $q:X\to Y=X/G$, and put
$\mathcal O'=q_*^G\mathcal O$. If $\mathcal F$ is a
$G$-$\mathcal O$-module, then there are canonical functorial isomorphisms
of underlying abelian sheaves
$$
R^n\left(q_*^G:\textit{Mod}_G(\mathcal O)
\longrightarrow\textit{Mod}(\mathcal O')\right)(\mathcal F)
\longrightarrow
\mathcal H^n(G,\mathcal F)
$$
for every $n\geq 0$. On the right the derived functor is computed after
forgetting the $\mathcal O$-module structure, as in Definition
\ref{definition-orbit-space-equivariant-cohomology}.

If $\mathcal A,\mathcal B$ are $G$-$\mathcal O$-modules and $\mathcal B$
is injective in $\textit{Mod}_G(\mathcal O)$, then
$$
\SheafHom_{\mathcal O}(\mathcal A,\mathcal B)
$$
is right acyclic for invariant direct image on $G$-$\mathcal O$-modules.
\end{lemma}

\begin{proof}
Write $v$ for the functor which forgets the $\mathcal O$-module structure.
Let $\mathcal I$ be injective in $\textit{Mod}_G(\mathcal O)$. For an open
$U\subset Y$, put $X_U=q^{-1}(U)$. Equivariant extension by zero along
$X_U\subset X$ is an exact left adjoint to restriction, as in the proof of
Lemma \ref{lemma-localize-equivariant-ext}. Hence
$\mathcal I|_{X_U}$ is injective in
$\textit{Mod}_G(\mathcal O|_{X_U})$.

Evaluation at $1$ identifies
$$
\SheafHom_{\mathcal O|_{X_U}}
(\mathcal O|_{X_U},\mathcal I|_{X_U})
=
\mathcal I|_{X_U}
$$
as abelian $G$-sheaves. Part (3) of Lemma
\ref{lemma-equivariant-hom-injective-acyclic}, applied on $X_U$, therefore
gives
$$
H^m(X_U;G,v\mathcal I|_{X_U})=0
$$
for $m>0$. Lemma \ref{lemma-localize-higher-invariant-direct-image}
shows that $v\mathcal I$ is right acyclic for invariant direct image on
abelian $G$-sheaves.

Now choose an injective resolution
$\mathcal F[0]\to\mathcal I^\bullet$ in
$\textit{Mod}_G(\mathcal O)$. The preceding paragraph and Derived
Categories, Lemma \ref{derived-lemma-leray-acyclicity}, show that
$v\mathcal I^\bullet$ computes invariant direct image after forgetting
$\mathcal O$. Term by term the two invariant direct images agree, and the
forgetful functor on the target is exact. Taking cohomology gives the first
assertion.

Finally, suppose that $\mathcal B$ is injective and put
$\mathcal H=\SheafHom_{\mathcal O}(\mathcal A,\mathcal B)$. For every
$U\subset Y$, restriction preserves the injectivity of $\mathcal B$, and
part (3) of Lemma \ref{lemma-equivariant-hom-injective-acyclic}, applied on
$X_U$, gives
$$
H^m(X_U;G,\mathcal H|_{X_U})=0
$$
for $m>0$. The sheafification assertion of Lemma
\ref{lemma-localize-higher-invariant-direct-image} and the comparison just
proved give $R^mq_*^G\mathcal H=0$ for $m>0$ in
$\textit{Mod}(\mathcal O')$.
\end{proof}


\begin{lemma}
\label{lemma-equivariant-underlying-ext}
\begin{reference}
\cite[Chapter V, \S 5.6, pp. 216--217]{Tohoku}
\end{reference}
Let $(X,\mathcal O)$ be a $G$-ringed space and let
$\mathcal A,\mathcal B$ be $G$-$\mathcal O$-modules. For every $q\geq 0$
there are canonical functorial isomorphisms
$$
R^qh_{\mathcal O,\mathcal A}(\mathcal B)
\longrightarrow
\Ext^q_{\textit{Mod}(\mathcal O)}(\mathcal A,\mathcal B)
$$
of abelian $G$-modules and
$$
R^q\mathbf h_{\mathcal O,\mathcal A}(\mathcal B)
\longrightarrow
\SheafExt^q_{\mathcal O}(\mathcal A,\mathcal B)
$$
of $G$-$\mathcal O$-modules. On the right hand sides the $G$-actions are those
induced by the equivariant structures on $\mathcal A$ and $\mathcal B$.
\end{lemma}

\begin{proof}
Choose an injective resolution
$$
\mathcal B[0]\longrightarrow\mathcal I^\bullet
$$
in $\textit{Mod}_G(\mathcal O)$. By Injectives, Lemma
\ref{injectives-lemma-equivariant-injective-underlying}, every
$\mathcal I^n$ is injective after forgetting the $G$-action. Hence the
complexes
$$
\Hom_{\mathcal O}(\mathcal A,\mathcal I^\bullet)
\quad\text{and}\quad
\SheafHom_{\mathcal O}(\mathcal A,\mathcal I^\bullet)
$$
compute, respectively, the ordinary global Ext groups and ordinary sheaf
Ext. They also compute the two derived functors on the left by Lemma
\ref{lemma-equivariant-hom-functors}. The conjugation actions on the two
complexes commute with their differentials, and the resulting actions on
cohomology are independent of the chosen equivariant injective resolution.
The second complex and its cohomology retain their $\mathcal O$-module
structures. This proves both isomorphisms and their functoriality.
\end{proof}


\begin{lemma}
\label{lemma-equivariant-ext-spectral-sequences}
\begin{reference}
\cite[Chapter V, \S 5.6, pp. 216--217]{Tohoku}
\end{reference}
\begin{history}
The official print has $Y$ in the first spectral sequence below. This is
forced by the types: $\SheafExt^q_{\mathcal O,G}$ is a sheaf on
$Y=X/G$. A later local correction changed $Y$ to $X$; we retain the
type-correct authority reading.
\end{history}
Let $(X,\mathcal O)$ be a $G$-ringed space, let $q:X\to Y=X/G$ be the
orbit map, and let $\mathcal A,\mathcal B$ be $G$-$\mathcal O$-modules.
For fixed $\mathcal A$, there are three additive first-quadrant spectral
sequences, functorial in $\mathcal B$ from the $E_2$-page onward,
$$
{}^{\mathrm I}E_2^{p,q}
=
H^p\left(
Y,\SheafExt^q_{\mathcal O,G}(\mathcal A,\mathcal B)
\right)
\Longrightarrow
\Ext^{p+q}_{\mathcal O,G}(X;\mathcal A,\mathcal B),
$$
$$
{}^{\mathrm {II}}E_2^{p,q}
=
H^p\left(
G,\Ext^q_{\textit{Mod}(\mathcal O)}(\mathcal A,\mathcal B)
\right)
\Longrightarrow
\Ext^{p+q}_{\mathcal O,G}(X;\mathcal A,\mathcal B),
$$
and
$$
{}^{\mathrm {III}}E_2^{p,q}
=
H^p\left(
X;G,\SheafExt^q_{\mathcal O}(\mathcal A,\mathcal B)
\right)
\Longrightarrow
\Ext^{p+q}_{\mathcal O,G}(X;\mathcal A,\mathcal B).
$$
Each abutment carries the canonical finite filtration supplied by the
spectral sequence. The three spectral sequences have their canonical edge
homomorphisms and associated five-term exact sequences.
\end{lemma}

\begin{proof}
Apply Lemma \ref{derived-lemma-grothendieck-spectral-sequence} to each
of the three factorizations in Lemma
\ref{lemma-equivariant-hom-functors}.

For the first factorization
$$
h_{\mathcal O,G,\mathcal A}
=
\Gamma(Y,-)\circ\mathbf h_{\mathcal O,G,\mathcal A},
$$
Lemma \ref{lemma-equivariant-hom-injective-acyclic}, part (1), says that
the inner functor sends injectives to flasque $\mathcal O'$-modules on
$Y$. These are right acyclic for global sections by Lemma
\ref{lemma-flasque-acyclic}. The inner derived functors are
$\SheafExt^q_{\mathcal O,G}$ by Definition
\ref{definition-equivariant-ext}, giving spectral sequence $\mathrm I$.

For the second factorization
$$
h_{\mathcal O,G,\mathcal A}
=
(-)^G\circ h_{\mathcal O,\mathcal A},
$$
part (2) of Lemma \ref{lemma-equivariant-hom-injective-acyclic} supplies
the acyclicity hypothesis for invariants. Lemma
\ref{lemma-equivariant-underlying-ext} identifies the inner derived
functors with
$\Ext^q_{\textit{Mod}(\mathcal O)}(\mathcal A,\mathcal B)$, including
their $G$-actions. This gives spectral sequence $\mathrm {II}$.

For the third factorization
$$
h_{\mathcal O,G,\mathcal A}
=
\Gamma_X^G\circ\mathbf h_{\mathcal O,\mathcal A},
$$
part (3) of Lemma \ref{lemma-equivariant-hom-injective-acyclic} gives the
required $\Gamma_X^G$-acyclicity. The second isomorphism of Lemma
\ref{lemma-equivariant-underlying-ext} identifies the inner derived
functors with $\SheafExt^q_{\mathcal O}$. This gives spectral sequence
$\mathrm {III}$.

In all three cases the abutment is the right derived functor of
$h_{\mathcal O,G,\mathcal A}$, hence is
$\Ext^{p+q}_{\mathcal O,G}(X;\mathcal A,\mathcal B)$ by Definition
\ref{definition-equivariant-ext}. Additivity, functoriality, and the finite
filtrations are those of the Grothendieck spectral sequence.
\end{proof}


\begin{lemma}
\label{lemma-equivariant-sheaf-ext-spectral-sequence}
\begin{reference}
\cite[Chapter V, \S 5.6, formulas (5.6.4) and (5.6.6)--(5.6.8),
pp. 214--215]{Tohoku}
\end{reference}
Let $(X,\mathcal O)$ be a $G$-ringed space, let $q:X\to Y=X/G$, put
$\mathcal O'=q_*^G\mathcal O$, and fix a $G$-$\mathcal O$-module
$\mathcal A$. There is an additive first-quadrant spectral sequence of
$\mathcal O'$-modules, functorial in the $G$-$\mathcal O$-module
$\mathcal B$ from the $E_2$-page onward, with
$$
E_2^{p,q}
=
R^pq_*^G\left(\SheafExt^q_{\mathcal O}
(\mathcal A,\mathcal B)\right)
=
\mathcal H^p\left(G,\SheafExt^q_{\mathcal O}
(\mathcal A,\mathcal B)\right)
$$
The final term has the canonical $\mathcal O'$-module structure transported
through the comparison of Lemma
\ref{lemma-equivariant-higher-invariant-direct-image-modules}. The spectral
sequence converges to
$$
\SheafExt^{p+q}_{\mathcal O,G}(\mathcal A,\mathcal B).
$$
It induces a finite filtration on every abutment sheaf.
\end{lemma}

\begin{proof}
Use the factorization
$$
\mathbf h_{\mathcal O,G,\mathcal A}
=
q_*^G\circ\mathbf h_{\mathcal O,\mathcal A}
$$
of Lemma \ref{lemma-equivariant-hom-functors}. The category
$\textit{Mod}_G(\mathcal O)$ has enough injectives by Injectives, Lemma
\ref{injectives-lemma-equivariant-modules-enough-injectives}, and invariant
direct image is additive left exact by Modules, Lemma
\ref{modules-lemma-invariant-direct-image-G-O-modules}. Lemma
\ref{lemma-equivariant-higher-invariant-direct-image-modules} gives the
required acyclicity on the image of every injective object. Derived
Categories, Lemma \ref{derived-lemma-grothendieck-spectral-sequence}
therefore applies.

Lemma \ref{lemma-equivariant-underlying-ext} identifies the inner derived
functors with $\SheafExt^q_{\mathcal O}(\mathcal A,\mathcal B)$ as
$G$-$\mathcal O$-modules. Lemma
\ref{lemma-equivariant-higher-invariant-direct-image-modules} identifies the
outer derived functors with the displayed $\mathcal H^p$, while Definition
\ref{definition-equivariant-ext} identifies the abutment. The remaining
assertions follow from the Grothendieck spectral sequence.
\end{proof}


\begin{lemma}
\label{lemma-equivariant-ext-locally-free}
\begin{reference}
\cite[Chapter V, \S 5.6, p. 218]{Tohoku}
\end{reference}
Let $(X,\mathcal O)$ be a $G$-ringed space and let
$\mathcal A,\mathcal B$ be $G$-$\mathcal O$-modules. Assume that, for some
$r\geq 0$, the underlying $\mathcal O$-module $\mathcal A$ is locally
isomorphic to $\mathcal O^r$. Then the third spectral sequence of Lemma
\ref{lemma-equivariant-ext-spectral-sequences} degenerates and its edge maps
give canonical isomorphisms, functorial in $\mathcal B$,
$$
H^n\left(X;G,\SheafHom_{\mathcal O}(\mathcal A,\mathcal B)\right)
\longrightarrow
\Ext^n_{\mathcal O,G}(X;\mathcal A,\mathcal B)
$$
for every $n\geq 0$. In particular, there are canonical isomorphisms,
functorial in $\mathcal B$,
$$
H^n(X;G,\mathcal B)
\longrightarrow
\Ext^n_{\mathcal O,G}(X;\mathcal O,\mathcal B)
$$
for every $n\geq 0$.
\end{lemma}

\begin{proof}
The functor $\SheafHom_{\mathcal O}(\mathcal A,-)$ is exact: exactness is
local, and on an open where $\mathcal A\cong\mathcal O^r$ this functor is
isomorphic to $(-)^r$. Hence
$$
\SheafExt^q_{\mathcal O}(\mathcal A,\mathcal B)=0
\quad\text{for }q>0
$$
and the third spectral sequence is concentrated in the row $q=0$.
Its bottom edge gives the first isomorphism. For
$\mathcal A=\mathcal O$, evaluation at $1$ gives a canonical isomorphism
$$
\SheafHom_{\mathcal O}(\mathcal O,\mathcal B)
\longrightarrow
\mathcal B
$$
of $G$-$\mathcal O$-modules: indeed,
$(g\varphi)(1)=g\varphi(g^{-1}1)=g\varphi(1)$. This gives the second
isomorphism.
\end{proof}


\begin{lemma}[Equivariant extension classes for finite locally free modules]
\label{lemma-equivariant-extension-classes-finite-locally-free}
\begin{reference}
\cite[Chapter V, \S 5.6, pp. 218--219]{Tohoku}
\end{reference}
Let $(X,\mathcal O)$ be a commutative $G$-ringed space and let
$\mathcal E,\mathcal F$ be $G$-$\mathcal O$-modules. Assume that the
underlying $\mathcal O$-module $\mathcal E$ is locally isomorphic to
$\mathcal O^r$ for some $r\geq 0$. Then equivalence classes of short exact
sequences
$$
0 \longrightarrow \mathcal F \longrightarrow \mathcal V
\longrightarrow \mathcal E \longrightarrow 0
$$
in $\textit{Mod}_G(\mathcal O)$ are canonically in bijection with
$\Ext^1_{\mathcal O,G}(X;\mathcal E,\mathcal F)$. Moreover, there are
canonical isomorphisms
$$
\Ext^1_{\mathcal O,G}(X;\mathcal E,\mathcal F)
\longrightarrow
H^1\left(X;G,\SheafHom_{\mathcal O}(\mathcal E,\mathcal F)\right)
$$
and
$$
\mathcal E^\vee\otimes_{\mathcal O}\mathcal F
\longrightarrow
\SheafHom_{\mathcal O}(\mathcal E,\mathcal F),
\qquad
\mathcal E^\vee=
\SheafHom_{\mathcal O}(\mathcal E,\mathcal O),
$$
where the second arrow is the evaluation isomorphism. Consequently,
$$
\Ext^1_{\mathcal O,G}(X;\mathcal E,\mathcal F)
\cong
H^1\left(X;G,\mathcal E^\vee\otimes_{\mathcal O}\mathcal F\right).
$$
\end{lemma}

\begin{proof}
The category $\textit{Mod}_G(\mathcal O)$ is abelian and has enough
injectives. The degree-one Yoneda description of Ext in an abelian
category identifies its degree-one Ext group with equivalence classes of
degree-one Yoneda extensions, which are exactly the displayed short exact sequences. The
definition of $\Ext^1_{\mathcal O,G}$ and Lemma
\ref{lemma-equivariant-ext-properties} identify that group with
$\Ext^1_{\textit{Mod}_G(\mathcal O)}(\mathcal E,\mathcal F)$.

\medskip\noindent
Since $\mathcal E$ is locally isomorphic to $\mathcal O^r$, Lemma
\ref{lemma-equivariant-ext-locally-free}, with $n=1$, gives the first
displayed isomorphism. The internal-hom action of Lemma
\ref{lemma-equivariant-internal-hom} makes
$\mathcal E^\vee$ a $G$-$\mathcal O$-module. The evaluation map
$$
\mathcal E^\vee\otimes_{\mathcal O}\mathcal F
\longrightarrow
\SheafHom_{\mathcal O}(\mathcal E,\mathcal F),
\qquad
\lambda\otimes f\longmapsto(e\longmapsto\lambda(e)f),
$$
is $G$-equivariant by the conjugation formula in that lemma. Locally on
$X$ it is the standard isomorphism
$(\mathcal O^r)^\vee\otimes\mathcal F\cong
\SheafHom_{\mathcal O}(\mathcal O^r,\mathcal F)$, hence it is an
isomorphism globally. Applying $H^1(X;G,-)$ and composing with the first
isomorphism proves the final assertion.
\end{proof}

\begin{remark}[Analytic boundary]
This lemma records the abstract equivariant Ext bridge. Identifying
holomorphic bundles with analytic $G$-$\mathcal O$-modules, or invoking
GAGA, requires separate analytic hypotheses and is not claimed here.
\label{remark-equivariant-extension-classes-analytic-boundary}
\end{remark}

\begin{lemma}[Holomorphic bundles and locally free analytic modules]
\label{lemma-holomorphic-vector-bundles-locally-free-analytic-modules}
\begin{reference}
\cite[Chapter V, \S 5.6, p. 219]{Tohoku}
\end{reference}
\begin{history}
This is the analytic bundle--module correspondence. It is separate from
the scheme-theoretic vector-bundle construction and does not use GAGA.
\end{history}
Let $X$ be a complex analytic space and let $n>0$. Sending a holomorphic
vector bundle $E$ of rank $n$ to its sheaf $\mathcal O(E)$ of holomorphic
sections defines an equivalence between
\begin{enumerate}
\item holomorphic vector bundles of rank $n$ on $X$, and
\item analytic $\mathcal O_X$-modules locally isomorphic to
$\mathcal O_X^n$.
\end{enumerate}
Every module in (2) is coherent. Under the inverse equivalence, the fibre
at $x\in X$ is canonically
$$
(E_{\mathcal L})_x
\cong
\mathcal L_x\otimes_{\mathcal O_{X,x}}\mathbf C,
$$
where $\mathbf C=\mathcal O_{X,x}/\mathfrak m_x$ is regarded as an
$\mathcal O_{X,x}$-module through the residue, or evaluation, map.
\end{lemma}

\begin{proof}
Choose a trivializing open covering $X=\bigcup U_i$ of a holomorphic vector
bundle $E$. On every $U_i$ its sheaf of holomorphic sections is
$\mathcal O_{U_i}^n$. In bundle trivializations, a holomorphic bundle map is
given by a matrix of holomorphic functions. The same matrix gives the
corresponding $\mathcal O_X$-linear map on section sheaves, and every such
module map is locally given by such a matrix. Hence the section-sheaf
functor is fully faithful.

Conversely, let $\mathcal L$ be an analytic $\mathcal O_X$-module locally
isomorphic to $\mathcal O_X^n$. Choose an open covering and bases
$$
\mathcal L|_{U_i}\longrightarrow\mathcal O_{U_i}^n.
$$
On $U_i\cap U_j$ the changes of basis are matrices
$$
A_{ij}\in\operatorname{GL}_n(\mathcal O_X(U_i\cap U_j))
$$
satisfying the cocycle identities. Gluing the trivial bundles
$U_i\times\mathbf C^n$ with these transition matrices produces a
holomorphic vector bundle $E_{\mathcal L}$, and the chosen bases glue to a
canonical isomorphism
$$
\mathcal O(E_{\mathcal L})\longrightarrow\mathcal L.
$$
In a chosen local basis, the fibre of the glued trivial bundle at $x$ is
$\mathbf C^n$, while tensoring the corresponding stalk basis with
$\mathcal O_{X,x}/\mathfrak m_x$ gives the same vector space. The
transition matrices on both sides are obtained by evaluation at $x$, so
these local identifications glue to the displayed canonical fibre
isomorphism.
Changing the covering or the bases gives the uniquely isomorphic bundle
corresponding under full faithfulness. Thus the section-sheaf functor is an
equivalence. Finally, an analytic module locally isomorphic to
$\mathcal O_X^n$ is coherent by Oka coherence; see
\cite[p. 16]{Grauert-direct-image}.
\end{proof}

\begin{lemma}[Equivariant holomorphic bundles and locally free modules]
\label{lemma-equivariant-holomorphic-vector-bundles-locally-free-modules}
\begin{reference}
\cite[Chapter V, \S 5.6, p. 219]{Tohoku}
\end{reference}
\begin{history}
This is the equivariant analytic bundle--module correspondence, not merely
the algebraic definition of an equivariant module. It does not assert the
following classification of bundle extensions, a cohomological
identification, descent to the quotient, or a GAGA comparison.
\end{history}
Let $X$ be a complex analytic space, let $n>0$, and let a group $G$ act on
$X$ by analytic automorphisms. The equivalence of the preceding lemma
restricts to an equivalence between holomorphic vector bundles of rank $n$
with a compatible $G$-action and $G$-$\mathcal O_X$-modules whose underlying
$\mathcal O_X$-module is locally isomorphic to $\mathcal O_X^n$.
Here the action on a bundle is by holomorphic, fibrewise linear
isomorphisms over the action on $X$, and the module action is in the sense
of Modules, Definition \ref{modules-definition-G-O-module}.
\end{lemma}

\begin{proof}
A compatible action on $E$ sends a
section $s$ over $U$ to the section over $gU$ defined by
$$
(\rho_gs)(gx)=g(s(x)).
$$
These maps obey the group law and
$$
\rho_g(as)=\rho_g(a)\rho_g(s),
$$
so $\mathcal O(E)$ is a $G$-$\mathcal O_X$-module.

In the other direction, let $\mathcal L$ be a locally free
$G$-$\mathcal O_X$-module. In local bases on $U$ and $gU$, each semilinear
map
$$
\rho_g:\mathcal L(U)\longrightarrow\mathcal L(gU)
$$
is represented by a matrix of holomorphic functions. The matrix is
invertible because $\rho_{g^{-1}}$ is its inverse after transport of the
base. It therefore defines a holomorphic fibrewise linear bundle
isomorphism
$$
E_{\mathcal L}|_U\longrightarrow E_{\mathcal L}|_{gU}
$$
over $x\mapsto gx$. Independence of bases follows from the transition
matrices, and the identities $\rho_{gh}=\rho_g\rho_h$ make these local maps
an action of $G$ on $E_{\mathcal L}$. The two constructions are inverse,
and full faithfulness identifies equivariant morphisms on both sides.
\end{proof}

\begin{lemma}[Coherent equivariant analytic extension classes]
\label{lemma-coherent-equivariant-analytic-extension-classes}
\begin{reference}
\cite[Chapter V, \S 5.6, p. 219]{Tohoku}
\end{reference}
Let $X$ be a complex analytic space on which a group $G$ acts by analytic
automorphisms. Let $\mathcal A$ and $\mathcal B$ be $G$-$\mathcal O_X$-modules
whose underlying $\mathcal O_X$-modules are coherent. For every short exact
sequence
$$
0 \longrightarrow \mathcal B \longrightarrow \mathcal M
\longrightarrow \mathcal A \longrightarrow 0
$$
in $\textit{Mod}_G(\mathcal O_X)$, the underlying $\mathcal O_X$-module
$\mathcal M$ is coherent. Moreover, equivalence classes of these extensions,
with fixed endpoints $\mathcal A$ and $\mathcal B$, are canonically identified
with
$$
\Ext^1_{\mathcal O_X,G}(X;\mathcal A,\mathcal B).
$$
This identification is compatible with Baer addition. The group on the right,
and hence the set of extension classes, is naturally a complex vector space.
\end{lemma}

\begin{proof}
By Modules, Lemma
\ref{modules-lemma-AB5-AB3-star-generator-G-O-modules}, kernels and cokernels
in $\textit{Mod}_G(\mathcal O_X)$ are computed on the underlying
$\mathcal O_X$-modules. Thus the displayed sequence remains exact after
forgetting $G$. Since its two endpoints are coherent, Modules, Lemma
\ref{modules-lemma-coherent-abelian} shows that $\mathcal M$ is coherent.

The category $\textit{Mod}_G(\mathcal O_X)$ is abelian and has enough
injectives. Derived Categories, Lemma \ref{derived-lemma-ext-1} identifies
degree-one Ext with equivalence classes of fixed-endpoint short exact
sequences, with the group law given by Baer addition. Cohomology, Lemma
\ref{lemma-equivariant-ext-properties} gives the canonical isomorphism
$$
\Ext^1_{\textit{Mod}_G(\mathcal O_X)}(\mathcal A,\mathcal B)
\cong
\Ext^1_{\mathcal O_X,G}(X;\mathcal A,\mathcal B),
$$
which proves the classification and its compatibility with addition.

Finally, analytic automorphisms fix the constant holomorphic functions. Hence
multiplication by every $\lambda\in\mathbf C$ is an equivariant
$\mathcal O_X$-linear endomorphism of any $G$-$\mathcal O_X$-module. If
$\mathcal B[0]\to\mathcal I^\bullet$ is an injective resolution, then
$$
\Hom_{\textit{Mod}_G(\mathcal O_X)}
(\mathcal A,\mathcal I^\bullet)
$$
is consequently a complex of complex vector spaces with complex-linear
differentials. Its first cohomology is the displayed equivariant Ext group by
the proof of Lemma \ref{lemma-equivariant-ext-properties}. This supplies the
asserted complex vector space structure; under the Yoneda identification,
scalar multiplication is equivalently pushout along
$\lambda\operatorname{id}_{\mathcal B}$.
\end{proof}

\begin{lemma}[Equivariant holomorphic bundle-extension classes]
\label{lemma-equivariant-holomorphic-vector-bundle-extension-classes}
\begin{reference}
\cite[Chapter V, \S 5.6, p. 219]{Tohoku}
\end{reference}
Let $X$ be a complex analytic space on which a group $G$ acts by analytic
automorphisms. Let $E$ and $F$ be holomorphic vector bundles on $X$ with
compatible $G$-actions. Equivalence classes of short exact sequences
$$
0 \longrightarrow F \longrightarrow V \longrightarrow E \longrightarrow 0
$$
of $G$-equivariant holomorphic vector bundles are canonically in bijection
with
$$
\Ext^1_{\mathcal O_X,G}
(X;\mathcal O(E),\mathcal O(F)).
$$
The bijection is compatible with Baer sum. Composing it with Lemma
\ref{lemma-equivariant-ext-locally-free} gives the canonical identification
$$
\Ext^1_{\mathcal O_X,G}
(X;\mathcal O(E),\mathcal O(F))
\cong
H^1\left(X;G,
\SheafHom_{\mathcal O_X}(\mathcal O(E),\mathcal O(F))\right).
$$
Together with the canonical evaluation identification, this is canonically
isomorphic to
$$
H^1\left(X;G,\mathcal O(E^\vee\otimes F)\right).
$$
\end{lemma}

\begin{proof}
Apply the section-sheaf functor of Lemma
\ref{lemma-equivariant-holomorphic-vector-bundles-locally-free-modules}.
A short exact sequence of equivariant holomorphic vector bundles is locally
split after forgetting the $G$-action, and hence gives a short exact sequence
$$
0 \longrightarrow \mathcal O(F) \longrightarrow \mathcal O(V)
\longrightarrow \mathcal O(E) \longrightarrow 0
$$
in $\textit{Mod}_G(\mathcal O_X)$.

Conversely, consider an extension in $\textit{Mod}_G(\mathcal O_X)$
$$
0 \longrightarrow \mathcal O(F) \longrightarrow \mathcal M
\longrightarrow \mathcal O(E) \longrightarrow 0.
$$
After forgetting the $G$-action, this sequence splits locally. Indeed, near
any point we may identify $\mathcal O(E)$ with $\mathcal O_X^r$; after
shrinking, local lifts in $\mathcal M$ of its standard basis define a
splitting. Thus locally
$$
\mathcal M \cong \mathcal O(F)\oplus\mathcal O(E),
$$
so $\mathcal M$ is finite locally free. The equivariant bundle--module
equivalence therefore supplies a $G$-equivariant holomorphic vector bundle
$V$ with $\mathcal O(V)\cong\mathcal M$. Full faithfulness transports the
two maps, and their locally split exact section sequence shows that
$0\to F\to V\to E\to0$ is exact.

These constructions are inverse on extension classes and preserve the
pullbacks and pushouts defining Baer sum. Lemma
\ref{lemma-equivariant-extension-classes-finite-locally-free} now gives the
first displayed bijection, and Lemma
\ref{lemma-equivariant-ext-locally-free} gives the first cohomological
identification. Finally, duals and tensor products commute with the
equivariant bundle--module equivalence, so the natural evaluation map gives
the $G$-equivariant isomorphism
$$
\mathcal O(E^\vee\otimes F)
\cong
\SheafHom_{\mathcal O_X}(\mathcal O(E),\mathcal O(F)).
$$
Applying $H^1(X;G,-)$ gives the last identification.
\end{proof}

\begin{lemma}[Coherence on a discontinuous analytic quotient]
\label{lemma-discontinuous-analytic-invariant-direct-image-coherent}
\begin{reference}
\cite[Chapter V, \S 5.6, p. 219]{Tohoku}
\cite[Theorems 1 and 2, pp. 1--6]{Cartan-analytic-quotient}
\cite[Hauptsatz I, p. 57]{Grauert-direct-image}
\end{reference}
\begin{history}
The notation $\mathcal L^G$ in the cited source means the invariant direct
image $q_*^G\mathcal L$ on the quotient, not a subsheaf on $X$. Coherence is
a local finite-analytic-quotient statement and does not require the
characteristic-zero hypothesis used there for the cohomology comparison. No
GAGA comparison, global properness, compactness, or finite-dimensionality
assertion is made.
\end{history}
Let $X$ be a complex analytic space and let $G$ act discontinuously on $X$
by analytic automorphisms. Let
$$
q:X\longrightarrow Y=X/G
$$
be the separated complex analytic quotient, and write $q_*^G$ for invariant
direct image. We use the standard local finite-quotient structure: for every
$x\in X$, with $H=G_x$ and $y=q(x)$, an $H$-invariant slice $V$ can be
chosen such that $U=q(V)$ is an open neighbourhood of $y$, the quotient map
identifies $U$ with $V/H$ as a complex analytic space, the induced map
$p:V\to U$ is finite, and
$$
\mathcal O_U=(p_*\mathcal O_V)^H.
$$
If $\mathcal L$ is a coherent analytic $G$-$\mathcal O_X$-module,
then $q_*^G\mathcal L$ is a coherent $\mathcal O_Y$-module.
Consequently, if $E$ and $F$ are holomorphic vector bundles on $X$ with
compatible $G$-actions, then
$$
q_*^G\mathcal O(E^\vee\otimes F)
$$
is a coherent analytic sheaf on $Y$.
\end{lemma}

\begin{proof}
We work locally on $Y$. Fix $y\in Y$, choose $x\in q^{-1}(y)$, and put
$H=G_x$. Choose the finite analytic quotient chart $p:V\to U$ from the
statement; its underlying slice is the one supplied by Sheaves, Lemma
\ref{sheaves-lemma-discontinuous-group-action-neighbourhoods}. Since $p$ is
finite, it is proper. Grauert's direct-image theorem therefore shows that
$$
\mathcal M=p_*(\mathcal L|_V)
$$
is a coherent $\mathcal O_U$-module.

The sectionwise restriction isomorphism in the degree-zero case of Lemma
\ref{lemma-higher-invariant-direct-image-slice} is compatible with invariant
holomorphic functions. It therefore gives a canonical isomorphism of
$\mathcal O_U$-modules
$$
(q_*^G\mathcal L)|_U=\mathcal M^H.
$$
Indeed, if $a$ is a local section of
$\mathcal O_U=(p_*\mathcal O_V)^H$, then
$$
h(as)=h(a)h(s)=a h(s)
$$
for $h\in H$. Thus the $H$-action on $\mathcal M$ is
$\mathcal O_U$-linear. Since $H$ is finite, its invariant subsheaf is the
kernel of the $\mathcal O_U$-linear morphism
$$
\mathcal M\longrightarrow\bigoplus_{h\in H}\mathcal M,
\qquad
s\longmapsto(hs-s)_{h\in H}.
$$
The source and target are coherent, so their kernel is coherent by Modules,
Lemma \ref{modules-lemma-coherent-abelian}. Hence $q_*^G\mathcal L$ is
coherent near $y$, and therefore on $Y$.

For the final assertion, compatible actions on $E$ and $F$ induce the dual
and tensor action on $E^\vee\otimes F$. Its sheaf of holomorphic sections is
finite locally free over $\mathcal O_X$, hence coherent by Oka coherence for
complex spaces; see \cite[p. 16]{Grauert-direct-image}. It is therefore a
coherent analytic $G$-$\mathcal O_X$-module, so the result applies.
\end{proof}

\begin{lemma}[Coherent cohomology on a compact complex analytic space]
\label{lemma-compact-complex-analytic-coherent-cohomology-finite-dimensional}
\begin{reference}
\cite[Hauptsatz I, p. 57]{Grauert-direct-image}
\end{reference}
Let $Y$ be a compact complex analytic space and let $\mathcal M$ be a
coherent analytic $\mathcal O_Y$-module. Then
$$
H^p(Y,\mathcal M)
$$
is a finite dimensional complex vector space for every $p\geq 0$.
\end{lemma}

\begin{proof}
The structure morphism
$$
\pi:Y\longrightarrow \mathop{\rm pt}
$$
is proper. Grauert's proper direct-image theorem shows that
$R^p\pi_*\mathcal M$ is a coherent analytic module on the point. Its module
of sections is $H^p(Y,\mathcal M)$. Since
$\mathcal O_{\mathop{\rm pt}}=\mathbf C$, coherence on the point is precisely
finite dimensionality over $\mathbf C$.
\end{proof}

\begin{lemma}[Coherence of analytic sheaf Ext]
\label{lemma-coherent-analytic-sheaf-ext}
Let $X$ be a complex analytic space and let $\mathcal A,\mathcal B$ be
coherent analytic $\mathcal O_X$-modules. Then
$$
\SheafExt^q_{\mathcal O_X}(\mathcal A,\mathcal B)
$$
is a coherent analytic $\mathcal O_X$-module for every $q\geq 0$.
\end{lemma}

\begin{proof}
The assertion is local on $X$. Fix $q$ and a point of $X$. Oka coherence for
complex spaces and coherence of $\mathcal A$ allow us, after shrinking, to
construct as long a finite-free resolution as is needed in degree $q$,
$$
\mathcal P_{q+1}\longrightarrow\mathcal P_q\longrightarrow\cdots
\longrightarrow\mathcal P_0\longrightarrow\mathcal A\longrightarrow 0,
$$
where every $\mathcal P_i$ has finite rank and every successive syzygy is
coherent. No termination of this resolution is asserted.

Applying $\SheafHom_{\mathcal O_X}(-,\mathcal B)$ computes sheaf Ext through
degree $q$, since internal Hom from a finite locally free module is exact.
Every term of the resulting complex is a finite direct sum of $\mathcal B$
and hence is coherent. Its kernels, images, and cohomology sheaves are
coherent by Modules, Lemma \ref{modules-lemma-coherent-abelian}. Thus the
degree-$q$ cohomology sheaf is coherent near the chosen point. Since the point
and $q$ were arbitrary, the result follows.
\end{proof}

\begin{lemma}[Finiteness of equivariant analytic Ext over a compact quotient]
\label{lemma-compact-analytic-quotient-equivariant-ext-finite-dimensional}
\begin{reference}
\cite[Chapter V, \S 5.6, p. 219]{Tohoku}
\cite[Hauptsatz I, p. 57]{Grauert-direct-image}
\end{reference}
\begin{history}
Compactness is imposed on the quotient $Y=X/G$, not on $X$. The conclusion
is finite dimensionality of the vector space of extension classes, not
finiteness of the set. No global finiteness or properness of $q$, freeness or
finiteness of $G$, algebraization, or GAGA assertion is made.
\end{history}
Let $X$ be a complex analytic space and let $G$ act discontinuously on $X$
by analytic automorphisms. Let
$$
q:X\longrightarrow Y=X/G
$$
be the separated complex analytic quotient with the local finite-quotient
structure of Lemma
\ref{lemma-discontinuous-analytic-invariant-direct-image-coherent}, and
assume that $Y$ is compact. If $\mathcal A$ and $\mathcal B$ are coherent
analytic $G$-$\mathcal O_X$-modules, then
$$
\Ext^n_{\mathcal O_X,G}(X;\mathcal A,\mathcal B)
$$
is a finite dimensional complex vector space for every $n\geq 0$.

If $E$ and $F$ are holomorphic vector bundles on $X$ with compatible
$G$-actions, write $\operatorname{ExtBun}^1_G(E,F)$ for equivalence classes,
with fixed endpoints, of equivariant short exact sequences of holomorphic
vector bundles
$$
0\longrightarrow F\longrightarrow V\longrightarrow E\longrightarrow 0.
$$
This set has a canonical complex vector space structure, is finite
dimensional, and there are canonical isomorphisms
$$
\begin{aligned}
\operatorname{ExtBun}^1_G(E,F)
&\longrightarrow
\Ext^1_{\mathcal O_X,G}
  (X;\mathcal O(E),\mathcal O(F)) \\
&\longrightarrow
H^1\left(Y,q_*^G\mathcal O(E^\vee\otimes F)\right).
\end{aligned}
$$
\end{lemma}

\begin{proof}
Put
$$
\mathcal C^q=
\SheafExt^q_{\mathcal O_X}(\mathcal A,\mathcal B).
$$
This is a coherent analytic $\mathcal O_X$-module by Lemma
\ref{lemma-coherent-analytic-sheaf-ext}, and Lemma
\ref{lemma-equivariant-underlying-ext} gives it its canonical
$G$-$\mathcal O_X$-module structure. For $x\in X$, multiplication by the
finite stabilizer order $|G_x|$ is an automorphism of the complex vector
space $\mathcal C^q_x$. Lemmas
\ref{lemma-finite-group-cohomology-invertible-order} and
\ref{lemma-higher-invariant-direct-image-stalk} therefore give
$$
R^p q_*^G\mathcal C^q=0
\quad\text{for }p>0.
$$
The spectral sequence of Lemma
\ref{lemma-equivariant-sheaf-ext-spectral-sequence} consequently degenerates
and its edge map gives a canonical isomorphism
$$
q_*^G\mathcal C^q
\longrightarrow
\SheafExt^q_{\mathcal O_X,G}(\mathcal A,\mathcal B).
$$
The left hand side is coherent by Lemma
\ref{lemma-discontinuous-analytic-invariant-direct-image-coherent}; hence so
is the sheaf on the right.

The first spectral sequence of Lemma
\ref{lemma-equivariant-ext-spectral-sequences} has
$$
E_2^{p,q}=
H^p\left(Y,
\SheafExt^q_{\mathcal O_X,G}(\mathcal A,\mathcal B)\right)
\Longrightarrow
\Ext^{p+q}_{\mathcal O_X,G}(X;\mathcal A,\mathcal B).
$$
Every term on the $E_2$-page is finite dimensional by Lemma
\ref{lemma-compact-complex-analytic-coherent-cohomology-finite-dimensional}.
For each total degree there are only finitely many terms, and the abutment
has the finite filtration supplied by that spectral sequence. Its successive
quotients are subquotients of finite dimensional vector spaces. This proves
the first assertion.

We finish with the bundle assertion. The section sheaves $\mathcal O(E)$ and
$\mathcal O(F)$ are finite locally free coherent analytic
$G$-$\mathcal O_X$-modules. Conversely, in an extension of
$\mathcal O(E)$ by $\mathcal O(F)$ the quotient $\mathcal O(E)$ is locally
free, so the underlying module sequence splits locally. The middle term is
therefore finite locally free. Varying the rank in Lemmas
\ref{lemma-holomorphic-vector-bundles-locally-free-analytic-modules} and
\ref{lemma-equivariant-holomorphic-vector-bundles-locally-free-modules}, and
using the same local-matrix full-faithfulness argument for maps of different
ranks, identifies module extensions and their equivalences with equivariant
holomorphic bundle extensions with fixed endpoints.

Lemma \ref{lemma-equivariant-extension-classes-finite-locally-free} now gives
the first displayed isomorphism and identifies its target with
$$
H^1\left(X;G,
\mathcal O(E)^\vee\otimes_{\mathcal O_X}\mathcal O(F)\right).
$$
The local trivializations used in the bundle--module equivalence identify the
coefficient module canonically with $\mathcal O(E^\vee\otimes F)$. Since its
stalks are complex vector spaces, Lemma
\ref{lemma-orbit-space-equivariant-cohomology-stabilizer-orders} gives the
second displayed isomorphism. Finally,
$q_*^G\mathcal O(E^\vee\otimes F)$ is coherent by Lemma
\ref{lemma-discontinuous-analytic-invariant-direct-image-coherent}, so its
first cohomology is finite dimensional by Lemma
\ref{lemma-compact-complex-analytic-coherent-cohomology-finite-dimensional}.
\end{proof}


\begin{lemma}
\label{lemma-equivariant-ext-crossed-product-spectral-sequence}
\begin{reference}
\cite[Chapter V, \S 5.6, spectral sequence IV, p. 218]{Tohoku}
\end{reference}
Let $G$ act on a topological space $X$, let $R$ be a commutative ring
endowed with an action of $G$, and give the constant sheaf
$\underline R_X$ its induced $G$-ring structure. Let $U$ be the skew group
ring generated by $R$ and symbols $[g]$, $g\in G$, with
relations
$$
[g][h]=[gh],
\qquad
[g]r=g(r)[g].
$$
Let $M$ be an $R$-module with a compatible semilinear $G$-action, viewed
equivalently as a left $U$-module, and let $\underline M_X$ be the
corresponding constant $G$-$\underline R_X$-module. For every
$G$-$\underline R_X$-module $\mathcal B$ there is a canonical isomorphism
$$
\Hom_{\textit{Mod}_G(\underline R_X)}
(\underline M_X,\mathcal B)
=
\Hom_U(M,\Gamma(X,\mathcal B)).
$$
Moreover, there is an additive first-quadrant spectral sequence, functorial
in $\mathcal B$ from the $E_2$-page onward,
$$
{}^{\mathrm{IV}}E_2^{p,q}(\mathcal B)
=
\Ext_U^p\left(M,H^q(X,\mathcal B)\right)
\Longrightarrow
\Ext_{\underline R_X,G}^{p+q}
(X;\underline M_X,\mathcal B).
$$
Here $H^q(X,\mathcal B)$ carries its canonical semilinear $G$-action and
hence its canonical $U$-module structure. The spectral sequence induces a
finite filtration on every abutment group.
\end{lemma}

\begin{proof}
The constant-sheaf functor gives an exact functor
$$
U\text{-}\textit{Mod}
\longrightarrow
\textit{Mod}_G(\underline R_X),
\qquad
N\longmapsto\underline N_X.
$$
Indeed, its underlying functor is exact by Sites, Modules, Lemma
\ref{sites-modules-lemma-constant-exact}, and the $U$-module structure is
equivalently an $R$-module structure together with a compatible semilinear
$G$-action. The usual constant-sheaf adjunction of Sites, Example
\ref{sites-example-constant-sheaf} restricts to an adjunction
$$
\Hom_{\textit{Mod}_G(\underline R_X)}(\underline N_X,\mathcal B)
=
\Hom_U(N,\Gamma(X,\mathcal B)).
$$
Taking $N=M$ gives the first displayed isomorphism. Since the left adjoint
is exact, Homology, Lemma
\ref{homology-lemma-adjoint-preserve-injectives} shows that
$$
\Gamma(X,-):\textit{Mod}_G(\underline R_X)\longrightarrow U\text{-}\textit{Mod}
$$
takes injective objects to injective objects.

Apply Derived Categories, Lemma
\ref{derived-lemma-grothendieck-spectral-sequence} to the composite
$$
\textit{Mod}_G(\underline R_X)
\xrightarrow{\ \Gamma(X,-)\ }
U\text{-}\textit{Mod}
\xrightarrow{\ \Hom_U(M,-)\ }
\textit{Ab}.
$$
Both module categories have enough injectives, and the preceding paragraph
gives the required acyclicity hypothesis. By Injectives, Lemma
\ref{injectives-lemma-equivariant-injective-underlying}, an injective
resolution in $\textit{Mod}_G(\underline R_X)$ is an injective resolution
after forgetting $G$. Consequently the inner derived functors are ordinary
sheaf cohomology,
$$
R^q\Gamma(X,\mathcal B)=H^q(X,\mathcal B),
$$
with the induced semilinear $G$-action. The outer derived functors are
$\Ext_U^p(M,-)$, while the adjunction identifies the composite with the
degree-zero equivariant Hom functor. Definition
\ref{definition-equivariant-ext} identifies its derived functors with the
displayed abutment. The remaining assertions follow from the Grothendieck
spectral sequence.
\end{proof}


\begin{lemma}
\label{lemma-equivariant-ext-trivial-action-spectral-sequence}
\begin{reference}
\cite[Chapter V, \S 5.6, formulas (5.6.5 ter) and (5.6.11 ter),
p. 218]{Tohoku}
\end{reference}
\begin{history}
The cited source assumes that $G$ acts trivially on $M$, but it allows a
nontrivial action on $R$. Under those hypotheses the fixed subgroup $N^G$
of a semilinear $R$-module need not be stable under $R$, so the asserted
$R$-module structure and the resulting Ext groups are not defined in
general. We add the hypothesis that $G$ acts trivially on $R$. The source
also prints $A$ in place of the variable $B$ in (5.6.11 ter); the statement
below uses $B$.
\end{history}
Let $G$ act on a topological space $X$, let $R$ be a commutative ring on
which $G$ acts trivially, and give the constant sheaf $\underline R_X$ its
induced $G$-ring structure. Let $M$ be an $R$-module on which $G$ acts
trivially, and let $\underline M_X$ be the corresponding constant
$G$-$\underline R_X$-module. For every $G$-$\underline R_X$-module
$\mathcal B$ there is a canonical isomorphism
$$
\Hom_{\textit{Mod}_G(\underline R_X)}
(\underline M_X,\mathcal B)
=
\Hom_R\left(M,\Gamma_X^G(\mathcal B)\right).
$$
Moreover, the groups $H^q(X;G,\mathcal B)$ have canonical $R$-module
structures, and there is an additive first-quadrant spectral sequence,
functorial in $\mathcal B$ from the $E_2$-page onward,
$$
{}^{\mathrm V}E_2^{p,q}(\mathcal B)
=
\Ext_R^p\left(M,H^q(X;G,\mathcal B)\right)
\Longrightarrow
\Ext_{\underline R_X,G}^{p+q}
(X;\underline M_X,\mathcal B).
$$
The spectral sequence induces a finite filtration on every abutment group.
\end{lemma}

\begin{proof}
Since the action on $R$ is trivial, the skew group ring of Lemma
\ref{lemma-equivariant-ext-crossed-product-spectral-sequence} is the group
ring $U=R[G]$. The functor
$$
R\text{-}\textit{Mod}\longrightarrow U\text{-}\textit{Mod}
$$
which equips an $R$-module with the trivial $G$-action is exact and is left
adjoint to
$$
(-)^G:U\text{-}\textit{Mod}\longrightarrow R\text{-}\textit{Mod}.
$$
Indeed, for an $R$-module $L$ with trivial $G$-action and a $U$-module $N$
there is a canonical equality
$$
\Hom_U(L,N)=\Hom_R(L,N^G).
$$
Consequently $(-)^G$ takes injective $U$-modules to injective $R$-modules
by Homology, Lemma
\ref{homology-lemma-adjoint-preserve-injectives}.

Put
$$
F(\mathcal B)=\Gamma(X,\mathcal B)^G.
$$
This is a left exact functor from
$\textit{Mod}_G(\underline R_X)$ to $R\text{-}\textit{Mod}$. The proof of
Lemma \ref{lemma-equivariant-ext-crossed-product-spectral-sequence} shows
that $\Gamma(X,-)$ takes injective objects of
$\textit{Mod}_G(\underline R_X)$ to injective $U$-modules. Hence the
preceding paragraph shows that $F$ takes injective objects to injective
$R$-modules. The same lemma and the displayed adjunction, applied to $M$,
give
$$
\begin{aligned}
\Hom_{\textit{Mod}_G(\underline R_X)}
(\underline M_X,\mathcal B)
&=\Hom_U(M,\Gamma(X,\mathcal B))\\
&=\Hom_R(M,F(\mathcal B)),
\end{aligned}
$$
which proves the first assertion.

We next compare the $R$-linear right derived functors of $F$ with the
abelian-group-valued equivariant cohomology of Definition
\ref{definition-orbit-space-equivariant-cohomology}. Let
$$
v:\textit{Mod}_G(\underline R_X)
\longrightarrow
\textit{Mod}_G(\underline{\mathbf{Z}}_X)
$$
be the exact forgetful functor, and let $\mathcal I$ be injective in
$\textit{Mod}_G(\underline R_X)$. By Injectives, Lemma
\ref{injectives-lemma-equivariant-injective-underlying}, the underlying
$\underline R_X$-module of $\mathcal I$ is injective. It is therefore
flasque by Lemma \ref{lemma-injective-flasque}, and hence
$$
H^q(X,v\mathcal I)=0\quad\text{for }q>0
$$
by Lemma \ref{lemma-flasque-acyclic}.

On the other hand, as shown in the proof of Lemma
\ref{lemma-equivariant-ext-crossed-product-spectral-sequence},
$\Gamma(X,\mathcal I)$ is an injective $R[G]$-module. Therefore its higher
derived functors for invariants vanish in the category of $R$-$G$-modules.
By \'Etale Cohomology, Lemma
\ref{etale-cohomology-lemma-modules-abelian}, the same is true after
forgetting the $R$-module structure:
$$
H^p\left(G,\Gamma(X,\mathcal I)\right)=0
\quad\text{for }p>0.
$$
Apply the second spectral sequence of Lemma
\ref{lemma-orbit-space-equivariant-spectral-sequences} to $v\mathcal I$.
Its $E_2$-page is
$$
H^p\left(G,H^q(X,v\mathcal I)\right).
$$
The preceding two vanishing statements show that this page is zero away
from $(p,q)=(0,0)$. Hence
$$
H^n(X;G,v\mathcal I)=0\quad\text{for }n>0.
$$
Thus $v\mathcal I$ is right acyclic for invariant global sections.

Choose an injective resolution
$\mathcal B\to\mathcal I^\bullet$ in
$\textit{Mod}_G(\underline R_X)$, which has enough injectives by
Injectives, Lemma
\ref{injectives-lemma-equivariant-modules-enough-injectives}. Exactness of
$v$ makes $v\mathcal I^\bullet$ a resolution of $v\mathcal B$, and the
preceding paragraph makes every term right acyclic for invariant global
sections. Leray's acyclicity lemma, Derived Categories, Lemma
\ref{derived-lemma-leray-acyclicity}, therefore shows that this resolution
computes $H^q(X;G,v\mathcal B)$. After forgetting $R$, the complex
$F(\mathcal I^\bullet)$ is exactly
$$
\Gamma(X,v\mathcal I^\bullet)^G.
$$
Consequently there are canonical functorial isomorphisms
$$
R^qF(\mathcal B)\cong H^q(X;G,v\mathcal B)
$$
of underlying abelian groups. These isomorphisms give the right hand side,
for which we suppress $v$ from the notation, its canonical $R$-module
structure.

Finally apply Derived Categories, Lemma
\ref{derived-lemma-grothendieck-spectral-sequence} to the composite
$$
\textit{Mod}_G(\underline R_X)
\xrightarrow{\ F\ }
R\text{-}\textit{Mod}
\xrightarrow{\ \Hom_R(M,-)\ }
\textit{Ab}.
$$
The source category and $R\text{-}\textit{Mod}$ have enough injectives, and
$F$ takes injectives to injectives, so the required acyclicity hypothesis
holds. The inner derived functors are the $R$-modules
$H^q(X;G,\mathcal B)$ by the preceding paragraph, and the outer derived
functors are $\Ext_R^p(M,-)$. The first part of the proof identifies the
composite with the degree-zero equivariant Hom functor, whose derived
functors are
$$
\Ext_{\underline R_X,G}^{p+q}
(X;\underline M_X,\mathcal B)
$$
by Definition \ref{definition-equivariant-ext}. The remaining assertions
follow from the Grothendieck spectral sequence.
\end{proof}


\section{Equivariant cohomology with supports}
\label{section-equivariant-cohomology-supports}

\noindent
Throughout this section a $G$-sheaf means an abelian $G$-sheaf. No extension
to modules over a general equivariant sheaf of rings is asserted.

\begin{definition}
\label{definition-equivariant-cohomology-family-supports}
\begin{reference}
\cite[Chapter V, \S 5.7, p. 219]{Tohoku}
\end{reference}
Let $G$ act on a topological space $X$, and let $\Phi$ be a family of
supports on $X$ in the sense of Definition
\ref{definition-family-supports}. For an abelian $G$-sheaf $\mathcal A$ on
$X$, set
$$
\Gamma^G_\Phi(X,\mathcal A)
=
\{s\in\Gamma(X,\mathcal A)^G\mid
\text{Supp}_X(s)\in\Phi\}.
$$
Thus the right hand side means the intersection of the invariant global
sections with the sections supported in $\Phi$; it does not require
$\Gamma_\Phi(X,\mathcal A)$ to be stable under $G$.

This defines an additive left exact functor
$$
\Gamma^G_\Phi(X,-):
\textit{Mod}_G(\underline{\mathbf Z}_X)
\longrightarrow
\textit{Ab};
$$
see Lemma \ref{lemma-equivariant-family-supports-invariant-saturation}.
The source category has enough injectives by Injectives, Lemma
\ref{injectives-lemma-equivariant-modules-enough-injectives}. We define
$$
H^n_\Phi(X;G,\mathcal A)
=
R^n\Gamma^G_\Phi(X,-)(\mathcal A),
\qquad n\geq 0,
$$
and put these groups equal to zero in negative degrees. If $\Phi$ is the
family of all closed subsets of $X$, this recovers
$H^n(X;G,\mathcal A)$ of Definition
\ref{definition-orbit-space-equivariant-cohomology}.
\end{definition}

\begin{lemma}[Invariant saturation of a family of supports]
\label{lemma-equivariant-family-supports-invariant-saturation}
\begin{reference}
\cite[Chapter V, \S 5.7, p. 219]{Tohoku}
\end{reference}
Let $G$ act on $X$, and let $\Phi$ be a family of supports on $X$. Put
$$
\Phi^G=
\left\{
Z\subset X\ \middle|\
\begin{aligned}
&Z\text{ is closed, and there is }W\in\Phi\text{ with }Z\subset W,\\
&gW=W\text{ for every }g\in G
\end{aligned}
\right\}.
$$
Then $\Phi^G$ is a family of supports, it is stable under the action of
$G$, and $\Phi^G\subset\Phi$. For every abelian $G$-sheaf $\mathcal A$
there is a functorial equality
$$
\Gamma^G_\Phi(X,\mathcal A)
=
\Gamma^G_{\Phi^G}(X,\mathcal A).
$$
For the family on the right, supported global sections form a
$\mathbf Z[G]$-module and
$$
\Gamma^G_{\Phi^G}(X,-)
=
(-)^G\circ\Gamma_{\Phi^G}(X,-).
$$
In particular, $\Gamma^G_\Phi(X,-)$ is additive and left exact.
\end{lemma}

\begin{proof}
Every family of supports contains the empty set. Hence the empty set is in
$\Phi^G$. If $Z_i\subset W_i$ exhibits $Z_i\in\Phi^G$ for $i=1,2$, then
$W_1\cup W_2$ is an invariant member of $\Phi$ containing
$Z_1\cup Z_2$. Downward closure is immediate. This proves that $\Phi^G$
is a family of supports. If $Z\subset W$ is such an exhibition, then
$gZ\subset W$ for every $g\in G$, so the family is $G$-stable. Downward
closure of $\Phi$ also gives $\Phi^G\subset\Phi$.

Let $s\in\Gamma(X,\mathcal A)^G$. The action isomorphisms on the stalks
carry $s_x$ to $s_{gx}$. Thus
$$
g\text{Supp}_X(s)=\text{Supp}_X(s)
$$
for every $g\in G$. If $\text{Supp}_X(s)\in\Phi$, this support itself is
an invariant member of $\Phi$, and hence it lies in $\Phi^G$. The converse
follows from $\Phi^G\subset\Phi$. This proves the functorial equality.

Because $\Phi^G$ is $G$-stable, the group action on global sections
preserves $\Gamma_{\Phi^G}(X,\mathcal A)$. This gives the last displayed
factorization. The functor $\Gamma_{\Phi^G}(X,-)$ is additive and left
exact by Lemma \ref{lemma-cohomology-family-supports}, and taking
$G$-invariants is additive and left exact. The final assertion follows.
\end{proof}

\begin{lemma}[Equivariant cohomology with supports on the orbit space]
\label{lemma-orbit-space-equivariant-support-spectral-sequence}
\begin{reference}
\cite[Chapter V, \S 5.7, pp. 219--220]{Tohoku}
\end{reference}
\begin{history}
The support subscript on the abutment in the cited source is essential. It
is the subscript restored in corrected-source decision
\texttt{TOH-SRC-P220-004}; omitting it would contradict the degree-zero
functor identity proved below.
\end{history}
Let $G$ act on $X$, let $Y=X/G$, and let $q:X\to Y$ be the orbit map. Let
$\Psi$ be a family of supports on $Y$, and put
$$
q^\dagger\Psi=
\{Z\subset X\mid Z\text{ is closed and }
\overline{q(Z)}\in\Psi\}.
$$
The closure in this formula is part of the definition; the image under an
arbitrary orbit map of a closed subset need not be closed. Then:
\begin{enumerate}
\item $q^\dagger\Psi$ is a family of supports on $X$, and it equals its
invariant saturation $(q^\dagger\Psi)^G$.
\item If $\mathcal A$ is an abelian $G$-sheaf on $X$ and an invariant
section $s\in\Gamma(X,\mathcal A)^G$ corresponds to
$\widetilde s\in\Gamma(Y,q_*^G\mathcal A)$, then
$$
\text{Supp}_Y(\widetilde s)
=
\overline{q(\text{Supp}_X(s))}.
$$
Consequently there is a natural equality of left exact functors
$$
\Gamma^G_{q^\dagger\Psi}(X,-)
=
\Gamma_\Psi(Y,-)\circ q_*^G.
$$
\item For every abelian $G$-sheaf $\mathcal A$ there is a functorial
first-quadrant spectral sequence
$$
E_2^{p,q}
=
H^p_\Psi\bigl(Y,\mathcal H^q(G,\mathcal A)\bigr)
\Longrightarrow
H^{p+q}_{q^\dagger\Psi}(X;G,\mathcal A).
$$
It is additive and functorial from the $E_2$-page onward and gives a
finite filtration on every abutment group.
\end{enumerate}
\end{lemma}

\begin{proof}
The empty set belongs to $q^\dagger\Psi$. If $Z'\subset Z$ are closed and
$Z\in q^\dagger\Psi$, then
$$
\overline{q(Z')}\subset\overline{q(Z)},
$$
so downward closure of $\Psi$ gives $Z'\in q^\dagger\Psi$. For two closed
subsets $Z_1,Z_2\subset X$, we have
$$
\overline{q(Z_1\cup Z_2)}
=
\overline{q(Z_1)}\cup\overline{q(Z_2)}.
$$
Thus closure under finite unions follows from the corresponding property
of $\Psi$. This proves that $q^\dagger\Psi$ is a family of supports.

Let $Z\in q^\dagger\Psi$, put $C=\overline{q(Z)}$, and set
$$
W=q^{-1}(C).
$$
The subset $W$ is closed, $G$-invariant, and contains $Z$. Since $q$ is
surjective, $q(W)=C$, and hence $W\in q^\dagger\Psi$. Therefore every
member of $q^\dagger\Psi$ is contained in an invariant member of the same
family. Lemma
\ref{lemma-equivariant-family-supports-invariant-saturation} gives
$(q^\dagger\Psi)^G=q^\dagger\Psi$.

We prove the support identity. Put $S=\text{Supp}_X(s)$. For $y\in Y$, the
definition of a stalk and of invariant direct image gives
\begin{align*}
\widetilde s_y=0
&\Longleftrightarrow
\text{there is an open }V\ni y\text{ such that }
s|_{q^{-1}(V)}=0\\
&\Longleftrightarrow
\text{there is an open }V\ni y\text{ such that }
V\cap q(S)=\emptyset\\
&\Longleftrightarrow
y\notin\overline{q(S)}.
\end{align*}
Hence $\text{Supp}_Y(\widetilde s)=\overline{q(S)}$. By Definition
\ref{definition-family-supports}, the section $\widetilde s$ is supported
in $\Psi$ if and only if $s$ is supported in $q^\dagger\Psi$. This proves
the functor identity.

Set
$$
F=q_*^G:
\textit{Mod}_G(\underline{\mathbf Z}_X)
\longrightarrow\textit{Ab}(Y)
\quad\text{and}\quad
T=\Gamma_\Psi(Y,-):\textit{Ab}(Y)\longrightarrow\textit{Ab}.
$$
These are additive left exact functors by Modules, Lemma
\ref{modules-lemma-invariant-direct-image-G-O-modules} and Lemma
\ref{lemma-cohomology-family-supports}. If $\mathcal I$ is an injective
abelian $G$-sheaf on $X$, then $F(\mathcal I)=q_*^G\mathcal I$ is flasque
by Injectives, Lemma
\ref{injectives-lemma-equivariant-injective-sections-invariants}. It is
therefore right acyclic for $T$ by Lemma
\ref{lemma-family-supports-flasque-soft-acyclic}.

The Grothendieck spectral sequence, Derived Categories, Lemma
\ref{derived-lemma-grothendieck-spectral-sequence}, now applies to
$T\circ F$. Definition
\ref{definition-orbit-space-equivariant-cohomology} identifies
$R^qF(\mathcal A)$ with $\mathcal H^q(G,\mathcal A)$, while Lemma
\ref{lemma-cohomology-family-supports} identifies $R^pT$ with
$H^p_\Psi(Y,-)$. The degree-zero functor identity above and Definition
\ref{definition-equivariant-cohomology-family-supports} identify the
derived composite with
$H^{p+q}_{q^\dagger\Psi}(X;G,\mathcal A)$. This proves (3).
\end{proof}

\begin{remark}
\label{remark-equivariant-support-second-spectral-sequence}
\begin{reference}
\cite[Chapter V, \S 5.7, p. 220]{Tohoku}
\end{reference}
For the invariant family $q^\dagger\Psi$, there is also a degree-zero
factorization
$$
\Gamma^G_{q^\dagger\Psi}(X,-)
=
(-)^G\circ\Gamma_{q^\dagger\Psi}(X,-).
$$
This factorization alone does not produce a second Grothendieck spectral
sequence. One would have to prove, in particular, that
$\Gamma_{q^\dagger\Psi}(X,\mathcal I)$ is right acyclic for $(-)^G$ for
every injective abelian $G$-sheaf $\mathcal I$; the preceding results do not
give that assertion. This is only a failure of the present hypotheses to
verify the Grothendieck acyclicity condition; it is not a nonexistence
claim. Thus we do not assert a spectral sequence with a putative initial page
$$
H^p\bigl(G,H^q_{q^\dagger\Psi}(X,\mathcal A)\bigr).
$$
The cited source next assumes, exactly, that every member of $\Psi$ has a
neighbourhood which is again a member of $\Psi$. Under that hypothesis it
passes through extension by zero and a filtered system of neighbourhoods in
order to construct a different second spectral sequence. We record the
hypothesis here because it marks the end of the obstruction statement, but
we do not yet identify supported cohomology with that filtered colimit,
construct the resulting spectral sequence, or interchange group cohomology
with a filtered colimit. Those are separate assertions requiring the
additional work carried out after the cited passage.
\end{remark}

\begin{lemma}[Filtered extension by zero for equivariant supports]
\label{lemma-equivariant-support-filtered-extension-by-zero}
\begin{reference}
\cite[Chapter V, \S 5.7, formulae (5.7.5) and (5.7.6),
p. 220]{Tohoku}
\end{reference}
\begin{history}
The cited passage assumes that every member of the support family on the
orbit space has a neighbourhood in that family and then asserts the derived
filtered-colimit formula (5.7.6). The neighbourhood hypothesis proves the
degree-zero formula (5.7.5), but by itself it does not prove (5.7.6).
Indeed, the filtered-colimit theorem for right derived functors produces
the right derived functors of the composites
$\mathcal A\mapsto\Gamma_X^G(\mathcal A_U)$. Identifying these with
$H^n(X;G,\mathcal A_U)$ requires the extension-by-zero acyclicity stated
explicitly below. Exactness of extension by zero alone does not supply that
acyclicity. Moreover, the ordinary supported-cohomology comparison invoked
later in the cited passage follows from Corollary 1 to Proposition 3.10.1
under the additional hypothesis that the induced family on $X$ is
paracompactifying. We keep these two requirements separate and make no
claim that either follows from the neighbourhood hypothesis.
\end{history}
Let $G$ act on a topological space $X$, let $Y=X/G$, and let
$q:X\to Y$ be the orbit map. Let $\Psi$ be a family of supports on $Y$,
put
$$
\Phi=q^\dagger\Psi,
$$
as in Lemma
\ref{lemma-orbit-space-equivariant-support-spectral-sequence}, and assume
that every $C\in\Psi$ has a closed neighbourhood belonging to $\Psi$; in
other words, there is a $B\in\Psi$ such that
$$
C\subset\mathop{\rm Int}_Y(B).
$$
Let $\mathcal U_\Psi$ be the set of open subsets $U\subset Y$ such that
$\overline U\in\Psi$, ordered by inclusion. For $U\in\mathcal U_\Psi$ put
$$
X_U=q^{-1}(U),\qquad j_U:X_U\longrightarrow X,
\qquad
E_U(\mathcal A)=\mathcal A_U=(j_U)_!(\mathcal A|_{X_U}).
$$
Then:
\begin{enumerate}
\item The set $\mathcal U_\Psi$ is directed. Each $E_U$ is an exact
endofunctor of $\textit{Mod}_G(\underline{\mathbf Z}_X)$, and for
$U\subset V$ there is a functorial monomorphism
$E_U\to E_V$.
\item For every abelian $G$-sheaf $\mathcal A$ on $X$, there is a
functorial equality
$$
\Gamma^G_\Phi(X,\mathcal A)
=
\colim_{U\in\mathcal U_\Psi}\Gamma_X^G(\mathcal A_U).
$$
\item Assume the following extension-by-zero acyclicity condition:
for every $U\in\mathcal U_\Psi$ and every injective abelian $G$-sheaf
$\mathcal I$ on $X$, the sheaf $\mathcal I_U$ is right acyclic for
$\Gamma_X^G$, or equivalently
\begin{equation}
\label{equation-equivariant-support-extension-zero-acyclic}
H^n(X;G,\mathcal I_U)=0
\quad\text{for all }n>0.
\end{equation}
Then the canonical maps
$$
\colim_{U\in\mathcal U_\Psi}H^n(X;G,\mathcal A_U)
\longrightarrow
H^n_\Phi(X;G,\mathcal A)
$$
are isomorphisms for every $n\in\mathbf Z$. In nonnegative degrees these
isomorphisms form an isomorphism of cohomological $\delta$-functors.
\item If, in addition, $\Phi$ is paracompactifying, then the canonical
maps of $G$-modules
$$
\colim_{U\in\mathcal U_\Psi}H^q(X,\mathcal A_U)
\longrightarrow
H^q_\Phi(X,\mathcal A)
$$
are isomorphisms for every $q\in\mathbf Z$, functorially in $\mathcal A$
and compatibly with connecting maps.
\end{enumerate}
\end{lemma}

\begin{proof}
The empty set belongs to $\mathcal U_\Psi$. If
$U,V\in\mathcal U_\Psi$, then
$$
\overline{U\cup V}=\overline U\cup\overline V\in\Psi,
$$
so $U\cup V\in\mathcal U_\Psi$. Thus this set is directed.
The open subset $X_U$ is $G$-invariant, and hence extension by zero gives
$\mathcal A_U$ its natural $G$-sheaf structure. Restriction and extension
by zero are exact; for the latter assertion see Lemma
\ref{lemma-extension-by-zero-locally-closed}. If $U\subset V$, the
canonical map $\mathcal A_U\to\mathcal A_V$ is a monomorphism, as is seen
on stalks. This proves (1).

The counit maps $\mathcal A_U\to\mathcal A$ give an injective map
$$
\colim_{U\in\mathcal U_\Psi}\Gamma(X,\mathcal A_U)
\longrightarrow
\Gamma(X,\mathcal A),
$$
because all transition maps are injective. A section in the source which
comes from $U$ has support $S\subset X_U$. Consequently
$$
\overline{q(S)}\subset\overline U\in\Psi,
$$
so $S\in\Phi$.

Conversely, let $s\in\Gamma_\Phi(X,\mathcal A)$, put
$S=\text{Supp}_X(s)$, and set $C=\overline{q(S)}$. Thus $C\in\Psi$.
Choose a closed neighbourhood $B\in\Psi$ of $C$ and put
$U=\mathop{\rm Int}_Y(B)$. Then
$$
C\subset U,
\qquad
\overline U\subset B,
$$
and hence $U\in\mathcal U_\Psi$ and $S\subset X_U$. The section $s$
therefore comes from a unique section of $\mathcal A_U$. We have proved a
functorial equality of $G$-modules
\begin{equation}
\label{equation-equivariant-support-sections-colimit}
\Gamma_\Phi(X,\mathcal A)
=
\colim_{U\in\mathcal U_\Psi}\Gamma(X,\mathcal A_U).
\end{equation}
The system on the right has injective transition maps. It may therefore be
viewed as a directed union of $G$-submodules. Taking invariants commutes
with such a union for an arbitrary group $G$: an element of the union is
fixed if and only if it is already fixed in any member containing it.
Taking invariants in
(\ref{equation-equivariant-support-sections-colimit}) proves (2).

For $U\in\mathcal U_\Psi$ consider the additive functor
$$
F_U:
\textit{Mod}_G(\underline{\mathbf Z}_X)\longrightarrow\textit{Ab},
\qquad
F_U(\mathcal A)=\Gamma_X^G(E_U(\mathcal A)).
$$
These functors form a directed system, and part (2) says that their
pointwise colimit is $\Gamma_\Phi^G(X,-)$. Let
$\mathcal A\to\mathcal I^\bullet$ be a nonnegative injective resolution in
$\textit{Mod}_G(\underline{\mathbf Z}_X)$. Such resolutions exist by
Injectives, Lemma
\ref{injectives-lemma-equivariant-modules-enough-injectives}.
Exactness of $E_U$ gives a resolution
$$
\mathcal A_U\longrightarrow E_U(\mathcal I^\bullet).
$$
Under hypothesis
(\ref{equation-equivariant-support-extension-zero-acyclic}) its terms are
right acyclic for $\Gamma_X^G$.
Leray's acyclicity lemma, Derived Categories, Lemma
\ref{derived-lemma-leray-acyclicity}, therefore gives
\begin{equation}
\label{equation-equivariant-support-derived-extension-zero}
R^nF_U(\mathcal A)
=
H^n(X;G,\mathcal A_U).
\end{equation}
The source category has enough injectives and abelian groups satisfy AB5.
Derived Categories, Lemma
\ref{derived-lemma-filtered-colimit-right-derived-functors}, applied to the
system $(F_U)$, now gives
$$
\begin{aligned}
\colim_{U\in\mathcal U_\Psi}H^n(X;G,\mathcal A_U)
&=\colim_{U\in\mathcal U_\Psi}R^nF_U(\mathcal A)\\
&=R^n\Gamma_\Phi^G(X,\mathcal A)\\
&=H^n_\Phi(X;G,\mathcal A).
\end{aligned}
$$
The same lemma supplies functoriality and compatibility with connecting
maps. This proves (3).

Assume finally that $\Phi$ is paracompactifying. Let $\mathcal V_\Phi$ be
the directed set of open subsets $V\subset X$ such that
$\overline V\in\Phi$. Lemma
\ref{lemma-family-supports-colimit-extension-by-zero} gives
\begin{equation}
\label{equation-equivariant-support-ordinary-all-opens-colimit}
\colim_{V\in\mathcal V_\Phi}H^q(X,\mathcal A_V)
=H^q_\Phi(X,\mathcal A).
\end{equation}
The opens $X_U$, for $U\in\mathcal U_\Psi$, form a cofinal subsystem of
$\mathcal V_\Phi$. First,
$$
q\left(\overline{q^{-1}(U)}\right)\subset\overline U,
$$
so $\overline{q^{-1}(U)}\in\Phi$. Conversely, let
$V\in\mathcal V_\Phi$ and put
$$
C=\overline{q(\overline V)}\in\Psi.
$$
Choose a closed neighbourhood $B\in\Psi$ of $C$ and put
$U=\mathop{\rm Int}_Y(B)$. Then $U\in\mathcal U_\Psi$ and
$$
V\subset q^{-1}(U).
$$
Thus cofinality changes the left hand side of
(\ref{equation-equivariant-support-ordinary-all-opens-colimit}) into the
colimit asserted in part (4).
Every $X_U$ is $G$-invariant, all transition maps are equivariant, and the
isomorphism
(\ref{equation-equivariant-support-ordinary-all-opens-colimit}) is natural
for the action isomorphisms defining the $G$-sheaf $\mathcal A$. Hence the
cofinality isomorphism is an isomorphism of $G$-modules. This proves (4).
\end{proof}


\begin{lemma}[Filtered second spectral sequence with supports]
\label{lemma-equivariant-support-filtered-second-spectral-sequence}
\begin{reference}
\cite[Chapter V, \S 5.7, formulae (5.7.7) and (5.7.7 bis),
p. 220]{Tohoku}
\end{reference}
\begin{history}
The three conditions listed after formula (5.7.7 bis) in the cited source
are three alternative sufficient conditions for commuting group cohomology
with the indicated filtered colimit; they are not simultaneous hypotheses.
They concern the identification of the $E_2$-page only and do not imply the
extension-by-zero acyclicity
(\ref{equation-equivariant-support-extension-zero-acyclic}) of Lemma
\ref{lemma-equivariant-support-filtered-extension-by-zero}. Accordingly,
the latter acyclicity remains an explicit hypothesis here.
\end{history}
In the situation of Lemma
\ref{lemma-equivariant-support-filtered-extension-by-zero}, assume its
extension-by-zero acyclicity condition
(\ref{equation-equivariant-support-extension-zero-acyclic}). For every
abelian $G$-sheaf $\mathcal A$ on $X$, there is an additive first-quadrant
spectral sequence,
functorial in $\mathcal A$ from the $E_2$-page onward, with
\begin{equation}
\label{equation-equivariant-support-filtered-second-spectral-sequence}
{}^{\mathrm {II}}E_2^{p,q}(\mathcal A)
=
\colim_{U\in\mathcal U_\Psi}
H^p\bigl(G,H^q(X,\mathcal A_U)\bigr)
\quad\Longrightarrow\quad
H^{p+q}_\Phi(X;G,\mathcal A).
\end{equation}
It induces a finite filtration on every abutment group.

Assume in addition that $\Phi$ is paracompactifying. Then the $E_2$-page
in
(\ref{equation-equivariant-support-filtered-second-spectral-sequence}) has
the usual form
\begin{equation}
\label{equation-equivariant-support-usual-second-page}
{}^{\mathrm {II}}E_2^{p,q}(\mathcal A)
=
H^p\bigl(G,H^q_\Phi(X,\mathcal A)\bigr)
\end{equation}
if any one of the following three conditions holds:
\begin{enumerate}
\item $G$ is finite;
\item for every $q\geq 0$ and $U\in\mathcal U_\Psi$, the action of $G$
on $H^q(X,\mathcal A_U)$ is trivial, and every group
$H_r(G,\mathbf Z)$ is finitely generated;
\item for every $q\geq 0$, there is a cofinal subsystem of
$\mathcal U_\Psi$ on which all transition maps
$$
H^q(X,\mathcal A_U)\longrightarrow H^q(X,\mathcal A_V)
$$
are isomorphisms.
\end{enumerate}
\end{lemma}

\begin{proof}
For each $U\in\mathcal U_\Psi$, apply the second spectral sequence of
Lemma \ref{lemma-orbit-space-equivariant-spectral-sequences} to the
abelian $G$-sheaf $\mathcal A_U$. This gives
\begin{equation}
\label{equation-equivariant-support-local-second-spectral-sequence}
{}^{\mathrm {II}}E_2^{p,q}(\mathcal A,U)
=H^p\bigl(G,H^q(X,\mathcal A_U)\bigr)
\quad\Longrightarrow\quad
H^{p+q}(X;G,\mathcal A_U).
\end{equation}
For $U\subset V$, the equivariant monomorphism
$\mathcal A_U\to\mathcal A_V$ and functoriality of that lemma give a
morphism from the spectral sequence for $U$ to the one for $V$. Hence the
spectral sequences
(\ref{equation-equivariant-support-local-second-spectral-sequence}) form a
directed system from their $E_2$-pages onward.

For $r\geq 2$, take the filtered colimit of the $r$th pages and their
differentials. Filtered colimits of abelian groups are exact by Algebra,
Lemma \ref{algebra-lemma-directed-colimit-exact}; consequently they commute
with kernels, images, and the cohomology of complexes. Therefore
$$
H\left(\colim_U{}^{\mathrm {II}}E_r(\mathcal A,U),
\colim_U d_r\right)
=
\colim_U{}^{\mathrm {II}}E_{r+1}(\mathcal A,U).
$$
These colimit pages and differentials thus form a spectral sequence, and
its $E_2$-page is the one displayed in
(\ref{equation-equivariant-support-filtered-second-spectral-sequence}).

We spell out convergence. In a first-quadrant cohomological spectral
sequence, the incoming differential at $(p,q)$ is zero for $r>p$, and the
outgoing differential is zero for $r>q+1$. Hence the terms at $(p,q)$
stabilize, uniformly in $U$, once
$$
r>\max\{p,q+1\}.
$$
It follows that the stable term of the colimit spectral sequence is the
colimit of the stable terms. For a fixed total degree $n$, let
$F_U^pH^n(X;G,\mathcal A_U)$ be the finite filtration supplied by
(\ref{equation-equivariant-support-local-second-spectral-sequence}),
and set
$$
F^p\left(\colim_UH^n(X;G,\mathcal A_U)\right)
=
\colim_UF_U^pH^n(X;G,\mathcal A_U).
$$
Exactness of filtered colimits identifies the successive quotients with
the corresponding colimits of the successive quotients, hence with
${}^{\mathrm {II}}E_\infty^{p,n-p}(\mathcal A)$. The filtration is finite,
since $0\leq p\leq n$. Part (3) of Lemma
\ref{lemma-equivariant-support-filtered-extension-by-zero} identifies its
underlying abutment with $H^n_\Phi(X;G,\mathcal A)$. This proves convergence
and the filtration assertion. Additivity and functoriality are inherited
from
(\ref{equation-equivariant-support-local-second-spectral-sequence}) and
from filtered colimits.

It remains to prove the three interchange criteria. Put
$$
M_U^q=H^q(X,\mathcal A_U).
$$
For every $p,q$ there is a canonical comparison map
\begin{equation}
\label{equation-equivariant-support-group-colimit-comparison}
\colim_UH^p(G,M_U^q)
\longrightarrow
H^p\left(G,\colim_UM_U^q\right).
\end{equation}

If $G$ is finite, its bar resolution has in every degree a finitely
generated free $\mathbf Z[G]$-module. Applying
$\Hom_{\mathbf Z[G]}(-,M)$ gives a cochain complex whose terms commute with
filtered colimits. Exactness of filtered colimits then shows that its
cohomology commutes with filtered colimits. Thus
(\ref{equation-equivariant-support-group-colimit-comparison}) is an
isomorphism in case (1).

Assume (2). For a trivial $G$-module $M$, the universal coefficient
sequence gives, naturally in $M$,
\begin{equation}
\label{equation-equivariant-support-group-universal-coefficient}
0\longrightarrow
\Ext^1_{\mathbf Z}(H_{p-1}(G,\mathbf Z),M)
\longrightarrow H^p(G,M)
\longrightarrow
\Hom_{\mathbf Z}(H_p(G,\mathbf Z),M)
\longrightarrow0,
\end{equation}
where the first term is zero for $p=0$. A finitely generated abelian group
is finitely presented, and both $\Hom_{\mathbf Z}(N,-)$ and
$\Ext^1_{\mathbf Z}(N,-)$ commute with filtered colimits when $N$ is
finitely generated. Taking filtered colimits of
(\ref{equation-equivariant-support-group-universal-coefficient}), and using
their exactness, proves that
(\ref{equation-equivariant-support-group-colimit-comparison}) is an
isomorphism in case (2).

Assume (3). On the stated cofinal subsystem every map $M_U^q\to M_V^q$
is an isomorphism. The map from any one of its terms to its colimit is
therefore an isomorphism, and the same is true after applying
$H^p(G,-)$. Thus
(\ref{equation-equivariant-support-group-colimit-comparison}) is an
isomorphism in case (3), without any condition on $G$.

Under any one of the three conditions,
(\ref{equation-equivariant-support-group-colimit-comparison}) identifies
the $E_2$-page in
(\ref{equation-equivariant-support-filtered-second-spectral-sequence}) with
$$
H^p\left(G,\colim_UM_U^q\right).
$$
Part (4) of Lemma
\ref{lemma-equivariant-support-filtered-extension-by-zero} identifies the
inner colimit, as a $G$-module, with $H^q_\Phi(X,\mathcal A)$. This gives
(\ref{equation-equivariant-support-usual-second-page}) and completes the
proof.
\end{proof}

\section{Ext with support}
\label{section-ext-support}

\begin{definition}
\label{definition-ext-with-support}
\begin{reference}
\cite[Chapter IV, \S 4.2, Remark 1, p. 190]{Tohoku}
\end{reference}
\begin{history}
The cited source uses the term ``antifilter of closed subsets'' for what
we call a family of supports in Definition
\ref{definition-family-supports}.
\end{history}
Let $(X, \mathcal{O}_X)$ be a ringed space and let $\Phi$ be a family of
supports on $X$. If $K, L \in D(\mathcal{O}_X)$ and
$R\SheafHom(K, L)$ is bounded below, then we set
$$
R\Hom_{X, \Phi}(K, L) =
R\Gamma_\Phi(X, R\SheafHom(K, L))
$$
and
$$
\Ext^n_{X, \Phi}(K, L) = H^n(R\Hom_{X, \Phi}(K, L)).
$$
If $Z \subset X$ is closed, then for the family $\Phi_Z$ of Definition
\ref{definition-family-supports} we write
$$
R\Hom_{X, Z}(K, L) = R\Hom_{X, \Phi_Z}(K, L)
\quad\text{and}\quad
\Ext^n_{X, Z}(K, L) = \Ext^n_{X, \Phi_Z}(K, L).
$$
\end{definition}

\begin{lemma}
\label{lemma-ext-with-support}
\begin{reference}
\cite[Chapter IV, Theorem 4.2.1, p. 188, and \S 4.2, Remark 1,
p. 190]{Tohoku}
\end{reference}
Let $(X, \mathcal{O}_X)$ be a ringed space and let $\Phi$ be a family of
supports on $X$. The following statements hold.
\begin{enumerate}
\item If $\mathcal{F}$ and $\mathcal{G}$ are $\mathcal{O}_X$-modules,
then $\Ext^n_{X, \Phi}(\mathcal{F}, \mathcal{G})$, for $n \geq 0$, is the
$n$th right derived functor, with respect to $\mathcal{G}$, of
$$
\mathcal{G} \longmapsto
\Gamma_\Phi(X, \SheafHom_{\mathcal{O}_X}(\mathcal{F}, \mathcal{G})).
$$
\item If $K$ is bounded above and $L$ is bounded below, then there is a
bounded spectral sequence
$$
E_2^{p, q} =
H_\Phi^p(X, \SheafExt^q(K, L))
\quad\Longrightarrow\quad
\Ext^{p + q}_{X, \Phi}(K, L).
$$
\end{enumerate}
\end{lemma}

\begin{proof}
Let $\mathcal{I}$ be an injective $\mathcal{O}_X$-module. We first show
that $\SheafHom_{\mathcal{O}_X}(\mathcal{F}, \mathcal{I})$ is flasque.
Let $j : U \to X$ be the inclusion of an open subspace. A morphism
$\mathcal{F}|_U \to \mathcal{I}|_U$ corresponds to a morphism
$j_!(\mathcal{F}|_U) \to \mathcal{I}$ by
Sheaves, Lemma \ref{sheaves-lemma-j-shriek-modules}. The canonical map
$j_!(\mathcal{F}|_U) \to \mathcal{F}$ is injective, as may be checked on
stalks using the same lemma. Hence the morphism extends to
$\mathcal{F} \to \mathcal{I}$ by injectivity of $\mathcal{I}$. This proves
the claim.

Lemma \ref{lemma-family-supports-flasque-soft-acyclic} shows that a
flasque module is right acyclic for $\Gamma_\Phi(X, -)$. If
$\mathcal{G} \to \mathcal{I}^\bullet$ is a bounded below injective
resolution, then this complex is K-injective by
Derived Categories, Lemma
\ref{derived-lemma-bounded-below-injectives-K-injective}. Thus
$\SheafHom_{\mathcal{O}_X}(\mathcal{F}, \mathcal{I}^\bullet)$ represents
$R\SheafHom(\mathcal{F}, \mathcal{G})$ by the construction in
Section \ref{section-internal-hom}. Its terms are right acyclic for
$\Gamma_\Phi(X, -)$ by the preceding paragraph. Leray's acyclicity lemma,
Derived Categories, Lemma \ref{derived-lemma-leray-acyclicity}, now shows
that applying $\Gamma_\Phi(X, -)$ to this complex computes
$R\Hom_{X, \Phi}(\mathcal{F}, \mathcal{G})$. This proves (1).

For (2), choose a bounded above complex $\mathcal{K}^\bullet$
representing $K$ and a bounded below complex $\mathcal{I}^\bullet$ of
injectives representing $L$. The latter is K-injective by
Derived Categories, Lemma
\ref{derived-lemma-bounded-below-injectives-K-injective}. Hence
$$
\SheafHom^\bullet(\mathcal{K}^\bullet, \mathcal{I}^\bullet)
$$
represents $R\SheafHom(K, L)$ by
Section \ref{section-internal-hom}, and it is bounded below. Apply
Derived Categories, Lemma \ref{derived-lemma-two-ss-complex-functor} to
the left exact functor $\Gamma_\Phi(X, -)$ and this bounded below complex.
Its second spectral sequence is exactly the spectral sequence in (2),
and its abutment is the cohomology of
$R\Gamma_\Phi(X, R\SheafHom(K, L))$.
\end{proof}





\section{Global derived hom}
\label{section-global-RHom}

\noindent
Let $(X, \mathcal{O}_X)$ be a ringed space. Let $K, L \in D(\mathcal{O}_X)$.
Using the construction of the internal hom in the derived category we
obtain a well defined object
$$
R\Hom_X(K, L) = R\Gamma(X, R\SheafHom(K, L))
$$
in $D(\Gamma(X, \mathcal{O}_X))$. We will sometimes write
$R\Hom_{\mathcal{O}_X}(K, L)$ for this object.
By Lemma \ref{lemma-section-RHom-over-U}
we have
$$
H^0(R\Hom_X(K, L)) = \Hom_{D(\mathcal{O}_X)}(K, L),
\quad
H^p(R\Hom_X(K, L)) = \Ext_{D(\mathcal{O}_X)}^p(K, L)
$$
If $f : Y \to X$ is a morphism of ringed spaces, then there is
a canonical map
$$
R\Hom_X(K, L) \longrightarrow R\Hom_Y(Lf^*K, Lf^*L)
$$
in $D(\Gamma(X, \mathcal{O}_X))$ by taking global sections of the map
defined in Remark \ref{remark-prepare-fancy-base-change}.








\section{Glueing complexes}
\label{section-glueing-complexes}

\noindent
We can glue complexes! More precisely, in certain circumstances we can
glue locally given objects of the derived category to a global object.
We first prove some easy cases and then we'll prove the very general
\cite[Theorem 3.2.4]{BBD}
in the setting of topological spaces and open coverings.

\begin{lemma}
\label{lemma-glue}
Let $(X, \mathcal{O}_X)$ be a ringed space. Let $X = U \cup V$ be
the union of two open subspaces of $X$. Suppose given
\begin{enumerate}
\item an object $A$ of $D(\mathcal{O}_U)$,
\item an object $B$ of $D(\mathcal{O}_V)$, and
\item an isomorphism $c : A|_{U \cap V} \to B|_{U \cap V}$.
\end{enumerate}
Then there exists an object $F$ of $D(\mathcal{O}_X)$
and isomorphisms $f : F|_U \to A$, $g : F|_V \to B$ such
that $c = g|_{U \cap V} \circ f^{-1}|_{U \cap V}$.
Moreover, given
\begin{enumerate}
\item an object $E$ of $D(\mathcal{O}_X)$,
\item a morphism $a : A \to E|_U$ of $D(\mathcal{O}_U)$,
\item a morphism $b : B \to E|_V$ of $D(\mathcal{O}_V)$,
\end{enumerate}
such that
$$
a|_{U \cap V}  = b|_{U \cap V} \circ c.
$$
Then there exists a morphism $F \to E$ in $D(\mathcal{O}_X)$
whose restriction to $U$ is $a \circ f$
and whose restriction to $V$ is $b \circ g$.
\end{lemma}

\begin{proof}
Denote $j_U$, $j_V$, $j_{U \cap V}$ the corresponding open immersions.
Choose a distinguished triangle
$$
F \to Rj_{U, *}A \oplus Rj_{V, *}B \to Rj_{U \cap V, *}(B|_{U \cap V})
\to F[1]
$$
where the map $Rj_{V, *}B \to Rj_{U \cap V, *}(B|_{U \cap V})$ is the
obvious one and where
$Rj_{U, *}A \to Rj_{U \cap V, *}(B|_{U \cap V})$
is the composition of
$Rj_{U, *}A \to Rj_{U \cap V, *}(A|_{U \cap V})$
with $Rj_{U \cap V, *}c$. Restricting to $U$ we obtain
$$
F|_U \to A \oplus (Rj_{V, *}B)|_U \to (Rj_{U \cap V, *}(B|_{U \cap V}))|_U
\to F|_U[1]
$$
Denote $j : U \cap V \to U$. Compatibility of restriction to opens and
cohomology shows that both
$(Rj_{V, *}B)|_U$ and $(Rj_{U \cap V, *}(B|_{U \cap V}))|_U$
are canonically isomorphic to $Rj_*(B|_{U \cap V})$.
Hence the second arrow of the last displayed diagram has
a section, and we conclude that the morphism $F|_U \to A$ is
an isomorphism. Similarly, the morphism $F|_V \to B$ is an
isomorphism. The existence of the morphism $F \to E$ follows
from the Mayer-Vietoris sequence for $\Hom$, see
Lemma \ref{lemma-mayer-vietoris-hom}.
\end{proof}

\begin{lemma}
\label{lemma-vanishing-and-glueing}
Let $f : (X, \mathcal{O}_X) \to (Y, \mathcal{O}_Y)$ be a morphism
of ringed spaces. Let $\mathcal{B}$ be a basis for the topology on $Y$.
\begin{enumerate}
\item Assume $K$ is in $D(\mathcal{O}_X)$ such that
for $V \in \mathcal{B}$ we have $H^i(f^{-1}(V), K) = 0$ for $i < 0$.
Then $Rf_*K$ has vanishing cohomology sheaves in negative degrees,
$H^i(f^{-1}(V), K) = 0$ for $i < 0$ for all opens $V \subset Y$, and
the rule $V \mapsto H^0(f^{-1}V, K)$ is a sheaf on $Y$.
\item Assume $K, L$ are in $D(\mathcal{O}_X)$ such that
for $V \in \mathcal{B}$ we have
$\Ext^i(K|_{f^{-1}V}, L|_{f^{-1}V}) = 0$ for $i < 0$.
Then $\Ext^i(K|_{f^{-1}V}, L|_{f^{-1}V}) = 0$ for $i < 0$
for all opens $V \subset Y$ and
the rule $V \mapsto \Hom(K|_{f^{-1}V}, L|_{f^{-1}V})$ is a sheaf on $Y$.
\end{enumerate}
\end{lemma}

\begin{proof}
Lemma \ref{lemma-unbounded-describe-higher-direct-images} tells us
$H^i(Rf_*K)$ is the sheaf associated to the presheaf
$V \mapsto H^i(f^{-1}(V), K) = H^i(V, Rf_*K)$.
The assumptions in (1) imply that $Rf_*K$ has vanishing cohomology
sheaves in degrees $< 0$. We conclude that for any open $V \subset Y$
the cohomology group $H^i(V, Rf_*K)$ is zero for $i < 0$ and is equal to
$H^0(V, H^0(Rf_*K))$ for $i = 0$. This proves (1).

\medskip\noindent
To prove (2) apply (1) to the complex $R\SheafHom(K, L)$ using
Lemma \ref{lemma-section-RHom-over-U} to do the translation.
\end{proof}

\begin{situation}
\label{situation-locally-given}
Let $(X, \mathcal{O}_X)$ be a ringed space. We are given
\begin{enumerate}
\item a collection of opens $\mathcal{B}$ of $X$,
\item for $U \in \mathcal{B}$ an object $K_U$ in $D(\mathcal{O}_U)$,
\item for $V \subset U$ with $V, U \in \mathcal{B}$ an isomorphism
$\rho^U_V : K_U|_V \to K_V$ in $D(\mathcal{O}_V)$,
\end{enumerate}
such that whenever we have $W \subset V \subset U$ with $U, V, W$ in
$\mathcal{B}$, then $\rho^U_W = \rho^V_W \circ \rho ^U_V|_W$.
\end{situation}

\noindent
We won't be able to prove anything about this without making more
assumptions. An interesting case is where $\mathcal{B}$ is a basis
for the topology on $X$. Another is the case where we have a morphism
$f : X \to Y$ of topological spaces and the elements of $\mathcal{B}$
are the inverse images of the elements of a basis for the topology of $Y$.

\medskip\noindent
In Situation \ref{situation-locally-given} a {\it solution}
will be a pair $(K, \rho_U)$ where $K$ is an object of $D(\mathcal{O}_X)$
and $\rho_U : K|_U \to K_U$, $U \in \mathcal{B}$
are isomorphisms such that
we have $\rho^U_V \circ \rho_U|_V = \rho_V$ for all $V \subset U$,
$U, V \in \mathcal{B}$. In certain cases solutions are unique.

\begin{lemma}
\label{lemma-uniqueness}
In Situation \ref{situation-locally-given} assume
\begin{enumerate}
\item $X = \bigcup_{U \in \mathcal{B}} U$ and
for $U, V \in \mathcal{B}$ we have
$U \cap V = \bigcup_{W \in \mathcal{B}, W \subset U \cap V} W$,
\item for any $U \in \mathcal{B}$ we have $\Ext^i(K_U, K_U) = 0$
for $i < 0$.
\end{enumerate}
If a solution $(K, \rho_U)$ exists, then it is unique up to unique isomorphism
and moreover $\Ext^i(K, K) = 0$ for $i < 0$.
\end{lemma}

\begin{proof}
Let $(K, \rho_U)$ and $(K', \rho'_U)$ be a pair of solutions.
Let $f : X \to Y$ be the continuous map constructed
in Topology, Lemma \ref{topology-lemma-create-map-from-subcollection}.
Set $\mathcal{O}_Y = f_*\mathcal{O}_X$.
Then $K, K'$ and $\mathcal{B}$ are as in
Lemma \ref{lemma-vanishing-and-glueing} part (2).
Hence we obtain the vanishing of negative exts for $K$ and we see that
the rule
$$
V \longmapsto \Hom(K|_{f^{-1}V}, K'|_{f^{-1}V})
$$
is a sheaf on $Y$. As both $(K, \rho_U)$ and $(K', \rho'_U)$ are solutions
the maps
$$
(\rho'_U)^{-1} \circ \rho_U : K|_U \longrightarrow K'|_U
$$
over $U = f^{-1}(f(U))$ agree on overlaps. Hence we get a unique global
section of the sheaf above which defines the desired isomorphism
$K \to K'$ compatible with all structure available.
\end{proof}

\begin{remark}
\label{remark-uniqueness}
With notation and assumptions as in Lemma \ref{lemma-uniqueness}.
Suppose that $U, V \in \mathcal{B}$. Let $\mathcal{B}'$ be the set of
elements of $\mathcal{B}$ contained in $U \cap V$. Then
$$
(\{K_{U'}\}_{U' \in \mathcal{B}'},
\{\rho_{V'}^{U'}\}_{V' \subset U'\text{ with }U', V' \in \mathcal{B}'})
$$
is a system on the ringed space $U \cap V$
satisfying the assumptions of Lemma \ref{lemma-uniqueness}.
Moreover, both $(K_U|_{U \cap V}, \rho^U_{U'})$ and
$(K_V|_{U \cap V}, \rho^V_{U'})$ are solutions to this system.
By the lemma we find a unique isomorphism
$$
\rho_{U, V} : K_U|_{U \cap V} \longrightarrow K_V|_{U \cap V}
$$
such that for every $U' \subset U \cap V$, $U' \in \mathcal{B}$ the
diagram
$$
\xymatrix{
K_U|_{U'} \ar[rr]_{\rho_{U, V}|_{U'}} \ar[rd]_{\rho^U_{U'}} & &
K_V|_{U'} \ar[ld]^{\rho^V_{U'}} \\
& K_{U'}
}
$$
commutes. Pick a third element $W \in \mathcal{B}$. We obtain isomorphisms
$\rho_{U, W} : K_U|_{U \cap W} \to K_W|_{U \cap W}$ and
$\rho_{V, W} : K_U|_{V \cap W} \to K_W|_{V \cap W}$ satisfying
similar properties to those of $\rho_{U, V}$. Finally,
we have
$$
\rho_{U, W}|_{U \cap V \cap W} =
\rho_{V, W}|_{U \cap V \cap W} \circ \rho_{U, V}|_{U \cap V \cap W}
$$
This is true by the uniqueness in the lemma
because both sides of the equality are the unique isomorphism
compatible with the maps $\rho^U_{U''}$ and $\rho^W_{U''}$
for $U'' \subset U \cap V \cap W$, $U'' \in \mathcal{B}$.
Some minor details omitted.
The collection $(K_U, \rho_{U, V})$ is a descent datum
in the derived category for the open covering
$\mathcal{U} : X = \bigcup_{U \in \mathcal{B}} U$ of $X$.
In this language we are looking for ``effectiveness of the descent datum''
when we look for the existence of a solution.
\end{remark}

\begin{lemma}
\label{lemma-solution-in-finite-case}
In Situation \ref{situation-locally-given} assume
\begin{enumerate}
\item $X = U_1 \cup \ldots \cup U_n$ with $U_i \in \mathcal{B}$,
\item for $U, V \in \mathcal{B}$ we have
$U \cap V = \bigcup_{W \in \mathcal{B}, W \subset U \cap V} W$,
\item for any $U \in \mathcal{B}$ we have $\Ext^i(K_U, K_U) = 0$
for $i < 0$.
\end{enumerate}
Then a solution exists and is unique up to unique isomorphism.
\end{lemma}

\begin{proof}
Uniqueness was seen in Lemma \ref{lemma-uniqueness}. We may prove the lemma
by induction on $n$. The case $n = 1$ is immediate.

\medskip\noindent
The case $n = 2$.
Consider the isomorphism
$\rho_{U_1, U_2} : K_{U_1}|_{U_1 \cap U_2} \to K_{U_2}|_{U_1 \cap U_2}$
constructed in Remark \ref{remark-uniqueness}.
By Lemma \ref{lemma-glue} we obtain an object $K$ in $D(\mathcal{O}_X)$
and isomorphisms $\rho_{U_1} : K|_{U_1} \to K_{U_1}$ and
$\rho_{U_2} : K|_{U_2} \to K_{U_2}$ compatible with $\rho_{U_1, U_2}$.
Take $U \in \mathcal{B}$. We will construct an isomorphism
$\rho_U : K|_U \to K_U$ and we will leave it to the reader to verify
that $(K, \rho_U)$ is a solution. Consider the set $\mathcal{B}'$
of elements of $\mathcal{B}$ contained in either $U \cap U_1$ or contained in
$U \cap U_2$. Then $(K_U, \rho^U_{U'})$ is a solution for the system
$(\{K_{U'}\}_{U' \in \mathcal{B}'},
\{\rho_{V'}^{U'}\}_{V' \subset U'\text{ with }U', V' \in \mathcal{B}'})$
on the ringed space $U$.
We claim that $(K|_U, \tau_{U'})$ is another solution where
$\tau_{U'}$ for $U' \in \mathcal{B}'$ is chosen as follows:
if $U' \subset U_1$ then we take the composition
$$
K|_{U'} \xrightarrow{\rho_{U_1}|_{U'}}
K_{U_1}|_{U'} \xrightarrow{\rho^{U_1}_{U'}}
K_{U'}
$$
and if $U' \subset U_2$ then we take the composition
$$
K|_{U'} \xrightarrow{\rho_{U_2}|_{U'}}
K_{U_2}|_{U'} \xrightarrow{\rho^{U_2}_{U'}}
K_{U'}.
$$
To verify this is a solution use the property of the map $\rho_{U_1, U_2}$
described in Remark \ref{remark-uniqueness} and the compatibility of
$\rho_{U_1}$ and $\rho_{U_2}$ with $\rho_{U_1, U_2}$. Having said this
we apply Lemma \ref{lemma-uniqueness} to see that we obtain a unique
isomorphism $K|_{U'} \to K_{U'}$ compatible with the maps $\tau_{U'}$ and
$\rho^U_{U'}$ for $U' \in \mathcal{B}'$.

\medskip\noindent
The case $n > 2$. Consider the open subspace
$X' = U_1 \cup \ldots \cup U_{n - 1}$ and let $\mathcal{B}'$ be the set of
elements of $\mathcal{B}$ contained in $X'$. Then we find a system
$(\{K_U\}_{U \in \mathcal{B}'}, \{\rho_V^U\}_{U, V \in \mathcal{B}'})$
on the ringed space $X'$ to which we may apply our induction hypothesis.
We find a solution $(K_{X'}, \rho^{X'}_U)$.
Then we can consider the collection
$\mathcal{B}^* = \mathcal{B} \cup \{X'\}$ of opens of $X$ and we see that
we obtain a system
$(\{K_U\}_{U \in \mathcal{B}^*},
\{\rho_V^U\}_{V \subset U\text{ with }U, V \in \mathcal{B}^*})$.
Note that this new system also satisfies condition (3)
by Lemma \ref{lemma-uniqueness} applied to the solution $K_{X'}$.
For this system we have $X = X' \cup U_n$.
This reduces us to the case $n = 2$ we worked out above.
\end{proof}

\begin{lemma}
\label{lemma-glueing-increasing-union}
Let $X$ be a ringed space. Let $E$ be a well ordered set and let
$$
X = \bigcup\nolimits_{\alpha \in E} W_\alpha
$$
be an open covering with $W_\alpha \subset W_{\alpha + 1}$
and $W_\alpha = \bigcup_{\beta < \alpha} W_\beta$ if $\alpha$ is not
a successor. Let $K_\alpha$ be an object of $D(\mathcal{O}_{W_\alpha})$
with $\Ext^i(K_\alpha, K_\alpha) = 0$ for $i < 0$.
Assume given isomorphisms
$\rho_\beta^\alpha :  K_\alpha|_{W_\beta} \to K_\beta$ in
$D(\mathcal{O}_{W_\beta})$ for all $\beta < \alpha$ with
$\rho_\gamma^\alpha = \rho_\gamma^\beta \circ \rho^\alpha_\beta|_{W_\gamma}$
for $\gamma < \beta < \alpha$.
Then there exists an object
$K$ in $D(\mathcal{O}_X)$ and isomorphisms
$K|_{W_\alpha} \to K_\alpha$ for $\alpha \in E$
compatible with the isomorphisms $\rho_\beta^\alpha$.
\end{lemma}

\begin{proof}
In this proof $\alpha, \beta, \gamma, \ldots$ represent elements of $E$.
Choose a K-injective complex
$I_\alpha^\bullet$ on $W_\alpha$ representing $K_\alpha$.
For $\beta < \alpha$ denote $j_{\beta, \alpha} : W_\beta \to W_\alpha$
the inclusion morphism. Using transfinite recursion we will construct for all
$\beta < \alpha$ a map of complexes
$$
\tau_{\beta, \alpha} :
(j_{\beta, \alpha})_!I_\beta^\bullet
\longrightarrow
I_\alpha^\bullet
$$
representing the adjoint to the inverse of the isomorphism
$\rho^\alpha_\beta : K_\alpha|_{W_\beta} \to K_\beta$.
Moreover, we will do this in such that for
$\gamma < \beta < \alpha$ we have
$$
\tau_{\gamma, \alpha} = \tau_{\beta, \alpha} \circ
(j_{\beta, \alpha})_!\tau_{\gamma, \beta}
$$
as maps of complexes. Namely, suppose already given $\tau_{\gamma, \beta}$
composing correctly for all $\gamma < \beta < \alpha$.
If $\alpha = \alpha' + 1$ is a successor, then we choose any map of complexes
$$
(j_{\alpha', \alpha})_!I_{\alpha'}^\bullet \to I_\alpha^\bullet
$$
which is adjoint to the inverse of the isomorphism
$\rho^\alpha_{\alpha'} : K_\alpha|_{W_{\alpha'}} \to K_{\alpha'}$
(possible because $I_\alpha^\bullet$ is K-injective)
and for any $\beta < \alpha'$ we set
$$
\tau_{\beta, \alpha} = \tau_{\alpha', \alpha} \circ
(j_{\alpha', \alpha})_!\tau_{\beta, \alpha'}
$$
If $\alpha$ is not a successor, then
we can consider the complex on $W_\alpha$ given by
$$
C^\bullet = \colim_{\beta < \alpha} (j_{\beta, \alpha})_!I_\beta^\bullet
$$
(termwise colimit) where the transition maps of the sequence
are given by the maps $\tau_{\beta', \beta}$ for
$\beta' < \beta < \alpha$. We claim that $C^\bullet$
represents $K_\alpha$. Namely, for $\beta < \alpha$ the restriction
of the coprojection $(j_{\beta, \alpha})_!I_\beta^\bullet \to C^\bullet$
gives a map
$$
\sigma_\beta : I_\beta^\bullet \longrightarrow C^\bullet|_{W_\beta}
$$
which is a quasi-isomorphism: if $x \in W_\beta$ then looking
at stalks we get
$$
(C^\bullet)_x =
\colim_{\beta' < \alpha}
\left((j_{\beta', \alpha})_!I_{\beta'}^\bullet\right)_x =
\colim_{\beta \leq \beta' < \alpha} (I_{\beta'}^\bullet)_x
\longleftarrow
(I_\beta^\bullet)_x
$$
which is a quasi-isomorphism. Here we used that taking stalks
commutes with colimits, that filtered colimits are exact, and
that the maps $(I_\beta^\bullet)_x \to (I_{\beta'}^\bullet)_x$
are quasi-isomorphisms for $\beta \leq \beta' < \alpha$.
Hence $(C^\bullet, \sigma_\beta^{-1})$ is a solution to the
system $(\{K_\beta\}_{\beta < \alpha},
\{\rho^\beta_{\beta'}\}_{\beta' < \beta < \alpha})$.
Since $(K_\alpha, \rho^\alpha_\beta)$ is another solution
we obtain a unique isomorphism $\sigma : K_\alpha \to C^\bullet$
in $D(\mathcal{O}_{W_\alpha})$ compatible with all our maps, see
Lemma \ref{lemma-solution-in-finite-case}
(this is where we use the vanishing of negative ext groups).
Choose a morphism $\tau : C^\bullet \to I_\alpha^\bullet$
of complexes representing $\sigma$. Then we set
$$
\tau_{\beta, \alpha} = \tau|_{W_\beta} \circ \sigma_\beta
$$
to get the desired maps. Finally, we take $K$ to be the object of the derived
category represented by the complex
$$
K^\bullet = \colim_{\alpha \in E} (W_\alpha \to X)_!I_\alpha^\bullet
$$
where the transition maps are given by our carefully constructed
maps $\tau_{\beta, \alpha}$ for $\beta < \alpha$.
Arguing exactly as above we see that for all $\alpha$
the restriction of the coprojection determines an isomorphism
$$
K|_{W_\alpha} \longrightarrow K_\alpha
$$
compatible with the given maps $\rho^\alpha_\beta$.
\end{proof}

\noindent
Using transfinite induction we can prove the result in the general case.

\begin{theorem}[BBD gluing lemma]
\label{theorem-glueing-bbd-general}
\begin{reference}
Special case of \cite[Theorem 3.2.4]{BBD}
without boundedness assumption.
\end{reference}
In Situation \ref{situation-locally-given} assume
\begin{enumerate}
\item $X = \bigcup_{U \in \mathcal{B}} U$,
\item for $U, V \in \mathcal{B}$ we have
$U \cap V = \bigcup_{W \in \mathcal{B}, W \subset U \cap V} W$,
\item for any $U \in \mathcal{B}$ we have $\Ext^i(K_U, K_U) = 0$
for $i < 0$.
\end{enumerate}
Then there exists an object $K$ of $D(\mathcal{O}_X)$
and isomorphisms $\rho_U : K|_U \to K_U$ in $D(\mathcal{O}_U)$ for
$U \in \mathcal{B}$ such that $\rho^U_V \circ \rho_U|_V = \rho_V$
for all $V \subset U$ with $U, V \in \mathcal{B}$.
The pair $(K, \rho_U)$ is unique up to unique isomorphism.
\end{theorem}

\begin{proof}
A pair $(K, \rho_U)$ is called a solution in the text above.
The uniqueness follows from Lemma \ref{lemma-uniqueness}.
If $X$ has a finite covering by elements of $\mathcal{B}$
(for example if $X$ is quasi-compact), then the theorem
is a consequence of Lemma \ref{lemma-solution-in-finite-case}.
In the general case we argue in exactly the same manner,
using transfinite induction and
Lemma \ref{lemma-glueing-increasing-union}.

\medskip\noindent
First we use transfinite recursion to choose opens $W_\alpha \subset X$
for any ordinal $\alpha$. Namely, we set $W_0 = \emptyset$.
If $\alpha = \beta + 1$ is a successor, then either $W_\beta = X$
and we set $W_\alpha = X$ or $W_\beta \not = X$ and we set
$W_\alpha = W_\beta \cup U_\alpha$ where
$U_\alpha \in \mathcal{B}$ is not contained in $W_\beta$.
If $\alpha$ is a limit ordinal we set
$W_\alpha = \bigcup_{\beta < \alpha} W_\beta$.
Then for large enough $\alpha$ we have $W_\alpha = X$.
Observe that for every $\alpha$ the open $W_\alpha$ is
a union of elements of $\mathcal{B}$. Hence if
$\mathcal{B}_\alpha = \{U \in \mathcal{B}, U \subset W_\alpha\}$, then
$$
S_\alpha = (\{K_U\}_{U \in \mathcal{B}_\alpha},
\{\rho_V^U\}_{V \subset U\text{ with }U, V \in \mathcal{B}_\alpha})
$$
is a system as in Lemma \ref{lemma-uniqueness} on the ringed space $W_\alpha$.

\medskip\noindent
We will show by transfinite induction that for every $\alpha$
the system $S_\alpha$ has a solution. This will prove the theorem
as this system is the system given in the theorem for large $\alpha$.

\medskip\noindent
The case where $\alpha = \beta + 1$ is a successor ordinal.
(This case was already treated in the proof of the lemma above
but for clarity we repeat the argument.)
Recall that $W_\alpha = W_\beta \cup U_\alpha$ for some
$U_\alpha \in \mathcal{B}$ in this case.
By induction hypothesis we have a solution
$(K_{W_\beta}, \{\rho^{W_\beta}_U\}_{U \in \mathcal{B}_\beta})$
for the system $S_\beta$.
Then we can consider the collection
$\mathcal{B}_\alpha^* = \mathcal{B}_\alpha \cup \{W_\beta\}$
of opens of $W_\alpha$ and we see that we obtain a system
$(\{K_U\}_{U \in \mathcal{B}_\alpha^*},
\{\rho_V^U\}_{V \subset U\text{ with }U, V \in \mathcal{B}_\alpha^*})$.
Note that this new system also satisfies condition (3)
by Lemma \ref{lemma-uniqueness} applied to the solution $K_{W_\beta}$.
For this system we have $W_\alpha = W_\beta \cup U_\alpha$.
This reduces us to the case handled in
Lemma \ref{lemma-solution-in-finite-case}.

\medskip\noindent
The case where $\alpha$ is a limit ordinal. Recall that
$W_\alpha = \bigcup_{\beta < \alpha} W_\beta$ in this case.
For $\beta < \alpha$ let
$(K_{W_\beta}, \{\rho^{W_\beta}_U\}_{U \in \mathcal{B}_\beta})$
be the solution for $S_\beta$.
For $\gamma < \beta < \alpha$ the restriction
$K_{W_\beta}|_{W_\gamma}$ endowed with the maps
$\rho^{W_\beta}_U$, $U \in \mathcal{B}_\gamma$
is a solution for $S_\gamma$. By uniqueness we get unique isomorphisms
$\rho_{W_\gamma}^{W_\beta} : K_{W_\beta}|_{W_\gamma} \to K_{W_\gamma}$
compatible with the maps $\rho^{W_\beta}_U$ and $\rho^{W_\gamma}_U$
for $U \in \mathcal{B}_\gamma$. These maps compose in the correct manner,
i.e., $\rho_{W_\delta}^{W_\gamma} \circ \rho_{W_\gamma}^{W_\beta}|_{W_\delta}
= \rho^{W_\delta}_{W_\beta}$ for $\delta < \gamma < \beta < \alpha$.
Thus we may apply Lemma \ref{lemma-glueing-increasing-union}
(note that the vanishing of negative exts is true for
$K_{W_\beta}$ by Lemma \ref{lemma-uniqueness} applied
to the solution $K_{W_\beta}$)
to obtain $K_{W_\alpha}$ and isomorphisms
$$
\rho_{W_\beta}^{W_\alpha} :
K_{W_\alpha}|_{W_\beta}
\longrightarrow
K_{W_\beta}
$$
compatible with the maps $\rho_{W_\gamma}^{W_\beta}$ for
$\gamma < \beta < \alpha$.

\medskip\noindent
To show that $K_{W_\alpha}$ is a solution we still need to construct the
isomorphisms $\rho_U^{W_\alpha} : K_{W_\alpha}|_U \to K_U$ for
$U \in \mathcal{B}_\alpha$ satisfying certain compatibilities.
We choose $\rho_U^{W_\alpha}$ to be the unique map such that
for any $\beta < \alpha$ and any $V \in \mathcal{B}_\beta$
with $V \subset U$ the diagram
$$
\xymatrix{
K_{W_\alpha}|_V \ar[r]_{\rho_U^{W_\alpha}|_V}
\ar[d]_{\rho_{W_\beta}^{W_\alpha}|_V}
& K_U|_V \ar[d]^{\rho_U^V} \\
K_{W_\beta} \ar[r]^{\rho_V^{W_\beta}}
& K_V
}
$$
commutes. This makes sense because
$$
(\{K_V\}_{V \subset U, V \in \mathcal{B}_\beta\text{ for some }\beta < \alpha},
\{\rho_V^{V'}\}_{V \subset V'\text{ with }V, V' \subset U
\text{ and }V, V' \in \mathcal{B}_\beta\text{ for some }\beta < \alpha})
$$
is a system as in Lemma \ref{lemma-uniqueness} on the ringed space $U$
and because $(K_U, \rho^U_V)$ and
$(K_{W_\alpha}|_U,  \rho_V^{W_\beta}\circ \rho_{W_\beta}^{W_\alpha}|_V)$
are both solutions for this system. This gives existence and uniqueness.
We omit the proof that these
maps satisfy the desired compatibilities (it is just bookkeeping).
\end{proof}



\section{Strictly perfect complexes}
\label{section-strictly-perfect}

\noindent
Strictly perfect complexes of modules are used to define the notions
of pseudo-coherent and perfect complexes later on. They are defined
as follows.

\begin{definition}
\label{definition-strictly-perfect}
Let $(X, \mathcal{O}_X)$ be a ringed space.
Let $\mathcal{E}^\bullet$ be a complex of $\mathcal{O}_X$-modules.
We say $\mathcal{E}^\bullet$ is {\it strictly perfect}
if $\mathcal{E}^i$ is zero for all but finitely many $i$ and
$\mathcal{E}^i$ is a direct summand of a finite free
$\mathcal{O}_X$-module for all $i$.
\end{definition}

\noindent
Warning: Since we do not assume that $X$ is a locally ringed space,
it may not be true that a direct summand of a finite free
$\mathcal{O}_X$-module is finite locally free.

\begin{lemma}
\label{lemma-cone}
The cone on a morphism of strictly perfect complexes is
strictly perfect.
\end{lemma}

\begin{proof}
This is immediate from the definitions.
\end{proof}

\begin{lemma}
\label{lemma-tensor}
The total complex associated to the tensor product of two
strictly perfect complexes is strictly perfect.
\end{lemma}

\begin{proof}
Omitted.
\end{proof}

\begin{lemma}
\label{lemma-strictly-perfect-pullback}
Let $f : (X, \mathcal{O}_X) \to (Y, \mathcal{O}_Y)$
be a morphism of ringed spaces. If $\mathcal{F}^\bullet$ is a strictly
perfect complex of $\mathcal{O}_Y$-modules, then
$f^*\mathcal{F}^\bullet$ is a strictly perfect complex of
$\mathcal{O}_X$-modules.
\end{lemma}

\begin{proof}
The pullback of a finite free module is finite free. The functor
$f^*$ is additive functor hence preserves direct summands. The lemma follows.
\end{proof}

\begin{lemma}
\label{lemma-local-lift-map}
Let $(X, \mathcal{O}_X)$ be a ringed space.
Given a solid diagram of $\mathcal{O}_X$-modules
$$
\xymatrix{
\mathcal{E} \ar@{..>}[dr] \ar[r] & \mathcal{F} \\
& \mathcal{G} \ar[u]_p
}
$$
with $\mathcal{E}$ a direct summand of a finite free
$\mathcal{O}_X$-module and $p$ surjective, then a dotted arrow
making the diagram commute exists locally on $X$.
\end{lemma}

\begin{proof}
We may assume $\mathcal{E} = \mathcal{O}_X^{\oplus n}$ for some $n$.
In this case finding the dotted arrow is equivalent to lifting the
images of the basis elements in $\Gamma(X, \mathcal{F})$. This is
locally possible by the characterization of surjective maps of
sheaves (Sheaves, Section \ref{sheaves-section-exactness-points}).
\end{proof}

\begin{lemma}
\label{lemma-local-homotopy}
Let $(X, \mathcal{O}_X)$ be a ringed space.
\begin{enumerate}
\item Let $\alpha : \mathcal{E}^\bullet \to \mathcal{F}^\bullet$
be a morphism of complexes of $\mathcal{O}_X$-modules
with $\mathcal{E}^\bullet$ strictly perfect and $\mathcal{F}^\bullet$
acyclic. Then $\alpha$ is locally on $X$ homotopic to zero.
\item Let $\alpha : \mathcal{E}^\bullet \to \mathcal{F}^\bullet$
be a morphism of complexes of $\mathcal{O}_X$-modules
with $\mathcal{E}^\bullet$ strictly perfect, $\mathcal{E}^i = 0$
for $i < a$, and $H^i(\mathcal{F}^\bullet) = 0$ for $i \geq a$.
Then $\alpha$ is locally on $X$ homotopic to zero.
\end{enumerate}
\end{lemma}

\begin{proof}
The first statement follows from the second, hence we only prove (2).
We will prove this by induction on the length of the complex
$\mathcal{E}^\bullet$. If $\mathcal{E}^\bullet \cong \mathcal{E}[-n]$
for some direct summand $\mathcal{E}$ of a finite free
$\mathcal{O}_X$-module and integer $n \geq a$, then the result follows from
Lemma \ref{lemma-local-lift-map} and the fact that
$\mathcal{F}^{n - 1} \to \Ker(\mathcal{F}^n \to \mathcal{F}^{n + 1})$
is surjective by the assumed vanishing of $H^n(\mathcal{F}^\bullet)$.
If $\mathcal{E}^i$ is zero except for $i \in [a, b]$, then we have a
split exact sequence of complexes
$$
0 \to \mathcal{E}^b[-b] \to \mathcal{E}^\bullet \to
\sigma_{\leq b - 1}\mathcal{E}^\bullet \to 0
$$
which determines a distinguished triangle in
$K(\mathcal{O}_X)$. Hence an exact sequence
$$
\Hom_{K(\mathcal{O}_X)}(
\sigma_{\leq b - 1}\mathcal{E}^\bullet, \mathcal{F}^\bullet)
\to
\Hom_{K(\mathcal{O}_X)}(\mathcal{E}^\bullet, \mathcal{F}^\bullet)
\to
\Hom_{K(\mathcal{O}_X)}(\mathcal{E}^b[-b], \mathcal{F}^\bullet)
$$
by the axioms of triangulated categories. The composition
$\mathcal{E}^b[-b] \to \mathcal{F}^\bullet$ is locally homotopic to
zero, whence we may assume our map comes from an element in the
left hand side of the displayed exact sequence above. This element
is locally zero by induction hypothesis.
\end{proof}

\begin{lemma}
\label{lemma-lift-through-quasi-isomorphism}
Let $(X, \mathcal{O}_X)$ be a ringed space.
Given a solid diagram of complexes of $\mathcal{O}_X$-modules
$$
\xymatrix{
\mathcal{E}^\bullet \ar@{..>}[dr] \ar[r]_\alpha & \mathcal{F}^\bullet \\
& \mathcal{G}^\bullet \ar[u]_f
}
$$
with $\mathcal{E}^\bullet$ strictly perfect, $\mathcal{E}^j = 0$ for
$j < a$ and $H^j(f)$ an isomorphism for $j > a$ and surjective for $j = a$,
then a dotted arrow making the diagram commute up to homotopy
exists locally on $X$.
\end{lemma}

\begin{proof}
Our assumptions on $f$ imply the cone $C(f)^\bullet$ has vanishing
cohomology sheaves in degrees $\geq a$.
Hence Lemma \ref{lemma-local-homotopy} guarantees there is an open
covering $X = \bigcup U_i$ such that the composition
$\mathcal{E}^\bullet \to \mathcal{F}^\bullet \to C(f)^\bullet$
is homotopic to zero over $U_i$. Since
$$
\mathcal{G}^\bullet \to \mathcal{F}^\bullet \to C(f)^\bullet \to
\mathcal{G}^\bullet[1]
$$
restricts to a distinguished triangle in $K(\mathcal{O}_{U_i})$
we see that we can lift $\alpha|_{U_i}$ up to homotopy to a map
$\alpha_i : \mathcal{E}^\bullet|_{U_i} \to \mathcal{G}^\bullet|_{U_i}$
as desired.
\end{proof}

\begin{lemma}
\label{lemma-local-actual}
Let $(X, \mathcal{O}_X)$ be a ringed space.
Let $\mathcal{E}^\bullet$, $\mathcal{F}^\bullet$ be complexes
of $\mathcal{O}_X$-modules with $\mathcal{E}^\bullet$ strictly perfect.
\begin{enumerate}
\item For any element
$\alpha \in \Hom_{D(\mathcal{O}_X)}(\mathcal{E}^\bullet, \mathcal{F}^\bullet)$
there exists an open covering $X = \bigcup U_i$ such that
$\alpha|_{U_i}$ is given by a morphism of complexes
$\alpha_i : \mathcal{E}^\bullet|_{U_i} \to \mathcal{F}^\bullet|_{U_i}$.
\item Given a morphism of complexes
$\alpha : \mathcal{E}^\bullet \to \mathcal{F}^\bullet$
whose image in the group
$\Hom_{D(\mathcal{O}_X)}(\mathcal{E}^\bullet, \mathcal{F}^\bullet)$
is zero, there exists an open covering $X = \bigcup U_i$ such that
$\alpha|_{U_i}$ is homotopic to zero.
\end{enumerate}
\end{lemma}

\begin{proof}
Proof of (1).
By the construction of the derived category we can find a quasi-isomorphism
$f : \mathcal{F}^\bullet \to \mathcal{G}^\bullet$ and a map of complexes
$\beta : \mathcal{E}^\bullet \to \mathcal{G}^\bullet$ such that
$\alpha = f^{-1}\beta$. Thus the result follows from
Lemma \ref{lemma-lift-through-quasi-isomorphism}.
We omit the proof of (2).
\end{proof}

\begin{lemma}
\label{lemma-Rhom-strictly-perfect}
Let $(X, \mathcal{O}_X)$ be a ringed space.
Let $\mathcal{E}^\bullet$, $\mathcal{F}^\bullet$ be complexes
of $\mathcal{O}_X$-modules with $\mathcal{E}^\bullet$ strictly perfect.
Then the internal hom $R\SheafHom(\mathcal{E}^\bullet, \mathcal{F}^\bullet)$
is represented by the complex $\mathcal{H}^\bullet$ with terms
$$
\mathcal{H}^n =
\bigoplus\nolimits_{n = p + q}
\SheafHom_{\mathcal{O}_X}(\mathcal{E}^{-q}, \mathcal{F}^p)
$$
and differential as described in Section \ref{section-hom-complexes}.
\end{lemma}

\begin{proof}
Choose a quasi-isomorphism $\mathcal{F}^\bullet \to \mathcal{I}^\bullet$
into a K-injective complex. Let $(\mathcal{H}')^\bullet$ be the
complex with terms
$$
(\mathcal{H}')^n =
\prod\nolimits_{n = p + q}
\SheafHom_{\mathcal{O}_X}(\mathcal{E}^{-q}, \mathcal{I}^p)
$$
which represents $R\SheafHom(\mathcal{E}^\bullet, \mathcal{F}^\bullet)$
by the construction in Section \ref{section-internal-hom}. It suffices
to show that the map
$$
\mathcal{H}^\bullet \longrightarrow (\mathcal{H}')^\bullet
$$
is a quasi-isomorphism. Given an open $U \subset X$ we have
by inspection
$$
H^0(\mathcal{H}^\bullet(U)) =
\Hom_{K(\mathcal{O}_U)}(\mathcal{E}^\bullet|_U, \mathcal{I}^\bullet|_U)
\to
H^0((\mathcal{H}')^\bullet(U)) =
\Hom_{D(\mathcal{O}_U)}(\mathcal{E}^\bullet|_U, \mathcal{I}^\bullet|_U)
$$
By Lemma \ref{lemma-local-actual} the sheafification of
$U \mapsto H^0(\mathcal{H}^\bullet(U))$
is equal to the sheafification of
$U \mapsto H^0((\mathcal{H}')^\bullet(U))$. A similar argument can be
given for the other cohomology sheaves. Thus $\mathcal{H}^\bullet$
is quasi-isomorphic to $(\mathcal{H}')^\bullet$ which proves the lemma.
\end{proof}

\begin{lemma}
\label{lemma-Rhom-strictly-perfect-K-flat}
In the situation of Lemma \ref{lemma-Rhom-strictly-perfect}
if $\mathcal{F}^\bullet$ is K-flat, then $\mathcal{H}^\bullet$ is K-flat.
\end{lemma}

\begin{proof}
Observe that $\mathcal{H}^\bullet$ is simply the hom complex
$\SheafHom^\bullet(\mathcal{E}^\bullet, \mathcal{F}^\bullet)$
since the boundedness of the strictly prefect complex
$\mathcal{E}^\bullet$ insures that the products in the definition
of the hom complex turn into direct sums.
Let $\mathcal{K}^\bullet$ be an acyclic complex of
$\mathcal{O}_X$-modules. Consider the map
$$
\gamma :
\text{Tot}(\mathcal{K}^\bullet \otimes
\SheafHom^\bullet(\mathcal{E}^\bullet, \mathcal{F}^\bullet))
\longrightarrow
\SheafHom^\bullet(\mathcal{E}^\bullet,
\text{Tot}(\mathcal{K}^\bullet \otimes \mathcal{F}^\bullet))
$$
of Lemma \ref{lemma-diagonal-better}. Since $\mathcal{F}^\bullet$ is K-flat,
the complex $\text{Tot}(\mathcal{K}^\bullet \otimes \mathcal{F}^\bullet)$
is acyclic, and hence by Lemma \ref{lemma-local-actual}
(or Lemma \ref{lemma-Rhom-strictly-perfect} if you like)
the target of $\gamma$ is acyclic too. Hence to prove the lemma it suffices to
show that $\gamma$ is an isomorphism of complexes. To see this, we may argue
by induction on the length of the complex $\mathcal{E}^\bullet$.
If the length is $\leq 1$ then the $\mathcal{E}^\bullet$
is a direct summand of $\mathcal{O}_X^{\oplus n}[k]$ for some
$n \geq 0$ and $k \in \mathbf{Z}$ and in this case the result
follows by inspection. If the length is $> 1$, then we reduce
to smaller length by considering the termwise split short exact
sequence of complexes
$$
0 \to \sigma_{\geq a + 1} \mathcal{E}^\bullet \to
\mathcal{E}^\bullet \to
\sigma_{\leq a} \mathcal{E}^\bullet \to 0
$$
for a suitable $a \in \mathbf{Z}$, see
Homology, Section \ref{homology-section-truncations}.
Then $\gamma$ fits into a morphism of termwise split
short exact sequences of complexes.
By induction $\gamma$ is an isomorphism for
$\sigma_{\geq a + 1} \mathcal{E}^\bullet$
and $\sigma_{\leq a} \mathcal{E}^\bullet$ and hence the result for
$\mathcal{E}^\bullet$ follows. Some details omitted.
\end{proof}

\begin{lemma}
\label{lemma-Rhom-complex-of-direct-summands-finite-free}
Let $(X, \mathcal{O}_X)$ be a ringed space.
Let $\mathcal{E}^\bullet$, $\mathcal{F}^\bullet$ be complexes
of $\mathcal{O}_X$-modules with
\begin{enumerate}
\item $\mathcal{F}^n = 0$ for $n \ll 0$,
\item $\mathcal{E}^n = 0$ for $n \gg 0$, and
\item $\mathcal{E}^n$ isomorphic to a direct summand of a finite
free $\mathcal{O}_X$-module.
\end{enumerate}
Then the internal hom $R\SheafHom(\mathcal{E}^\bullet, \mathcal{F}^\bullet)$
is represented by the complex $\mathcal{H}^\bullet$ with terms
$$
\mathcal{H}^n =
\bigoplus\nolimits_{n = p + q}
\SheafHom_{\mathcal{O}_X}(\mathcal{E}^{-q}, \mathcal{F}^p)
$$
and differential as described in Section \ref{section-internal-hom}.
\end{lemma}

\begin{proof}
Choose a quasi-isomorphism $\mathcal{F}^\bullet \to \mathcal{I}^\bullet$
where $\mathcal{I}^\bullet$ is a bounded below complex of injectives.
Note that $\mathcal{I}^\bullet$ is K-injective
(Derived Categories, Lemma
\ref{derived-lemma-bounded-below-injectives-K-injective}).
Hence the construction in Section \ref{section-internal-hom}
shows that
$R\SheafHom(\mathcal{E}^\bullet, \mathcal{F}^\bullet)$ is
represented by the complex $(\mathcal{H}')^\bullet$ with terms
$$
(\mathcal{H}')^n =
\prod\nolimits_{n = p + q}
\SheafHom_{\mathcal{O}_X}(\mathcal{E}^{-q}, \mathcal{I}^p) =
\bigoplus\nolimits_{n = p + q}
\SheafHom_{\mathcal{O}_X}(\mathcal{E}^{-q}, \mathcal{I}^p)
$$
(equality because there are only finitely many nonzero terms).
Note that $\mathcal{H}^\bullet$ is the total complex associated to
the double complex with terms
$\SheafHom_{\mathcal{O}_X}(\mathcal{E}^{-q}, \mathcal{F}^p)$
and similarly for $(\mathcal{H}')^\bullet$.
The natural map $(\mathcal{H}')^\bullet \to \mathcal{H}^\bullet$
comes from a map of double complexes.
Thus to show this map is a quasi-isomorphism, we may use the spectral
sequence of a double complex
(Homology, Lemma \ref{homology-lemma-first-quadrant-ss})
$$
{}'E_1^{p, q} =
H^p(\SheafHom_{\mathcal{O}_X}(\mathcal{E}^{-q}, \mathcal{F}^\bullet))
$$
converging to $H^{p + q}(\mathcal{H}^\bullet)$ and similarly for
$(\mathcal{H}')^\bullet$. To finish the proof of the lemma it
suffices to show that $\mathcal{F}^\bullet \to \mathcal{I}^\bullet$
induces an isomorphism
$$
H^p(\SheafHom_{\mathcal{O}_X}(\mathcal{E}, \mathcal{F}^\bullet))
\longrightarrow
H^p(\SheafHom_{\mathcal{O}_X}(\mathcal{E}, \mathcal{I}^\bullet))
$$
on cohomology sheaves whenever $\mathcal{E}$ is a direct summand of a
finite free $\mathcal{O}_X$-module. Since this is clear when $\mathcal{E}$
is finite free the result follows.
\end{proof}





\section{Pseudo-coherent modules}
\label{section-pseudo-coherent}

\noindent
In this section we discuss pseudo-coherent complexes.

\begin{definition}
\label{definition-pseudo-coherent}
Let $(X, \mathcal{O}_X)$ be a ringed space. Let $\mathcal{E}^\bullet$
be a complex of $\mathcal{O}_X$-modules. Let $m \in \mathbf{Z}$.
\begin{enumerate}
\item We say $\mathcal{E}^\bullet$ is {\it $m$-pseudo-coherent}
if there exists an open covering $X = \bigcup U_i$ and for each $i$
a morphism of complexes
$\alpha_i : \mathcal{E}_i^\bullet \to \mathcal{E}^\bullet|_{U_i}$
where $\mathcal{E}_i^\bullet$ is strictly perfect on $U_i$ and
$H^j(\alpha_i)$ is an isomorphism for $j > m$ and $H^m(\alpha_i)$
is surjective.
\item We say $\mathcal{E}^\bullet$ is {\it pseudo-coherent}
if it is $m$-pseudo-coherent for all $m$.
\item We say an object $E$ of $D(\mathcal{O}_X)$ is
{\it $m$-pseudo-coherent} (resp.\ {\it pseudo-coherent})
if and only if it can be represented by a $m$-pseudo-coherent
(resp.\ pseudo-coherent) complex of $\mathcal{O}_X$-modules.
\end{enumerate}
\end{definition}

\noindent
If $X$ is quasi-compact, then an $m$-pseudo-coherent object
of $D(\mathcal{O}_X)$ is in $D^-(\mathcal{O}_X)$. But this need
not be the case if $X$ is not quasi-compact.

\begin{lemma}
\label{lemma-pseudo-coherent-independent-representative}
Let $(X, \mathcal{O}_X)$ be a ringed space. Let $E$ be an object
of $D(\mathcal{O}_X)$.
\begin{enumerate}
\item If there exists an open covering $X = \bigcup U_i$,
strictly perfect complexes $\mathcal{E}_i^\bullet$ on $U_i$, and
maps $\alpha_i : \mathcal{E}_i^\bullet \to E|_{U_i}$ in
$D(\mathcal{O}_{U_i})$ with $H^j(\alpha_i)$ an isomorphism for $j > m$
and $H^m(\alpha_i)$ surjective, then $E$ is $m$-pseudo-coherent.
\item If $E$ is $m$-pseudo-coherent, then any complex representing
$E$ is $m$-pseudo-coherent.
\end{enumerate}
\end{lemma}

\begin{proof}
Let $\mathcal{F}^\bullet$ be any complex representing $E$
and let $X = \bigcup U_i$ and
$\alpha_i : \mathcal{E}_i^\bullet \to E|_{U_i}$ be as in (1).
We will show that $\mathcal{F}^\bullet$ is $m$-pseudo-coherent
as a complex, which will prove (1) and (2) simultaneously.
By Lemma \ref{lemma-local-actual}
we can after refining the open covering $X = \bigcup U_i$
represent the maps $\alpha_i$ by maps of complexes
$\alpha_i : \mathcal{E}_i^\bullet \to \mathcal{F}^\bullet|_{U_i}$.
By assumption
$H^j(\alpha_i)$ are isomorphisms for $j > m$, and $H^m(\alpha_i)$
is surjective whence $\mathcal{F}^\bullet$ is
$m$-pseudo-coherent.
\end{proof}

\begin{lemma}
\label{lemma-pseudo-coherent-pullback}
Let $f : (X, \mathcal{O}_X) \to (Y, \mathcal{O}_Y)$
be a morphism of ringed spaces. Let $E$ be an object of
$D(\mathcal{O}_Y)$. If $E$ is $m$-pseudo-coherent,
then $Lf^*E$ is $m$-pseudo-coherent.
\end{lemma}

\begin{proof}
Represent $E$ by a complex $\mathcal{E}^\bullet$ of $\mathcal{O}_Y$-modules
and choose an open covering $Y = \bigcup V_i$ and
$\alpha_i : \mathcal{E}_i^\bullet \to \mathcal{E}^\bullet|_{V_i}$
as in Definition \ref{definition-pseudo-coherent}.
Set $U_i = f^{-1}(V_i)$.
By Lemma \ref{lemma-pseudo-coherent-independent-representative}
it suffices to show that $Lf^*\mathcal{E}^\bullet|_{U_i}$ is
$m$-pseudo-coherent.
Choose a distinguished triangle
$$
\mathcal{E}_i^\bullet \to
\mathcal{E}^\bullet|_{V_i} \to
C \to
\mathcal{E}_i^\bullet[1]
$$
The assumption on $\alpha_i$ means exactly that the cohomology sheaves
$H^j(C)$ are zero for all $j \geq m$. Denote $f_i : U_i \to V_i$
the restriction of $f$. Note that $Lf^*\mathcal{E}^\bullet|_{U_i} =
Lf_i^*(\mathcal{E}|_{V_i})$. Applying $Lf_i^*$ we obtain
the distinguished triangle
$$
Lf_i^*\mathcal{E}_i^\bullet \to
Lf_i^*\mathcal{E}|_{V_i} \to
Lf_i^*C \to
Lf_i^*\mathcal{E}_i^\bullet[1]
$$
By the construction of $Lf_i^*$ as a left derived functor we see that
$H^j(Lf_i^*C) = 0$ for $j \geq m$ (by the dual of Derived Categories, Lemma
\ref{derived-lemma-negative-vanishing}). Hence $H^j(Lf_i^*\alpha_i)$ is an
isomorphism for $j > m$ and $H^m(Lf^*\alpha_i)$ is surjective.
On the other hand,
$Lf_i^*\mathcal{E}_i^\bullet = f_i^*\mathcal{E}_i^\bullet$.
is strictly perfect by Lemma \ref{lemma-strictly-perfect-pullback}.
Thus we conclude.
\end{proof}

\begin{lemma}
\label{lemma-cone-pseudo-coherent}
Let $(X, \mathcal{O}_X)$ be a ringed space and $m \in \mathbf{Z}$.
Let $(K, L, M, f, g, h)$ be a distinguished triangle in $D(\mathcal{O}_X)$.
\begin{enumerate}
\item If $K$ is $(m + 1)$-pseudo-coherent and $L$ is $m$-pseudo-coherent
then $M$ is $m$-pseudo-coherent.
\item If $K$ and $M$ are $m$-pseudo-coherent, then $L$ is $m$-pseudo-coherent.
\item If $L$ is $(m + 1)$-pseudo-coherent and $M$
is $m$-pseudo-coherent, then $K$ is $(m + 1)$-pseudo-coherent.
\end{enumerate}
\end{lemma}

\begin{proof}
Proof of (1). Choose an open covering $X = \bigcup U_i$ and
maps $\alpha_i : \mathcal{K}_i^\bullet \to K|_{U_i}$ in $D(\mathcal{O}_{U_i})$
with $\mathcal{K}_i^\bullet$ strictly perfect and $H^j(\alpha_i)$
isomorphisms for $j > m + 1$ and surjective for $j = m + 1$.
We may replace $\mathcal{K}_i^\bullet$ by
$\sigma_{\geq m + 1}\mathcal{K}_i^\bullet$
and hence we may assume that $\mathcal{K}_i^j = 0$
for $j < m + 1$. After refining the open covering we may choose
maps $\beta_i : \mathcal{L}_i^\bullet \to L|_{U_i}$ in $D(\mathcal{O}_{U_i})$
with $\mathcal{L}_i^\bullet$ strictly perfect such that
$H^j(\beta)$ is an isomorphism for $j > m$ and
surjective for $j = m$. By
Lemma \ref{lemma-lift-through-quasi-isomorphism}
we can, after refining the covering, find maps of complexes
$\gamma_i : \mathcal{K}^\bullet \to \mathcal{L}^\bullet$
such that the diagrams
$$
\xymatrix{
K|_{U_i} \ar[r] & L|_{U_i} \\
\mathcal{K}_i^\bullet \ar[u]^{\alpha_i} \ar[r]^{\gamma_i} &
\mathcal{L}_i^\bullet \ar[u]_{\beta_i}
}
$$
are commutative in $D(\mathcal{O}_{U_i})$ (this requires representing the
maps $\alpha_i$, $\beta_i$ and $K|_{U_i} \to L|_{U_i}$
by actual maps of complexes; some details omitted).
The cone $C(\gamma_i)^\bullet$ is strictly perfect (Lemma \ref{lemma-cone}).
The commutativity of the diagram implies that there exists a morphism
of distinguished triangles
$$
(\mathcal{K}_i^\bullet, \mathcal{L}_i^\bullet, C(\gamma_i)^\bullet)
\longrightarrow
(K|_{U_i}, L|_{U_i}, M|_{U_i}).
$$
It follows from the induced map on long exact cohomology sequences and
Homology, Lemmas \ref{homology-lemma-four-lemma} and
\ref{homology-lemma-five-lemma}
that $C(\gamma_i)^\bullet \to M|_{U_i}$ induces an isomorphism
on cohomology in degrees $> m$ and a surjection in degree $m$.
Hence $M$ is $m$-pseudo-coherent by
Lemma \ref{lemma-pseudo-coherent-independent-representative}.

\medskip\noindent
Assertions (2) and (3) follow from (1) by rotating the distinguished
triangle.
\end{proof}

\begin{lemma}
\label{lemma-tensor-pseudo-coherent}
Let $(X, \mathcal{O}_X)$ be a ringed space. Let $K, L$ be objects
of $D(\mathcal{O}_X)$.
\begin{enumerate}
\item If $K$ is $n$-pseudo-coherent and $H^i(K) = 0$ for $i > a$
and $L$ is $m$-pseudo-coherent and $H^j(L) = 0$ for $j > b$, then
$K \otimes_{\mathcal{O}_X}^\mathbf{L} L$ is $t$-pseudo-coherent
with $t = \max(m + a, n + b)$.
\item If $K$ and $L$ are pseudo-coherent, then
$K \otimes_{\mathcal{O}_X}^\mathbf{L} L$ is pseudo-coherent.
\end{enumerate}
\end{lemma}

\begin{proof}
Proof of (1). By replacing $X$ by the members of an open covering
we may assume there exist strictly perfect complexes $\mathcal{K}^\bullet$
and $\mathcal{L}^\bullet$ and maps
$\alpha : \mathcal{K}^\bullet \to K$ and
$\beta : \mathcal{L}^\bullet \to L$ with $H^i(\alpha)$ and isomorphism
for $i > n$ and surjective for $i = n$ and with
$H^i(\beta)$ and isomorphism for $i > m$ and surjective for $i = m$.
Then the map
$$
\alpha \otimes^\mathbf{L} \beta :
\text{Tot}(\mathcal{K}^\bullet \otimes_{\mathcal{O}_X} \mathcal{L}^\bullet)
\to K \otimes_{\mathcal{O}_X}^\mathbf{L} L
$$
induces isomorphisms on cohomology sheaves in degree $i$ for
$i > t$ and a surjection for $i = t$. This follows from the
spectral sequence of tors (details omitted).

\medskip\noindent
Proof of (2). We may first replace $X$ by the members of an open
covering to reduce to the case that $K$ and $L$ are bounded above.
Then the statement follows immediately from case (1).
\end{proof}

\begin{lemma}
\label{lemma-summands-pseudo-coherent}
Let $(X, \mathcal{O}_X)$ be a ringed space. Let $m \in \mathbf{Z}$.
If $K \oplus L$ is $m$-pseudo-coherent (resp.\ pseudo-coherent)
in $D(\mathcal{O}_X)$ so are $K$ and $L$.
\end{lemma}

\begin{proof}
Assume that $K \oplus L$ is $m$-pseudo-coherent.
After replacing $X$ by the members of an open covering we may
assume $K \oplus L \in D^-(\mathcal{O}_X)$, hence
$L \in D^-(\mathcal{O}_X)$.
Note that there is a distinguished triangle
$$
(K \oplus L, K \oplus L, L \oplus L[1]) =
(K, K, 0) \oplus (L, L, L \oplus L[1])
$$
see
Derived Categories, Lemma \ref{derived-lemma-direct-sum-triangles}.
By
Lemma \ref{lemma-cone-pseudo-coherent}
we see that $L \oplus L[1]$ is $m$-pseudo-coherent.
Hence also $L[1] \oplus L[2]$ is $m$-pseudo-coherent.
By induction $L[n] \oplus L[n + 1]$ is $m$-pseudo-coherent.
Since $L$ is bounded above we see that $L[n]$ is $m$-pseudo-coherent
for large $n$. Hence working backwards, using the distinguished triangles
$$
(L[n], L[n] \oplus L[n - 1], L[n - 1])
$$
we conclude that $L[n - 1], L[n - 2], \ldots, L$ are $m$-pseudo-coherent
as desired.
\end{proof}

\begin{lemma}
\label{lemma-complex-pseudo-coherent-modules}
Let $(X, \mathcal{O}_X)$ be a ringed space.
Let $m \in \mathbf{Z}$. Let $\mathcal{F}^\bullet$ be a (locally) bounded
above complex of $\mathcal{O}_X$-modules such that
$\mathcal{F}^i$ is $(m - i)$-pseudo-coherent for all $i$.
Then $\mathcal{F}^\bullet$ is $m$-pseudo-coherent.
\end{lemma}

\begin{proof}
Omitted. Hint: use Lemma \ref{lemma-cone-pseudo-coherent} and truncations
as in the proof of
More on Algebra, Lemma \ref{more-algebra-lemma-complex-pseudo-coherent-modules}.
\end{proof}

\begin{lemma}
\label{lemma-cohomology-pseudo-coherent}
Let $(X, \mathcal{O}_X)$ be a ringed space. Let $m \in \mathbf{Z}$. Let
$E$ be an object of $D(\mathcal{O}_X)$. If $E$ is (locally) bounded above
and $H^i(E)$ is $(m - i)$-pseudo-coherent for all $i$, then
$E$ is $m$-pseudo-coherent.
\end{lemma}

\begin{proof}
Omitted. Hint: use Lemma \ref{lemma-cone-pseudo-coherent} and truncations
as in the proof of
More on Algebra, Lemma \ref{more-algebra-lemma-cohomology-pseudo-coherent}.
\end{proof}

\begin{lemma}
\label{lemma-finite-cohomology}
Let $(X, \mathcal{O}_X)$ be a ringed space.
Let $K$ be an object of $D(\mathcal{O}_X)$.
Let $m \in \mathbf{Z}$.
\begin{enumerate}
\item If $K$ is $m$-pseudo-coherent and $H^i(K) = 0$
for $i > m$, then $H^m(K)$ is a finite type $\mathcal{O}_X$-module.
\item If $K$ is $m$-pseudo-coherent and $H^i(K) = 0$
for $i > m + 1$, then $H^{m + 1}(K)$ is a finitely presented
$\mathcal{O}_X$-module.
\end{enumerate}
\end{lemma}

\begin{proof}
Proof of (1). We may work locally on $X$. Hence we may assume there exists
a strictly perfect complex $\mathcal{E}^\bullet$ and a map
$\alpha : \mathcal{E}^\bullet \to K$ which induces
an isomorphism on cohomology in degrees $> m$ and a surjection in degree $m$.
It suffices to prove the result for $\mathcal{E}^\bullet$.
Let $n$ be the largest integer such that $\mathcal{E}^n \not = 0$.
If $n = m$, then $H^m(\mathcal{E}^\bullet)$ is a quotient of
$\mathcal{E}^n$ and the result is clear.
If $n > m$, then $\mathcal{E}^{n - 1} \to \mathcal{E}^n$ is surjective as
$H^n(E^\bullet) = 0$. By Lemma \ref{lemma-local-lift-map}
we can locally find a section of this surjection and write
$\mathcal{E}^{n - 1} = \mathcal{E}' \oplus \mathcal{E}^n$.
Hence it suffices to prove the result for the complex
$(\mathcal{E}')^\bullet$ which is the same as $\mathcal{E}^\bullet$
except has $\mathcal{E}'$ in degree $n - 1$ and $0$ in degree $n$.
We win by induction on $n$.

\medskip\noindent
Proof of (2). We may work locally on $X$. Hence we may assume there exists
a strictly perfect complex $\mathcal{E}^\bullet$ and a map
$\alpha : \mathcal{E}^\bullet \to K$ which induces
an isomorphism on cohomology in degrees $> m$ and a surjection in degree $m$.
As in the proof of (1) we can reduce to the case that $\mathcal{E}^i = 0$
for $i > m + 1$. Then we see that
$H^{m + 1}(K) \cong H^{m + 1}(\mathcal{E}^\bullet) =
\Coker(\mathcal{E}^m \to \mathcal{E}^{m + 1})$
which is of finite presentation.
\end{proof}

\begin{lemma}
\label{lemma-n-pseudo-module}
Let $(X, \mathcal{O}_X)$ be a ringed space. Let $\mathcal{F}$ be a sheaf
of $\mathcal{O}_X$-modules.
\begin{enumerate}
\item $\mathcal{F}$ viewed as an object of $D(\mathcal{O}_X)$ is
$0$-pseudo-coherent if and only if $\mathcal{F}$ is a finite type
$\mathcal{O}_X$-module, and
\item $\mathcal{F}$ viewed as an object of $D(\mathcal{O}_X)$ is
$(-1)$-pseudo-coherent if and only if $\mathcal{F}$ is an
$\mathcal{O}_X$-module of finite presentation.
\end{enumerate}
\end{lemma}

\begin{proof}
Use Lemma \ref{lemma-finite-cohomology}
to prove the implications in one direction and
Lemma \ref{lemma-cohomology-pseudo-coherent} for the other.
\end{proof}





\section{Tor dimension}
\label{section-tor}

\noindent
In this section we take a closer look at resolutions by flat modules.

\begin{definition}
\label{definition-tor-amplitude}
Let $(X, \mathcal{O}_X)$ be a ringed space.
Let $E$ be an object of $D(\mathcal{O}_X)$.
Let $a, b \in \mathbf{Z}$ with $a \leq b$.
\begin{enumerate}
\item We say $E$ has {\it tor-amplitude in $[a, b]$}
if $H^i(E \otimes_{\mathcal{O}_X}^\mathbf{L} \mathcal{F}) = 0$
for all $\mathcal{O}_X$-modules $\mathcal{F}$ and all $i \not \in [a, b]$.
\item We say $E$ has {\it finite tor dimension}
if it has tor-amplitude in $[a, b]$ for some $a, b$.
\item We say $E$ {\it locally has finite tor dimension}
if there exists an open covering $X = \bigcup U_i$ such that
$E|_{U_i}$ has finite tor dimension for all $i$.
\end{enumerate}
An $\mathcal{O}_X$-module $\mathcal{F}$ has {\it tor dimension $\leq d$}
if $\mathcal{F}[0]$ viewed as an object of $D(\mathcal{O}_X)$ has
tor-amplitude in $[-d, 0]$.
\end{definition}

\noindent
Note that if $E$ as in the definition
has finite tor dimension, then $E$ is an object of
$D^b(\mathcal{O}_X)$ as can be seen by taking $\mathcal{F} = \mathcal{O}_X$
in the definition above.

\begin{lemma}
\label{lemma-last-one-flat}
Let $(X, \mathcal{O}_X)$ be a ringed space.
Let $\mathcal{E}^\bullet$ be a bounded above complex of flat
$\mathcal{O}_X$-modules with tor-amplitude in $[a, b]$.
Then $\Coker(d_{\mathcal{E}^\bullet}^{a - 1})$ is a flat
$\mathcal{O}_X$-module.
\end{lemma}

\begin{proof}
As $\mathcal{E}^\bullet$ is a bounded above complex of flat modules we see that
$\mathcal{E}^\bullet \otimes_{\mathcal{O}_X} \mathcal{F} =
\mathcal{E}^\bullet \otimes_{\mathcal{O}_X}^{\mathbf{L}} \mathcal{F}$
for any $\mathcal{O}_X$-module $\mathcal{F}$.
Hence for every $\mathcal{O}_X$-module $\mathcal{F}$ the sequence
$$
\mathcal{E}^{a - 2} \otimes_{\mathcal{O}_X} \mathcal{F} \to
\mathcal{E}^{a - 1} \otimes_{\mathcal{O}_X} \mathcal{F} \to
\mathcal{E}^a \otimes_{\mathcal{O}_X} \mathcal{F}
$$
is exact in the middle. Since
$\mathcal{E}^{a - 2} \to \mathcal{E}^{a - 1} \to \mathcal{E}^a \to
\Coker(d^{a - 1}) \to 0$
is a flat resolution this implies that
$\text{Tor}_1^{\mathcal{O}_X}(\Coker(d^{a - 1}), \mathcal{F}) = 0$
for all $\mathcal{O}_X$-modules $\mathcal{F}$. This means that
$\Coker(d^{a - 1})$ is flat, see Lemma \ref{lemma-flat-tor-zero}.
\end{proof}

\begin{lemma}
\label{lemma-tor-amplitude}
Let $(X, \mathcal{O}_X)$ be a ringed space.
Let $E$ be an object of $D(\mathcal{O}_X)$.
Let $a, b \in \mathbf{Z}$ with $a \leq b$. The following are equivalent
\begin{enumerate}
\item $E$ has tor-amplitude in $[a, b]$.
\item $E$ is represented by a complex
$\mathcal{E}^\bullet$ of flat $\mathcal{O}_X$-modules with
$\mathcal{E}^i = 0$ for $i \not \in [a, b]$.
\end{enumerate}
\end{lemma}

\begin{proof}
If (2) holds, then we may compute
$E \otimes_{\mathcal{O}_X}^\mathbf{L} \mathcal{F} =
\mathcal{E}^\bullet \otimes_{\mathcal{O}_X} \mathcal{F}$
and it is clear that (1) holds.

\medskip\noindent
Assume that (1) holds. We may represent $E$ by a bounded above complex
of flat $\mathcal{O}_X$-modules $\mathcal{K}^\bullet$, see
Section \ref{section-flat}.
Let $n$ be the largest integer such that $\mathcal{K}^n \not = 0$.
If $n > b$, then $\mathcal{K}^{n - 1} \to \mathcal{K}^n$ is surjective as
$H^n(\mathcal{K}^\bullet) = 0$. As $\mathcal{K}^n$ is flat we see that
$\Ker(\mathcal{K}^{n - 1} \to \mathcal{K}^n)$ is flat
(Modules, Lemma \ref{modules-lemma-flat-ses}).
Hence we may replace $\mathcal{K}^\bullet$ by
$\tau_{\leq n - 1}\mathcal{K}^\bullet$. Thus, by induction on $n$, we
reduce to the case that $K^\bullet$ is a complex of flat
$\mathcal{O}_X$-modules with $\mathcal{K}^i = 0$ for $i > b$.

\medskip\noindent
Set $\mathcal{E}^\bullet = \tau_{\geq a}\mathcal{K}^\bullet$.
Everything is clear except that $\mathcal{E}^a$ is flat
which follows immediately from Lemma \ref{lemma-last-one-flat}
and the definitions.
\end{proof}

\begin{lemma}
\label{lemma-tor-amplitude-pullback}
Let $f : (X, \mathcal{O}_X) \to (Y, \mathcal{O}_Y)$ be a morphism of ringed
spaces. Let $E$ be an object of $D(\mathcal{O}_Y)$.
If $E$ has tor amplitude in $[a, b]$, then $Lf^*E$ has tor amplitude in
$[a, b]$.
\end{lemma}

\begin{proof}
Assume $E$ has tor amplitude in $[a, b]$. By
Lemma \ref{lemma-tor-amplitude}
we can represent $E$ by a complex of
$\mathcal{E}^\bullet$ of flat $\mathcal{O}$-modules with
$\mathcal{E}^i = 0$ for $i \not \in [a, b]$. Then
$Lf^*E$ is represented by $f^*\mathcal{E}^\bullet$.
By Modules, Lemma \ref{modules-lemma-base-change-flat}
the modules $f^*\mathcal{E}^i$ are flat.
Thus by Lemma \ref{lemma-tor-amplitude}
we conclude that $Lf^*E$ has tor amplitude in $[a, b]$.
\end{proof}

\begin{lemma}
\label{lemma-tor-amplitude-stalk}
Let $(X, \mathcal{O}_X)$ be a ringed space.
Let $E$ be an object of $D(\mathcal{O}_X)$.
Let $a, b \in \mathbf{Z}$ with $a \leq b$. The following are equivalent
\begin{enumerate}
\item $E$ has tor-amplitude in $[a, b]$.
\item for every $x \in X$ the object $E_x$ of $D(\mathcal{O}_{X, x})$
has tor-amplitude in $[a, b]$.
\end{enumerate}
\end{lemma}

\begin{proof}
Taking stalks at $x$ is the same thing as pulling back by the
morphism of ringed spaces $(x, \mathcal{O}_{X, x}) \to (X, \mathcal{O}_X)$.
Hence the implication (1) $\Rightarrow$ (2) follows from
Lemma \ref{lemma-tor-amplitude-pullback}.
For the converse, note that taking stalks commutes with tensor
products (Modules, Lemma \ref{modules-lemma-stalk-tensor-product}).
Hence
$$
(E \otimes_{\mathcal{O}_X}^\mathbf{L} \mathcal{F})_x =
E_x \otimes_{\mathcal{O}_{X, x}}^\mathbf{L} \mathcal{F}_x
$$
On the other hand, taking stalks is exact, so
$$
H^i(E \otimes_{\mathcal{O}_X}^\mathbf{L} \mathcal{F})_x =
H^i((E \otimes_{\mathcal{O}_X}^\mathbf{L} \mathcal{F})_x) =
H^i(E_x \otimes_{\mathcal{O}_{X, x}}^\mathbf{L} \mathcal{F}_x)
$$
and we can check whether
$H^i(E \otimes_{\mathcal{O}_X}^\mathbf{L} \mathcal{F})$ is zero
by checking whether all of its stalks are zero
(Modules, Lemma \ref{modules-lemma-abelian}). Thus (2) implies (1).
\end{proof}

\begin{lemma}
\label{lemma-cone-tor-amplitude}
Let $(X, \mathcal{O}_X)$ be a ringed space.
Let $(K, L, M, f, g, h)$ be a distinguished
triangle in $D(\mathcal{O}_X)$. Let $a, b \in \mathbf{Z}$.
\begin{enumerate}
\item If $K$ has tor-amplitude in $[a + 1, b + 1]$ and
$L$ has tor-amplitude in $[a, b]$ then $M$ has
tor-amplitude in $[a, b]$.
\item If $K$ and $M$ have tor-amplitude in $[a, b]$, then
$L$ has tor-amplitude in $[a, b]$.
\item If $L$ has tor-amplitude in $[a + 1, b + 1]$
and $M$ has tor-amplitude in $[a, b]$, then
$K$ has tor-amplitude in $[a + 1, b + 1]$.
\end{enumerate}
\end{lemma}

\begin{proof}
Omitted. Hint: This just follows from the long exact cohomology sequence
associated to a distinguished triangle and the fact that
$- \otimes_{\mathcal{O}_X}^{\mathbf{L}} \mathcal{F}$
preserves distinguished triangles.
The easiest one to prove is (2) and the others follow from it by
translation.
\end{proof}

\begin{lemma}
\label{lemma-tensor-tor-amplitude}
Let $(X, \mathcal{O}_X)$ be a ringed space. Let $K, L$ be objects of
$D(\mathcal{O}_X)$. If $K$ has tor-amplitude in $[a, b]$ and
$L$ has tor-amplitude in $[c, d]$ then $K \otimes_{\mathcal{O}_X}^\mathbf{L} L$
has tor amplitude in $[a + c, b + d]$.
\end{lemma}

\begin{proof}
Omitted. Hint: use the spectral sequence for tors.
\end{proof}

\begin{lemma}
\label{lemma-summands-tor-amplitude}
Let $(X, \mathcal{O}_X)$ be a ringed space. Let $a, b \in \mathbf{Z}$.
For $K$, $L$ objects of $D(\mathcal{O}_X)$ if $K \oplus L$ has tor
amplitude in $[a, b]$ so do $K$ and $L$.
\end{lemma}

\begin{proof}
Clear from the fact that the Tor functors are additive.
\end{proof}






\section{Perfect complexes}
\label{section-perfect}

\noindent
In this section we discuss properties of perfect complexes on
ringed spaces.

\begin{definition}
\label{definition-perfect}
Let $(X, \mathcal{O}_X)$ be a ringed space.
Let $\mathcal{E}^\bullet$ be a complex of $\mathcal{O}_X$-modules.
We say $\mathcal{E}^\bullet$ is {\it perfect} if there exists
an open covering $X = \bigcup U_i$ such that for each $i$
there exists a morphism of complexes
$\mathcal{E}_i^\bullet \to \mathcal{E}^\bullet|_{U_i}$
which is a quasi-isomorphism with $\mathcal{E}_i^\bullet$
a strictly perfect complex of $\mathcal{O}_{U_i}$-modules.
An object $E$ of $D(\mathcal{O}_X)$ is {\it perfect}
if it can be represented by a perfect complex of $\mathcal{O}_X$-modules.
\end{definition}

\noindent
If $X$ is quasi-compact, then a perfect object of $D(\mathcal{O}_X)$
is in $D^b(\mathcal{O}_X)$. But this need not be the case if
$X$ is not quasi-compact.

\begin{lemma}
\label{lemma-perfect-independent-representative}
Let $(X, \mathcal{O}_X)$ be a ringed space.
Let $E$ be an object of $D(\mathcal{O}_X)$.
\begin{enumerate}
\item If there exists an open covering $X = \bigcup U_i$ and
strictly perfect complexes $\mathcal{E}_i^\bullet$ on $U_i$
such that $\mathcal{E}_i^\bullet$ represents $E|_{U_i}$ in
$D(\mathcal{O}_{U_i})$, then $E$ is perfect.
\item If $E$ is perfect, then any complex representing $E$ is perfect.
\end{enumerate}
\end{lemma}

\begin{proof}
Identical to the proof of
Lemma \ref{lemma-pseudo-coherent-independent-representative}.
\end{proof}

\begin{lemma}
\label{lemma-perfect-on-locally-ringed}
Let $(X, \mathcal{O}_X)$ be a ringed space. Let $E$ be an object of
$D(\mathcal{O}_X)$. Assume that all stalks $\mathcal{O}_{X, x}$
are local rings. Then the following are equivalent
\begin{enumerate}
\item $E$ is perfect,
\item there exists an open covering $X = \bigcup U_i$ such that
$E|_{U_i}$ can be represented by a finite complex of finite locally
free $\mathcal{O}_{U_i}$-modules, and
\item there exists an open covering $X = \bigcup U_i$ such that
$E|_{U_i}$ can be represented by a finite complex of finite
free $\mathcal{O}_{U_i}$-modules.
\end{enumerate}
\end{lemma}

\begin{proof}
This follows from Lemma \ref{lemma-perfect-independent-representative}
and the fact that on $X$ every direct summand of a finite free module
is finite locally free. See Modules, Lemma
\ref{modules-lemma-direct-summand-of-locally-free-is-locally-free}.
\end{proof}

\begin{lemma}
\label{lemma-perfect-precise}
Let $(X, \mathcal{O}_X)$ be a ringed space.
Let $E$ be an object of $D(\mathcal{O}_X)$.
Let $a \leq b$ be integers. If $E$ has tor amplitude in $[a, b]$
and is $(a - 1)$-pseudo-coherent, then $E$ is perfect.
\end{lemma}

\begin{proof}
After replacing $X$ by the members of an open covering we may assume there
exists a strictly perfect complex $\mathcal{E}^\bullet$ and a map
$\alpha : \mathcal{E}^\bullet \to E$ such that $H^i(\alpha)$ is an isomorphism
for $i \geq a$. We may and do replace $\mathcal{E}^\bullet$ by
$\sigma_{\geq a - 1}\mathcal{E}^\bullet$. Choose a distinguished triangle
$$
\mathcal{E}^\bullet \to E \to C \to \mathcal{E}^\bullet[1]
$$
From the vanishing of cohomology sheaves of $E$ and $\mathcal{E}^\bullet$
and the assumption on $\alpha$ we obtain $C \cong \mathcal{K}[2 - a]$ with
$\mathcal{K} = \Ker(\mathcal{E}^{a - 1} \to \mathcal{E}^a)$.
Let $\mathcal{F}$ be an $\mathcal{O}_X$-module.
Applying $- \otimes_{\mathcal{O}_X}^\mathbf{L} \mathcal{F}$
the assumption that $E$ has tor amplitude in $[a, b]$
implies $\mathcal{K} \otimes_{\mathcal{O}_X} \mathcal{F} \to
\mathcal{E}^{a - 1} \otimes_{\mathcal{O}_X} \mathcal{F}$ has image
$\Ker(\mathcal{E}^{a - 1} \otimes_{\mathcal{O}_X} \mathcal{F}
\to \mathcal{E}^a \otimes_{\mathcal{O}_X} \mathcal{F})$.
It follows that $\text{Tor}_1^{\mathcal{O}_X}(\mathcal{E}', \mathcal{F}) = 0$
where $\mathcal{E}' = \Coker(\mathcal{E}^{a - 1} \to \mathcal{E}^a)$.
Hence $\mathcal{E}'$ is flat (Lemma \ref{lemma-flat-tor-zero}).
Thus $\mathcal{E}'$ is locally a direct summand of a finite free module by
Modules, Lemma \ref{modules-lemma-flat-locally-finite-presentation}.
Thus locally the complex
$$
\mathcal{E}' \to \mathcal{E}^{a + 1} \to \ldots \to \mathcal{E}^b
$$
is quasi-isomorphic to $E$ and $E$ is perfect.
\end{proof}

\begin{lemma}
\label{lemma-perfect}
Let $(X, \mathcal{O}_X)$ be a ringed space.
Let $E$ be an object of $D(\mathcal{O}_X)$.
The following are equivalent
\begin{enumerate}
\item $E$ is perfect, and
\item $E$ is pseudo-coherent and locally has finite tor dimension.
\end{enumerate}
\end{lemma}

\begin{proof}
Assume (1). By definition this means there exists an open covering
$X = \bigcup U_i$ such that $E|_{U_i}$ is represented by a
strictly perfect complex. Thus $E$ is pseudo-coherent (i.e.,
$m$-pseudo-coherent for all $m$) by
Lemma \ref{lemma-pseudo-coherent-independent-representative}.
Moreover, a direct summand of a finite free module is flat, hence
$E|_{U_i}$ has finite Tor dimension by
Lemma \ref{lemma-tor-amplitude}. Thus (2) holds.

\medskip\noindent
Assume (2). After replacing $X$ by the members of an open covering
we may assume there exist integers $a \leq b$ such that $E$
has tor amplitude in $[a, b]$. Since $E$ is $m$-pseudo-coherent
for all $m$ we conclude using Lemma \ref{lemma-perfect-precise}.
\end{proof}

\begin{lemma}
\label{lemma-perfect-pullback}
Let $f : (X, \mathcal{O}_X) \to (Y, \mathcal{O}_Y)$ be a morphism of ringed
spaces. Let $E$ be an object of $D(\mathcal{O}_Y)$. If $E$ is perfect in
$D(\mathcal{O}_Y)$, then $Lf^*E$ is perfect in $D(\mathcal{O}_X)$.
\end{lemma}

\begin{proof}
This follows from Lemma \ref{lemma-perfect},
\ref{lemma-tor-amplitude-pullback}, and
\ref{lemma-pseudo-coherent-pullback}.
(An alternative proof is to copy the proof of
Lemma \ref{lemma-pseudo-coherent-pullback}.)
\end{proof}

\begin{lemma}
\label{lemma-two-out-of-three-perfect}
Let $(X, \mathcal{O}_X)$ be a ringed space. Let $(K, L, M, f, g, h)$
be a distinguished triangle in $D(\mathcal{O}_X)$. If two out of three of
$K, L, M$ are perfect then the third is also perfect.
\end{lemma}

\begin{proof}
First proof: Combine
Lemmas \ref{lemma-perfect}, \ref{lemma-cone-pseudo-coherent}, and
\ref{lemma-cone-tor-amplitude}.
Second proof (sketch): Say $K$ and $L$ are perfect. After replacing
$X$ by the members of an open covering we may assume that $K$ and $L$
are represented by strictly perfect complexes $\mathcal{K}^\bullet$
and $\mathcal{L}^\bullet$. After replacing $X$ by the members
of an open covering we may assume the map $K \to L$ is given by
a map of complexes $\alpha : \mathcal{K}^\bullet \to \mathcal{L}^\bullet$,
see Lemma \ref{lemma-local-actual}.
Then $M$ is isomorphic to the cone of $\alpha$ which is strictly
perfect by Lemma \ref{lemma-cone}.
\end{proof}

\begin{lemma}
\label{lemma-tensor-perfect}
Let $(X, \mathcal{O}_X)$ be a ringed space.
If $K, L$ are perfect objects of $D(\mathcal{O}_X)$, then
so is $K \otimes_{\mathcal{O}_X}^\mathbf{L} L$.
\end{lemma}

\begin{proof}
Follows from
Lemmas \ref{lemma-perfect}, \ref{lemma-tensor-pseudo-coherent}, and
\ref{lemma-tensor-tor-amplitude}.
\end{proof}

\begin{lemma}
\label{lemma-summands-perfect}
Let $(X, \mathcal{O}_X)$ be a ringed space.
If $K \oplus L$ is a perfect object of $D(\mathcal{O}_X)$, then
so are $K$ and $L$.
\end{lemma}

\begin{proof}
Follows from
Lemmas \ref{lemma-perfect}, \ref{lemma-summands-pseudo-coherent}, and
\ref{lemma-summands-tor-amplitude}.
\end{proof}

\begin{lemma}
\label{lemma-pushforward-perfect}
Let $(X, \mathcal{O}_X)$ be a ringed space. Let $j : U \to X$ be an
open subspace. Let $E$ be a perfect object of $D(\mathcal{O}_U)$
whose cohomology
sheaves are supported on a closed subset $T \subset U$ with $j(T)$
closed in $X$. Then $Rj_*E$ is a perfect object of $D(\mathcal{O}_X)$.
\end{lemma}

\begin{proof}
Being a perfect complex is local on $X$. Thus it suffices to check that
$Rj_*E$ is perfect when restricted to $U$ and $V = X \setminus j(T)$.
We have $Rj_*E|_U = E$ which is perfect. We have
 $Rj_*E|_V = 0$ because $E|_{U \setminus T} = 0$.
\end{proof}

\begin{lemma}
\label{lemma-perfect-max-open-coh-loc-free}
Let $(X, \mathcal{O}_X)$ be a ringed space. Let $E$ in $D(\mathcal{O}_X)$
be perfect. Assume that all stalks $\mathcal{O}_{X, x}$ are local rings.
Then the set
$$
U =
\{x \in X \mid
H^i(E)_x\text{ is a finite free }
\mathcal{O}_{X, x}\text{-module for all }i\in \mathbf{Z}\}
$$
is open in $X$ and is the maximal open set $U \subset X$ such that
$H^i(E)|_U$ is finite locally free for all $i \in \mathbf{Z}$.
\end{lemma}

\begin{proof}
Note that if $V \subset X$ is some open such that $H^i(E)|_V$
is finite locally free for all $i \in \mathbf{Z}$ then $V \subset U$.
Let $x \in U$. We will show that an open neighbourhood of $x$
is contained in $U$ and that $H^i(E)$ is finite locally free
on this neighbourhood for all $i$. This will finish the proof.
During the proof we may (finitely many times) replace
$X$ by an open neighbourhood of $x$.
Hence we may assume $E$ is represented
by a strictly perfect complex $\mathcal{E}^\bullet$. Say
$\mathcal{E}^i = 0$ for $i \not \in [a, b]$. We will
prove the result by induction on $b - a$. The module
$H^b(E) = \Coker(d^{b - 1} : \mathcal{E}^{b - 1} \to \mathcal{E}^b)$
is of finite presentation. Since $H^b(E)_x$ is finite free,
we conclude $H^b(E)$ is finite free in an open neighbourhood of $x$ by
Modules, Lemma \ref{modules-lemma-finite-presentation-stalk-free}.
Thus after replacing $X$ by a (possibly smaller) open
neighbourhood we may assume we have a direct sum decomposition
$\mathcal{E}^b = \Im(d^{b - 1}) \oplus H^b(E)$ and
$H^b(E)$ is finite free, see Lemma \ref{lemma-local-lift-map}.
Doing the same argument again, we see that we may
assume $\mathcal{E}^{b - 1} = \Ker(d^{b - 1}) \oplus \Im(d^{b - 1})$.
The complex
$\mathcal{E}^a \to \ldots \to \mathcal{E}^{b - 2} \to \Ker(d^{b - 1})$
is a strictly perfect complex representing a perfect
object $E'$ with $H^i(E) = H^i(E')$ for $i \not = b$.
Hence we conclude by our induction hypothesis.
\end{proof}






\section{Duals}
\label{section-duals}

\noindent
In this section we characterize the dualizable objects of
the category of complexes and of the derived category.
In particular, we will see that an object of $D(\mathcal{O}_X)$
has a dual if and only if it is perfect (this follows from
Example \ref{example-dual-derived} and
Lemma \ref{lemma-left-dual-derived}).

\begin{lemma}
\label{lemma-symmetric-monoidal-cat-complexes}
Let $(X, \mathcal{O}_X)$ be a ringed space. The category of complexes
of $\mathcal{O}_X$-modules with tensor product defined by
$\mathcal{F}^\bullet \otimes \mathcal{G}^\bullet =
\text{Tot}(\mathcal{F}^\bullet \otimes_{\mathcal{O}_X} \mathcal{G}^\bullet)$
is a symmetric monoidal category (for sign rules, see
More on Algebra, Section \ref{more-algebra-section-sign-rules}).
\end{lemma}

\begin{proof}
Omitted. Hints: as unit $\mathbf{1}$ we take the complex having
$\mathcal{O}_X$ in degree $0$ and zero in other degrees with
obvious isomorphisms
$\text{Tot}(\mathbf{1} \otimes_{\mathcal{O}_X} \mathcal{G}^\bullet) =
\mathcal{G}^\bullet$ and
$\text{Tot}(\mathcal{F}^\bullet \otimes_{\mathcal{O}_X} \mathbf{1}) =
\mathcal{F}^\bullet$.
to prove the lemma you have to check the commutativity
of various diagrams, see Categories, Definitions
\ref{categories-definition-monoidal-category} and
\ref{categories-definition-symmetric-monoidal-category}.
The verifications are straightforward in each case.
\end{proof}

\begin{example}
\label{example-dual}
Let $(X, \mathcal{O}_X)$ be a ringed space. Let $\mathcal{F}^\bullet$
be a locally bounded complex of $\mathcal{O}_X$-modules such that each
$\mathcal{F}^n$ is locally a direct summand of a finite
free $\mathcal{O}_X$-module. In other words, there is an open covering
$X = \bigcup U_i$ such that $\mathcal{F}^\bullet|_{U_i}$ is a strictly
perfect complex. Consider the complex
$$
\mathcal{G}^\bullet = \SheafHom^\bullet(\mathcal{F}^\bullet, \mathcal{O}_X)
$$
as in Section \ref{section-hom-complexes}. Let
$$
\eta :
\mathcal{O}_X
\to
\text{Tot}(\mathcal{F}^\bullet \otimes_{\mathcal{O}_X} \mathcal{G}^\bullet)
\quad\text{and}\quad
\epsilon :
\text{Tot}(\mathcal{G}^\bullet \otimes_{\mathcal{O}_X} \mathcal{F}^\bullet)
\to
\mathcal{O}_X
$$
be $\eta = \sum \eta_n$ and $\epsilon = \sum \epsilon_n$
where $\eta_n : \mathcal{O}_X \to
\mathcal{F}^n \otimes_{\mathcal{O}_X} \mathcal{G}^{-n}$
and
$\epsilon_n : \mathcal{G}^{-n} \otimes_{\mathcal{O}_X} \mathcal{F}^n
\to \mathcal{O}_X$ are as in Modules, Example \ref{modules-example-dual}.
Then $\mathcal{G}^\bullet, \eta, \epsilon$
is a left dual for $\mathcal{F}^\bullet$ as in
Categories, Definition \ref{categories-definition-dual}.
We omit the verification that
$(1 \otimes \epsilon) \circ (\eta \otimes 1) = \text{id}_{\mathcal{F}^\bullet}$
and
$(\epsilon \otimes 1) \circ (1 \otimes \eta) =
\text{id}_{\mathcal{G}^\bullet}$. Please compare with
More on Algebra, Lemma \ref{more-algebra-lemma-left-dual-complex}.
\end{example}

\begin{lemma}
\label{lemma-left-dual-complex}
Let $(X, \mathcal{O}_X)$ be a ringed space. Let $\mathcal{F}^\bullet$
be a complex of $\mathcal{O}_X$-modules. If $\mathcal{F}^\bullet$
has a left dual in the monoidal category of complexes of
$\mathcal{O}_X$-modules
(Categories, Definition \ref{categories-definition-dual})
then $\mathcal{F}^\bullet$ is a locally bounded complex whose terms are
locally direct summands of finite free $\mathcal{O}_X$-modules
and the left dual is as constructed in Example \ref{example-dual}.
\end{lemma}

\begin{proof}
By uniqueness of left duals
(Categories, Remark \ref{categories-remark-left-dual-adjoint})
we get the final statement provided we show that $\mathcal{F}^\bullet$
is as stated. Let $\mathcal{G}^\bullet, \eta, \epsilon$ be a left dual.
Write $\eta = \sum \eta_n$ and $\epsilon = \sum \epsilon_n$
where $\eta_n : \mathcal{O}_X \to
\mathcal{F}^n \otimes_{\mathcal{O}_X} \mathcal{G}^{-n}$
and
$\epsilon_n : \mathcal{G}^{-n} \otimes_{\mathcal{O}_X} \mathcal{F}^n
\to \mathcal{O}_X$. Since
$(1 \otimes \epsilon) \circ (\eta \otimes 1) = \text{id}_{\mathcal{F}^\bullet}$
and
$(\epsilon \otimes 1) \circ (1 \otimes \eta) = \text{id}_{\mathcal{G}^\bullet}$
by Categories, Definition \ref{categories-definition-dual} we see immediately
that we have
$(1 \otimes \epsilon_n) \circ (\eta_n \otimes 1) = \text{id}_{\mathcal{F}^n}$
and
$(\epsilon_n \otimes 1) \circ (1 \otimes \eta_n) =
\text{id}_{\mathcal{G}^{-n}}$.
Hence we see that $\mathcal{F}^n$ is locally a direct summand of a finite
free $\mathcal{O}_X$-module by
Modules, Lemma \ref{modules-lemma-left-dual-module}.
Since the sum $\eta = \sum \eta_n$ is locally finite, we conclude that
$\mathcal{F}^\bullet$ is locally bounded.
\end{proof}

\begin{lemma}
\label{lemma-internal-hom-evaluate-isom}
Let $(X, \mathcal{O}_X)$ be a ringed space. Let $K, L, M \in D(\mathcal{O}_X)$.
If $K$ is perfect, then the map
$$
R\SheafHom(L, M) \otimes_{\mathcal{O}_X}^\mathbf{L} K
\longrightarrow
R\SheafHom(R\SheafHom(K, L), M)
$$
of Lemma \ref{lemma-internal-hom-evaluate} is an isomorphism.
\end{lemma}

\begin{proof}
Since the map is globally defined and since formation of the right and
left hand side commute with localization
(see Lemma \ref{lemma-restriction-RHom-to-U}), to prove this we may work
locally on $X$. Thus we may assume $K$ is represented by a strictly
perfect complex $\mathcal{E}^\bullet$.

\medskip\noindent
If $K_1 \to K_2 \to K_3$ is a distinguished triangle in $D(\mathcal{O}_X)$,
then we get distinguished triangles
$$
R\SheafHom(L, M) \otimes_{\mathcal{O}_X}^\mathbf{L} K_1 \to
R\SheafHom(L, M) \otimes_{\mathcal{O}_X}^\mathbf{L} K_2 \to
R\SheafHom(L, M) \otimes_{\mathcal{O}_X}^\mathbf{L} K_3
$$
and
$$
R\SheafHom(R\SheafHom(K_1, L), M) \to
R\SheafHom(R\SheafHom(K_2, L), M)
R\SheafHom(R\SheafHom(K_3, L), M)
$$
See Section \ref{section-flat} and
Lemma \ref{lemma-RHom-triangulated}.
The arrow of Lemma \ref{lemma-internal-hom-evaluate} is functorial in $K$
hence we get a morphism between these distinguished triangles.
Thus, if the result holds for $K_1$ and $K_3$, then the result holds for
$K_2$ by Derived Categories, Lemma
\ref{derived-lemma-third-isomorphism-triangle}.

\medskip\noindent
Combining the remarks above with the distinguished triangles
$$
\sigma_{\geq n}\mathcal{E}^\bullet \to \mathcal{E}^\bullet \to
\sigma_{\leq n - 1}\mathcal{E}^\bullet
$$
of stupid trunctions, we reduce to the case where $K$ consists
of a direct summand of a finite free $\mathcal{O}_X$-module placed
in some degree. By an obvious compatibility of the problem with direct sums
(similar to what was said above) and shifts this reduces us to the case
where $K = \mathcal{O}_X^{\oplus n}$ for some integer $n$.
This case is clear.
\end{proof}

\begin{lemma}
\label{lemma-dual-perfect-complex}
Let $(X, \mathcal{O}_X)$ be a ringed space. Let $K$ be a perfect object of
$D(\mathcal{O}_X)$. Then $K^\vee = R\SheafHom(K, \mathcal{O}_X)$ is a
perfect object too and $(K^\vee)^\vee \cong K$. There are
functorial isomorphisms
$$
M \otimes^\mathbf{L}_{\mathcal{O}_X} K^\vee = R\SheafHom(K, M)
$$
and
$$
H^0(X, M \otimes^\mathbf{L}_{\mathcal{O}_X} K^\vee) =
\Hom_{D(\mathcal{O}_X)}(K, M)
$$
for $M$ in $D(\mathcal{O}_X)$.
\end{lemma}

\begin{proof}
By Lemma \ref{lemma-internal-hom-evaluate} there is a canonical map
$$
K = R\SheafHom(\mathcal{O}_X, \mathcal{O}_X)
\otimes_{\mathcal{O}_X}^\mathbf{L} K \longrightarrow
R\SheafHom(R\SheafHom(K, \mathcal{O}_X), \mathcal{O}_X) =
(K^\vee)^\vee
$$
which is an isomorphism by Lemma \ref{lemma-internal-hom-evaluate-isom}.
To check the other statements we will use without further mention that
formation of internal hom commutes with restriction to opens
(Lemma \ref{lemma-restriction-RHom-to-U}).
We may check $K^\vee$ is perfect locally on $X$.
By Lemma \ref{lemma-dual}
to see the final statement it suffices to check that the map
(\ref{equation-eval})
$$
M \otimes^\mathbf{L}_{\mathcal{O}_X} K^\vee
\longrightarrow
R\SheafHom(K, M)
$$
is an isomorphism. This is local on $X$ as well.
Hence it suffices to prove these two statements $K$ is represented
by a strictly perfect complex.

\medskip\noindent
Assume $K$ is represented by the strictly perfect complex
$\mathcal{E}^\bullet$. Then it follows from
Lemma \ref{lemma-Rhom-strictly-perfect}
that $K^\vee$ is represented by the complex whose terms are
$(\mathcal{E}^{-n})^\vee =
\SheafHom_{\mathcal{O}_X}(\mathcal{E}^{-n}, \mathcal{O}_X)$
in degree $n$. Since $\mathcal{E}^{-n}$ is a direct summand of a finite
free $\mathcal{O}_X$-module, so is $(\mathcal{E}^{-n})^\vee$.
Hence $K^\vee$ is represented by a strictly perfect complex too
and we see that $K^\vee$ is perfect.
To see that (\ref{equation-eval}) is an isomorphism, represent
$M$ by a complex $\mathcal{F}^\bullet$.
By Lemma \ref{lemma-Rhom-strictly-perfect} the complex
$R\SheafHom(K, M)$ is represented by the complex with terms
$$
\bigoplus\nolimits_{n = p + q}
\SheafHom_{\mathcal{O}_X}(\mathcal{E}^{-q}, \mathcal{F}^p)
$$
On the other hand, the object $M \otimes^\mathbf{L}_{\mathcal{O}_X} K^\vee$
is represented by the complex with terms
$$
\bigoplus\nolimits_{n = p + q}
\mathcal{F}^p \otimes_{\mathcal{O}_X} (\mathcal{E}^{-q})^\vee
$$
Thus the assertion that (\ref{equation-eval}) is an isomorphism
reduces to the assertion that the canonical map
$$
\mathcal{F}
\otimes_{\mathcal{O}_X}
\SheafHom_{\mathcal{O}_X}(\mathcal{E}, \mathcal{O}_X)
\longrightarrow
\SheafHom_{\mathcal{O}_X}(\mathcal{E}, \mathcal{F})
$$
is an isomorphism when $\mathcal{E}$ is a direct summand of a finite
free $\mathcal{O}_X$-module and $\mathcal{F}$ is any $\mathcal{O}_X$-module.
This follows immediately from the corresponding statement when
$\mathcal{E}$ is finite free.
\end{proof}

\begin{lemma}
\label{lemma-symmetric-monoidal-derived}
Let $(X, \mathcal{O}_X)$ be a ringed space. The derived category
$D(\mathcal{O}_X)$ is a symmetric monoidal category with tensor product
given by derived tensor product with usual associativity and
commutativity constraints (for sign rules, see
More on Algebra, Section \ref{more-algebra-section-sign-rules}).
\end{lemma}

\begin{proof}
Omitted. Compare with Lemma \ref{lemma-symmetric-monoidal-cat-complexes}.
\end{proof}

\begin{example}
\label{example-dual-derived}
Let $(X, \mathcal{O}_X)$ be a ringed space. Let $K$ be a perfect object
of $D(\mathcal{O}_X)$. Set $K^\vee = R\SheafHom(K, \mathcal{O}_X)$
as in Lemma \ref{lemma-dual-perfect-complex}.
Then the map
$$
K \otimes_{\mathcal{O}_X}^\mathbf{L} K^\vee \longrightarrow R\SheafHom(K, K)
$$
is an isomorphism (by the lemma). Denote
$$
\eta :
\mathcal{O}_X
\longrightarrow
K \otimes_{\mathcal{O}_X}^\mathbf{L} K^\vee
$$
the map sending $1$ to the section corresponding to
$\text{id}_K$ under the isomorphism above.
Denote
$$
\epsilon :
K^\vee
\otimes_{\mathcal{O}_X}^\mathbf{L} K
\longrightarrow
\mathcal{O}_X
$$
the evaluation map (to construct it you can use
Lemma \ref{lemma-internal-hom-composition} for example). Then
$K^\vee, \eta, \epsilon$ is a left dual for $K$ as in
Categories, Definition \ref{categories-definition-dual}.
We omit the verification that
$(1 \otimes \epsilon) \circ (\eta \otimes 1) = \text{id}_K$
and
$(\epsilon \otimes 1) \circ (1 \otimes \eta) =
\text{id}_{K^\vee}$.
\end{example}

\begin{lemma}
\label{lemma-left-dual-derived}
Let $(X, \mathcal{O}_X)$ be a ringed space. Let $M$ be an object
of $D(\mathcal{O}_X)$. If $M$ has a left dual in the monoidal category
$D(\mathcal{O}_X)$ (Categories, Definition \ref{categories-definition-dual})
then $M$ is perfect and the left dual is as constructed in
Example \ref{example-dual-derived}.
\end{lemma}

\begin{proof}
Let $x \in X$. It suffices to find an open neighbourhood $U$ of $x$
such that $M$ restricts to a perfect complex over $U$. Hence during the
proof we can (finitely often) replace $X$ by an open neighbourhood of $x$.
Let $N, \eta, \epsilon$ be a left dual.

\medskip\noindent
We are going to use the following argument several times. Choose any
complex $\mathcal{M}^\bullet$
of $\mathcal{O}_X$-modules representing $M$. Choose a K-flat complex
$\mathcal{N}^\bullet$ representing $N$ whose terms are flat
$\mathcal{O}_X$-modules, see Lemma \ref{lemma-K-flat-resolution}.
Consider the map
$$
\eta : \mathcal{O}_X \to
\text{Tot}(\mathcal{M}^\bullet \otimes_{\mathcal{O}_X} \mathcal{N}^\bullet)
$$
After shrinking $X$ we can find an integer $N$ and for
$i = 1, \ldots, N$ integers $n_i \in \mathbf{Z}$ and sections
$f_i$ and $g_i$ of $\mathcal{M}^{n_i}$ and $\mathcal{N}^{-n_i}$
such that
$$
\eta(1) = \sum\nolimits_i f_i \otimes g_i
$$
Let $\mathcal{K}^\bullet \subset \mathcal{M}^\bullet$ be any subcomplex
of $\mathcal{O}_X$-modules containing the sections $f_i$
for $i = 1, \ldots, N$.
Since
$\text{Tot}(\mathcal{K}^\bullet \otimes_{\mathcal{O}_X} \mathcal{N}^\bullet)
\subset
\text{Tot}(\mathcal{M}^\bullet \otimes_{\mathcal{O}_X} \mathcal{N}^\bullet)$
by flatness of the modules $\mathcal{N}^n$, we see that $\eta$ factors through
$$
\tilde \eta :
\mathcal{O}_X \to
\text{Tot}(\mathcal{K}^\bullet \otimes_{\mathcal{O}_X} \mathcal{N}^\bullet)
$$
Denoting $K$ the object of $D(\mathcal{O}_X)$ represented by
$\mathcal{K}^\bullet$ we find a commutative diagram
$$
\xymatrix{
M \ar[rr]_-{\eta \otimes 1} \ar[rrd]_{\tilde \eta \otimes 1} & &
M \otimes^\mathbf{L} N \otimes^\mathbf{L} M
\ar[r]_-{1 \otimes \epsilon} &
M \\
& &
K \otimes^\mathbf{L} N \otimes^\mathbf{L} M
\ar[u] \ar[r]^-{1 \otimes \epsilon} &
K \ar[u]
}
$$
Since the composition of the upper row is the identity on $M$
we conclude that $M$ is a direct summand of $K$ in $D(\mathcal{O}_X)$.

\medskip\noindent
As a first use of the argument above, we can choose the subcomplex
$\mathcal{K}^\bullet = \sigma_{\geq a} \tau_{\leq b}\mathcal{M}^\bullet$
with $a < n_i < b$ for $i = 1, \ldots, N$. Thus $M$ is a direct
summand in $D(\mathcal{O}_X)$ of a bounded complex and we conclude
we may assume $M$ is in $D^b(\mathcal{O}_X)$. (Recall that the process
above involves shrinking $X$.)

\medskip\noindent
Since $M$ is in $D^b(\mathcal{O}_X)$ we may choose
$\mathcal{M}^\bullet$ to be a bounded above complex of
flat modules (by Modules, Lemma \ref{modules-lemma-module-quotient-flat} and
Derived Categories, Lemma \ref{derived-lemma-subcategory-left-resolution}).
Then we can choose $\mathcal{K}^\bullet = \sigma_{\geq a}\mathcal{M}^\bullet$
with $a < n_i$ for $i = 1, \ldots, N$ in the argument above.
Thus we find that we may assume $M$ is a direct summand in
$D(\mathcal{O}_X)$ of a bounded complex of flat modules.
In particular,  $M$ has finite tor amplitude.

\medskip\noindent
Say $M$ has tor amplitude in $[a, b]$. Assuming $M$ is $m$-pseudo-coherent
we are going to show that (after shrinking $X$) we may assume $M$
is $(m - 1)$-pseudo-coherent. This will finish the proof by
Lemma \ref{lemma-perfect-precise} and the fact that
$M$ is $(b + 1)$-pseudo-coherent in any case.
After shrinking $X$ we may assume there exists a strictly perfect
complex $\mathcal{E}^\bullet$ and a map $\alpha : \mathcal{E}^\bullet \to M$
in $D(\mathcal{O}_X)$ such that $H^i(\alpha)$ is an isomorphism for
$i > m$ and surjective for $i = m$. We may and do assume
that $\mathcal{E}^i = 0$ for $i < m$. Choose a distinguished triangle
$$
\mathcal{E}^\bullet \to M \to L \to \mathcal{E}^\bullet[1]
$$
Observe that $H^i(L) = 0$ for $i \geq m$. Thus we may represent
$L$ by a complex $\mathcal{L}^\bullet$ with $\mathcal{L}^i = 0$
for $i \geq m$. The map $L \to \mathcal{E}^\bullet[1]$
is given by a map of complexes
$\mathcal{L}^\bullet \to \mathcal{E}^\bullet[1]$
which is zero in all degrees except in degree $m - 1$
where we obtain a map $\mathcal{L}^{m - 1} \to \mathcal{E}^m$, see
Derived Categories, Lemma \ref{derived-lemma-negative-exts}.
Then $M$ is represented by the complex
$$
\mathcal{M}^\bullet :
\ldots \to
\mathcal{L}^{m - 2} \to
\mathcal{L}^{m - 1} \to
\mathcal{E}^m \to
\mathcal{E}^{m + 1} \to \ldots
$$
Apply the discussion in the second paragraph to this complex to get
sections $f_i$ of $\mathcal{M}^{n_i}$ for $i = 1, \ldots, N$.
For $n < m$ let $\mathcal{K}^n \subset \mathcal{L}^n$
be the $\mathcal{O}_X$-submodule generated by the sections
$f_i$ for $n_i = n$ and $d(f_i)$ for $n_i = n - 1$.
For $n \geq m$ set $\mathcal{K}^n = \mathcal{E}^n$.
Clearly, we have a morphism of
distinguished triangles
$$
\xymatrix{
\mathcal{E}^\bullet \ar[r] &
\mathcal{M}^\bullet \ar[r] &
\mathcal{L}^\bullet \ar[r] &
\mathcal{E}^\bullet[1] \\
\mathcal{E}^\bullet \ar[r] \ar[u] &
\mathcal{K}^\bullet \ar[r] \ar[u] &
\sigma_{\leq m - 1}\mathcal{K}^\bullet \ar[r] \ar[u] &
\mathcal{E}^\bullet[1] \ar[u]
}
$$
where all the morphisms are as indicated above.
Denote $K$ the object of $D(\mathcal{O}_X)$ corresponding to the complex
$\mathcal{K}^\bullet$.
By the arguments in the second paragraph of the proof we obtain
a morphism $s : M \to K$ in $D(\mathcal{O}_X)$ such that the composition
$M \to K \to M$ is the identity on $M$. We don't know that the
diagram
$$
\xymatrix{
\mathcal{E}^\bullet \ar[r] &
\mathcal{K}^\bullet \ar@{=}[r] &
K \\
\mathcal{E}^\bullet \ar[u]^{\text{id}} \ar[r]^i &
\mathcal{M}^\bullet \ar@{=}[r] &
M \ar[u]_s
}
$$
commutes, but we do know it commutes after composing with the
map $K \to M$. By Lemma \ref{lemma-local-actual} after shrinking $X$ we may
assume that $s \circ i$ is given by a map of complexes
$\sigma : \mathcal{E}^\bullet \to \mathcal{K}^\bullet$.
By the same lemma we may assume the composition of $\sigma$
with the inclusion $\mathcal{K}^\bullet \subset \mathcal{M}^\bullet$
is homotopic to zero by some homotopy
$\{h^i : \mathcal{E}^i \to \mathcal{M}^{i - 1}\}$.
Thus, after replacing $\mathcal{K}^{m - 1}$ by
$\mathcal{K}^{m - 1} + \Im(h^m)$ (note that after doing this
it is still the case that $\mathcal{K}^{m - 1}$ is generated
by finitely many global sections), we see that
$\sigma$ itself is homotopic to zero!
This means that we have a commutative solid diagram
$$
\xymatrix{
\mathcal{E}^\bullet \ar[r] &
M \ar[r] &
\mathcal{L}^\bullet \ar[r] &
\mathcal{E}^\bullet[1] \\
\mathcal{E}^\bullet \ar[r] \ar[u] &
K \ar[r] \ar[u] &
\sigma_{\leq m - 1}\mathcal{K}^\bullet \ar[r] \ar[u] &
\mathcal{E}^\bullet[1] \ar[u] \\
\mathcal{E}^\bullet \ar[r] \ar[u] &
M \ar[r] \ar[u]^s &
\mathcal{L}^\bullet \ar[r] \ar@{..>}[u] &
\mathcal{E}^\bullet[1] \ar[u]
}
$$
By the axioms of triangulated categories we obtain a dotted
arrow fitting into the diagram.
Looking at cohomology sheaves in degree $m - 1$ we see that we obtain
$$
\xymatrix{
H^{m - 1}(M) \ar[r] &
H^{m - 1}(\mathcal{L}^\bullet) \ar[r] &
H^m(\mathcal{E}^\bullet) \\
H^{m - 1}(K) \ar[r] \ar[u] &
H^{m - 1}(\sigma_{\leq m - 1}\mathcal{K}^\bullet) \ar[r] \ar[u] &
H^m(\mathcal{E}^\bullet) \ar[u] \\
H^{m - 1}(M) \ar[r] \ar[u] &
H^{m - 1}(\mathcal{L}^\bullet) \ar[r] \ar[u] &
H^m(\mathcal{E}^\bullet) \ar[u]
}
$$
Since the vertical compositions are the identity in both the
left and right column, we conclude the vertical composition
$H^{m - 1}(\mathcal{L}^\bullet) \to
H^{m - 1}(\sigma_{\leq m - 1}\mathcal{K}^\bullet) \to
H^{m - 1}(\mathcal{L}^\bullet)$ in the middle is surjective!
In particular $H^{m - 1}(\sigma_{\leq m - 1}\mathcal{K}^\bullet) \to
H^{m - 1}(\mathcal{L}^\bullet)$ is surjective.
Using the induced map of long exact sequences of cohomology
sheaves from the morphism of triangles above, a diagram chase
shows this implies $H^i(K) \to H^i(M)$ is an isomorphism
for $i \geq m$ and surjective for $i = m - 1$.
By construction we can choose an $r \geq 0$ and a surjection
$\mathcal{O}_X^{\oplus r} \to \mathcal{K}^{m - 1}$. Then the
composition
$$
(\mathcal{O}_X^{\oplus r} \to \mathcal{E}^m \to
\mathcal{E}^{m + 1} \to \ldots ) \longrightarrow
K \longrightarrow M
$$
induces an isomorphism on cohomology sheaves in degrees $\geq m$ and
a surjection in degree $m - 1$ and the proof is complete.
\end{proof}





\section{Miscellany}
\label{section-misc}

\noindent
Some results which do not fit anywhere else.

\begin{lemma}
\label{lemma-colim-and-lim-of-duals}
Let $(X, \mathcal{O}_X)$ be a ringed space. Let
$(K_n)_{n \in \mathbf{N}}$ be a system of perfect objects of $D(\mathcal{O}_X)$.
Let $K = \text{hocolim} K_n$ be the derived colimit
(Derived Categories, Definition \ref{derived-definition-derived-colimit}).
Then for any object $E$ of $D(\mathcal{O}_X)$ we have
$$
R\SheafHom(K, E) = R\lim E \otimes^\mathbf{L}_{\mathcal{O}_X} K_n^\vee
$$
where $(K_n^\vee)$ is the inverse system of dual perfect complexes.
\end{lemma}

\begin{proof}
By Lemma \ref{lemma-dual-perfect-complex} we have
$R\lim E \otimes^\mathbf{L}_{\mathcal{O}_X} K_n^\vee =
R\lim R\SheafHom(K_n, E)$
which fits into the distinguished triangle
$$
R\lim R\SheafHom(K_n, E) \to
\prod R\SheafHom(K_n, E) \to
\prod R\SheafHom(K_n, E)
$$
Because $K$ similarly fits into the distinguished triangle
$\bigoplus K_n \to \bigoplus K_n \to K$ it suffices to show that
$\prod R\SheafHom(K_n, E) = R\SheafHom(\bigoplus K_n, E)$.
This is a formal consequence of (\ref{equation-internal-hom})
and the fact that derived tensor product commutes with direct sums.
\end{proof}

\begin{lemma}
\label{lemma-ext-composition-is-cup}
Let $(X, \mathcal{O}_X)$ be a ringed space. Let $K$ and $E$ be objects
of $D(\mathcal{O}_X)$ with $E$ perfect. The diagram
$$
\xymatrix{
H^0(X, K \otimes_{\mathcal{O}_X}^\mathbf{L} E^\vee) \times H^0(X, E)
\ar[r] \ar[d] &
H^0(X, K \otimes_{\mathcal{O}_X}^\mathbf{L} E^\vee
\otimes_{\mathcal{O}_X}^\mathbf{L} E) \ar[d] \\
\Hom_X(E, K) \times H^0(X, E) \ar[r] &
H^0(X, K)
}
$$
commutes where the top horizontal arrow is the cup product, the
right vertical arrow uses
$\epsilon : E^\vee \otimes_{\mathcal{O}_X}^\mathbf{L} E \to \mathcal{O}_X$
(Example \ref{example-dual-derived}), the left vertical arrow uses
Lemma \ref{lemma-dual-perfect-complex}, and the bottom horizontal
arrow is the obvious one.
\end{lemma}

\begin{proof}
We will abbreviate $\otimes = \otimes_{\mathcal{O}_X}^\mathbf{L}$
and $\mathcal{O} = \mathcal{O}_X$. We will identify $E$ and $K$
with $R\SheafHom(\mathcal{O}, E)$ and $R\SheafHom(\mathcal{O}, K)$
and we will identify $E^\vee$ with $R\SheafHom(E, \mathcal{O})$.

\medskip\noindent
Let $\xi \in H^0(X, K \otimes E^\vee)$ and $\eta \in H^0(X, E)$.
Denote $\tilde \xi : \mathcal{O} \to K \otimes E^\vee$ and
$\tilde \eta : \mathcal{O} \to E$ the corresponding maps in
$D(\mathcal{O})$. By Lemma \ref{lemma-second-cup-equals-first}
the cup product $\xi \cup \eta$ corresponds to
$\tilde \xi \otimes \tilde \eta : \mathcal{O} \to
K \otimes E^\vee \otimes E$.

\medskip\noindent
We claim the map $\xi' : E \to K$ corresponding to $\xi$ by
Lemma \ref{lemma-dual-perfect-complex} is the composition
$$
E = \mathcal{O} \otimes E
\xrightarrow{\tilde \xi \otimes 1_E}
K \otimes E^\vee \otimes E
\xrightarrow{1_K \otimes \epsilon}
K
$$
The construction in Lemma \ref{lemma-dual-perfect-complex}
uses the evaluation map (\ref{equation-eval}) which in turn
is constructed using the identification of $E$ with
$R\SheafHom(\mathcal{O}, E)$ and the composition
$\underline{\circ}$ constructed
in Lemma \ref{lemma-internal-hom-composition}.
Hence $\xi'$ is the composition
\begin{align*}
E = \mathcal{O} \otimes
R\SheafHom(\mathcal{O}, E)
& \xrightarrow{\tilde \xi \otimes 1}
R\SheafHom(\mathcal{O}, K) \otimes
R\SheafHom(E, \mathcal{O}) \otimes
R\SheafHom(\mathcal{O}, E) \\
& \xrightarrow{\underline{\circ} \otimes 1}
R\SheafHom(E, K) \otimes R\SheafHom(\mathcal{O}, E) \\
& \xrightarrow{\underline{\circ}}
R\SheafHom(\mathcal{O}, K) = K
\end{align*}
The claim follows immediately from this and the fact that
the composition $\underline{\circ}$ constructed in
Lemma \ref{lemma-internal-hom-composition} is associative
(insert future reference here) and the fact that $\epsilon$
is defined as the composition
$\underline{\circ} : E^\vee \otimes E \to \mathcal{O}$ in
Example \ref{example-dual-derived}.

\medskip\noindent
Using the results from the previous two paragraphs, we find
the statement of the lemma is that
$(1_K \otimes \epsilon) \circ (\tilde \xi \otimes \tilde \eta)$
is equal to
$(1_K \otimes \epsilon) \circ (\tilde \xi \otimes 1_E)
\circ (1_\mathcal{O} \otimes \tilde \eta)$
which is immediate.
\end{proof}

\begin{lemma}
\label{lemma-pullback-internal-hom}
Let $h : X \to Y$ be a morphism of ringed spaces.
Let $K, M$ be objects of $D(\mathcal{O}_Y)$. The
canonical map
$$
Lh^*R\SheafHom(K, M) \longrightarrow R\SheafHom(Lh^*K, Lh^*M)
$$
of Remark \ref{remark-prepare-fancy-base-change}
is an isomorphism in the following cases
\begin{enumerate}
\item $K$ is perfect,
\item $h$ is flat, $K$ is pseudo-coherent, and $M$ is (locally) bounded below,
\item $\mathcal{O}_X$ has finite tor dimension over $h^{-1}\mathcal{O}_Y$,
$K$ is pseudo-coherent, and $M$ is (locally) bounded below,
\end{enumerate}
\end{lemma}

\begin{proof}
Proof of (1). The question is local on $Y$, hence we may assume that
$K$ is represented by a strictly perfect complex
$\mathcal{E}^\bullet$, see Section \ref{section-perfect}.
Choose a K-flat complex $\mathcal{F}^\bullet$ representing $M$.
Apply Lemma \ref{lemma-Rhom-strictly-perfect} to see that
$R\SheafHom(K, L)$ is represented by the complex
$\mathcal{H}^\bullet =
\SheafHom^\bullet(\mathcal{E}^\bullet, \mathcal{F}^\bullet)$
with terms $\mathcal{H}^n = \bigoplus\nolimits_{n = p + q}
\SheafHom_{\mathcal{O}_X}(\mathcal{E}^{-q}, \mathcal{F}^p)$.
By the construction of $Lh^*$ in Section \ref{section-derived-pullback}
we see that $Lh^*K$ is represented by the strictly perfect complex
$h^*\mathcal{E}^\bullet$ (Lemma \ref{lemma-strictly-perfect-pullback}).
Similarly, the object $Lh^*M$ is represented by
the complex $h^*\mathcal{F}^\bullet$.
Finally, the object $Lh^*R\SheafHom(K, M)$
is represented by $h^*\mathcal{H}^\bullet$ as
$\mathcal{H}^\bullet$ is K-flat by
Lemma \ref{lemma-Rhom-strictly-perfect-K-flat}.
Thus to finish the proof it suffices to show that
$h^*\mathcal{H}^\bullet =
\SheafHom^\bullet(h^*\mathcal{E}^\bullet, h^*\mathcal{F}^\bullet)$.
For this it suffices to note that
$h^*\SheafHom(\mathcal{E}, \mathcal{F}) =
\SheafHom(h^*\mathcal{E}, \mathcal{F})$
whenever $\mathcal{E}$ is a direct summand of a finite free
$\mathcal{O}_X$-module.

\medskip\noindent
Proof of (2). Since $h$ is flat, we can compute $Lh^*$
by simply using $h^*$ on any complex of $\mathcal{O}_Y$-modules.
In particular we have $H^i(Lh^*K) = h^*H^i(K)$ for all $i \in \mathbf{Z}$.
Say $H^i(M) = 0$ for $i < a$. Let $K' \to K$ be a morphism
of $D(\mathcal{O}_Y)$ which defines an isomorphism
$H^i(K') \to H^i(K)$ for all $i \geq b$. Then the corresponding maps
$$
R\SheafHom(K, M) \to R\SheafHom(K', M)
$$
and
$$
R\SheafHom(Lh^*K, Lh^*M) \to R\SheafHom(Lh^*K', Lh^*M)
$$
are isomorphisms on cohomology sheaves in degrees $< a - b$ (details omitted).
Thus to prove the map in the statement of the lemma induces an
isomorphism on cohomology sheaves in degrees $< a - b$ it suffices to
prove the result for $K'$ in those degrees. Also, as in the proof
of part (1) the question is local on $Y$. Thus we may assume $K$ is
represented by a strictly perfect complex, see
Section \ref{section-pseudo-coherent}. This reduces us to case (1).

\medskip\noindent
Proof of (3). The proof is the same as the proof of (2) except one
uses that $Lh^*$ has bounded cohomological dimension to get the
desired vanishing. We omit the details.
\end{proof}

\begin{lemma}
\label{lemma-stalk-internal-hom}
Let $X$ be a ringed space. Let $K, M$ be objects of $D(\mathcal{O}_X)$.
Let $x \in X$. The canonical map
$$
R\SheafHom(K, M)_x \longrightarrow
R\Hom_{\mathcal{O}_{X, x}}(K_x, M_x)
$$
is an isomorphism in the following cases
\begin{enumerate}
\item $K$ is perfect,
\item $K$ is pseudo-coherent and $M$ is (locally) bounded below.
\end{enumerate}
\end{lemma}

\begin{proof}
Let $Y = \{x\}$ be the singleton ringed space with structure sheaf
given by $\mathcal{O}_{X, x}$. Then apply
Lemma \ref{lemma-pullback-internal-hom}
to the flat inclusion morphism $Y \to X$.
\end{proof}

\begin{lemma}
\label{lemma-stalk-ext-finite-resolution}
Let $X$ be a ringed space. Let $\mathcal{F}$ and $\mathcal{G}$ be
$\mathcal{O}_X$-modules, let $x \in X$, and let $n \geq 0$. Suppose that
on an open neighbourhood $U$ of $x$ there is an exact sequence
$$
\mathcal{F}_{n + 1} \longrightarrow \mathcal{F}_n \longrightarrow
\ldots \longrightarrow \mathcal{F}_0 \longrightarrow
\mathcal{F}|_U \longrightarrow 0
$$
where $\mathcal{F}_0, \ldots, \mathcal{F}_{n + 1}$ are finite free
$\mathcal{O}_U$-modules. Then the canonical maps
$$
\SheafExt^q_{\mathcal{O}_X}(\mathcal{F}, \mathcal{G})_x
\longrightarrow
\Ext^q_{\mathcal{O}_{X, x}}(\mathcal{F}_x, \mathcal{G}_x)
$$
are isomorphisms for $0 \leq q \leq n$.
\end{lemma}

\begin{proof}
We may replace $X$ by $U$. We argue by induction on $n$. If $n = 0$,
then $\mathcal{F}$ is of finite presentation and the assertion is
Modules, Lemma \ref{modules-lemma-stalk-internal-hom}.

Suppose $n > 0$ and set
$\mathcal{K} = \Ker(\mathcal{F}_0 \to \mathcal{F})$. The sequence
$$
0 \longrightarrow \mathcal{K} \longrightarrow \mathcal{F}_0
\longrightarrow \mathcal{F} \longrightarrow 0
$$
and its sequence of stalks give compatible long exact sequences for
internal Ext and module Ext. The truncated sequence
$$
\mathcal{F}_{n + 1} \longrightarrow \ldots \longrightarrow
\mathcal{F}_1 \longrightarrow \mathcal{K} \longrightarrow 0
$$
satisfies the induction hypothesis. Moreover, internal Ext and module Ext
of the finite free module $\mathcal{F}_0$ vanish in positive degrees, and
the comparison map in degree zero is an isomorphism. The comparison maps
for $\mathcal{F}$ are therefore isomorphisms in degrees $0, \ldots, n$
by the long exact sequences and the induction hypothesis.
\end{proof}

\begin{remark}
\label{remark-stalk-ext-finite-presentation}
\begin{reference}
Compare with \cite[Expos\'e 149, Proposition 2, pp. 170--171]{FGA}.
\end{reference}
For $n = 0$ the hypothesis of
Lemma \ref{lemma-stalk-ext-finite-resolution} is finite presentation, and
the conclusion concerns degree zero only. Comparison in every degree
requires finite free partial resolutions of arbitrary length; for example,
it follows when $\mathcal{F}[0]$ is pseudo-coherent by
Lemma \ref{lemma-stalk-internal-hom}. Thus a single finite presentation does
not supply the syzygy hypotheses used for all higher Ext groups. On a
Noetherian scheme a coherent module is pseudo-coherent by Derived Categories
of Schemes, Lemma
\ref{perfect-lemma-identify-pseudo-coherent-noetherian}, so the coherent
algebraic case of the cited result is covered.
\end{remark}







\section{Invertible objects in the derived category}
\label{section-invertible-D-or-R}

\noindent
We characterize invertible objects in the derived category of
a ringed space (both in the case where the stalks of the
structure sheaf are local and where not).

\begin{lemma}
\label{lemma-category-summands-finite-free}
Let $(X, \mathcal{O}_X)$ be a ringed space.
Set $R = \Gamma(X, \mathcal{O}_X)$. The category of
$\mathcal{O}_X$-modules which are summands of finite free
$\mathcal{O}_X$-modules is equivalent to the category of
finite projective $R$-modules.
\end{lemma}

\begin{proof}
Observe that a finite projective $R$-module is the same thing
as a summand of a finite free $R$-module.
The equivalence is given by the functor $\mathcal{E} \mapsto
\Gamma(X, \mathcal{E})$. The inverse functor is given by the construction of
Modules, Lemma \ref{modules-lemma-construct-quasi-coherent-sheaves}.
\end{proof}

\begin{lemma}
\label{lemma-invertible-derived}
Let $(X, \mathcal{O}_X)$ be a ringed space. Let $M$ be an object
of $D(\mathcal{O}_X)$. The following are equivalent
\begin{enumerate}
\item $M$ is invertible in $D(\mathcal{O}_X)$, see
Categories, Definition \ref{categories-definition-invertible}, and
\item there is a locally finite direct product decomposition
$$
\mathcal{O}_X = \prod\nolimits_{n \in \mathbf{Z}} \mathcal{O}_n
$$
and for each $n$ there is an invertible $\mathcal{O}_n$-module
$\mathcal{H}^n$ (Modules, Definition \ref{modules-definition-invertible})
and $M = \bigoplus \mathcal{H}^n[-n]$ in $D(\mathcal{O}_X)$.
\end{enumerate}
If (1) and (2) hold, then $M$ is a perfect object of $D(\mathcal{O}_X)$. If
$\mathcal{O}_{X, x}$ is a local ring for all $x \in X$ these condition
are also equivalent to
\begin{enumerate}
\item[(3)] there exists an open covering $X = \bigcup U_i$
and for each $i$ an integer $n_i$ such that $M|_{U_i}$
is represented by an invertible $\mathcal{O}_{U_i}$-module
placed in degree $n_i$.
\end{enumerate}
\end{lemma}

\begin{proof}
Assume (2). Consider the object $R\SheafHom(M, \mathcal{O}_X)$
and the composition map
$$
R\SheafHom(M, \mathcal{O}_X) \otimes_{\mathcal{O}_X}^\mathbf{L} M \to
\mathcal{O}_X
$$
To prove this is an isomorphism, we may work locally. Thus we may
assume $\mathcal{O}_X = \prod_{a \leq n \leq b} \mathcal{O}_n$
and $M = \bigoplus_{a \leq n \leq b} \mathcal{H}^n[-n]$.
Then it suffices to show that
$$
R\SheafHom(\mathcal{H}^m, \mathcal{O}_X)
\otimes_{\mathcal{O}_X}^\mathbf{L} \mathcal{H}^n
$$
is zero if $n \not = m$ and equal to $\mathcal{O}_n$ if $n = m$.
The case $n \not = m$ follows from the fact that $\mathcal{O}_n$ and
$\mathcal{O}_m$ are flat $\mathcal{O}_X$-algebras with
$\mathcal{O}_n \otimes_{\mathcal{O}_X} \mathcal{O}_m = 0$.
Using the local structure of invertible $\mathcal{O}_X$-modules
(Modules, Lemma \ref{modules-lemma-invertible}) and working locally
the isomorphism in case $n = m$ follows in a straightforward manner;
we omit the details. Because $D(\mathcal{O}_X)$ is symmetric monoidal,
we conclude that $M$ is invertible.

\medskip\noindent
Assume (1). The description in (2) shows that we have a candidate
for $\mathcal{O}_n$, namely,
$\SheafHom_{\mathcal{O}_X}(H^n(M), H^n(M))$.
If this is a locally finite family of sheaves of rings
and if $\mathcal{O}_X = \prod \mathcal{O}_n$, then we immediately
obtain the direct sum decomposition $M = \bigoplus H^n(M)[-n]$
using the idempotents in $\mathcal{O}_X$ coming from the product
decomposition.
This shows that in order to prove (2) we may work locally on $X$.

\medskip\noindent
Choose an object $N$ of $D(\mathcal{O}_X)$
and an isomorphism
$M \otimes_{\mathcal{O}_X}^\mathbf{L} N \cong \mathcal{O}_X$.
Let $x \in X$.
Then $N$ is a left dual for $M$ in the monoidal category
$D(\mathcal{O}_X)$ and we conclude that $M$ is perfect by
Lemma \ref{lemma-left-dual-derived}. By symmetry we see that
$N$ is perfect. After replacing $X$ by an open neighbourhood of $x$,
we may assume $M$ and $N$ are represented by a strictly perfect
complexes $\mathcal{E}^\bullet$ and $\mathcal{F}^\bullet$.
Then $M \otimes_{\mathcal{O}_X}^\mathbf{L} N$ is represented by
$\text{Tot}(\mathcal{E}^\bullet \otimes_{\mathcal{O}_X} \mathcal{F}^\bullet)$.
After another shinking of $X$ we may assume the mutually inverse
isomorphisms
$\mathcal{O}_X \to M \otimes_{\mathcal{O}_X}^\mathbf{L} N$ and
$M \otimes_{\mathcal{O}_X}^\mathbf{L} N \to \mathcal{O}_X$
are given by maps of complexes
$$
\alpha : \mathcal{O}_X \to
\text{Tot}(\mathcal{E}^\bullet \otimes_{\mathcal{O}_X} \mathcal{F}^\bullet)
\quad\text{and}\quad
\beta :
\text{Tot}(\mathcal{E}^\bullet \otimes_{\mathcal{O}_X} \mathcal{F}^\bullet)
\to \mathcal{O}_X
$$
See Lemma \ref{lemma-local-actual}. Then $\beta \circ \alpha = 1$
as maps of complexes and $\alpha \circ \beta = 1$ as a morphism
in $D(\mathcal{O}_X)$. After shrinking $X$
we may assume the composition $\alpha \circ \beta$ is homotopic to $1$
by some homotopy $\theta$ with components
$$
\theta^n :
\text{Tot}^n(\mathcal{E}^\bullet \otimes_{\mathcal{O}_X} \mathcal{F}^\bullet)
\to
\text{Tot}^{n - 1}(
\mathcal{E}^\bullet \otimes_{\mathcal{O}_X} \mathcal{F}^\bullet)
$$
by the same lemma as before. Set $R = \Gamma(X, \mathcal{O}_X)$. By
Lemma \ref{lemma-category-summands-finite-free}
we find that we obtain
\begin{enumerate}
\item $M^\bullet = \Gamma(X, \mathcal{E}^\bullet)$ is a bounded complex
of finite projective $R$-modules,
\item $N^\bullet = \Gamma(X, \mathcal{F}^\bullet)$ is a bounded complex
of finite projective $R$-modules,
\item $\alpha$ and $\beta$ correspond to maps of complexes
$a : R \to \text{Tot}(M^\bullet \otimes_R N^\bullet)$ and
$b : \text{Tot}(M^\bullet \otimes_R N^\bullet) \to R$,
\item $\theta^n$ corresponds to a map
$h^n : \text{Tot}^n(M^\bullet \otimes_R N^\bullet) \to
\text{Tot}^{n - 1}(M^\bullet \otimes_R N^\bullet)$, and
\item $b \circ a = 1$ and $b \circ a - 1 = dh + hd$,
\end{enumerate}
It follows that $M^\bullet$ and $N^\bullet$ define
mutually inverse objects of $D(R)$. By
More on Algebra, Lemma \ref{more-algebra-lemma-invertible-derived}
we find a product decomposition $R = \prod_{a \leq n \leq b} R_n$
and invertible $R_n$-modules $H^n$ such
that $M^\bullet \cong \bigoplus_{a \leq n \leq b} H^n[-n]$.
This isomorphism in $D(R)$ can be lifted to an morphism
$$
\bigoplus H^n[-n] \longrightarrow M^\bullet
$$
of complexes because each $H^n$ is projective as an $R$-module.
Correspondingly, using Lemma \ref{lemma-category-summands-finite-free} again,
we obtain an morphism
$$
\bigoplus H^n \otimes_R \mathcal{O}_X[-n] \to \mathcal{E}^\bullet
$$
which is an isomorphism in $D(\mathcal{O}_X)$. Setting
$\mathcal{O}_n = R_n \otimes_R \mathcal{O}_X$ we conclude (2) is true.

\medskip\noindent
If all stalks of $\mathcal{O}_X$ are local, then it is straightforward
to prove the equivalence of (2) and (3). We omit the details.
\end{proof}






\section{Compact objects}
\label{section-compact}

\noindent
In this section we study compact objects in the derived category of modules
on a ringed space. We recall that compact objects are defined in
Derived Categories, Definition \ref{derived-definition-compact-object}.
On suitable ringed spaces the perfect objects are compact.

\begin{lemma}
\label{lemma-when-jshriek-compact}
Let $X$ be a ringed space. Let $j : U \to X$ be the
inclusion of an open. The $\mathcal{O}_X$-module $j_!\mathcal{O}_U$ is a
compact object of $D(\mathcal{O}_X)$ if there exists an integer $d$ such that
\begin{enumerate}
\item $H^p(U, \mathcal{F}) = 0$ for all $p > d$, and
\item the functors $\mathcal{F} \mapsto H^p(U, \mathcal{F})$
commute with direct sums.
\end{enumerate}
\end{lemma}

\begin{proof}
Assume (1) and (2). Since
$\Hom(j_!\mathcal{O}_U, \mathcal{F}) = \mathcal{F}(U)$
by Sheaves, Lemma \ref{sheaves-lemma-j-shriek-modules}
we have $\Hom(j_!\mathcal{O}_U, K) = R\Gamma(U, K)$ for
$K$ in $D(\mathcal{O}_X)$. Thus we have to show that $R\Gamma(U, -)$
commutes with direct sums. The first assumption means that the functor
$F = H^0(U, -)$ has finite cohomological dimension. Moreover, the second
assumption implies any direct sum of injective modules is acyclic for $F$.
Let $K_i$ be a family of objects of $D(\mathcal{O}_X)$.
Choose K-injective representatives $I_i^\bullet$ with injective terms
representing $K_i$, see Injectives, Theorem
\ref{injectives-theorem-K-injective-embedding-grothendieck}.
Since we may compute $RF$ by applying $F$ to any complex of acyclics
(Derived Categories, Lemma \ref{derived-lemma-unbounded-right-derived})
and since $\bigoplus K_i$ is represented by $\bigoplus I_i^\bullet$
(Injectives, Lemma \ref{injectives-lemma-derived-products})
we conclude that $R\Gamma(U, \bigoplus K_i)$ is represented by
$\bigoplus H^0(U, I_i^\bullet)$. Hence $R\Gamma(U, -)$ commutes
with direct sums as desired.
\end{proof}

\begin{lemma}
\label{lemma-perfect-is-compact}
Let $X$ be a ringed space. Assume that the underlying topological space
of $X$ has the following properties:
\begin{enumerate}
\item $X$ is quasi-compact,
\item there exists a basis of quasi-compact open subsets, and
\item the intersection of any two quasi-compact opens is quasi-compact.
\end{enumerate}
Let $K$ be a perfect object of $D(\mathcal{O}_X)$. Then
\begin{enumerate}
\item[(a)] $K$ is a compact object of $D^+(\mathcal{O}_X)$
in the following sense: if $M = \bigoplus_{i \in I} M_i$ is
bounded below, then $\Hom(K, M) = \bigoplus_{i \in I} \Hom(K, M_i)$.
\item[(b)] If $X$ has finite cohomological dimension, i.e., if there exists
a $d$ such that $H^i(X, \mathcal{F}) = 0$ for $i > d$, then
$K$ is a compact object of $D(\mathcal{O}_X)$.
\end{enumerate}
\end{lemma}

\begin{proof}
Let $K^\vee$ be the dual of $K$, see
Lemma \ref{lemma-dual-perfect-complex}. Then we have
$$
\Hom_{D(\mathcal{O}_X)}(K, M) =
H^0(X, K^\vee \otimes_{\mathcal{O}_X}^\mathbf{L} M)
$$
functorially in $M$ in $D(\mathcal{O}_X)$.
Since $K^\vee \otimes_{\mathcal{O}_X}^\mathbf{L} -$ commutes with
direct sums it suffices
to show that $R\Gamma(X, -)$ commutes with the relevant direct sums.

\medskip\noindent
Proof of (b). Since $R\Gamma(X, K) = R\Hom(\mathcal{O}_X, K)$
and since $H^p(X, -)$ commutes with direct sums by
Lemma \ref{lemma-quasi-separated-cohomology-colimit}
this is a special case of
Lemma \ref{lemma-when-jshriek-compact}

\medskip\noindent
Proof of (a). Let $\mathcal{I}_i$, $i \in I$ be a collection of injective
$\mathcal{O}_X$-modules. By Lemma \ref{lemma-quasi-separated-cohomology-colimit}
we see that
$$
H^p(X, \bigoplus\nolimits_{i \in I} \mathcal{I}_i) =
\bigoplus\nolimits_{i \in I} H^p(X, \mathcal{I}_i) = 0
$$
for all $p$. Now if $M = \bigoplus M_i$ is as in (a), then we
see that there exists an $a \in \mathbf{Z}$ such that $H^n(M_i) = 0$
for $n < a$. Thus we can choose complexes of injective $\mathcal{O}_X$-modules
$\mathcal{I}_i^\bullet$ representing $M_i$
with $\mathcal{I}_i^n = 0$ for $n < a$, see
Derived Categories, Lemma \ref{derived-lemma-injective-resolutions-exist}.
By Injectives, Lemma \ref{injectives-lemma-derived-products}
we see that the direct sum complex $\bigoplus \mathcal{I}_i^\bullet$
represents $M$. By Leray acyclicity
(Derived Categories, Lemma \ref{derived-lemma-leray-acyclicity})
we see that
$$
R\Gamma(X, M) = \Gamma(X, \bigoplus \mathcal{I}_i^\bullet) =
\bigoplus \Gamma(X, \bigoplus \mathcal{I}_i^\bullet) =
\bigoplus R\Gamma(X, M_i)
$$
as desired.
\end{proof}


















\section{Projection formula}
\label{section-projection-formula}

\noindent
In this section we collect variants of the projection formula.
The most basic version is Lemma \ref{lemma-projection-formula}.
After we state and prove it, we discuss a more general version
involving perfect complexes.

\begin{lemma}
\label{lemma-injective-tensor-finite-locally-free}
Let $X$ be a ringed space.
Let $\mathcal{I}$ be an injective $\mathcal{O}_X$-module.
Let $\mathcal{E}$ be an $\mathcal{O}_X$-module.
Assume $\mathcal{E}$ is finite locally free on $X$, see
Modules, Definition \ref{modules-definition-locally-free}.
Then $\mathcal{E} \otimes_{\mathcal{O}_X} \mathcal{I}$ is
an injective $\mathcal{O}_X$-module.
\end{lemma}

\begin{proof}
This is true because under the assumptions of the lemma we have
$$
\Hom_{\mathcal{O}_X}(\mathcal{F},
\mathcal{E} \otimes_{\mathcal{O}_X} \mathcal{I})
=
\Hom_{\mathcal{O}_X}(
\mathcal{F} \otimes_{\mathcal{O}_X} \mathcal{E}^\vee, \mathcal{I})
$$
where
$\mathcal{E}^\vee = \SheafHom_{\mathcal{O}_X}(\mathcal{E}, \mathcal{O}_X)$
is the dual of $\mathcal{E}$ which is finite locally free also. Since tensoring
with a finite locally free sheaf is an exact functor we win by
Homology, Lemma \ref{homology-lemma-characterize-injectives}.
\end{proof}

\begin{lemma}
\label{lemma-projection-formula}
Let $f : X \to Y$ be a morphism of ringed spaces.
Let $\mathcal{F}$ be an $\mathcal{O}_X$-module.
Let $\mathcal{E}$ be an $\mathcal{O}_Y$-module.
Assume $\mathcal{E}$ is finite locally free on $Y$, see
Modules, Definition \ref{modules-definition-locally-free}.
Then there exist isomorphisms
$$
\mathcal{E} \otimes_{\mathcal{O}_Y} R^qf_*\mathcal{F}
\longrightarrow
R^qf_*(f^*\mathcal{E} \otimes_{\mathcal{O}_X} \mathcal{F})
$$
for all $q \geq 0$. In fact there exists an isomorphism
$$
\mathcal{E} \otimes_{\mathcal{O}_Y} Rf_*\mathcal{F}
\longrightarrow
Rf_*(f^*\mathcal{E} \otimes_{\mathcal{O}_X} \mathcal{F})
$$
in $D^{+}(Y)$ functorial in $\mathcal{F}$.
\end{lemma}

\begin{proof}
Choose an injective resolution $\mathcal{F} \to \mathcal{I}^\bullet$
on $X$. Note that $f^*\mathcal{E}$ is finite locally free also, hence
we get a resolution
$$
f^*\mathcal{E} \otimes_{\mathcal{O}_X} \mathcal{F}
\longrightarrow
f^*\mathcal{E} \otimes_{\mathcal{O}_X} \mathcal{I}^\bullet
$$
which is an injective resolution by
Lemma \ref{lemma-injective-tensor-finite-locally-free}.
Apply $f_*$ to see that
$$
Rf_*(f^*\mathcal{E} \otimes_{\mathcal{O}_X} \mathcal{F})
=
f_*(f^*\mathcal{E} \otimes_{\mathcal{O}_X} \mathcal{I}^\bullet).
$$
Hence the lemma follows if we can show that
$f_*(f^*\mathcal{E} \otimes_{\mathcal{O}_X} \mathcal{F}) =
\mathcal{E} \otimes_{\mathcal{O}_Y} f_*(\mathcal{F})$ functorially
in the $\mathcal{O}_X$-module $\mathcal{F}$. This is clear when
$\mathcal{E} = \mathcal{O}_Y^{\oplus n}$, and follows in general
by working locally on $Y$. Details omitted.
\end{proof}

\noindent
Let $f : X \to Y$ be a morphism of ringed spaces.
Let $E \in D(\mathcal{O}_X)$ and $K \in D(\mathcal{O}_Y)$.
Without any further assumptions there is a map
\begin{equation}
\label{equation-projection-formula-map}
Rf_*E \otimes^\mathbf{L}_{\mathcal{O}_Y} K
\longrightarrow
Rf_*(E \otimes^\mathbf{L}_{\mathcal{O}_X} Lf^*K)
\end{equation}
Namely, it is the adjoint to the canonical map
$$
Lf^*(Rf_*E \otimes^\mathbf{L}_{\mathcal{O}_Y} K) =
Lf^*Rf_*E \otimes^\mathbf{L}_{\mathcal{O}_X} Lf^*K
\longrightarrow
E \otimes^\mathbf{L}_{\mathcal{O}_X} Lf^*K
$$
coming from the map $Lf^*Rf_*E \to E$ and Lemmas
\ref{lemma-pullback-tensor-product} and \ref{lemma-adjoint}.
A reasonably general version of the projection formula is the following.

\begin{lemma}
\label{lemma-projection-formula-perfect}
Let $f : X \to Y$ be a morphism of ringed spaces.
Let $E \in D(\mathcal{O}_X)$ and $K \in D(\mathcal{O}_Y)$.
If $K$ is perfect, then
$$
Rf_*E \otimes^\mathbf{L}_{\mathcal{O}_Y} K =
Rf_*(E \otimes^\mathbf{L}_{\mathcal{O}_X} Lf^*K)
$$
in $D(\mathcal{O}_Y)$.
\end{lemma}

\begin{proof}
To check (\ref{equation-projection-formula-map}) is an isomorphism
we may work locally on $Y$, i.e., we have to find an open covering
$\{V_j \to Y\}$
such that the map restricts to an isomorphism on $V_j$. By definition
of perfect objects, this means we may assume $K$ is represented by
a strictly perfect complex of $\mathcal{O}_Y$-modules.
Note that, completely generally, the statement is true for
$K = K_1 \oplus K_2$, if and only if the statement is true for
$K_1$ and $K_2$. Hence we may assume $K$ is a finite
complex of finite free $\mathcal{O}_Y$-modules.
In this case a simple argument involving stupid truncations reduces
the statement to the case where $K$ is represented by a finite
free $\mathcal{O}_Y$-module. Since the statement is invariant
under finite direct summands in the $K$ variable, we conclude
it suffices to prove it for $K = \mathcal{O}_Y[n]$
in which case it is trivial.
\end{proof}

\noindent
Here is a case where the projection formula is true in complete
generality.

\begin{lemma}
\label{lemma-projection-formula-closed-immersion}
Let $f : X \to Y$ be a morphism of ringed spaces such that $f$ is a
homeomorphism onto a closed subset. Then
(\ref{equation-projection-formula-map}) is an isomorphism always.
\end{lemma}

\begin{proof}
Since $f$ is a homeomorphism onto a closed subset, the functor $f_*$
is exact (Modules, Lemma \ref{modules-lemma-i-star-exact}). Hence
$Rf_*$ is computed by applying $f_*$ to any representative complex.
Choose a K-flat complex $\mathcal{K}^\bullet$ of $\mathcal{O}_Y$-modules
representing $K$ and choose any complex $\mathcal{E}^\bullet$
of $\mathcal{O}_X$-modules representing $E$. Then
$Lf^*K$ is represented by $f^*\mathcal{K}^\bullet$ which is
a K-flat complex of $\mathcal{O}_X$-modules
(Lemma \ref{lemma-pullback-K-flat}). Thus the right hand side of
(\ref{equation-projection-formula-map}) is represented by
$$
f_*\text{Tot}(\mathcal{E}^\bullet
\otimes_{\mathcal{O}_X} f^*\mathcal{K}^\bullet)
$$
By the same reasoning we see that the left hand side is represented by
$$
\text{Tot}(f_*\mathcal{E}^\bullet \otimes_{\mathcal{O}_Y} \mathcal{K}^\bullet)
$$
Since $f_*$ commutes with direct sums
(Modules, Lemma \ref{modules-lemma-i-star-right-adjoint})
it suffices to show that
$$
f_*(\mathcal{E} \otimes_{\mathcal{O}_X} f^*\mathcal{K}) =
f_*\mathcal{E} \otimes_{\mathcal{O}_Y} \mathcal{K}
$$
for any $\mathcal{O}_X$-module $\mathcal{E}$ and $\mathcal{O}_Y$-module
$\mathcal{K}$. We will check this by checking on stalks.
Let $y \in Y$. If $y \not \in f(X)$, then the stalks
of both sides are zero. If $y = f(x)$, then we see that we have to show
$$
\mathcal{E}_x \otimes_{\mathcal{O}_{X, x}}
(\mathcal{O}_{X, x} \otimes_{\mathcal{O}_{Y, y}} \mathcal{F}_y) =
\mathcal{E}_x \otimes_{\mathcal{O}_{Y, y}} \mathcal{F}_y
$$
(using Sheaves, Lemma \ref{sheaves-lemma-stalks-closed-pushforward}
and Lemma \ref{sheaves-lemma-stalk-pullback-modules}).
This equality holds and therefore the lemma has been proved.
\end{proof}

\begin{remark}
\label{remark-compatible-with-diagram}
The map (\ref{equation-projection-formula-map}) is compatible with the
base change map of Remark \ref{remark-base-change} in the following sense.
Namely, suppose that
$$
\xymatrix{
X' \ar[r]_{g'} \ar[d]_{f'} &
X \ar[d]^f \\
Y' \ar[r]^g &
Y
}
$$
is a commutative diagram of ringed spaces.
Let $E \in D(\mathcal{O}_X)$ and $K \in D(\mathcal{O}_Y)$.
Then the diagram
$$
\xymatrix{
Lg^*(Rf_*E \otimes^\mathbf{L}_{\mathcal{O}_Y} K) \ar[r]_p \ar[d]_t &
Lg^*Rf_*(E \otimes^\mathbf{L}_{\mathcal{O}_X} Lf^*K) \ar[d]_b \\
Lg^*Rf_*E \otimes^\mathbf{L}_{\mathcal{O}_{Y'}} Lg^*K \ar[d]_b &
Rf'_*L(g')^*(E \otimes^\mathbf{L}_{\mathcal{O}_X} Lf^*K) \ar[d]_t \\
Rf'_*L(g')^*E \otimes^\mathbf{L}_{\mathcal{O}_{Y'}} Lg^*K \ar[rd]_p &
Rf'_*(L(g')^*E \otimes^\mathbf{L}_{\mathcal{O}_{Y'}} L(g')^*Lf^*K) \ar[d]_c \\
& Rf'_*(L(g')^*E \otimes^\mathbf{L}_{\mathcal{O}_{Y'}} L(f')^*Lg^*K)
}
$$
is commutative. Here arrows labeled $t$ are gotten by an application of
Lemma \ref{lemma-pullback-tensor-product}, arrows labeled $b$ by an
application of Remark \ref{remark-base-change}, arrows labeled $p$
by an application of (\ref{equation-projection-formula-map}), and
$c$ comes from $L(g')^* \circ Lf^* = L(f')^* \circ Lg^*$.
We omit the verification.
\end{remark}















\section{An operator introduced by Berthelot and Ogus}
\label{section-eta}

\noindent
This section continuous the discussion started in
More on Algebra, Section \ref{more-algebra-section-eta}.
We strongly encourage the reader to read that section first.

\begin{lemma}
\label{lemma-invertible-ideal-sheaf}
Let $(X, \mathcal{O}_X)$ be a ringed space. Let
$\mathcal{I} \subset \mathcal{O}_X$ be a sheaf of ideals.
Consider the following two conditions
\begin{enumerate}
\item for every $x \in X$ there exists an open neighbourhood
$U \subset X$ of $x$ and $f \in \mathcal{I}(U)$ such that
$\mathcal{I}|_U = \mathcal{O}_U \cdot f$ and
$f : \mathcal{O}_U \to \mathcal{O}_U$ is injective, and
\item $\mathcal{I}$ is invertible as an $\mathcal{O}_X$-module.
\end{enumerate}
Then (1) implies (2) and the converse is true if all stalks
$\mathcal{O}_{X, x}$ of the structure sheaf are local rings.
\end{lemma}

\begin{proof}
Omitted. Hint: Use Modules, Lemma
\ref{modules-lemma-invertible-is-locally-free-rank-1}.
\end{proof}

\begin{situation}
\label{situation-eta}
Let $(X, \mathcal{O}_X)$ be a ringed space. Let
$\mathcal{I} \subset \mathcal{O}_X$ be a sheaf of ideals
satisfying condition (1) of
Lemma \ref{lemma-invertible-ideal-sheaf}\footnote{The discussion
in this section can be generalized to the case where all we
require is that $\mathcal{I}$ is an invertible $\mathcal{O}_X$-module
as defined in Modules, Section \ref{modules-section-invertible}.}.
\end{situation}

\begin{lemma}
\label{lemma-I-torsion-free}
In Situation \ref{situation-eta}
let $\mathcal{F}$ be an $\mathcal{O}_X$-module.
The following are equivalent
\begin{enumerate}
\item the subsheaf $\mathcal{F}[\mathcal{I}] \subset \mathcal{F}$
of sections annihilated by $\mathcal{I}$ is zero,
\item the subsheaf $\mathcal{F}[\mathcal{I}^n]$ is zero for all $n \geq 1$,
\item the multiplication map
$\mathcal{I} \otimes_{\mathcal{O}_X} \mathcal{F} \to \mathcal{F}$
is injective,
\item for every open $U \subset X$ such that
$\mathcal{I}|_U = \mathcal{O}_U \cdot f$
for some $f \in \mathcal{I}(U)$
the map $f : \mathcal{F}|_U \to \mathcal{F}|_U$ is injective,
\item for every $x \in X$ and generator $f$ of the ideal
$\mathcal{I}_x \subset \mathcal{O}_{X, x}$ the element $f$
is a nonzerodivisor on the stalk $\mathcal{F}_x$.
\end{enumerate}
\end{lemma}

\begin{proof}
Omitted.
\end{proof}

\noindent
In Situation \ref{situation-eta}
let $\mathcal{F}$ be an $\mathcal{O}_X$-module.
If the equivalent conditions of Lemma \ref{lemma-I-torsion-free} hold,
then we will say that $\mathcal{F}$ is {\it $\mathcal{I}$-torsion free}.
If so, then for any $i \in \mathbf{Z}$ we will denote
$$
\mathcal{I}^i\mathcal{F} =
\mathcal{I}^{\otimes i} \otimes_{\mathcal{O}_X} \mathcal{F}
$$
so that we have inclusions
$$
\ldots \subset
\mathcal{I}^{i + 1}\mathcal{F} \subset
\mathcal{I}^i\mathcal{F} \subset
\mathcal{I}^{i - 1}\mathcal{F} \subset \ldots
$$
The modules $\mathcal{I}^i\mathcal{F}$ are locally isomorphic to
$\mathcal{F}$ as $\mathcal{O}_X$-modules, but not globally.

\medskip\noindent
Let $\mathcal{F}^\bullet$ be a complex of $\mathcal{I}$-torsion free
$\mathcal{O}_X$-modules with
differentials $d^i : \mathcal{F}^i \to \mathcal{F}^{i + 1}$.
In this case we define $\eta_\mathcal{I}\mathcal{F}^\bullet$
to be the complex with terms
\begin{align*}
(\eta_\mathcal{I}\mathcal{F})^i
& =
\Ker\left(
d^i, -1 :
\mathcal{I}^i\mathcal{F}^i \oplus \mathcal{I}^{i + 1}\mathcal{F}^{i + 1}
\to
\mathcal{I}^i\mathcal{F}^{i + 1}
\right) \\
& =
\Ker\left(d^i :
\mathcal{I}^i\mathcal{F}^i
\to
\mathcal{I}^i\mathcal{F}^{i + 1}/
\mathcal{I}^{i + 1}\mathcal{F}^{i + 1}
\right)
\end{align*}
and differential induced by $d^i$. In other words, a local section
$s$ of $(\eta_\mathcal{I}\mathcal{F})^i$ is the same thing as a local section
$s$ of $\mathcal{I}^i\mathcal{F}^i$ such that its image $d^i(s)$
in $\mathcal{I}^i\mathcal{F}^{i + 1}$ is in the subsheaf
$\mathcal{I}^{i + 1}\mathcal{F}^{i + 1}$.
Observe that $\eta_\mathcal{I}\mathcal{F}^\bullet$
is another complex of $\mathcal{I}$-torsion free modules.

\medskip\noindent
Let $a^\bullet : \mathcal{F}^\bullet \to \mathcal{G}^\bullet$ be a map of
complexes of $\mathcal{I}$-torsion free $\mathcal{O}_X$-modules.
Then we obtain a map of complexes
$$
\eta_\mathcal{I} a^\bullet :
\eta_\mathcal{I}\mathcal{F}^\bullet
\longrightarrow
\eta_\mathcal{I}\mathcal{G}^\bullet
$$
induced by the maps
$\mathcal{I}^i\mathcal{F}^i \to \mathcal{I}^i\mathcal{G}^i$.
The reader checks that we obtain
an endo-functor on the category of complexes of
$\mathcal{I}$-torsion free $\mathcal{O}_X$-modules.

\medskip\noindent
If $a^\bullet, b^\bullet : \mathcal{F}^\bullet \to \mathcal{G}^\bullet$
are two maps of
complexes of $\mathcal{I}$-torsion free $\mathcal{O}_X$-modules
and $h = \{h^i : \mathcal{F}^i \to \mathcal{G}^{i - 1}\}$ is a homotopy
between $a^\bullet$ and $b^\bullet$, then we define
$\eta_\mathcal{I}h$ to be the family of maps
$(\eta_\mathcal{I}h)^i : (\eta_\mathcal{I}\mathcal{F})^i \to
(\eta_\mathcal{I}\mathcal{G})^{i - 1}$
which sends $x$ to $h^i(x)$; this makes sense as
$x$ a local section of $\mathcal{I}^i\mathcal{F}^i$
implies $h^i(x)$ is a local section of $\mathcal{I}^i\mathcal{G}^{i - 1}$
which is certainly contained in $(\eta_\mathcal{I}\mathcal{G})^{i - 1}$.
The reader checks that $\eta_\mathcal{I}h$ is a homotopy
between $\eta_\mathcal{I}a^\bullet$ and $\eta_\mathcal{I}b^\bullet$.
All in all we see that we obtain a functor
$$
\eta_f :
K(\mathcal{I}\text{-torsion free }\mathcal{O}_X\text{-modules})
\longrightarrow
K(\mathcal{I}\text{-torsion free }\mathcal{O}_X\text{-modules})
$$
on the homotopy category
(Derived Categories, Section \ref{derived-section-homotopy})
of the additive category of $\mathcal{I}$-torsion free
$\mathcal{O}_X$-modules.
There is no sense in which $\eta_\mathcal{I}$ is an exact functor of
triangulated categories; compare with More on Algebra,
Example \ref{more-algebra-example-eta-not-distinguished}.

\begin{lemma}
\label{lemma-eta-stalks}
In Situation \ref{situation-eta}
let $\mathcal{F}^\bullet$ be a complex of $\mathcal{I}$-torsion free
$\mathcal{O}_X$-modules.
For $x \in X$ choose a generator $f \in \mathcal{I}_x$. Then
the stalk $(\eta_\mathcal{I}\mathcal{F}^\bullet)_x$ is canonically
isomorphic to the complex $\eta_f\mathcal{F}^\bullet_x$ constructed
in More on Algebra, Section \ref{more-algebra-section-eta}.
\end{lemma}

\begin{proof}
Omitted.
\end{proof}

\begin{lemma}
\label{lemma-eta-first-property}
In Situation \ref{situation-eta}
let $\mathcal{F}^\bullet$ be a complex of $\mathcal{I}$-torsion free
$\mathcal{O}_X$-modules. There is a canonical isomorphism
$$
\mathcal{I}^{\otimes i} \otimes_{\mathcal{O}_X}
\left(
H^i(\mathcal{F}^\bullet)/H^i(\mathcal{F}^\bullet)[\mathcal{I}]
\right)
\longrightarrow H^i(\eta_\mathcal{I}\mathcal{F}^\bullet)
$$
of cohomology sheaves.
\end{lemma}

\begin{proof}
We define a map
$$
\mathcal{I}^{\otimes i} \otimes_{\mathcal{O}_X} H^i(\mathcal{F}^\bullet)
\longrightarrow
H^i(\eta_\mathcal{I}\mathcal{F}^\bullet)
$$
as follows. Let $g$ be a local section of $\mathcal{I}^{\otimes i}$
and let $\overline{s}$ be a local section of $H^i(\mathcal{F}^\bullet)$.
Then $\overline{s}$ is (locally) the class of a local section $s$ of
$\Ker(d^i : \mathcal{F}^i \to \mathcal{F}^{i + 1})$.
Then we send $g \otimes \overline{s}$ to the local section
$gs$ of $(\eta_\mathcal{I}\mathcal{F})^i \subset \mathcal{I}^i\mathcal{F}$.
Of course $gs$ is in the kernel of $d^i$ on
$\eta_\mathcal{I}\mathcal{F}^\bullet$ and hence defines a local
section of $H^i(\eta_\mathcal{I}\mathcal{F}^\bullet)$.
Checking that this is well defined is without problems.
We claim that this map factors through an isomorphism
as given in the lemma. This we my check on stalks and hence
via Lemma \ref{lemma-eta-stalks} this translates into the result of
More on Algebra, Lemma \ref{more-algebra-lemma-eta-first-property}.
\end{proof}

\begin{lemma}
\label{lemma-eta-qis}
In Situation \ref{situation-eta}
let $\mathcal{F}^\bullet \to \mathcal{G}^\bullet$ be a map of
complexes of $\mathcal{I}$-torsion free $\mathcal{O}_X$-modules.
Then the induced map
$\eta_\mathcal{I}\mathcal{F}^\bullet \to \eta_\mathcal{I}\mathcal{G}^\bullet$
is a quasi-isomorphism too.
\end{lemma}

\begin{proof}
This is true because the isomorphisms of Lemma \ref{lemma-eta-first-property}
are compatible with maps of complexes.
\end{proof}

\begin{lemma}
\label{lemma-Leta}
In Situation \ref{situation-eta} there is an additive
functor\footnote{Beware that this functor isn't exact, i.e.,
does not transform distinguished triangles into distinguished triangles.}
$L\eta_\mathcal{I} : D(\mathcal{O}_X) \to D(\mathcal{O}_X)$
such that if $M$ in $D(\mathcal{O}_X)$ is represented by a complex
$\mathcal{F}^\bullet$ of $\mathcal{I}$-torsion free $\mathcal{O}_X$-modules,
then $L\eta_\mathcal{I}M = \eta_\mathcal{I}\mathcal{F}^\bullet$.
Similarly for morphisms.
\end{lemma}

\begin{proof}
Denote $\mathcal{T} \subset \textit{Mod}(\mathcal{O}_X)$ the full subcategory
of $\mathcal{I}$-torsion free $\mathcal{O}_X$-modules.
We have a corresponding inclusion
$$
K(\mathcal{T})
\quad\subset\quad
K(\textit{Mod}(\mathcal{O}_X)) = K(\mathcal{O}_X)
$$
of $K(\mathcal{T})$ as a full triangulated subcategory of $K(\mathcal{O}_X)$.
Let $S \subset \text{Arrows}(K(\mathcal{T}))$ be the quasi-isomorphisms.
We will apply
Derived Categories, Lemma \ref{derived-lemma-localization-subcategory}
to show that the map
$$
S^{-1}K(\mathcal{T}) \longrightarrow D(\mathcal{O}_X)
$$
is an equivalence of triangulated categories. The lemma shows that
it suffices to prove: given a complex $\mathcal{G}^\bullet$ of
$\mathcal{O}_X$-modules,
there exists a quasi-isomorphism $\mathcal{F}^\bullet \to \mathcal{G}^\bullet$
with $\mathcal{F}^\bullet$
a complex of $\mathcal{I}$-torsion free $\mathcal{O}_X$-modules.
By Lemma \ref{lemma-K-flat-resolution} we can find a quasi-isomorphism
$\mathcal{F}^\bullet \to \mathcal{G}^\bullet$ such that the complex
$\mathcal{F}^\bullet$ is K-flat (we won't use this) and
consists of flat $\mathcal{O}_X$-modules $\mathcal{F}^i$.
By the third characterization of Lemma \ref{lemma-I-torsion-free}
we see that a flat $\mathcal{O}_X$-module is an
$\mathcal{I}$-torsion free $\mathcal{O}_X$-module
and we are done.

\medskip\noindent
With these preliminaries out of the way we can define $L\eta_f$.
Namely, by the discussion following Lemma \ref{lemma-I-torsion-free}
this section we have already a well defined functor
$$
K(\mathcal{T}) \xrightarrow{\eta_f} K(\mathcal{T}) \to
K(\mathcal{O}_X) \to D(\mathcal{O}_X)
$$
which according to Lemma \ref{lemma-eta-qis} sends quasi-isomorphisms
to quasi-isomorphisms. Hence this functor factors over
$S^{-1}K(\mathcal{T}) = D(\mathcal{O}_X)$ by
Categories, Lemma \ref{categories-lemma-properties-left-localization}.
\end{proof}

\noindent
In Situation \ref{situation-eta} let us construct the Bockstein operators.
First we observe that there is a commutative diagram
$$
\xymatrix{
0 \ar[r] &
\mathcal{I}^{i + 1} \ar[r] \ar[d] &
\mathcal{I}^i \ar[r] \ar[d] &
\mathcal{I}^i/\mathcal{I}^{i + 1}
\ar[r] \ar@{=}[d] &
0 \\
0 \ar[r] &
\mathcal{I}^{i + 1}/\mathcal{I}^{i + 2} \ar[r] &
\mathcal{I}^i/\mathcal{I}^{i + 2} \ar[r] &
\mathcal{I}^i/\mathcal{I}^{i + 1} \ar[r] &
0
}
$$
whose rows are short exact sequences of $\mathcal{O}_X$-modules.
Let $M$ be an object of $D(\mathcal{O}_X)$.
Tensoring the above diagram with $M$ gives a morphism
$$
\xymatrix{
M \otimes^\mathbf{L} \mathcal{I}^{i + 1} \ar[r] \ar[d] &
M \otimes^\mathbf{L} \mathcal{I}^i \ar[r] \ar[d] &
M \otimes^\mathbf{L} \mathcal{I}^i/\mathcal{I}^{i + 1} \ar[d]^{\text{id}} \\
M \otimes^\mathbf{L} \mathcal{I}^{i + 1}/\mathcal{I}^{i + 2} \ar[r] &
M \otimes^\mathbf{L} \mathcal{I}^i/\mathcal{I}^{i + 2} \ar[r] &
M \otimes^\mathbf{L} \mathcal{I}^i/\mathcal{I}^{i + 1}
}
$$
of distinguished triangles. The long exact sequence of
cohomology sheaves associated
the bottom triangle in particular determines the
{\it Bockstein operator}
$$
\beta = \beta^i :
H^i(M \otimes^\mathbf{L} \mathcal{I}^i/\mathcal{I}^{i + 1})
\longrightarrow
H^{i + 1}(M \otimes^\mathbf{L} \mathcal{I}^{i + 1}/\mathcal{I}^{i + 2})
$$
for all $i \in \mathbf{Z}$. For later use we record here that by
the commutative diagram above there is a factorization
\begin{equation}
\label{equation-factorization-bockstein}
\vcenter{
\xymatrix{
H^i(M \otimes^\mathbf{L} \mathcal{I}^i/\mathcal{I}^{i + 1})
\ar[r]_\delta \ar[rd]_\beta &
H^{i + 1}(M \otimes^\mathbf{L} \mathcal{I}^{i + 1}) \ar[d] \\
&
H^{i + 1}(M \otimes^\mathbf{L} \mathcal{I}^{i + 1}/\mathcal{I}^{i + 2})
}
}
\end{equation}
of the Bockstein operator where $\delta$ is the boundary operator
coming from the top distinguished triangle in the commutative diagram above.
We obtain a complex
\begin{equation}
\label{equation-complex-bocksteins}
H^\bullet(M/\mathcal{I}) =
\left[
\begin{matrix}
\ldots \\
\downarrow \\
H^{i - 1}(M \otimes^\mathbf{L} \mathcal{I}^{i - 1}/\mathcal{I}^i) \\
\downarrow \beta \\
H^i(M \otimes^\mathbf{L} \mathcal{I}^i/\mathcal{I}^{i + 1}) \\
\downarrow \beta \\
H^{i + 1}(M \otimes^\mathbf{L} \mathcal{I}^{i + 1}/\mathcal{I}^{i + 2}) \\
\downarrow \\
\ldots
\end{matrix}
\right]
\end{equation}
i.e., that $\beta \circ \beta = 0$. Namely, we can check this on stalks
and in this case we can deduce it from the corresponding result in algebra
shown in More on Algebra, Section \ref{more-algebra-section-eta}.
Alternative proof: the short exact sequences
$0 \to \mathcal{I}^{i + 1}/\mathcal{I}^{i + 2}
\to \mathcal{I}^i/\mathcal{I}^{i + 2}
\to \mathcal{I}^i/\mathcal{I}^{i + 1} \to 0$
define maps
$b^i : \mathcal{I}^i/\mathcal{I}^{i + 1} \to
(\mathcal{I}^{i + 1}/\mathcal{I}^{i + 2})[1]$
in $D(\mathcal{O}_X)$
which induce the maps $\beta$ above by tensoring with $M$
and taking cohomology sheaves.
Then one shows that the composition
$b^{i + 1}[1] \circ b^i : \mathcal{I}^i/\mathcal{I}^{i + 1} \to
(\mathcal{I}^{i + 1}/\mathcal{I}^{i + 2})[1] \to
(\mathcal{I}^{i + 2}/\mathcal{I}^{i + 3})[2]$
is zero in $D(\mathcal{O}_X)$ by using the criterion in
Derived Categories, Lemma \ref{derived-lemma-cup-ext-1-zero}
using that the module $\mathcal{I}^i/\mathcal{I}^{i + 3}$
is an extension of $\mathcal{I}^{i + 1}/\mathcal{I}^{i + 3}$
by $\mathcal{I}^i/\mathcal{I}^{i + 1}$.

\begin{lemma}
\label{lemma-eta-second-property}
In Situation \ref{situation-eta} let $M$ be an object of
$D(\mathcal{O}_X)$. There is a canonical isomorphism
$$
L\eta_\mathcal{I}M \otimes^\mathbf{L} \mathcal{O}_X/\mathcal{I}
\longrightarrow
H^\bullet(M/\mathcal{I})
$$
in $D(\mathcal{O}_X)$ where the right hand side is the complex
(\ref{equation-complex-bocksteins}).
\end{lemma}

\begin{proof}
By the construction of $L\eta_\mathcal{I}$ in Lemma \ref{lemma-eta-qis}
we may assume $M$ is represented by a complex of $\mathcal{I}$-torsion
free $\mathcal{O}_X$-modules $\mathcal{F}^\bullet$. Then
$L\eta_\mathcal{I}M$ is represented by the complex
$\eta_\mathcal{I}\mathcal{F}^\bullet$ which is a
complex of $\mathcal{I}$-torsion free $\mathcal{O}_X$-modules as well.
Thus $L\eta_\mathcal{I}M \otimes^\mathbf{L} \mathcal{O}_X/\mathcal{I}$
is represented by the complex
$\eta_\mathcal{I}\mathcal{F}^\bullet \otimes \mathcal{O}_X/\mathcal{I}$.
Similarly, the complex $H^\bullet(M/\mathcal{I})$ has terms
$H^i(\mathcal{F}^\bullet \otimes \mathcal{I}^i/\mathcal{I}^{i + 1})$.

\medskip\noindent
Let $f$ be a local generator for $\mathcal{I}$.
Let $s$ be a local section of $(\eta_\mathcal{I}\mathcal{F})^i$.
Then we can write $s = f^is'$ for a local section $s'$ of
$\mathcal{F}^i$ and similarly $d^i(s) = f^{i + 1}t$ for a local
section $t$ of $\mathcal{F}^{i + 1}$. Thus $d^i$ maps $f^is'$
to zero in $\mathcal{F}^{i + 1} \otimes \mathcal{I}^i/\mathcal{I}^{i + 1}$.
Hence we may map $s$ to the class of $f^is'$ in
$H^i(\mathcal{F}^\bullet \otimes \mathcal{I}^i/\mathcal{I}^{i + 1})$.
This rule defines a map
$$
(\eta_\mathcal{I}\mathcal{F})^i \otimes \mathcal{O}_X/\mathcal{I}
\longrightarrow
H^i(\mathcal{F}^\bullet \otimes \mathcal{I}^i/\mathcal{I}^{i + 1})
$$
of $\mathcal{O}_X$-modules. A calculation shows that these maps
are compatible with differentials (essentially because $\beta$
sends the class of $f^is'$ to the class of $f^{i + 1}t$), whence
a map of complexes representing the arrow in the statement of the lemma.

\medskip\noindent
To finish the proof, we observe that the construction given
in the previous paragraph agrees on stalks with the maps
constructed in More on Algebra, Lemma
\ref{more-algebra-lemma-eta-second-property}
hence we conclude.
\end{proof}

\begin{lemma}
\label{lemma-eta-tensor-invertible}
In Situation \ref{situation-eta}
let $\mathcal{F}^\bullet$ be a complex of
$\mathcal{I}$-torsion free $\mathcal{O}_X$-modules.
Let $\mathcal{L}$ be an invertible $\mathcal{O}_X$-module.
Then $\eta_\mathcal{I}(\mathcal{F}^\bullet \otimes \mathcal{L}) =
(\eta_\mathcal{I}\mathcal{F}^\bullet) \otimes \mathcal{L}$.
\end{lemma}

\begin{proof}
Immediate from the construction.
\end{proof}

\begin{lemma}
\label{lemma-eta-cohomology-locally-free}
In Situation \ref{situation-eta} let $M$ be an object of $D(\mathcal{O}_X)$.
Let $x \in X$ with $\mathcal{O}_{X, x}$ nonzero. If $H^i(M)_x$
is finite free over $\mathcal{O}_{X, x}$, then $H^i(L\eta_\mathcal{I}M)_x$
is finite free over $\mathcal{O}_{X, x}$ of the same rank.
\end{lemma}

\begin{proof}
Namely, say $f \in \mathcal{O}_{X, x}$ generates the stalk $\mathcal{I}_x$.
Then $f$ is a nonzerodivisor in $\mathcal{O}_{X, x}$ and hence
$H^i(M)_x[f] = 0$. Thus by
Lemma \ref{lemma-eta-first-property}
we see that $H^i(L\eta_\mathcal{I}M)_x$
is isomorphic to
$\mathcal{I}^i_x \otimes_{\mathcal{O}_{X, x}} H^i(M)_x$
which is free of the same rank as desired.
\end{proof}









\input{chapters}

\bibliography{my}
\bibliographystyle{amsalpha}

\end{document}
